% hol-kanren.tex — main file; compile with XeLaTeX: tectonic --keep-logs --keep-intermediates hol-kanren.tex, run TWICE.
% A single run from scratch stops before the page breaks settle (the second run reads the first run's .aux and .toc).
\PassOptionsToPackage{dvipsnames}{xcolor}
\documentclass[a4paper,justified,nofonts,nobib,nols]{tufte-book}
% preamble.tex — shared by every chapter of hol-kanren.tex; tufte-book class, compile with XeLaTeX
% (tectonic hol-kanren.tex, or latexmk -xelatex). The class is loaded in hol-kanren.tex as
%   \PassOptionsToPackage{dvipsnames}{xcolor}
%   \documentclass[a4paper,justified,nofonts,nobib,nols]{tufte-book}
% tufte-book already loads fontspec (through xltxtra), xcolor, hyperref, titlesec, changepage and geometry.
%
% Layout: a 107mm text column and a 49.4mm margin column. The margin holds \sidenote, \marginnote and the
% captions of figure/table floats. Code listings and the reading boxes span the full width (text + margin),
% so that code lines keep their length; wide tables and rule figures use fullwidth, figure* or table*.
\usepackage{amsmath,amsthm}
\usepackage{unicode-math}
\setmainfont{LibertinusSerif}[Extension=.otf,UprightFont=*-Regular,BoldFont=*-Bold,ItalicFont=*-Italic,BoldItalicFont=*-BoldItalic]
\setsansfont{LibertinusSans}[Extension=.otf,UprightFont=*-Regular,BoldFont=*-Bold,ItalicFont=*-Italic]
\setmathfont{LibertinusMath-Regular.otf}
\setmonofont{DejaVu Sans Mono}[Scale=MatchLowercase]
\usepackage{mathpartir}
\usepackage[most]{tcolorbox}
\usepackage{booktabs,longtable,array,tabularx}
\usepackage{enumitem}
\usepackage{csquotes}
\usepackage{microtype}
\hypersetup{hidelinks}
\usepackage[capitalise,nameinlink]{cleveref}

% ---------------------------------------------------------------- tufte adjustments
% Number chapters, sections and subsections (tufte numbers nothing), so that \cref targets read "Section 9.8".
\setcounter{secnumdepth}{2}
\setcounter{tocdepth}{1}
% The width of text + margin, and the part of it that lies in the margin.
\newlength{\fullwidthlen}
\setlength{\fullwidthlen}{\textwidth}
\addtolength{\fullwidthlen}{\marginparsep}
\addtolength{\fullwidthlen}{\marginparwidth}
\newlength{\overhang}
\setlength{\overhang}{\marginparsep}
\addtolength{\overhang}{\marginparwidth}
% A longtable starts at the left edge and may run into the margin: give its columns widths in \fullwidthlen.
\setlength{\LTleft}{0pt}
\setlength{\LTright}{0pt plus 1fill minus \overhang}
% Inside a reading box, code and tables keep to the box (which is already full width).
\newif\ifinbox \inboxfalse

% ---------------------------------------------------------------- code
% fvextra, not listings: listings scrambles code points above 255 (λ, α, ∀, ...) under XeTeX.
% Environments: lprolog (λProlog source), scheme (aux Scheme source), answers (actual output of a run);
% each takes optional fvextra settings, e.g. \begin{scheme}[fontsize=\footnotesize]. Outside a box they
% span the full width. Inline: \lp{...} and \scm{...} (balanced braces; not inside section titles or
% captions, use \texttt there).
\usepackage{fvextra}
\fvset{fontsize=\small,breaklines=true,frame=leftline,framerule=1.2pt,xleftmargin=1.2em,framesep=0.6em,tabsize=2}
\newcommand{\codewidebegin}{\ifinbox\else\fvset{xrightmargin=-\overhang}\fi}% local to the environment
\newcommand{\codewideend}{}
\DefineVerbatimEnvironment{lprologV}{Verbatim}{rulecolor=\color{Plum!60}}
\DefineVerbatimEnvironment{schemeV}{Verbatim}{rulecolor=\color{OliveGreen!60}}
\DefineVerbatimEnvironment{answersV}{Verbatim}{fontsize=\footnotesize,formatcom=\color{black!75},rulecolor=\color{Gray!30}}
% \VerbatimEnvironment comes first, so that fancyvrb waits for \end{scheme}, not \end{fullwidth}; the inner environment
% parses the optional [...] itself (a wrapper with its own optional argument would eat the line end).
\newenvironment{lprolog}{\VerbatimEnvironment\codewidebegin\begin{lprologV}}{\end{lprologV}\codewideend}
\newenvironment{scheme}{\VerbatimEnvironment\codewidebegin\begin{schemeV}}{\end{schemeV}\codewideend}
\newenvironment{answers}{\VerbatimEnvironment\codewidebegin\begin{answersV}}{\end{answersV}\codewideend}
\newcommand{\lp}{\Verb[breaklines=true,fontsize=\small,formatcom=\color{Plum!70!black}]}
\newcommand{\scm}{\Verb[breaklines=true,fontsize=\small,formatcom=\color{OliveGreen!70!black}]}

% ---------------------------------------------------------------- theorem-like environments
\theoremstyle{definition}
\newtheorem{definition}{Definition}[chapter]
\newtheorem{example}[definition]{Example}
\theoremstyle{plain}
\newtheorem{theorem}[definition]{Theorem}
\newtheorem{proposition}[definition]{Proposition}
\newtheorem{lemma}[definition]{Lemma}
\newtheorem{corollary}[definition]{Corollary}
\theoremstyle{remark}
\newtheorem{remark}[definition]{Remark}

% ---------------------------------------------------------------- reading boxes
% tutorial: informal walk-through; kanren: how the concept maps to (aux kanren micro); pitfall: a trap the
% book stresses. They span the full width; inside them code and tables keep to the box. No \sidenote or
% \marginnote inside a box (the margin is taken by the box).
\tcbset{hkbox/.style={enhanced,breakable,grow to right by=\overhang,fonttitle=\bfseries\sffamily,
  boxrule=0.6pt,arc=1.5pt,before upper={\inboxtrue}}}
\newtcolorbox{tutorial}[1][]{hkbox,colback=blue!3,colframe=blue!45!black,
  title={Tutorial\ifx\relax#1\relax\else: #1\fi}}
\newtcolorbox{kanren}[1][]{hkbox,colback=green!3,colframe=green!35!black,
  title={Kanren angle\ifx\relax#1\relax\else: #1\fi}}
\newtcolorbox{pitfall}[1][]{hkbox,colback=orange!4,colframe=orange!60!black,
  title={Pitfall\ifx\relax#1\relax\else: #1\fi}}

% ---------------------------------------------------------------- notation
\newcommand{\book}[1]{\textup{[MN~p.\,#1]}}          % a PRINTED page of Miller & Nadathur (= PDF page - 16)
\newcommand{\imp}{\supset}                              % implication, as in the book
\newcommand{\seq}[3]{#1;#2\longrightarrow#3}            % Σ; P ⟶ G, a sequent of uniform proof search
\newcommand{\Sig}{\Sigma}
\newcommand{\Prog}{\mathcal{P}}
\newcommand{\Llam}{L_\lambda}
\newcommand{\fohc}{\textit{fohc}}  \newcommand{\fohh}{\textit{fohh}}
\newcommand{\hohc}{\textit{hohc}}  \newcommand{\hohh}{\textit{hohh}}
\newcommand{\typ}[1]{\mathsf{#1}}                        % a type or sort, e.g. \typ{o}, \typ{int}
\newcommand{\cst}[1]{\mathsf{#1}}                        % a constant of a signature
\newcommand{\abs}[2]{\lambda #1.\,#2}
\newcommand{\eqb}{=_{\beta\eta}}
\newcommand{\subst}[3]{#1[#2/#3]}                        % t[s/x]
\newcommand{\rulename}[1]{\textsc{#1}}
\newcommand{\aux}{\textsf{aux}}
\newcommand{\micro}{\textsf{(aux kanren micro)}}


\title[Higher-Order Logic Programming for (aux kanren micro)]%
  {Higher-Order Logic Programming for \mbox{\texttt{(aux kanren micro)}}}
\author[Prepared for Massimo Nocentini, with Claude (Anthropic)]%
  {Prepared for Massimo Nocentini, with Claude (Anthropic)}
\date{October 2026}

% Running heads and PDF metadata take these plain-text forms (no macros).
\renewcommand{\plaintitle}{Higher-Order Logic Programming for (aux kanren micro)}
\renewcommand{\plainauthor}{Prepared for Massimo Nocentini, with Claude (Anthropic)}
\hypersetup{pdftitle={Higher-Order Logic Programming for (aux kanren micro)},
  pdfauthor={Prepared for Massimo Nocentini, with Claude (Anthropic)},bookmarksdepth=subsection}

% Table of contents: tufte's section entries give the number a 2em box, too narrow for "13.10" or "A.10".
\titlecontents{section}[0em]
  {\vspace{0\baselineskip}\begin{fullwidth}\Large\rmfamily\itshape}
  {\hspace*{2.6em}\contentslabel{2.6em}}
  {\hspace*{2.6em}}
  {\rmfamily\upshape\qquad\thecontentspage}
  [\end{fullwidth}]

% Part pages in the style of tufte's chapter headings (italic, flush left, across the full width), in place of the
% book class's bold centered ones.
\titleformat{\part}[display]
  {\relax\begin{fullwidth}}
  {\itshape\huge\partname~\thepart}
  {1.5ex}
  {\Huge\rmfamily\itshape}
  [\end{fullwidth}]
% \part* (only Appendices): once titlesec formats parts, it also writes their ToC line, with the part page's number.
\titleformat{name=\part,numberless}[display]
  {\relax\begin{fullwidth}}
  {}
  {0pt}
  {\Huge\rmfamily\itshape}
  [\end{fullwidth}]

% Chapter titles that wrap: ragged right across the full width (the class option `justified' would
% stretch the first line of the title).
\titleformat{\chapter}[display]
  {\relax\begin{fullwidth}\RaggedRight}
  {\itshape\huge\thechapter}
  {0pt}
  {\huge\rmfamily\itshape}
  [\end{fullwidth}]

% Margin material (sidenotes, margin notes, captions) ragged right, as tufte sets it by default: the option
% `justified' is meant for the text column, and a justified 49mm column opens wide gaps.
\makeatletter
\renewcommand*{\@tufte@sidenote@justification}{\RaggedRight}
\renewcommand*{\@tufte@marginnote@justification}{\RaggedRight}
\renewcommand*{\@tufte@caption@justification}{\RaggedRight}
\makeatother

% The book class's thebibliography sets its running head in capitals; tufte sets the others as small caps of the title.
\makeatletter
\patchcmd{\thebibliography}{\@mkboth{\MakeUppercase\bibname}{\MakeUppercase\bibname}}{\@mkboth{\bibname}{\bibname}}{}%
  {\PackageWarningNoLine{hol-kanren}{could not patch the running head of thebibliography}}
\makeatother

% A paragraph that finds no break within the tolerance (a long cross-reference or inline formula in the
% 107mm column) may spread a little extra space over its lines instead of running into the margin.
\setlength{\emergencystretch}{1.5em}

\begin{document}

\frontmatter

% ---------------------------------------------------------------- title page (tufte has no \subtitle)
\begin{titlepage}
\begin{fullwidth}
\vspace*{0.12\textheight}
\noindent{\fontsize{30}{36}\selectfont\itshape Higher-Order Logic Programming\\[2pt]
  for \texttt{(aux kanren micro)}\par}
\vspace{2.5\baselineskip}
\noindent{\Large\itshape A reading of Miller \& Nadathur's \enquote{Programming with Higher-Order Logic}
  and of λKanren, with a proposal for \texttt{(aux kanren lambda)}\par}
\vfill
\noindent{\large Prepared for Massimo Nocentini, with Claude (Anthropic)\par}
\vspace{0.6\baselineskip}
\noindent{\large\itshape October 2026\par}
\vspace*{0.05\textheight}
\end{fullwidth}
\end{titlepage}

% ---------------------------------------------------------------- abstract
\chapter*{Abstract}
\noindent \texttt{(aux kanren micro)} is a small, fast miniKanren whose terms are first-order Scheme data and
whose goals are first-order Horn clauses. λProlog, the language of Miller and Nadathur's
\emph{Programming with Higher-Order Logic}, runs a much larger logic: goals may say \enquote{for every new
name $x$} and \enquote{assuming $D$}, terms may contain λ-binders, and unification works modulo α, β
and η, so that programs over syntax with binders need no code for renaming, fresh names or
capture-avoiding substitution. This document asks what it would take to give \texttt{(aux kanren micro)}
the same power. Part~I summarizes the book chapter by chapter, with its core examples; Part~II surveys
the miniKanren literature on these features, centered on Ma and Friedman's higher-order unification
algorithm for λKanren; Part~III describes the engine as it is today and proposes a concrete answer: a
generic extension protocol for the engine and a client module, \texttt{(aux kanren lambda)}, providing
λ-terms, the goals $\forall$ and $\supset$, and higher-order pattern unification with universe levels and
delayed non-patterns, prototyped, tested and benchmarked.

\tableofcontents

\mainmatter

\chapter{Introduction}\label{ch:intro}

\micro{} is a small, fast miniKanren. Its terms are first-order Scheme data, its unifier is first-order
unification with an occurs check, and its goals are conjunctions, disjunctions and existentials. In
logical terms it runs first-order Horn clauses (\cref{sec:horn-micro}). λProlog, the language of Miller and
Nadathur's \emph{Programming with Higher-Order Logic}~\cite{MillerNadathur2012}, runs a much larger logic.
Goals may say \enquote{for every new name $x$} and \enquote{assuming $D$}, terms may contain λ-binders, and
unification works modulo α, β and η. With these three features, programs over syntax with binders (type
checkers, evaluators, provers, program transformations) contain no code for renaming, fresh names or
capture-avoiding substitution, because the logic does that work.

This document asks what it would take to give \micro{} the same power, and it ends with a concrete answer:
a module, \texttt{(aux kanren lambda)}, prototyped, tested and benchmarked. This chapter says what the
document is for (\cref{sec:intro-goal}), runs one example through the whole story
(\cref{sec:intro-brief}), maps the chapters (\cref{sec:intro-map}), proposes three ways to read them
(\cref{sec:intro-paths}), and fixes the conventions (\cref{sec:intro-conventions}).

% ==================================================================================================
\section{What this document is for}\label{sec:intro-goal}

\paragraph{The goal.} The goal is to bring three ideas of λProlog into \micro:
\begin{enumerate}[nosep]
  \item \emergencystretch=2em \textbf{hereditary Harrop goals}, that is the universal goal $\forall x\,G$ (\lp{pi x\ G}, proved
        with a new scoped constant) and the hypothetical goal $D\imp G$ (\lp{D => G}, proved with the
        program temporarily extended by the clause $D$);
  \item \textbf{λ-tree syntax}: object-level binders represented by meta-level λ-ab\-strac\-tions, so that
        α-equivalence is built in and object-level substitution is meta-level β-reduction;
  \item \textbf{higher-order pattern unification} ($\Llam$): unification of λ-terms restricted to the
        fragment where it is decidable and has most general unifiers, with the other equations delayed.
\end{enumerate}
\noindent The three are not independent. \lp{pi} is what lets a program walk under a λ-binder, and
pattern unification is what lets it build a λ-term back. The book's programs use all three together.

\paragraph{The sources.} There are two. The first is the book itself, which Part~I summarizes chapter by
chapter, with its core examples run on \micro{} wherever that is possible today. The second is the
miniKanren literature on these features, which Part~II surveys. Its center is
Ma and Friedman's \emph{A New Higher-order Unification Algorithm for λKanren}~\cite{MaFriedman2021}, read
together with its predecessor λKanren~\cite{MaLuFriedman2020}, nominal unification and
αKanren~\cite{UrbanPittsGabbay2004,ByrdFriedman2007}, the eigenvariable and negation
work~\cite{HemannByrdFriedman2013,Moiseenko2019,JinEtAl2021}, and the techniques of real λProlog systems:
universe levels~\cite{NadathurJayaramanKwon1995}, Miller's and Nipkow's pattern
unification~\cite{Miller1991,Nipkow1993}, delayed equations~\cite{DowekHardinKirchnerPfenning1996},
Teyjus~\cite{NadathurLinnell2005} and ELPI~\cite{DunchevEtAl2015}.

\paragraph{The outcome.} Part~III is a proposal in two pieces.
\begin{itemize}[nosep]
  \item \emergencystretch=2em A \textbf{generic extension protocol} in \micro{} (\texttt{micro.diff}, $+235/{-}41$ lines,
        comments included; 214 of the added lines are not blank). The state gains one slot that \micro{} never reads, and a registry of hooks opens the
        points that are closed today: the unifier's structural clause, the occurs check, the step after each
        \scm{=°}, the disequality, type and absento stores, and reification. Nothing in the edit mentions λ,
        levels or clauses.
  \item A \textbf{client module}, \texttt{(aux kanren lambda)}, 1255 lines, which registers one extension.
        It provides λ-terms (\scm{ƛ}, \scm{@}), the goals \scm{∀°} and \scm{⊃°} with
        \scm{define-dynamic-relation}, and $\Llam$ unification with universe levels and delayed
        non-patterns. Through the hooks it also makes \micro's own \scm{≠°}, \scm{absent°},
        \scm{symbol°} and \scm{number°} λ-aware; it exports no second version of them.
\end{itemize}
\noindent The design was not chosen on paper. Three prototypes were built: one that leaves \micro{}
untouched, one with a hook protocol, one with λ-terms native in the core. They were compared by three
reviews, on engineering, on fidelity to the book and on pedagogy. The hook design, ranked first by two of
the three, became the base, and about twenty grafts proposed from the other two were each adopted or
explicitly declined. Three adversarial review rounds
then attacked the result and sent thirteen, eight and twelve reports. All were real. Every defect but
one was fixed and pinned by a regression case\sidenote{A cost reported by the third round among them.}; the
exception, the cost of the scope check on first-order search once a \scm{∀°} has run, reported by the
second round, is measured and documented as a limitation. The final suite has
83 cases with 320 hand-written assertions, plus 6 cases that pin the known limits, and \micro's own 125
cases pass unchanged, with and without the module loaded: 339 cases, no failures. The protocol costs about
1\% on first-order search when no extension is loaded. All of this lives in
\texttt{doc/hol-kanren/prototype/}, and after merging it would be \texttt{src/aux.kanren.lambda.scm} and
\texttt{src/test/microkanren-lambda*.scm}.

% ==================================================================================================
\section{The whole story in brief}\label{sec:intro-brief}

This section follows one program from the book through the whole document: type inference for the untyped
λ-calculus \book{190--191}. It is two clauses of λProlog, and it uses every feature at once.

\subsection{The program, twice}

Object-level λ-terms are built with two constructors, \lp{app} for application and \lp{abs} for
abstraction. \lp{abs} takes a \emph{meta-level} abstraction: the identity is \lp{abs x\ x}. Types are
built with \lp{arr}. The relation \lp{typeof M T} says that term \lp{M} has type \lp{T}:
\begin{lprolog}
typeof (app M N) A :- typeof M (arr B A), typeof N B.
typeof (abs M) (arr A B) :- pi x\ typeof x A => typeof (M x) B.
\end{lprolog}
The second clause reads: \enquote{to type \lp{abs M}, take a new constant \lp{x}, assume that \lp{x} has
type \lp{A}, and type the body \lp{(M x)} with \lp{B}}. There is no clause for variables and no context
argument. The hypothesis \lp{typeof x A} \emph{is} the context, and it disappears when the proof of the
body is over. In the proposed module the same program reads:
\begin{scheme}
(define-dynamic-relation (typeof° m t)
  (cond°
    ((fresh° (M N B) (=° m `(app ,M ,N)) (typeof° M `(arr ,B ,t)) (typeof° N B)))
    ((fresh° (R A B) (=° m `(abs ,R)) (=° t `(arr ,A ,B))
       (∀° (x) (⊃° ((typeof° x A)) (typeof° (@ R x) B)))))))
\end{scheme}
\noindent The dictionary is short. \lp{pi x\ G} is \scm{(∀° (x) G)}, \lp{D => G} is \scm{(⊃° (D) G)}, the
application \lp{(M x)} of a logic variable is \scm{(@ R x)}, and \lp{x\ t} is \scm{(ƛ (x) t)}. The relation
is made by \scm{define-dynamic-relation} instead of \scm{define-relation} because it receives hypotheses.
The first-order structure (\scm{abs}, \scm{app}, \scm{arr}) stays ordinary Scheme data, as everywhere in
\micro.

\subsection{Running it forwards}

The book types the identity and fails on self-application \book{191}:
\begin{lprolog}
?- typeof (abs x\x) Ty.          Ty = arr T1 T1
?- typeof (abs x\ app x x) Ty.   no
\end{lprolog}
The module gives the same answers. \scm{°->list/datum} lists the answers as plain data, and \scm{α} is a
variable that the answer leaves unbound:
\begin{scheme}
(°->list/datum (fresh° (q) (typeof° `(abs ,(ƛ (x) x)) q)))
(°->list/datum (fresh° (q) (typeof° `(abs ,(ƛ (x) `(app ,x ,x))) q)))
\end{scheme}
\begin{answers}
((arr α α))
()
\end{answers}

Two rules of the book's proof search are at work, GENERIC for $\forall$ and AUGMENT for $\imp$
(\cref{sec:hh-rules}). In the sequent $\seq{\Sig}{\Prog}{G}$, $\Sig$ is the signature, the constants in
scope, and $\Prog$ is the program:
\begin{mathpar}
\inferrule*[right=\rulename{generic}]{\seq{\Sig\cup\{c\}}{\Prog}{\subst{G}{c}{x}}}{\seq{\Sig}{\Prog}{\forall x\,G}}
\quad (c\notin\Sig)
\and
\inferrule*[right=\rulename{augment}]{\seq{\Sig}{\Prog\cup\{D\}}{G}}{\seq{\Sig}{\Prog}{D\imp G}}
\end{mathpar}
Read bottom-up, GENERIC proves a universal goal by inventing a constant, and AUGMENT proves an implication
by assuming its antecedent. Both changes are local: they hold for the proof of the premise only.

\begin{tutorial}[What happens to the state while \texttt{typeof°} types λx.x]
Write $d$ for the current \emph{depth}, the number of enclosing \scm{∀°}, and $\Gamma$ for the assumed
clauses. Both live in the search state, in the slot that the protocol adds; the substitution is
\micro's, as before. Each row of the table below is one step towards the first answer, and two of them
carry the ideas. At step~5, $R$ is already bound to $\abs{x}{x}$, so \texttt{(@ R c)} reduces by β to $c$
itself: object-level substitution is meta-level β. At step~6 the dynamic relation tries its hypotheses
first, most recent first, and then its two static clauses, which both fail on $c$, since an
eigenvariable is neither an \texttt{app} nor an \texttt{abs}. The hypothesis \texttt{typeof° c A}
mentions the variable $A$ of the enclosing clause, and $A$ is \emph{shared}, not renamed. That is how
typing the body tells the outer goal what $A$ is (\cref{sec:hh-link}).
\begin{center}\small
\begin{tabular}{@{}r>{\raggedright\arraybackslash}p{5.6cm}cl@{}}
\toprule
& goal & $d$ & effect \\
\midrule
1 & \texttt{typeof° (abs (ƛ (x) x)) q} & 0 & the \texttt{app} clause fails; the \texttt{abs} clause applies\\
2 & \texttt{(abs R) = (abs (ƛ (x) x))} & 0 & first-order: $R := \abs{x}{x}$, \texttt{q := (arr A B)}\\
3 & \texttt{∀° (x)} & 1 & a new eigenvariable $c$, of level 1\\
4 & \texttt{⊃° ((typeof° c A))} & 1 & $\Gamma = \{\texttt{typeof° c A}\}$\\
5 & \texttt{typeof° (@ R c) B} & 1 & \texttt{(@ R c)} head-normalizes by β to $c$\\
6 & backchain on the hypothesis & 1 & \texttt{typeof° c A = typeof° c B}: $A$, $B$ unified\\
7 & answer leaves \texttt{⊃°} and \texttt{∀°} & 0 & $\Gamma=\emptyset$; \texttt{q = (arr A A)}, printed \texttt{(arr α α)}\\
\bottomrule
\end{tabular}
\end{center}
Self-application fails at the same step. For \texttt{(abs (ƛ (x) (app x x)))}, the body needs
$\texttt{typeof° c (arr B$'$ B)}$ and $\texttt{typeof° c B$'$}$, so the hypothesis forces
$A = \texttt{(arr B$'$ B)}$ and $A = B'$: first-order unification fails on its occurs check, as in the book.
\end{tutorial}
\noindent The book explains this program in \cref{sec:lambdatree-typeof}; \cref{sec:module-typeof} traces
it again at the level of the module's state, and \cref{sec:examples-typeof} runs it on the book's whole
session, the combinator S included.

\subsection{Running it backwards}

Ask instead for terms of type $i\to i$. λProlog's depth-first search loops on this query, because it tries
the \lp{app} clause first and descends forever. The interleaving search of \micro{} does not:
\begin{scheme}
(°->list/datum (take° 3 (fresh° (q) (typeof° q '(arr i i)))))
\end{scheme}
\begin{answers}
((abs (ƛ (x) x))
 (abs (ƛ (x) (app (abs (ƛ (x1) x1)) x)))
 (abs (ƛ (x) (app (abs (ƛ (x1) x)) x))))
\end{answers}
\noindent The second answer is $\lambda x.\,(\lambda x_1.x_1)\,x$. The printer renames the inner binder
to \texttt{x1}, apart from the outer one. The first answer is the instructive one. Now $R$ is unknown when
step~5 is reached, so backchaining on the hypothesis must solve the equation $R\,c \doteq c$, where $R$ was
created outside the \scm{∀°} and $c$ inside it. That is a \emph{higher-order} unification problem, and it
is one of the easy ones.

\begin{definition}[Levels and patterns]\label{def:intro-levels}
Every logic variable $X$ has a \emph{level} $\ell(X)$, the depth at which it was created. Every
eigenvariable $c$ has the level $\ell(c)$ of the \scm{∀°} that made it, one more than the depth outside.
\begin{enumerate}[nosep]
  \item \emph{Scope check.} A binding $X := t$, made by any \scm{=°}, is accepted only if every
        eigenvariable $c$ in $t$ has $\ell(c)\le\ell(X)$. Each unbound variable $Y$ of $t$ with
        $\ell(Y)>\ell(X)$ is then \emph{lowered}: replaced by a new variable of level $\ell(X)$ or, when
        $Y$ occurs applied, \emph{pruned} of the arguments that $X$ may not see.
  \item \emph{Pattern.} An application $F\,a_1\dots a_n$ of an unbound variable $F$ is a \emph{pattern}
        when its arguments, after head normalization, are distinct eigenvariables, each with
        $\ell(a_i)>\ell(F)$.
\end{enumerate}
The full rule, including the cases that must wait for a later binding, is in \cref{sec:module-unif}.
\end{definition}
\noindent The first condition is the \emph{scope check}. It is the operational content of the order of
quantifiers in $\exists R\,\forall c$: $R$ was chosen before $c$ existed, so its value cannot mention $c$
(\cref{sec:hh-scope,sec:lambda-unifprob}). Levels replace the book's explicit quantifier prefix, as Teyjus
does~\cite{NadathurJayaramanKwon1995}. The second condition defines Miller's fragment $\Llam$, in which
unification is decidable and has most general unifiers~\cite{Miller1991} (\cref{sec:hounif-llambda}).

\begin{example}[Scope and patterns]\label{ex:intro-scope}
Here $\ell(R)=0$ and $\ell(c)=1$. The equation $R\doteq c$ violates the scope check and fails. The equation
$R\,c\doteq c$ is a pattern: abstracting $c$ out of the right-hand side gives the most general unifier
$R := \abs{x}{x}$. In the module, with \scm{q} in the role of $R$:
\begin{scheme}
(°->list #f (fresh° (q) (∀° (x) (=° q x))))
(°->list/datum (fresh° (q) (∀° (x) (=° (@ q x) x))))
\end{scheme}
\begin{answers}
()
((ƛ (x) x))
\end{answers}
The first is the book's \lp{sigma x\ pi y\ x = y}, which is not provable \book{85}: no single term
equals every new constant.
\end{example}

The other two answers of the backwards run come from the \lp{app} clause under the binder: the body
\scm{(@ R c)} must unify with \scm{(app M N)}, where $M$ and $N$ were created inside the \scm{∀°}, at level~1.
Variable elimination (\cref{sec:hounif-algo}) builds $R := \abs{y}{\texttt{(app (M$'$ y) (N$'$ y))}}$ and
\emph{lowers} $M$ and $N$ to new variables $M'$, $N'$ of level~0, applied to $c$: $M := M'\,c$,
$N := N'\,c$. Scheme code built none of these λ-terms; unification did.

\paragraph{When an equation is not a pattern.} Some equations have no most general unifier. The book's
$\forall x.\,F\,x\,x\doteq x$ \book{208} has two incomparable solutions, $\abs{a}{\abs{b}{a}}$ and
$\abs{a}{\abs{b}{b}}$, and $F\,x\,x$ is not a pattern, since its arguments repeat. The module does what
Teyjus does, \enquote{dynamic $\Llam$} (\cref{sec:hounif-dynamic}): it suspends the equation, retries
it after every later binding, which may have turned it into a pattern, and reports it in the answer if it
is still there at the end. Under \scm{°->pretty}, which prints constraints in the style of αKanren, the residual equation is
closed over $x$ when it leaves the scope of the \scm{∀°}. A later binding decides it:
\begin{scheme}
(°->pretty (fresh° (q) (∀° (x) (=° (@ q x x) x))))
(°->list/datum (fresh° (q) (∀° (x) (=° (@ q x x) x)) (=° q (ƛ (a b) b))))
\end{scheme}
\begin{answers}
((α : (= (ƛ (x) (@ α x x)) (ƛ (x) x))))
((ƛ (a b) b))
\end{answers}

\subsection{What the engine had to learn}

None of the above runs on today's \micro. \Cref{tab:intro-concepts} lists what the example needed, where
the book explains it, and where this document builds it.
\begin{table*}[h]
\small
\begin{tabularx}{\linewidth}{@{}>{\raggedright\arraybackslash}p{4.4cm}>{\raggedright\arraybackslash}X>{\raggedright\arraybackslash}p{4.3cm}@{}}
\toprule
concept & in the example & theory \,/\, in the module\\
\midrule
λ-tree syntax, α-equivalence & \texttt{(abs (ƛ (x) x))}, binders printed apart & \cref{ch:lambdatree} \,/\, \cref{ch:module}\\
β-reduction, application & \texttt{(@ R x)} is \texttt{x} once \texttt{R := (ƛ (x) x)} & \cref{ch:lambda} \,/\, \cref{ch:module}\\
GENERIC, eigenvariables & \texttt{(∀° (x) …)} & \cref{ch:hh} \,/\, \cref{ch:design,ch:module}\\
AUGMENT, dynamic relations & \texttt{(⊃° ((typeof° x A)) …)} & \cref{ch:hh} \,/\, \cref{ch:module}\\
scope check and levels & \texttt{R = x} fails & \cref{ch:hh,ch:hounif} \,/\, \cref{ch:module}\\
pattern unification & \texttt{R x = x} gives \texttt{R := λx.x} & \cref{ch:hounif} \,/\, \cref{ch:module}\\
dynamic $\Llam$, residuals & \texttt{F x x = x} waits & \cref{ch:hounif} \,/\, \cref{ch:module}\\
interleaving search & typing backwards terminates & \cref{ch:horn} \,/\, \cref{ch:micro}\\
\bottomrule
\end{tabularx}
\caption{The running example, concept by concept.}\label{tab:intro-concepts}
\end{table*}

The state of \micro{} has a fixed set of fields, and its unifier, walks and reifier are closed
(\cref{sec:micro-ext}). The proposal opens them once, generically, and puts everything λ-specific in a
client module. In the state's new slot the module keeps the depth, the level of each variable, the
assumed clauses $\Gamma$, the suspended constraints (equations, and types or absentos on flexible
terms), and the disequalities that involve λ-terms or eigenvariables. Through the hooks it receives every
unification problem that contains one of its terms, it checks the scope of the bindings that \scm{=°}
makes, and it takes part in printing. Because the hooks sit inside \micro's own \scm{=°}, every relation
built on \scm{=°}, such as \micro's \scm{cons°} or the user's own, respects the scope of \scm{∀°} with no
change. A prototype that wrapped \micro{} from the outside instead let eigenvariables escape through exactly
such relations (\cref{ch:design}).

\begin{kanren}[Why not just use Scheme closures?]
A Scheme procedure can stand for a binder as long as it is only \emph{called}: \lp{mappred} and continuation
passing work on today's \micro{} (\cref{sec:hoprog-free}). The backwards run above needs more: unification
must look \emph{inside} the abstraction, compare it modulo α, and build one. A closure can do none of these,
so \scm{ƛ} builds a data structure, a locally nameless term whose binders are de Bruijn
indices~\cite{deBruijn1972,Chargueraud2012}, which prints with the user's names.
\end{kanren}

\begin{pitfall}[Hypotheses need a dynamic relation]
In λProlog every predicate can receive clauses through \lp{=>}. In the module only a relation made by
\scm{define-dynamic-relation} can, because a relation made by \scm{define-relation} is a closed Scheme
procedure. Assuming a clause for it is an error, not a silent no-op:
\begin{scheme}
(define-relation (static° x) (=° x 1))
(define-dynamic-relation (dynamic° x) (=° x 1))
(°->list/datum (fresh° (q) (⊃° ((static° 2)) (static° q))))   ; an error
(°->list/datum (fresh° (q) (⊃° ((dynamic° 2)) (dynamic° q))))
\end{scheme}
\begin{answers}
Error: (⊃°) not a dynamic relation, see define-dynamic-relation: #<procedure (static° x)>
(2 1)
\end{answers}
\noindent The first line is what \texttt{csi} prints for the error, which stops the script; the last
query was run on its own. In the second answer the hypothesis comes first, then the
static clause, as in the book \book{78}.
\end{pitfall}

% ==================================================================================================
\section{Map of the document}\label{sec:intro-map}

The document has three parts and two appendices. Part~I follows the book, Part~II the literature, and
Part~III builds the proposal.

\subsection*{Part I: the book}

\paragraph{\Cref{ch:horn}, Logic programming as proof search.} The book's Introduction and Chapters~1--2:
typed first-order terms, first-order unification, the search semantics of goals, and backchaining on Horn
clauses (\fohc). It establishes that pure miniKanren \emph{is} \fohc{} without types and with an interleaving
search, and it shows where the interleaving finds answers that depth-first search misses.

\paragraph{\Cref{ch:hh}, Hereditary Harrop formulas.} Chapter~3 of the book: AUGMENT and GENERIC, hypothetical
reasoning on a course database, scoped constants, why $\exists\forall$ and $\forall\exists$
differ, and how assumed clauses share variables with their context. The logic is intuitionistic, indeed
minimal, not classical. Its last section lists what $\forall$ and $\imp$ ask of a miniKanren: the program
context must live in the search state, and scoped constants need a scope check.

\paragraph{\Cref{ch:lambda}, Typed λ-terms and unification problems.} Chapter~4: simply typed λ-terms,
conversion by α, β and η as the equality of the logic, normal forms, unification problems as quantified formulas,
raising, and why higher-order unification is undecidable (Post's correspondence problem, Hilbert's tenth
problem).

\paragraph{\Cref{ch:hoprog}, Higher-order programming.} {\emergencystretch=2em Chapters~5--6: rigid and flexible atoms, the
languages \hohc{} and \hohh{}, predicate variables (\lp{mappred}, folds, continuation passing),
function variables solved by unification (functional difference lists), and modules as logic. It sorts
the book's programs into those that run on today's \micro{} (they only \emph{call} a predicate) and those
that need λ-terms (they look \emph{inside} one).\par}

\paragraph{\Cref{ch:lambdatree}, λ-tree syntax.} {\emergencystretch=2em Chapter~7, the chapter that motivates the proposal:
object-level binders as meta-level abstractions, substitution as β, the \emph{mobility} of binders
(term-level λ, then \lp{pi}, then an eigenvariable), and the book's programs over λ-terms and formulas:
evaluators, type inference, de~Bruijn translation, normal forms, prenex forms, copy clauses, and
substitution run backwards. It ends with the $\Llam$ subset.\par}

\paragraph{\Cref{ch:hounif}, Unification of λ-terms and higher-order patterns.} Chapter~8: Huet's procedure
and its matching trees, the pattern fragment, the pattern unification algorithm as one walk of variable
elimination (inversion, pruning, occurs and scope checks), traces of the book's examples, Nipkow's
de~Bruijn algorithm with a running Scheme transliteration, and the pragmatics of real systems: universe
levels in place of raising, and delayed non-patterns.

\paragraph{\Cref{ch:apps}, Applications.} Chapters~9--11: provers for intuitionistic and classical logic,
natural deduction with proof terms, tactics, typing and evaluation of functional programs, partial
evaluation, CPS conversion, and the π-calculus. For each application it records which of $\forall$, $\imp$
and λ-term unification it needs, and which parts run on \micro{} today.

\subsection*{Part II: the literature}

\paragraph{\Cref{ch:lkanren}, Higher-order miniKanren.} Ma and Friedman 2021 in full: grafting versus
substitution, binder merging and its \enquote{true = false} hazard, static closures, and a critical assessment
with three counter-examples. Then λKanren 2020, nominal unification and αKanren, eigenvariables and negation
inside miniKanren, and the implementation techniques of Miller, Nipkow, Dowek et al., Nadathur et al.,
Teyjus, ELPI and Makam. It ends with a comparison of the design options (\cref{sec:lkanren-compare}) and
deliberately does not choose.

\subsection*{Part III: the proposal}

\paragraph{\Cref{ch:micro}, Anatomy of \micro.} {\emergencystretch=2em The engine as it is today, at commit \texttt{f53bdcb}: terms
and variables, the six-field state, the walk, unification, the stores D (disequalities), A (type tags) and T
(absento), the interleaving search, reification and aggregation. It closes with what is open to extension and
what is closed, and with three tricks that need no edit (including working AUGMENT and GENERIC prototypes)
together with their limits.\par}

\paragraph{\Cref{ch:design}: requirements, design space, protocol.} The nine acceptance items A1--A9
collected from Parts~I and~II (the scope of $\forall$ at every binding, AUGMENT with the book's clause
order, type inference, λ-tree syntax, $\Llam$ with delay, λProlog's answers where Ma and Friedman differ,
the book's programs, \micro's stores on λ-terms, and \micro's own suites unchanged). Then the design space,
the four meanings of \scm{=°} on λ-terms left open by \cref{sec:lkanren-compare}, of which pattern
unification with levels is retained; the three prototypes (a layer on top of \micro, hooks in \micro, a
native engine) and the three reviews that ranked them; and the extension protocol that came out of it:
the state's seventh field, the registry, the thirteen hooks with their contracts and call sites, the
invariants every client must keep, and its cost. It ends with the history of the adversarial review rounds.

\paragraph{\Cref{ch:module}: the module.} \texttt{(aux kanren lambda)} in detail: locally nameless terms
built by \scm{ƛ} and \scm{@}; the extension's state (depth, levels, $\Gamma$, suspended constraints,
λ-disequalities), copied and never mutated; the unification rules with variable elimination, the eager
scope check, the fast paths and the wake-up loop, with worked traces; \scm{∀°} and \scm{⊃°}, including
the discharge of constraints at the exit of a scope and backchaining in \scm{define-dynamic-relation};
disequality, absento and types on λ-terms; reification; and the API reference.

\paragraph{\Cref{ch:examples}: the book's examples in the module.} The core programs of Part~I, each beside its
λProlog original and with its actual output: α, β and shadowing; the scope of eigenvariables; the course
database, jars and bugs, and reverse with assumed clauses; type inference; call-by-name evaluation, copy and
substitution, existential generalization; de~Bruijn translation and prenex forms; \lp{mappred} and
palindromes with functional difference lists; Fischer's CPS transformation; and the unification examples of
the book's Chapter~8 one by one. It also runs Ma and Friedman's queries and explains each difference, and
closes with a table of all the examples.

\paragraph{\Cref{ch:roadmap}: tests, limitations, performance, merge plan.} {\emergencystretch=2em How the proposal is tested
(three layers of tests, and \micro's suites run twice); nine documented limitations and deviations from
λProlog, each pinned by a test (among them: only dynamic relations take hypotheses, terms are untyped,
non-patterns are suspended and never enumerated, answers come in \micro's interleaving order); the
performance measurements; a merge in three stages (first the protocol, then the module and its suites,
then the show extension); six later work items; and nine questions only the maintainer of \aux{} can
answer.\par}

\subsection*{Appendices}

\paragraph{\Cref{app:rules}, Reference card.} The formal rules on a few pages: λProlog syntax and its
logical reading, the goal and clause grammars, uniform proofs, typing, λ-conversion, raising, first-order
unification, Huet's procedure, pattern unification and universe levels, each with its page in the book.

\paragraph{\Cref{app:glossary}, Glossary.} The vocabulary of the three communities (proof theory,
unification, miniKanren) in one alphabetical list, with pointers to the chapters and to the book, and a
bridge between their names for the same idea.

% ==================================================================================================
\section{Reading paths}\label{sec:intro-paths}

The chapters are written to be read in order, but they need not be. \Cref{tab:intro-paths} summarizes the
three paths below; page counts are approximate.

\paragraph{The fast path.} The minimum to understand the proposal, about 45 pages before
\cref{ch:design,ch:module}:
\begin{enumerate}[nosep]
  \item \cref{sec:intro-brief} of this chapter (4 pages);
  \item the two rules, the scope condition, and how quantification links goals and clauses:
        \cref{sec:hh-rules,sec:hh-scope,sec:hh-link} (5 pages);
  \item λ-tree syntax:
        \cref{sec:lambdatree-encoding,sec:lambdatree-beta,sec:lambdatree-mobility,sec:lambdatree-typeof}
        (9 pages);
  \item what rigid and flexible terms and the normal form of a unification problem are:
        \cref{sec:hounif-intro,sec:hounif-problems} (the introduction says enough about imitation and
        projection to skip \cref{sec:hounif-huet}); then patterns, their unification, the worked traces
        and what real systems do: \cref{sec:hounif-llambda,sec:hounif-algo,sec:hounif-traces,sec:hounif-pragmatics}
        (13 pages in all);
  \item the options: \cref{sec:lkanren-compare,sec:lkanren-summary} (3 pages; the summary defines each
        term the comparison uses);
  \item the engine's state, its constraint stores and its extension points:
        \cref{sec:micro-state,sec:micro-stores,sec:micro-ext} (8 pages; skip \cref{sec:micro-stores} if you
        know micro's D/A/T chain);
  \item then \cref{ch:design,ch:module} (about 60 pages), the limitations in \cref{sec:roadmap-limits},
        and the staged merge in \cref{sec:roadmap-merge} (2 pages), which says how to build and load
        the edited micro.
\end{enumerate}

\paragraph{The λProlog path.} Part~I alone, \cref{ch:horn} to \cref{ch:apps}, about 185 pages, is a
self-contained reading of the book. The \emph{Kanren angle} boxes can be skipped. \Cref{app:rules} serves as
a reference card along the way.

\paragraph{The implementer path.} For whoever will merge or extend the code:
\cref{ch:hounif} (the algorithm), \cref{ch:micro} (the engine), \cref{ch:design} (the protocol),
\cref{ch:module} (the module), and \cref{ch:roadmap} (tests, costs, merge plan), about 120 pages, with
\cref{sec:rules-pattern,sec:rules-levels} at hand. \Cref{ch:examples} then serves as a set of worked tests.

\begin{table}[htb]
\centering\small
\caption{Three reading paths. A filled bullet means read in full, a hollow one read the sections named in
the text. Pages are those of this document.}\label{tab:intro-paths}
\begin{tabular}{@{}lrccc@{}}
\toprule
chapter & pages & fast & λProlog & implementer\\
\midrule
\cref{ch:intro}, Introduction & 12 & $\circ$ & $\bullet$ & $\bullet$\\
\cref{ch:horn}, proof search & 28 & & $\bullet$ & \\
\cref{ch:hh}, hereditary Harrop & 24 & $\circ$ & $\bullet$ & \\
\cref{ch:lambda}, λ-terms & 22 & & $\bullet$ & \\
\cref{ch:hoprog}, higher-order programming & 30 & & $\bullet$ & \\
\cref{ch:lambdatree}, λ-tree syntax & 30 & $\circ$ & $\bullet$ & \\
\cref{ch:hounif}, pattern unification & 24 & $\circ$ & $\bullet$ & $\bullet$\\
\cref{ch:apps}, applications & 28 & & $\bullet$ & \\
\cref{ch:lkanren}, the literature & 30 & $\circ$ & & \\
\cref{ch:micro}, the engine & 22 & $\circ$ & & $\bullet$\\
\cref{ch:design}, design and protocol & 26 & $\bullet$ & & $\bullet$\\
\cref{ch:module}, the module & 32 & $\bullet$ & & $\bullet$\\
\cref{ch:examples}, the examples & 28 & & & $\bullet$\\
\cref{ch:roadmap}, status and roadmap & 18 & $\circ$ & & $\bullet$\\
\cref{app:rules}, reference card & 12 & & $\circ$ & $\circ$\\
\bottomrule
\end{tabular}
\end{table}


% ==================================================================================================
\section{Conventions}\label{sec:intro-conventions}

\paragraph{Boxes.} Three kinds of box interrupt the text, and each can be skipped on a first reading.
A \textbf{Tutorial} box walks through a computation step by step: a trace of unification, of a $\forall$ or
$\imp$ step, of the changes of a state (see the box in \cref{sec:intro-brief}). A \textbf{Kanren angle} box
says how a concept maps to \micro, what runs there today, and what the proposal changes. A \textbf{Pitfall}
box marks a trap: one the book stresses, or one that a Scheme programmer is likely to fall into.

\paragraph{Code.} Three environments, told apart by the color of the rule on their left:
\begin{lprolog}
typeof (abs M) (arr A B) :- pi x\ typeof x A => typeof (M x) B.   % λProlog
\end{lprolog}
\begin{scheme}
(°->list/datum (fresh° (q) (typeof° `(abs ,(ƛ (x) x)) q)))       ; aux Scheme
\end{scheme}
\begin{answers}
((arr α α))
\end{answers}
\noindent The last one, \texttt{answers}, is always the \emph{actual output} of a run of the Scheme code
above it, pasted verbatim, never written by hand. Inline, \lp{pi x\ G} is λProlog and \scm{(∀° (x) g)} is
Scheme. Answers come in three forms. \scm{°->list/datum} gives plain data, comparable with \scm{equal?}.
\scm{°->list #f} gives \micro's own form, a closed λ-expression whose parameters are the answer's unbound
variables, such as \scm{(λ () (ƛ (x) x))}. \scm{°->pretty}, from \texttt{(aux kanren micro show)}, gives
the αKanren-style view with constraints, such as \scm{(α : (= …))}. Reified variables print as \scm{α}, \scm{β}, \scm{γ}; a working variable that a
constraint keeps prints as \scm{_3}. Names follow \aux: goals end in \scm{°}, their applicative forms in
\scm{ª}, term constructors carry no suffix (\scm{ƛ}, \scm{@}), and internals carry the prefix
\scm{μkanren-}.

\paragraph{Pages of the book.} \book{n} refers to the \emph{printed} page $n$ of~\cite{MillerNadathur2012}, the
number at the top of the page, which is the PDF page minus 16.
\begin{pitfall}[Two page numberings]
The comments in the prototype's files (\texttt{examples.scm}, the test suite, the module's header) cite the
book as \texttt{p.N} with the \emph{PDF} page. So \texttt{p.207} in \texttt{examples.scm}, next to the
\texttt{typeof} queries, is \book{191} here.
\end{pitfall}

\paragraph{Notation.} A sequent $\seq{\Sig}{\Prog}{G}$ asks for a proof of goal $G$ from program $\Prog$ over
signature $\Sig$. Implication is written $\imp$, as in the book, and $D\imp G$ is \lp{D => G}. The logics are
\fohc{} and \fohh{} (first-order Horn and hereditary Harrop) and \hohc{} and \hohh{} (their higher-order
versions). $\Llam$ is the pattern fragment, $\eqb$ is λ-convertibility, $\subst{t}{s}{x}$ substitutes $s$ for $x$
in $t$, and $\doteq$ is an equation to be solved. $\ell(X)$ and $\ell(c)$ are the levels of
\cref{def:intro-levels}. Logic variables are upper case ($F$, $X$, $R$), eigenvariables lower case ($c$,
$x$, $y$). Rules of proof search are set in small capitals (\rulename{generic}, \rulename{augment},
\rulename{backchain}).

\paragraph{Where the outputs come from.} {\emergencystretch=2em Snippets in Parts~I and~II, and in \cref{ch:micro}, run on today's
\micro, at commit \texttt{f53bdcb} of the \texttt{aux-\allowbreak lambdakanren} branch, unless a box says otherwise; some
define a small throwaway interpreter first (\cref{sec:hh-interp}). The outputs of the proposal, in this
chapter and from \cref{ch:design} on, come from the prototype as described in Part~III: \micro{} with the
protocol applied (\texttt{micro.diff}), the show module with \texttt{show.diff}, and
\texttt{(aux kanren lambda)}, all in \texttt{doc/\allowbreak hol-kanren/\allowbreak prototype/}. Every snippet was run by
\texttt{csi} from CHICKEN~6.0.1pre1 (rev \texttt{55615603}), as installed; the
benchmarks were compiled with \mbox{\texttt{csc -O2 -d1}}.\marginnote{The document was typeset with XeLaTeX, through Tectonic
0.17.0, in Libertinus and DejaVu Sans Mono.}\par}

% ==================================================================================================
\section{Summary}\label{sec:intro-summary}
\begin{itemize}
  \item The goal is to bring λProlog's universal and hypothetical goals, λ-tree syntax and higher-order pattern
        unification into \micro, following Miller and Nadathur's book (Part~I) and the miniKanren literature,
        Ma and Friedman 2021 first (Part~II).
  \item The outcome (Part~III) is a generic extension protocol in \micro{} plus the client module
        \texttt{(aux kanren lambda)}. Three prototypes were compared by three reviews, and three adversarial
        rounds hardened the result: 89 λ cases and \micro's 125 cases, run with and without the module,
        pass, and the protocol costs about 1\% on first-order search.
  \item The running example is \lp{typeof}. \scm{∀°} gives the body an eigenvariable one level deeper, and
        \scm{⊃°} assumes its type. Backchaining on that hypothesis unifies the types, and the assumption
        disappears with the answer.
  \item Run backwards, the same program needs $R\,c\doteq c$: a pattern, solved by $R:=\abs{x}{x}$, while
        $R\doteq c$ fails the scope check. Non-patterns wait and are reported as residual equations.
  \item Three reading paths: a fast path of about 45 pages before the design and the module, the λProlog
        path (Part~I), and the implementer path.
  \item Answers are actual outputs; \book{n} is a printed page; the prototype's comments use PDF pages
        (printed $+16$).
\end{itemize}


\part{The book: Programming with Higher-Order Logic}
\chapter{Logic programming as proof search}\label{ch:horn}

\Cref{ch:intro} gave the overall plan of this document: read Miller and Nadathur's \emph{Programming with
Higher-Order Logic}~\cite{MillerNadathur2012} with a miniKanren programmer's eyes, and find out what
\micro{} would need in order to host its ideas. This chapter lays the first-order foundation. It covers the
book's Introduction and its Chapters~1 and~2, plus the parts of the Teyjus appendix needed to run the
examples. Nothing here goes beyond Prolog in expressive power. What is new, compared with the usual
Prolog textbook, is the \emph{point of view}. A logic program is a set of formulas, a query is a formula,
and running the query means \emph{searching for a proof} in a sequent calculus. Each logical connective
has a fixed search meaning, independent of the program. Atomic goals are solved by \emph{backchaining}
on program clauses. A language counts as a logic programming language when this goal-directed search
is \emph{complete}.

For a miniKanren programmer the punchline comes early. Pure miniKanren, with \scm{=°},
\scm{fresh°}, \scm{cond°} and recursive relations, \emph{is} first-order Horn-clause logic programming, except for its
search strategy and its lack of types. This chapter makes that statement precise, so that later
chapters can say exactly what Horn clauses \emph{lack}: local constants and local clauses
(\cref{ch:hh}), and $\lambda$-terms as data (\cref{ch:lambda,ch:lambdatree}).

\section{Computation as deduction}\label{sec:horn-deduction}

The Introduction distinguishes two ways of using logic in computer science \book{1--3}. In
\emph{computation-as-model}, a computation is a mathematical structure (states, transitions), and
logic talks \emph{about} it from outside, as Hoare triples or temporal modalities do. In
\emph{computation-as-deduction}, the logical objects themselves (formulas, terms, proofs) \emph{are} the
elements of the computation. This second view splits in two again:
\begin{itemize}
  \item \textbf{proof normalization}: the state is a proof (a $\lambda$-term), and computing means reducing
    it to normal form. This is the basis of functional programming;
  \item \textbf{proof search}: the state is a \emph{sequent}, that is, a goal together with its assumptions,
    and computing means searching for a derivation of it. The way sequents change during the search
    \emph{is} the dynamics of the computation. This is the basis of logic programming.
\end{itemize}

Proof search in logic programming is not the proof search of mathematicians. Mathematicians use
\emph{lemmas}, and the proof rule for lemmas is the \emph{cut rule}: to prove $B$, state a lemma $C$, then
prove $C$ and prove $B$ from $C$. Gentzen's cut-elimination theorem~\cite{Gentzen1935} says that cut can
always be eliminated, for classical and for intuitionistic logic. Cut-free proofs are huge and of
little use to a mathematician, but they are exactly what a machine can look for. So in logic
programming ``the cut rule is excluded from the execution of logic programs, and computation is based on
the search for cut-free proofs''; cut and cut elimination remain useful to \emph{reason about} programs
\book{3}. We come back to this in \cref{sec:horn-logic}.

\paragraph{Programs, goals, sequents.} A formula is either part of a \emph{logic program}, which gives
meaning to the non-logical constants it mentions, or a \emph{goal}, which we try to derive from the
program \book{3}. The state of the search is a collection of sequents
\[
  \seq{\Sig}{\Prog}{G}
\]
where $\Sig$ is a \emph{signature} (the typed non-logical constants available), $\Prog$ is the program and
$G$ is the goal. The Introduction states the principle that drives the whole book \book{3--4}:
\begin{itemize}
  \item if $G$ is \emph{not atomic}, the search commits to decomposing its top-level connective. The
    \emph{search semantics} of a connective is fixed and does not depend on the program;
  \item if $G$ is \emph{atomic}, the program is consulted. The search \emph{backchains}: it looks in
    $\Prog$ for an implication whose conclusion matches the atom, and then tries to prove the premise.
\end{itemize}
``Logic programming can be seen abstractly as a logical framework in which a strategy that alternates
between goal reduction and backchaining is complete'' \book{4}.

\paragraph{How can a sequent change?} The Introduction then asks a qualitative question \book{4}: if
proving $\seq{\Sig}{\Prog}{A}$ leads to proving $\seq{\Sig'}{\Prog'}{A'}$, what can differ between the two?
For \emph{Horn clauses}, $\Sig'=\Sig$ and $\Prog'=\Prog$: signature and program are global and fixed, and
only the atom changes. For \emph{hereditary Harrop formulas}, $\Sig'\supseteq\Sig$ and
$\Prog'\supseteq\Prog$: the search can add constants and clauses that are visible only in part of the
computation, which gives modular programming and data abstraction. Since more of the dynamics then
comes from the logic and not from the particular predicates, ``the metatheory of the logic can be of
more value in proving properties of those programs'' \book{4}. This chapter treats the first case; the
second is \cref{ch:hh}.

\section{Typed first-order terms}\label{sec:horn-terms}

Before writing relations we need data. Book Chapter~1 fixes the language of \emph{terms}. It has no
predicates and no programs yet. It uses actual \lp{λProlog} syntax and is the $\lambda$-free special
case of what \cref{ch:lambda} generalises.

\subsection{Sorts, type constructors and kinds}

λProlog has built-in \emph{sorts} (atomic types) \book{10--11}: \lp{int}, \lp{real}, \lp{string},
\lp{in_stream}, \lp{out_stream}, and \lp{o}, the type of formulas, which has a special status
(\cref{sec:horn-formulas}). It also has the built-in \emph{type constructor} \lp{list}, so \lp{(list int)}
is a type. Type constructors are classified by \emph{kinds}:
\[
  \langle\text{kind}\rangle ::= \typ{type} \mid \typ{type}\to\langle\text{kind}\rangle .
\]
A kind is therefore just an arity: \lp{int : type}, \lp{list : type -> type}. New type constructors are
declared with \lp{kind} \book{12}:
\begin{lprolog}
kind  pair  type -> type -> type.
\end{lprolog}
after which \lp{(pair string int)} is a type. Type expressions are type variables (uppercase tokens),
arrows $\tau_1\to\tau_2$ (written \lp{->}, right associative), and type constructors applied to exactly as
many arguments as their arity \book{12--13}.

\begin{definition}[Target type, order {\book{13}}]\label{def:horn-order}
Every type can be written uniquely as $\tau=\tau_1\to\cdots\to\tau_n\to\tau_0$ with $\tau_0$ not an arrow.
$\tau_0$ is the \emph{target type} and $\tau_1,\dots,\tau_n$ are the \emph{argument types}; $\tau$ is
\emph{functional} when $n\ge1$. The \emph{order} of a type is
\[
  \operatorname{ord}(\tau)=0 \text{ if $\tau$ is not functional},\qquad
  \operatorname{ord}(\tau_1\to\tau_2)=\max(\operatorname{ord}(\tau_1)+1,\ \operatorname{ord}(\tau_2)).
\]
\end{definition}
The order counts how deeply the arrow is nested \emph{to the left}. For example \lp{int} and
\lp{list (int -> int)} have order~0 (the order does not look inside constructor arguments),
\lp{int -> string -> o} has order~1, and \lp{int -> (int -> o) -> o} has order~2 \book{14}. In the
first-order language every declared constant has a type of order at most~1 that does not mention
\lp{o}, every variable has a type of order~0, and type variables may only be instantiated with order-0
types \book{13--16}. The last condition is needed because order is not stable under type substitution:
\lp{A -> list A -> o} has order~1, but instantiating \lp{A} with \lp{int -> int} gives a type of
order~2 \book{13--14}.

\subsection{Terms, polymorphism and typing}

All symbols are \emph{constants}; a constant of functional type is a \emph{value constructor}. Lists use
the polymorphic constants \lp{nil : list A} and \lp{:: : A -> list A -> list A}, where \lp{::} is infix and
right associative, so \lp{2 :: 1 :: nil} has type \lp{list int} \book{15}. Constants are declared with
\lp{type} \book{16}:
\begin{lprolog}
type  pr  A -> B -> pair A B.
\end{lprolog}
A crucial asymmetry governs polymorphism \book{17}:
\begin{itemize}
  \item each \emph{occurrence of a constant} may take its own instance of the declared type. In
    \lp{(2 :: 1 :: nil) :: (1 :: nil) :: nil} the first two \lp{nil}s have type \lp{list int} and the last
    one has type \lp{list (list int)};
  \item a \emph{variable} has one type at all its occurrences. So \lp{((X :: nil) :: X :: nil)} cannot be
    typed, because $X$ would need both a type $\tau$ and the type \lp{list} $\tau$.
\end{itemize}
The rules are those of the type assignment calculus of the book's Figure~1.2 \book{18}, with judgments
$\Sig;\Gamma\Vdash_f t:\tau$, where $\Sig$ types the constants and $\Gamma$ the variables:
\begin{mathpar}
  \inferrule{c:\sigma\in\Sig \\ \tau\vartriangleleft_f\sigma}{\Sig;\Gamma\Vdash_f c:\tau}
  \and
  \inferrule{x:\tau\in\Gamma}{\Sig;\Gamma\Vdash_f x:\tau}
  \and
  \inferrule{\Sig;\Gamma\Vdash_f g:\tau_1\to\tau_2 \\ \Sig;\Gamma\Vdash_f t:\tau_1}{\Sig;\Gamma\Vdash_f (g\ t):\tau_2}
\end{mathpar}
Here $\tau\vartriangleleft_f\sigma$ (``$\tau$ is an instance of $\sigma$'') holds when $\tau$ is obtained from
$\sigma$ by substituting order-0 types other than $\typ{o}$ for type variables. A \emph{first-order term} is a
$t$ with $\Sig;\Gamma\Vdash_f t:\tau$ for some $\tau$ of order~0 different from $\typ{o}$ \book{18}. The
first rule is where constants are polymorphic: the instance is chosen per occurrence. The second rule is
where variables are monomorphic: $\Gamma$ gives one type per variable.

The book states a principal typing property \book{18--19}: if a term has some typing, it has a most general
one, unique up to renaming of type variables. For example \lp{(X :: nil) :: nil} can be typed with
$X:\typ{int}$ or $X:\typ{list}\ \typ{int}$, and most generally with $X:A$, the whole term having type
\lp{list (list A)}. Application is left associative. In canonical form a term is a variable, a constant,
or $(f\ t_1\ \dots\ t_n)$ whose \emph{head} $f$ is always a constant \book{19}: first-order variables cannot
be applied.

\subsection{Representing data}\label{sec:horn-data}

The recipe for representing a class of objects is \book{19}: pick a type for it; declare one value
constructor per way of building an object, all with that common target type; let recursive components
have the target type as argument type. Binary trees \book{20}:
\begin{lprolog}
kind  btree  type -> type.
type  empty  btree A.
type  node   A -> btree A -> btree A -> btree A.
\end{lprolog}
\lp{(node 3 (node 1 empty empty) (node 4 empty empty))} has type \lp{btree int}, while a tree that mixes
\lp{3} and \lp{"cat"} is ill typed. This declaration is \emph{parametric}: every type variable of an
argument type also occurs in the target type. λProlog does not require this; the declaration
\lp{type node A -> btree -> btree -> btree.} gives heterogeneous trees \book{21}. Unlike ML datatypes, the
set of constructors of a type is \emph{open} (declarations can be added later), and terms can contain
variables, so a term describes a whole class of objects \book{21--22}.

The same recipe encodes the abstract syntax of a language. First-order formulas \book{22--23}:
\begin{lprolog}
kind  term, form  type.
type  a, b  term.
type  f     term -> term -> term.
type  var   string -> term.
% f(a, f(x, b))   is represented as   (f a (f (var "x") b))

type  p  term -> term -> form.
type  q  term -> form.

type  ff,        % encoding the false proposition
      tt         form.     % encoding the true proposition
type  &&,        % encoding conjunction
      !!,        % encoding disjunction
      ==>        form -> form -> form. % encoding implication
infixl  &&   5.
infixl  !!   4.
infixr  ==>  3.
\end{lprolog}
With the precedences above, \lp{(p a b) && (q a) !! (q (f a b)) ==> (p b (f b a))} stands for
$((p(a,b)\wedge q(a))\vee q(f(a,b)))\imp p(b,f(b,a))$.

The declarations \lp{infixl}, \lp{infixr} (and \lp{infix}, \lp{prefix}, \lp{postfix}, \dots) are
\emph{operator declarations} $\langle\text{fixity}\rangle\ c_1,\dots,c_n\ \langle\text{precedence}\rangle$:
they change only the concrete syntax, a higher number binding more tightly. For instance, after
\lp{infixl pr 5.} the expression \lp{(3 pr 4 pr "three")} stands for \lp{((pr ((pr 3) 4)) "three")}
\book{16}.

Imperative programs are encoded in the same way \book{24--25}. The book deliberately uses only two sorts,
\lp{expr} and \lp{stmt}, and one \emph{type-neutral} sort for all expressions, whether integer or boolean
valued; object-level typing is left to a later analysis written as a logic program:
\begin{lprolog}
kind stmt, expr  type.
type id     string -> expr.
type c      int -> expr.
type t, f   expr.
type &&, !!,
     plus, minus,
     mult, <    expr -> expr -> expr.
type :=     expr -> expr -> stmt.        % for assignment
type cond   expr -> stmt -> stmt -> stmt. % for conditionals
type while  expr -> stmt -> stmt.        % for while loops
type seq    stmt -> stmt -> stmt.        % for composition
infixl &&, mult   5.
infixl !!, plus, minus  4.
infix  <    3.
infix  :=   2.
\end{lprolog}
The C-like program \texttt{v = 1; i = n; while (0 < i) \{ v = v * i; i = i - 1; \}} is then the term
\begin{lprolog}
(seq ((id "v") := (c 1))
  (seq ((id "i") := (id "n"))
    (while ((c 0) < (id "i"))
      (seq ((id "v") := (id "v") mult (id "i"))
           ((id "i") := (id "i") minus (c 1))))))
\end{lprolog}
Because \lp{:=} has the lowest precedence, \lp{(id "v") := (id "v") mult (id "i")} parses as
\lp{(:= (id "v") (mult (id "v") (id "i")))}. The term \emph{is} the program's abstract syntax tree:
``first-order terms are a generalization of the labeled trees that are used in depicting abstract
syntax'' \book{25}. Types do not express every well-formedness condition (the left side of \lp{:=} must
encode an identifier); such conditions are checked by logic programs over the representation.

\begin{pitfall}[First-order syntax does not capture binding]
Adding \lp{type all term -> form -> form.} lets us write $\forall x\,(p(a,x)\wedge q(x))$ as
\lp{(all (var "x") ((p a (var "x")) && (q (var "x"))))}. The term has the right tree shape, but nothing
in it says that \lp{(var "x")} is \emph{bound} by \lp{all}: renaming, capture-avoiding substitution and
scoping must be programmed by hand in every computation over the syntax. ``This is a fundamental
limitation of first-order approaches to the treatment of syntax'' \book{23}. Representing object
variables by meta-level logic variables instead makes object scoping uncontrollable \book{22}. The book
removes the limitation with $\lambda$-tree syntax (\cref{ch:lambdatree}).
\end{pitfall}

\section{First-order unification}\label{sec:horn-unif}

Unification is the engine under backchaining. The book presents it as a nondeterministic rewriting of a
\emph{multiset of equations} \book{26--29}, the Martelli--Montanari style. A miniKanren programmer knows
the algorithm; the point of this section is the precise vocabulary (mgu, solved form, occurs check) and
the role of types.

\begin{remark}
The book is deliberately ``not a formal development of logic'' \book{6--7}: it states its results in
running prose, without numbered definitions or theorems. The definition, theorem and proposition
environments in this chapter are our packaging of those statements, each with the page where the book
makes it.
\end{remark}

Before the definitions, note the three ways a program uses unification \book{26--27}. Take
$X:\typ{int}$, $L:\typ{list}\ \typ{int}$.
\begin{itemize}
  \item \emph{Decomposition}: $\{(X\mathbin{:\!:}L)=(1\mathbin{:\!:}2\mathbin{:\!:}\cst{nil})\}$ has the unifier
    $\{\langle X,1\rangle,\langle L,2\mathbin{:\!:}\cst{nil}\rangle\}$: the left side is a \emph{pattern} that
    extracts head and tail.
  \item \emph{Recognition}: the same pattern unifies only with non-empty lists, and with $T$ repeated
    the pattern $(\cst{node}\ El\ T\ T)$ recognises the non-empty trees whose two subtrees are
    identical.
  \item \emph{Construction}: in
    \[\{(X\mathbin{:\!:}L_1)=(1\mathbin{:\!:}\cst{nil}),\ L_2=(2\mathbin{:\!:}\cst{nil}),\
    (X\mathbin{:\!:}L_2)=L_3\}\]
    the first equation takes an input apart, the last builds a new term from its
    pieces, and the ``output'' variable $L_3$ receives $1\mathbin{:\!:}2\mathbin{:\!:}\cst{nil}$.
\end{itemize}
The construction case is how a clause head such as \lp{append (X :: L1) L2 (X :: L3)} computes.

\begin{definition}[Unification problem, unifier, mgu {\book{26--28}}]\label{def:horn-mgu}
A \emph{unification problem} is a finite multiset of equations $t=s$ between first-order terms, the two
sides of each equation having the same type. A \emph{unifier} is a type-preserving substitution that
makes the two sides of every equation identical. A substitution $\theta$ is \emph{more general} than
$\varphi$ if $\varphi$ is obtained from $\theta$ by further substitution. A \emph{most general unifier} (mgu)
is a unifier more general than every other.
\end{definition}
Substitutions are \emph{constraints} on the values of variables. $\{\langle X,Y\rangle\}$ says only that
$X$ and $Y$ are equal; $\{\langle X,1\rangle,\langle Y,1\rangle\}$ satisfies that constraint but says more
\book{27}. A good unifier is as general as possible, because later steps of the computation will add
more constraints, and an over-specific choice could make them fail.

\begin{theorem}[mgu {\book{28}}]\label{thm:horn-mgu}
Typed first-order unification problems have most general unifiers whenever they have unifiers, and these
are unique up to substitutions that rename variables.
\end{theorem}
For instance the problem $\{(X\mathbin{:\!:}L)=(Y\mathbin{:\!:}Z\mathbin{:\!:}\cst{nil})\}$ has the two mgu's
\[
\{\langle X,Y\rangle,\langle L,Z\mathbin{:\!:}\cst{nil}\rangle\}
\quad\text{and}\quad
\{\langle Y,X\rangle,\langle L,Z\mathbin{:\!:}\cst{nil}\rangle\},
\]
which differ by the renaming $X\leftrightarrow Y$, and infinitely many less general ground unifiers \book{27--28}.

\paragraph{The transformation system.} The book gives three transformations, \emph{term reduction},
\emph{reorientation} and \emph{variable elimination} \book{28}. We split them into six rules
(this split is ours, the content is the book's). $E\uplus\{\dots\}$ is multiset union, $\bot$ is a problem
with no unifier, $\mathcal V$ the set of variables, and $f^\sigma$ means an occurrence of $f$ at type $\sigma$.
\begin{mathpar}
  \inferrule{E\uplus\{f^{\sigma}\,t_1\ldots t_n = f^{\sigma}\,s_1\ldots s_n\}}{E\uplus\{t_1=s_1,\ldots,t_n=s_n\}}\;\rulename{Decompose}
  \and
  \inferrule{E\uplus\{f^{\sigma}\,\vec t = g^{\sigma'}\,\vec s\}}{\bot}\;\rulename{Clash}\ (f\neq g\text{ or }\sigma\neq\sigma')
  \\
  \inferrule{E\uplus\{t=x\}}{E\uplus\{x=t\}}\;\rulename{Orient}\ (t\notin\mathcal V)
  \and
  \inferrule{E\uplus\{x=x\}}{E}\;\rulename{Delete}
  \and
  \inferrule{E\uplus\{x=t\}}{\bot}\;\rulename{Check}\ (x\in\operatorname{fv}(t),\ t\neq x)
  \\
  \inferrule{E\uplus\{x=t\}}{\subst{E}{t}{x}\uplus\{x=t\}}\;\rulename{Eliminate}\ (x\notin\operatorname{fv}(t),\ x\in\operatorname{fv}(E))
\end{mathpar}
\rulename{Decompose} and \rulename{Clash} are term reduction. They are justified because a \emph{rigid} term
$(f\ t_1\dots t_n)$ keeps its head under any substitution: different heads, \emph{or the same name at two
different types}, can never be made equal \book{28--29}. \rulename{Orient} is bookkeeping.
\rulename{Delete}, \rulename{Check} and \rulename{Eliminate} are variable elimination, and they rest on the
finiteness of terms: $x=t$ with $x$ strictly inside $t$ has no finite solution. This failure is the
\emph{occurs check}. The side condition of \rulename{Clash} is stated here for types without type
variables; with polymorphic constants, ``$\sigma\neq\sigma'$'' becomes ``$\sigma$ and $\sigma'$ have no common
instance'', and \rulename{Decompose} first unifies $\sigma$ with $\sigma'$ as generally as possible
\book{30--31}.

\begin{definition}[Solved form {\book{29}}]
$E$ is in \emph{solved form} if every equation has a variable on the left, and each such variable occurs
nowhere else in $E$. A solved form is read as a substitution, and it is its own mgu.
\end{definition}
\begin{theorem}[Termination and correctness {\book{29}}]
If variable elimination is never applied twice to the same unchanged equation, every sequence of
transformations terminates in $\bot$ or in a solved form, and every transformation preserves the set of
unifiers. The transformations thus define a nondeterministic algorithm for computing mgu's.
\end{theorem}

\begin{tutorial}[Three unification traces {\book{29--30}}]
\emph{Success.} The repeated variable $T$ makes the pattern recognise trees with two identical subtrees.
Write $n_2$ for $(\cst{node}\ 2\ \cst{empty}\ \cst{empty})$:
\begin{align*}
 & \{(\cst{node}\ El\ T\ T) = (\cst{node}\ 1\ n_2\ n_2)\}\\
 \longrightarrow\ & \{El=1,\ T=n_2,\ T=n_2\} && \text{term reduction}\\
 \longrightarrow\ & \{El=1,\ T=n_2,\ n_2=n_2\} && \text{variable elimination}\\
 \longrightarrow\ & \{El=1,\ T=n_2,\ 2=2,\ \cst{empty}=\cst{empty},\ \cst{empty}=\cst{empty}\} && \text{term reduction}\\
 \longrightarrow\ & \{El=1,\ T=n_2\} && \text{term reduction}\times 3
\end{align*}
The last step removes equations between constants: \rulename{Decompose} with $n=0$ arguments. The
result is a solved form, hence the mgu.

\emph{Constant clash.} If the right subtree is $(\cst{node}\ 3\ \cst{empty}\ \cst{empty})$ instead, the same
steps reach the equation $2=3$, and term reduction yields $\bot$.

\emph{Occurs check.} $\{T=(\cst{bt}\ 1\ T\ T)\}$ fails at once, because $T$ occurs in the right-hand side.
\end{tutorial}

\paragraph{Types in unification.} Types enter only in term reduction, when two heads have the same name:
their types must agree \book{30--31}. With polymorphic constants the type of an occurrence is not
determined by its name, so types cannot simply be checked once in advance. The book's example: with
$c:A\to A\to i$ and $a:i$, the problem $\{(c\ 1\ Y)=(c\ X\ a)\}$ is well typed (both sides have type $i$),
but in $(c\ 1\ Y)$ the constant $c$ occurs at type $\typ{int}\to\typ{int}\to i$ and in $(c\ X\ a)$ at
$i\to i\to i$. Ignoring types yields the \emph{ill-typed} unifier $\{\langle X,1\rangle,\langle Y,a\rangle\}$
\book{31}. The book stops there; completing the example, term reduction with types sees two occurrences
of $c$ at the types $\typ{int}\to\typ{int}\to i$ and $i\to i\to i$, which have no common instance, and
fails: the typed problem has no unifier. Types are
themselves first-order terms, so the same algorithm unifies them. The book concludes \book{31} that types
can be erased safely when the original problem's typing guarantees that the intermediate problems stay
well typed, in particular when every type variable in a constructor's argument types also appears in its
target type (the parametric condition of \cref{sec:horn-data}).

\paragraph{Logical reading and the occurs-check policy.} Constants behave like universally quantified
variables and logic variables like existentially quantified ones, so a unification problem is a formula
$\forall x_1\ldots\forall x_n\,\exists y_1\ldots\exists y_m\,[t_1=s_1\wedge\cdots\wedge t_p=s_p]$ \book{31--32}.
Later chapters mix the quantifier prefix (\cref{ch:hh,ch:hounif}). Many Prologs omit the occurs check for
speed; the book does not: ``we will require failure in any case where the occurs-check test is
satisfied'', because ``finiteness of terms, enforced through this test, will be important to encoding
quantifier dependencies'' \book{32}. Without it, deduction can be unsound \book{33}.

\begin{kanren}[Unification in \micro]
\micro's unifier already is the unifier of this section, minus types. Terms are Scheme data, so pairs,
vectors and records decompose structurally (\rulename{Decompose}), and different atoms or different shapes
fail (\rulename{Clash}). The substitution is \emph{triangular}: binding $x\mapsto t$ never rewrites the
other equations, and walking a term on demand plays the role of \rulename{Eliminate}. The occurs check is
on. Here are the book's problems, run on today's \micro{} (file \texttt{unify.scm}; each \scm{q} collects
the variables of interest, and each result is printed after its label):
\begin{scheme}
; U1  {(X :: L) = (1 :: 2 :: nil)}
(μkanren-run (q -1 #t) (fresh° (x l) (=° q (list x l)) (=° (cons x l) '(1 2))))
; U2  {(node El T T) = (node 1 (node 2 empty empty) (node 2 empty empty))}
(μkanren-run (q -1 #t)
  (fresh° (el t) (=° q (list el t))
    (=° `(node ,el ,t ,t) '(node 1 (node 2 empty empty) (node 2 empty empty)))))
; U3  the same with a 3 in the right subtree: constant clash
(μkanren-run (q -1 #t)
  (fresh° (el t) (=° q (list el t))
    (=° `(node ,el ,t ,t) '(node 1 (node 2 empty empty) (node 3 empty empty)))))
; U4  {(X :: L1) = (1 :: nil), L2 = (2 :: nil), (X :: L2) = L3}: construction
(μkanren-run (q -1 #t)
  (fresh° (x l1 l2) (=° (cons x l1) '(1)) (=° l2 '(2)) (=° (cons x l2) q)))
; U5  {(X :: L) = (Y :: Z :: nil)}: an mgu, shown up to renaming
(μkanren-run (q -1 #f)
  (fresh° (x l y z) (=° q (list x l y z)) (=° (cons x l) (list y z))))
; U6  {T = (bt 1 T T)}: occurs check
(μkanren-run (q -1 #t) (=° q `(bt 1 ,q ,q)))
; U7  {(c 1 Y) = (c X a)}: untyped unification returns the ill-typed unifier
(μkanren-run (q -1 #t) (fresh° (x y) (=° q (list x y)) (=° `(c 1 ,y) `(c ,x a))))
\end{scheme}
\begin{answers}
U1 ((1 (2)))
U2 ((1 (node 2 empty empty)))
U3 ()
U4 ((1 2))
U5 ((λ (α β) (cons α (cons (cons β (quote ())) (cons α (cons β (quote ())))))))
U6 ()
U7 ((1 a))
\end{answers}
In U4, \scm{q} plays the output variable $L_3$. U5 is run with \scm{#f} as the third argument of \scm{μkanren-run}, which prints the raw reified answer,
a \scm{λ} over the unbound variables, instead of evaluating it as \scm{#t} does; it reads $X=\alpha$, $L=(\beta)$, $Y=\alpha$, $Z=\beta$. Reification picks one
representative of the mgu ``up to renaming'' (\cref{thm:horn-mgu}). U7 is the warning of the book come
true: an untyped engine happily returns the ill-typed unifier. A typed layer would require typed
constant occurrences, or a static check plus erasure for parametric signatures. The proposal stays untyped
(\cref{sec:roadmap-lim-untyped}); simple types are a later work item (\cref{sec:roadmap-later-types}).
\end{kanren}

\section{First-order formulas}\label{sec:horn-formulas}

Book Chapter~2 relaxes the signature: constants may now have target type \lp{o}. They are the
\emph{predicate symbols}, and an expression of type \lp{o} headed by one of them is an \emph{atomic
formula} \book{34--35}:
\begin{lprolog}
type memb    A -> list A -> o.
type append  list A -> list A -> list A -> o.
\end{lprolog}
The logical constants and their λProlog spellings are \book{35}:
\begin{center}
\begin{tabular}{@{}lllll@{}}
\toprule
math & λProlog & type & fixity & conventional use\\
\midrule
$\top$ & \texttt{true} & \texttt{o} & & \\
$\imp$ & \texttt{=>} & \texttt{o -> o -> o} & infix, right assoc. & implication in goals\\
$\wedge$ & \texttt{\&} & \texttt{o -> o -> o} & infix, right assoc. & joins program clauses\\
$\wedge$ & \texttt{,} & \texttt{o -> o -> o} & infix, left assoc. & joins goals\\
$\vee$ & \texttt{;} & \texttt{o -> o -> o} & infix, left assoc. & \\
$\imp$ (reversed) & \texttt{:-} & \texttt{o -> o -> o} & infix, non assoc. & ``is implied by''\\
$\forall_\tau x\,F$ & \texttt{pi x\textbackslash\ F} & & binder & \\
$\exists_\tau x\,F$ & \texttt{sigma x\textbackslash\ F} & & binder & \\
\bottomrule
\end{tabular}
\end{center}
The two spellings of $\wedge$ and of $\imp$ are a stylistic convention \book{47--48}; breaking it does not
affect well-formedness. All these operators bind less tightly than application \book{35}, and \lp{&}
binds more tightly than \lp{:-} \book{48}, so \lp{A & B :- G} reads as $G\imp(A\wedge B)$. Only the
logical constants may take arguments of type \lp{o} in the first-order setting \book{35}. Quantifiers carry an optional type, \lp{(pi (x:T)\ F)}, inferred most generally when
omitted; the scope of \lp{x\ } extends as far to the right as possible, and only the binding pattern
matters, so \lp{pi x\ (p x) => sigma y\ (q x y, pi x\ (q y x))} and
\lp{pi z\ (p z) => sigma y\ (q z y, pi v\ (q y v))} are the same formula \book{36--37}. The typing rules
of the quantifiers (Figure~2.1, \book{38}) extend those of terms:
\begin{mathpar}
  \inferrule{\Sig;\Gamma,x:\tau\Vdash_f B:\typ{o}}{\Sig;\Gamma\Vdash_f \forall_\tau x\,B:\typ{o}}
  \and
  \inferrule{\Sig;\Gamma,x:\tau\Vdash_f B:\typ{o}}{\Sig;\Gamma\Vdash_f \exists_\tau x\,B:\typ{o}}
  \and
  \inferrule{\Sig;\Gamma\Vdash_f B:\tau}{\Sig;\Gamma\Vdash_f C:\tau}\ (B, C \text{ differ only in bound names})
\end{mathpar}
with the proviso, for the first two, that $x$ is not typed by $\Sig$ or $\Gamma$; the third rule renames
bound variables so that the proviso can be met. A \emph{first-order $\Sig$-formula} is an $F$ with
$\Sig;\emptyset\Vdash_f F:\typ{o}$; such formulas are closed \book{38}.

\section{The search semantics of goals}\label{sec:horn-search}

A logic programming language is fixed by four ingredients \book{38}: signatures, program clauses, goals, and
a calculus for building proofs. A sequent $\seq{\Sig}{\Prog}{G}$ has a set $\Prog$ of $\Sig$-formulas as its
program, and it states the judgment ``$G$ is provable from $\Prog$ and $\Sig$''.

\begin{definition}[Search semantics of the logical symbols {\book{39}}]\label{def:horn-reduction}
Read as steps of an interpreter:
\begin{description}[font=\normalfont\scshape]
  \item[and] reduce $\seq{\Sig}{\Prog}{B_1\wedge B_2}$ to both $\seq{\Sig}{\Prog}{B_1}$ and $\seq{\Sig}{\Prog}{B_2}$;
  \item[or] reduce $\seq{\Sig}{\Prog}{B_1\vee B_2}$ to either $\seq{\Sig}{\Prog}{B_1}$ or $\seq{\Sig}{\Prog}{B_2}$;
  \item[instan] reduce $\seq{\Sig}{\Prog}{\exists_\tau x\,B}$ to $\seq{\Sig}{\Prog}{\subst{B}{t}{x}}$ for some $\Sig$-term $t$ of type $\tau$;
  \item[augment] reduce $\seq{\Sig}{\Prog}{B_1\imp B_2}$ to $\seq{\Sig}{\Prog,B_1}{B_2}$;
  \item[generic] reduce $\seq{\Sig}{\Prog}{\forall_\tau x\,B}$ to $\seq{c:\tau,\Sig}{\Prog}{\subst{B}{c}{x}}$, where $c$ is a new constant, not in $\Sig$;
  \item[true] $\seq{\Sig}{\Prog}{\top}$ is provable at once.
\end{description}
\end{definition}
None of these looks at $\Prog$. Read downwards, they are the \emph{right-introduction rules} of a sequent
calculus, the book's Figure~2.2 \book{40}:
\begin{mathpar}
  \inferrule{ }{\seq{\Sig}{\Prog}{\top}}\;\top R
  \and
  \inferrule{\seq{\Sig}{\Prog}{B_1} \\ \seq{\Sig}{\Prog}{B_2}}{\seq{\Sig}{\Prog}{B_1\wedge B_2}}\;\wedge R
  \and
  \inferrule{\seq{\Sig}{\Prog}{B_1}}{\seq{\Sig}{\Prog}{B_1\vee B_2}}\;\vee R
  \and
  \inferrule{\seq{\Sig}{\Prog}{B_2}}{\seq{\Sig}{\Prog}{B_1\vee B_2}}\;\vee R
  \\
  \inferrule{\seq{\Sig}{\Prog,B_1}{B_2}}{\seq{\Sig}{\Prog}{B_1\imp B_2}}\;{\imp}R
  \and
  \inferrule{\seq{\Sig}{\Prog}{\subst{B}{t}{x}} \\ \Sig;\emptyset\Vdash_f t:\tau}{\seq{\Sig}{\Prog}{\exists_\tau x\,B}}\;\exists R
  \and
  \inferrule{\seq{c:\tau,\Sig}{\Prog}{\subst{B}{c}{x}}}{\seq{\Sig}{\Prog}{\forall_\tau x\,B}}\;\forall R
\end{mathpar}
($c$ not declared in $\Sig$ in $\forall R$). Each rule is sound in ``any interesting logic''. The question is
\emph{completeness} \book{41}: if a sequent with a non-atomic goal is provable, must the premise(s) of the
corresponding rule be provable too? The answer depends on which formulas are allowed and on the logic.

\begin{example}[Why goal-directedness can fail {\book{41--42}}]\label{ex:horn-counter}
Let $\Sig=\{p:\typ{o},q:\typ{o},r:i\to\typ{o},a:i,b:i\}$.
\begin{enumerate}
  \item $\seq{\Sig}{p\vee q}{q\vee p}$ is provable, but \rulename{or} leaves $\seq{\Sig}{p\vee q}{q}$ or
    $\seq{\Sig}{p\vee q}{p}$, and neither is provable. A disjunction in the \emph{program} breaks
    goal-directed search.
  \item $\seq{\Sig}{\emptyset}{p\vee(p\imp q)}$ is a classical tautology (excluded middle), but \rulename{or}
    leaves $p$, unprovable, or $p\imp q$, which \rulename{augment} reduces to the unprovable $\seq{\Sig}{p}{q}$.
  \item $\seq{\Sig}{(r\,a\wedge r\,b)\imp q}{\exists x\,(r\,x\imp q)}$ is classically provable by cases on
    $r\,a$, but every choice of witness $t$ leaves an unprovable sequent.
\end{enumerate}
\end{example}
The book concludes \book{42} that programs must not contain disjunctions (they must be \emph{definite},
``because they do not contain indefinite information''), and that the logic should be
\emph{intuitionistic}, where excluded middle does not hold. Counterexamples~2 and~3 need $\imp$ in the
\emph{goal}, which Horn clauses forbid; the switch of logic pays off only from \cref{ch:hh} on.

Atomic goals need rules that look at the program; the general framework leaves them open on purpose
\book{42}. When the search succeeds, the system reports an \emph{answer substitution}: the free variables
of a query are read as existentially quantified at its head, and the answer maps them to terms that make
the computation succeed \book{42--43}.

\section{Horn clauses and backchaining}\label{sec:horn-fohc}

\begin{definition}[fohc {\book{43}}]\label{def:horn-fohc}
With $A$ ranging over atomic formulas, the goals $G$ and program clauses $D$ of \emph{first-order Horn
clauses} (\fohc) are
\[
  G ::= \top \mid A \mid G\wedge G \mid G\vee G \mid \exists_\tau x\,G
  \qquad
  D ::= A \mid G\imp D \mid D\wedge D \mid \forall_\tau x\,D .
\]
\end{definition}
Only $\imp$ and $\forall$ are missing at the top of goals, and only $\vee$ and $\exists$ at the top of
clauses. The family contains the familiar Prolog shape $\forall x_1\ldots\forall x_m\,(A_1\wedge\cdots\wedge
A_n\imp A_0)$, but is more liberal: the head may sit under further implications, conjunctions and
universals \book{43--44}.

Backchaining is first explained informally \book{44}: if $\Prog$ contains $A$, succeed; if it contains $G'\imp
A$, reduce the goal $A$ to $G'$. Formally it uses a \emph{focused} sequent $\Sig;\Prog\xrightarrow{D}A$, ``try
to prove the atom $A$ by backchaining on $D$''. The rules of the book's Figure~2.3 \book{45} are
\begin{mathpar}
  \inferrule{\Sig;\Prog\xrightarrow{D}A}{\seq{\Sig}{\Prog}{A}}\;\rulename{decide}\ (D\in\Prog)
  \and
  \inferrule{ }{\Sig;\Prog\xrightarrow{A}A}\;\rulename{initial}
  \and
  \inferrule{\Sig;\Prog\xrightarrow{D}A \\ \seq{\Sig}{\Prog}{G}}{\Sig;\Prog\xrightarrow{G\imp D}A}\;{\imp}L
  \\
  \inferrule{\Sig;\Prog\xrightarrow{D_1}A}{\Sig;\Prog\xrightarrow{D_1\wedge D_2}A}\;\wedge L
  \and
  \inferrule{\Sig;\Prog\xrightarrow{D_2}A}{\Sig;\Prog\xrightarrow{D_1\wedge D_2}A}\;\wedge L
  \and
  \inferrule{\Sig;\Prog\xrightarrow{\subst{D}{t}{x}}A \\ \Sig;\emptyset\Vdash_f t:\tau}{\Sig;\Prog\xrightarrow{\forall_\tau x\,D}A}\;\forall L
\end{mathpar}
\rulename{decide} picks a clause; \rulename{initial} closes the branch when the focused atom is
\emph{identical} to the goal; ${\imp}L$ proves the body $G$ and continues on $D$; $\wedge L$ picks a
conjunct; $\forall L$ instantiates. The formula over the arrow is on the left of the sequent, in focus, so
these are \emph{left}-introduction rules complementing Figure~2.2. For a Prolog-shaped clause
$D=\forall\bar x\,(A_1\wedge\cdots\wedge A_n\imp A_0)$ the rules compose into one step \book{45}:
\begin{mathpar}
  \inferrule{\seq{\Sig}{\Prog}{A_1\theta} \\ \cdots \\ \seq{\Sig}{\Prog}{A_n\theta}}{\Sig;\Prog\xrightarrow{D}A}
\end{mathpar}
where $\theta$ maps each $x_i$ to a $\Sig$-term of its type and $A=A_0\theta$: the instance
$A_1\theta\wedge\cdots\wedge A_n\theta\imp A_0\theta$ follows from $\Prog$, and modus ponens gives $A_0\theta$.
This grouping of small rules into large ones is a first instance of Andreoli's \emph{focusing}
\cite{Andreoli1992}, which the book cites as the general setting \book{73}.

\begin{definition}[O-proof {\book{45}}]
An \emph{O-proof} is a proof built with the rules of the book's Figures~1.2, 2.2 and~2.3.
\end{definition}
\begin{theorem}[Completeness of goal-directed search for fohc {\book{45}}]\label{thm:horn-oproof}
If $\Prog$ is an \fohc{} program and $G$ an \fohc{} goal, then $\seq{\Sig}{\Prog}{G}$ has an O-proof iff it is
provable in classical logic iff it is provable in intuitionistic logic.
\end{theorem}
\begin{proposition}[Flatness {\book{46}}]\label{prop:horn-flat}
Every sequent in an O-proof of $\seq{\Sig}{\Prog}{G}$ has the form $\seq{\Sig}{\Prog}{G'}$ or
$\Sig;\Prog\xrightarrow{D}A$, with the \emph{same} $\Sig$ and $\Prog$.
\end{proposition}
So in \fohc{} the program and signature are flat and global, and \rulename{augment} and \rulename{generic}
never fire. These two rules are exactly what \cref{ch:hh} adds. The book's term for a logic in which this
kind of goal-directed (``uniform'') search is complete is \emph{abstract logic programming
language}~\cite{MillerNadathurPfenningScedrov1991,MillerNadathur2012} \book{73}; \fohc{} is the base case.

\begin{tutorial}[An O-proof for a query with two answers]
Take the program $\Prog=\{D_1,D_2\}$ of \lp{append} (\cref{sec:horn-syntax}), with
$D_1=\forall l\,(\cst{app}\ \cst{nil}\ l\ l)$ and $D_2=\forall x\forall l_1\forall l_2\forall l_3\,(\cst{app}\ l_1\
l_2\ l_3\imp\cst{app}\ (x\mathbin{:\!:}l_1)\ l_2\ (x\mathbin{:\!:}l_3))$, and the goal $\exists X\exists Y\,\cst{app}\ X\ Y\
(1\mathbin{:\!:}\cst{nil})$. Instead of guessing terms in $\exists R$ and $\forall L$ we use logic variables and let
\rulename{initial} unify (this is the practical version of the rules, \cref{sec:horn-pragmatics}).
\begin{enumerate}
  \item $\exists R$ twice: logic variables $X,Y$; goal $\cst{app}\ X\ Y\ (1\mathbin{:\!:}\cst{nil})$, atomic.
  \item \rulename{decide} $D_1$; $\forall L$ with a new $L$; \rulename{initial} unifies
    $\cst{app}\ \cst{nil}\ L\ L$ with the goal: $X=\cst{nil}$, $L=Y$, $Y=1\mathbin{:\!:}\cst{nil}$.
    \emph{First answer}: $X=\cst{nil}$, $Y=1\mathbin{:\!:}\cst{nil}$.
  \item Back to step 2, \rulename{decide} $D_2$ instead. $\forall L$ four times with new $X',L_1,L_2,L_3$,
    then ${\imp}L$ splits into a focused premise and a goal premise. \rulename{initial} on the focused
    premise unifies $\cst{app}\ (X'\mathbin{:\!:}L_1)\ L_2\ (X'\mathbin{:\!:}L_3)$ with the goal: $X=X'\mathbin{:\!:}L_1$, $Y=L_2$,
    $X'=1$, $L_3=\cst{nil}$. The goal premise is now $\cst{app}\ L_1\ Y\ \cst{nil}$.
  \item For $\cst{app}\ L_1\ Y\ \cst{nil}$: \rulename{decide} $D_1$ succeeds with $L_1=\cst{nil}$,
    $Y=\cst{nil}$. \emph{Second answer}: $X=1\mathbin{:\!:}\cst{nil}$, $Y=\cst{nil}$. \rulename{decide} $D_2$ fails:
    its head needs the third argument to be $x\mathbin{:\!:}l_3$, which clashes with $\cst{nil}$.
\end{enumerate}
The O-proof for the second answer, with the unifier already applied. Write $A_0$ for
$\cst{app}\ (1\mathbin{:\!:}\cst{nil})\ \cst{nil}\ (1\mathbin{:\!:}\cst{nil})$ and $A_1$ for $\cst{app}\ \cst{nil}\ \cst{nil}\ \cst{nil}$:
\begin{mathpar}
\inferrule*[right=$\exists R\times 2$]{
 \inferrule*[right=decide]{
  \inferrule*[right=$\forall L\times 4$]{
   \inferrule*[right=${\imp}L$]{
     \inferrule*[right=initial]{ }{\Sig;\Prog\xrightarrow{A_0}A_0}
     \\
     \inferrule*[right=decide$+\forall L+$initial]{ }{\seq{\Sig}{\Prog}{A_1}}
   }{\Sig;\Prog\xrightarrow{A_1\imp A_0}A_0}
  }{\Sig;\Prog\xrightarrow{D_2}A_0}
 }{\seq{\Sig}{\Prog}{A_0}}
}{\seq{\Sig}{\Prog}{\exists X\exists Y\,\cst{app}\ X\ Y\ (1\mathbin{:\!:}\cst{nil})}}
\end{mathpar}
Every sequent in it has the same $\Sig$ and $\Prog$, as \cref{prop:horn-flat} promises.
\end{tutorial}

\section{λProlog concrete syntax and interaction}\label{sec:horn-syntax}

In program clauses, an unbound token beginning with an uppercase letter is a variable universally
quantified over the whole clause; \lp{_} is an anonymous variable, each occurrence distinct \book{46--47}.
So \book{47}
\begin{lprolog}
type append  list A -> list A -> list A -> o.
append nil L L.
append (X :: L1) L2 (X :: L3) :- append L1 L2 L3.
\end{lprolog}
abbreviates \lp{pi L\ append nil L L.} and
\lp{pi X\ pi L1\ pi L2\ pi L3\ append (X :: L1) L2 (X :: L3) :- append L1 L2 L3.}; the types of the
variables are inferred (\lp{X : A}, the lists \lp{list A}). Anonymous variables save names for values
that are never used again \book{47}:
\begin{lprolog}
type sublist  list A -> list A -> o.
sublist L K :- append _ T K, append L _ T.
\end{lprolog}
is the same clause as \lp{sublist L K :- append U T K, append L V T.}: $L$ is a sublist of $K$ if $K$
has a suffix $T$ of which $L$ is a prefix.

\paragraph{Equivalent groupings of clauses.} Because the D-grammar allows $\wedge$ and $\imp$ to nest, one
predicate can be written in several equivalent ways \book{48--49}. With
\lp{kind bool type.}, constructors \lp{neg}, \lp{and}, \lp{or}, \lp{imp} and \lp{type ident bool -> bool -> o.},
the four clauses of Form~1 can be compressed into a single clause, Form~3:
\begin{lprolog}
ident (neg B) (neg D) :- ident B D.
ident (and B C) (and D E) :- ident B D, ident C E.
ident (or  B C) (or  D E) :- ident B D, ident C E.
ident (imp B C) (imp D E) :- ident B D, ident C E.
\end{lprolog}
\begin{lprolog}
ident (neg B) (neg D) &
 (ident (and B C) (and D E) &
  ident (or  B C) (or  D E) &
  ident (imp B C) (imp D E) :- ident C E) :- ident B D.
\end{lprolog}
How does Form~3 parse and run? Since \lp{&} binds more tightly than \lp{:-}, it is the single D-formula
\[
  G_1\imp\bigl(D_{\cst{neg}}\wedge(G_2\imp(D_{\cst{and}}\wedge D_{\cst{or}}\wedge D_{\cst{imp}}))\bigr)
\]
(universally closed), where $G_1$ stands for the goal \lp{ident B D}, $G_2$ for
\lp{ident C E}, and $D_{\cst{neg}},\dots$ for the four heads. It is legal because the D-grammar allows
$G\imp D$ and $D\wedge D$ to nest. To solve an atom \lp{ident (and B0 C0) X}, \rulename{decide} picks this
clause, $\forall L$ instantiates its variables, ${\imp}L$ sets the body \lp{ident B D} aside as a goal
premise and keeps backchaining on the conjunction, $\wedge L$ chooses the right conjunct, ${\imp}L$ sets
\lp{ident C E} aside, $\wedge L$ chooses $D_{\cst{and}}$, and \rulename{initial} unifies it with the goal.
The two goals set aside are exactly the body of the second clause of Form~1. The book claims that the
forms are equivalent in derivable goals and in observable behaviour, and that a smart implementation
preprocesses Form~3 into Form~1 \book{49}. (Figure~2.3 does not say in which order the two premises of
${\imp}L$ are explored; for the behaviour to match Form~1, whose body is \lp{ident B D, ident C E}, the
outer premise must be solved first. This remark is ours.) (As printed, \lp{ident} has no base case,
so it proves no closed goal; the example is about clause structure only.)

\paragraph{The read--prove--print loop.} Programs live in modules (\cref{sec:hoprog-modules}). Teyjus compiles a
module with \texttt{tjcc}, links it with \texttt{tjlink} and runs it with \texttt{tjsim}, which shows the
prompt \lp{[lists] ?-} \book{281--282}. In \emph{queries} the convention for variables is dual: an
uppercase token is \emph{existentially} quantified, with outermost scope \book{50}. The prove phase looks
for an O-proof, with three possible outcomes: a proof, no proof, or non-termination. Non-termination cannot
be ruled out by any interpreter, because provability in \fohc{} is undecidable \book{51}. The print phase
shows the bindings of the named query variables only \book{51--53}:
\begin{lprolog}
[lists] ?- append (1::nil) (2::nil) (3::nil).
no
[lists] ?- append (1::nil) (2::nil) (1::2::nil).
solved
[lists] ?- append (1::nil) (2::nil) X.
X = (1::2::nil)
[lists] ?- append X Y (1::nil).
X = nil
Y = 1::nil;
X = 1::nil
Y = nil;
no
\end{lprolog}
Typing \lp{;} asks for another proof. The queries \lp{sigma X\ append (1::nil) (2::nil) X} and
\lp{append (1::nil) (2::nil) _} attempt the same sequent as the third query, but print just
\lp{solved}: \lp{sigma}-bound and anonymous variables are not shown \book{51--52}.

\begin{kanren}[The same session in \micro]
A \fohc{} predicate becomes a \scm{define-relation} whose body is a \scm{cond°} with one line per clause:
\scm{fresh°} for the clause variables, \scm{=°} for head unification, then the body goals. File
\texttt{append.scm}:
\begin{scheme}
;; append nil L L.
;; append (X :: L1) L2 (X :: L3) :- append L1 L2 L3.
(define-relation (append° l k m)
  (cond° ((=° l '()) (=° k m))
         ((fresh° (x l1 l3)
            (=° l (cons x l1)) (=° m (cons x l3))
            (append° l1 k l3)))))

(μkanren-run (q -1 #t) (append° '(1) '(2) '(3)))
(μkanren-run (q -1 #t) (append° '(1) '(2) '(1 2)))
(μkanren-run (q -1 #t) (append° '(1) '(2) q))
(μkanren-run (q -1 #t) (fresh° (x y) (=° q (list x y)) (append° x y '(1))))
\end{scheme}
\begin{answers}
append (1) (2) (3)    => ()
append (1) (2) (1 2)  => (α)
append (1) (2) X      => ((1 2))
append X Y (1)        => ((() (1)) ((1) ()))
\end{answers}
\scm{()} is \lp{no}. The answer \scm{(α)} is \lp{solved}: one proof, with the query variable left
unconstrained. The run argument \scm{-1} asks for all answers, what typing \lp{;} repeatedly does in
Teyjus. The variables of an inner \scm{fresh°} are hidden like \lp{sigma}-bound ones.
\end{kanren}

\section{Core examples}\label{sec:horn-examples}

The book illustrates three programming styles \book{53--59}: search problems, relations over recursive
data, and relations over abstract syntax.

\subsection{Search: finite-state machines}\label{sec:horn-fsm}

The book's Figure~2.4 \book{54} describes acceptance by a finite-state machine whose arcs are labelled by
\emph{words}:
\begin{lprolog}
kind state, letter        type.
type q1, q2, q3, q4, q5   state.
type a, b                 letter.
type start, final         state -> o.
type path, trans          state -> list letter -> state -> o.
type accept               list letter -> o.

path S nil S.
path S Letters T :- trans S Arc M, append Arc Rest Letters,
                    path M Rest T.
accept W :- start S, path S W F, final F.
\end{lprolog}
The goal \lp{append Arc Rest Letters} consumes the arc's word as a prefix of the input. Figure~2.5 is a
nondeterministic machine for $(a+b)^*(ab+ba)$, the words ending in $ab$ or $ba$:
\begin{lprolog}
start q1 & final q2 & final q3.
trans q1 (a::nil) q1      & trans q1 (b::nil) q1.
trans q1 (a::b::nil) q2   & trans q1 (b::a::nil) q3.
\end{lprolog}
To enumerate accepted words, the book first generates lists of increasing length \book{55}:
\begin{lprolog}
type lists  list A -> o.
lists nil.
lists (_::L) :- lists L.
\end{lprolog}
and reports \book{55--56}
\begin{lprolog}
[fsm1] ?- accept (b::b::a::b::nil).
solved
[fsm1] ?- accept (b::a::X::Y::nil).
X = a
Y = b ;
X = b
Y = a ;
no
[fsm1] ?- lists L.
L = nil ;
L = T1 :: nil ;
L = T1 :: T2 :: nil ;
L = T1 :: T2 :: T3 :: nil
[fsm1] ?- lists L, accept L.
L = a :: b :: nil ;
L = b :: a :: nil ;
L = a :: a :: b :: nil ;
L = a :: b :: a :: nil ;
L = b :: a :: b :: nil
\end{lprolog}
Of the last two queries the book notes that ``both these queries have an unbounded number of successful
proofs''; a depth-first engine ``will systematically produce lists of increasing length'' \book{55--56}.
Note the answer \lp{L = T1 :: nil}: an answer substitution may contain variables, and then every instance
of it is an answer. Why is the generator needed? Under
depth-first search, \lp{accept W} with \lp{W} unbound tries the first \lp{trans} clause
\lp{trans q1 (a::nil) q1} first; \lp{append} then binds \lp{W = a :: Rest}, and \lp{path q1 Rest F} is
in the same situation again, forever. No answer is ever printed. \lp{lists L} breaks the loop by giving
\lp{accept} lists of fixed length, one length at a time.

\begin{pitfall}[An erratum in Figure~2.6]
The book also gives a deterministic five-state machine (Figure~2.6, \book{54}), printed as
\begin{lprolog}
start q1 & final q4 & final q5.
trans q1 (a::nil) q2   &   trans q1 (b::nil) q3.
trans q2 (a::nil) q1   &   trans q2 (b::nil) q4.
trans q3 (a::nil) q5   &   trans q3 (b::nil) q1.
trans q4 (a::nil) q5   &   trans q4 (b::nil) q3.
trans q5 (a::nil) q2   &   trans q5 (b::nil) q4.
\end{lprolog}
and says that exactly the same behaviour would be observed with it \book{56}. As printed it does not
accept the stated language: the
clauses \lp{trans q2 (a::nil) q1} and \lp{trans q3 (b::nil) q1} send the machine back to the start, so it
rejects \lp{aab}, \lp{bba}, and so on. Replacing them by \lp{trans q2 (a::nil) q2} and
\lp{trans q3 (b::nil) q3} fixes it. A brute-force comparison with the regular expression on all words of
length up to~8 confirms both the defect and the fix.
\end{pitfall}

\subsection{Recursive data: insertion into a search tree}

Relations over recursive data follow the structure of the data \book{57}:
\begin{lprolog}
type insert  int -> btree int -> btree int -> o.
insert X empty (node X empty empty).
insert X (node A L R) (node A NL R) :- X < A,  insert X L NL.
insert X (node A L R) (node A L NR) :- X >= A, insert X R NR.
\end{lprolog}
\begin{lprolog}
[btree] ?- insert 4 (node 3 (node 2 empty empty) empty) T.
T = node 3 (node 2 empty empty) (node 4 empty empty)
\end{lprolog}
The comparisons \lp{<} and \lp{>=} are built-ins ``correctly used only when both arguments are
instantiated to numbers'' \book{49}, so \lp{insert} is well defined only when its first two arguments are
closed terms \book{57}. This is a \emph{mode} restriction: the relation is not usable in every direction.

\subsection{Abstract syntax: a sequent prover}\label{sec:horn-prv}

The third style uses \fohc{} to animate an inference system. The book encodes a classical propositional
sequent calculus with multisets $\Gamma\longrightarrow\Delta$ on both sides (Figure~2.7, \book{58}), over the
formulas of \cref{sec:horn-data}:
\begin{mathpar}
  \inferrule{ }{\Gamma,\bot\longrightarrow\Delta}\;\bot L
  \and
  \inferrule{ }{\Gamma,A\longrightarrow\Delta,A}\;\rulename{initial}
  \\
  \inferrule{\Gamma,A\longrightarrow\Delta \\ \Gamma,B\longrightarrow\Delta}{\Gamma,A\vee B\longrightarrow\Delta}\;\vee L
  \and
  \inferrule{\Gamma\longrightarrow\Delta,A,B}{\Gamma\longrightarrow\Delta,A\vee B}\;\vee R
  \\
  \inferrule{\Gamma,A,B\longrightarrow\Delta}{\Gamma,A\wedge B\longrightarrow\Delta}\;\wedge L
  \and
  \inferrule{\Gamma\longrightarrow\Delta,A \\ \Gamma\longrightarrow\Delta,B}{\Gamma\longrightarrow\Delta,A\wedge B}\;\wedge R
  \\
  \inferrule{\Gamma\longrightarrow\Delta,A \\ \Gamma,B\longrightarrow\Delta}{\Gamma,A\imp B\longrightarrow\Delta}\;{\imp}L
  \and
  \inferrule{\Gamma,A\longrightarrow\Delta,B}{\Gamma\longrightarrow\Delta,A\imp B}\;{\imp}R
\end{mathpar}
Several formulas on the right make the calculus classical. Here $\Gamma,\Delta,A,B$ are \emph{schema}
variables of the object calculus, not logic variables of λProlog. Multisets are lists, and a rule matches
``$\Gamma, P$'' by selecting $P$ and keeping the rest:
\begin{lprolog}
type memb_and_rest  A -> list A -> list A -> o.
memb_and_rest X (X :: L) L.
memb_and_rest X (Y :: L) (Y :: L1) :- memb_and_rest X L L1.
\end{lprolog}
The translation recipe is mechanical \book{59}: \emph{a rule's conclusion becomes the head of a clause, its
premises become the body, and the schema variables become clause variables}. The eight clauses, one per
rule in the order $\bot L$, \rulename{initial}, $\wedge L$, $\vee L$, ${\imp}L$, $\wedge R$, $\vee R$,
${\imp}R$, as printed:
\begin{lprolog}
type prv  list form -> list form -> o.

prv Gamma Delta :- memb_and_rest ff Gamma _.
prv Gamma Delta :- memb_and_rest A Gamma _,
                   memb_and_rest A Delta _.
prv Gamma Delta :-
      memb_and_rest (A && B) Gamma Gamma',
      prv (A :: B :: Gamma') Delta.
prv Gamma Delta :-
      memb_and_rest (A !! B) Gamma Gamma',
      prv (A :: Gamma') Delta, prv (B :: Gamma') Delta.
prv Gamma Delta :-
      memb_and_rest (A ==> B) Gamma Gamma',
      prv Gamma' (A :: Delta), prv (B :: Gamma') Delta.
prv Gamma Delta :-
      memb_and_rest (A && B) Delta Delta',
      prv Gamma (A :: Delta'), prv Gamma (B :: Delta').
prv Gamma Delta :-
      memb_and_rest (A !! B) Delta Delta',
      prv Gamma (A :: B :: Delta').
prv Gamma Delta :-
      memb_and_rest (A ==> B) Delta Delta',
      prv (A :: Delta) (B :: Gamma').
\end{lprolog}
Read the $\vee L$ clause against its rule: the conclusion $\Gamma,A\vee B\longrightarrow\Delta$ becomes the
head \lp{prv Gamma Delta} together with the selection \lp{memb_and_rest (A !! B) Gamma Gamma'} (so the
clause's \lp{Gamma'} is the rule's $\Gamma$), and the two premises become the two body goals. The
\rulename{initial} clause needs no selection of the rest: it only asks that some $A$ be on both sides.
\begin{lprolog}
[logic] ?- prv nil (((p a b) !! ((p a b) ==> (q a a))) :: nil).
solved
\end{lprolog}
Note the \emph{level shift}. This is $p\vee(p\imp q)$, the formula that \cref{ex:horn-counter} showed
goal-directed search cannot prove (it is not even an \fohc{} goal, since it contains $\imp$). Here it is
object-level \emph{data}, proved by an object-level classical calculus that is \emph{programmed} in
\fohc: the search proves the meta-level atom \lp{prv nil (...)}, by $\vee R$ and then ${\imp}R$ and
\rulename{initial} on \lp{p a b}.

\begin{pitfall}[An erratum in the last \texttt{prv} clause]
The clause for ${\imp}R$ is printed \book{59} with the body \lp{prv (A :: Delta) (B :: Gamma')}. It should be
\lp{prv (A :: Gamma) (B :: Delta')}. As printed, \lp{Gamma'} is a singleton variable, the right-hand side
becomes an open list, and the initial clause can unify a hypothesis into that open tail. The printed
program then ``proves'' $p\,a\,b\imp q\,a\,a$, which is not valid. A \micro{} port of the printed clause
confirms it: with the printed clause, the query below for \lp{(p a b) ==> (q a a)} returns \scm{(α)}
instead of \scm{()} (file \texttt{prv-printed.scm}).
\end{pitfall}

\begin{kanren}[The prover in \micro]
The port keeps the clause-per-line shape (file \texttt{prv.scm}, last clause corrected). Formulas are
lists such as \scm{(!! A B)}:
\begin{scheme}
(define-relation (memb-and-rest° x l r)
  (cond° ((=° l (cons x r)))
         ((fresh° (y l1 r1)
            (=° l (cons y l1)) (=° r (cons y r1)) (memb-and-rest° x l1 r1)))))

(define-relation (prv° Γ Δ)
  (cond°
    ((fresh° (_) (memb-and-rest° 'ff Γ _)))                                   ; ⊥L
    ((fresh° (a _1 _2) (memb-and-rest° a Γ _1) (memb-and-rest° a Δ _2)))      ; initial
    ((fresh° (a b Γ*) (memb-and-rest° `(&& ,a ,b) Γ Γ*)                       ; ∧L
       (prv° `(,a ,b . ,Γ*) Δ)))
    ((fresh° (a b Γ*) (memb-and-rest° `(!! ,a ,b) Γ Γ*)                       ; ∨L
       (prv° (cons a Γ*) Δ) (prv° (cons b Γ*) Δ)))
    ((fresh° (a b Γ*) (memb-and-rest° `(==> ,a ,b) Γ Γ*)                      ; ⊃L
       (prv° Γ* (cons a Δ)) (prv° (cons b Γ*) Δ)))
    ((fresh° (a b Δ*) (memb-and-rest° `(&& ,a ,b) Δ Δ*)                       ; ∧R
       (prv° Γ (cons a Δ*)) (prv° Γ (cons b Δ*))))
    ((fresh° (a b Δ*) (memb-and-rest° `(!! ,a ,b) Δ Δ*)                       ; ∨R
       (prv° Γ `(,a ,b . ,Δ*))))
    ((fresh° (a b Δ*) (memb-and-rest° `(==> ,a ,b) Δ Δ*)                      ; ⊃R
       (prv° (cons a Γ) (cons b Δ*))))))

(define pab '(p a b))
(define qaa '(q a a))
(μkanren-run (q 1 #t) (prv° '() `((!! ,pab (==> ,pab ,qaa)))))
(μkanren-run (q 1 #t) (prv° '() `((==> ,pab ,qaa))))
\end{scheme}
\begin{answers}
prv nil (p a b !! (p a b ==> q a a))   => (α)
prv nil (p a b ==> q a a)              => ()
\end{answers}
\end{kanren}

\section{Pragmatics of the search}\label{sec:horn-pragmatics}

The rules of \cref{sec:horn-search,sec:horn-fohc} are nondeterministic. λProlog resolves the choices with
a ``simple and rigid search strategy'' so that a programmer can predict behaviour and resource usage
\book{60}:
\begin{enumerate}
  \item for $\vee R$ and $\wedge L$, the left subformula is tried before the right one;
  \item for $\wedge R$, the left premise is derived first;
  \item for \rulename{decide}, $\Prog$ is a \emph{list}, not a set: clauses are tried in textual order, and
    order and multiplicity matter;
  \item the terms of $\exists R$ and $\forall L$ are not enumerated (most types are infinite) but replaced
    by \emph{logic variables}, instantiated by unification \book{61}.
\end{enumerate}
The price is \emph{incompleteness}: ``there may be sequents that have derivations but for which no
derivations will be found because of the strategy. In fact, the chosen strategy even can cause indefinite
looping when a more flexible search strategy might be able to find a proof'' \book{60}. Since the strategy
is known, programs can usually be restructured, as \lp{lists L} did for \lp{accept}.

With a logic variable $X$ that is new in the search, the quantifier rules and \rulename{initial} become
\book{61}
\begin{mathpar}
  \inferrule{\Sig;\emptyset\Vdash_f X:\tau \\ \seq{\Sig}{\Prog}{\subst{B}{X}{x}}}{\seq{\Sig}{\Prog}{\exists_\tau x\,B}}
  \and
  \inferrule{\Sig;\Prog\xrightarrow{\subst{D}{X}{x}}A \\ \Sig;\emptyset\Vdash_f X:\tau}{\Sig;\Prog\xrightarrow{\forall_\tau x\,D}A}
  \and
  \inferrule{ }{\Sig;\Prog\xrightarrow{A'}A}\;\rulename{initial}\ (A, A'\text{ unifiable})
\end{mathpar}
where the unifier must be ``percolated all over the derivation that has been constructed thus far'', and
typed unification guarantees that the delayed typing premises hold. Finally the program can be precompiled
\book{61--62}: applying ${\imp}L$ and $\wedge L$ in advance turns it into a list of clauses
$\forall\bar x\,(G_1\wedge\cdots\wedge G_n\imp A)$, which are then partitioned by predicate name, compiled,
and indexed on their arguments, as in the Warren Abstract Machine.

\paragraph{Cut, negation as failure, arithmetic.} Book Chapter~2 has no pruning operator; it only
stresses that the cut \emph{rule} ``has no connection to Prolog's pruning operator !'' \book{63}. Teyjus
provides \lp{!}, \lp{fail} and \lp{not} as metalogical predicates, with the warning that
``negation-by-failure has a problematic semantics in a situation where implicational goals are used in an
essential way'' \book{287}. The book later defines \book{137}
\begin{lprolog}
type not  o -> o.
not P :- P, !, fail.
not P.
\end{lprolog}
and shows how far this is from logic: \lp{?- X = 2, not (1 = X).} succeeds, whereas
\lp{?- not (1 = X), X = 2.} fails, so ``even a basic property such as commutativity of conjunctions becomes
invalid'' \book{137--138} (\cref{ch:hoprog}). Arithmetic is \emph{intensional}: \lp{2 + 3} is a term, so the goal
\lp{(2 + 3 = 5)} fails \book{66}, and only the non-logical predicate \lp{is} evaluates
\book{67,286}. The comparisons used by \lp{insert} are of the same non-logical kind.

\section{Logical notions: cut and presentations of fohc}\label{sec:horn-logic}

\Cref{thm:horn-oproof} gives \fohc{} the duality the book wants: solving a goal means both ``it follows
logically'' and ``a certain derivation exists'' \book{62}. The cut rule \book{63} is
\begin{mathpar}
  \inferrule{\seq{\Sig}{\Prog}{B} \\ \seq{\Sig}{\Prog,B}{G}}{\seq{\Sig}{\Prog}{G}}\;\rulename{cut}
\end{mathpar}
It uses a lemma $B$. Choosing lemmas takes ingenuity, so cut is not used to compute, but it is used to
\emph{reason about} computation. Cut elimination says that adding cut to classical or intuitionistic
logic does not change the derivable sequents, and the completeness of O-proofs means that an adapted cut is
admissible for them. Consequence: replacing a subformula of a program by a logically equivalent one gives a
program that proves the same goals, ``in the sense of O-provability, although maybe not under a
depth-first search strategy'' \book{63}.

This licenses three equivalent presentations of \fohc{} clauses \book{63--64}: the liberal one of
\cref{def:horn-fohc}; the common Prolog one, $F ::= A\mid F\wedge F$, $D ::= A\mid F\imp A\mid\forall_\tau
x\,D$; and a compact one, $D ::= A\mid A\imp D\mid\forall_\tau x\,D$. A clause of one presentation is
equivalent to a set of clauses of another through equivalences such as
\[
\forall x\,(B_1\wedge B_2)\equiv(\forall x\,B_1)\wedge(\forall x\,B_2),\qquad
B_1\imp(B_2\imp B_3)\equiv(B_1\wedge B_2)\imp B_3,
\]
\[
(B_1\vee B_2)\imp B_3\equiv(B_1\imp B_3)\wedge(B_2\imp B_3),\qquad
B_1\imp\forall x\,B_2\equiv\forall x\,(B_1\imp B_2)\ (x\notin\operatorname{fv}(B_1)),
\]
together with distributivity of $\wedge$ and $\vee$ and the others listed on \book{64}.

\begin{example}[Size and behaviour {\book{64--66}}]\label{ex:horn-size}
In the common presentation, $((p\vee r)\wedge(q\vee t))\imp s$ must become four clauses
\[
(p\wedge q)\imp s,\quad (r\wedge q)\imp s,\quad (p\wedge t)\imp s,\quad (r\wedge t)\imp s,
\]
which duplicates $p, q, r, t$, exponentially in general. It also changes behaviour: with the original
clause $p$ is derived at most once per use, with the four clauses it is derived again for each clause
that mentions it. Introducing new predicates $pr, qt$ for $p\vee r$ and $q\vee t$ instead,
\[
p\imp pr,\quad r\imp pr,\quad q\imp qt,\quad t\imp qt,\quad (pr\wedge qt)\imp s,
\]
grows the program only linearly. This is not a logical equivalence (there are new constants), but it is
\emph{conservative}: for every goal $G$ over the original signature $\Sig_1$,
$\seq{\Sig_1}{\Prog,D_1}{G}$ is provable iff $\seq{\Sig_2}{\Prog,D_2}{G}$ is.
\end{example}

\section{The role of types}\label{sec:horn-types}

\paragraph{Types classify expressions.} In the first-order setting the only distinction types
\emph{force} is between propositions (\lp{o}) and everything else \book{66}. Equality is syntactic: ``two
closed expressions are equal only if they are exactly the same'', so \lp{2 + 3}, \lp{3 + 2} and \lp{5}
are different elements of \lp{int}, and \lp{o} is a set of formulas, not of truth values
\book{66--67}.

\paragraph{Polymorphism, transparency, determinacy.} The declaration \lp{type nil list A.} is read as
$\forall A\,(\mathit{type}\ \mathit{nil}\ (\mathit{list}\ A))$: \lp{nil} has many types \book{67}.
\begin{definition}[Transparent, determinate {\book{67--68}}]
For an order-1 type $\tau_1\to\cdots\to\tau_n\to\tau_0$, a type variable occurring in $\tau_0$ is
\emph{transparent}; the type is \emph{determinate} if all its type variables are transparent.
\end{definition}
Knowing the target type of a determinate constructor fixes the types of its arguments: all elements of a
\lp{list A} have the same type. The constructor \lp{type cons A -> lst -> lst.} is not determinate and
builds heterogeneous lists \book{68}.

\paragraph{Checking and inference.} λProlog is strongly typed: types of constants are declared, types of
variables are inferred most generally, and checking is static \book{68--69}. The asymmetry of
\cref{sec:horn-terms} returns: \lp{append (1::nil) (2::nil) X, append ("abc"::nil) ("efg"::nil) Y.} is well
typed, with the two \lp{append} occurrences at different instances, but replacing \lp{Y} by \lp{X} makes it
ill typed, since a variable has one type \book{68}.

\paragraph{Types at run time.} With polymorphism and non-determinate constructors, types cannot always be
compiled away \book{69}. The book's Figure~2.8 builds heterogeneous lists with the non-determinate
\lp{cons} and dispatches on the run-time type of their elements \book{70}:
\begin{lprolog}
kind lst       type.
type null      lst.
type cons      A -> lst -> lst.
type separate  lst -> list int -> list real -> o.
separate (cons (X:int)  L) (X::K) M :- separate L K M.
separate (cons (X:real) L) K (X::M) :- separate L K M.
separate null nil nil.
\end{lprolog}
The annotation \lp{(X:int)} must be matched, during the search, against the type of the head element of
the incoming list: this is run-time type processing. Figure~2.9 achieves the same with
\emph{injections}, so that the dispatch is on \emph{terms} \book{70}:
\begin{lprolog}
kind numb      type.
type inj_int   int  -> numb.
type inj_real  real -> numb.
type separate  list numb -> list int -> list real -> o.
separate ((inj_int  X)::L) (X::K) M :- separate L K M.
separate ((inj_real X)::L) K (X::M) :- separate L K M.
separate nil nil nil.
\end{lprolog}
The book recommends determinate constructors and injections \book{71--72}: static typing then catches more
errors, and less type information has to be carried during proof search. Giving constants the type \lp{A}
would make nonsense well typed. The book poses the queries
\begin{lprolog}
separate (cons 1.0 (cons 2 (cons 3.0 null))) L K
separate ((inj_real 1.0)::(inj_int 2)::(inj_real 3.0)::nil) L K
\end{lprolog}
to the two figures and says that both bind ``K to (2 :: nil) and L to (1.0 :: 3.0 :: nil)'' \book{70}.
That swaps the answers: the second argument collects the integers, so \lp{L = (2 :: nil)} and
\lp{K = (1.0 :: 3.0 :: nil)}. The \micro{} port below confirms it.

\begin{kanren}[Types and \micro]
\micro{} has one universe of Scheme data. The book observes that in the first-order setting the only
distinction types \emph{force} is between \lp{o} and the rest, so all other sorts could be collapsed into
one \book{66}; \micro{} is that collapsed case. What the collapse gives up is static error detection,
typed unification (U7 above) and the run-time dispatch of Figure~2.8. Injections (Figure~2.9) are ordinary
miniKanren style, tagged lists such as \scm{(inj_int 2)}, and port clause by clause (file
\texttt{separate.scm}):
\begin{scheme}
(define-relation (separate° ns is rs)
  (cond° ((fresh° (x l k) (=° ns `((inj_int ,x) . ,l)) (=° is (cons x k)) (separate° l k rs)))
         ((fresh° (x l m) (=° ns `((inj_real ,x) . ,l)) (=° rs (cons x m)) (separate° l is m)))
         ((=° ns '()) (=° is '()) (=° rs '()))))

(μkanren-run (q -1 #t)
  (fresh° (l k) (=° q (list l k))
    (separate° '((inj_real 1.0) (inj_int 2) (inj_real 3.0)) l k)))
\end{scheme}
\begin{answers}
separate (...) L K => (((2) (1.0 3.0)))
\end{answers}
The answer is the pair $(L, K)$: the integers first, as the types demand. Where a type layer could live (a static
checker written as a relation over declared signatures, since type unification is just \scm{=°} on type
terms) is left open: the proposal stays untyped (\cref{sec:roadmap-lim-untyped}); simple types are a later
work item (\cref{sec:roadmap-later-types}).
\end{kanren}

\section{From fohc to \texorpdfstring{\texttt{(aux kanren micro)}}{(aux kanren micro)}}\label{sec:horn-micro}

We can now state precisely how \fohc{} maps onto \micro{} (described in \cref{ch:micro}):

\begin{center}
\small
\begin{tabular}{@{}>{\raggedright\arraybackslash}p{0.27\linewidth}>{\raggedright\arraybackslash}p{0.33\linewidth}>{\raggedright\arraybackslash}p{0.33\linewidth}@{}}
\toprule
\fohc{} (book) & \micro & remark\\
\midrule
$\top$, \texttt{true} ($\top R$) & \texttt{✓°} & \\
$G_1\wedge G_2$, \texttt{,} ($\wedge R$) & \texttt{(and° g1 g2)} & bind interleaves answers\\
$G_1\vee G_2$, \texttt{;} ($\vee R$) & \texttt{(or° g1 g2)}, \texttt{cond°} & interleaving, not left first\\
$\exists x\,G$, \texttt{sigma} ($\exists R$, logic variable) & \texttt{(fresh° (x) g)} & exactly the rule of \book{61}\\
atomic goal, \rulename{decide} & call \texttt{(p° t ...)} & clause choice compiled into the relation\\
clauses of $p$, $\forall\bar x\,(G\imp p\,\bar t)$ & \texttt{define-relation} with one \texttt{cond°} line per clause & precompiled and indexed by name, \book{61--62}\\
$\forall L$ with logic variable & \texttt{fresh°} opening each line & renaming apart is automatic\\
\rulename{initial} with unification & \texttt{=°} against the clause head & the threaded substitution percolates $\theta$\\
occurs check & built in & as the book requires, \book{32}\\
answer substitution & reified answer, \texttt{α}, \texttt{β}, \dots & up to renaming\\
typed unification & none & see U7\\
\texttt{<}, \texttt{>=}, \texttt{is} & \texttt{project°} + Scheme & non-logical, moded\\
\texttt{!} (pruning) & no hard cut; soft cut \texttt{if°}, \texttt{take°} & control, not logic\\
\texttt{not} & \texttt{(if° g ✗° ✓°)} & same unsoundness\\
${\imp}R$ (\rulename{augment}), $\forall R$ (\rulename{generic}) & absent & not needed in \fohc{}; \cref{ch:hh}\\
\bottomrule
\end{tabular}
\end{center}

So pure miniKanren \emph{is} \fohc, with two caveats. First, the shape of a relation (one disjunction of
existentially quantified equations per predicate) is the precompiled form of \book{61--62}, and the
built-in \scm{=°} is definable in \fohc{} by the single clause $\forall x\,(\cst{eq}\ x\ x)$.
Second, a Scheme-hosted miniKanren is strictly \emph{larger} than \fohc: goals are first-class Scheme
values, and \scm{project°} leaves the logic altogether. The \fohc{} fragment is a discipline, not something
the host enforces.

\begin{kanren}[The FSM, side by side with its λProlog original]
File \texttt{fsm.scm} ports Figures~2.4 and~2.5 clause by clause (\scm{append°} as above):
\begin{scheme}
;; Figure 2.4
(define-relation (path° s letters t)
  (cond° ((=° letters '()) (=° s t))
         ((fresh° (arc m rest)
            (trans° s arc m) (append° arc rest letters) (path° m rest t)))))
(define-relation (accept° w)
  (fresh° (s f) (start° s) (path° s w f) (final° f)))

;; Figure 2.5
(define-relation (start° s) (=° s 'q1))
(define-relation (final° s) (cond° ((=° s 'q2)) ((=° s 'q3))))
(define-relation (trans° s arc t)
  (cond° ((=° (list s arc t) '(q1 (a) q1)))
         ((=° (list s arc t) '(q1 (b) q1)))
         ((=° (list s arc t) '(q1 (a b) q2)))
         ((=° (list s arc t) '(q1 (b a) q3)))))

(define-relation (lists° l)
  (cond° ((=° l '())) ((fresh° (h t) (=° l (cons h t)) (lists° t)))))

(μkanren-run (q -1 #t) (accept° '(b b a b)))
(μkanren-run (q -1 #t) (fresh° (x y) (=° q (list x y)) (accept° `(b a ,x ,y))))
(μkanren-run (q 4 #t) (lists° q))
(μkanren-run (q 5 #t) (lists° q) (accept° q))
(μkanren-run (q 5 #t) (accept° q))
\end{scheme}
\begin{answers}
accept (b b a b)         => (α)
accept (b a X Y)         => ((a b) (b a))
lists L          (4)     => (() (α) (α β) (α β γ))
lists L, accept L (5)    => ((a b) (b a) (a a b) (a b a) (b a b))
accept L          (5)    => ((a b) (b a) (a a b) (a b a) (b a b))
\end{answers}
The first four lines agree with the book (\lp{lists L} gives \lp{nil}, \lp{T1 :: nil}, \dots). The last
line is the interesting one: \scm{(accept° q)} \emph{without} the generator produces answers, where
depth-first search loops on the $a$-arc. The order happens to coincide with the book's here; in general
one should not expect \micro{} to reproduce λProlog's answer order.
\end{kanren}

\subsection{Depth-first search versus interleaving}\label{sec:horn-dfs}

λProlog commits to the depth-first, left-to-right, clause-order strategy of \cref{sec:horn-pragmatics}.
Teyjus compiles it to an abstract machine \book{277}. \micro{} keeps the same goals but changes the
strategy, as μKanren does~\cite{HemannFriedman2013}: \scm{or°} interleaves its two answer streams, switching
on each suspension, and a relation call suspends. A clause that recurses forever therefore gets only a
share of the search, and the remaining clauses still run. That is why \scm{(accept° q)} terminates
answer by answer.

What does this buy, logically? \Cref{thm:horn-oproof} is a statement about the proof system: O-proofs exist
whenever proofs exist. A search strategy is complete if it finds a proof whenever one exists, and, for answers, if
every answer substitution of an O-proof is eventually produced (up to instantiation). Depth-first search
is not \book{60}. For miniKanren's interleaving search, Rozplokhas, Vyatkin and
Boulytchev~\cite{RozplokhasVyatkinBoulytchev2020} proved soundness and completeness with respect to a
denotational semantics: every answer is eventually produced. \micro's particular interleaving has not
been checked against the hypotheses of that proof, so read it as a strong expectation for \micro, not a
theorem about it. Interleaving is fair for \emph{disjunction} only. Termination is not
promised: a query with no answers, or the rest of a query after its last answer, can still diverge,
which is unavoidable since \fohc{} provability is undecidable \book{51}. Conjunction remains
order-sensitive. For example
\begin{scheme}
(μkanren-run (q 1 #t) (lists° q) (=° q 'not-a-list))
\end{scheme}
never returns: \scm{lists°} offers infinitely many lists, and each is rejected by the second conjunct. We
ran it with a 10-second timeout, and it produced no output. With the conjuncts swapped it fails at once.

\begin{pitfall}[Answer order is not part of the logic]
The book stresses that programs are lists, so the order and multiplicity of clauses are observable under
depth-first search \book{60}, and logically equivalent transformations can change the number of
derivations (\cref{ex:horn-size}). Under interleaving the same holds, with a different order. A port of a
book example to \micro{} should compare answer \emph{sets}, unless the order has been checked.
\end{pitfall}

\begin{kanren}[Negation as failure and moded built-ins in \micro]
\micro{} has a soft cut, \scm{(if° g? gt gf)}: if \scm{g?} has answers, all of them continue with \scm{gt},
otherwise \scm{gf} runs. With it, \lp{not} is one line, and the book's commutativity failure \book{137--138}
reproduces exactly (file \texttt{naf.scm}). For contrast, the last two lines use the disequality constraint
\scm{≠°}, which is not part of \fohc{} (\cref{ch:micro}) but is commutative:
\begin{scheme}
(define (not° g) (if° g ✗° ✓°))

(μkanren-run (q -1 #t) (=° q 2) (not° (=° 1 q)))
(μkanren-run (q -1 #t) (not° (=° 1 q)) (=° q 2))
(μkanren-run (q -1 #t) (=° q 2) (≠° 1 q))
(μkanren-run (q -1 #t) (≠° 1 q) (=° q 2))
\end{scheme}
\begin{answers}
X = 2, not (1 = X)   => (2)
not (1 = X), X = 2   => ()
X = 2, X =/= 1       => (2)
X =/= 1, X = 2       => (2)
\end{answers}
Moded built-ins behave as in λProlog. The port of \lp{insert} guards with \scm{project°}, which hands the
current values to Scheme (file \texttt{insert.scm}):
\begin{scheme}
(define-relation (insert° x t t*)
  (cond° ((=° t 'empty) (=° t* `(node ,x empty empty)))
         ((fresh° (a l r nl)
            (=° t `(node ,a ,l ,r)) (=° t* `(node ,a ,nl ,r))
            (project° ((x° x) (a° a)) (if (< x° a°) ✓° ✗°))
            (insert° x l nl)))
         ((fresh° (a l r nr)
            (=° t `(node ,a ,l ,r)) (=° t* `(node ,a ,l ,nr))
            (project° ((x° x) (a° a)) (if (>= x° a°) ✓° ✗°))
            (insert° x r nr)))))

(μkanren-run (q -1 #t) (insert° 4 '(node 3 (node 2 empty empty) empty) q))
(μkanren-run (q -1 #t) (fresh° (x) (insert° x '(node 3 empty empty) q)))   ; mode broken
\end{scheme}
\begin{answers}
insert 4 (node 3 (node 2 empty empty) empty) T => ((node 3 (node 2 empty empty) (node 4 empty empty)))
insert X (node 3 empty empty) T => 
Error: (<) bad argument type - not a number: _1
\end{answers}
The second query breaks the mode: Scheme's \scm{<} receives the unbound logic variable \scm{_1} and the run
aborts with the error shown (printed on standard error). Likewise intensional
arithmetic (file \texttt{arith.scm}):
\begin{scheme}
(μkanren-run (q -1 #t) (=° '(+ 2 3) 5))
(μkanren-run (q -1 #t) (=° q '(+ 2 3)))
(μkanren-run (q -1 #t)
  (fresh° (e) (=° e '(+ 2 3)) (project° ((e* e)) (=° q (apply + (cdr e*))))))
\end{scheme}
\begin{answers}
2 + 3 = 5        => ()
X = 2 + 3        => ((+ 2 3))
X is 2 + 3       => (5)
\end{answers}
\end{kanren}

\begin{kanren}[What \textup{fohc} does not need, and later chapters will]
By \cref{prop:horn-flat}, nothing in this chapter needs the program or the signature to change during
the search, so \micro{} can keep relations as closed Scheme procedures and needs no signature object.
\Cref{ch:hh} adds \rulename{augment} (clauses that are visible only along one branch of the search) and
\rulename{generic} (fresh constants, together with a scoping check in unification that the occurs-check
discussion of \book{32} anticipates). \Cref{ch:lambda,ch:hounif} add $\lambda$-terms as data. The
corresponding requirements on \micro{} are collected in \cref{ch:design}.
\end{kanren}

\section{Summary}\label{sec:horn-summary}
\begin{itemize}
  \item Logic programming is \emph{computation as proof search}: the state is a sequent
    $\seq{\Sig}{\Prog}{G}$, computing means looking for a \emph{cut-free} proof, and cut is used to reason
    about programs, not to run them \book{3,63}.
  \item Data are typed first-order terms. Kinds are arities. Each occurrence of a polymorphic constant
    takes its own instance of the type, but a variable has one type. Principal types exist. First-order
    encodings of syntax do not capture binding \book{10--25}.
  \item First-order unification is a terminating transformation system on multisets of equations. It computes
    mgu's, unique up to renaming. The occurs check is mandatory in the book. Types matter only when two
    heads have the same name, and they can be erased for parametric signatures \book{26--32}.
  \item Connectives in goals have a fixed search semantics, the right rules (\rulename{and}, \rulename{or},
    \rulename{instan}, \rulename{augment}, \rulename{generic}, \rulename{true}); atoms are solved by
    \emph{backchaining}, the focused left rules \rulename{decide}, \rulename{initial}, ${\imp}L$, $\wedge L$,
    $\forall L$ \book{39--45}.
  \item In \fohc{} ($G ::= \top\mid A\mid G\wedge G\mid G\vee G\mid\exists x\,G$; $D ::= A\mid G\imp D\mid D\wedge
    D\mid\forall x\,D$) O-proofs are complete for classical and intuitionistic logic, and the program and
    signature never change during the search \book{43--46}.
  \item λProlog fixes depth-first, left-to-right, clause-order search with logic variables and
    unification. It is predictable but incomplete. \lp{!}, \lp{not}, \lp{is} and comparisons are
    non-logical \book{60--62,137,286}.
  \item Logically equivalent presentations of clauses prove the same goals but may differ in size and in the
    number of derivations; introducing predicates is conservative \book{63--66}.
  \item Types classify expressions, not values; equality is syntactic. Determinate constructors and
    injections are preferred to run-time type dispatch \book{66--72}.
  \item Pure miniKanren is \fohc{}, without types and with an interleaving search. The interleaving search
    finds answers that depth-first search misses (\scm{(accept° q)}), but conjunction order still decides
    termination. \micro{} reproduces the book's examples and its warnings (ill-typed unifier,
    non-commutative \lp{not}, moded built-ins).
\end{itemize}

\chapter{Hereditary Harrop formulas: implication and universal goals}\label{ch:hh}

% local notation: a focused (backchaining) sequent  Σ; P --D--> A
\newcommand{\hhfoc}[4]{#1;#2\xrightarrow{\,#3\,}#4}

In \cref{ch:horn} we met the logic of first-order Horn clauses, \fohc, the logic behind Prolog.
Its search semantics was given by a handful of reduction rules, one per logical symbol, and two
of them never fired: AUGMENT, for goals of the form $D\imp G$, and GENERIC, for goals of the form
$\forall x\,G$. The reason is purely syntactic: in \fohc{} neither goals nor clause bodies may
contain implications or universal quantifiers, so the signature $\Sig$ and the program $\Prog$ stay
fixed for the whole computation. Miller and Nadathur \cite{MillerNadathur2012} observe that such a language ``does not make
very strong use of the structure afforded by logic'' \book{75}.

This chapter lifts the restriction. The logic of \emph{first-order hereditary Harrop formulas},
\fohh, admits $D\imp G$ and $\forall x\,G$ in goals and hence, at any depth, in clause bodies. The
two new rules have a simple operational reading:
\begin{itemize}
  \item to prove $D\imp G$, prove $G$ with $D$ \emph{temporarily added to the program};
  \item to prove $\forall x\,G$, prove $G[c/x]$ with a \emph{new constant} $c$ temporarily added to
        the signature.
\end{itemize}
So ``programs and signatures can grow dynamically in the course of searching for a proof''
\book{75}, and they grow and shrink with a stack discipline. Three things follow, and they organise
the chapter: hypothetical reasoning becomes a programming technique (\cref{sec:hh-augment}); the
interplay of $\forall$ with logic variables imposes a \emph{scope condition} on unification
(\cref{sec:hh-generic,sec:hh-link}); and the declarative reading of the language is no longer
classical logic but intuitionistic, indeed minimal, logic (\cref{sec:hh-logic}). The full power of
$\forall$, binding \emph{names inside data}, ``must await the introduction of higher-order features''
\book{75}; it returns in \cref{ch:lambdatree}. Here $\forall$ shows its scoping face. The last section
(\cref{sec:hh-kanren}) collects what all this demands of a miniKanren, and of \micro{} in
particular.

%=====================================================================================================
\section{Syntax: goals, clauses and polarity}\label{sec:hh-syntax}

\begin{definition}[\fohh{} goals and clauses \book{75}]\label{def:hh-grammar}
Let $A$ range over first-order atomic formulas. Goals $G$ and program clauses $D$ are given by the
mutually recursive rules
\[
\begin{array}{rcl}
G &::=& \top \mid A \mid G\wedge G \mid G\vee G \mid \exists_\tau x\,G \mid D\imp G \mid \forall_\tau x\,G\\[2pt]
D &::=& A \mid G\imp D \mid D\wedge D \mid \forall_\tau x\,D
\end{array}
\]
A $D$-formula is called a \emph{first-order hereditary Harrop formula}; \fohh{} names both this class
of formulas and the logic programming framework built on it \book{76}.
\end{definition}

Compare with \fohc{} (\cref{ch:horn}), where $G::=\top\mid A\mid G\wedge G\mid G\vee G\mid\exists x\,G$
and $D$ has the same grammar. Two remarks make the definition easy to remember.
\begin{itemize}
  \item \emph{The $D$ grammar did not change.} The restrictions on clauses ``are, in fact, identical to
        the ones present in the \fohc{} setting'' \book{76}: no top-level $\vee$ or $\exists$, and the
        premise of a top-level implication is a goal. The class of clauses is larger only because goals,
        which appear as clause bodies, are now richer.
  \item \emph{The two grammars swap roles across $\imp$.} The premise of an implication \emph{goal} is a
        clause (it will be assumed); the premise of an implication \emph{clause} is a goal (it will be
        proved). Goals still are not freely generated: $(p\vee q)\imp r$ is not a goal, because
        $p\vee q$ is not a clause.
\end{itemize}

\paragraph{Concrete syntax.} In λProlog, $D\imp G$ in a goal is written \lp{D => G}; a clause
$G\imp A$ is written \lp{A :- G}; \lp{,} and \lp{&} both denote $\wedge$ (the chapter uses \lp{&}
inside assumed conjunctions of clauses), \lp{;} is $\vee$, \lp{pi x\ B} is $\forall x\,B$ and
\lp{sigma x\ B} is $\exists x\,B$. Free capitalised variables of a clause are implicitly universally
quantified at the front of the clause; free variables of a query are implicitly existential. The scope
of \lp{x\} extends as far to the right as possible.

\begin{pitfall}[the precedence of \texttt{=>}]
The connectives are listed in the order \lp{=>}, \lp{&}, \lp{,}, \lp{;}, \lp{:-}; their precedences
follow that order, so \lp{=>} binds \emph{tightest}; \lp{=>} and \lp{&} are right associative
\book{35}. Hence \lp{D => G1, G2} parses as \lp{(D => G1), G2}, and the hypothesis does \emph{not}
scope over \lp{G2}. This is why the parentheses in the queries of this chapter, such as
\lp{fact (finished X Y) => (fact (graduates X), fact (cs_major X))} and \lp{(p 2 & p 3) => p X}, are
essential, while \lp{p 2 => p 3 => p X} means $p\,2\imp(p\,3\imp p\,X)$.
\end{pitfall}

\subsection{Positive and negative occurrences}

With implications inside formulas it becomes meaningful to ask on which side of an arrow a
subformula sits.

\begin{definition}[polarity \book{76}]\label{def:hh-polarity}
$B$ is a positive subformula occurrence of $B$. If $C$ is a positive (negative) occurrence in $B$, then
it is a positive (negative) occurrence in $B\wedge B'$, $B'\wedge B$, $B\vee B'$, $B'\vee B$,
$B'\imp B$, $\forall x\,B$ and $\exists x\,B$, and it is a negative (positive) occurrence in $B\imp B'$.
Equivalently: $C$ is positive in $B$ if it lies to the left of an even number of implications, negative if
odd.
\end{definition}

\begin{proposition}[\book{76}]\label{prop:hh-polarity}
Positive subformulas of goals are goals and negative subformulas of goals are clauses; dually, positive
subformulas of clauses are clauses and negative ones are goals. A \fohc{} goal has no negative
subformulas. A hereditary Harrop formula is exactly a formula in which no positive subformula occurrence
is a disjunction or an existential quantification.
\end{proposition}

The last characterisation explains the name. The negative occurrences are the places where a formula
is \emph{used as a hypothesis}. Requiring that no positive part of a clause be a disjunction or an
existential says that whatever we assume, now or later in the search, never forces us to reason by cases
or to invent an anonymous witness. The adjective \emph{hereditary} refers to this
``at every positive position'' requirement (\cref{sec:hh-logic}).

\begin{example}
In the clause $\forall J\,((\forall x\,(\cst{bug}\,x\imp\cst{in}\,x\,J\imp\cst{dead}\,x))\imp
\cst{sterile}\,J)$ (\cref{sec:hh-generic}) the body $\forall x(\dots)$ is negative, hence a goal, and the
subformulas $\cst{bug}\,x$ and $\cst{in}\,x\,J$ are positive again (each lies to the left of exactly two
implications): they are clauses, and indeed they will be \emph{assumed}.
\end{example}

\subsection{Clausal order}

\begin{definition}[clausal order \book{76--77}]\label{def:hh-clausal}
\[
\begin{array}{r@{\;=\;}l}
\mathit{clausal}(A) & 0 \qquad (A\text{ atomic or }\top)\\
\mathit{clausal}(B_1\wedge B_2) = \mathit{clausal}(B_1\vee B_2) & \max(\mathit{clausal}(B_1),\mathit{clausal}(B_2))\\
\mathit{clausal}(B_1\imp B_2) & \max(\mathit{clausal}(B_1)+1,\ \mathit{clausal}(B_2))\\
\mathit{clausal}(\forall x\,B) = \mathit{clausal}(\exists x\,B) & \mathit{clausal}(B)
\end{array}
\]
\end{definition}

In \fohc{} goals have clausal order~0 and clauses order 0 or~1; in \fohh{} both are unbounded \book{77}.
Reading $\imp$ as the type arrow $\to$, the definition is the same as that of the order of a type
(\cref{ch:intro}); the book stresses that this is ``not accidental'' and will exploit it in its
Section~7.6 \book{77}. Worked values: \lp{q :- (q => q)}, that is $(q\imp q)\imp q$,
has order~2; the $\cst{sterile}$ clause above has order~2; the clause of \cref{fig:hh-rev3} below has
order~3.

%=====================================================================================================
\section{Operational semantics: AUGMENT and GENERIC}\label{sec:hh-rules}

Nothing new has to be invented. The framework of \cref{ch:horn} already contains the rules for
implications and universal quantifiers at the top of goals: ``These are to be treated by the AUGMENT and
GENERIC reduction rules'' \book{77}. Atomic goals are still solved by backchaining, which is possible
because $D$-formulas have the same shape as in \fohc{}, modulo the richer goals. For reference,
\cref{fig:hh-rules} collects the rules: the right-introduction rules for goals and the backchaining rules
for a clause $D$ ``in focus'', written $\hhfoc{\Sig}{\Prog}{D}{A}$.

\begin{figure}[t]
\small
\begin{mathpar}
\inferrule*[right=$\top$R]{ }{\seq{\Sig}{\Prog}{\top}}
\and
\inferrule*[right=$\wedge$R]{\seq{\Sig}{\Prog}{B_1} \\ \seq{\Sig}{\Prog}{B_2}}{\seq{\Sig}{\Prog}{B_1\wedge B_2}}
\and
\inferrule*[right=$\vee$R]{\seq{\Sig}{\Prog}{B_i}}{\seq{\Sig}{\Prog}{B_1\vee B_2}}
\and
\inferrule*[right=$\exists$R]{\seq{\Sig}{\Prog}{B[t/x]} \\ \Sig;\emptyset\vdash_f t:\tau}{\seq{\Sig}{\Prog}{\exists_\tau x\,B}}
\\
\inferrule*[right={$\imp$R\ (\rulename{augment})}]{\seq{\Sig}{\Prog,B_1}{B_2}}{\seq{\Sig}{\Prog}{B_1\imp B_2}}
\and
\inferrule*[right={$\forall$R\ (\rulename{generic})}]{\seq{c:\tau,\Sig}{\Prog}{B[c/x]}}{\seq{\Sig}{\Prog}{\forall_\tau x\,B}}\ \ c\notin\Sig
\\
\inferrule*[right=decide]{\hhfoc{\Sig}{\Prog}{D}{A}}{\seq{\Sig}{\Prog}{A}}\ \ D\in\Prog
\and
\inferrule*[right=initial]{ }{\hhfoc{\Sig}{\Prog}{A}{A}}
\and
\inferrule*[right=$\imp$L]{\hhfoc{\Sig}{\Prog}{D}{A} \\ \seq{\Sig}{\Prog}{G}}{\hhfoc{\Sig}{\Prog}{G\imp D}{A}}
\\
\inferrule*[right=$\wedge$L]{\hhfoc{\Sig}{\Prog}{D_i}{A}}{\hhfoc{\Sig}{\Prog}{D_1\wedge D_2}{A}}
\and
\inferrule*[right=$\forall$L]{\hhfoc{\Sig}{\Prog}{D[t/x]}{A} \\ \Sig;\emptyset\vdash_f t:\tau}{\hhfoc{\Sig}{\Prog}{\forall_\tau x\,D}{A}}
\end{mathpar}
\caption{The rules of O-proofs: right rules (book Figure 2.2) and backchaining rules (book Figure 2.3).
In \fohc{} the two rules of the second row are never used. With logic variables, $\exists$R and
$\forall$L introduce a new variable $X$ instead of guessing $t$, and initial unifies.}
\label{fig:hh-rules}
\end{figure}

\begin{definition}[O-proof \book{77}]\label{def:hh-oproof}
An \emph{O-proof} in \fohh{} is a derivation built with the typing rules for first-order terms and the
rules of \cref{fig:hh-rules} (book Figures 1.2, 2.2 and 2.3).
\end{definition}

In the logic-variable versions of the rules (\cref{ch:horn}, \book{61}) the typing premise of
$\exists$R and $\forall$L is delayed and initial unifies. \Cref{sec:hh-generic} shows that once
\textsc{generic} is present the delayed choice must still respect the signature in which the variable was
created.

\begin{tutorial}[programs grow and shrink like a stack]
Take the goal $(D_0\imp((D_1\imp G_1)\wedge(D_2\imp G_2)))\wedge G_3$ and the program $\Prog$ \book{77}.
\rulename{and} splits it in two. The left conjunct is an implication: \rulename{augment} moves $D_0$ into
the program, and \rulename{and} splits again; each conjunct is an implication again. In the end we must
prove
\[
G_1 \text{ from } \{D_1,D_0\}\cup\Prog,\qquad
G_2 \text{ from } \{D_2,D_0\}\cup\Prog,\qquad
G_3 \text{ from } \Prog.
\]
$D_1$ is never visible to $G_2$, and $D_0$ is never visible to $G_3$. Read as an interpreter state,
the program is pushed when we enter the antecedent's scope and popped when the subproof is done:
``implicational goals allow logic programs to grow and shrink following a stack-based discipline''
\book{78}. Notice that in the proof \emph{tree} nothing is popped: each branch simply carries its own
program. Popping is what a sequential interpreter does when it walks the tree.
\end{tutorial}

\paragraph{Programs are lists.} Under a depth-first strategy the order of clauses matters, so programs
are ``best thought of not as sets but rather as lists of clauses'' \book{78}. λProlog adds assumed
clauses \emph{at the front}: ``the most recently added clauses are the first to be used in backchaining''
\book{78}. Inside an assumed conjunction \lp{D1 & D2} the textual order is kept. With a program
containing only \lp{p 1}:
\begin{lprolog}
?- p 2 => p 3 => p X.
X = 3;
X = 2;
X = 1
?- (p 2 & p 3) => p X.
X = 2;
X = 3;
X = 1
\end{lprolog}

\begin{kanren}[answer order is a depth-first artefact]
Running the same two queries through the small interpreter of \cref{sec:hh-interp}, which uses the
interleaving \scm{or°} of \micro{} for the choice among clauses, gives the same \emph{sets} of answers
in a different order:
\begin{answers}
p 2 => p 3 => p X    : (3 2 1)
(p 2 & p 3) => p X   : (2 1 3)
\end{answers}
In the second query the $\wedge$L choice inside \lp{p 2 & p 3} is interleaved with the decide choice
over the rest of the program. Tests of this interpreter should not compare ordered answer lists with
λProlog's; the order of the program still matters, since front clauses are tried first. The proposed
module does reproduce the book's order here, \scm{(3 2 1)} and \scm{(2 3 1)}, because a dynamic relation
is one flat disjunction of its hypotheses, most recent first (\cref{def:design-a2},
\cref{sec:examples-order}). Compare ordered lists only where every clause answers without further search.
\end{kanren}

%=====================================================================================================
\section{Implicational goals and hypothetical reasoning}\label{sec:hh-augment}

\subsection{Entailment between Horn clauses}\label{sec:hh-entail}

Let $q,r,s,t,u$ be propositions (type $\typ{o}$) and consider the program \book{78}
\begin{lprolog}
s :- r, q.
t :- q, u.
q :- r.
\end{lprolog}
It has no atomic consequences at all, yet clauses follow from it: \lp{t :- r, u.} and \lp{s :- r.}
\book{79}. In \fohh{} we can ask this directly. The query \lp{?- s :- r.} ``is better written as''
\lp{?- r => s.} \book{79}: \rulename{augment} adds \lp{r} and ordinary Horn reasoning proves \lp{s}.

\begin{tutorial}[the O-proof of \texttt{r => s}]
Write $\Prog_0$ for the three clauses and $\Prog'=\Prog_0,r$; the signature is omitted. The search, read
goal-first (each line is reduced to the next ones):
\begin{enumerate}
  \item $\Prog_0\longrightarrow r\imp s$: by $\imp$R (\textsc{augment}) it suffices to show
        $\Prog'\longrightarrow s$.
  \item $\Prog'\longrightarrow s$: decide on the clause $r\wedge q\imp s$; by $\imp$L we need the head to
        match, $\Prog'\xrightarrow{\,s\,}s$ (initial), and the body, $\Prog'\longrightarrow r\wedge q$.
  \item $\Prog'\longrightarrow r\wedge q$: by $\wedge$R, both $\Prog'\longrightarrow r$ and
        $\Prog'\longrightarrow q$.
  \item $\Prog'\longrightarrow r$: decide on the assumed clause $r$, then initial.
  \item $\Prog'\longrightarrow q$: decide on $r\imp q$; initial for the head, and the body
        $\Prog'\longrightarrow r$ as in step~4.
\end{enumerate}
The only new step is the first one; the
rest is the Horn backchaining of \cref{ch:horn}, run in a larger program. In the same way
\lp{(r => u) => (r => t)} is provable \book{79}: it assumes the clauses \lp{r => u} and \lp{r}, and then
\lp{t} follows through \lp{q} and \lp{u}. (This is a different entailment from \lp{t :- r, u}, which as a
query is \lp{r => u => t}.)
\end{tutorial}

The last example of the section is a warning \book{79}:
\begin{lprolog}
?- (q :- (q => q)) => (q :- (q => q)).
\end{lprolog}
It has the form $B\imp B$, so it has an O-proof from the empty program, but a depth-first interpreter
``will not terminate on the given query''. To see why, unfold it. \rulename{augment} assumes
$(q\imp q)\imp q$; the goal is again $(q\imp q)\imp q$, so \rulename{augment} assumes $q\imp q$, i.e.\ the
clause \lp{q :- q}, \emph{at the front}. Now the program is \lp{[q :- q, q :- (q => q)]} and the goal
is \lp{q}. Depth-first search picks \lp{q :- q} and loops for ever. The O-proof picks the other clause:
its body \lp{q => q} assumes \lp{q} and proves \lp{q} at once.

\begin{pitfall}[provable, yet the interpreter loops]
Depth-first backchaining is incomplete, and AUGMENT makes this easier to hit, because assumed clauses
go to the front. The book makes the same point with conjunction: \lp{loop, fail} loops while
\lp{fail, loop} fails, with \lp{loop :- loop.} \book{88}. In the interleaving search of miniKanren the
clause choice is fair, and the interpreter of \cref{sec:hh-interp} does find the proof of $B\imp B$, in fact
infinitely many of them (\scm{run 3} returns \texttt{(α α α)}); conjunction, however, remains as
order-sensitive as in λProlog.
\end{pitfall}

\subsection{Hypothetical reasoning on a database}\label{sec:hh-db}

Implications in goals express \emph{hypothetical} questions: ``what if this were true?''.
\Cref{fig:hh-db} is the book's course database: facts about which students finished which courses, rules
for being a CS major and for graduating, and a constraint saying it is inconsistent to have finished both
210 and 250.

\begin{figure}[t]
\begin{lprolog}
kind entry               type.
type fact                entry -> o.
type false               o.
kind person              type.
type kim, dana           person.
type finished            person -> int -> entry.
type cs_major, graduates person -> entry.

fact (finished kim  102) & fact (finished dana 101).
fact (finished kim  210) & fact (finished dana 250).

fact (cs_major X)  :-
     (fact (finished X 101); fact (finished X 102)),
     fact (finished X 250), fact (finished X 301).
fact (graduates X) :-
     (fact (finished X 101); fact (finished X 102)),
     (fact (finished X 210); fact (finished X 250)),
     fact (finished X 301).

false :- fact (finished X 210), fact (finished X 250).
\end{lprolog}
\caption{An encoding of a small course-related database (book Figure 3.1, printed p.\,80).}
\label{fig:hh-db}
\end{figure}

Note the $D\wedge D$ clauses on the two fact lines, and that \lp{false} is an ordinary proposition: it
has a clause, so it can be proved, but it has no logical meaning (\cref{sec:hh-minimal}). The queries
and answers of the book \book{80--81}:
\begin{lprolog}
?- fact (finished dana X) => false.
X = 210;
no
?- fact (finished X Y) => (fact (graduates X), fact (cs_major X)).
X = dana
Y = 301;
no
?- fact (finished kim X) => fact (finished kim Y) =>
               (fact (graduates kim), fact (cs_major kim)).
X = 250
Y = 301
?- fact (finished kim 250) => fact (finished kim 301) => false.
solved
\end{lprolog}
The first asks whether some course would make the database inconsistent if Dana finished it (210, and
only 210). The second asks: ``If person X took just one more course Y, will he or she graduate with a CS
degree?'' (only Dana, with 301); this ``indirectly implies that there is no single course'' Kim can take
\book{80}, hence the third query, with two hypotheses: Kim needs 250 and 301. The fourth shows that this
plan is inconsistent, so ``there is no consistent way for Kim to graduate as a computer science major''
\book{81}.

\begin{tutorial}[a logic variable inside an assumed clause]
Follow the first query. Its free variable $X$ is existential, so $\exists$R makes it a logic variable
\emph{before} \rulename{augment} runs. The assumed clause is therefore \lp{fact (finished dana X)} with
$X$ \emph{not} universally quantified: it is one specific, still unknown fact. Now the goal \lp{false}
backchains on the constraint clause, whose own variable (call it $Z$) is renamed by $\forall$L:
\begin{enumerate}
  \item \lp{fact (finished Z 210)}: decide on the assumed clause unifies $Z=\cst{dana}$, $X=210$;
  \item \lp{fact (finished dana 250)}: true in the database.
\end{enumerate}
So proving \lp{false} \emph{binds the variable inside the hypothesis}, and the binding is reported to
the query. In the second query the single hypothesis \lp{fact (finished X Y)} likewise has one binding
for $X$ and $Y$, shared by every use of the clause. \Cref{sec:hh-link} turns this into a
programming technique.
\end{tutorial}

\begin{kanren}[the four queries on \micro]
The interpreter of \cref{sec:hh-interp} reproduces the four answers: \texttt{(210)} for Dana's
inconsistency, \texttt{((dana 301))} for the one-course question, a success for the inconsistency of Kim's
plan, and, asked for \emph{all} answers, \texttt{(250 301)} and the symmetric \texttt{(301 250)}, each
twice, for the two-course question: the book shows only the first answer, and the others correspond to
different proofs: the symmetric pair because either hypothesis can supply either course, the duplicates
because the ``210 or 250'' disjunct in the body of \lp{graduates} can be proved both from the database
(210) and from the hypothesis that supplies 250. As for what this demands of
the engine: the program must be part of what a goal sees, it must be extended for the duration of a
subgoal only, and an assumed clause must be able to \emph{mention logic variables of the surrounding
computation without renaming them} (\cref{sec:hh-kanren-context}).
\end{kanren}

\subsection{An interactive ``what if'' database}

\Cref{fig:hh-dbloop} turns hypothetical reasoning into a small interactive program; it uses the
non-logical \lp{print}, \lp{read} and cut \lp{!} \book{81}.

\begin{figure}[t]
\begin{lprolog}
type db                           o.
kind command                      type.
type do                           command -> o.
type enter, query, whatif, check  entry -> command.
type quit, consis                 command.

db :- print "Command?", read Command, do Command.

do quit.
do (enter Fact) :- fact Fact => db.
do (query Q) :- (fact Q, !, print "yes\n"; print "no\n"), db.
do (whatif Conjecture) :- (fact Conjecture => db),
                          print "Resuming\n", db.
do consis :- false, print "no\n", !; print "yes\n".
do (check Entry) :- (fact Entry, print "yes\n", !;
                    fact Entry => false, print "no\n", !;
                    print "no, but it could be true\n"),  db.
\end{lprolog}
\caption{A program that allows hypothetical queries against an encoded database (book Figure 3.2,
printed p.\,81).}
\label{fig:hh-dbloop}
\end{figure}

The book's session \book{81--82} is reproduced in \cref{fig:hh-session}.

\begin{figure}[t]
\noindent\begin{minipage}[t]{0.49\linewidth}
\begin{answers}
?- db.
Command? query (graduates dana).
no
Command? whatif (finished dana 301).
Command? query (graduates dana).
yes
Command? query (cs_major dana).
yes
Command? quit.
Resuming.
\end{answers}
\end{minipage}\hfill
\begin{minipage}[t]{0.49\linewidth}
\begin{answers}
Command? query (finished kim 101).
no
Command? query (graduates kim).
no
Command? whatif (finished kim 301).
Command? query (graduates kim).
yes
Command? query (cs_major kim).
no
Command? whatif (finished kim 250).
Command? query (cs_major kim).
yes
Command? consis.
no
Command?
\end{answers}
\end{minipage}
\caption{The book's session with the program of \cref{fig:hh-dbloop} (printed pp.\,81--82); read the
left column first.}
\label{fig:hh-session}
\end{figure}


The interesting point is where the ``database state'' lives: \emph{in the stack of nested
\rulename{augment}s}. \lp{enter F} runs the rest of the session inside \lp{fact F => db}, so $F$
persists; \lp{whatif F} runs a \emph{sub-session} inside \lp{fact F => db}, and when \lp{quit} ends that
sub-session the clause continues with \lp{print "Resuming\n", db} in the outer program, where $F$ is gone.
That is why Dana stops graduating after \lp{Resuming}. The Kim hypotheses nest: the innermost session sees
both 301 and 250, so \lp{cs_major kim} holds and \lp{consis} reports an inconsistency (210 and 250 together
prove \lp{false}). No \lp{assert}/\lp{retract} is involved: undoing a hypothesis is just leaving a scope.

\begin{remark}
The transcript is slightly idealised: \lp{do consis} is the only command besides \lp{quit} whose clause
does not end with \lp{db}, so after \lp{consis} the current sub-session ends and the enclosing \lp{whatif}
should print \texttt{Resuming} before the next prompt.
\end{remark}

%=====================================================================================================
\section{Universal goals and scoped constants}\label{sec:hh-generic}

\textsc{generic} transforms the goal $\forall_\tau x\,G(x)$ in the context $\langle\Sig,\Prog\rangle$ into
$G[c/x]$ in the context $\langle\Sig\cup\{c:\tau\},\Prog\rangle$ for a token $c$ not in $\Sig$; which token
is chosen does not matter \book{82}.

\begin{definition}[eigenvariable, scoped constant \book{82}]\label{def:hh-eigen}
The new token $c$ introduced by \textsc{generic} is an \emph{eigenvariable}. Since it is never
instantiated during computation, the book also calls it a \emph{scoped constant}.
\end{definition}

The parallel with implication is exact: implication augments the \emph{program} for part of the
computation, $\forall$ augments the \emph{signature} \book{82}.

\begin{tutorial}[sterile jars \book{83}]
A jar is sterile if every bug in it is dead; a bug in a heated jar is dead; jar \lp{j} has been heated:
\begin{lprolog}
kind jar, bug        type.
type j               jar.
type sterile, heated jar -> o.
type dead, bug       bug -> o.
type in              bug -> jar -> o.

sterile J :- pi x\ bug x => in x J => dead x.
dead B    :- heated J, in B J, bug B.
heated j.
\end{lprolog}
The signature has \emph{no constants of type $\typ{bug}$}: the program assumes nothing about the existence
of bugs. Now solve \lp{sterile j}.
\begin{enumerate}
  \item Backchain on the first clause: the goal becomes \lp{pi x\ bug x => in x j => dead x}.
  \item \textsc{generic}: pick a fresh \lp{g}; the goal becomes \lp{bug g => in g j => dead g}, and
        \lp{g : bug} joins the signature.
  \item \textsc{augment} twice: \lp{bug g} and \lp{in g j} join the program; the goal is \lp{dead g}.
  \item Backchain on the second clause: \lp{heated J} succeeds with $J=\cst{j}$, then \lp{in g j} and
        \lp{bug g} succeed by the two assumptions.
\end{enumerate}
``After this goal succeeds, the constant \lp{g} is removed from the signature, and the two clauses
\lp{bug g} and \lp{in g j} are similarly removed from the program'' \book{83}. The proof talks about
an arbitrary bug, never about a particular one.
\end{tutorial}

\subsection{An intensional universal quantifier}

The jar proof illustrates that \fohh's $\forall$ is \emph{intensional}: ``Proofs of universal goals do not
make any assumptions about the structure of the domain of quantification, and in fact, the same generic
proof must work for every element of the domain'' \book{83}. The alternative, \emph{extensional}, reading
would prove every instance, possibly in different ways, typically by induction; but ``the fohh logic does
not encompass rules for induction'' \book{83--84}. With
\begin{lprolog}
kind nat   type.
type zero  nat.
type succ  nat -> nat.
type plus  nat -> nat -> nat -> o.

plus zero L L.
plus (succ N) M (succ P) :- plus N M P.
\end{lprolog}
the query \lp{?- pi N\ plus N zero N.} is \emph{not} provable, although every closed instance is: the
generic goal \lp{plus c zero c} matches no clause head, since \lp{c} is neither \lp{zero} nor a
\lp{succ}. The base and step cases of the obvious induction are provable,
\begin{lprolog}
?- plus zero zero zero,
   pi N\ plus N zero N =>
         plus (succ N) zero (succ N).
\end{lprolog}
but no rule of \fohh{} combines them into \lp{pi N\ plus N zero N} \book{84}.

\begin{pitfall}[$\forall$ is not ``for every value'']
Do not expect \lp{pi} to enumerate a domain, and do not confuse it with the existential query
\lp{plus N zero N}, which enumerates \lp{zero}, \lp{succ zero}, \dots{} for ever. An empty type is harmless
for $\forall$: the jar proof works although no bug exists. The bibliographic notes call the two readings
\emph{generic} and \emph{universal} judgments; they coincide under the open-world assumption of \fohh,
and under a closed world Miller and Tiu's $\nabla$ quantifier captures the generic one \book{94},
\cite{MillerTiu2005}.

The dual remark concerns $\exists$. The typing premise $\Sig;\emptyset\vdash_f t:\tau$ of $\exists$R is,
with logic variables, \emph{delayed}: ``all the remaining, delayed typing judgments must be provable for
instances of the logic variables'' \book{61}. A query that ends with an unbound logic variable of an empty
type (say a \lp{sigma} over \lp{bug} in the jar signature, with a body that does not constrain it) has no
O-proof, because no closed term of type \lp{bug} exists, although an interpreter reports success. Inside
a \textsc{generic} scope over \lp{bug} the type is inhabited by the scoped constant, which is why the jar
proof is fine. An untyped miniKanren silently assumes every unbound answer variable has an instance.
\end{pitfall}

\subsection{Substitution must avoid capture}

In \fohc{} a clause is $\forall x_1\dots\forall x_m(A_1\wedge\dots\wedge A_n\imp A_0)$ with all quantifiers
outermost, so instantiating them is plain replacement \book{84}. In \fohh{} quantifiers occur inside
clauses, and plain replacement can capture.

\begin{definition}[free for \book{85}]\label{def:hh-freefor}
Let $x$ be a free variable of $B$ and $t$ a term of the same type. $t$ is \emph{free for} $x$ in $B$ if no
free occurrence of $x$ in $B$ is in the scope of a quantifier that binds a variable free in $t$.
\end{definition}
If $t$ is free for $x$ in $B$, replacement realises $B[t/x]$ soundly; otherwise one first renames bound
variables of $B$, and ``we shall assume that substitution is always performed in this logically sound
fashion'' \book{85}.

\begin{example}[\book{85}]
Substituting \lp{(f y)} for \lp{X} in \lp{p X :- pi y\ q X y} by naive replacement gives
\lp{p (f y) :- pi y\ q (f y) y}, where the constant \lp{y} has been captured: there are now two bound
occurrences where there was one. The correct result is \lp{p (f y) :- pi z\ q (f y) z}.
\end{example}

\subsection{Quantifier order matters}\label{sec:hh-scope}

``Possibly the simplest example illustrating the relevance of quantifier scopes in computation'' is
\book{85}
\begin{lprolog}
?- sigma x\ pi y\ x = y.
\end{lprolog}
It is not provable: \lp{pi y\ t = y} has no proof for any term $t$ in which $y$ does not occur. It is
\emph{true} in a model with a single element, but not every model is like that \book{85--86}. Proof
search agrees: \textsc{generic} creates a new constant $c$ and must prove $c=t$, which fails. Choosing
$y$ itself for $x$ does not help, because capture-avoiding substitution turns the goal into
\lp{pi z\ y = z} \book{86}. Swapping the quantifiers changes everything: \lp{pi y\ sigma x\ x = y}
\emph{is} provable, with witness the scoped constant itself.

\begin{tutorial}[the scope condition, derived from the rules]
Where does the asymmetry come from in the rules of \cref{fig:hh-rules}? From the side condition of
$\exists$R: the witness $t$ must satisfy $\Sig;\emptyset\vdash_f t:\tau$, where $\Sig$ is the signature
\emph{at the moment the existential is reduced}.
\begin{itemize}
  \item In $\exists x\,\forall y\,(x=y)$, $\exists$R is applied first, with $\Sig$ not containing the
        constant $c$ that \textsc{generic} will create later. No $\Sig$-term equals $c$.
  \item In $\forall y\,\exists x\,(x=y)$, \textsc{generic} runs first, the signature becomes $c,\Sig$, and
        $t=c$ is a legal witness.
\end{itemize}
With logic variables the witness is chosen \emph{later}, by unification. The side condition does not go
away; it becomes a constraint on unification: \emph{a logic variable created when the signature was
$\Sig$ may only be bound to terms built from $\Sig$} (and from other variables bound under the same
restriction). This is the scope condition behind the remark that unification ``must be modified \dots so
as to ensure that quantifier scopes are respected'' when \fohh{} computations alternate essential
existential and universal quantifiers \book{94}.
\end{tutorial}

The condition has a subtle consequence when a variable is bound to a term that contains \emph{other},
younger variables. Nadathur, Jayaraman and Kwon's implementation of $\forall$ goals tags every constant and
logic variable with a \emph{universe index}: \textsc{generic} increments the index and tags the new
constant with it, a logic variable is tagged with the index current at its creation, and $X$ (tag $i$) may
be bound to $t$ only if (1) the occurs check passes, (2) $t$ contains no constant with tag $>i$, and (3)
before binding, the tags of the variables of $t$ that are $>i$ are lowered to~$i$
\cite{NadathurJayaramanKwon1995}. Rule (3) is the subtle one. Their example: with the program clause
$\forall z\,(q(z)\imp p(d(z)))$, consider the goal $\exists x\,\forall y\,(q(y)\imp p(x))$.
\begin{enumerate}
  \item $\exists$: a logic variable $X$ with tag~1. $\forall$: a constant $c$ with tag~2. \textsc{augment}:
        assume $q(c)$; the goal is $p(X)$.
  \item Backchain on the clause: $\forall$L creates $Z$ with the current tag~2; unify $p(X)=p(d(Z))$,
        binding $X:=d(Z)$, which is allowed: $d$ belongs to the original signature (tag~1).
  \item The body $q(Z)$ must now be proved from the assumption $q(c)$, i.e.\ $Z:=c$. If $Z$ still had tag~2
        this would succeed, and $X$ would be $d(c)$, a term mentioning a constant that did not exist when
        $X$ was created. Rule (3) lowered $Z$ to tag~1 when $X$ was bound, so $Z:=c$ fails, as it should.
\end{enumerate}
The goal is not provable, since $x$ must be fixed before $y$, and only the lowering detects it.

\begin{kanren}[what $\forall$ demands of the unifier]
A miniKanren that wants \fohh's $\forall$ needs two things that \micro{} does not have
\cite{MaLuFriedman2020,HemannByrdFriedman2013,NadathurJayaramanKwon1995}:
\begin{enumerate}
  \item \emph{A new kind of term}, the scoped (eigen) constant: unique, never bound, equal only to
        itself. It must be distinct from every existing term, including records with equal slots,
        since \micro{} unifies records slot-wise.
  \item \emph{A scope check}: no variable that existed before the constant was created may ever reach it
        through the substitution, including through younger variables (rule (3) above). Known designs:
        an explicit quantifier prefix consulted by unification, as in λKanren 2020
        \cite{MaLuFriedman2020}; numeric universe levels with lowering, as in Teyjus and ELPI
        \cite{NadathurJayaramanKwon1995,DunchevEtAl2015}; an absento-like ``eigen'' constraint between the
        constants and the list of variables in scope, as in miniKanren with symbolic constraints
        \cite{HemannByrdFriedman2013}; or a check of the answers when the scope is left, as the
        interpreter of \cref{sec:hh-interp} does.
\end{enumerate}
Which of these fits \micro, and how variable creation numbers relate to levels, is a design question
taken up in \cref{ch:design}. Note also that αKanren's names do \emph{not} play this role: a nom ``only
unifies with itself or with an unassociated variable'' \cite{Byrd2009}, so nothing prevents an older
variable from being bound to it. \Cref{sec:hh-kanren-eigen} runs three tempting shortcuts and shows each
one proving $\exists x\,\forall y\,(x=y)$.
\end{kanren}

%=====================================================================================================
\section{Quantification links goals and clauses}\label{sec:hh-link}

In Prolog the program is fixed, so logic variables only ever occur in goals. In λProlog an
implication goal can assume a clause that mentions a logic variable created by an enclosing $\exists$ or by
the head of the clause being used; logic variables then ``may appear eventually in program clauses''
\book{86}. We met this in the Dana query. Here it becomes a technique.

\Cref{fig:hh-rev3} defines \lp{reverse} through an accumulator predicate \lp{rev} whose clauses exist only
while a \lp{reverse} goal is being solved.

\begin{figure}[t]
\begin{lprolog}
type reverse    list A -> list A -> o.
type rev        list A -> list A -> list A -> o.

reverse L K :-
   ((pi L\ rev nil L L) &
    (pi X\ pi L\ pi K\ pi M\ rev (X::L) K M :- rev L K (X::M)))
        => rev L K nil.
\end{lprolog}
\begin{lprolog}
type reverse, rev    list A -> list A -> o.

reverse L K :-
   (rev nil K &
    (pi X\ pi L\ pi K\ rev (X::L) K :- rev L (X::K)))
        => rev L nil.
\end{lprolog}
\caption{Two implementations of reverse with locally assumed clauses: with a three-place auxiliary
predicate (top, book Figure 3.3) and with a two-place one (bottom, book Figure 3.4), printed p.\,87.}
\label{fig:hh-rev3}
\label{fig:hh-rev4}
\end{figure}

With the first version, \lp{?- reverse (1::2::nil) P.} reduces to \lp{?- rev (1::2::nil) P nil.} in the
program augmented with
\begin{lprolog}
rev nil L L.
rev (X::L) K M :- rev L K (X::M).
\end{lprolog}
and eventually binds \lp{P} to \lp{(2::1::nil)} \book{86--87}. The explicit \lp{pi}s ``cannot be
dropped: Doing so will change the scope of the quantification and hence also the meaning'' \book{86}:
without \lp{pi L\} the clause \lp{rev nil L L} would mention the \emph{head's} \lp{L}, the list being
reversed, instead of being a general fact.

The second version starts from a different observation \book{87}. From the program
\begin{lprolog}
rv nil (c::b::a::nil).
rv (X::N) M :- rv N (X::M).
\end{lprolog}
one can prove \lp{rv (a::b::c::nil) nil}; replacing \lp{(c::b::a::nil)} by any list \lp{L}, the goal
\lp{rv K nil} is provable if and only if \lp{L} and \lp{K} are reverses of each other. This cannot be
written in \fohc, ``because to do so we require the ability to `tie' the binding for a variable in the
specification of the reverse predicate to one that appears in a goal at a particular point in the
computation'' \book{87}. In \fohh{} it is the second program of \cref{fig:hh-rev4}.

\begin{tutorial}[running the two-place reverse]
Solve \lp{?- reverse (1::2::nil) P.} with the second program.
\begin{enumerate}
  \item Backchain on the \lp{reverse} clause: $\forall$L renames its \lp{L}, \lp{K} to fresh logic
        variables; unification binds the first to \lp{1::2::nil} and identifies the second with \lp{P}.
  \item \textsc{augment} adds, at the front,
\begin{lprolog}
rev nil K.
rev (X::L) K :- rev L (X::K).
\end{lprolog}
        Here the \lp{K} of the first clause is the logic variable shared with \lp{P}: ``not like the
        usual implicitly universally quantified variables'' \book{88}. The \lp{K} of the second clause is
        bound by its own \lp{pi K\} and is renamed at every use.
  \item The goal \lp{rev (1::2::nil) nil} becomes \lp{rev (2::nil) (1::nil)}, then
        \lp{rev nil (2::1::nil)}.
  \item This last goal matches \lp{rev nil K} by binding the shared \lp{K}, ``and hence \lp{P}'', to
        \lp{2::1::nil} \book{88}.
\end{enumerate}
The base case of the recursion was written \emph{before} the recursion ran, with a hole for the answer;
the computation fills the hole from inside the scope.
\end{tutorial}

\begin{pitfall}[implicit quantification lies inside assumed clauses]
The display convention ``actually becomes somewhat misleading when such variables can appear in
clauses'' \book{88}: in \lp{rev nil K} the capital \lp{K} looks universally quantified but is a single
shared unknown. Two practical consequences. First, one may not copy an assumed clause wholesale at each
use (the \lp{copy_term} reflex of Prolog meta-interpreters): only the variables bound by the clause's own
quantifiers are renamed, the others stay shared \cite{NadathurJayaramanKwon1995}. Second, \lp{rev} is an
ordinary global predicate name, so the local clauses mix with any other assumption about \lp{rev}. Under a
stray assumption \lp{rev nil Z}, the interpreter of \cref{sec:hh-interp} returns two answers to the reverse
query, the second one produced by the stray clause:
\begin{answers}
under rev () Z       : (((2 1) α) (α (2 1)))
\end{answers}
Making \lp{rev} truly local requires quantifying over the predicate itself, a higher-order feature
(\cref{ch:hoprog}).
\end{pitfall}

\begin{kanren}[shared variables come for free with closures]
In a miniKanren embedded in Scheme, an assumed clause $\forall\bar x\,(G\imp p(\bar t))$ can be represented
by a host closure that, when called with the arguments of a goal, runs \scm{fresh°} for $\bar x$, unifies
the arguments with $\bar t$ and then runs $G$. Every use renames $\bar x$, as $\forall$L requires, while
every logic variable that is merely \emph{free} in the clause, like \lp{K} above or Dana's \lp{X}, is
captured by the closure and therefore shared, exactly the semantics of this section, with no copying
machinery. What closures do not provide is the scope information of \cref{sec:hh-scope} for the renamed
$\bar x$: in the universe-index scheme each such variable is tagged with the index current at the time of use
\cite{NadathurJayaramanKwon1995}, and that index must come from the search state. See \cref{sec:module-hh}.
\end{kanren}

%=====================================================================================================
\section{The logic behind fohh}\label{sec:hh-logic}

Logic is useful for programs because logically equivalent formulas, $B\equiv C$, ought to be
interchangeable. With side effects or incomplete search they are not: \lp{print "yes", print "no"}
differs from its swap, and \lp{loop, fail} from \lp{fail, loop}. ``Nevertheless \dots if we rule out
programs with side effects and consider complete proof-search strategies, then logical equivalences can be
used meaningfully in reasoning about logic programs'' \book{88}. The question is then: \emph{which} logic?

\subsection{Not classical logic}

For \fohc{} the rules are complete for both classical and intuitionistic logic (\cref{ch:horn}). Not so
for \fohh.

\begin{theorem}[\book{89}]\label{thm:hh-intuitionistic}
The operational semantics for sequents in \fohh{} is sound and complete for intuitionistic logic, but
classical logic is not sound for this operational semantics.
\end{theorem}

Take $p\vee(p\imp q)$ from the empty program. By $\vee$R it is provable iff $p$ is provable from $\emptyset$
(no) or $q$ is provable from $\{p\}$ (no). So there is no O-proof, in agreement with intuitionistic logic.
Classically, however \book{89},
\[
B_1\vee(B_2\imp B_3)\ \equiv\ B_1\vee\neg B_2\vee B_3\ \equiv\ (B_2\imp B_1)\vee B_3\ \equiv\ B_2\imp(B_1\vee B_3),
\]
so $p\vee(p\imp q)\equiv p\imp(p\vee q)$, which is clearly provable. Classical logic lets the scope of an
implication be ``extruded'' over a disjunction (the term comes from the π-calculus, \book{95}), and so
``the interpretation we desire for implication as a scoping mechanism is not valid in classical logic''
\book{89}. Three more classically valid goals without O-proofs \book{89}:
\[
(r\,a\wedge r\,b\imp q)\imp\exists x\,(r\,x\imp q)
\qquad\qquad
((p\imp q)\imp p)\imp p
\qquad\qquad
\exists x\,\forall y\,(p\,x\imp p\,y)
\]
The middle one is Peirce's formula. ``It is clear from all these examples that classical logic cannot
provide the declarative semantics of λProlog'' \book{89}.

\begin{tutorial}[why the three formulas fail]
\textbf{First.} \textsc{augment} assumes the clause \lp{q :- r a, r b}; $\exists$R makes $X$;
\textsc{augment} assumes $r\,X$; the goal $q$ backchains and needs $r\,a$ and $r\,b$. The only $r$-fact is
$r\,X$: the first subgoal binds $X:=a$ and then $r\,b$ fails. Classically one could argue ``by cases on
$x$''; a single O-proof must commit to one witness.

\textbf{Peirce.} \textsc{augment} assumes $(p\imp q)\imp p$, i.e.\ \lp{p :- (p => q)},
goal $p$. Backchaining gives the goal $p\imp q$: assume $p$, prove $q$. Nothing proves $q$.

\textbf{Drinker.} (The third formula is a form of the classical ``drinker paradox''; the name is not the
book's.) $\exists$R makes $X$, \textsc{generic} makes $c$, \textsc{augment} assumes $p\,X$, goal
$p\,c$. The only way is to unify $p\,X=p\,c$, i.e.\ $X:=c$, which the scope condition forbids. This is
exactly the failure of $\exists x\,\forall y\,(x=y)$, and indeed an engine without the scope check
``proves'' it (\cref{sec:hh-interp}).
\end{tutorial}

\subsection{Minimal and intuitionistic negation}\label{sec:hh-minimal}

Define $\neg B$ as $B\imp\bot$. \emph{Minimal logic} treats $\bot$ as essentially a non-logical constant;
\emph{intuitionistic logic} adds \emph{ex falso quodlibet} \book{89--90}:
\begin{mathpar}
\inferrule*[right=$\bot$R]{\seq{\Sig}{\Prog}{\bot}}{\seq{\Sig}{\Prog}{B}}
\end{mathpar}
A program is \emph{inconsistent} if it entails $\bot$; intuitionistically, it then entails everything
\book{90}. In λProlog one picks a proposition, say \lp{false}, for $\bot$; since it is non-logical, ex falso
is unavailable, and the negation obtained is \emph{minimal logic negation} \book{90}. That is what the
database constraint of \cref{fig:hh-db} uses. The interpreter can be extended to intuitionistic negation:
``before failing on a given query, the interpreter would try to determine if \lp{false} is provable. If it
is, then it would have to succeed on the query in question. Such a check for inconsistency would have to be
coordinated with each (sub)query that leads to the program being augmented'' \book{90}.

\begin{example}[where minimal and intuitionistic negation part ways]
Consider the goal \lp{p => (p => false) => q}, i.e.\ $p\imp\neg p\imp q$. \textsc{augment} twice puts
\lp{p} and the clause \lp{false :- p} in the program; the goal \lp{q} matches no clause head, so there is
no O-proof, and the formula is not provable in minimal logic either. If \lp{false} is read as $\bot$, it is
intuitionistically provable: \lp{false} follows by backchaining on \lp{false :- p}, and ex falso gives
\lp{q}. The extended interpreter described above finds this proof: before failing on \lp{q} it tries
\lp{false}, in the augmented program, and succeeds. The check must be made in the augmented program,
which is why it has to be ``coordinated'' with \textsc{augment}: the empty program is consistent, the
program $\{p,\ \neg p\}$ is not. (The interpreter of \cref{sec:hh-interp}, which has no ex falso, returns
\texttt{()} for this goal.)
\end{example}

Minimal negation is weak but not useless. These are provable \book{90}:
\begin{lprolog}
(p => q) => ((q => false) => (p => false)).
p => ((p => false) => false).
\end{lprolog}
Their converses are provable neither minimally nor intuitionistically, but are classically. The instance of
excluded middle \lp{?- p; (p => false).} is not provable; its double negation
\lp{?- ((p; (p => false)) => false) => false.} is.

\begin{tutorial}[a proof that uses a hypothesis twice \book{91}]
In the display below, the formulas assumed so far are written to the left of \lp{?-}:
\begin{lprolog}
(p; (p => false)) => false ?- false.
(p; (p => false)) => false ?- p; (p => false).
(p; (p => false)) => false ?- p => false.
p, (p; (p => false)) => false ?- false.
p, (p; (p => false)) => false ?- p; (p => false).
p, (p; (p => false)) => false ?- p.
\end{lprolog}
In line 1 \textsc{augment} has assumed the clause \lp{false :- p; (p => false)}. From line 1 to 2 we
backchain on it; from 2 to 3 we choose the right disjunct; from 3 to 4 \textsc{augment} assumes \lp{p};
from 4 to 5 we backchain on the \emph{same} clause again, now in a larger program; from 5 to 6 we choose the
left disjunct, and ``This last query succeeds immediately because the goal \lp{p} is also in the program at
that point.'' The pattern, a
hypothesis used twice, the second time in a bigger world, is typical of intuitionistic proofs.
\end{tutorial}

\begin{theorem}[\book{91}]\label{thm:hh-minimal}
If $\Prog$ is a \fohh{} program and $G$ a \fohh{} goal over $\Sig$, then $\seq{\Sig}{\Prog}{G}$ has an
O-proof if and only if it is provable in minimal logic.
\end{theorem}

Since \fohh{} formulas contain no $\bot$ (only the non-logical \lp{false}), minimal and intuitionistic
provability coincide on them, ``for this reason, we shall say that the declarative semantics of λProlog is
also given by intuitionistic logic'' \book{91}.

\begin{kanren}[\texttt{=> false} is not negation as failure]
Minimal negation needs nothing from the engine: \lp{false} is an ordinary predicate, and
$\neg D$ as a goal is ``assume $D$, prove \lp{false}''. It must not be confused with the soft-cut negation
\scm{(if° g ✗° ✓°)} available in \micro, which is negation as failure and unsound on non-ground goals, nor
with the disequality constraint \scm{≠°}. Under the open-world reading of \fohh{} the book notes that
$A\imp\bot$ for atomic $A$ ``can never be true because there is always some future world in which $A$ is
true'' \book{94}. Classical principles (excluded middle, Peirce, scope extrusion) must \emph{not} be built
into a \fohh{} engine: they are unsound for its operational semantics.
\end{kanren}

\subsection{Why ``hereditary Harrop''}

Harrop showed that if all assumptions (the formulas left of the sequent arrow) are of the form
$H::=\top\mid A\mid H\wedge H\mid B\imp H\mid\forall_\tau x\,H$, with $B$ arbitrary, then intuitionistic
logic has the \emph{disjunction property} (a proof of $F_1\vee F_2$ yields a proof of $F_1$ or of $F_2$)
and the \emph{existence property} (a proof of $\exists x\,F$ yields a witness $t$ with a proof of
$F[t/x]$) \cite{Harrop1960}, \book{93}. Intuitively, these are what goal-directed search relies on: when
the goal is a disjunction or an existential, committing to one disjunct or one witness loses nothing. Harrop's condition
gives them only at the root of a proof; to have them at every point of goal-directed search, disjunctions and existentials
are prohibited ``hereditarily at all positive locations'' \book{93}, whence hereditary Harrop formulas
\cite{MillerNadathurPfenningScedrov1991}. Semantically, intuitionistic logic is modelled by Kripke models
over possible worlds; identifying worlds with signature--program pairs, a \fohh{} computation moves to larger
worlds (the open-world assumption) \book{94}.

%=====================================================================================================
\section{Notable subsets of fohh}\label{sec:hh-subsets}

Can the grammar be simplified without losing expressiveness? The book gives two restrictions \book{91--92}:
\begin{align}
G &::= \top \mid A \mid G\wedge G \mid G\vee G \mid \exists x\,G \mid D\imp G \mid \forall x\,G
\qquad D ::= A \mid G\imp D \mid D\wedge D \mid \forall x\,D \tag{3.1}\\
D &::= A \mid G\imp A \mid D\wedge D \mid \forall x\,D \tag{3.2}\\
D &::= A \mid D\imp D \mid D\wedge D \mid \forall x\,D \tag{3.3}
\end{align}

\paragraph{Atomic heads, (3.2).} Clauses of (3.1) can be preprocessed, ``without changing their essential
operational interpretation'', into clauses whose top-level implications have atomic conclusions
\book{91}; declaratively, every (3.1) clause is intuitionistically (and minimally) equivalent to a
\emph{set} of (3.2) clauses \book{92}. The equivalences doing the work are those of \fohc{} normalisation
(\cref{ch:horn}, \book{64}):
\begin{gather*}
\forall x(B_1\wedge B_2)\equiv\forall x B_1\wedge\forall x B_2,\qquad
B_1\imp(B_2\imp B_3)\equiv(B_1\wedge B_2)\imp B_3,\\
B_1\imp(B_2\wedge B_3)\equiv(B_1\imp B_2)\wedge(B_1\imp B_3),
\end{gather*}
and $B_1\imp\forall x\,B_2\equiv\forall x(B_1\imp B_2)$ when $x$ is not free in $B_1$.

\begin{example}[normalising a clause to atomic heads]
The (3.1) clause $\forall x\,(g\,x\imp(p\,x\wedge(h\,x\imp q\,x)))$ has a conjunction and an implication
in the conclusion of its top-level implication. Distributing $\imp$ over $\wedge$, then $\forall$ over
$\wedge$, and finally currying the inner implication gives the two (3.2) clauses
\[
\forall x\,(g\,x\imp p\,x)\qquad\text{and}\qquad\forall x\,((g\,x\wedge h\,x)\imp q\,x),
\]
that is, \lp{p X :- g X.} and \lp{q X :- g X, h X.} Note that the body $g\,x$ now occurs twice. The book
warns, for \fohc, that this kind of preprocessing can make programs exponentially larger (from
$G\imp(D_1\wedge D_2)$ to $(G\imp D_1)\wedge(G\imp D_2)$ the goal $G$ is duplicated) and need not preserve
operational behaviour such as the number of derivations \book{64--65}. Which of $g\,x$ and $h\,x$ is
tried first in the second clause is an operational choice that the logic does not fix.
\end{example}

\paragraph{No $\vee$, no $\exists$, (3.3).} Dropping disjunction and existentials from goals gives
$G::=A\mid G\wedge G\mid D\imp G\mid\forall_\tau x\,G$ with $D$ as in (3.2); ``now the G- and D-formulas are
identical'' (up to the normalisation above), so one grammar (3.3) serves for both: atoms closed under
\emph{unrestricted} $\imp$, $\wedge$ and $\forall$ \book{92}.

\begin{pitfall}[equivalent as programs, not as formulas]
The passage from (3.1)/(3.2) to (3.3) is \emph{not} a logical equivalence \book{92}. The clause
$(p\imp(q\vee r))\imp s$ is legal in (3.2) but is not intuitionistically equivalent to any set of (3.3)
formulas; it implies $((p\imp q)\imp s)\wedge((p\imp r)\imp s)$, but not conversely (classically the two
are equivalent). Yet ``as program clauses, however, these two formulas prove exactly the same atomic
formulas'' \book{92}. The book does not spell out why; one way to see it is operational. To use the
clause on the goal $s$ we must prove the goal $p\imp(q\vee r)$, i.e.\ assume $p$ and then, in the
augmented program, commit to $q$ or to $r$: exactly what the two split clauses test. The intuitionistic
difference would only show if $p\imp(q\vee r)$ could be \emph{assumed}, and in \fohh{} it never is,
because it is not a clause.
\end{pitfall}

Another way to get rid of $\vee$ and $\exists$ is to define them as non-logical predicates with higher-order
Horn clauses (\cref{ch:hoprog}).

%=====================================================================================================
\section{What implication and universal goals ask of a miniKanren}\label{sec:hh-kanren}

We can now say precisely what is missing from \micro{} (\cref{ch:micro}). A \micro{} goal maps a state
(a substitution indexed by variable creation number, plus constraint stores) to a lazy stream of states;
relations are Scheme procedures, so \emph{the program is the fixed Scheme environment}: exactly the
\fohc{} situation this chapter started from.

\begin{center}
\small
\begin{tabularx}{\linewidth}{@{}lX@{}}
\toprule
\fohh & \micro{} today\\
\midrule
$\top$, $G\wedge G$, $G\vee G$ & \scm{✓°}, \scm{and°}, \scm{or°}/\scm{cond°} (fair interleaving)\\
$\exists x\,G$ ($\exists$R with a logic variable) & \scm{fresh°}\\
$\forall x\,D$ in a clause ($\forall$L) & \scm{fresh°} in a relation body\\
decide and $\wedge$L over the clauses of a predicate & the \scm{cond°} of a \scm{define-relation}\\
initial, by unification & \scm{=°} against the head\\
$D\imp G$ (\textsc{augment}) & missing: the program is not part of the state\\
$\forall x\,G$ (\textsc{generic}) & missing: no scoped constants, no scope check\\
\bottomrule
\end{tabularx}
\end{center}

\subsection{The program context belongs to the search state}\label{sec:hh-kanren-context}

The sequent $\seq{\Sig}{\Prog}{G}$ shows the program as an \emph{inherited, read-only} context: no rule
returns a modified program to its caller. In a miniKanren this translates into four requirements.
\begin{enumerate}
  \item \emph{A goal must be able to read the current program} when it reaches an atomic goal, so that
        assumed clauses are tried together with (and, in λProlog, before) the static ones.
  \item \emph{The extension is local to a subgoal.} After $G$ in $D\imp G$ produces an answer, the
        conjunction that follows must see the old program again. If the program lives in the state that
        \scm{and°} threads from one conjunct to the next, each answer state of $G$ must have its program
        restored.
  \item \emph{The extension is local to a branch.} Two disjuncts with different hypotheses must not see
        each other's clauses, even though the interleaving search alternates between them.
  \item \emph{Free logic variables of an assumed clause are shared}, not renamed
        (\cref{sec:hh-link}).
\end{enumerate}
Requirements 2 and 3 are what make a host-language \emph{dynamic binding} (\texttt{parameterize},
\texttt{fluid-let}, or a global variable set and restored) tempting, and wrong. A \micro{} goal does almost
nothing when applied to a state: \scm{fresh°} and \scm{define-relation} wrap the real work in a promise
(\scm{δ}), which is forced later by whoever consumes the stream, after the dynamic extent of the binding
has ended. The following file puts a ``current program'' in a parameter and tries it three ways.
\begin{scheme}[fontsize=\footnotesize]
(import scheme (chicken base) (only (scheme base) make-parameter parameterize)
        (aux base) (aux stream) (aux kanren micro))

;; The "current program" in a host-language parameter: a list of clauses (λ (x) goal).
(define program (make-parameter '()))

(define (any-clause° cs x)                     ; try every clause, most recent first
  (if (null? cs) ✗° (or° ((car cs) x) (any-clause° (cdr cs) x))))

(define-relation (hyp° x)                      ; a predicate whose clauses live in `program`
  (λ (s) ((any-clause° (program) x) s)))

(define (assume° clause g)                     ; D ⊃ G by dynamic binding: WRONG
  (λ (s) (parameterize ((program (cons clause (program)))) (g s))))

(define (fact a) (λ (y) (=° y a)))             ; the clause  hyp a.

(define (show x) (display x) (newline))
(show (μkanren-run (q -1 #t) (assume° (fact 'a) (hyp° q))))                  ; 1
(show (μkanren-run (q -1 #t) (parameterize ((program (list (fact 'a)))) (hyp° q)))) ; 2
(show (parameterize ((program (list (fact 'a)))) (μkanren-run (q -1 #t) (hyp° q)))) ; 3
\end{scheme}
\begin{answers}
()
()
(a)
\end{answers}
Run 1 binds the parameter around the \emph{application} of the goal to the state; run 2 around the
\emph{construction} of the goal. Both lose the hypothesis: \scm{(hyp° q)} returns a promise, and the
program is read only when the promise is forced, outside \texttt{parameterize}. Only run 3, which forces the
\emph{entire} search inside the binding, works, and it is useless for scoping, since the ``scope'' is then the
whole query. The parameter-based \texttt{assume} in the code accompanying λKanren 2020 loses hypotheses in
the same way for relation calls nested in delayed bodies; the paper also describes a state-threaded design
without the problem \cite{MaLuFriedman2020}. Hence: \emph{hypotheses must live in the search state} (or be
passed as an explicit argument, as in \cref{sec:hh-interp}), never in host dynamic state; see
\cref{ch:design}.

\subsection{Scoped constants need a scope check}\label{sec:hh-kanren-eigen}

Three shortcuts suggest themselves for $\forall y\,G$ in a miniKanren. Each one ``proves''
$\exists x\,\forall y\,(x=y)$, which has no proof (\cref{sec:hh-scope}):
\begin{scheme}[fontsize=\footnotesize]
(import scheme (chicken base) (aux base) (aux stream) (aux kanren micro))
(define (show x) (display x) (newline))

;; Three wrong encodings of the goal  ∃x ∀y (x = y),  which has no proof.

; (a) ∀ read as ∃: y is just another logic variable
(show (μkanren-run (x -1 #t) (fresh° (y) (=° x y))))

; (b) ∀ read as "a brand new constant", with no scope check
(show (μkanren-run (x -1 #t) (let ((c (gensym 'c))) (=° x c))))

; (c) ∀ read as "a fresh variable that must still be unbound at the end"
(define (naive-∀° f)
  (fresh° (y) (f y) (project° ((v y)) (if (μkanren-var? v) ✓° ✗°))))
(show (μkanren-run (x -1 #t) (naive-∀° (λ (y) (=° x y)))))
(show (μkanren-run (x -1 #t) (naive-∀° (λ (y) (=° x (list 'f y))))))
\end{scheme}
\begin{answers}
(α)
(c117)
(α)
((f α))
\end{answers}
(a) confuses the quantifiers. (b) gets a rigid constant but lets it escape into the answer: this is the
``scope extrusion'' that the scope condition forbids. (c) is the subtle one: \micro{} binds the older
variable \scm{x} to the younger \scm{y}, so \scm{y} stays unbound and the check passes; and with
\scm{(=° x (list 'f y))} it is \scm{x} that gets bound in \emph{any} miniKanren, so checking \scm{y} can
never detect the escape. A rigid constant \emph{plus} a check on the variables older than it is required,
as stated in the kanren box of \cref{sec:hh-scope}.

\begin{pitfall}[scoped constants and the constraint stores]
\micro{} also has disequality, type and absento constraints, and the book's language has none of them.
Moiseenko remarks that ``there is no sound implementation of eigen variables with the support of
disequality constraints'' \cite{Moiseenko2019}, and the \texttt{eigen} of miniKanren with symbolic
constraints is documented as unsupported with constraints other than \texttt{==}
\cite{HemannByrdFriedman2013}. Take $\forall x\,(x\neq\alpha)$, with $\alpha$ outer. Under an
extensional reading of $\forall$ \cite{JinEtAl2021} it should fail, since $x$ could be $\alpha$'s value;
under \rulename{generic}'s reading of $x$ as a new name, as in λProlog's \lp{pi x\ not (x = A)}, it holds.
A naive eigen constant gives the second answer by accident, without a scope check. The proposal implements
both readings consistently, the generic one by default (\cref{def:roadmap-readings},
\cref{sec:module-stores}).
\end{pitfall}

\subsection{A throwaway interpreter of the rules}\label{sec:hh-interp}

To check the examples of this chapter on today's \micro, and to see the rules of \cref{fig:hh-rules}
at work, here is a small interpreter of O-proofs. It is \emph{not} a proposal for the new module
(\cref{ch:design,ch:module} do that): formulas are data, it re-examines every clause at every atomic goal,
and the program is an explicit argument \scm{P}. But it is faithful to the rules, and it is complete with
respect to the clause choice thanks to the interleaving \scm{or°}.
\begin{scheme}[fontsize=\footnotesize]
(import scheme (chicken base) (aux base) (aux stream) (aux kanren micro))

;; A throwaway interpreter of O-proofs over (aux kanren micro): formulas are data,
;; binders are Scheme closures.
;;   G ::= #t | A | (and G ...) | (or G ...) | (some ,(λ (x) G)) | (eq t u)
;;       | (imp D G) | (all ,(λ (x) G))
;;   D ::= A | (if G D) | (both D D) | (all ,(λ (x) D))
(define (tag? t F) (and (pair? F) (eq? (car F) t)))

(define (all-of° gs) (if (null? gs) ✓° (and° (car gs) (μ s ((all-of° (cdr gs)) s)))))
(define (one-of° gs) (if (null? gs) ✗° (or° (car gs) (μ s ((one-of° (cdr gs)) s)))))

(define (solve° P G)                                  ; Σ; P ⟶ G
  (cond
    ((eq? G #t)       ✓°)                                           ; ⊤R
    ((tag? 'and G)    (all-of° (map (λ (g) (solve° P g)) (cdr G))))  ; ∧R
    ((tag? 'or G)     (one-of° (map (λ (g) (solve° P g)) (cdr G))))  ; ∨R
    ((tag? 'some G)   (fresh° (x) (solve° P ((cadr G) x))))          ; ∃R, logic variable
    ((tag? 'eq G)     (=° (cadr G) (caddr G)))
    ((tag? 'imp G)    (solve° (cons (cadr G) P) (caddr G)))         ; AUGMENT, at the front
    ((tag? 'all G)    (generic° P (cadr G) #t))                     ; GENERIC
    ((tag? 'all/unchecked G) (generic° P (cadr G) #f))
    (else             (decide° P P G))))                            ; atomic goal

(define (decide° P Ds A)                              ; decide: some D ∈ P, front first
  (if (null? Ds) ✗° (or° (bc° P (car Ds) A) (μ s ((decide° P (cdr Ds) A) s)))))

(define (bc° P D A)                                   ; Σ; P —D→ A
  (cond
    ((tag? 'if D)   (and° (bc° P (caddr D) A) (μ s ((solve° P (cadr D)) s))))  ; ⊃L
    ((tag? 'both D) (or° (bc° P (cadr D) A) (bc° P (caddr D) A)))            ; ∧L
    ((tag? 'all D)  (fresh° (x) (bc° P ((cadr D) x) A)))                      ; ∀L
    (else           (=° D A))))                                               ; initial

;; GENERIC: a new constant c; an answer of G[c/x] survives only if no variable
;; created before c (index < n) reaches c through the substitution.
(define (mentions? c w)                              ; does the walked term w contain c?
  (or (eq? w c) (and (pair? w) (or (mentions? c (car w)) (mentions? c (cdr w))))))

(define (generic° P body check?)
  (λ (s)
    (let* ((n (μkanren-state-vars-count s))
           (c (gensym 'c))
           (ok? (λ (s*) (let loop ((i 0))
                          (or (= i n)
                              (and (not (mentions? c (μkanren-state-find/value (make-μkanren-var i) s*)))
                                   (loop (add1 i)))))))
           (answers ((solve° P (body c)) s)))
      (if check? (filter§ ok? answers) answers))))

(define-syntax ask                                    ; (ask n (q) P G): first n answers, -1 = all
  (syntax-rules () ((_ n (q) P G) (μkanren-run (q n #t) (solve° P G)))))
(define (show label x) (display label) (display x) (newline))
\end{scheme}

How the code mirrors the rules: \textsc{augment} is \scm{(solve° (cons D P) G)}, so the new clause goes to
the front, the scope ends with the Scheme call, and each branch carries its own \scm{P} (requirements 2
and 3 of \cref{sec:hh-kanren-context} hold by construction). Binders are Scheme closures, so $\forall$L and
$\exists$R call \scm{fresh°} when the clause or goal is used: a clause's own variables are renamed at each
use, while a logic variable merely mentioned in it is shared (requirement 4). \textsc{generic} creates a
\scm{gensym} constant (\micro{} compares symbols with \scm{eq?}) and enforces the scope condition lazily:
an answer of $G[c/x]$ is kept only if no variable with creation index below \scm{n}, i.e.\ older than
\scm{c}, reaches \scm{c}; the deep walk also catches escapes through younger variables. It prunes late,
which is acceptable for a test bench, not for an engine. \scm{all/unchecked} omits the check.

\paragraph{The examples.} The file below begins with \scm{(include "hh.scm")}, the listing above, and
runs a selection of the chapter's queries, in the book's order.
\begin{scheme}[fontsize=\footnotesize]
(include "hh.scm")
;; Section 3.2 of the book: clause order, Horn entailment, B ⊃ B
(show "p 2 => p 3 => p X    : " (ask -1 (x) '((p 1)) `(imp (p 2) (imp (p 3) (p ,x)))))
(show "(p 2 & p 3) => p X   : " (ask -1 (x) '((p 1)) `(imp (both (p 2) (p 3)) (p ,x))))
(define P0 '((if (and r q) s) (if (and q u) t) (if r q)))
(show "r => s ; t           : " (list (ask -1 (q) P0 '(imp r s)) (ask -1 (q) P0 't)))
(show "B => B, 3 proofs     : " (ask 3 (q) '() '(imp (if (imp q q) q) (imp (if q q) q))))
;; Figure 3.1, the course database
(define (fin x c) `(fact (finished ,x ,c)))
(define DB
  `((both ,(fin 'kim 102) ,(fin 'dana 101))
    (both ,(fin 'kim 210) ,(fin 'dana 250))
    (all ,(λ (X) `(if (and (or ,(fin X 101) ,(fin X 102)) ,(fin X 250) ,(fin X 301))
                      (fact (cs_major ,X)))))
    (all ,(λ (X) `(if (and (or ,(fin X 101) ,(fin X 102)) (or ,(fin X 210) ,(fin X 250))
                           ,(fin X 301))
                      (fact (graduates ,X)))))
    (all ,(λ (X) `(if (and ,(fin X 210) ,(fin X 250)) false)))))
(show "dana X => false      : " (ask -1 (x) DB `(imp ,(fin 'dana x) false)))
(show "X Y => grad, major   : "
      (μkanren-run (q -1 #t)
        (fresh° (x y) (=° q (list x y))
          (solve° DB `(imp ,(fin x y) (and (fact (graduates ,x)) (fact (cs_major ,x))))))))
(show "kim X => kim Y => .. : "
      (μkanren-run (q -1 #t)
        (fresh° (x y) (=° q (list x y))
          (solve° DB `(imp ,(fin 'kim x) (imp ,(fin 'kim y)
                         (and (fact (graduates kim)) (fact (cs_major kim)))))))))
(show "kim 250, 301 => false: " (ask -1 (q) DB `(imp ,(fin 'kim 250) (imp ,(fin 'kim 301) false))))
;; Section 3.3: jars, quantifier order, levels
(define JAR
  `((all ,(λ (J) `(if (all ,(λ (x) `(imp (bug ,x) (imp (in ,x ,J) (dead ,x))))) (sterile ,J))))
    (all ,(λ (B) `(all ,(λ (J) `(if (and (heated ,J) (in ,B ,J) (bug ,B)) (dead ,B))))))
    (heated j)))
(show "sterile Q            : " (ask -1 (q) JAR `(sterile ,q)))
(show "pi y\\ X = f y        : " (ask -1 (x) '() `(all ,(λ (y) `(eq ,x (f ,y))))))
(show "  ... unchecked      : " (ask -1 (x) '() `(all/unchecked ,(λ (y) `(eq ,x (f ,y))))))
(show "pi y\\ sigma x\\ x = y : " (ask -1 (q) '() `(all ,(λ (y) `(some ,(λ (x) `(eq ,x ,y)))))))
(define PQ `((all ,(λ (z) `(if (q ,z) (p (d ,z)))))))
(show "pi y\\ (q y => p X)   : " (ask -1 (x) PQ `(all ,(λ (y) `(imp (q ,y) (p ,x))))))
(show "  ... unchecked      : " (ask -1 (x) PQ `(all/unchecked ,(λ (y) `(imp (q ,y) (p ,x))))))
;; Figure 3.4: reverse; the base clause captures the outer K
(define REV
  `((all ,(λ (L) `(all ,(λ (K)
      `(if (imp (both (rev () ,K)
                      (all ,(λ (X) `(all ,(λ (L) `(all ,(λ (K)
                          `(if (rev ,L (,X . ,K)) (rev (,X . ,L) ,K)))))))))
                (rev ,L ()))
           (reverse ,L ,K))))))))
(show "reverse (1 2) P      : " (ask -1 (p) REV `(reverse (1 2) ,p)))
(show "reverse P (1 2 3)    : " (ask 1 (p) REV `(reverse ,p (1 2 3))))
(show "under rev () Z       : "
      (μkanren-run (q -1 #t)
        (fresh° (p z) (=° q (list p z)) (solve° REV `(imp (rev () ,z) (reverse (1 2) ,p))))))
;; Section 3.4: Peirce, the drinker, double negation of excluded middle
(show "Peirce               : " (ask -1 (q) '() '(imp (if (imp p q) p) p)))
(show "drinker              : " (ask -1 (q) '() `(some ,(λ (x) `(all ,(λ (y) `(imp (p ,x) (p ,y))))))))
(show "  ... unchecked      : "
      (ask -1 (q) '() `(some ,(λ (x) `(all/unchecked ,(λ (y) `(imp (p ,x) (p ,y))))))))
(show "p ; (p => false)     : " (ask -1 (q) '() '(or p (imp p false))))
(show "not not (p ; not p)  : " (ask 1 (q) '() '(imp (if (or p (imp p false)) false) false)))
(show "p => not p => q      : " (ask -1 (q) '() '(imp p (imp (if p false) q))))
\end{scheme}
\begin{answers}
p 2 => p 3 => p X    : (3 2 1)
(p 2 & p 3) => p X   : (2 1 3)
r => s ; t           : ((α) ())
B => B, 3 proofs     : (α α α)
dana X => false      : (210)
X Y => grad, major   : ((dana 301))
kim X => kim Y => .. : ((301 250) (250 301) (250 301) (301 250))
kim 250, 301 => false: (α)
sterile Q            : (j)
pi y\ X = f y        : ()
  ... unchecked      : ((f c488))
pi y\ sigma x\ x = y : (α)
pi y\ (q y => p X)   : ()
  ... unchecked      : ((d c563))
reverse (1 2) P      : ((2 1))
reverse P (1 2 3)    : ((3 2 1))
under rev () Z       : (((2 1) α) (α (2 1)))
Peirce               : ()
drinker              : ()
  ... unchecked      : (α)
p ; (p => false)     : ()
not not (p ; not p)  : (α)
p => not p => q      : ()
\end{answers}
A provable closed goal answers \texttt{(α)} (a proof, nothing to report); an unprovable one answers
\texttt{()}. Every line agrees with the book; the Kim query lists all answers, while the book shows only
the first. The goal \texttt{t} on its own fails, since the Horn program has no atomic consequences, and the
last line is the minimal-logic failure of $p\imp\neg p\imp q$ discussed in \cref{sec:hh-minimal}. In \texttt{sterile Q} the outer variable is bound to the constant \lp{j}, which is fine: \lp{j}
belongs to the signature in which \lp{Q} was created. In the reverse example, \scm{(rev () ,K)} is built
with the \emph{outer} \scm{K}, a logic variable, which the closure captures: the base case shares the
answer with the query. The three \texttt{unchecked} lines are the instructive ones. Without the scope
condition a scoped constant leaks into an answer, once directly and once through the clause variable $Z$ of
the level-lowering example, and the drinker formula $\exists x\,\forall y\,(p\,x\imp p\,y)$, valid
classically but not intuitionistically, becomes ``provable''. Dropping the scope check does not only produce
odd answers: it changes the logic.

%=====================================================================================================
\section{Summary}\label{sec:hh-summary}

\begin{itemize}
  \item \fohh{} extends \fohc{} with $D\imp G$ and $\forall x\,G$ in goals; the clause grammar is unchanged,
        and the two grammars swap roles across $\imp$. Hereditary Harrop formulas are those with no positive
        $\vee$ or $\exists$ (\cref{def:hh-grammar,prop:hh-polarity}).
  \item \textsc{augment} proves $D\imp G$ with $D$ added at the front of the program; \textsc{generic} proves
        $\forall x\,G$ with a new scoped constant in the signature. Both follow a stack discipline
        (\cref{fig:hh-rules}).
  \item Implication gives hypothetical reasoning: entailment between clauses, ``what if'' queries, and
        a database whose state is the stack of open hypotheses (\cref{sec:hh-augment}).
  \item $\forall$ is intensional: one generic proof for all elements, no induction, empty types harmless
        (\cref{sec:hh-generic}).
  \item With logic variables, $\exists\forall$ and $\forall\exists$ differ through a \emph{scope condition}: a
        variable may only be bound to terms over the signature of its creation, and binding must propagate this
        restriction to the variables of the bound term (universe levels and lowering, \cref{sec:hh-scope}).
  \item Assumed clauses may contain logic variables of the surrounding computation; these are shared, not
        renamed, which lets a clause ``tie'' its binding to a goal (reverse, \cref{sec:hh-link}).
  \item O-proofs are sound and complete for intuitionistic logic and coincide with minimal provability on
        \fohh; classical logic is unsound for them (scope extrusion, Peirce, drinker). \lp{false} gives
        minimal negation, not negation as failure (\cref{thm:hh-intuitionistic,thm:hh-minimal}).
  \item Clauses can be normalised to atomic heads (3.2); removing $\vee$ and $\exists$ (3.3) preserves atomic
        consequences but not logical equivalence (\cref{sec:hh-subsets}).
  \item For a miniKanren: the program must live in the search state (per branch, restored after each answer of
        the scoped subgoal), never in host dynamic binding, because goals run lazily; $\forall$ needs rigid
        scoped constants and a scope check in, or around, unification (\cref{sec:hh-kanren}, and
        \cref{ch:design} for the design).
\end{itemize}

\chapter{Typed λ-terms and unification problems}\label{ch:lambda}

\Cref{ch:horn,ch:hh} did logic programming over \emph{first-order} terms. \Cref{ch:hh} added
universal quantifiers and implications to goals (\fohh), and with them eigenvariables: constants made
up during the search that no logic variable created earlier may capture. Chapter~4 of the book
\cite{MillerNadathur2012} keeps the logic and changes the \emph{terms}. First-order terms give way to
\emph{simply typed λ-terms}. This chapter studies those terms on their own and leaves the higher-order
logic built on them to \cref{ch:hoprog} \book{96}.

The change has three consequences, and each of them touches a different part of a miniKanren:
\begin{enumerate}[nosep]
  \item \emph{What equality means.} Two terms are equal when they are λ-convertible (α, β, η), not
        only when they are identical. This affects the term representation and the unifier.
  \item \emph{What a logic variable may be bound to.} A unification problem becomes a
        \emph{quantified formula}. Its prefix mixes ∀ and ∃, and the order of the quantifiers decides
        which symbols a variable may use. This affects the state, which has to track scope.
  \item \emph{How hard unification is.} There may be several, infinitely many, or no most general
        unifiers, and the problem is undecidable. This affects the search: unification has to become a
        stream of answers.
\end{enumerate}
The chapter follows the book's order. \Cref{sec:lambda-syntax,sec:lambda-conversion,sec:lambda-normal}
cover the calculus. \Cref{sec:lambda-unifprob,sec:lambda-raising,sec:lambda-solutions,sec:lambda-examples}
recast unification as quantified equalities. \Cref{sec:lambda-vocab} fixes the vocabulary that
\cref{ch:hounif} uses for the unification procedure, and \cref{sec:lambda-undecidable} proves
undecidability.

% ======================================================================================
\section{Syntax and types of λ-terms}\label{sec:lambda-syntax}

Typed λ-terms are built from typed constants and variables with two operations. The first is
\emph{application}, written as juxtaposition and associating to the left, exactly as for first-order terms.
The second is \emph{abstraction}, which is new: from an expression $t$ in which a variable $x{:}\alpha$
may occur free, it forms the function $\lambda(x{:}\alpha)\,t$ \book{97}. In λProlog this is written
\lp{(x:A)\ T}. The type annotation can be dropped (\lp{x\ T}), and the type is then inferred.

\paragraph{Parsing.} The body of an abstraction extends as far to the right as possible. So
application binds tighter than abstraction, and abstraction associates to the right \book{97}:
\begin{center}
\begin{tabular}{@{}ll@{}}
\toprule
mathematical & λProlog concrete syntax\\
\midrule
$\lambda f\,\lambda x\,(f\,(f\,(f\,x)))$ & \lp{f\x\ f (f (f x))}\\
$\lambda x\,(f\,(g\,\lambda y\,(h\,x\,y))\,x)$ & \lp{x\ f (g y\ h x y) x}\\
$\lambda x\,\lambda y\,x$ & \lp{x\y\ x}\\
\bottomrule
\end{tabular}
\end{center}
A backslash cannot occur inside an identifier, so no spaces are needed around it.

\begin{definition}[Types and terms]\label{def:lambda-syntax}
Types are generated from primitive types (sorts such as $\typ{i}$, and the type $\typ{o}$ of
propositions), type variables $A$, type constructors such as $\typ{list}\,A$, and the arrow
$\tau_1\to\tau_2$, which associates to the right. Terms are
\[
  t \;::=\; c \;\mid\; x \;\mid\; (t\;t) \;\mid\; \lambda(x{:}\tau)\,t .
\]
(The book introduces the grammar in prose \book{97--99}. This compact form is ours.)
\end{definition}

\paragraph{Typing.} The judgement $\Sig;\Gamma\Vdash t:\tau$ (in the book, with a double turnstile) says
that $t$ has type $\tau$ when $\Sig$ gives the types of the constants and $\Gamma$ those of the
variables. It is defined by the rules of Figure~4.1 \book{98}:
\begin{mathpar}
\inferrule{c:\sigma\in\Sig \\ \tau\triangleleft\sigma}{\Sig;\Gamma\Vdash c:\tau}
\and
\inferrule{x:\tau\in\Gamma}{\Sig;\Gamma\Vdash x:\tau}
\and
\inferrule{\Sig;\Gamma\Vdash g:\tau_1\to\tau_2 \\ \Sig;\Gamma\Vdash t:\tau_1}{\Sig;\Gamma\Vdash (g\;t):\tau_2}
\and
\inferrule{\Sig;\Gamma,x:\tau\Vdash t:\sigma}{\Sig;\Gamma\Vdash \lambda(x{:}\tau)\,t:\tau\to\sigma}\;(\dagger)
\and
\inferrule{\Sig;\Gamma\Vdash B:\tau}{\Sig;\Gamma\Vdash C:\tau}\;(\ddagger)
\end{mathpar}
Proviso $(\dagger)$: $x$ is not declared as a type or kind in $\Sig$ or $\Gamma$. Proviso $(\ddagger)$:
$B$ and $C$ differ only in the names of bound variables. $\tau\triangleleft\sigma$ means that $\tau$ is a
substitution instance of $\sigma$, so a polymorphic constant is instantiated afresh at every use
\book{97}. Compared with the first-order rules of \cref{ch:horn}, the restrictions on the types in
$\Sig$ and $\Gamma$ are gone: a constant may have any type, including a functional type that ends in
$\typ{o}$. There is one new rule, for abstraction. And there are \emph{no special rules for the
quantifiers}, because the quantifiers have become ordinary constants (see below) \book{97--98}.

\begin{definition}[Well-formed, closed, simply typed {\book{98}}]\label{def:lambda-wf}
A term $t$ is a \emph{well-formed λ-term over $\Sig$} if $\Sig;\Gamma\Vdash t:\tau$ is derivable for
some $\Gamma$ and $\tau$. It is \emph{closed} if this holds with $\Gamma$ empty, and \emph{open}
otherwise. It is \emph{simply typed} if neither its type nor the types of its subterms contain type
variables.
\end{definition}

When every bound variable carries its annotation, the type of a well-formed term is unique up to
renaming of type variables. When annotations are omitted, type inference gives each bound variable a
fresh type variable and refines it \enquote{in as minimal a way as possible}. The type found this
way is again unique up to renaming, the principal type \book{98--99}. The unification theory of
\cref{ch:hounif} is stated for simply typed terms.

\begin{example}[A typing derivation]\label{ex:lambda-typing}
The Church numeral \lp{f\x\ f (f x)} has type $(\typ{i}\to\typ{i})\to\typ{i}\to\typ{i}$. Write
$\Gamma=f{:}\typ{i}\to\typ{i},\,x{:}\typ{i}$:
\begin{mathpar}
\inferrule*[right=($\dagger$)]{
  \inferrule*[right=($\dagger$)]{
    \inferrule*[right=app]{
      \inferrule*{f:\typ{i}\to\typ{i}\in\Gamma}{\Sig;\Gamma\Vdash f:\typ{i}\to\typ{i}} \\
      \inferrule*[right=app]{\Sig;\Gamma\Vdash f:\typ{i}\to\typ{i} \\ \Sig;\Gamma\Vdash x:\typ{i}}{\Sig;\Gamma\Vdash f\,x:\typ{i}}
    }{\Sig;\Gamma\Vdash f\,(f\,x):\typ{i}}
  }{\Sig;f{:}\typ{i}\to\typ{i}\Vdash \lambda x\,f\,(f\,x):\typ{i}\to\typ{i}}
}{\Sig;\emptyset\Vdash \lambda f\,\lambda x\,f\,(f\,x):(\typ{i}\to\typ{i})\to\typ{i}\to\typ{i}}
\end{mathpar}
Read the derivation bottom up. Each abstraction rule $(\dagger)$ moves one binder into the context.
The application rule \textsc{app} then checks $f$ against a functional type twice, and the leaves are
instances of the variable rule (the two leaves above $f\,x$ are left implicit). Because the derivation
ends with an empty context, the term is closed. No constant was used, so this derivation works for
every $\Sig$, including one with no constant of type $\typ{i}$, a case that matters in
\cref{sec:lambda-church,sec:lambda-solutions}.
\end{example}

\subsection{Formulas are terms of type o}\label{sec:lambda-formulas}

A \emph{formula} is a well-formed term of type $\typ{o}$. The connectives are constants such as
$\wedge:\typ{o}\to\typ{o}\to\typ{o}$. The quantifiers become constants too
\book{99}:
\[
  \forall : (A\to\typ{o})\to\typ{o}, \qquad \exists : (A\to\typ{o})\to\typ{o}
  \qquad\text{(concretely \lp{pi} and \lp{sigma})},
\]
and $\forall x\,B$ abbreviates $\forall(\lambda x\,B)$, written $\forall_\tau x\,B$ when the type of $x$
matters. A quantified formula is typed with the ordinary rules. From $\Sig;\Gamma,x{:}\sigma\Vdash B:\typ{o}$
the abstraction rule gives $\lambda x\,B:\sigma\to\typ{o}$, the instance
$\forall:(\sigma\to\typ{o})\to\typ{o}$ of the constant gives the rest, and so $\forall(\lambda x\,B):\typ{o}$.
For example, applying \lp{pi} at type \lp{(list int -> o) -> o} to an abstraction of type
\lp{list int -> o} \book{99}:
\begin{lprolog}
y\ append (1::2::nil) y X.
pi y\ append (1::2::nil) y X.
\end{lprolog}
Since signatures are no longer restricted, logical symbols may now appear \emph{inside the arguments of
non-logical ones} \book{99--100}:
\begin{lprolog}
type foreach  (A -> o) -> list A -> o.
foreach (x\ x > 5, x < 9) (3::10::6::8::nil).
\end{lprolog}
Here the conjunction \lp{,} sits inside an argument of \lp{foreach}. Higher-order programming
(\cref{ch:hoprog}) rests on this possibility. One more fact matters a lot later. \emph{The logic has a
single binder, λ-abstraction}: ∀, ∃ and every object-level binder (\cref{ch:lambdatree}) are encoded
with it \book{100}.

\begin{kanren}[terms are data, goals are procedures]
In \micro{} a goal is a Scheme procedure from a state to a stream of states. A procedure is opaque to
the unifier: it cannot be walked, compared or reified. So a formula on which one wants to
\emph{unify}, for example to solve for a predicate variable, has to be represented as
\emph{data}, such as a record for $\wedge$ or for \lp{pi} applied to an abstraction, and then
interpreted by a meta-level relation. This is the design of \cref{ch:hoprog,ch:lkanren}. For
\lp{foreach} with a ground Scheme predicate, ordinary higher-order Scheme is enough. Note also that
the term language now needs a representation of λ itself, which Scheme closures cannot provide:
\cref{sec:lambda-normal,sec:lambda-unifprob} use de Bruijn terms \cite{deBruijn1972} (lists in the first, records in the second), and \cref{ch:design}
discusses the options.
\end{kanren}

% ======================================================================================
\section{λ-conversion: the equality of the logic}\label{sec:lambda-conversion}

An abstraction is meant to denote a function, and an application of an abstraction is meant to denote
the result of the function. The conversion rules make this intended meaning part of the equality.
They rely on substitution, which must not capture variables.

\begin{definition}[Free for; substitution {\book{100}}]\label{def:lambda-subst}
$t$ is \emph{free for} $x$ in $s$ if no free occurrence of $x$ in $s$ lies in the scope of an
abstraction that binds a free variable of $t$. If $t$ is free for $x$ in $s$ and both have the same
type, $\subst{s}{t}{x}$ is $s$ with every free occurrence of $x$ replaced by $t$.
\end{definition}
For example, $(f\,x)$ is free for $u$ in $\lambda w\,(g\,u\,w)$, but $(f\,w)$ is not: substituting it
would put its free $w$ under the binder $\lambda w$. Ignoring the proviso is logically unsound, for the
same reason as with quantified formulas in \cref{ch:hh}.

\begin{definition}[The rules of λ-conversion {\book{100--101}}]\label{def:lambda-conv}
The following rewrites may be applied to any subterm:
\begin{center}
\begin{tabular}{@{}lll@{}}
\toprule
rule & redex $\Rightarrow$ contractum & proviso\\
\midrule
α-rewriting & $\lambda x\,s\;\Rightarrow\;\lambda y\,\subst{s}{y}{x}$ & $y$ free for $x$ in $s$, $y$ not free in $s$\\
β-contraction & $(\lambda x\,s)\,t\;\Rightarrow\;\subst{s}{t}{x}$ & $t$ free for $x$ in $s$ (lifted below)\\
η-contraction & $\lambda x\,(s\,x)\;\Rightarrow\;s$ & $x$ not free in $s$\\
\bottomrule
\end{tabular}
\end{center}
The converses are β-expansion and η-expansion. \emph{α-conversion} is the reflexive, symmetric,
transitive closure of α-rewriting. \emph{β-reduction} is the reflexive, transitive closure of
α-conversion together with β-contraction, and \emph{β-conversion} is its symmetric, transitive closure.
η-reduction and η-conversion are defined in the same way. \emph{λ-conversion} is the transitive
closure of α-, β- and η-conversion.
\end{definition}

λ-conversion is an equivalence and a congruence, and \emph{it is the equality of higher-order logic}
\book{101}. Two consequences stand out. First, terms related by λ-conversion cannot be told apart in
the logic. Second, \emph{no specification can determine the name of a bound variable}, because that
name can always be changed by α-conversion without changing the term. \Cref{ch:lambdatree} builds
λ-tree syntax on this second fact.

\begin{example}[Four equal terms {\book{101}}]\label{ex:lambda-four}
\begin{lprolog}
x\y\ f (g x) y
X\Y\ f (g X) Y
x\ f (g x)
x\y\ f ((u\v\v) (2 + 3) (g x)) y
\end{lprolog}
The first two terms are α-convertible. The third is the first after an η-contraction:
\[\lambda y\,(f\,(g\,x)\,y)\;\Rightarrow\;f\,(g\,x).\]
A β-contraction of the fourth gives
\begin{lprolog}
x\y\ f ((v\v) (g x)) y
\end{lprolog}
and a second one gives back the first. Note that \lp{(2 + 3)} disappears entirely, because $u$ does
not occur in the body \lp{v\v}.
\end{example}

\begin{pitfall}[capital letters are not always logic variables]
In \lp{X\Y\ f (g X) Y} the tokens \lp{X} and \lp{Y} are \emph{bound} names, not logic variables. In
λProlog only \emph{free} capitalised tokens are logic variables. A Scheme reader of λProlog code has to
look at where a name is bound before deciding what it is.
\end{pitfall}

\paragraph{Liberalised substitution.} From now on $\subst{s}{t}{x}$ is defined even when $t$ is not free
for $x$ in $s$. First α-rename $s$ to some $s'$ in which $t$ is free for $x$, then substitute. The result
depends on the renaming only up to α, so modulo equality it is well defined. β-contraction then becomes
$(\lambda x\,s)\,t\Rightarrow\subst{s}{t}{x}$ with no proviso \book{101}. This is ordinary
capture-avoiding substitution.

\begin{tutorial}[capture-avoiding β, by hand]
Contract $(\lambda u\,\lambda w\,(g\,u\,w))\,(f\,w)$. The argument has a free $w$. A naive textual
replacement would give $\lambda w\,(g\,(f\,w)\,w)$, where the outer $w$ has been \emph{captured} and now
refers to the bound variable. This is a different function. Liberalised substitution first renames the
binder, $\lambda u\,\lambda w'\,(g\,u\,w')$, and then substitutes:
\[
  (\lambda u\,\lambda w\,(g\,u\,w))\,(f\,w)\;\Rightarrow\;\lambda w'\,(g\,(f\,w)\,w').
\]
With de Bruijn indices (\cref{sec:lambda-normal}) the renaming becomes index arithmetic: the argument
is shifted by one when it goes under a binder.
\end{tutorial}

% ======================================================================================
\section{Normal forms and the power of β}\label{sec:lambda-normal}

A \emph{β-redex} is a term $(\lambda x\,s)\,t$. An \emph{η-redex} is a term $\lambda x\,(t\,x)$ with $x$
not free in $t$. A term with no β-redexes is in \emph{β-normal form}. If it has no η-redexes either, it
is in \emph{λ-normal form} (also called βη-normal form) \book{101--102}.

\begin{theorem}[Normal forms {\book{102}}]\label{thm:lambda-nf}
Every term of the typed language has a λ-normal form, unique up to α-conversion. It is written
$\mathrm{λnorm}(s)$. It can be computed by contracting β-redexes until none is left and then
contracting η-redexes.
\end{theorem}

\begin{corollary}[Decidable equality {\book{102}}]\label{cor:lambda-eq}
Two terms of the same type are λ-convertible if and only if their λ-normal forms are α-equivalent.
\end{corollary}

\begin{remark}[Background the book does not spell out]
The book states \cref{thm:lambda-nf} as \enquote{well known}. In the literature it rests on two
classical properties of the simply typed λ-calculus. \emph{Strong normalisation} says that every
reduction sequence terminates, which gives existence. The \emph{Church–Rosser property} (confluence)
says that diverging reductions can be joined again, which gives uniqueness. The algorithm also
relies on a standard fact the book uses without proof: η-contracting a β-normal term creates no new
β-redex, so doing η after β is enough.
\end{remark}

\begin{pitfall}[equality is intensional]
Logical equivalences do not affect term equality. With $p,q:\typ{i}\to\typ{o}$, the term
$\lambda x\,(p\,x)\wedge(q\,x)$ is \emph{not} equal to $\lambda x\,(q\,x)\wedge(p\,x)$ \book{102}. Types
are inhabited by \emph{expressions}, not by sets. A syntactic unifier gets this right automatically. A
unifier that knew about commutativity of $\wedge$ would be wrong.
\end{pitfall}

\subsection{Computing with β: Church numerals}\label{sec:lambda-church}

The rule $(\lambda x\,s)\,t\Rightarrow\subst{s}{t}{x}$ can copy $t$ many times or throw it away. Even when
$s$ and $t$ are both normal, the result need not be, because $x$ may sit in head position
\book{102}. So normalisation is a model of computation. Let the sort $\typ{i}$ have \emph{no constants}.
Then the closed λ-normal forms of type $(\typ{i}\to\typ{i})\to\typ{i}\to\typ{i}$ are exactly
the \emph{Church numerals} $\ulcorner n\urcorner=\lambda f\,\lambda x\,(f^n\,x)$ \book{102}. Arithmetic
is then definable \book{103}:
\[
  \mathit{succ}=\lambda n\,\lambda f\,\lambda x\,f\,(n\,f\,x),\quad
  \mathit{add}=\lambda n\,\lambda m\,\lambda f\,\lambda x\,n\,f\,(m\,f\,x),\quad
  \mathit{mul}=\lambda n\,\lambda m\,\lambda f\,\lambda x\,n\,(m\,f)\,x .
\]
λProlog prints λ-normal forms, so it can run these functions (footnote: the ambient signature has the
sort \lp{i}) \book{103}:
\begin{lprolog}
?- N = ((n\m\f\x\ n (m f) x) ((f:i -> i)\x\ f (f x)) (f\x\ f (f x))).
N = f\x\ f (f (f (f x)))
?-
\end{lprolog}
The book adds that the functions definable this way are \enquote{limited to the polynomial ones}
\book{103}. (The phrasing is informal. The precise class, a result outside the book due to
Schwichtenberg and Statman, is the \emph{extended} polynomials, which also allow a test for zero.)

\paragraph{Blow-up.} Define the \emph{size} of a term as its number of applications. Normalising can
increase the size enormously. The book's series \book{104} is
\begin{lprolog}
(g\e\ e)           (e\f\ e (e f)) (f\x\ f (f x)).
(g\e\ g e)         (e\f\ e (e f)) (f\x\ f (f x)).
(g\e\ g (g e))     (e\f\ e (e f)) (f\x\ f (f x)).
(g\e\ g (g (g e))) (e\f\ e (e f)) (f\x\ f (f x)).
\end{lprolog}
with \lp{x : i}, \lp{f : i -> i}, \lp{e : (i -> i) -> i -> i} and
\lp{g : ((i -> i) -> i -> i) -> (i -> i) -> i -> i}. The \lp{g\e\} parts are Church numerals one type
level up. \lp{g\e\ g (g e)} has a \emph{fourth-order} type. The $(n{+}1)$-th term has size $n+6$, and its
normal forms are the numerals 2, 4, 16 and 256 \book{104}.

\begin{remark}[An erratum in the size formula]\label{rem:lambda-erratum-size}
The book says the $(n{+}1)$-th normal form has size $2^{2^{\cdot^{\cdot^{2}}}}$ with $n+1$ twos
\book{104}. The term \lp{e\f\ e (e f)} squares a numeral, and the outer numeral applies it $n$ times to
$\ulcorner 2\urcorner$. So the normal form is $\ulcorner 2^{2^n}\urcorner$, of size $2^{2^n}$: 2, 4, 16,
256, 65536. This matches the book's own \enquote{256} for the fourth term. A right-nested tower with
four twos is 65536, though, so the printed formula is right only if the tower is read
left-nested, $((2^2)^2)^{\cdots}$. The run below confirms the numbers. The point of the passage is
unaffected: sizes can explode (Statman \cite{Statman1979} showed that λ-conversion is not elementary
recursive, as the book's notes cite \book{116}).
\end{remark}

\begin{kanren}[de Bruijn terms and a β-normaliser in plain Scheme]
With de Bruijn indices \cite{deBruijn1972} a bound variable is the number of binders between its
occurrence and its binder. α-equivalent terms are then \emph{the same data structure}, and
capture-avoidance turns into the index arithmetic of \scm{shift}. The following illustration is plain
Scheme, not a proposed API (the design is in \cref{ch:design}). It converts named terms to de Bruijn
form and computes β-normal forms (file \texttt{norm-core.scm}):
\begin{scheme}
;; A tiny β-normaliser: de Bruijn inside, names outside.
;; Input: (λ (x y ...) body), applications (f a b ...), symbols.
(define (->db t env)                    ; named term -> de Bruijn
  (cond
    ((symbol? t) (let ((i (index-of t env))) (if i i t)))
    ((eq? (car t) 'λ)
     (let loop ((xs (cadr t)) (env env))
       (if (null? xs)
           (->db (caddr t) env)
           (list 'lam (loop (cdr xs) (cons (car xs) env))))))
    (else
     (let loop ((f (->db (car t) env)) (as (cdr t)))
       (if (null? as)
           f
           (loop (list 'app f (->db (car as) env)) (cdr as)))))))
(define (index-of x l)
  (let loop ((l l) (i 0))
    (cond ((null? l) #f)
          ((eq? (car l) x) i)
          (else (loop (cdr l) (+ i 1))))))
(define (shift d c t)                   ; add d to the indices >= c
  (cond ((integer? t) (if (>= t c) (+ t d) t))
        ((symbol? t) t)
        ((eq? (car t) 'lam) (list 'lam (shift d (+ c 1) (cadr t))))
        (else (list 'app (shift d c (cadr t)) (shift d c (caddr t))))))
(define (subst t j s)                   ; t[s/j], capture-free
  (cond ((integer? t) (if (= t j) s t))
        ((symbol? t) t)
        ((eq? (car t) 'lam)
         (list 'lam (subst (cadr t) (+ j 1) (shift 1 0 s))))
        (else (list 'app (subst (cadr t) j s) (subst (caddr t) j s)))))
(define (beta body arg) (shift -1 0 (subst body 0 (shift 1 0 arg))))
(define (nf t)                          ; β-normal form
  (cond ((or (integer? t) (symbol? t)) t)
        ((eq? (car t) 'lam) (list 'lam (nf (cadr t))))
        (else (let ((f (nf (cadr t))))
                (if (and (pair? f) (eq? (car f) 'lam))
                    (nf (beta (cadr f) (caddr t)))
                    (list 'app f (nf (caddr t))))))))
(define (size t)                        ; number of applications
  (cond ((not (pair? t)) 0)
        ((eq? (car t) 'lam) (size (cadr t)))
        (else (+ 1 (size (cadr t)) (size (caddr t))))))
(define (church t)                      ; λf λx f^n x -> n, else #f
  (and (pair? t) (eq? (car t) 'lam)
       (pair? (cadr t)) (eq? (car (cadr t)) 'lam)
       (let loop ((b (cadr (cadr t))) (n 0))
         (cond ((eqv? b 0) n)
               ((and (pair? b) (eq? (car b) 'app) (eqv? (cadr b) 1))
                (loop (caddr b) (+ n 1)))
               (else #f)))))
(define (run t) (church (nf (->db t '()))))
\end{scheme}
Running Church arithmetic and the blow-up series (file \texttt{church.scm}):
\begin{scheme}
(include "norm-core.scm")
(define two '(λ (f x) (f (f x))))
(define three '(λ (f x) (f (f (f x)))))
(define (show . xs) (display xs) (newline))
(show 'succ-3  (run `((λ (n f x) (f (n f x))) ,three)))
(show 'add-2-3 (run `((λ (n m f x) (n f (m f x))) ,two ,three)))
(show 'mul-2-2 (run `((λ (n m f x) (n (m f) x)) ,two ,two)))
(define (series n)        ; (g\e\ g^n e) (e\f\ e (e f)) (f\x\ f (f x))
  `((λ (g e) ,(let loop ((k n)) (if (= k 0) 'e `(g ,(loop (- k 1))))))
    (λ (e f) (e (e f)))
    ,two))
(do ((n 0 (+ n 1))) ((> n 4))
  (let* ((t (->db (series n) '())) (v (nf t)))
    (show 'n n 'size (size t) 'nf-size (size v) 'numeral (church v))))
\end{scheme}
\begin{answers}
(succ-3 4)
(add-2-3 5)
(mul-2-2 4)
(n 0 size 6 nf-size 2 numeral 2)
(n 1 size 7 nf-size 4 numeral 4)
(n 2 size 8 nf-size 16 numeral 16)
(n 3 size 9 nf-size 256 numeral 256)
(n 4 size 10 nf-size 65536 numeral 65536)
\end{answers}
The input sizes are $n+6$, as the book says, and the normal-form sizes are $2^{2^n}$, as
\cref{rem:lambda-erratum-size} says. For the unifier, the useful lesson is that a de Bruijn λ-term
built from Scheme records or lists is \emph{already} handled by \micro's structural \scm{=°} up to α.
What \scm{=°} cannot do is β: binding a logic variable that sits in head position creates a redex, so
some walk-and-normalise step has to be added (\cref{def:module-hnf}).
\end{kanren}

\paragraph{Why logic programs do not explode.} The book gives three reasons \book{104--105}. (1)
Integers and data structures are built-in or made from non-logical constructors, not Church-encoded.
(2) Bound variables usually have primitive (order-0) types. If every bound variable in a term has a
primitive type and $s,t$ are λ-normal, then $\subst{s}{t}{x}$ is λ-normal: it may contain many copies of
$t$, but there are no new redexes. (3) In the sublanguage $\Llam$ (\cref{ch:lambdatree}), β is
restricted to \emph{β\textsubscript{0}-contraction} $(\lambda x\,s)\,x\Rightarrow s$, which only renames,
and normalisation then \emph{shrinks} terms.

\begin{remark}[Order of a type]
The book uses \enquote{$n$-th order} without giving a formula. The usual definition, consistent with all
its uses, is $\mathrm{ord}(\tau)=0$ for primitive $\tau$ and
$\mathrm{ord}(\tau_1\to\tau_2)=\max(\mathrm{ord}(\tau_1)+1,\mathrm{ord}(\tau_2))$. Checks: the type
$(\typ{i}\to\typ{i})\to\typ{i}\to\typ{i}$ is called second-order \book{102}, and the type of
\lp{g\e\ g (g e)} fourth-order \book{104}. A unification problem whose variables have order at most
2 is called \emph{third-order} \book{116}.
\end{remark}

% ======================================================================================
\section{Unification problems as quantified equalities}\label{sec:lambda-unifprob}

Section~1.5 of the book (\cref{ch:intro}) defined a unification problem as a multiset $\{t_1=s_1,\dots,t_p=s_p\}$ and related it
to the formula $\forall y_1\dots\forall y_m\,\exists x_1\dots\exists x_n\,[t_1=s_1\wedge\dots\wedge
t_p=s_p]$, with the $y$'s the constants and the $x$'s the logic variables \book{105}. Following Miller
\cite{Miller1991,Miller1992}, the book now makes the \emph{formula} the official presentation and
allows any quantifier prefix.

\begin{definition}[Generalised unification problem {\book{105}}]\label{def:lambda-unifprob}
A unification problem is a formula
\[
  Q_1x_1\dots Q_nx_n\,[\,t_1=s_1\wedge\dots\wedge t_m=s_m\,]\qquad(n,m\ge 0,\ Q_i\in\{\forall,\exists\}),
\]
with $t_j$ and $s_j$ λ-terms of the same type. The bound variables $x_i$ may have types of any order.
\end{definition}

Why mix the quantifiers? The book gives two reasons \book{106}, each with a query that fails.

\emph{First, \fohh{} proof search creates mixed prefixes.} A \lp{pi} in a goal introduces an
eigenvariable \emph{after} the logic variables already present:
\begin{lprolog}
?- pi y\ X = y.
no
?-
\end{lprolog}
The problem is $\exists x\,\forall y\,[x=y]$. It fails because no instance of \lp{X} may contain the
eigenvariable that stands for $y$.

\emph{Second, binders in terms behave like inner universal quantifiers.}
\begin{lprolog}
?- (y\ X) = (y\ y).
no
?-
\end{lprolog}
The problem is $\exists x\,[(\lambda y\,x)=(\lambda y\,y)]$. The equation $(\lambda y\,t)=(\lambda y\,s)$
is logically equivalent to $\forall y\,[t=s]$, so this is again $\exists x\,\forall y\,[x=y]$, and it
fails for the same reason: \emph{no instance of \lp{X} may contain a variable bound by an abstraction
inside whose body \lp{X} occurs} \book{106}.

\begin{tutorial}[reading a quantifier prefix]
Every problem in this chapter can be read with one rule.
\begin{quote}
\emph{An existential variable may be instantiated only with a term built from symbols quantified to its
\textbf{left}} (together with fresh logic variables when one is content with a unifier rather than a
solution, \cref{sec:lambda-solutions}).
\end{quote}
Here is where each quantifier of a λProlog problem comes from.
\begin{itemize}[nosep]
\item \textbf{Constants of the signature} are outermost universals. They are to the left of everything,
      so every logic variable may use them.
\item \textbf{Free variables of the query} (\lp{X}) are existentials.
\item \textbf{A \lp{pi} in a goal} adds a universal \emph{at that point}, to the right of the logic
      variables that already exist. Variables created later, for example by backchaining inside the
      scope of the \lp{pi}, sit to its right and may use it (\cref{ch:hh}).
\item \textbf{A λ inside an equation} is an innermost universal around that equation.
\end{itemize}
Four readings follow, with $g:\typ{i}\to\typ{i}\to\typ{i}$:
\begin{center}
\begin{tabular}{@{}lll@{}}
\toprule
problem & what the ∃ may use & verdict\\
\midrule
$\forall y\,\exists x\,[x=y]$ & $y$ & solved by $x:=y$\\
$\exists x\,\forall y\,[x=y]$ & nothing that mentions $y$ & no solution\\
$\forall a\,\exists F\,[(F\,a)=(g\,a\,a)]$ & $g$ and $a$ & four solutions\\
$\exists F\,\forall a\,[(F\,a)=(g\,a\,a)]$ & $g$ only & one solution\\
\bottomrule
\end{tabular}
\end{center}
The last two rows are worked out in \cref{sec:lambda-examples}. The same mechanism explains the
second failing query. In \lp{(y\ X) = (y\ y)} the logic variable \lp{X} was created
\emph{outside} the binder, so it is to the left of the inner $\forall y$. It may be instantiated with
anything except a term that mentions that $y$.
\end{tutorial}

\paragraph{From ∀∃ to ∀∃∀.} In \fohc{} the natural class is ∀∃ problems: constants first, then logic
variables. With λ-terms the natural class is \textbf{∀∃∀} \book{106}. The outer ∀'s are the
constants. The inner ∀'s come from the binders of the equations: an equation at type $\tau\to\sigma$
is replaced by $\forall(w{:}\tau)$ around an equation at type $\sigma$, and those quantifiers are moved
out over the whole conjunction.

\begin{example}[Turning λ-equations into a ∀∃∀ prefix {\book{106--107}}]\label{ex:lambda-prenex}
Let $f:\typ{i}\to\typ{i}$ and $a:\typ{i}$, and consider
\[
  \exists G\,\exists H\,[(\lambda w\,f\,a)=(\lambda w\,f\,(G\,w))\;\wedge\;(\lambda w\,f\,w)=(\lambda w\,H\,(H\,w))],
\]
where $G,H:\typ{i}\to\typ{i}$. The transformation has three steps.
\begin{enumerate}[nosep]
\item Make the constants explicit:
  $\forall f\,\forall a\,\exists G\,\exists H\,[\dots]$.
\item Replace each λ-equation by a quantified one:
  \[\forall f\,\forall a\,\exists G\,\exists H\,[\forall w\,[(f\,a)=(f\,(G\,w))]\wedge\forall w\,[(f\,w)=(H\,(H\,w))]].\]
\item Prenex (one possibility):
  \[\forall f\,\forall a\,\exists G\,\exists H\,\forall w\,\forall w'\,[(f\,a)=(f\,(G\,w))\wedge(f\,w')=(H\,(H\,w'))].\]
\end{enumerate}
Moving quantifiers to the prefix can introduce dependencies that are not really part of the problem:
$\forall w$ and $\forall w'$ now both scope over the whole conjunction, although each was attached to one equation. The book calls these \enquote{logically harmless}
and assumes prenex form from now on \book{107}.
\end{example}
The example only illustrates the transformation. By a short check of ours (the book does not
discuss it), it has no solution. The first equation forces $G=\lambda w\,a$. But a closed $H$ over
$f,a$ is $\lambda z\,f^k z$ or $\lambda z\,f^k a$, so $H\,(H\,w')$ is $f^{2k}w'$ or $f^k a$, never $f\,w'$.

\begin{kanren}[miniKanren has only ∃, and binders leak]
\scm{fresh°} is an existential, and there is no universal. The second failing query shows what this
costs once binders are present. Encode \lp{(y\ X) = (y\ y)} with de Bruijn records (illustrative
records, file \texttt{scope.scm}):
\begin{scheme}
(import (aux kanren micro) (aux base) (aux stream))
;; de Bruijn records: α-equivalent terms are equal data
(define-record lam body)
(define-record dbi n)
(define (show x) (display x) (newline))
;; λx.λy.x and λu.λv.u are the same record tree
(show (°->list #f (fresh° (q) (=° (make-lam (make-lam (make-dbi 1)))
                                   (make-lam (make-lam (make-dbi 1))))
                               (=° q 'equal))))
;; ?- (y\ X) = (y\ y).   λProlog answers no; plain =° answers yes
(show (°->list #f (fresh° (q) (=° (make-lam q) (make-lam (make-dbi 0))))))
;; ∃x[x = x]: the answer is a unifier, not a closed witness
(show (°->list #f (fresh° (q) (=° q q))))
\end{scheme}
\begin{answers}
((λ () (quote equal)))
((λ () (make-record-instance (quote dbi) 0)))
((λ (α) α))
\end{answers}
The first line is what we want: α for free. The second is a \emph{wrong yes}. \scm{q} is bound to a
dangling index, which is exactly the forbidden capture of the inner $\forall y$. (The third line is
discussed in \cref{sec:lambda-solutions}.) For the prefix discipline, a miniKanren therefore needs
(i) eigenvariables, that is, fresh constants that are never bound, and (ii) a \emph{scope check} on
every binding: a variable may not receive a symbol quantified to its right. In \micro{} the monotone
\scm{vars-count} index already records the order in which things are created, which is almost the
prefix order. The open question is how a variable inherits the restriction of the variable it is
bound to. λProlog implementations answer it with \emph{universe levels}
\cite{NadathurJayaramanKwon1995,DunchevEtAl2015}, and \cref{ch:design} spells out the proposal.
\end{kanren}

% ======================================================================================
\section{Simplifying prefixes: raising}\label{sec:lambda-raising}

A prefix can alternate as often as it likes. Raising is the tool that brings every prefix to the
∀∃∀ shape, or even to ∃∀. It moves an existential to the left of a universal and pays for it by making
the existential's type higher-order.

\begin{proposition}[Raising {\book{107}}]\label{prop:lambda-raising}
Let $x:\tau$, $y:\sigma$ and $h:\sigma\to\tau$. Then
\[
  \forall y\,\exists x\,G\quad\text{is provable iff}\quad\exists h\,\forall y\,\subst{G}{(h\,y)}{x}.
\]
\end{proposition}
\begin{proof}[Proof (the book's argument)]
($\Rightarrow$) A proof of the left formula replaces $y$ by a new constant $c$ and $x$ by a term $t$,
which may contain $c$. Prove the right formula with $h:=\lambda c\,t$ and $y:=c$. Note that $c$ is not
free in $\lambda c\,t$, as the order $\exists h\,\forall y$ requires. ($\Leftarrow$) A proof of the right
formula replaces $h$ by some $s$ and $y$ by a new $c$, and $c$ cannot be free in $s$. Prove the left
formula with $y:=c$ and $x:=(s\,c)$.
\end{proof}

Applied repeatedly, raising moves every existential to the left of the universals before it. As a rule,
for a block of universals:
\[
  \forall y_1{:}\sigma_1\dots\forall y_k{:}\sigma_k\;\exists x{:}\tau\;G
  \;\rightsquigarrow\;
  \exists h{:}\sigma_1\to\dots\to\sigma_k\to\tau\;\forall y_1\dots\forall y_k\;\subst{G}{(h\,y_1\dots y_k)}{x}.
\]
The \enquote{cost} is that the type of the existential is \emph{raised} by the types of the universals it
crosses \book{108}.

\begin{tutorial}[raising, step by step]
\textbf{Step 0: a problem whose answer depends on order.} Let $g:\typ{i}\to\typ{i}\to\typ{i}$ and
$a:\typ{i}$ be constants, and compare
\[
  (\mathrm{P}_1)\;\forall w\,\exists z\,[z=g\,a\,w]
  \qquad\text{and}\qquad
  (\mathrm{P}_2)\;\exists z\,\forall w\,[z=g\,a\,w].
\]
By the reading rule, in P\textsubscript{1} $z$ may use $w$ and $z:=g\,a\,w$ solves it. In
P\textsubscript{2} $z$ may not use $w$, and no term without $w$ equals $g\,a\,w$, so P\textsubscript{2}
has no solution.

\textbf{Step 1: raise P\textsubscript{1}.} Move $\exists z$ across $\forall w$. The new variable
$h:\typ{i}\to\typ{i}$ replaces $z$ by $(h\,w)$:
\[
  \exists h\,\forall w\,[h\,w=g\,a\,w].
\]
Now $h$ comes first and may not mention $w$. But it is \emph{applied} to $w$, so it gets $w$ back as an
argument. The solution is $h:=\lambda c\,g\,a\,c$, and the check is $(\lambda c\,g\,a\,c)\,w\to_\beta g\,a\,w$.
Translating back with the proof of \cref{prop:lambda-raising} gives $z:=h\,w=g\,a\,w$, the solution of
Step~0. P\textsubscript{2} cannot be raised, because its existential is already on the left. The
difference between the two problems is now visible in the syntax: $h$ is applied to $w$ in one, and
$z$ stands alone in the other.

\textbf{Step 2: raise over constants as well (∃∀ form).} Constants are outer universals, so one can
also raise across them. Take $\forall a\,\exists F\,[(F\,a)=(g\,a\,a)]$ with $F:\typ{i}\to\typ{i}$.
Raising $F$ across $\forall a$ gives $F':\typ{i}\to\typ{i}\to\typ{i}$, with $F$ replaced by $(F'\,a)$:
\[
  \exists F'\,\forall a\,[(F'\,a\,a)=(g\,a\,a)].
\]
$F'$ may not use $a$. It must be $\lambda y\,\lambda x\,(g\,s_1\,s_2)$ with each $s_i\in\{x,y\}$, so there
are \emph{four} solutions. Mapped back by $F:=F'\,a$, they are the four solutions of the original
problem, which \cref{sec:lambda-examples} lists. For instance, $F'=\lambda y\,\lambda x\,g\,y\,x$ gives
$F=\lambda x\,g\,a\,x$.

\textbf{What to remember.} Raising is an exact logical equivalence. It replaces \enquote{may use the
universals to my left} by \enquote{is a closed function applied to those universals}. A logic variable
under binders, \lp{X} inside \lp{y\ ...}, can always be read as a raised variable applied to the bound
variables, \lp{(H y)}. When the arguments are \emph{distinct} bound variables, as in $h\,w$, the
application is a \emph{higher-order pattern} (\cref{ch:lambdatree,ch:hounif}), and pattern problems
have most general unifiers. In Step~2 the argument $a$ is repeated, $F'\,a\,a$, which is not a pattern,
and indeed there is no most general solution.
\end{tutorial}

\paragraph{Skolemization, the dual.} Skolemizing an \emph{assumption} $\forall x\,\exists y\,D$ with
$x:\sigma$ and $y:\tau$ gives $\forall x\,\subst{D}{(f\,x)}{y}$, where $f:\sigma\to\tau$ is a new
\emph{Skolem constant}. For a \emph{goal} the roles swap: $\exists x\,\forall y\,G$ becomes
$\exists x\,\subst{G}{(f\,x)}{y}$, that is, $\forall f\,\exists x\,\subst{G}{(f\,x)}{y}$ \book{108}. The
two techniques are dual in two senses. Skolemization moves an \emph{∃ to a smaller scope} and
introduces a new \emph{constant} of raised type. Raising moves a \emph{∀ to a smaller scope} and
introduces a new \emph{existential variable} of raised type. Logically they differ: \enquote{relating
unifiers for Skolemized unification problems to original, un-Skolemized unification problems is
problematic, especially when unification involves variables of higher-order type} \book{108}. Raising,
in contrast, is an equivalence (\cref{prop:lambda-raising}), which is why the book uses it to normalise
prefixes.

\begin{pitfall}[raising creates higher-order variables]
Raising a first-order problem gives a higher-order one: in Step~1, $z:\typ{i}$ became
$h:\typ{i}\to\typ{i}$. A first-order miniKanren with ∀-goals cannot implement scope by raising unless
it also has λ-terms and β. Prenexing (\cref{ex:lambda-prenex}) has a related cost: it may add spurious
dependencies. These are logically harmless, but a scope-tracking implementation becomes less
permissive than it needs to be.
\end{pitfall}

\begin{kanren}[raising versus levels]
Raising is how the book \emph{explains} logic variables under \lp{pi} and λ. Implementations usually
do not raise up front. Nadathur, Jayaraman and Kwon tag each variable and constant with a \emph{universe index}
\cite{NadathurJayaramanKwon1995}. A binding is allowed only if the term's constants have tags no larger
than the variable's, and the tags of the term's variables are lowered to match. ELPI uses the same
\emph{levels} \cite{DunchevEtAl2015}. Nadathur and Linnell \cite{NadathurLinnell2005} show how to erase
the prefix and raise \emph{on the fly}, only over the universals that actually occur. For \micro{},
the choice is between raising, which needs only terms (λ, application, β) but makes every existential
higher-order, and levels, which need a new per-variable piece of state and keep variables first-order.
The two must agree on the examples of this chapter. \Cref{ch:design} makes the choice.
\end{kanren}

% ======================================================================================
\section{Unifiers, solutions and empty types}\label{sec:lambda-solutions}

Once a problem is a formula, solving it should mean proving it. The connection is subtler than it
looks. The higher-order version of \rulename{$\exists$R} \book{108--109} is
\begin{mathpar}
\inferrule{\seq{\Sig}{\Prog}{\subst{B}{t}{x}} \\ \Sig;\emptyset\Vdash t:\tau}{\seq{\Sig}{\Prog}{\exists_\tau x\,B}}\;\rulename{$\exists$R}
\end{mathpar}
The witness $t$ must be \emph{closed}, because its context is empty. Proving $\exists x\,[x=x]$ with
$x:\typ{i}$ therefore needs a closed term of type $\typ{i}$, and whether one exists depends on
$\Sig$. In the same way, $\exists x\,\exists y\,\exists z\,[x=(f\,y)\wedge y=(f\,z)]$ with
$f:\typ{i}\to\typ{i}$ has the most general unifier $\{\langle X,f\,(f\,Z)\rangle,\langle Y,f\,Z\rangle\}$
in the sense of the book's Section~1.5 (\cref{ch:intro}). The formula, however, is provable only if some closed term of type
$\typ{i}$ exists \book{109}.

\begin{definition}[Legitimate substitution, unifier, solution {\book{109}}]\label{def:lambda-solution}
Consider the ∀∃∀ problem
$\forall u_1\dots\forall u_p\,\exists x_1\dots\exists x_n\,\forall w_1\dots\forall w_q\,[t_1=s_1\wedge\dots\wedge t_m=s_m]$
and let $\theta$ be a substitution for $x_1,\dots,x_n$.
\begin{itemize}[nosep]
\item $\theta$ is \emph{legitimate} if no $w_j$ occurs free in any $\theta(x_i)$.
\item A legitimate $\theta$ is a \emph{unifier} if $\theta(t_j)$ and $\theta(s_j)$ are λ-convertible for
      every $j$.
\item A unifier is a \emph{solution} if the only free variables in the terms $\theta(x_i)$ are among
      $u_1,\dots,u_p$.
\end{itemize}
A problem can have a unifier and no solution. \emph{It is a solution that makes the quantified formula
provable.}
\end{definition}

In the book's Section~1.5 (\cref{ch:intro}) \enquote{solution} and \enquote{unifier} meant the same thing. They have to be separated
now because, unlike first-order classical logic, higher-order intuitionistic logic does not assume that
every type is \emph{inhabited}. The Church numerals show why this matters. They are the closed terms of
type $(\typ{i}\to\typ{i})\to\typ{i}\to\typ{i}$ only if $\typ{i}$ has \emph{no} closed term $t$, since
otherwise $\lambda f\,\lambda x\,t$ would be a further, non-numeral inhabitant \book{109--110}.

\begin{example}[Unifier but no solution]
Let $\Sig$ contain only $u,v:\typ{i}\to\typ{i}$, so $\typ{i}$ is empty. The problem $\exists x\,[x=x]$
has the unifier $\{x\mapsto X\}$ with $X$ a fresh logic variable. It has no solution, because there is
no closed term of type $\typ{i}$. Add a constant $a:\typ{i}$ and $\{x\mapsto a\}$ becomes a solution.
\end{example}

\begin{kanren}[miniKanren answers are unifiers]
The third answer of \texttt{scope.scm} above, \scm{((λ (α) α))} for $\exists x\,[x=x]$, is a
\emph{unifier} in the sense of \cref{def:lambda-solution}. The reified \scm{α} is a fresh logic
variable, not a closed witness. This is also λProlog's practice: answers may contain variables, and
leftover flex–flex constraints (\cref{sec:lambda-vocab}). An untyped miniKanren cannot see empty types
at all. Obtaining \emph{solutions} would need an explicit generator of closed inhabitants of a type
over a signature, itself a proof-search relation. \Cref{sec:lambda-undecidable} shows that the gap
between unifiers and solutions is exactly where undecidability lives.
\end{kanren}

% ======================================================================================
\section{Example problems and their solutions}\label{sec:lambda-examples}

These examples \book{110--111} are the core intuition pump of the chapter. Each one shows a way in
which higher-order unification differs from the first-order kind.

\subsection{Quantifier order changes the answers}\label{sec:lambda-order}

Let $g:\typ{i}\to\typ{i}\to\typ{i}$ be a constructor of $\typ{i}$.
\begin{itemize}
\item $\forall a\,\exists F\,[(F\,a)=(g\,a\,a)]$ has \textbf{four} solutions:
      $F\in\{\lambda x\,g\,a\,a,\ \lambda x\,g\,x\,a,\ \lambda x\,g\,a\,x,\ \lambda x\,g\,x\,x\}$.
\item $\exists F\,\forall a\,[(F\,a)=(g\,a\,a)]$ has the \textbf{unique} solution $F=\lambda a\,g\,a\,a$.
\end{itemize}
\emph{Walk-through.} $F:\typ{i}\to\typ{i}$ is closed, so $F=\lambda x\,b$ with $b$ built from $g$, $x$ and
whatever $F$ may use. The β-normal form of $F\,a$ is $\subst{b}{a}{x}$, and it must be $g\,a\,a$. So $b$ is
$g\,s_1\,s_2$ where each $s_i$ either \emph{is} $a$ or \emph{becomes} $a$ under the substitution, that is,
$s_i\in\{a,x\}$. In the first problem $a$ is to the left of $F$, so both choices are allowed for each
position: $2\times 2=4$. In the second, $F$ may not mention $a$, only $s_i=x$ remains, and there is one
solution. (\Cref{sec:lambda-raising} reached the same count of four by raising.)

\begin{kanren}[a first-order simulation of both prefixes]
In this example $F$ is applied to a single argument and its body is first-order, so the problem can be
simulated in today's \micro. Represent $F$'s body as a first-order term in which \scm{x} stands for the
bound variable, and relate a body to its instance $\subst{b}{a}{x}$ with a relation \scm{inst°}. The
universal $a$ is the symbol \scm{a}. The ∃F∀a prefix, \enquote{$F$ may not mention $a$}, is then the
constraint \scm{(absent° 'a body)}. This is an illustration, not general higher-order unification (file
\texttt{order.scm}):
\begin{scheme}
(import (aux kanren micro) (aux base) (aux stream))
;; F = λx.body, body a first-order term over g, a and x = (dbi 0)
(define-record dbi n)
(set-record-printer! dbi (λ (r port) (display "x" port)))
(define x (make-dbi 0))
(define (show x) (display x) (newline))
;; (inst° body arg out): out = body[arg/x], the normal form of (F arg)
(define-relation (inst° body arg out)
  (cond°
    ((=° body x) (=° out arg))
    ((symbol° body) (=° out body))
    ((null° body) (null° out))
    ((fresh° (h t h* t*)
       (=° body (cons h t)) (=° out (cons h* t*))
       (inst° h arg h*) (inst° t arg t*)))))
;; ∀a ∃F [(F a) = (g a a)]: F's body may mention the constant a
(show (°->list/ground
        (fresh° (body) (inst° body 'a '(g a a)))))
;; ∃F ∀a [(F a) = (g a a)]: a comes after F, F's body must avoid it
(show (°->list/ground
        (fresh° (body) (absent° 'a body) (inst° body 'a '(g a a)))))
\end{scheme}
\begin{answers}
((g x x) (g a x) (g x a) (g a a))
((g x x))
\end{answers}
These are the book's four answers and its one answer. Two points carry over to the real design. First,
the scope check is a \emph{constraint} on what a variable may contain, like \scm{absent°}, and it is
attached at the moment the variable is created. Second, a relation that runs substitution
\emph{backwards} (\scm{inst°} with \scm{out} known) is what imitation and projection do in Huet's
procedure (\cref{ch:hounif}): each \scm{a} in the output is either imitated or projected from the
argument.
\end{kanren}

\subsection{Words over a two-letter alphabet}\label{sec:lambda-words}

Let the only constructors with target $\typ{i}$ be $u,v:\typ{i}\to\typ{i}$. Then $\typ{i}$ has no closed
terms, and the closed terms of type $\typ{i}\to\typ{i}$ are \emph{words}: $\lambda w\,w$ is the empty word
$\varepsilon$, $\lambda w\,u\,w$ is \enquote{u}, and $\lambda w\,u\,(v\,(u\,w))$ is \enquote{uvu}. Function
composition is concatenation \book{110}. Three problems, all with $F,G:\typ{i}\to\typ{i}$:
\begin{enumerate}
\item $\exists F\,\exists G\,[(\lambda w\,F\,(G\,w))=(\lambda w\,u\,(v\,(u\,w)))]$ has \textbf{four}
      solutions, the four ways of splitting uvu:
      $(F,G)\in\{(\mathrm{uvu},\varepsilon),(\mathrm{uv},\mathrm{u}),(\mathrm{u},\mathrm{vu}),(\varepsilon,\mathrm{uvu})\}$.
      Replacing the right-hand side by any other closed term asks for all ways of splitting that word.
\item $\exists F\,[(\lambda w\,u\,(F\,(u\,w)))=(\lambda w\,u\,(v\,(v\,(u\,w))))]$ has the \textbf{only}
      solution $F=\lambda w\,v\,(v\,w)$. The left-hand side is a \emph{pattern} that checks whether a word
      starts and ends with u.
\item $\exists F\,[(\lambda w\,u\,(F\,w))=(\lambda w\,F\,(u\,w))]$ is the word equation
      $\mathrm{u}\cdot F=F\cdot\mathrm{u}$. Its solutions are \textbf{infinitely many}:
      $F\in\{\lambda w\,w,\ \lambda w\,u\,w,\ \lambda w\,u\,(u\,w),\dots\}$, that is, $\mathrm{u}^k$ for every
      $k\ge 0$. They are closed and pairwise distinct, so none is an instance of another: \emph{there is
      no most general solution}.
\end{enumerate}

\begin{kanren}[words are lists, composition is append]
Because $\typ{i}$ is empty, a closed term of type $\typ{i}\to\typ{i}$ \emph{is} a word, and in this one case
higher-order unification becomes ordinary relational list programming (file \texttt{words.scm}):
\begin{scheme}
(import (aux kanren micro) (aux base) (aux stream))
;; closed terms of type i -> i over u, v : i -> i are words,
;; and composition is concatenation
(define-relation (append° r s rs)
  (cond°
    ((null° r) (=° s rs))
    ((fresh° (a d c) (cons° a d r) (cons° a c rs) (append° d s c)))))
(define (show x) (display x) (newline))
;; ∃F ∃G [(λw F (G w)) = (λw u (v (u w)))]
(show (μkanren-run (q -1 #t)
        (fresh° (F G) (append° F G '(u v u)) (=° q (list F G)))))
;; ∃F [(λw u (F (u w))) = (λw u (v (v (u w))))]
(show (μkanren-run (q -1 #t)
        (fresh° (uF) (append° '(u) q uF) (append° uF '(u) '(u v v u)))))
;; ∃F [(λw u (F w)) = (λw (F (u w)))]: five of infinitely many
(show (μkanren-run (q 5 #t)
        (fresh° (W) (append° '(u) q W) (append° q '(u) W))))
\end{scheme}
\begin{answers}
((() (u v u)) ((u) (v u)) ((u v) (u)) ((u v u) ()))
((v v))
(() (u) (u u) (u u u) (u u u u))
\end{answers}
The third query shows the right attitude toward problems without an mgu. A miniKanren does not need
a single answer: it returns a \emph{stream} of them, and \scm{take°} (here \scm{μkanren-run} with
$n=5$) cuts the stream. This is the shape a general higher-order unifier in \micro{} must have
(\cref{sec:lambda-undecidable}).
\end{kanren}

\subsection{Weak functionality}\label{sec:lambda-weak}

Let $\typ{i}$ have at least the constants $a,b,c,d:\typ{i}$. Then
$\exists F\,[(F\,a)=b\wedge(F\,c)=d]$ has \textbf{no solution}. The first equation is solved only by
$F=\lambda x\,b$: $F$ cannot produce $b$ by projecting its argument $a$, so it must be the constant
function. That $F$ fails the second equation \book{110--111}. A function that maps $a\mapsto b$ and
$c\mapsto d$ exists as a set-theoretic object, but it is \emph{not expressible} by these λ-terms, which
have no case analysis on constants. The book calls this \enquote{a relatively weak notion of
functionality}.

\begin{pitfall}[what a first-order intuition gets wrong]
\begin{itemize}[nosep]
\item \emph{Quantifier order is semantic}: $\forall a\,\exists F$ and $\exists F\,\forall a$ give four
      answers and one answer.
\item \emph{There may be several incomparable solutions} (four splittings of uvu), \emph{or
      infinitely many with no mgu} (u·F = F·u). A unifier returning one substitution, as
      \scm{μkanren-state-unify} does, cannot express this.
\item \emph{λ-terms are not arbitrary functions}: no case analysis, so $F\,a=b\wedge F\,c=d$ is
      unsolvable.
\end{itemize}
\end{pitfall}

% ======================================================================================
\section{Vocabulary for a unification procedure}\label{sec:lambda-vocab}

\Cref{ch:hounif} presents Huet's procedure \cite{Huet1975}. The book prepares it with some vocabulary
and a few structural observations \book{111--113}.

\begin{definition}[Binder, head, arguments, body {\book{111}}]\label{def:lambda-head}
Every β-normal term has the shape
\[
  \lambda x_1\dots\lambda x_n\,(h\;t_1\cdots t_p)\qquad(n,p\ge 0),
\]
where $h$ is a constant or a variable and every $t_i$ is β-normal. $x_1,\dots,x_n$ is the
\emph{binder}, $h$ the \emph{head}, $t_1,\dots,t_p$ the \emph{arguments}, and $(h\,t_1\cdots t_p)$ the
\emph{body}.
\end{definition}

\begin{definition}[βη-long normal form {\book{111}}]\label{def:lambda-long}
A term is in \emph{βη-long normal form} if it has the shape above, every $t_i$ is βη-long, and the body
has a non-functional type. For simply typed terms, that means a primitive type. Every β-normal form
can be η-expanded into an equal βη-long form.
\end{definition}
\begin{example}[η-expanding to long form {\book{111}}]
Let $f:(\typ{i}\to\typ{i})\to\typ{i}$. The term $f$ is β-normal but not long, since its body has
functional type. One η-expansion gives $\lambda w\,f\,w$ with $w:\typ{i}\to\typ{i}$. This is still not
long, because the argument $w$ has a non-primitive type. Expanding $w$ gives $\lambda w\,f\,(\lambda u\,w\,u)$
with $u:\typ{i}$, which is βη-long. Note that deciding when to stop requires knowing the \emph{types}.
\end{example}

\begin{definition}[Rigid and flexible {\book{111--113}}]\label{def:lambda-rigid}
In a problem over β-normal terms, an occurrence $\lambda\bar x\,(h\,t_1\cdots t_p)$ is \emph{rigid} if
$h$ is ∀-bound in the prefix (constants included) or belongs to the binder $\bar x$. It is
\emph{flexible} if $h$ is ∃-bound in the prefix. Equations are rigid–rigid, rigid–flexible,
flexible–rigid or flexible–flexible, and a problem is \emph{flexible–flexible} if all of its equations
are.
\end{definition}
The names are meant to suggest what substitution can do. Substitution never changes the head of a rigid
term: every instance has a β-normal form with the same head. It can change the head of a flexible term
\book{112}.

\paragraph{Observation 1: rigid heads must agree.} By α and η, two rigid terms of the same type can be
given the same binder: $\lambda\bar x\,(h\,t_1\cdots t_p)$ and $\lambda\bar x\,(k\,s_1\cdots s_q)$. Since
neither head can change, a unifier can exist only if $h$ and $k$ are the same symbol \book{112}. With $F,G$
logic variables and $c,d$ constants, none of the following has a unifier:
\begin{gather*}
  \{(\lambda x\,(d\,(c\,(F\,x))))=(\lambda x\,(c\,(d\,(G\,x))))\}\qquad
  \{(\lambda x\,(x\,(F\,x)))=(\lambda x\,(c\,(G\,x)))\}\\
  \{(\lambda x\,\lambda y\,(x\,(F\,x\,y)))=(\lambda x\,\lambda y\,(y\,(G\,x\,y)))\}
\end{gather*}
The first one, written precisely, is $\forall c\,\forall d\,\exists F\,\exists G\,[\lambda x\,(d\,(c\,(F\,x)))=\lambda x\,(c\,(d\,(G\,x)))]$.
Note that in the second and third problems one of the heads is a \emph{bound} variable, which counts as
rigid.

\paragraph{Observation 2: substitution cannot create captured occurrences.} If a binder variable has
no free occurrence in the body of a β-normal form, it has none in any instance either, because
substitution avoids capture. So
\[
  \{(\lambda x\,\lambda y\,x\,(F\,x\,y))=(\lambda x\,\lambda y\,G\,y\,y)\}
\]
has no unifier \book{112--113}. Every instance of the left side has head $x$. No instance of the right
side can contain the bound $x$: $G$ may be instantiated with a term containing a \emph{free} $x$, but
that is a different $x$.

\paragraph{Observation 3: flexible terms are limited only in their use of the binder.} A flexible term
can become almost any term of its type, \emph{except} in how it uses the binder variables. For
instance, $\lambda x\,\lambda y\,F\,(f\,x\,(G\,x))$ can become any $\lambda x\,\lambda y\,t$ with $x,y$ not
free in $t$: take $F:=\lambda z\,t$, which throws its argument away. Some instances also contain $x$,
such as $\lambda x\,\lambda y\,(f\,x\,x)$, obtained with $F:=\lambda z\,z$ and $G:=\lambda z\,z$. But no
instance contains $y$ \book{113}, because $y$ is passed to neither $F$ nor $G$, and by Observation~2
substitution cannot introduce it.

\paragraph{Observation 4: a flex–flex pair.} Unify $\lambda x\,\lambda y\,G\,x$ with
$\lambda x\,\lambda y\,H\,y$. Instances of the first ignore $y$, and instances of the second ignore $x$, so a
common instance is $\lambda x\,\lambda y\,t$ with neither variable free in $t$. The \enquote{most general}
unifier, derived in \cref{ch:hounif}, is $G:=\lambda w\,T$ and $H:=\lambda w\,T$, with $T$ a new variable
\book{113}. (Pattern unification calls this step \emph{pruning}, \cref{ch:hounif}.)

\begin{tutorial}[classifying heads in practice]
Take the prefix $\forall c\,\exists F\,\exists G$ and three occurrences:
\begin{center}
\begin{tabular}{@{}lll@{}}
\toprule
term & head & status\\
\midrule
$\lambda x\,(c\,(F\,x))$ & $c$, ∀-bound & rigid: every instance starts with $c$\\
$\lambda x\,(x\,(F\,x))$ & $x$, in the binder & rigid: every instance starts with the bound $x$\\
$\lambda x\,(F\,(c\,x))$ & $F$, ∃-bound & flexible: $F:=\lambda z\,z$ makes the head $c$\\
\bottomrule
\end{tabular}
\end{center}
A rigid–rigid equation either fails immediately or decomposes into equations between arguments
(\cref{ch:hounif}). A flex–rigid equation is where the search branches (imitate the rigid head, or
project an argument). A flex–flex equation always has a unifier, as in Observation~4, so Huet's
procedure \emph{postpones} it (\cref{ch:hounif}). It need not have a \emph{solution}, as
\cref{sec:lambda-undecidable} shows.
\end{tutorial}

\begin{kanren}[what the vocabulary means for the unifier]
\micro's \scm{=°} already performs \enquote{rigid–rigid}: two records or pairs with different
constructors fail at once, and the same constructor decomposes slot-wise. A de Bruijn representation
whose application node exposes head and arguments, as in \cref{def:lambda-head}, makes Observation~1
a constructor clash. Bound-variable heads (\scm{dbi} records) are rigid automatically. Everything about
\emph{flexible} heads is new: the unifier must recognise a logic variable in head position, normalise
after binding it, and either solve the equation deterministically (patterns, \cref{ch:lambdatree}) or
branch and delay it. Delayed flex–flex equations need a store in the state, next to the disequality
store, and \cref{ch:design} proposes one.
\end{kanren}

% ======================================================================================
\section{Undecidability}\label{sec:lambda-undecidable}

Derivability from \fohc{} programs is already undecidable, and so is unifiability of typed λ-terms
modulo λ-conversion \book{113}. The book proves the second fact by two reductions. It also stresses that
the specifications of later chapters produce \enquote{much tamer} problems.

\subsection{The Post correspondence problem}\label{sec:lambda-pcp}

\begin{definition}[PCP {\book{114}}]
Given $n\ge 1$ and two lists $s_1,\dots,s_n$ and $t_1,\dots,t_n$ of words over $\{u,v\}$, decide whether
there is a \emph{non-empty} sequence $i_1,\dots,i_k$ with $1\le i_j\le n$ such that
$s_{i_1}\cdots s_{i_k}=t_{i_1}\cdots t_{i_k}$. This problem is undecidable.
\end{definition}

Encode a word $r=r_1\cdots r_m$ as $\hat r=\lambda w\,(r_1\,(\cdots(r_m\,w)\cdots))$, as in
\cref{sec:lambda-words}. Concatenation is then $\lambda w\,(\hat r\,(\hat q\,w))$. Take a sort $\typ{i}$
\emph{with no constructors}, and introduce $u,v:\typ{i}\to\typ{i}$ by universal quantifiers
\book{114}:
\[
  \exists F\,\exists G\,\forall u\,\forall v\,
  [\,(F\,\hat s_1\cdots\hat s_n)=(F\,\hat t_1\cdots\hat t_n)\;\wedge\;(F\,u\cdots u)=\lambda w\,(u\,(G\,u\,w))\,],
\]
with $G:(\typ{i}\to\typ{i})\to\typ{i}\to\typ{i}$. For the lists $\mathrm{uv},\mathrm{u}$ and
$\mathrm{u},\mathrm{vu}$ \book{114}:
\[
  \exists F\,\exists G\,\forall u\,\forall v\,
  [\,(F\,(\lambda w\,u\,(v\,w))\,u)=(F\,u\,(\lambda w\,v\,(u\,w)))\;\wedge\;(F\,u\,u)=\lambda w\,(u\,(G\,u\,w))\,],
\]
where $F$ has the second-order type $(\typ{i}\to\typ{i})\to(\typ{i}\to\typ{i})\to\typ{i}\to\typ{i}$.

\begin{tutorial}[why the encoding is correct]
\textbf{Read the prefix.} $F$ and $G$ come \emph{before} $\forall u\,\forall v$, so they may not mention
$u$ or $v$. They can only rearrange their arguments. $\typ{i}$ has no constructors, so a closed $F$
must be $\lambda x_1\dots\lambda x_n\,\lambda w\,d$ with $d:\typ{i}$ built only from $x_1,\dots,x_n,w$.
The only possible shape is $d=x_{i_1}(\cdots(x_{i_k}\,w)\cdots)$ with $k\ge 0$ \book{114}.

\textbf{Read off the index sequence.} $F\,\hat s_1\cdots\hat s_n$ normalises to the word
$s_{i_1}\cdots s_{i_k}$, and the right side to $t_{i_1}\cdots t_{i_k}$. So the first equation says
exactly that $i_1,\dots,i_k$ solves the PCP, except in the degenerate case $k=0$, where both sides are
$\lambda w\,w$.

\textbf{Exclude $k=0$.} The second equation says that $F\,u\cdots u$ starts with $u$, which forces
$k\ge 1$. In the other direction, a PCP solution gives a unification solution directly
\book{114--115}. The argument extends to \emph{unifiers} and to types $\typ{i}$ with constructors
\book{115}.

\textbf{Check the instance.} For $\mathrm{uv},\mathrm{u}$ / $\mathrm{u},\mathrm{vu}$, the sequence
$1,2$ works: $\mathrm{uv}\cdot\mathrm{u}=\mathrm{u}\cdot\mathrm{vu}=\mathrm{uvu}$. This suggests
$F=\lambda x_1\,\lambda x_2\,\lambda w\,x_1\,(x_2\,w)$ and $G=\lambda y\,\lambda w\,y\,w$. $G$ may not mention
$u$ and has to obtain it from its argument. The book does not print this solution. The normaliser of
\cref{sec:lambda-normal}, with $u,v$ as free symbols, confirms both equations, and also the Diophantine
gadget of the next subsection (file \texttt{checks.scm}):
\begin{scheme}
(include "norm-core.scm")
(define (same? s t) (equal? (nf (->db s '())) (nf (->db t '()))))
(define (show . xs) (display xs) (newline))
;; PCP instance uv,u / u,vu; u and v are free symbols
(define F '(λ (x1 x2 w) (x1 (x2 w))))
(define G '(λ (y w) (y w)))
(show 'pcp (same? `(,F (λ (w) (u (v w))) u) `(,F u (λ (w) (v (u w)))))
           (same? `(,F u u) `(λ (w) (u (,G u w)))))
;; gadget for x = 1: λf λx N f (N f x) = λf λx N f (f x)
(define (numeral n)
  `(λ (f x) ,(let loop ((k n)) (if (= k 0) 'x `(f ,(loop (- k 1)))))))
(define (gadget N)
  (same? `(λ (f x) (,N f (,N f x))) `(λ (f x) (,N f (f x)))))
(show 'numerals-0..3 (map (lambda (n) (gadget (numeral n))) '(0 1 2 3)))
(show 'λf.λx.U (gadget '(λ (f x) U)))
\end{scheme}
\begin{answers}
(pcp #t #t)
(numerals-0..3 (#f #t #f #f))
(λf.λx.U #t)
\end{answers}
\end{tutorial}

Historically, Huet \cite{Huet1973} gave the PCP reduction, which shows that \emph{third-order}
unification is undecidable \book{116}.

\subsection{Hilbert's tenth problem}\label{sec:lambda-hilbert}

Hilbert's tenth problem asks for an algorithm that decides whether a system of Diophantine equations
(variables, numerals, $+$, $\times$ over $\mathbb N$) has a solution. No such algorithm exists
\book{115}. It is enough to treat equations of the forms $x=1$, $x+y=z$ and $x\times y=z$. Take a sort
$\typ{i}$ without constructors, so that the closed terms of $(\typ{i}\to\typ{i})\to\typ{i}\to\typ{i}$ are
exactly the Church numerals $\ulcorner n\urcorner$, and give every existential this type. Then
\book{115}:
\begin{enumerate}
\item The flex–flex problem $\exists N\,[(\lambda f\,\lambda x\,N\,f\,(N\,f\,x))=(\lambda f\,\lambda x\,(N\,f\,(f\,x)))]$
      has the unique solution $N=\ulcorner 1\urcorner$. With $N=\ulcorner n\urcorner$ the sides are
      $f^{2n}x$ and $f^{n+1}x$, so $2n=n+1$.
\item The flex–flex problem $\exists M\,\exists N\,\exists P\,[(\lambda f\,\lambda x\,N\,f\,(M\,f\,x))=(\lambda f\,\lambda x\,P\,f\,x)]$
      has the solution $\{(N,r),(M,s),(P,t)\}$ iff $r=\ulcorner n\urcorner$, $s=\ulcorner m\urcorner$ and
      $t=\ulcorner n+m\urcorner$ for some $n,m\ge 0$.
\item The flex–flex problem $\exists M\,\exists N\,\exists P\,[(\lambda f\,\lambda x\,N\,(M\,f)\,x)=(\lambda f\,\lambda x\,P\,f\,x)]$
      has the solution $\{(N,r),(M,s),(P,t)\}$ iff $r=\ulcorner n\urcorner$, $s=\ulcorner m\urcorner$ and
      $t=\ulcorner n\times m\urcorner$.
\end{enumerate}

\begin{remark}[An erratum in the addition gadget]
The book prints the left side of item 2 as $\lambda f\,\lambda x\,(N\,f)\,(M\,f)\,x$ \book{115}. This parses
as $((N\,f)\,(M\,f))\,x$, which is ill-typed, since $N\,f:\typ{i}\to\typ{i}$ is applied to
$M\,f:\typ{i}\to\typ{i}$. The intended term is the addition combinator of \cref{sec:lambda-church},
$\lambda f\,\lambda x\,N\,f\,(M\,f\,x)$, and that is the term shown above.
\end{remark}

\begin{theorem}[{\book{116}}]\label{thm:lambda-hilbert}
A system of such equations becomes a flex–flex problem with one existential variable per
$\mathbb N$-variable and one equation per equation. The problem has a \emph{solution} iff the system has
a solution in $\mathbb N$. The variables have order at most 2, so the problem is third-order, and the
construction extends to arbitrary Diophantine equations. Hence \textbf{it is undecidable whether a
third-order flex–flex problem has a solution.}
\end{theorem}

\paragraph{Unifiers versus solutions, once more.} The encoding \emph{always} has a unifier: pick a new
variable $U:\typ{i}$ and set every existential to $\lambda f\,\lambda x\,U$. Both sides of every gadget then
normalise to $\lambda f\,\lambda x\,U$. This term is not closed, so it is not a solution \book{116}. The run
above confirms both facts for gadget~1. Among $\ulcorner 0\urcorner,\dots,\ulcorner 3\urcorner$ only
$\ulcorner 1\urcorner$ works, and $\lambda f\,\lambda x\,U$ works too. So the undecidable question concerns
\emph{solutions}. Asking only for unifiers, a flex–flex problem is trivially unifiable, and this is why
Huet's procedure can leave flex–flex pairs for last. In the bibliographic notes \book{116--117} the book
credits this reduction to Miller \cite{Miller1992}. It differs from Goldfarb's \cite{Goldfarb1981}, which
already shows that \emph{second-order} unification, for unifiers, is undecidable.

\begin{kanren}[undecidability and the search]
Since unification is undecidable, no complete procedure terminates on every input. A complete
higher-order unifier is a \emph{semi-decision procedure}: a search, some of whose branches never
finish. In a miniKanren this is a familiar situation. Such a unifier would be a \emph{goal} that produces a
stream of unifiers. Huet's branching steps would become disjunctions, and \micro's \emph{interleaving}
\scm{orª} would explore them fairly, which a depth-first exploration does not. Flex–flex pairs, which always
have unifiers, would be left in a constraint store and reified with the answer, which then is a unifier and
not necessarily a solution. The proposal keeps only the deterministic part: pattern problems are solved by
one \scm{=°}, and non-patterns, flex–flex pairs included, wait in a store and are reported
(\cref{sec:module-wake}). Huet's branching is left to an optional search over that store, which needs
types (\cref{sec:roadmap-later-huet}).
\end{kanren}

\begin{remark}[Polymorphism]
All the formal results of the book's chapter concern \emph{simply typed} terms. Type variables are
allowed for convenience, for instance one declaration for \lp{pi}, but \emph{unification does not
instantiate type variables}: they behave like sorts, so the simply typed results still apply
\book{117}.
\end{remark}

% ======================================================================================
\section{Summary}\label{sec:lambda-summary}

\begin{itemize}
\item Terms are simply typed λ-terms (Figure~4.1). Formulas are terms of type $\typ{o}$, the quantifiers
      are constants $\forall,\exists:(A\to\typ{o})\to\typ{o}$ applied to abstractions, and λ is the
      \emph{only} binder.
\item The equality of the logic is λ-conversion (α, β, η). Bound names are unobservable, substitution
      avoids capture, and equality ignores logical equivalences.
\item Every term has a λ-normal form, unique up to α, so equality is decidable. β is powerful (Church
      numerals) and normal forms can be huge ($2^{2^n}$ in the book's series, see the erratum).
      Practical programs avoid this through built-in data, order-0 binders and $\Llam$.
\item A unification problem is a quantified formula $Q_1x_1\dots Q_nx_n\,[\bigwedge t_j=s_j]$.
      \emph{An existential may use only what is quantified to its left.} \lp{pi} and λ both introduce
      universals, which is why \lp{pi y\ X = y} and \lp{(y\ X) = (y\ y)} fail.
\item \emph{Raising} ($\forall y\,\exists x\,G\rightsquigarrow\exists h\,\forall y\,\subst{G}{(h\,y)}{x}$) is an
      exact equivalence that brings any prefix to ∀∃∀ or ∃∀, at the cost of higher-order types.
      Skolemization is its dual but does not preserve unifiers well.
\item A \emph{unifier} may leave free variables, while a \emph{solution} uses only the outer universals.
      Only solutions give proofs, and empty types make the difference real. miniKanren answers are
      unifiers.
\item Higher-order unification has several, infinitely many, or no solutions, and no mgu in general.
      Quantifier order changes the count (4 versus 1), and λ-terms cannot do case analysis.
\item Vocabulary for \cref{ch:hounif}: binder, head, arguments, body; βη-long form; rigid and
      flexible terms; rigid heads must agree; vacuous binders stay vacuous; flex–flex pairs always
      have unifiers.
\item Unifiability is undecidable (PCP, third order), and so is the existence of solutions of
      third-order flex–flex problems (Hilbert's tenth problem). Hence a miniKanren higher-order unifier
      must be a goal producing a fair stream of unifiers, with flex–flex pairs postponed.
\end{itemize}

\chapter{Higher-order programming}\label{ch:hoprog}

% local notation: the sequent of Figure 6.5 that processes a multiset of E-formulas
\newcommand{\hoprogE}[4]{#1;#2\mathrel{-\!\langle #3\rangle\!\!\longrightarrow}#4}
\newcommand{\hoprogT}[1]{\lBrack #1\rBrack}

\Cref{ch:lambda} replaced first-order terms by simply typed $\lambda$-terms and allowed quantification at
every type, including predicate types (types whose target is $\typ{o}$); it also showed that unification of
such terms is undecidable and may have infinitely many incomparable solutions. This chapter, which follows
Chapter~5 of the book \cite{MillerNadathur2012} and the core of its Chapter~6, asks the practical
question: \emph{what logic programming languages do we get, and what do they buy us?}

Two kinds of higher-order variable will appear, and it is worth keeping them apart from the start.
\begin{itemize}
\item \textbf{Predicate variables} (target type $\typ{o}$) are instantiated by relations and then
  \emph{invoked as goals}. They give higher-order programming in the style a Scheme programmer knows:
  \lp{mappred}, folds, closures of relations, continuation passing.
\item \textbf{Non-predicate function variables} (target type not $\typ{o}$) are instantiated by
  $\lambda$-terms that are only ever \emph{evaluated by $\beta\eta$-conversion}. Their values can also be
  \emph{found by solving unification problems}; the book calls this ``perhaps the single most novel aspect
  of higher-order logic programming'' \book{138}.
\end{itemize}
We first fix which formulas may be goals and which may be clauses (\cref{sec:hoprog-languages}), then
program with predicate variables (\cref{sec:hoprog-hoprog,sec:hoprog-flex,sec:hoprog-reason,sec:hoprog-logconst,sec:hoprog-naf}),
then with function variables (\cref{sec:hoprog-functions,sec:hoprog-panacea}), compare with functional
programming (\cref{sec:hoprog-fp}), and finally show how the same two logical devices that hide a
predicate, $\forall$ in goals and $\imp$ in goals, also give a modules language
(\cref{sec:hoprog-modules}).

A question runs through the whole chapter, and the kanren boxes answer it case by case: \emph{in a
miniKanren hosted in Scheme, where a relation is a Scheme procedure, which of these features are
free, and which really need $\lambda$-terms in the term language?} \Cref{sec:hoprog-free} collects the
answers.

\begin{tutorial}[How to run the Scheme listings]
Every Scheme listing in this chapter was run on today's \micro{} and its output is shown verbatim
below it. Each file begins with \scm{(include "prelude.scm")}, which loads this prelude (the include line is
left out of the listings):
\begin{scheme}
(import scheme (chicken base) (chicken pretty-print)
        (aux base) (aux stream) (aux kanren micro))
(define-syntax-rule (show label expr)
  (begin (display label) (display " ") (write expr) (newline)))
\end{scheme}
\scm{(μkanren-run (q n grounded) g ...)} returns up to \scm{n} answers for \scm{q} (all answers
when \scm{n} is negative). With \scm{grounded} set to \scm{#t} it returns them as plain data; with
\scm{#f} it returns them as reified Scheme $\lambda$-expressions.
\end{tutorial}

% ======================================================================================================
\section{Which formulas are goals and which are clauses}\label{sec:hoprog-languages}

At first order (\cref{ch:horn,ch:hh}) goals and clauses of \fohh{} were
\begin{align*}
G &::= \top \mid A \mid G\wedge G \mid G\vee G \mid \forall x\,G \mid \exists x\,G \mid D\imp G\\
D &::= A \mid G\imp D \mid \forall x\,D \mid D\wedge D
\end{align*}
and \fohc{} is the fragment without $\forall$ and $\imp$ in goals \book{118}. The higher-order languages
keep this top-level shape. What changes is \emph{what an atom $A$ may be}: once quantification over
predicates is allowed, an atom may contain $\lambda$-terms and predicate variables, and its head may itself
be a variable.

\subsection{Rigid and flexible atoms}

\begin{definition}[Rigid and flexible atoms {\book{119}}]\label{def:hoprog-rigid}
The $\lambda$-normal form of a higher-order atomic formula has the shape $(h\,t_1\ldots t_n)$ where $h$ is
either a variable or a non-logical constant. The atom is \emph{rigid} if $h$ is a constant and
\emph{flexible} if $h$ is a variable.
\end{definition}

So \lp{age bob 23} and \lp{mappred P L K} are rigid, while \lp{P X Y} is flexible.
The atom \lp{(x\ age x 23) bob} is rigid too, because its $\lambda$-normal form is \lp{age bob 23}.
Flexible atoms in \emph{goals} are what make higher-order programming possible: in the body of
\lp{mappred P (X::L) (Y::K) :- P X Y, mappred P L K.} the goal \lp{P X Y} is flexible.
Flexible atoms as \emph{clause heads} are a different matter.

\subsection{Why clause heads must be rigid}\label{sec:hoprog-rigidheads}

The book gives two reasons for rejecting flexible heads \book{119--120}.

\paragraph{Consistency.} A program is \emph{inconsistent} if every formula is provable from it. Every
\fohc{} and \fohh{} program is consistent. This matters in practice: otherwise a complete interpreter would
have to check the program for consistency before it could answer ``no''. With flexible heads,
inconsistent programs are easy to write:
\begin{lprolog}
P 5 :- q.
q.
\end{lprolog}
To prove any goal $B$, backchain on the first clause with $P := \abs{x}{B}$; what is left is \lp{q}, which
is provable. The single clause $\forall p\,p$ is an even shorter example \book{120}.

\paragraph{Modularity.} A clause $\forall x_1\ldots\forall x_m(A_1\wedge\cdots\wedge A_n\imp A_0)$ is read as
part of the definition of the relation named by the head of $A_0$. If that head is a variable, the clause
adds meaning to \emph{every} predicate. That defeats modular reading of code, and it also defeats
indexing clauses by predicate name.

Two natural uses are lost, and the book argues that losing them is acceptable. \emph{Leibniz
equality} (two terms are equal if they satisfy the same properties) would turn $x+0=x$ into
\begin{lprolog}
P (X + 0) :- P X.
P X :- P (X + 0).
\end{lprolog}
This is logically meaningful but operationally a disaster. A goal containing $n$ occurrences of, say,
\lp{5 + 0} admits $2^n$ abstractions for \lp{P}, and even a goal with no occurrence can always
backchain on the first clause, so the search loops \book{120}. \emph{Computed clauses}, as in
\lp{?- convert d P, P => g.}, put a variable where a clause is expected. There is no simple way to
guarantee statically that the value bound to \lp{P} will be a legal clause \book{121}.

\subsection{No implications inside terms}\label{sec:hoprog-noimp}

Logical symbols inside the arguments of atoms are useful (the symmetric closure
\lp{x\y\ adj x y ; adj y x} is a term containing $\vee$), but an \emph{implication} inside a term breaks
goal-directed search. The book's counterexample \book{121} uses constants $r:\typ{o}\to\typ{o}$ and
$s,t:\typ{o}$, and a variable $Q:\typ{o}\to\typ{o}\to\typ{o}$:
\[
\exists Q\,\big[\,\forall p\,\forall q\,[\,r\,(p\imp q)\imp r\,(Q\,p\,q)\,]\ \wedge\ Q\,(t\vee s)\,(s\vee t)\,\big].
\]
The first conjunct forces $Q := \abs{x}{\abs{y}{x\imp y}}$: after the $\forall$ and $\imp$ introductions
the goal $r\,(Q\,p\,q)$ must match the hypothesis $r\,(p\imp q)$. The second conjunct then becomes
$(t\vee s)\imp(s\vee t)$, which is provable in any reasonable logic, but \emph{not} by the goal-directed
rules of \cref{ch:horn}: the instance of $Q$ moved a disjunction from a goal position into a hypothesis,
where a left rule for $\vee$ would be needed.

\begin{definition}[Polarity {\book{121--122}}]\label{def:hoprog-polarity}
When implications are excluded from terms, an occurrence of a logical symbol is \emph{positive} if and
only if it is either within an atomic formula or is in the scope of only logical constants and to the
left of an even number of implications. All other occurrences are \emph{negative}.
\end{definition}

The point of the definition is an invariant: during cut-free proof search, positive occurrences stay
positive and negative ones stay negative. Goal-directed search relies on it. In the counterexample both
disjunctions start out positive, and instantiating $Q$ makes one of them negative. (The ban is a design
choice, not the only one possible: the bibliographic notes recall Felty's result that $\imp$ may occur
in terms if $\vee$ and $\exists$ are banned from formulas and terms instead \book{148}.)

\subsection{Herbrand universes, hohc and hohh}\label{sec:hoprog-hohc}

The quantifier rules $\exists R$ and $\forall L$ substitute terms. To keep goals goals and clauses clauses,
the terms that may be substituted have to be restricted.

\begin{definition}[Herbrand universes {\book{122}}]\label{def:hoprog-herbrand}
Let $\Sig$ contain at least the logical constants. $\mathcal{H}_1^\Sig$ is the set of all (possibly open)
$\lambda$-normal terms over $\Sig$ that contain neither $\forall$ nor $\imp$. $\mathcal{H}_2^\Sig$ is the
set of all (possibly open) $\lambda$-normal terms over $\Sig$ that do not contain $\imp$.
\end{definition}

Terms of $\mathcal{H}_1^\Sig$ may contain $\top,\wedge,\vee,\exists$, and the types of their subterms are
unrestricted. $\mathcal{H}_2^\Sig$ also allows $\forall$.

\begin{definition}[\hohc{}: higher-order Horn clauses {\book{122}}]\label{def:hoprog-hohc}
Let $A$ range over atoms in $\mathcal{H}_1^\Sig$ and $A_r$ over \emph{rigid} atoms in
$\mathcal{H}_1^\Sig$. Goals and clauses are the formulas whose $\lambda$-normal forms are
\begin{align*}
G &::= \top \mid A \mid G\wedge G \mid G\vee G \mid \exists_\tau x\,G\\
D &::= A_r \mid G\imp D \mid D\wedge D \mid \forall_\tau x\,D
\end{align*}
with no restriction on $\tau$. The $G$-formulas are exactly the formulas of $\mathcal{H}_1^\Sig$ of type $\typ{o}$.
\end{definition}

The operational semantics is the goal-directed calculus of \cref{ch:horn}, with three
changes \book{123}: formulas are taken in $\lambda$-normal form, \rulename{initial} works up to
$\alpha$-conversion, and the quantifier rules instantiate with $\lambda$-terms from the Herbrand universe
($\Vdash$ is the typing judgement for the instantiation term):
\begin{mathpar}
\inferrule{ }{\Sig;\Prog\xrightarrow{A'}A}\;\rulename{initial}\ (A'\equiv_\alpha A)
\and
\inferrule{\seq{\Sig}{\Prog}{B[t/x]} \\ \Sig;\emptyset\Vdash t:\tau \\ t\in\mathcal{H}_1^\Sig}{\seq{\Sig}{\Prog}{\exists_\tau x\,B}}\;\rulename{$\exists$R}
\and
\inferrule{\Sig;\Prog\xrightarrow{D[t/x]}A \\ \Sig;\emptyset\Vdash t:\tau \\ t\in\mathcal{H}_1^\Sig}{\Sig;\Prog\xrightarrow{\forall_\tau x\,D}A}\;\rulename{$\forall$L}
\end{mathpar}
Why is the restriction needed? Without it, the goal $\exists p\,p$ could be instantiated with
$\forall x\,(q\,x)$ for some $q\in\Sig$, which is not an \hohc{} goal. Substituting a term of
$\mathcal{H}_1^\Sig$ into a goal or a clause keeps it in its category. Intuitively, we never substitute
something that could not appear as the argument of an atom or in the structure of a goal \book{123}. In an
implementation, $t$ is not guessed: it is a logic variable that higher-order unification later fixes in
\rulename{initial}.

\begin{theorem}[Completeness of goal-directed search for \hohc{} {\book{123--124}, \book{148}}]\label{thm:hoprog-hohc}
Derivability of goals from program clauses in this calculus coincides exactly with provability in
classical higher-order logic. By the bibliographic notes, this holds for both the classical and the
intuitionistic versions of Church's Simple Theory of Types, without the axioms of extensionality,
infinity and choice \cite{NadathurMiller1990}.
\end{theorem}

Since $\forall R$ and ${\imp}R$ never fire in \hohc{}, the signature and the program stay fixed during a
computation \book{124}.

\begin{definition}[\hohh{}: higher-order hereditary Harrop formulas {\book{124}}]\label{def:hoprog-hohh}
Let $A$ range over atoms in $\mathcal{H}_2^\Sig$ and $A_r$ over rigid atoms in $\mathcal{H}_2^\Sig$:
\begin{align*}
G &::= \top \mid A \mid G\wedge G \mid G\vee G \mid \exists_\tau x\,G \mid D\imp G \mid \forall_\tau x\,G\\
D &::= A_r \mid G\imp D \mid D\wedge D \mid \forall_\tau x\,D
\end{align*}
The instances for $\exists R$ and $\forall L$ are terms of $\mathcal{H}_2^\Sig$, where $\Sig$ is the
signature \emph{current at that point of the derivation}.
\end{definition}

Goals are richer than $\mathcal{H}_2^\Sig$ (they may contain $\imp$; no term of $\mathcal{H}_2^\Sig$ does).
Derivability now coincides with provability in higher-order \emph{intuitionistic}
logic \cite{MillerNadathurPfenningScedrov1991} \book{124}. Unlike \hohc{}, $\forall R$
(\rulename{generic}) and ${\imp}R$ (\rulename{augment}) are active, so the signature and the program change
during computation, as in \cref{ch:hh}.

\subsection{\texorpdfstring{hohh$^+$}{hohh+}: hiding predicate names}\label{sec:hoprog-hohhplus}

\Cref{ch:hh} could limit the \emph{scope} of an auxiliary predicate's clauses with $\imp$, but not the
\emph{visibility of its name}. A context that already had clauses for \lp{rev} could corrupt a
\lp{reverse} that assumed its own \lp{rev}. One small liberalisation fixes this.

\begin{definition}[\hohh$^+$ {\book{125}}]\label{def:hoprog-hohhplus}
\hohh$^+$ is \hohh{} in which the head of an atomic program clause may also be a \emph{variable}, provided that
variable is bound by an \emph{essentially universal} quantifier that has the clause in its scope.
\end{definition}

In a goal, a bound variable is essentially universal if it is bound by a positive $\forall$, a negative
$\exists$, or a term-level $\lambda$. It is essentially existential otherwise. The definition is dualised
for clauses (\book{205}; see \cref{ch:lambdatree}). Essentially universal variables only ever become
eigenvariables. When search reaches the goal $\forall p\,(\ldots D\ldots\imp G)$, $\forall R$ has already
replaced $p$ by a \emph{new constant}, so by the time $D$ is added to the program it is a legal \hohh{} clause
with a rigid head. Hence \hohh$^+$ does not change computational behaviour; it only lets a program
introduce \emph{closed, hidden} predicate definitions. The book's two versions of \lp{reverse}
\book{126}:
\begin{lprolog}
type reverse   list A -> list A -> o.
reverse L K :-
  pi rev\
    (pi L\ rev nil L L) &
    (pi X\ pi L\ pi K\ pi M\ rev (X::L) K M :- rev L K (X::M))
    => rev L K nil.
\end{lprolog}
\begin{lprolog}
type reverse   list A -> list A -> o.
reverse L K :- pi rv\
  ( rv nil K &
    (pi X\ pi N\ pi M\ rv (X::N) M :- rv N (X::M)))
  => rv L nil.
\end{lprolog}

\begin{tutorial}[What the interpreter does with \texttt{reverse (1::2::3::nil) K}]
Using the first version \book{125}:
\begin{enumerate}
\item backchain on \lp{reverse}; the body is a \lp{pi} goal, so $\forall R$ generates a new predicate
  symbol, call it \lp{c}, and adds it to the signature;
\item the body is now an implication, so ${\imp}R$ adds the two clauses
\begin{lprolog}
pi L\ c nil L L.
pi X\ pi L\ pi K\ pi M\ c (X::L) K M :- c L K (X::M).
\end{lprolog}
\item it proves \lp{c (1::2::3::nil) K nil} by ordinary backchaining and obtains \lp{K = (3::2::1::nil)};
\item on leaving the goal, it removes \lp{c} and its clauses from the context.
\end{enumerate}
No clause outside \lp{reverse} can mention \lp{c}: the name did not exist when those clauses were written.
\end{tutorial}

\begin{pitfall}[Shadowing and shared variables in assumed clauses]
In the first version the inner \lp{pi L\} and \lp{pi K\} \emph{shadow} the outer \lp{L} and \lp{K}. In
the second version the assumed fact \lp{rv nil K} mentions the \emph{outer} logic variable \lp{K}, so the
hypothetical clause is not closed: \lp{K} is shared with the caller and is \emph{not} renamed at each use of
the clause.
\end{pitfall}

\begin{kanren}[Hidden predicates are lexical scope]
In a Scheme-hosted miniKanren a relation is a procedure, and Scheme already has a way to make a name
invisible: a local definition. The first \lp{reverse}, with \lp{rev} hidden, is one \scm{let}:
\begin{scheme}
;; Figure 5.1 with the hidden predicate rev as a lexically local relation
(define reverse°
  (let ()
    (define-relation (rev l k m)
      (cond° ((null° l) (=° k m))
             ((fresh° (x l*) (cons° x l* l) (rev l* k (cons x m))))))
    (λ (l k) (rev l k '()))))
(show 1 (μkanren-run (q -1 #t) (reverse° '(1 2 3) q)))
(show 2 (μkanren-run (q 1 #t) (reverse° q '(1 2 3))))
\end{scheme}
\begin{answers}
1 ((3 2 1))
2 ((3 2 1))
\end{answers}
No client of \scm{reverse°} can add clauses to \scm{rev}. This is the encapsulation that \hohh$^+$ provides,
and it costs nothing. (The second query asks for one answer: in that mode a request for all answers
does not terminate.) What lexical scope does \emph{not} give is a clause set that \emph{grows during the
search}, such as the memo table of \cref{ex:hoprog-fib}, or a hidden name that exists as a first-class
\emph{constant}, one that can be stored in terms and must not escape its scope. Both need new machinery
in the state: a dynamic clause store for $\imp$, and eigenvariables with a scope check for $\forall$.
See \cref{ch:design}.
\end{kanren}

% ======================================================================================================
\section{Higher-order relational programming}\label{sec:hoprog-hoprog}

\subsection{Mapping, filtering and quantifying over lists}

The basic higher-order relations of the book \book{127}:
\begin{lprolog}
type foreach, forsome   (A -> o) -> list A -> o.
type mappred            (A -> B -> o) -> list A -> list B -> o.
type sublist            (A -> o) -> list A -> list A -> o.

foreach P nil.
foreach P (X::L) :- P X, foreach P L.

forsome P (X::L) :- P X; forsome P L.

mappred P nil nil.
mappred P (X::L) (Y::K) :- P X Y, mappred P L K.

sublist P (X::L) (X::K) :- P X, sublist P L K.
sublist P (X::L) K :- sublist P L K.
sublist P nil nil.
\end{lprolog}
Read relationally \book{126}: \lp{mappred P L K} says the lists have equal length and are related pointwise
by \lp{P}; \lp{forsome P L} says some member satisfies \lp{P}; \lp{foreach P L} says all members do;
\lp{sublist P L K} says \lp{K} is an order-preserving sublist of \lp{L} whose members all satisfy \lp{P}.
Because of the clause order, the \emph{maximal} sublists come first. A small database:
\begin{lprolog}
kind name    type.
type bob, sue, ned   name.
type age     name -> int -> o.
type male, female    name -> o.

age bob 23 & age sue 24 & age ned 23.
male bob & female sue & male ned.
\end{lprolog}
\begin{example}[Queries over the database {\book{126--127}}]\label{ex:hoprog-mappred}
\begin{lprolog}
?- mappred age (ned::bob::sue::nil) L.
L = (23::23::24::nil)
?- mappred age L (23::24::nil).
L = (bob::sue::nil);
L = (ned::sue::nil)
?- sublist male (ned::bob::sue::nil) L.
L = ned::bob::nil;
L = ned::nil;
L = bob::nil;
L = nil;
no
?- forsome female (ned::bob::sue::nil).
solved
?- foreach female (ned::bob::sue::nil).
no
\end{lprolog}
The second query runs \lp{mappred} ``backwards'': \lp{age} is a relation, not a function, so the list of
names is computed from the list of ages.
\end{example}

\subsection{\texorpdfstring{$\lambda$}{Lambda}-terms of predicate type as arguments}

A predicate argument need not be a constant. \lp{x\ age jane x} has type \lp{int -> o} (it is
$\eta$-equivalent to \lp{age jane}), and \lp{n\ age n 24} has type \lp{name -> o} \book{129}. The book
reports:
\begin{lprolog}
?- (n\ age n 24) W.          % β-equal to ?- age W 24.
W = sue
?- mappred (x\y\ age x y) (ned::bob::sue::nil) L.
L = (23::23::24::nil)
?- mappred (x\y\ age y x) (23::24::nil) K.
K = (bob::sue::nil);
K = (ned::sue::nil)
?- foreach (x\ sigma y\ age x y) (ned::bob::sue::nil).
solved
?- foreach (x\ age x A) (ned::bob::sue::nil).
no
?- foreach (x\ age x A) (ned::bob::nil).
A = 23
\end{lprolog}

\begin{pitfall}[A logic variable inside a \texorpdfstring{$\lambda$}{lambda} is shared, not abstracted]
In \lp{x\ age x A}, the variable \lp{A} lies \emph{outside} the abstraction: the same \lp{A} is used by
every instance of the body. The last two queries therefore ask for a \emph{common} age. Compare
\lp{x\ sigma y\ age x y}, where each element gets its own \lp{y}.
\end{pitfall}

\begin{kanren}[Predicate variables with known values are free]
In a Scheme-hosted miniKanren, a predicate argument whose value is known when it is called is just a
procedure that returns a goal. \lp{x\y\ age y x} becomes \scm{(λ (x y) (age° y x))}, and
\lp{x\ sigma y\ age x y} becomes \scm{(λ (x) (fresh° (y) (age° x y)))}. A logic variable captured by the
Scheme \scm{λ} is shared by every call, exactly as in the pitfall above:
\begin{scheme}
;; Figure 5.4: the age database, one relation per predicate constant
(define-relation (age° n a)
  (cond° ((=° n 'bob) (=° a 23))
         ((=° n 'sue) (=° a 24))
         ((=° n 'ned) (=° a 23))))
(define-relation (male° n) (or° (=° n 'bob) (=° n 'ned)))
;; Figure 5.3: the argument P is a Scheme procedure returning a goal
(define-relation (mappred° P L K)
  (cond° ((null° L) (null° K))
         ((fresh° (X L* Y K*) (cons° X L* L) (cons° Y K* K)
            (P X Y) (mappred° P L* K*)))))
(define-relation (foreach° P L)
  (cond° ((null° L))
         ((fresh° (X L*) (cons° X L* L) (P X) (foreach° P L*)))))
(define-relation (sublist° P L K)
  (cond° ((fresh° (X L* K*) (cons° X L* L) (cons° X K* K)
            (P X) (sublist° P L* K*)))
         ((fresh° (X L*) (cons° X L* L) (sublist° P L* K)))
         ((null° L) (null° K))))

(show 1 (μkanren-run (q -1 #t) (mappred° age° '(ned bob sue) q)))
(show 2 (μkanren-run (q -1 #t) (mappred° age° q '(23 24))))
(show 3 (μkanren-run (q -1 #t) (sublist° male° '(ned bob sue) q)))
;; λ-terms of predicate type are Scheme lambdas
(show 4 (μkanren-run (q -1 #t) (mappred° (λ (x y) (age° y x)) '(23 24) q)))
(show 5 (μkanren-run (q -1 #t)
          (foreach° (λ (x) (fresh° (y) (age° x y))) '(ned bob sue))
          (=° q 'solved)))
(show 6 (μkanren-run (q -1 #t) (foreach° (λ (x) (age° x q)) '(ned bob sue))))
(show 7 (μkanren-run (q -1 #t) (foreach° (λ (x) (age° x q)) '(ned bob))))
\end{scheme}
\begin{answers}
1 ((23 23 24))
2 ((bob sue) (ned sue))
3 (() (bob) (ned) (ned bob))
4 ((bob sue) (ned sue))
5 (solved)
6 ()
7 (23)
\end{answers}
Every answer \emph{set} matches the book. The \emph{order} of the \lp{sublist} answers (line 3) does not,
because \scm{cond°} interleaves its branches while λProlog searches them depth-first in clause order.
Compare answer sets unless the search strategy is part of the specification.
\end{kanren}

\subsection{Closures, composition and folds}\label{sec:hoprog-closures}

Relations that take relations and return relations \book{127--128}:
\begin{lprolog}
type ref, sym, trans   (A -> A -> o) -> A -> A -> o.
ref   R X Y :- X = Y; R X Y.
sym   R X Y :- R X Y; R Y X.
trans R X Y :- R X Y.
trans R X Z :- R X Y, trans R Y Z.

kind node  type.
type a, b, c, d, e   node.
type adj   node -> node -> o.
adj a b & adj b c & adj b d & adj d c & adj c e.
\end{lprolog}
Both \lp{?- trans adj a d.} and \lp{?- sym adj b a.} succeed. Thanks to currying, the equivalence closure
is just the partial application \lp{(trans (sym (ref adj)))}. The non-recursive closures can also be
written as $\lambda$-terms, for example \lp{(x\y\ x = y ; adj x y)}. The book adds a warning
\book{128}. The transitive closure can be written \emph{declaratively} as a single term,
\begin{lprolog}
x\y\ pi p\
   (pi U\ pi V\ adj U V => p U V) =>
   (pi U\ pi V\ pi W\ adj U V => p V W => p U W) =>
   p x y
\end{lprolog}
but it contains $\imp$. It is therefore \emph{not} in $\mathcal{H}_2^\Sig$ and cannot be passed as an
argument (\cref{sec:hoprog-noimp}).

Union, composition and a left fold \book{130}:
\begin{lprolog}
type union    (A -> B -> o) -> (A -> B -> o) -> A -> B -> o.
type compose  (A -> B -> o) -> (B -> C -> o) -> A -> C -> o.
type foldl    (A -> B -> B -> o) -> list A -> B -> B -> o.

union   R S X Y :- R X Y; S X Y.
compose R S X Y :- R X Z, S Z Y.
foldl P nil X X.
foldl P (Z::L) X Y :- P Z X W, foldl P L W Y.
\end{lprolog}
If \lp{P : A -> B -> B -> o} is a family of binary relations on \lp{B} indexed by \lp{A}, then
\lp{(foldl P)} composes them one after another along a list. Reversing a list then means pushing every
element onto a stack and popping them all off \book{130--131}:
\begin{lprolog}
kind stack   type -> type.
type emp     stack A.
type stk     A -> stack A -> stack A.
type empty   stack A -> o.
type enter, remove   A -> stack A -> stack A -> o.
type reverse list A -> list A -> o.

empty emp.
enter X S (stk X S).
remove X (stk X S) S.
reverse L K :- compose (foldl enter L) (foldl remove K) emp emp.
\end{lprolog}

\begin{kanren}[Partial application is curried \texttt{define}]
λProlog is curried, while \scm{define-relation} has a fixed arity. A relation-returning relation is
written with CHICKEN's curried \scm{define}, \scm{(define ((sym° R) x y) ...)}; a recursive one defines a
local relation and returns it:
\begin{scheme}
(define-relation (adj° x y)
  (cond° ((=° x 'a) (=° y 'b)) ((=° x 'b) (=° y 'c)) ((=° x 'b) (=° y 'd))
         ((=° x 'd) (=° y 'c)) ((=° x 'c) (=° y 'e))))
;; curried closure builders: (sym° R) is itself a binary relation
(define ((ref° R) x y) (or° (=° x y) (R x y)))
(define ((sym° R) x y) (or° (R x y) (R y x)))
(define (trans° R)
  (define-relation (T x z) (cond° ((R x z)) ((fresh° (y) (R x y) (T y z)))))
  T)
(show 1 (μkanren-run (q -1 #t) ((trans° adj°) 'a 'd) (=° q 'solved)))
(show 2 (μkanren-run (q -1 #t) ((sym° adj°) 'b 'a) (=° q 'solved)))
(show 3 (μkanren-run (q -1 #t) ((trans° adj°) 'a q)))
(show 4 (μkanren-run (q 1 #t)
          ((trans° (sym° (ref° adj°))) 'e 'a) (=° q 'solved)))
\end{scheme}
\begin{answers}
1 (solved)
2 (solved)
3 (b c d e c e)
4 (solved)
\end{answers}
Line 3 shows that \lp{trans} is a relation, not a set: \scm{c} and \scm{e} come back once per path. Line 4
uses the equivalence closure built by partial application. Composition and the stack-based reverse:
\begin{scheme}
;; Figures 5.5 and 5.6: compose, foldl, and reverse through a stack
(define ((compose° R S) x y) (fresh° (z) (R x z) (S z y)))
(define-relation (foldl° P L X Y)
  (cond° ((null° L) (=° X Y))
         ((fresh° (Z L* W) (cons° Z L* L) (P Z X W) (foldl° P L* W Y)))))
(define-relation (enter° X S S*) (=° S* `(stk ,X ,S)))
(define-relation (remove° X S S*) (=° S `(stk ,X ,S*)))
(define-relation (reverse° L K)
  ((compose° (λ (s s*) (foldl° enter° L s s*))
             (λ (s s*) (foldl° remove° K s s*)))
   'emp 'emp))
(show 1 (μkanren-run (q -1 #t) (reverse° '(1 2 3) q)))
(show 2 (μkanren-run (q -1 #t) (foldl° enter° '(1 2 3) 'emp q)))
\end{scheme}
\begin{answers}
1 ((3 2 1))
2 ((stk 3 (stk 2 (stk 1 emp))))
\end{answers}
The partial application \lp{(foldl enter L)} becomes an explicit \scm{λ}. None of this needs anything
beyond Scheme closures.
\end{kanren}

\subsection{Binary clauses: continuation passing}\label{sec:hoprog-cps}

Predicate variables of type $\typ{o}$ are \emph{goal-valued}, so a program can pass its own continuation
\book{131}. For each predicate $p:\tau_1\to\cdots\to\tau_j\to\typ{o}$ introduce
$p':\tau_1\to\cdots\to\tau_j\to\typ{o}\to\typ{o}$, and write $A'$ for $(p'\,t_1\ldots t_j)$ when $A=(p\,t_1\ldots t_j)$.
Then rewrite every clause:
\begin{gather*}
\forall\bar x\,[A_1\wedge\dots\wedge A_n\imp A_0]\ (n>0)\ \rightsquigarrow\ \forall\bar x\,\forall K\,[\,(A'_1\,(\dots(A'_n\,K)\dots))\imp(A'_0\,K)\,]\\
\forall\bar x\,A_0\ \rightsquigarrow\ \forall\bar x\,\forall K\,[\,K\imp(A'_0\,K)\,]
\end{gather*}
with $K:\typ{o}$. Conjunction becomes explicit sequencing, and every clause body has at most one, possibly
complex, atom. For example
\begin{lprolog}
adj a b.
adj b c.
path X Y :- adj X Y.
path X Y :- adj X Z, path Z Y.
\end{lprolog}
becomes
\begin{lprolog}
adj' a b K :- K.
adj' b c K :- K.
path' X Y K :- adj' X Y K.
path' X Y K :- adj' X Z (path' Z Y K).
\end{lprolog}
and \lp{?- path X Y.} is provable if and only if \lp{?- path' X Y true.} is provable \book{132}.

\begin{kanren}[Goals are already first-class]
In μKanren a goal is a value, so the binary-clause form is an idiom, not a feature. The continuation is a
goal argument:
\begin{scheme}
;; Section 5.3: binary clauses, the continuation K is a goal
(define (adj/k° x y k)
  (cond° ((=° x 'a) (=° y 'b) k) ((=° x 'b) (=° y 'c) k)))
(define-relation (path/k° x y k)
  (cond° ((adj/k° x y k))
         ((fresh° (z) (adj/k° x z (path/k° z y k))))))
(show 1 (μkanren-run (q -1 #t)
          (fresh° (x y) (path/k° x y ✓°) (=° q (list x y)))))
(show 2 (μkanren-run (q -1 #t)
          (fresh° (x y) (path/k° x y (=° q (list x y))))))
\end{scheme}
\begin{answers}
1 ((a b) (b c) (a c))
2 ((a b) (b c) (a c))
\end{answers}
In the second query the continuation does the work of the query itself. Conjunction in μKanren
(\scm{andª}, which feeds every state of the first goal to the second) is this same continuation
composition.
\end{kanren}

\subsection{Memoisation with hidden, growing predicates}

\begin{example}[Memoised Fibonacci in \hohh$^+$ {\book{132}}]\label{ex:hoprog-fib}
\begin{lprolog}
type fib_memo   int -> ((int -> int -> o) -> o) -> o.
fib_memo N G :-
   pi memo\ pi loop\
     memo 0 0 => memo 1 1 =>
     (pi F1\ pi F2\ pi F3\ pi C\ pi C'\
        (loop C F1 F2 :- (C >  N, (G memo)) ;
                         (C =< N, C' is C + 1, F3 is F1 + F2,
                          memo C F3 => loop C' F2 F3))) =>
     loop 2 0 1.
\end{lprolog}
\lp{memo} and \lp{loop} are hidden predicates, bound by \lp{pi} and used as heads of assumed clauses. The
memo table \emph{grows} through \lp{memo C F3 => ...}. When \lp{C > N}, the continuation \lp{G}, a
predicate that takes a predicate, is called with the hidden name \lp{memo}, whose extension is then
$\{(0,0),(1,1),\ldots,(N,\mathit{fib}\,N)\}$. For the query
\begin{lprolog}
?- fib_memo 20 (fib\ sigma M\ fib N M, M is N * N).
\end{lprolog}
the book states, without listing them, that ``there are exactly three values for N that make
this atom provable'' \book{132}. A short computation gives $N\in\{0,1,12\}$: the continuation asks
for $\mathit{fib}\,N=N\cdot N$, and $\mathit{fib}\,0=0$, $\mathit{fib}\,1=1$, $\mathit{fib}\,12=144$.
\end{example}

This example mixes three things: higher-order programming, dynamic clause addition, and a hidden name
handed to a continuation. A Scheme \scm{letrec} cannot express the growing table, which is why the
kanren box of \cref{sec:hoprog-hohhplus} lists dynamic clauses as a genuine requirement.

% ======================================================================================================
\section{Flexible goals}\label{sec:hoprog-flex}

What happens when a goal's head is a variable that is \emph{still unbound} when the goal is reached? Take
\lp{P bob 23} against the age database. \lp{P := age} works, but so do \lp{x\y\ age x 23, age ned y},
\lp{x\y\ age sue 24}, \lp{x\y\ memb 4 (3::4::5::nil)}, and in fact \lp{x\y\ g} for \emph{any} provable
closed goal \lp{g} \book{132--133}. Such a goal is \emph{underconstrained}. Three treatments are
possible \book{134}:
\begin{enumerate}
\item \textbf{Suspend} the goal in the hope that other goals will instantiate its head, and resume it when
  it becomes rigid. If only flexible goals remain, solve each of them by instantiating its head with
  $\abs{x_1}{\ldots\abs{x_n}{\top}}$, the universally true relation of the right type, which is ``the
  (extensionally) largest one that works''.
\item \textbf{Eagerly} bind the head to the universally true relation.
\item Issue a \textbf{run-time error}.
\end{enumerate}
The first option is complete with respect to provability in intuitionistic higher-order logic. The other
two are incomplete. Yet experience with λProlog implementations says the third is the most useful while
developing code: a flexible goal met by depth-first search is usually a bug, for example a predicate
constant mistyped with a capital letter \book{134}. Option~1 is the same idea that unification-based
interpreters apply to flexible-flexible unification problems: the completeness results ``provide the
foundations for a unification-based interpreter \ldots\ that explicitly delays the consideration of
flexible-flexible unification problems as well as flexible goals'' \book{148}.

A predicate variable becomes meaningful when the program first \emph{constrains} it to a structured space
of candidates \book{133}:
\begin{lprolog}
kind i  type.
type jane, mary, john   i.
type mother, father, wife, husband   i -> i -> o.
type primrel, rel   (i -> i -> o) -> o.

primrel father & primrel mother & primrel wife & primrel husband.
rel R :- primrel R.
rel (x\y\ sigma z\ R x z, S z y) :- primrel R, primrel S.
mother jane mary & wife john jane.
\end{lprolog}
\lp{?- R john mary.} on its own is underconstrained. \lp{?- rel R, R john mary.} is meaningful:
\lp{rel} enumerates candidate relations, after which \lp{R john mary} is rigid. The only answer is
\lp{R = x\y\ sigma Z\ wife x Z, mother Z y}. Note that the head of the second \lp{rel} clause contains a
$\lambda$-term with $\exists$ and $\wedge$. This is legal because the term lies in $\mathcal{H}_1^\Sig$.
Whether a flexible goal is underconstrained thus ``depends on when the goal is invoked''
\book{133}: goal order matters operationally.

\begin{tutorial}[Depth-first search for \texttt{rel R, R john mary}]
Rename the clause variables of the second \lp{rel} clause to \lp{R1} and \lp{S1}, so that they are not
confused with the query's \lp{R}.
\begin{enumerate}
\item \lp{rel R} backchains on the first clause and \lp{primrel R} binds \lp{R := father}. The second
  goal is now the \emph{rigid} atom \lp{father john mary}, which has no clauses and fails. Backtracking
  tries \lp{mother}, \lp{wife} and \lp{husband} in turn; all fail the same way.
\item \lp{rel R} backchains on the second clause:
  \lp{R := x\y\ sigma z\ R1 x z, S1 z y}, with the subgoals \lp{primrel R1} and \lp{primrel S1}.
  After they bind \lp{R1} and \lp{S1}, the goal \lp{R john mary} is $\beta$-reduced to
  \lp{sigma z\ R1 john z, S1 z mary}, which is again built from rigid atoms.
\item With \lp{R1 := father} or \lp{R1 := mother} the atom \lp{R1 john z} fails for every choice of
  \lp{S1}. With \lp{R1 := wife} it gives \lp{z := jane}; then \lp{S1 := father} fails on
  \lp{father jane mary} and \lp{S1 := mother} succeeds on \lp{mother jane mary}.
\end{enumerate}
The answer is \lp{R} with \lp{R1} and \lp{S1} instantiated, which is the book's answer. The predicate
variable is never invoked while unbound: \lp{rel} always gives it a value first.
\end{tutorial}

\begin{kanren}[Binding a logic variable to a relation: possible, but opaque]
Any Scheme value can be bound to a logic variable, procedures included, so \lp{rel} can be transcribed.
The composed relation is built by \scm{project°} once \scm{r} and \scm{s} are known, and invoking the
value of \scm{R} means walking it and applying it:
\begin{scheme}
(import (chicken condition))
;; Figure 5.8: relations as values of logic variables
(define-relation (mother° x y) (=° x 'jane) (=° y 'mary))
(define-relation (wife° x y) (=° x 'john) (=° y 'jane))
(define-relation (father° x y) ✗°)
(define-relation (husband° x y) ✗°)
(define-relation (primrel° R)
  (cond° ((=° R father°)) ((=° R mother°)) ((=° R wife°)) ((=° R husband°))))
(define-relation (rel° R)
  (cond° ((primrel° R))
         ((fresh° (r s) (primrel° r) (primrel° s)
            (project° ((r* r) (s* s))
              (=° R (λ (x y) (fresh° (z) (r* x z) (s* z y)))))))))
;; invoking the value of R: walk it, then apply it
(define-relation (rel-john-mary° R)
  (rel° R) (project° ((R* R)) (R* 'john 'mary)))
(show 1 (μkanren-run (q -1 #f) (rel-john-mary° q)))
(show 2 (μkanren-run (q -1 #f) (fresh° (R) (rel-john-mary° R) (=° q 'found))))
;; a flexible goal: R is still unbound when it is invoked
(show 3 (handle-exceptions exn
            (list 'error ((condition-property-accessor 'exn 'message) exn))
          (μkanren-run (q -1 #t) (project° ((R* q)) (R* 'bob 23)))))
\end{scheme}
\begin{answers}
1 ((λ () #<procedure (? x y)>))
2 ((λ () (quote found)))
3 (error "call of non-procedure")
\end{answers}
The search finds the single answer (line 2), but the answer itself is \scm{#<procedure (? x y)>}, not
\lp{x\y\ sigma Z\ wife x Z, mother Z y}. Grounded reification (\scm{#t}) even fails on it, because it
\scm{eval}s the reified expression, and \scm{eval} rejects a procedure as an ``illegal non-atomic
object''. Line 3 is a flexible goal: \scm{project°} hands over an unbound variable record, and applying
it is a Scheme error. That is the book's option~3, but only by accident. Options~1 and~2 need a goal that
can be \emph{inspected} before it is run. Option~1 also needs a store of suspended goals that unification
wakes up, much as the disequality store is re-checked after each \scm{=°}. Printing \lp{R} as a
$\lambda$-term needs relations represented as terms. The proposal keeps relations as closures
(\cref{sec:roadmap-lim-dynamic}); predicate quantification is sketched as \cref{sec:roadmap-later-pred}.
\end{kanren}

% ======================================================================================================
\section{Reasoning about programs by substitution}\label{sec:hoprog-reason}

Logic gives a means of \emph{reasoning} about programs, and hidden predicates make it sharper
\book{134--136}. Claim: the second \lp{reverse} of \cref{sec:hoprog-hohhplus} is symmetric, that is, if
\lp{reverse L K} is provable then so is \lp{reverse K L}. The proof outline \book{135--136}:
\begin{enumerate}
\item The only way to prove \lp{reverse L K} is by backchaining, so its body is provable:
\begin{lprolog}
pi rv\ ( rv nil K &
         (pi X\ pi N\ pi M\ rv (X::N) M :- rv N (X::M)))
       => rv L nil
\end{lprolog}
\item A provable universal has provable instances. Instantiate \lp{rv} with \lp{x\y\ (not (rv y x))}, where
  \lp{not} is \emph{logical} negation, not negation as failure:
\begin{lprolog}
not(rv K nil) &
(pi X\ pi N\ pi M\ not(rv M (X::N)) :- not(rv (X::M) N))
=> not(rv nil L).
\end{lprolog}
\item Take the contrapositive of each implication, which is classically valid:
\begin{lprolog}
rv nil L & (pi X\ pi N\ pi M\ rv (X::M) N :- rv M (X::N))
=> rv K nil.
\end{lprolog}
\item This is a first-order Horn entailment, and on Horn clauses classical and intuitionistic provability
  coincide.
\item Generalise over \lp{rv} and rename \lp{M} and \lp{N}. The result is the body of \lp{reverse} with
  \lp{L} and \lp{K} swapped. \qed
\end{enumerate}
Intuitively, the trace of \lp{rv} on \lp{(a::b::c::nil)} is a table with rows
\lp{(a::b::c::nil) nil}, \lp{(b::c::nil) (a::nil)}, \lp{(c::nil) (b::a::nil)}, \lp{nil (c::b::a::nil)}.
Swapping the columns and reading the rows in reverse order gives a valid trace for the reverse of
\lp{(c::b::a::nil)} \book{135}.

Three subtleties are worth noting. The \lp{rv} inside the instantiating term \lp{x\y\ (not (rv y x))}
is a \emph{new} predicate variable that happens to reuse the name; step~5 generalises over it. The
instantiating term contains $\neg$, that is $\imp\bot$, so it is
\emph{not} in $\mathcal{H}_2^\Sig$. That is fine, because the argument uses the metatheory
(soundness, then instantiation, contraposition and completeness), not goal-directed search. And the argument
needs \lp{rv} to be hidden: no other clause can mention it, so any $\lambda$-term may be ``plugged in''
at its place.

\begin{kanren}[Not available to closures]
A relation that is a Scheme closure offers nothing to substitute into: the argument above treats the
program as a \emph{formula}. A port can at most turn the theorem into a property test (for random
lists, \lp{reverse L K} implies \lp{reverse K L}). Reasoning of this kind belongs to a deep
embedding, where clauses are data.
\end{kanren}

% ======================================================================================================
\section{Defining the logical constants}\label{sec:hoprog-logconst}

Predicate variables are powerful enough to \emph{define} logical constants inside the logic
\book{136--137}. In \hohc{}, clauses can specify the right-introduction rules:
\begin{lprolog}
type tt, ff   o.
type or       o -> o -> o.
type exists   (A -> o) -> o.
tt.                  % true
or P Q :- P.         % disjunction
or P Q :- Q.
exists B :- B T.     % existential quantifier
\end{lprolog}
\lp{ff} has no clauses, so proving it fails. Only the right rules are captured. The left rules (for
example the rule of cases for a disjunctive hypothesis) are not. In \hohh{} and \hohh$^+$, $\forall$
and $\imp$ give the standard impredicative encodings \book{136}:
\begin{gather*}
\bot \equiv \forall p\,p \qquad
\top \equiv \forall p\,(p\imp p)\\
B\vee C \equiv \forall p\,((B\imp p)\imp(C\imp p)\imp p)\qquad
\exists x\,(G\,x) \equiv \forall p\,((\forall x\,(G\,x\imp p))\imp p)
\end{gather*}
(the last three are \hohh$^+$ goals, since $p$ heads assumed clauses and is bound by a positive $\forall$).

\begin{tutorial}[Running the encodings]
Goal-directed search reads them as follows. For $\forall p\,p$, $p$ becomes a fresh constant with no
clauses, so the goal fails. For $\forall p\,((B\imp p)\imp(C\imp p)\imp p)$, a fresh $c$ is created, the
clauses $c$ \lp{:-} $B$ and $c$ \lp{:-} $C$ are assumed, and $c$ is proved by backchaining on either clause,
that is, by proving $B$ or $C$. For $\forall p\,((\forall x\,(G\,x\imp p))\imp p)$, a fresh $c$ is created,
the clause $\forall x\,(c$ \lp{:-} $G\,x)$ is assumed, and $c$ is proved by backchaining: $x$ becomes a logic
variable $X$ and the goal $G\,X$ remains.
\end{tutorial}

\begin{kanren}[Which definitions are free]
The \hohc{} definitions are one-liners over existing goals: \lp{tt} is \scm{✓°}, \lp{ff} is \scm{✗°},
\lp{or} is \scm{or°}, and \lp{exists B :- B T} is \scm{(fresh° (t) (B t))}. The \hohh{} encodings are not
free. They need a goal-level $\forall$ that creates a fresh constant and a goal-level $\imp$ that assumes
clauses for it. Those are the two state extensions of the kanren box in \cref{sec:hoprog-hohhplus}.
\end{kanren}

% ======================================================================================================
\section{Conditional and negation-as-failure}\label{sec:hoprog-naf}

With the cut \lp{!}, the book defines a conditional and negation-as-failure \book{137}:
\begin{lprolog}
type if   o -> o -> o -> o.
if P Q R :- P, !, Q.
if P Q R :- R.

type not  o -> o.
not P :- P, !, fail.
not P.
% alternatively
not P :- if P fail true.
\end{lprolog}
\lp{if P Q R} tries \lp{P}. If it succeeds, the cut commits to its first solution and \lp{Q} is proved;
only if \lp{P} fails is \lp{R} proved. This is not logic: \lp{?- X = 2, not (1 = X).} succeeds, while
\lp{?- not (1 = X), X = 2.} fails. ``Thus even a basic property such as commutativity of conjunctions
becomes invalid in this setting'', and the book uses these predicates ``extremely sparingly''
\book{137--138}.

\begin{kanren}[The cut as \texttt{take° 1}]
\micro{} has the soft cut \scm{if°}, which keeps \emph{every} answer of its test (like miniKanren's
\texttt{conda}), and \scm{take°}. The cut-based \lp{if} commits to the first answer, so the faithful
transcription of \lp{not P} is \scm{(if° (take° 1 P) ✗° ✓°)}. The \scm{take° 1} is not cosmetic: without it,
a test with infinitely many answers would make the conditional fail on each of them forever. The
non-commutativity reproduces exactly, and the disequality constraint \scm{≠°} is the declarative
alternative that commutes:
\begin{scheme}
;; Section 5.7: not P :- if P fail true, the cut read as "first answer only"
(show 1 (μkanren-run (q -1 #t) (=° q 2) (if° (take° 1 (=° 1 q)) ✗° ✓°)))
(show 2 (μkanren-run (q -1 #t) (if° (take° 1 (=° 1 q)) ✗° ✓°) (=° q 2)))
;; the declarative alternative: a disequality constraint commutes
(show 3 (μkanren-run (q -1 #t) (=° q 2) (≠° 1 q)))
(show 4 (μkanren-run (q -1 #t) (≠° 1 q) (=° q 2)))
\end{scheme}
\begin{answers}
1 (2)
2 ()
3 (2)
4 (2)
\end{answers}
\end{kanren}

% ======================================================================================================
\section{\texorpdfstring{$\lambda$}{Lambda}-terms as functions}\label{sec:hoprog-functions}

So far the higher-order variables have been predicate variables. Variables of non-predicate functional
type are instantiated by $\lambda$-terms that are only \emph{evaluated}, by $\beta\eta$-conversion. As
\cref{ch:lambda} showed, these functions are weak. ``This weakness, however can be an asset for a
different kind of computation: Rather than instantiating function variables with $\lambda$-terms, we can
think of finding values for them by solving unification problems'' \book{138}.

\subsection{\texttt{mapfun} and \texttt{reducefun}}

\begin{lprolog}
type mapfun     (A -> B) -> list A -> list B -> o.
type reducefun  (A -> B -> B) -> list A -> B -> B -> o.

mapfun F nil nil.
mapfun F (X::L) ((F X)::K) :- mapfun F L K.

reducefun F nil Z Z.
reducefun F (H::T) Z (F H R) :- reducefun F T Z R.
\end{lprolog}
Assume a type \lp{i}, constants \lp{a, b, c, d : i} and \lp{g : i -> i -> i}.

\paragraph{Forwards.} \lp{?- mapfun (x\ g a x) (a::b::nil) L.} has exactly one answer,
\lp{L = ((g a a)::(g a b)::nil)}, obtained by $\beta$-normalising \lp{((x\ g a x) a)} and
\lp{((x\ g a x) b)} \book{138}. \lp{mappred} runs arbitrary goals, so it is strictly more powerful. In fact
\lp{mapfun F L K :- mappred (x\y\ y = F x) L K.}

\paragraph{Backwards: function synthesis.} \book{139}
\begin{lprolog}
?- mapfun F (a::b::nil) ((g a a)::(g a b)::nil).
F = x\ g a x                        % the only answer
\end{lprolog}
The first element alone poses \lp{F a = g a a}, which has four unifiers:
\begin{lprolog}
(x\ g x x)   (x\ g a x)   (x\ g x a)   (x\ g a a)
\end{lprolog}
Only \lp{x\ g a x} also solves \lp{F b = g a b}, so an interpreter that tries another unifier first has to
backtrack. Weakness can also show up as failure: \lp{?- mapfun F (a::b::nil) (c::d::nil).} answers
\lp{no}. Simply typed $\lambda$-terms cannot do case analysis on constants, so no term maps
\lp{a} to \lp{c} and \lp{b} to \lp{d} (\cref{ch:lambda}) \book{139--140}.

\paragraph{\texttt{reducefun}.} Here \lp{+} is only a term constructor, with no arithmetic \book{140}:
\begin{lprolog}
?- reducefun (x\y\ x + y) (3::4::8::nil) 6 R.
R = 3 + (4 + (8 + 6));
?- reducefun F (4::8::nil) 6 (1 + (4 + (1 + (8 + 6)))).
F = x\y\ 1 + (4 + (1 + (8 + 6)));
F = x\y\ 1 + (x + (1 + (8 + 6)));
F = x\y\ 1 + (x + y);
no
?- pi z\ reducefun F (4::8::nil) z (1 + (4 + (1 + (8 + z)))).
F = x\y\ 1 + (x + y);
no
\end{lprolog}

\begin{tutorial}[Why three answers, and then one]
Unfolding the two clauses turns the second query into the single equation
$F\,4\,(F\,8\,6) = 1+(4+(1+(8+6)))$. Its solutions are the bodies obtained by abstracting some occurrences
of the constants of the target consistently. The constant function abstracts nothing. The second answer
abstracts the \lp{4}. The third abstracts both the \lp{4} and the inner \lp{1 + (8 + 6)}, which has the
shape \lp{x + y}. A candidate such as \lp{x\y\ 1 + (4 + y)} fails, because $F\,8\,6$ would give
$1+(4+6)\neq 1+(8+6)$. In the third query, $6$ is replaced by a new constant $z$ introduced by
\lp{pi} \emph{inside} the scope of $F$'s implicit $\exists$. Instantiations of $F$ may not contain $z$, and
that eigenvariable check removes the two answers that would have to copy the base case into $F$. Only the
real fold function is left.
\end{tutorial}

\begin{kanren}[Forwards is free, synthesising \texttt{F} is not]
With \lp{F} a Scheme procedure that \emph{builds a term}, \lp{mapfun} and \lp{reducefun} run
forwards unchanged. Since \lp{F} is applied to a logic variable while the goal is being built, they also
run backwards in the \emph{list} argument:
\begin{scheme}
;; Figure 5.10, forward mode only: F is a Scheme procedure that builds a term
(define-relation (mapfun° F L K)
  (cond° ((null° L) (null° K))
         ((fresh° (X L* K*) (cons° X L* L) (cons° (F X) K* K)
            (mapfun° F L* K*)))))
(define-relation (reducefun° F L Z R)
  (cond° ((null° L) (=° Z R))
         ((fresh° (H T R*) (cons° H T L) (=° R (F H R*))
            (reducefun° F T Z R*)))))
(show 1 (μkanren-run (q -1 #t) (mapfun° (λ (x) `(g a ,x)) '(a b) q)))
(show 2 (μkanren-run (q -1 #t)
          (reducefun° (λ (x y) `(+ ,x ,y)) '(3 4 8) 6 q)))
;; F is applied to a logic variable while the goal is built:
;; the term (g a X) contains X, so the list argument can be unknown
(show 3 (μkanren-run (q -1 #t)
          (mapfun° (λ (x) `(g a ,x)) q '((g a a) (g a b)))))
\end{scheme}
\begin{answers}
1 (((g a a) (g a b)))
2 ((+ 3 (+ 4 (+ 8 6))))
3 ((a b))
\end{answers}
What cannot be transcribed is the case where \lp{F} itself is unknown. A Scheme procedure is opaque, so
there is nothing to unify \lp{F a} against. Function synthesis needs $\lambda$-terms as data, with
$\beta$-normalisation, and a unifier that handles applications of a variable. The problems above,
\lp{F a = g a a} and \lp{F 4 R = ...}, are \emph{not} higher-order patterns in the sense of
\cite{Miller1991}, because \lp{F} is applied to constants. The first is a second-order
\emph{matching} problem (its right-hand side is closed); the second is close to one (the flexible side
also contains the unknown \lp{R}). Both need Huet-style branching \cite{Huet1975} (see \cref{ch:hounif}). \Cref{sec:roadmap-later-huet} discusses where that
belongs in \micro{}.
\end{kanren}

\subsection{Functional difference lists}\label{sec:hoprog-fdl}

A first-order difference list \lp{(dl L K)}, with \lp{K} a suffix of \lp{L}, denotes \lp{L} without
\lp{K}. For instance, both of the terms
\lp{(dl (1::2::3::4::5::nil) (4::5::nil))} and \lp{(dl (1::2::3::L) L)}
denote the list \lp{(1::2::3::nil)}, and with the second form concatenation needs no recursion. A \emph{functional}
difference list (due to Brisset and Ridoux \cite{BrissetRidoux1991}) represents the list as a function
from its tail to the whole list \book{140--142}. The book's Figure~5.11 collects the labels of a
binary tree into a list in two ways, with ordinary and with functional difference lists. Here are the
declarations and the functional version:
\begin{lprolog}
kind dlist    type -> type.
type dl       list A -> list A -> dlist A.

type concat   dlist A -> dlist A -> dlist A -> o.
concat (dl L1 L2) (dl L2 L3) (dl L1 L3).

kind btree    type -> type.
type empty    btree A.
type bt       A -> btree A -> btree A -> btree A.

kind fdlist   type -> type.
type fdl      (list A -> list A) -> fdlist A.

type collect'  btree A -> list A -> o.
type aux'      btree A -> fdlist A -> o.
collect' Bt (A nil) :- aux' Bt (fdl A).
aux' empty (fdl x\ x).
aux' (bt N L R) (fdl x\ A (N::(B x))) :- aux' L (fdl A),
                                         aux' R (fdl B).
\end{lprolog}
\lp{(fdl x\ (1::2::3::x))} denotes \lp{(1::2::3::nil)}, with a ``pointer'' to its tail. Each operation is
a single unification \book{142}:
\begin{itemize}
\item head and tail: unify with \lp{x\ (H::(T x))}; the head is \lp{H} and the tail is \lp{(fdl T)};
\item last and front: unify with \lp{x\ F (Y::x)}; the last element is \lp{Y} and the front is \lp{(fdl F)};
\item concatenation of \lp{(fdl A)} and \lp{(fdl B)}: \lp{(fdl x\ A (B x))}, that is, function composition;
\item conversion to a list: apply to \lp{nil}, as \lp{(A nil)} in \lp{collect'}.
\end{itemize}
Palindromes then need no explicit reversal \book{142}:
\begin{lprolog}
type palindrome   fdlist A -> o.
palindrome (fdl x\x).
palindrome (fdl x\ Y::x).
palindrome (fdl x\ Y::(F (Y::x))) :- palindrome (fdl F).
\end{lprolog}
\begin{lprolog}
?- palindrome (fdl x\ a::b::c::b::a::x).
solved
?- palindrome (fdl x\ a::b::a::(F x)).
F = x\ x;
F = x\ b::a::x;
F = x\ a::b::a::x;
F = x\ Y::a::b::a::x;
F = x\ Z::Z::a::b::a::x;
F = x\ Y::Z::Y::a::b::a::x;
F = x\ Y::Z::Z::Y::a::b::a::x;
F = x\ Y::Z::U::Z::Y::a::b::a::x;
F = x\ Y::Z::U::U::Z::Y::a::b::a::x
\end{lprolog}
The second query has infinitely many answers. It finds the \lp{F} whose list, appended after
\lp{a::b::a}, makes a palindrome.

\begin{pitfall}[Types do not enforce the invariants]
\lp{(dl (1::2::3::nil) (4::5::nil))} is well typed, but it does not denote a list in the intended sense.
Likewise a function that ignores its argument, as in \lp{(fdl x\ (1::2::3::nil))}, has no tail ``pointer''
\book{143}. (The book prints the second term as \lp{(fdl x\ (1::2::3))}, which is literally ill typed.)
\end{pitfall}

\begin{kanren}[A closure is an fdl you can only run]
In its intended mode (tree given, list wanted), \lp{collect'} needs no higher-order unification. Every
clause head is unified with a fresh variable, so the only requirement is $\lambda$-terms that can be
built and applied. A Scheme closure provides that, using \scm{project°} to compose once the parts are
known:
\begin{scheme}
;; Figure 5.11, collect' in its intended mode, an fdl as a Scheme closure
(define-relation (aux/fdl° bt F)
  (cond° ((=° bt 'empty) (=° F (λ (x) x)))
         ((fresh° (n l r a b) (=° bt `(bt ,n ,l ,r))
            (aux/fdl° l a) (aux/fdl° r b)
            (project° ((a* a) (b* b))
              (=° F (λ (x) (a* (cons n (b* x))))))))))
(define-relation (collect/fdl° bt L)
  (fresh° (a) (aux/fdl° bt a) (project° ((a* a)) (=° L (a* '())))))
(define tree '(bt 2 (bt 1 empty empty) (bt 3 empty empty)))
(show 1 (μkanren-run (q -1 #t) (collect/fdl° tree q)))
\end{scheme}
\begin{answers}
1 ((1 2 3))
\end{answers}
\scm{project°} makes this non-relational: it only works when the subtrees are known. The operations that
make functional difference lists interesting (head and tail by unification with \lp{x\ (H::(T x))}, last
and front with \lp{x\ F (Y::x)}, \lp{palindrome}, and \lp{collect'} run backwards) all unify a pattern
\emph{against} a $\lambda$-term. They need $\lambda$-terms in the term language, and \lp{x\ F (Y::x)}
already applies a variable to a non-variable argument (\cref{ch:hounif}).
\end{kanren}

% ======================================================================================================
\section{Higher-order unification is not a panacea}\label{sec:hoprog-panacea}

It is tempting to make higher-order unification do all the work. The book's examples show why that is
often a mistake \book{143--145}. Take \lp{kind i type. type a,b i. type f i -> i -> i.}. The aim is to
abstract all occurrences of \lp{a} in \lp{(f a (f a b))}, that is, to obtain \lp{(x\ f x (f x b))}. The
naive program:
\begin{lprolog}
type extract_a  i -> (i -> i) -> o.
extract_a (F a) F.

?- extract_a (f a (f a b)) F.
F = x\ f x (f x b);
F = x\ f x (f a b);
F = x\ f a (f x b);
F = x\ f a (f a b);
no
\end{lprolog}
With $n$ occurrences there are $2^n$ ``extractions''. A cut, \lp{extract_a (F a) F :- !.}, relies on the
order in which the unifiers are enumerated, so it is not declarative. Two better options:
\begin{itemize}
\item \textbf{Reorder the quantifiers} so that the eigenvariable check does the selection:
\begin{lprolog}
pi a\ sigma F\ (F a) = (f a (f a b)).   % four proofs
sigma F\ pi a\ (F a) = (f a (f a b)).   % exactly one: F = x\ (f x (f x b))
\end{lprolog}
In the second line \lp{a} is a constant that is new for \lp{F}, so \lp{F} cannot mention it and every
occurrence must be abstracted. (This is also a higher-order pattern, \lp{F} applied to a distinct bound
variable, which is what makes it well behaved; see \cref{ch:lambdatree}.)
\item \textbf{Recurse on the syntax}, using the signature of \lp{i}:
\begin{lprolog}
extract_a a x\x.
extract_a b x\b.
extract_a (f T S) (x\ f (U x) (V x)) :- extract_a T U, extract_a S V.
\end{lprolog}
\end{itemize}
A second example is rewriting by congruence \book{145}:
\begin{lprolog}
type rewrite  int -> int -> o.
rewrite (0 + X) X.
rewrite (1 * X) X.
rewrite (X - X) 0.
rewrite (C X) (C Y) :- rewrite X Y.
\end{lprolog}
The last clause uses \lp{C} as a ``context with a hole''. \lp{rewrite ((1 * 5) - 5) (5 - 5)} is derivable,
but it has \emph{infinitely many} derivations: the last clause can be used any number of times with
\lp{C = x\x} before \lp{C = x\ x - 5} is finally chosen. The remedy is again explicit recursion over the
structure of terms.

\begin{kanren}[Interleaving as a home for non-unitary unification]
Every example in this section is a case where unification returns a \emph{stream} of solutions. In a
miniKanren, alternative unifiers are naturally branches of a disjunction, and interleaving enumerates
infinite families fairly. That fits Huet's procedure \cite{Huet1975} well. The lesson of the book still
applies: design relations so that unification has at most one answer (patterns, structural recursion, the
right quantifier order), and treat branching unification as a last resort. \Cref{ch:hounif} and
\cref{sec:roadmap-later-huet} take this up.
\end{kanren}

% ======================================================================================================
\section{Comparison with functional programming}\label{sec:hoprog-fp}

Most higher-order idioms of Scheme or ML carry over to logic programming easily (\lp{mappred},
\lp{sublist}, folds). The book argues that the logic programming notion is stronger in two ways
\book{146--147}:
\begin{enumerate}
\item predicate variables are only one kind of functional variable. Non-predicate function variables
  give new kinds of computation (\cref{sec:hoprog-functions}), developed in \cref{ch:lambdatree} for
  syntax with binders;
\item predicate expressions can be \emph{compared}, intensionally:
\begin{lprolog}
type eq_pred (A -> o) -> (A -> o) -> o.
eq_pred R R.
\end{lprolog}
  succeeds if and only if its arguments are unifiable predicate \emph{expressions}. For example,
  \lp{?- eq_pred (x\ p x, q x) (x\ q x, p x).} fails whatever the clauses for \lp{p} and \lp{q}. The
  comparison is structural, modulo $\lambda$-conversion, not equality of sets. Functional languages forbid
  comparing functions because their semantics is extensional, and extensional equality is undecidable.
\end{enumerate}
A design remark follows \book{147}. If the only goal were functional-style higher-order programming, a
more restricted syntax would do, for example one that forbids quantification over higher-order
non-predicate types. What general higher-order quantification adds is $\lambda$-terms that represent
syntax with binding, together with functional quantification to manipulate them. That was ``much of the
impetus for the development of λProlog''. Nor is λProlog a merger of functional and logic
programming: typed $\lambda$-terms are too weak as functions for that \book{148}.

\begin{kanren}[Unification on closures is pointer identity]
\micro{} unifies procedures by \scm{eq?}, which is finer than λProlog's $\lambda$-convertibility:
\begin{scheme}
;; Section 5.10: unification on procedures is pointer identity (eq?)
(define-relation (age° n a)
  (cond° ((=° n 'bob) (=° a 23)) ((=° n 'sue) (=° a 24))))
(define (older-than n)
  (λ (x) (fresh° (a) (age° x a) (project° ((a* a)) (if (> a* n) ✓° ✗°)))))
(show 1 (μkanren-run (q -1 #t) (=° age° age°) (=° q 'same)))
(show 2 (μkanren-run (q -1 #t)
          (=° (older-than 20) (older-than 20)) (=° q 'same)))
(show 3 (μkanren-run (q -1 #t)
          (let ((p (older-than 20))) (=° p p)) (=° q 'same)))
\end{scheme}
\begin{answers}
1 (same)
2 ()
3 (same)
\end{answers}
Two evaluations of the same \scm{λ} with a captured value give distinct closures (line 2), whereas
\lp{eq_pred (x\ p x) (x\ p x)} succeeds in λProlog. Intensional comparison, and every feature that
inspects a predicate (printing it, matching on it, \lp{eq_pred}), needs predicates represented as terms.
\end{kanren}

% ======================================================================================================
\section{Modules: programming in the large with logic}\label{sec:hoprog-modules}

Languages usually start with programming-in-the-small, and later ``a second modules language is often
imposed on the previously existing core language''. Its constructs may have ``little or no connection''
with the semantics of the core \book{150}. Chapter~6 of the book argues the opposite for logic
programming: the abstractions needed for modularity are \emph{already present} in \hohh$^+$
\book{150--151}:
\begin{itemize}
\item composing code is \emph{conjunction} of program clauses;
\item visibility of names is \emph{existential quantification} over program clauses (in effect, $\forall$
  in goals);
\item using a module to answer a query is an \emph{implicational goal} $E\imp G$.
\end{itemize}
The modules language is designed to support abstraction, information hiding, parametrization,
interfaces and separate compilation. It must not complicate the meaning of the core language, and it must
let higher-order programming coexist with modules (in Prolog, getting \texttt{call/1} to interact
correctly with modules). Finally, it should support a notion of module equivalence (representation
independence) \book{150--151}.

\subsection{Modules, signatures and accumulation}

A module is a named sequence of kind, type and clause declarations, and a signature with the same name is
its interface \book{152--153}. The book's module \lp{smlists} declares and defines \lp{id},
\lp{memb}, \lp{member} (a deterministic membership test that uses the cut), an auxiliary
\lp{revapp}, \lp{reverse} (via \lp{revapp}), \lp{append} and \lp{memb_and_rest} (select an element and
return the rest). Its signature is
\begin{lprolog}
sig smlists.
type id              list A -> list A -> o.
type memb, member    A -> list A -> o.
type reverse         list A -> list A -> o.
type append          list A -> list A -> list A -> o.
type memb_and_rest   A -> list A -> list A -> o.
end
\end{lprolog}
The signature restricts the external visibility of the module's names to exactly those it mentions, so
\lp{revapp} is hidden. Hiding applies to predicates and also to constants and data constructors, and
``such hidden symbols remain hidden even through the results of computations'' \book{153}. A query
\lp{[smlists] ?- revapp ...} is \emph{ill formed}: a static error, not a run-time failure \book{154}.
\lp{accumulate m1, ..., mn.} conceptually inserts the code of the listed modules, after \emph{renaming}
their private constants, so that the accumulated copy is entirely local to the accumulating module
\book{154}. \lp{accum_sig} inserts signatures, which is how names are re-exported \book{155--156}.

A module \emph{matches} a signature when its implicit signature (its own declarations, plus the
\emph{explicit} signatures of the modules it accumulates) is \emph{mergeable} with it: kinds are
identical, types are identical up to renaming of type variables, and operator declarations agree
\book{156--158}. The book's running example \book{158}:
\begin{lprolog}
module m1.
kind   item   type.
type   p,q    item -> o.
type   a      item.
type   r      list item -> o.
p X :- q X.
r (a :: nil).
end

sig m1.
kind   item   type.
type   p,q    item -> o.
type   a      item.
end
\end{lprolog}
\begin{lprolog}
module m2.
kind   item   type.
type   q, r   item -> o.
type   a      item.
q X :- r X.
r a.
end

sig m2.
kind   item   type.
type   q, r   item -> o.
type   a      item.
end
\end{lprolog}
\begin{lprolog}
module m3.
accumulate m1, m2.
type   s, t   item -> o.
type   b      item.
s X :- p X.
t b.
end

sig m3.
kind   item   type.
type   s, t   item -> o.
type   a      item.
end
\end{lprolog}
\lp{r} is hidden in \lp{m1} at type \lp{list item -> o}, while \lp{m2} exports an \lp{r : item -> o}. The
two never meet. If \lp{sig m1} also exported \lp{r}, the signatures of \lp{m1} and \lp{m2} would not be
mergeable and \lp{m3} would be ill formed \book{159--160}. Hiding is what lets two modules use the same
name at incompatible types. The name \lp{q}, on the other hand, is exported by both \lp{m1} and
\lp{m2} at the same type, so in \lp{m3} it is \emph{one} constant. It is the only link between
\lp{m1}'s clause \lp{p X :- q X} and \lp{m2}'s clause \lp{q X :- r X}.

\subsection{The logical reading: \texorpdfstring{$\exists$}{exists} over clauses, \texorpdfstring{$\imp$}{implication} for queries}\label{sec:hoprog-eformulas}

Two equivalences, valid in most logics, motivate the construction \book{160}:
\[
\exists y\,\forall x\,(D(x)\imp G(y)) \equiv \exists y\,((\exists x\,D(x))\imp G(y))
\qquad
\exists x\,(D_1(x)\wedge D_2) \equiv (\exists x\,D_1(x))\wedge D_2
\]
(when $x$ is not free in $G(y)$, respectively in $D_2$). The first says that a universal in a goal makes
$x$ \emph{local to the clause} $D(x)$. The second lets the scope of such an $\exists$ be narrowed to part of
a large conjunction.

\begin{definition}[E-formulas {\book{161}}]\label{def:hoprog-eformula}
The goals of \hohh$^+$ are extended to
\begin{align*}
G &::= \top \mid A \mid G\wedge G \mid G\vee G \mid \exists x\,G \mid E\imp G \mid \forall x\,G\\
D &::= A_r \mid G\imp D \mid D\wedge D \mid \forall x\,D\\
E &::= D \mid E\wedge E \mid \exists x\,E
\end{align*}
A rigid head may now also be a variable bound by an $\exists$ of an enclosing E-formula. E-formulas are
not user syntax. They appear only in the translation of modules \book{162}.
\end{definition}

A program contains only $D$-formulas, so ${\imp}R$ is replaced by a small phase that decomposes the
antecedent. It uses a sequent $\hoprogE{\Sig}{\Prog}{\Delta}{G}$, where $\Delta$ is a multiset of E-formulas
\book{161}:
\begin{mathpar}
\inferrule{\hoprogE{\Sig}{\Prog}{E}{G}}{\seq{\Sig}{\Prog}{E\imp G}}\;\rulename{$\imp$R$'$}
\and
\inferrule{\seq{\Sig}{\Prog}{G}}{\hoprogE{\Sig}{\Prog}{}{G}}\;\rulename{finish}
\and
\inferrule{\hoprogE{\Sig}{\Prog,D}{\Delta}{G}}{\hoprogE{\Sig}{\Prog}{D,\Delta}{G}}\;\rulename{reclassify}
\and
\inferrule{\hoprogE{\Sig}{\Prog}{E_1,E_2,\Delta}{G}}{\hoprogE{\Sig}{\Prog}{E_1\wedge E_2,\Delta}{G}}\;\rulename{$\wedge$L$'$}
\and
\inferrule{\hoprogE{\Sig,y}{\Prog}{E[y/x],\Delta}{G}}{\hoprogE{\Sig}{\Prog}{\exists x\,E,\Delta}{G}}\;\rulename{$\exists$L$'$}\ (y\notin\Sig)
\end{mathpar}
All these rules are \emph{invertible}, so this phase never backtracks. ``During this new phase, the local
constants are introduced dynamically into computation, and the clauses using these new constants are
integrated with the ambient logic program'' \book{161}. Note that $\exists L'$ is the same mechanism as
$\forall R$ for goals: a fresh constant is added to $\Sig$.

\begin{definition}[A module as a formula {\book{162}}]\label{def:hoprog-modformula}
Let $h_1,\ldots,h_k$ be the hidden constants of a module $M$, that is, those in its implicit signature but
not in its explicit one. If $M$ accumulates $M_1,\ldots,M_n$ and has its own clauses $C_1,\ldots,C_m$, then
\[
\hoprogT{M} = \exists h_1\cdots\exists h_k\,\big(\hoprogT{M_1}\wedge\cdots\wedge\hoprogT{M_n}\wedge C_1\wedge\cdots\wedge C_m\big).
\]
A query \lp{[M] ?- g X.} is the goal $\exists x\,(\hoprogT{M}\imp g\,x)$ \book{163}.
\end{definition}

For the running example \book{162--163}:
\begin{align*}
\hoprogT{m1} &= \exists r\,\big((\forall x(q\,x\imp p\,x))\wedge (r\,(a::nil))\big)\qquad
\hoprogT{m2} = \forall x(r\,x\imp q\,x)\wedge(r\,a)\\
\hoprogT{m3} &= \exists p\exists q\exists r\exists b\Big(\exists r\big((\forall x(q\,x\imp p\,x))\wedge(r\,(a::nil))\big)\\
&\qquad\qquad\wedge(\forall x(r\,x\imp q\,x)\wedge(r\,a))\wedge(\forall x(p\,x\imp s\,x)\wedge(t\,b))\Big)
\end{align*}
The inner $\exists r$ (\lp{m1}'s private \lp{r}) shadows the outer one (\lp{m2}'s exported \lp{r}, hidden
again by \lp{m3}), as $\alpha$-conversion dictates. \lp{a}, \lp{s} and \lp{t} stay free, because they are
\lp{m3}'s exports.

\begin{tutorial}[Using a hidden constant versus exposing it]
The book reports \book{163--164}:
\begin{lprolog}
[m3] ?- sigma x\ t x.      % succeeds (x instantiated to the hidden constant b)
[m3] ?- t X.               % FAILS
\end{lprolog}
\lp{[m3] ?- p a.} is ill formed (\lp{p} is not in \lp{sig m3}), and \lp{[m3] ?- s a.} succeeds. The
difference between the first two lies in where the quantifiers sit.
\begin{itemize}
\item \lp{t X} is $\exists x\,(\hoprogT{m3}\imp t\,x)$. The logic variable $X$ is created \emph{first}, at a
  signature that does not contain \lp{b}. Then ${\imp}R'$ and $\exists L'$ introduce a fresh $b'$.
  Backchaining on \lp{t b'} needs $X := b'$, which is not a term over $X$'s signature, so the goal fails.
\item \lp{sigma x\ t x} is $\hoprogT{m3}\imp\exists x\,(t\,x)$. $\exists L'$ adds $b'$ to the signature
  \emph{before} the existential is processed, so $b'$ is an admissible witness and the goal succeeds.
\end{itemize}
In the book's words, this is ``the difference between using a constant in the course of searching for a
solution \ldots\ and attempting to expose the identity of this constant externally'' \book{164}.
\end{tutorial}

\begin{tutorial}[{The rules of the E-phase at work on \texttt{[m3] ?- s a.}}]
The query has no logic variables, so it is the goal $\hoprogT{m3}\imp s\,a$. Read the derivation
bottom-up, starting from the empty program:
\begin{enumerate}
\item ${\imp}R'$ moves $\hoprogT{m3}$ into the E-multiset: $\hoprogE{\Sig}{\emptyset}{\hoprogT{m3}}{s\,a}$.
\item $\exists L'$ four times introduces fresh constants $p',q',r',b'$ for \lp{m3}'s hidden names and
  adds them to $\Sig$.
\item $\wedge L'$ splits the body into its conjuncts: \lp{m1}'s formula
  $\exists r\,(\ldots)$, the clauses $\forall x(r'\,x\imp q'\,x)$ and $r'\,a$ from \lp{m2}, and the
  clauses $\forall x(p'\,x\imp s\,x)$ and $t\,b'$ of \lp{m3}.
\item $\exists L'$ once more introduces a fresh $r''$ for \lp{m1}'s private \lp{r}, and $\wedge L'$
  yields $\forall x(q'\,x\imp p'\,x)$ and $r''\,(a::nil)$.
\item \rulename{reclassify} moves each of these $D$-formulas into the program, and \rulename{finish}
  returns to an ordinary sequent $\seq{\Sig,p',q',r',b',r''}{\Prog_{m3}}{s\,a}$.
\item Backchaining now proves $s\,a$ from $p'\,a$, $p'\,a$ from $q'\,a$, $q'\,a$ from $r'\,a$, and
  $r'\,a$ is a fact.
\end{enumerate}
Steps 1--5 make no choices, which is what invertibility means in practice: loading a module never
backtracks. The fact $r''\,(a::nil)$ is never used, and it could not be: no clause outside \lp{m1}'s
own formula can mention $r''$.
\end{tutorial}

\paragraph{Elaboration.} Rather than replay \cref{def:hoprog-modformula} for every query, an implementation
inlines the accumulated modules and \emph{renames} their private names. The elaborated \lp{m3} contains
the clauses \lp{p X :- q X.}, \lp{r' (a :: nil).}, \lp{q X :- r X.}, \lp{r a.}, \lp{s X :- p X.} and
\lp{t b.}, where \lp{r'} is a new name for \lp{m1}'s local \lp{r} \book{165}. The book's implementation
section adds that, at run time, the machinery that already handles alternating $\exists/\forall$ in
goals also handles local constants, provided they are annotated properly when the module is used to
prove queries \book{172}.

\subsection{Abstract datatypes}

Exporting a type constructor while hiding its data constructors gives an abstract datatype \book{166}:
\begin{lprolog}
sig stack.
kind store           type -> type.
type init            store A -> o.
type add, remove     A -> store A -> store A -> o.
end
\end{lprolog}
\begin{lprolog}
module stack.
kind store           type -> type.
type emp             store A.
type stk             A -> store A -> store A.
type init            store A -> o.
type add, remove     A -> store A -> store A -> o.
init emp.
add    X S (stk X S).
remove X (stk X S) S.
end
\end{lprolog}
\lp{[stack] ?- init A.} \emph{fails}, because it tries to exhibit the representation of the empty stack.
\lp{[stack] ?- sigma A\ sigma B\ sigma C\ init A, add 1 A B, remove X B C.} uses the representation
anonymously and succeeds with \lp{X = 1} \book{166}. A client such as \lp{graph_search} can hold store
values in its variables, but only \lp{init}, \lp{add} and \lp{remove} can examine them \book{167}.

\paragraph{Extension and parametrization.} An exported predicate can be \emph{extended} by later modules.
A propositional prover \lp{proplogic} exports the type \lp{form} and the predicate \lp{prove}, and a
module \lp{quantlogic} that accumulates it adds both constructors and clauses \book{168--169}:
\begin{lprolog}
sig quantlogic.
accum_sig proplogic.
kind term        type.
type all, some   (term -> form) -> form.
end

module quantlogic.
accumulate proplogic, smlists.
kind term        type.
type all, some   (term -> form) -> form.

prove L F :- memb_and_rest (some P) L L',
             pi c \ prove ((P c)::L) F.
prove L (all P) :- pi c\ prove L (P c).
prove L (some P) :- prove L (P T).
prove L F :- memb_and_rest (all P) L L',
             append L' [all P] L'', prove ((P T)::L'') F.
end
\end{lprolog}
This works because \lp{form} and \lp{prove} are in \lp{proplogic}'s \emph{explicit} signature, so
accumulation neither renames nor hides them. The recursive calls inside \lp{proplogic}'s clauses reach
the new quantifier clauses too. Both provers accumulate \lp{smlists}, which gives two private copies. The
book calls this ``at least wasteful'' \book{169} and offers an alternative: accumulate the
\emph{signature} \lp{smlists} into both modules, and supply one implementation when everything is
accumulated at the top \book{170}. Parametrization is then done by shared names that are closed at the
top, not by passing arguments explicitly.

\subsection{Higher-order programming and predicate visibility}\label{sec:hoprog-callp}

The test case for the second desideratum \book{171}:
\begin{lprolog}
sig comblibrary.
type  call  o -> o.
end

module comblibrary.
type  call  o -> o.
type  p  list int -> o.
p (1 :: nil).
call Q :- Q.
end

sig test.
type  test  list int -> o.
end

module test.
accumulate comblibrary.
type  test  list int -> o.
type  p  list int -> o.
p (2 :: nil).
test X :- call (p X).
end
\end{lprolog}
For \lp{[test] ?- test X.}, which \lp{p} does \lp{call (p X)} invoke: the one visible where \lp{call} is
defined (\lp{X = 1 :: nil}), or the one visible where the name is written (\lp{X = 2 :: nil})? ``The second
interpretation is the commonly used resolution'', and it ``is a natural consequence of the semantics
that we have presented'' \book{171}. In the elaborated module, \lp{comblibrary}'s \lp{p} is renamed
\lp{p'}, so the only answer is \lp{2 :: nil} \book{171--172}. \lp{p X} is a \emph{term} of type
\lp{o} whose head is fixed when the program is parsed. It is not a name looked up at run time.

\begin{kanren}[Modules: what Scheme scoping gives, and what it does not]
Scheme's lexical scope already does module elaboration with renaming: a local definition is private, and
each evaluation of a scope creates fresh closures. Goals are closures, so Figure~6.10 resolves to the
book's ``second interpretation'' automatically. A Scheme-level abstract type is equally easy:
\begin{scheme}
;; Figure 6.10: a "module" is a lexical scope returning its exported relations
(define call°                            ; module comblibrary, exports call
  (let ()
    (define-relation (p l) (=° l '(1)))  ; hidden
    (define (call g) g)
    call))
(define test°                            ; module test, exports test
  (let ()
    (define-relation (p l) (=° l '(2)))
    (define-relation (test x) (call° (p x)))
    test))
(show 1 (μkanren-run (q -1 #t) (test° q)))
;; Figure 6.6: the stack "abstract datatype", constructors made by gensym
(define-values (init° add° remove°)
  (let ((emp (gensym 'emp)) (stk (gensym 'stk)))
    (define-relation (init s) (=° s emp))
    (define-relation (add x s s*) (=° s* (list stk x s)))
    (define-relation (remove x s s*) (=° s (list stk x s*)))
    (values init add remove)))
(show 2 (μkanren-run (q -1 #t)
          (fresh° (a b c) (init° a) (add° 1 a b) (remove° q b c))))
(show 3 (μkanren-run (q -1 #f) (init° q)))
\end{scheme}
\begin{answers}
1 ((2))
2 (1)
3 ((λ () (quote emp223)))
\end{answers}
Line 1 is the book's answer, and line 2 uses the stack anonymously, as the \lp{sigma} query does. Line 3
shows the honest difference. Scheme hides the \emph{name} \scm{emp}, but the answer to \lp{init A} still
\emph{exhibits} the representation, where λProlog fails. Matching the book needs hidden constructors that
are eigen-constants, introduced when the module is used, and a scope check in unification that stops a
query variable created earlier from being bound to them. That is the same machinery as $\forall$ in
goals (\cref{sec:hoprog-hohhplus}). Two more things are not free. One is \emph{extensible} predicates
(\lp{quantlogic} adding clauses to \lp{proplogic}'s \lp{prove}), because a \scm{define-relation} is
closed. The other is parametrization by shared names closed at the top. Signature matching is mostly
about types and kinds, which an untyped miniKanren can only approximate (names and arities). The proposal
is in \cref{ch:design,ch:module}. One rule follows directly from Figure~6.10: never resolve a relation from
a symbol at run time, because that reintroduces exactly the Prolog \texttt{call/1} ambiguity.
\end{kanren}

% ======================================================================================================
\section{What a Scheme-hosted miniKanren gets for free}\label{sec:hoprog-free}

\Cref{tab:hoprog-free} collects the kanren boxes of this chapter.

\begin{table}[ht]
\centering\small
\caption{Book features versus a miniKanren whose relations are Scheme procedures.}\label{tab:hoprog-free}
\begin{tabular}{@{}>{\raggedright\arraybackslash}p{0.40\textwidth}>{\raggedright\arraybackslash}p{0.17\textwidth}>{\raggedright\arraybackslash}p{0.36\textwidth}@{}}
\toprule
Feature & With closures & What is needed beyond closures\\
\midrule
\texttt{mappred}, \texttt{foreach}, \texttt{forsome}, \texttt{sublist}, folds, \texttt{compose} & free & nothing (curried \texttt{define} for partial application)\\
$\lambda$-terms of predicate type as arguments & free & nothing (Scheme \texttt{λ})\\
closures \texttt{ref}/\texttt{sym}/\texttt{trans}, binary clauses (CPS) & free & nothing\\
hidden predicates with fixed clauses (\texttt{reverse}) & free & lexical scope\\
\texttt{if}, \texttt{not} & free & \texttt{if°} with \texttt{take° 1}\\
\texttt{mapfun}/\texttt{reducefun} forwards, \texttt{collect'} forwards & free & closures that build terms; \texttt{project°}\\
module hiding, predicate visibility (Fig.~6.10) & free & lexical scope\\
\midrule
\texttt{rel R, R john mary} & opaque answers & relations as terms, to print or compare them\\
flexible goals & accidental error & inspectable goals; a store of suspended goals\\
growing hidden predicates (\texttt{fib\_memo}), $\forall$/$\imp$ encodings of $\vee$, $\exists$ & no & goal-level $\forall$ (eigenvariables, scope check) and $\imp$ (dynamic clauses)\\
abstract datatypes that stay hidden in answers & partial & eigen-constants and a scope check\\
extensible predicates across modules & no & open clause sets, linking\\
\texttt{eq\_pred}, reasoning by substitution & no & predicates and clauses as $\lambda$-terms\\
\texttt{mapfun}/\texttt{reducefun} backwards, \texttt{extract\_a}, \texttt{palindrome}, fdl destructuring & no & $\lambda$-terms, $\beta\eta$-normalisation, higher-order unification\\
\bottomrule
\end{tabular}
\end{table}

The dividing line is simple. Whatever only \emph{calls} a predicate or a function is free, because a
Scheme procedure can be called. Whatever must \emph{look inside} one (unify it, print it, compare it,
synthesise it) needs $\lambda$-terms in the term language. A third group, $\forall$ and $\imp$ in goals,
needs neither $\lambda$-terms nor inspection but does need new components in the search state:
eigenvariables with a scope check, and a dynamic clause store. \Cref{ch:design} builds on this
classification.

% ======================================================================================================
\section{Summary}\label{sec:hoprog-summary}
\begin{itemize}
\item Atoms are \emph{rigid} (constant head) or \emph{flexible} (variable head). Clause heads must be rigid,
  which keeps programs consistent and modular; flexible atoms are allowed in goals.
\item Implications may not occur inside terms. Instantiation terms are drawn from Herbrand universes:
  $\mathcal{H}_1$ (no $\forall,\imp$) for \hohc{} and $\mathcal{H}_2$ (no $\imp$) for \hohh{}. Goal-directed
  search is then complete: classical and intuitionistic for \hohc{}, intuitionistic for \hohh{}.
\item \hohh$^+$ lets a variable bound by an essentially universal quantifier head assumed clauses, which
  gives hidden, closed predicate definitions.
\item Predicate variables give \lp{mappred}, \lp{sublist}, closures, \lp{compose}, \lp{foldl} and the CPS
  transformation. Logic variables inside a $\lambda$ are shared.
\item An unbound flexible goal is underconstrained. Suspending it is complete; in practice it is reported
  as an error. Constrain predicate variables first (\lp{rel R, R john mary}).
\item Hidden predicates support reasoning by substitution (the symmetry of \lp{reverse}). Logical
  constants are definable. \lp{if} and \lp{not} via cut are non-logical.
\item Function variables are evaluated by $\beta\eta$ and can be synthesised by unification
  (\lp{mapfun}, \lp{reducefun}, functional difference lists). Higher-order unification is not unitary, so
  prefer structural recursion and the right quantifier order.
\item Unlike functions in functional programs, predicate expressions can be compared
  intensionally, with \lp{eq_pred}.
\item Modules are logic: conjunction composes code, $\exists$ over clauses hides names, and a query
  against a module is $\exists x\,(\hoprogT{M}\imp G)$. Hidden constants may be used but not exposed, and
  predicate names resolve statically.
\item In a Scheme-hosted miniKanren, calling is free and looking inside is not; $\forall$ and $\imp$ in
  goals need new state (\cref{tab:hoprog-free}, \cref{ch:design}).
\end{itemize}

% ch06-lambdatree.tex — λ-tree syntax: computing with binders (Miller & Nadathur, Chapter 7)
\chapter{\texorpdfstring{$\lambda$}{λ}-tree syntax: computing with binders}\label{ch:lambdatree}

The previous chapters gave us two kinds of machinery. \Cref{ch:lambda} and \cref{ch:hoprog} gave us
simply typed $\lambda$-terms as \emph{data}, with an equality that already contains $\alpha$-, $\beta$- and
$\eta$-conversion. \Cref{ch:hh} gave us the two goal connectives of hereditary Harrop formulas: the
universal goal \lp{pi x\ G}, proved with a fresh \emph{eigenvariable}, and the implicational goal
\lp{D => G}, proved with the program temporarily extended by \lp{D}. This chapter, which follows Chapter~7 of
the book \cite{MillerNadathur2012} (\book{175}--\book{210}), shows what the two kinds of machinery can do
together: they let a logic program compute over syntax that contains \emph{binders}.

This is the chapter that motivates the whole proposal of this document. Every program in it
manipulates formulas with quantifiers or $\lambda$-terms. None of them contains a single line about variable
names, renaming, or capture-avoiding substitution. All of that is \enquote{relegated to the meta-level}.
A miniKanren that wants to run these programs must therefore provide the meta-level: $\lambda$-terms as data,
eigenvariables with a scope discipline, and a scoped program extension. We will read the book's programs
closely, walk through how \lp{pi} and \lp{=>} make a binder \enquote{mobile}, and, in the \emph{Kanren
angle} boxes, record what each concept asks of \micro.

\section{The problem: syntax with binders}\label{sec:lambdatree-problem}

Many symbolic objects contain binders \book{175}--\book{176}: the $x$ in $\int f(x)\,dx$, in
$\frac{d}{dx}$ and in $\sum_{x=1}^{100}$; quantifiers; $\lambda$-abstractions; local declarations and formal
parameters in programming languages; type quantifiers of polymorphic types; name restriction in the
$\pi$-calculus. Programs that handle such objects (differentiators, compilers, interpreters, theorem
provers, type checkers) all need the same three notions \book{176}:
\begin{enumerate}[nosep]
  \item the distinction between \emph{free} and \emph{bound} occurrences of a variable;
  \item the irrelevance of bound names ($\alpha$-conversion): the formulas $\forall x.\,p(x)$ and
        $\forall y.\,p(y)$ are the same;
  \item \emph{capture-avoiding substitution}, needed both for instantiating quantifiers and for passing
        parameters.
\end{enumerate}
The book observes that \enquote{a computational treatment of such binding-related notions is complicated,
and programming language support for realizing them can be valuable} \book{176}. Anyone who has written a
substitution function over a named representation, with its renaming case, knows why.

The book's answer has three parts. We state them now and spend the chapter on them.
\begin{enumerate}[nosep]
  \item \textbf{Encoding.} Each object-level binder becomes a constant applied to a
        $\lambda$-abstraction of the meta-level. So \lp{all (x\ B)} encodes the formula $\forall x.\,B$, and
        \lp{abs (x\ t)} encodes the term $\lambda x.\,t$.
  \item \textbf{Substitution is $\beta$.} Instantiating the bound variable of an abstraction \lp{B} with
        \lp{T} is just writing the application \lp{(B T)}; the meta-level $\beta$-conversion does the rest.
  \item \textbf{Mobility.} To recurse \emph{under} a binder, the program turns the term-level $\lambda$ into a
        formula-level \lp{pi}, which proof search turns into a proof-level eigenvariable; facts about the
        new eigenvariable are assumed with \lp{=>}.
\end{enumerate}

\section{Encoding objects with binding structure}\label{sec:lambdatree-encoding}

\subsection{First-order formulas}\label{sec:lambdatree-fol}

The object logic is first-order logic over two sorts, \lp{term} and \lp{form}. The propositional
connectives are ordinary constants (re-listed on \book{177}):
\begin{lprolog}
type   ff,          % encoding the false proposition
       tt   form.   % encoding the true proposition
type   &&,          % encoding conjunction
       !!,          % encoding disjunction
       ==>  form -> form -> form.   % encoding implication
type   neg  form -> form.           % encoding negation
infixl &&  5.
infixl !!  4.
infixr ==> 3.
\end{lprolog}
A sample non-logical signature \book{177}:
\begin{lprolog}
type p  term -> form.
type q  term -> term -> form.
type f  term -> term.
type a  term.
\end{lprolog}
Terms of type \lp{term -> form} are \enquote{abstractions of a term over a formula} \book{178}, for example
\begin{lprolog}
x\ (p x) ==> (p (f x))
Z\ (p Z) ==> (p (f Z))
Z\ (p Z) && (p (f Z))
\end{lprolog}
The first two are \emph{the same} $\lambda$-term, since they differ only in the bound name. Note that the
capital \lp{Z} after \lp{Z\} is a bound variable, not a logic variable. The quantifiers are then
\emph{second-order constants}: their argument has a functional type \book{178}.
\begin{lprolog}
type all, some  (term -> form) -> form.
\end{lprolog}
\begin{lprolog}
all x\ (p x) ==> (p (f x))          % ∀x (p(x) ⊃ p(f(x)))
some Z\ (p Z) && (p (f Z))          % ∃Z (p(Z) ∧ p(f(Z)))
\end{lprolog}
The book states the principle behind this as follows: \enquote{the mechanism that we have described divides
the encoding of the object logic quantifiers into two parts: a second order constant that identifies the
quantifier and an abstraction that captures its binding effect} \book{178}. For instance
$\forall x\,\forall y\,(q(x,y)\imp\exists z\,(p(z)\wedge q(z,y)))$ is encoded as
\begin{lprolog}
all x\ all y\ (q x y) ==> some z\ p z && q z y.
\end{lprolog}
where \lp{\} extends as far to the right as possible and \lp{&&} (priority 5) binds tighter than
\lp{==>} (priority 3).

\begin{definition}[The encoding of formulas]\label{def:lambdatree-encoding}
For first-order terms and formulas over the signature above, define $\ulcorner\cdot\urcorner$ by
\[
\begin{aligned}
\ulcorner x\urcorner &= x, &
\ulcorner f(t_1,\dots,t_n)\urcorner &= f\,\ulcorner t_1\urcorner\cdots\ulcorner t_n\urcorner,\\
\ulcorner B\wedge C\urcorner &= \ulcorner B\urcorner \mathbin{\texttt{\&\&}} \ulcorner C\urcorner
\quad(\text{similarly for } \vee,\imp,\neg,\top,\bot), &
\ulcorner\forall x.\,B\urcorner &= \cst{all}\,(\abs{x}{\ulcorner B\urcorner}),\\
& & \ulcorner\exists x.\,B\urcorner &= \cst{some}\,(\abs{x}{\ulcorner B\urcorner}).
\end{aligned}
\]
\end{definition}

The book does not state $\ulcorner\cdot\urcorner$ as a definition or prove anything about it; this formal
version is ours, and so is the name. The chapter relies on two properties of it without stating them.
First, it is a bijection between closed object formulas over the signature, taken up to renaming of bound
variables, and closed terms of type \lp{form} over the same signature, taken up to $\alpha\beta\eta$-conversion.
A canonical representative of each class is the $\beta$-normal, $\eta$-long term, which is what
$\ulcorner\cdot\urcorner$ produces. Second, it commutes with substitution:
$\ulcorner \subst{B}{t}{x}\urcorner = \subst{\ulcorner B\urcorner}{\ulcorner t\urcorner}{x}$. The second
property is what makes \enquote{substitution is $\beta$} correct.

\subsection{Untyped \texorpdfstring{$\lambda$}{λ}-terms}\label{sec:lambdatree-untyped}

The second running object language is the untyped $\lambda$-calculus \book{178}--\book{179}:
\begin{lprolog}
kind tm   type.
type app  tm -> tm -> tm.
type abs  (tm -> tm) -> tm.
\end{lprolog}
\begin{center}
\begin{tabular}{@{}ll@{}}
\toprule
untyped term & encoding \\
\midrule
$\lambda x\, x$ & \lp{(abs x\x)} \\
$\lambda x\,(x\,x)$ & \lp{(abs x\ (app x x))} \\
$\lambda x\,\lambda y\, x$ & \lp{(abs x\ abs y\ x)} \\
$\lambda x\,\lambda y\, y$ & \lp{(abs x\ abs y\ y)} \\
$\lambda x\,\lambda y\,(y\,x)$ & \lp{(abs x\ abs y\ app y x)} \\
$\lambda x\,\lambda y\,\lambda z\,((x\,z)\,(y\,z))$ & \lp{(abs x\abs y\abs z\ (app (app x z) (app y z)))} \\
$(\lambda x\,(x\,x))\,(\lambda x\,(x\,x))$ & \lp{(app (abs x\ app x x) (abs x\ app x x))} \\
\bottomrule
\end{tabular}
\end{center}
Each object application becomes one \lp{app}; each object $\lambda$ becomes one \lp{abs} followed by one
meta-level $\lambda$. Two remarks help in reading the rest of the chapter.
\begin{itemize}[nosep]
  \item No closed term of type \lp{tm} can be built from \lp{app} alone: every closed term starts with an
        \lp{abs}, whose argument (of type \lp{tm -> tm}) is something like \lp{x\x} or \lp{x\ app x x}.
  \item \enquote{the type tm can be seen as admitting new constructors as one descends under binders}
        \book{179}. Below \lp{abs x\ abs y\} both \lp{x} and \lp{y} behave as extra constants of type
        \lp{tm}. Mobility (\cref{sec:lambdatree-mobility}) makes this literally true: under a binder the
        program works with a genuinely new constant.
\end{itemize}

\subsection{What the encoding gives for free}\label{sec:lambdatree-properties}

The abstractions over syntax live at arrow types (\lp{term -> form}, \lp{tm -> tm}), and the equality of
those types includes the three conversions of \cref{ch:lambda}.
\begin{description}[nosep,font=\normalfont\itshape]
  \item[$\alpha$.] The object language identifies formulas up to renaming of bound variables, and so does
        the meta-language. \enquote{Since $\alpha$-conversion is included within the logical notion of equality,
        it is impossible to access the names of bound variables} \book{180}. The book presents this as a
        strength: the programmer is spared what does not matter, and the implementer may drop names entirely
        (for instance with de Bruijn indices \cite{deBruijn1972}). The cost is that printing and parsing must
        invent names.
  \item[$\beta$.] This is the subject of \cref{sec:lambdatree-beta}.
  \item[$\eta$.] Its effect is \enquote{rather mild} \book{180}: for $f:\typ{i}\to\typ{i}$,
        $f=_\eta\abs{x}{f\,x}$. Hence $\forall x\,(p\,x)$ may be written \lp{(all x\ p x)} \emph{or}
        \lp{(all p)}: \enquote{not only does the hohh language abstract away from names of bindings, but it also
        sometimes allows one to abstract away from the presence of actual binders} \book{180}.
\end{description}

\begin{pitfall}[An $\eta$-short quantifier is still a quantifier]
A head pattern \lp{(all B)} must also match the term \lp{(all p)}, and \lp{(B T)} must then reduce to
\lp{(p T)}. An implementation either unifies modulo $\eta$ or keeps terms $\eta$-long. A representation that
stores the binder \enquote{shape} syntactically and expects every quantifier to carry an explicit
$\lambda$ will reject \lp{(all p)}.
\end{pitfall}

\section{Object-level substitution is meta-level \texorpdfstring{$\beta$}{β}}\label{sec:lambdatree-beta}

\subsection{Instantiating quantifiers}

The formula $\forall x\,(p\,x\imp\exists y\,(q\,x\,y))$ is \lp{all x\ p x ==> some y\ q x y}. To instantiate
it with $f(a)$ we simply apply the abstraction to \lp{(f a)} \book{180}--\book{181}:
\begin{lprolog}
(x\ p x ==> some y\ q x y) (f a)
\end{lprolog}
which $\beta$-converts to \lp{p (f a) ==> some y\ q (f a) y}. (The book prints \lp{p (h a)} here, a typo;
the page image confirms it, and it must be \lp{p (f a)}.) Note what was \emph{not} needed: no renaming of
the inner \lp{y}, no check that \lp{(f a)} has no free \lp{y}. Capture avoidance is the meta-level's job.

The predicate \lp{list_instan} relates a list of $n\ge 0$ terms and a formula with at least $n$ outer
universal quantifiers to the result of instantiating the first $n$ \book{181}:
\begin{lprolog}
type list_instan  list term -> form -> form -> o.
list_instan nil B B.
list_instan (T::Ts) (all B) C :- list_instan Ts (B T) C.
\end{lprolog}
The head \lp{(all B)} is $\eta$-equal to \lp{(all x\ B x)}: matching it binds \lp{B} to the abstraction
itself, and \enquote{the substitution for a bound variable is realized simply by writing down (B T)}
\book{181}.

\subsection{An interpreter for an object logic}\label{sec:lambdatree-interp}

A more substantial example is an interpreter for first-order Horn clauses, written in higher-order Horn
clauses (Figure~7.1, \book{181}). The object-level program is an explicit argument \lp{D}, a single formula.
\begin{lprolog}
type interp     form -> form -> o.
type backchain  form -> form -> form -> o.
type atom       form -> o.

interp D tt.
interp D (G1 && G2)  :- interp D G1, interp D G2.
interp D (G1 !! G2)  :- interp D G1; interp D G2.
interp D (some   G)  :- interp D (G X).
interp D A           :- atom A, backchain D D A.
backchain D A A.
backchain D (D1 && D2) A   :- backchain D D1 A; backchain D D2 A.
backchain D (all D1) A     :- backchain D (D1 X) A.
backchain D (G ==> D1) A   :- backchain D D1 A, interp D G.
\end{lprolog}
The two higher-order clauses are the interesting ones. An object goal \lp{(some G)} is solved by solving
\lp{(G X)} for a \emph{fresh logic variable} \lp{X}: the meta-level $\beta$ plugs the logic variable in for the
object-bound variable, and later unification will find a witness. An object clause \lp{(all D1)} is used by
instantiating it with a fresh logic variable as well. The sample object program of Figure~7.2 (\book{182})
is
\begin{lprolog}
type a, b, c    term.
type adj, path  term -> term -> form.
type prog       form -> o.
atom (adj _ _) & atom (path _ _).
prog ((adj a b) && (adj b c) &&
      (all x\ all y\ (adj x y) ==> (path x y)) &&
      (all x\ all y\ all z\ (adj x y) && (path y z) ==> (path x z))).
\end{lprolog}
Two notational points recur in the rest of the chapter. An \lp{&} between two clause heads, as in
\lp{atom (adj _ _) & atom (path _ _).}, is clause-level conjunction. The heads share one body (here an empty
one), so the line stands for two clauses written together. The underscore \lp{_} is an anonymous logic
variable. The predicate \lp{atom} is not part of the interpreter. It is defined for each object signature by
listing its predicate constants, and here it says that \lp{adj} and \lp{path} atoms are atomic.
The query \lp{?- prog P, interp P (path a X).} has two answers, \lp{X = b} and then \lp{X = c} \book{182}.

\begin{tutorial}[How \texttt{interp} finds \texttt{X = b}]
The goal \lp{interp P (path a X)} matches the last \lp{interp} clause, because \lp{atom (path a X)} holds. We
must prove \lp{backchain P P (path a X)}. The program \lp{P} is a left-nested \lp{&&} of four parts, so the
second \lp{backchain} clause descends into it. The two \lp{adj} facts do not match \lp{path a X}. The third part
is \lp{all x\ all y\ (adj x y) ==> (path x y)}: the \lp{(all D1)} clause applies it to a fresh \lp{X1}, giving
\lp{(D1 X1)} $=_\beta$ \lp{all y\ (adj X1 y) ==> (path X1 y)}; once more with \lp{Y1}, giving
\lp{(adj X1 Y1) ==> (path X1 Y1)}. The \lp{==>} clause then requires \lp{backchain P (path X1 Y1) (path a X)}
(the fact clause, binding \lp{X1 = a}, \lp{Y1 = X}) and then \lp{interp P (adj a X)}, which backchains to the fact
\lp{adj a b} with \lp{X = b}. On backtracking, the fourth part of the program yields \lp{X = c} through
\lp{adj a b} and \lp{path b c}. At no point did the interpreter rename a bound variable: each $\beta$-step
\lp{(D1 X1)} replaced an object-level bound variable by a logic variable.
\end{tutorial}

\subsection{Evaluators for the untyped \texorpdfstring{$\lambda$}{λ}-calculus}

Call-by-name and call-by-value evaluation are four clauses (Figure~7.3, \book{182}):
\begin{lprolog}
type cbn, cbv  tm -> tm -> o.
cbn (abs R) (abs R).
cbn (app M N) V :- cbn M (abs R), cbn (R N) V.
cbv (abs R) (abs R).
cbv (app M N) V :- cbv M (abs R), cbv N U, cbv (R U) V.
\end{lprolog}
\lp{(R N)} is the object-level $\beta$-contraction, done by the meta-level $\beta$. (The text on \book{182}
says \enquote{the term U such that eval N U is provable}; \lp{eval} is a loose name for \lp{cbv}.) The behaviour
is the expected one \book{183}. $\Omega=(\lambda x\,(x\,x))(\lambda x\,(x\,x))$ has no value under either
evaluator: the search for a derivation is infinite. On $(\lambda x\,\lambda w\,w)\,\Omega$, call-by-name
terminates and call-by-value loops:
\begin{lprolog}
?- cbn (app (abs x\ abs w\w) (app (abs x\ app x x) (abs x\ app x x))) V.
   % V = (abs w\w)
?- cbv (app (abs x\ abs w\w) (app (abs x\ app x x) (abs x\ app x x))) V.
   % infinite search
\end{lprolog}
The book draws the moral of this section in one paragraph \book{183}. The specifications are succinct and
declarative because \enquote{all the details regarding object-level bound variables and object-level
substitution has been relegated to the meta-level}. They are still formal and precise: the details
\enquote{have not gone away; they simply have become the burden of an implementation of a metalogic such as
$\lambda$Prolog}, and the programmer \enquote{is relieved of dealing with the details behind such
\textquote{concrete nonsense} as binder names, capture avoiding substitutions, etc.}

\begin{kanren}[What works today: abstractions as Scheme closures]
All the programs of this section use an abstraction only as an \emph{input}. The abstraction is already known
when \lp{(R N)} is formed. A miniKanren can imitate this with \emph{functional} higher-order abstract syntax:
represent \lp{abs R} by a list holding a Scheme procedure, and perform \lp{(R N)} with \scm{project°}. This
runs on today's \micro, unchanged:
\begin{scheme}
(define (abs f) `(abs ,f))
(define (app m n) `(app ,m ,n))
(define ω (abs (λ (x) (app x x))))
(define Ω (app ω ω))
(define K* (abs (λ (x) (abs (λ (w) w)))))          ; λx λw w

; call-by-name evaluation, Figure 7.3: (R N) is Scheme application, available only through project°
(define-relation (cbn° m v)
  (cond°
    ((fresh° (r) (=° m `(abs ,r)) (=° v m)))
    ((fresh° (m1 n r) (=° m `(app ,m1 ,n)) (cbn° m1 `(abs ,r))
       (project° ((r* r)) (cbn° (r* n) v))))))

; printing a HOAS term needs invented names: apply each closure to a generated symbol
(define (hoas->sexp t)
  (let loop ((t t) (k 1))
    (cond
      ((and (pair? t) (eq? (car t) 'abs))
       (let ((w (string->symbol (string-append "w" (number->string k)))))
         `(λ (,w) ,(loop ((cadr t) w) (+ k 1)))))
      ((and (pair? t) (eq? (car t) 'app)) `(,(loop (cadr t) k) ,(loop (caddr t) k)))
      (else t))))

(print (μkanren-run (q -1 #t) (fresh° (v) (cbn° (app K* Ω) v) (project° ((v* v)) (=° q (hoas->sexp v*))))))
; α-equivalent closures are different Scheme objects: unification fails
(print (μkanren-run (q -1 #t) (=° (abs (λ (x) x)) (abs (λ (y) y)))))
\end{scheme}
\begin{answers}
((λ (w1) w1))
()
\end{answers}
The first answer is the book's \lp{V = (abs w\w)}. The second answer is the problem. Two $\alpha$-equivalent
abstractions are distinct procedures, and \micro{} compares procedures with \scm{eq?}, so they do not unify. A
closure cannot be printed without being applied to invented names (that is what \scm{hoas->sexp} does). A
closure cannot be inspected except by applying it, and no unification can \emph{produce} one. Recognisers
such as \lp{term}, \lp{bnorm} and \lp{typeof} below could still be run on closures. They only apply an
input abstraction to a fresh constant, although they also need \lp{pi} and \lp{=>}. But every program that
\emph{builds} an abstraction as output (\lp{red1}, path synthesis, \lp{trans} backwards, \lp{nnf},
\lp{prenex}, \lp{copy}, \lp{subst} backwards), and every program that asks whether a body depends on its
bound variable (\lp{vacuous}, the $\eta$-redex), is out of reach. \Cref{sec:lambdatree-hoas} comes back to
this \enquote{functional} reading of higher-order abstract syntax.
\end{kanren}

\section{Mobility of binders}\label{sec:lambdatree-mobility}

\subsection{Why an abstraction cannot be taken apart}

Equal expressions should have equal parts: equal non-empty lists have equal heads and equal tails. Applied to
$\lambda$-terms, this principle says that \enquote{we should not be able to decompose an abstraction into the
variable that it binds and its body} \book{183}. The terms \lp{(all x\ some y\ p x ==> p y)} and
\lp{(all z\ some y\ p z ==> p y)} are equal, yet the names \lp{x} and \lp{z} differ, and so do the would-be
bodies \lp{(some y\ p x ==> p y)} and \lp{(some y\ p z ==> p y)}. \enquote{Such decompositions, in fact, have no
logical status within the hohh language: Because equality includes $\alpha$-conversion, the name for a bound
variable is a fiction, and an implementation may not represent it at all} \book{183}.

So how does a program recurse into the body of a binder? The book's answer \book{184}:
\enquote{during proof search, a term-level binding (i.e., a $\lambda$-abstraction) can be converted into a
formula-level binding (i.e., a quantifier), which then can be converted into a proof-level binding (i.e., an
eigenvariable).}
\[
\underbrace{\abs{y}{t}}_{\text{term level}}
\;\rightsquigarrow\;
\underbrace{\forall x.\,(\cdots\imp P\,((\abs{y}{t})\,x))}_{\text{formula level}}
\;=_\beta\;
\forall x.\,(\cdots\imp P\,(\subst{t}{x}{y}))
\;\rightsquigarrow\;
\underbrace{P\,(\subst{t}{c}{y})\ \text{ with } c \text{ fresh}}_{\text{proof level}}
\]

\subsection{The proof-search rules behind mobility}

Recall from \cref{ch:hh} the rules for the two connectives that make mobility work, together with the
backchaining rule. A goal is proved relative to a signature $\Sig$ (constants and eigenvariables in scope)
and a program $\Prog$.
\begin{mathpar}
\inferrule{\seq{\Sig\cup\{c\}}{\Prog}{\subst{G}{c}{x}} \\ c\notin\Sig}{\seq{\Sig}{\Prog}{\forall x.\,G}}\;\rulename{generic}
\and
\inferrule{\seq{\Sig}{\Prog\cup\{D\}}{G}}{\seq{\Sig}{\Prog}{D\imp G}}\;\rulename{augment}
\and
\inferrule{\seq{\Sig}{\Prog}{\subst{G}{t}{x}} \\ t \text{ a } \Sig\text{-term}}{\seq{\Sig}{\Prog}{\exists x.\,G}}\;\rulename{instance}
\end{mathpar}
For an atomic goal $A$, \rulename{backchain} picks a clause $\forall\bar x\,(G\imp A')$ of $\Prog$ and an instance
with $A'\theta\eqb A$, then proves $G\theta$. In an implementation the term $t$ of \rulename{instance} and the
substitution $\theta$ are found by unification with logic variables. Then the side condition $c\notin\Sig$
becomes the \textbf{scoping constraint}: a logic variable created \emph{before} the eigenvariable $c$ must never
be instantiated with a term containing $c$, because at the time it was created its witness had to be a
$\Sig$-term. This is why \lp{sigma X\ pi y\ X = y} fails while \lp{pi y\ sigma X\ X = y} succeeds. Chapter~7
never mentions the constraint, but several of its programs (\lp{vacuous}, the $\eta$-redex clause,
\lp{prenex}) are correct only because of it.

\subsection{A first example: recognising terms}\label{sec:lambdatree-term}

Because of typing, the predicate \lp{term : tm -> o} could be defined by the single clause \lp{term T.}. A more
instructive definition reflects the two constructors of \lp{tm} into two clauses \book{184}. Written so that
their logical content is visible:
\begin{lprolog}
pi M\ term M => pi N\ term N => term (app M N).
pi R\ (pi x\ term x => term (R x)) => term (abs R).
\end{lprolog}
or, equivalently, in the usual clause syntax:
\begin{lprolog}
term (app M N) :- term M, term N.
term (abs R)   :- pi x\ term x => term (R x).
\end{lprolog}

\begin{tutorial}[Following the binder of \texttt{(abs y\textbackslash\ app y y)}]
The book traces the query \lp{?- (term (abs y\ app y y)).} \book{184}--\book{185}.
\begin{lprolog}
?- (term (abs y\ app y y)).
?- pi x\ term x => term ((y\ app y y) x).     % backchain on abs clause
?- pi x\ term x => term (app x x).            % β-conversion: term-level y -> formula-level x
?- term (app c c).                            % eigenvariable c, program + (term c)
\end{lprolog}
\begin{enumerate}[nosep]
  \item \emph{Backchaining} on the \lp{abs} clause unifies \lp{R} with the abstraction \lp{y\ app y y}. Nothing
        has been decomposed: \lp{R} holds the whole abstraction.
  \item The body of the clause applies \lp{R} to the \lp{pi}-bound \lp{x}. The redex \lp{((y\ app y y) x)} is
        a $\beta_0$-redex (\cref{sec:lambdatree-llambda}): a $\lambda$ applied to a bound variable. Contracting
        it moves the term-level binder \lp{y} to the formula-level binder \lp{x}.
  \item \rulename{generic} replaces \lp{x} by a new eigenvariable \lp{c}. \rulename{augment} adds the clause
        \lp{term c} to the program. The formula-level binder has become a proof-level binder.
  \item \lp{term (app c c)} backchains on the \lp{app} clause, and both subgoals \lp{term c} are closed by the
        assumed clause.
\end{enumerate}
Here \lp{c} is exactly one of the \enquote{new constructors} of \lp{tm} that appear under a binder. The program
has no clause for variables, and it needs none: each variable gets its clause at the moment its binder is
crossed, and loses it when the proof of the body is complete.
\end{tutorial}

The book distils the example into a programming idiom. We record it as a definition because every program
of the chapter is an instance of it.

\begin{definition}[The binder-mobility idiom, \book{185}]\label{def:lambdatree-idiom}
\enquote{In order to continue a recursive analysis within the scope of a binder, first, apply that binder to a
universal quantified goal variable, for which the metalogic will substitute an eigenvariable, and second,
using an implicational goal, assume new clauses that extend the definition of various predicates so that
they can deal with the presence of this new eigenvariable.}
\end{definition}

In code, the idiom is always the shape
\lp{p (c R) ... :- pi x\ (hyp x => p (R x) ...).} The abstraction \lp{R} is never opened. It is applied to a
fresh eigenvariable, and the assumed \lp{hyp x} tells the predicates of the program what to do when they meet
it.

\begin{pitfall}[A fresh constant is not an eigenvariable]
Proving \lp{pi x\ G} with a fresh constant is only half of \rulename{generic}. The other half is the scoping
constraint. Here is what happens without it, in today's \micro, if a universal goal is \enquote{implemented}
by handing a \scm{gensym} to the body:
\begin{scheme}
; a NAIVE universal goal: a fresh Scheme constant, with no eigenvariable scope check
(define (naive-pi f) (f (gensym 'c)))

; λProlog: ?- sigma X\ pi y\ X = y.   must fail; the naive encoding succeeds
(print (μkanren-run (q -1 #t) (naive-pi (λ (y) (=° q y)))))
; λProlog: ?- pi y\ sigma X\ X = y.   must succeed (X is created inside the scope of y)
(print (μkanren-run (q -1 #t) (naive-pi (λ (y) (fresh° (x) (=° x y) (=° q 'ok))))))
\end{scheme}
\begin{answers}
(c102)
(ok)
\end{answers}
The first query should fail. The query variable \scm{q} was created before the constant, yet it now holds that
constant, which has escaped its scope. The same mistake in a unifier, letting a clause variable be bound to a
term that mentions a variable bound inside its scope, makes \lp{vacuous} (\cref{sec:lambdatree-fol-comp})
accept formulas whose quantifier is not vacuous. The constraint is not a matter of efficiency: it is part of
the meaning of $\forall$.
\end{pitfall}

\begin{kanren}[What mobility requires of a miniKanren]
Mobility needs three things that \micro{} lacks (see \cref{ch:micro} for the engine).
\begin{enumerate}[nosep]
  \item \emph{Eigenvariables}: constants created during the search, distinct from every other term, with a
        scope. The engine must enforce the scoping constraint: a logic variable older than an eigenvariable
        can never be bound to a term containing it, \emph{even indirectly} through younger variables (binding
        $X:=f(Z)$ and later $Z:=c$ leaks $c$ into $X$). λProlog implementations use universe levels
        \cite{NadathurJayaramanKwon1995}. The miniKanren with symbolic constraints uses an absento-like
        \enquote{eigen} constraint over the variables in scope \cite{HemannByrdFriedman2013}.
  \item \emph{Scoped program extension}: \lp{D => G} adds \lp{D} only while \lp{G} is being proved. Other branches
        of the search and the answers of the enclosing goal must not see \lp{D}. Because miniKanren's search
        interleaves branches, the extension must live in the search \emph{state}, not in a host-language
        dynamic variable. A relation must be able to consult the assumed clauses as well as its static ones.
  \item \emph{Application of an abstraction to an eigenvariable}, \lp{(R x)}, on both sides of the equation:
        as input (when \lp{R} is known) and as output (when \lp{R} is a logic variable to be solved, as in most
        programs below).
\end{enumerate}
How the proposal provides each of these is the subject of \cref{ch:design}.
\end{kanren}

\section{Computing with untyped \texorpdfstring{$\lambda$}{λ}-terms}\label{sec:lambdatree-untyped-comp}

This section applies the idiom to the untyped $\lambda$-calculus. The programs grow from pure recognisers
to programs that build abstractions as output, to programs that assume rules with bodies.

\subsection{\texorpdfstring{$\beta$}{β}-normal forms}\label{sec:lambdatree-bnorm}

\begin{definition}[$\beta$-normal, $\beta$-body-normal, \book{185}]\label{def:lambdatree-bnorm}
An untyped $\lambda$-term is in \emph{$\beta$-normal form} if it contains no $\beta$-redex. Equivalently and
positively, it has the shape $\lambda x_1\ldots\lambda x_n\,(v\,t_1\ldots t_m)$ with $n,m\ge0$, where $v$ is a
variable and every $t_i$ is $\beta$-normal. The inner part $v\,t_1\ldots t_m$ is called
\emph{$\beta$-body-normal}.
\end{definition}

The definition becomes three inference rules and three clauses \book{186}:
\begin{mathpar}
\inferrule{\forall x.\ \mathrm{bbnorm}(x)\imp\mathrm{bnorm}(M\,x)}{\mathrm{bnorm}(\cst{abs}\,M)}
\and
\inferrule{\mathrm{bbnorm}(H)}{\mathrm{bnorm}(H)}
\and
\inferrule{\mathrm{bbnorm}(M) \\ \mathrm{bnorm}(N)}{\mathrm{bbnorm}(\cst{app}\,M\,N)}
\end{mathpar}
\begin{lprolog}
type bnorm, bbnorm  tm -> o.

bnorm (abs M)    :- pi x\ bbnorm x => bnorm (M x).
bnorm H          :- bbnorm H.
bbnorm (app M N) :- bbnorm M, bnorm N.
\end{lprolog}
There is no static clause saying that a variable is $\beta$-body-normal, and there could not be one, since
bound variables are not terms of the program. A variable becomes body-normal only through the assumption
\lp{bbnorm x} made when its binder is crossed. A redex \lp{app (abs R) N} is rejected because
\lp{bbnorm (abs R)} has no clause.

\begin{tutorial}[A normal term and a redex]
This walk-through is ours; the book gives no query for \lp{bnorm}. Take first
\lp{?- bnorm (abs x\ app x x).}
\begin{enumerate}[nosep]
  \item The \lp{abs} clause matches with \lp{M = x\ app x x}. The body is
        \lp{pi x\ bbnorm x => bnorm (M x)}: an eigenvariable \lp{c} is created, \lp{bbnorm c} is assumed, and the
        goal is \lp{bnorm (app c c)} (after $\beta_0$).
  \item \lp{bnorm (app c c)} is not an \lp{abs}, so the second clause reduces it to \lp{bbnorm (app c c)}, and the
        third clause to \lp{bbnorm c} and \lp{bnorm c}.
  \item \lp{bbnorm c} is the assumption. \lp{bnorm c} goes through the second clause to \lp{bbnorm c} again, also
        the assumption. The query succeeds.
\end{enumerate}
Now take \lp{?- bnorm (app (abs x\x) (abs x\x)).} The \lp{abs} clause does not match. The second clause asks for
\lp{bbnorm (app (abs x\x) (abs x\x))}, and the third clause for \lp{bbnorm (abs x\x)}. No static clause has a
head \lp{bbnorm (abs _)}, and the only assumptions ever made are \lp{bbnorm c} on eigenvariables, so the query
fails. The structure of the clauses alone, \enquote{a variable is body-normal only when its binder said so},
rules the redex out.
\end{tutorial}

A pleasant consequence is a substitution lemma that costs nothing.

\begin{theorem}[Substitution lemma \enquote{for free}, \book{186}]\label{thm:lambdatree-bnorm-subst}
Let $\Gamma=\{\cst{bbnorm}\,c_1,\ldots,\cst{bbnorm}\,c_n\}$, $n\ge0$, with the $c_i$ eigenvariables. If
\lp{(bnorm (abs R))} and \lp{(bbnorm S)} are both derivable from the program above together with $\Gamma$, then
\lp{(bnorm (R S))} is derivable too. In object-level terms: if $\lambda x\,t$ is $\beta$-normal and $s$ is
$\beta$-body-normal, then $\subst{t}{s}{x}$ is $\beta$-normal.
\end{theorem}
\begin{proof}[Proof (book's argument, with the inversion step made explicit)]
A derivation of \lp{bnorm (abs R)} must end with the \lp{abs} clause, so
\lp{pi x\ bbnorm x => bnorm (R x)} is derivable. Derivability is intuitionistic, so the universal quantifier
can be instantiated with \lp{S}, giving \lp{bbnorm S => bnorm (R S)}, and modus ponens with \lp{bbnorm S}
gives \lp{bnorm (R S)}.
\end{proof}
The book says only \enquote{since derivability is based on intuitionistic logic}. Two facts sit behind that
phrase: cut admissibility, and the fact that an eigenvariable may be instantiated by any term. The inversion
step (\enquote{must end with the \lp{abs} clause}) also uses the shape of $\Gamma$. $\Gamma$ contains only
atoms \lp{bbnorm c} on eigenvariables, so no assumption can match \lp{bnorm (abs R)}. In a named
representation this lemma is a page of index bookkeeping. Here the logic proves it.

\subsection{Reduction to normal form}\label{sec:lambdatree-reduce}

Figure~7.4 (\book{187}) specifies one-step reduction and reduction to normal form:
\begin{lprolog}
type redex, red1, reduce  tm -> tm -> o.
redex (app (abs R) N) (R N).
red1 M N :- redex M N.
red1 (app M N) (app P N) & red1 (app N M) (app N P) :- red1 M P.
red1 (abs M) (abs N) :- pi x\ red1 (M x) (N x).
reduce M M :- bnorm M.
reduce M N :- red1 M P, reduce P N.

type repeat   (A -> A -> o) -> A -> A -> o.
type reduce'  tm -> tm -> o.
repeat Pred M N :- Pred M P, !, repeat Pred P N.
repeat Pred M M.
reduce' M N :- repeat red1 M N.
\end{lprolog}
\begin{itemize}[nosep]
  \item \lp{redex} relates a $\beta$-redex to its contractum, again by meta-level $\beta$: \lp{(R N)}.
  \item \lp{red1} contracts exactly one redex somewhere in the term. The \lp{abs} clause is new. Its
        \emph{output} is the abstraction \lp{N}, and the goal \lp{red1 (M x) (N x)} asks unification to solve
        \lp{(N x) = t} for an unknown \lp{N}: the answer is \lp{N} $:=$ \lp{x\ t}, with the eigenvariable abstracted
        out of \lp{t}. This is a \emph{higher-order pattern} problem (\cref{sec:lambdatree-llambda,ch:hounif}).
        There is no \lp{=>} here: \lp{red1} should simply fail on a bare variable, so no clause about \lp{x} is
        needed.
  \item \lp{reduce} uses \lp{bnorm} to decide when to stop. \lp{reduce'} uses the polymorphic higher-order
        \lp{repeat}, which iterates its argument until it fails. The cut \lp{!} \enquote{prune[s] away alternative
        solutions to the parent goal} \book{186}. (The text there says \lp{repeat'}; the figure's predicate is
        \lp{repeat}.)
\end{itemize}

\begin{tutorial}[One step under a binder: \texttt{red1 (abs x\textbackslash\ app (abs y\textbackslash y) x) N}]
This is our own example of an output abstraction, traced against the clauses of Figure~7.4 without the
$\eta$-clause below. The term is $\lambda x\,((\lambda y\,y)\,x)$, with a redex under the binder.
\begin{enumerate}[nosep]
  \item \lp{red1 M N :- redex M N} fails: the term is an \lp{abs}, not an \lp{app}. The two \lp{app} clauses do
        not match either.
  \item The \lp{abs} clause matches with \lp{M = x\ app (abs y\y) x} and \lp{N = abs N1} for a new logic
        variable \lp{N1 : tm -> tm}. It creates an eigenvariable \lp{c}. The goal is
        \lp{red1 (app (abs y\y) c) (N1 c)}.
  \item The \lp{redex} clause applies with \lp{R = y\y} and contractum \lp{(R c)}, which is \lp{c}. Now unification
        must solve \lp{(N1 c) = c}. This is a pattern equation: \lp{N1} is applied to the eigenvariable \lp{c}
        only. Its unique solution abstracts \lp{c} out of the right-hand side: \lp{N1 := x\ x}.
  \item The answer is \lp{N = abs x\x}. On backtracking, the \lp{app} clauses try to reduce \lp{(abs y\y)} or
        \lp{c} in one step. The first descends under \lp{y\} to the goal \lp{red1 d _} on a new eigenvariable
        \lp{d}, the second gives \lp{red1 c _}. Both fail, because no clause, static or assumed, applies to a
        bare eigenvariable. There is no other answer.
\end{enumerate}
The program contains no code for \enquote{rebuild the binder around the reduced body}. Unification does it in
step 3, and the scoping constraint guarantees that \lp{N1}, created before \lp{c}, receives \lp{c} only through
its argument.
\end{tutorial}

The $\eta$-redex needs one clause \book{187}:
\begin{lprolog}
redex (abs x\ app M x) M.
\end{lprolog}
It is correct because of the scoping constraint. \lp{M} is a clause variable, quantified \emph{outside} the
binder \lp{x}, so no unifier can let \lp{M} depend on \lp{x}. This is exactly the side condition of $\eta$:
$x$ not free in $M$. For instance, \lp{abs x\ app x x} is not an $\eta$-redex: matching it against
\lp{abs x\ app M x} would need \lp{M = x}, and the scoping constraint forbids that binding. With this clause
added, the term of the tutorial above is also an $\eta$-redex. Matching \lp{abs x\ app (abs y\y) x} against
\lp{abs x\ app M x} gives \lp{M = abs y\y}, which does not mention \lp{x}, so \lp{red1} now finds
\lp{abs y\y} first, through \lp{redex}, and then the $\alpha$-equal \lp{abs x\x} again through the \lp{abs}
clause.

The book adds a semantic remark \book{186}--\book{187}. A model of the untyped $\lambda$-calculus is a domain
$D$ with maps $f:D\to[D\to D]$ and $g:[D\to D]\to D$; \lp{tm} plays $D$, \lp{app} plays $f$, \lp{abs} plays
$g$. The $\beta$-clause states $(f\circ g)\,R=R$ and the $\eta$-clause states $(g\circ f)\,M=M$: \enquote{these
are the familiar properties stating that f and g are, in fact, retracts}.

\subsection{Reduction based on paths}\label{sec:lambdatree-paths}

A \emph{path} walks from the root of a term to an occurrence of a variable. It passes \emph{through} an
abstraction and chooses left or right at an application \book{187}. Paths contain binders too:
\begin{lprolog}
kind path type.
type bnd          (path -> path) -> path.
type left, right  path -> path.
type path         tm -> path -> o.
path (app M _) (left P) &
path (app _ M) (right P) :- path M P.
path (abs R) (bnd S) :- pi x\ pi p\ path x p => path (R x) (S p).
\end{lprolog}
The \lp{abs} clause introduces \emph{two} eigenvariables at once: a term variable \lp{x} and a path variable
\lp{p}. The assumption \lp{path x p} links them: the path to the occurrence of \lp{x} \emph{is} \lp{p}. The paths of
\lp{(abs x\ app x (abs y\ app y x))} are \book{188}
\begin{lprolog}
bnd u\ left u
bnd u\ right (bnd v\ left v)
bnd u\ right (bnd v\ right u)
\end{lprolog}
The relation also runs backwards: from a set of paths it synthesises the term, using \lp{foreach} from
\cref{ch:hoprog} (\lp{foreach P nil. foreach P (X::L) :- P X, foreach P L.}) \book{188}:
\begin{lprolog}
?- foreach (path N)
     ((bnd u\ left u) ::
      (bnd u\ right (bnd v\ left v))::
      (bnd u\ right (bnd v\ right u))::nil).
N = abs W1\ app W1 (abs W2\ app W2 W1);
no
\end{lprolog}

\begin{tutorial}[Synthesising a term from its paths]
Take the first path. \lp{path N (bnd u\ left u)} can only use the \lp{abs} clause, so \lp{N = abs R} and
\lp{S = u\ left u}. Under eigenvariables \lp{x}, \lp{p} with \lp{path x p} assumed, the goal is
\lp{path (R x) (left p)}. The \lp{left} clause requires \lp{(R x) = app M _}. Unification solves this pattern
problem with \lp{R := x\ app (M' x) (K x)}, two new logic variables applied to \lp{x}. Then \lp{path (M' x) p}
is solved by the assumption \lp{path x p} itself, giving \lp{M' := x\ x}. The second and third paths fill in
the right branch in the same way. Nothing in the program says \enquote{build an abstraction}. The abstraction
in the answer is produced by unification, which abstracts the eigenvariable \lp{x} out of a term that contains
it. This is the operation that Scheme closures could not perform.
\end{tutorial}

Paths can also cross a \emph{marked} $\beta$-redex. A path through \lp{(app (abs R) N)} goes through the body
\lp{R}; when it reaches the bound variable of \lp{R}, it continues into \lp{N}. Marked redexes are written
\lp{(beta N R)} (Figure~7.5, \book{189}):
\begin{lprolog}
type beta     tm -> (tm -> tm) -> tm.
type addbeta  tm -> tm -> o.
type bpath    tm -> path -> o.

bpath (app M _) (left P)  &
bpath (app _ M) (right P) :- bpath M P.
bpath (abs R)   (bnd S)   :- pi x\ pi p\ bpath x p =>
                                         bpath (R x) (S p).
bpath (beta N R) P :-
   pi x\ (pi Q\ bpath x Q :- bpath N Q) => bpath (R x) P.

addbeta (app (abs R) N) (beta M S) :- addbeta (abs R) (abs S),
                                      addbeta N M.
addbeta (app (app M N) P) (app O Q) :- addbeta (app M N) O,
                                       addbeta P Q.
addbeta (abs R) (abs S) :-
  pi x\ (pi M\ pi N\ addbeta (app x M) (app x N) :- addbeta M N) =>
        (addbeta x x) => addbeta (R x) (S x).
\end{lprolog}
Two features are new.
\begin{itemize}[nosep]
  \item The \lp{beta} clause assumes a \emph{rule}, not a fact: \lp{pi Q\ bpath x Q :- bpath N Q}, \enquote{the
        paths of \lp{x} are the paths of \lp{N}}. The assumed clause has a body and its own universally
        quantified variable \lp{Q}, renamed at each use. It also mentions \lp{N}, a variable of the
        enclosing clause, which it shares rather than renames.
  \item \lp{addbeta} marks every redex. It dispatches on the head of an application: an \lp{abs} head is
        marked, an \lp{app} head recurses, and a bound-variable head \lp{x} is handled by the rule assumed
        when \lp{x}'s binder was crossed. A bare variable is copied by the assumed fact \lp{addbeta x x}.
\end{itemize}
Normalisation then has two stages: collect the \lp{bpath}s, then synthesise with \lp{path} \book{189}:
\begin{lprolog}
?- sigma B\ addbeta (app (abs x\x) (abs x\x)) B, bpath B Path.
Path = bnd W1\ W1;
no
?- foreach (P\ path T P) (bnd (W1\ W1) :: nil).
T = abs W1\ W1;
no
\end{lprolog}
For $K\,(S\,K)$ with $K=\lambda x\lambda y\,x$ and $S=\lambda x\lambda y\lambda z\,((x\,z)\,(y\,z))$
\book{189}--\book{190}, the query
\begin{lprolog}
?- sigma K\ sigma S\ sigma B\
     K = (abs x\ abs y\ x),
     S = (abs x\ abs y\ abs z\ app (app x z) (app y z)),
     addbeta (app K (app S K)) B, bpath B Path.
\end{lprolog}
yields four paths, and synthesis gives
\begin{lprolog}
?- foreach (path T)
     ((bnd W1\ bnd W2\ bnd W3\ left (left (bnd W4\ bnd W5\ W4)))::
      (bnd W1\ bnd W2\ bnd W3\ left (right W3))::
      (bnd W1\ bnd W2\ bnd W3\ right (left W2))::
      (bnd W1\ bnd W2\ bnd W3\ right (right W3))::nil).
T = abs W1\ abs W2\ abs W3\
      app (app (abs W4\ abs W5\ W4) W3) (app W2 W3);
no
\end{lprolog}
The result $\lambda w_1\lambda w_2\lambda w_3\,(((\lambda w_4\lambda w_5\,w_4)\,w_3)\,(w_2\,w_3))$ is \emph{not}
$\beta$-normal. The redex $(K\,w_3)$ was \emph{created} by the reduction, so it was not marked. Iterating
\enquote{mark, collect, synthesise} reaches the normal form if it terminates, and in the untyped calculus it
need not terminate \book{190}.

\begin{pitfall}[Collecting all answers is outside the logic]
The two-stage method collects all solutions of \lp{bpath B Path} into a list. The book stresses that this
\enquote{is not formally possible in the simple proof-theoretic framework that we have described here}
\book{190}. Some Prolog systems provide \lp{bagof}, and persistent storage can simulate it. A miniKanren has
the same situation. A goal that consumes a whole sub-search, such as the aggregation goals of \micro{}
(\cref{ch:micro}), is available, but it is extra-logical. It also terminates only when the sub-search is
finite.
\end{pitfall}

\subsection{Type inference}\label{sec:lambdatree-typeof}

The book's type inference for the untyped calculus (\book{190}--\book{191}) is the best-known
instance of the idiom:
\begin{lprolog}
kind ty type.
type arr     ty -> ty -> ty.
type typeof  tm -> ty -> o.
typeof (app M N) A :- typeof M (arr B A), typeof N B.
typeof (abs M) (arr A B) :- pi x\ typeof x A => typeof (M x) B.
\end{lprolog}
Read as inference rules, the typing context $\Gamma$ is the set of assumed atoms \lp{typeof c A}:
\begin{mathpar}
\inferrule{\Gamma\vdash M:B\to A \\ \Gamma\vdash N:B}{\Gamma\vdash \cst{app}\,M\,N : A}
\and
\inferrule{\Gamma,\ c:A\vdash M\,c:B \\ c \text{ fresh}}{\Gamma\vdash\cst{abs}\,M : A\to B}
\end{mathpar}
The context is not a data structure of the program. It is the set of clauses that the \lp{=>} goals have added,
and the variable rule of the type system is \rulename{backchain} on one of them. The interaction is
\book{191}:
\begin{lprolog}
?- typeof (abs x\ abs y\ abs z\ app (app x z) (app y z)) Ty.
Ty = arr (arr T1 (arr T2 T3)) (arr (arr T1 T2) (arr T1 T3))
?- typeof (abs x\x) Ty.
Ty = arr T1 T1
?- typeof (abs x\ app x x) Ty.
no
\end{lprolog}

\begin{tutorial}[Why \texttt{typeof (abs x\textbackslash x) Ty} answers \texttt{arr T1 T1}]
\begin{enumerate}[nosep]
  \item Backchaining on the \lp{abs} clause binds \lp{Ty = arr A B} for fresh logic variables \lp{A}, \lp{B}, and
        \lp{M = x\x}.
  \item \rulename{generic} creates an eigenvariable \lp{c}. \rulename{augment} adds the clause \lp{typeof c A}.
        This assumed clause contains the \emph{logic variable} \lp{A}, which is shared with the rest of the
        derivation. Assumptions are not closed facts.
  \item The goal is now \lp{typeof ((x\x) c) B}, that is, \lp{typeof c B}. Neither static clause applies, since
        \lp{c} is neither an \lp{app} nor an \lp{abs}. The assumed clause applies and unifies \lp{B} with \lp{A}.
  \item So \lp{Ty = arr A A}, printed \lp{arr T1 T1}.
\end{enumerate}
For \lp{abs x\ app x x} the same steps give the assumption \lp{typeof c A}, and \lp{typeof (app c c) B} needs
\lp{typeof c (arr B0 B)} and \lp{typeof c B0}. With the assumption these are \lp{A = arr B0 B} and \lp{A = B0},
hence \lp{B0 = arr B0 B}. The occurs check rejects this, and the answer is \lp{no}. Without an occurs check the
answer would be a cyclic type.
\end{tutorial}

Two more details matter. The type \lp{ty} has only the constructor \lp{arr}, so it has \emph{no closed
inhabitants}. Inferred types are therefore necessarily open, which links back to the distinction between
unifiers and solutions of \cref{ch:lambda}. With an added constant \lp{type i ty.} one can ask closed
questions \book{191}:
\begin{lprolog}
?- typeof (abs x\x) (arr i i).
yes
?- typeof (abs x\x) (arr i Ty).
Ty = i
\end{lprolog}

\begin{theorem}[Subject reduction for call-by-name, \book{191}--\book{192}]\label{thm:lambdatree-subject}
If \lp{(cbn P V)} and \lp{(typeof P A)} are derivable, then \lp{(typeof V A)} is derivable.
\end{theorem}
\begin{proof}[Proof (book's argument)]
By induction on the derivation of \lp{cbn P V}. If it uses the first \lp{cbn} clause, $P=V$. Otherwise
\lp{P = (app M N)}, and there are shorter derivations of \lp{cbn M (abs R)} and \lp{cbn (R N) V}. From
\lp{typeof (app M N) A} there is a \lp{B} with \lp{typeof M (arr B A)} and \lp{typeof N B}. By the induction
hypothesis, \lp{typeof (abs R) (arr B A)}. Hence \lp{pi x\ typeof x B => typeof (R x) A} is derivable.
Instantiating \lp{x} with \lp{N} and using \lp{typeof N B} gives \lp{typeof (R N) A}. The induction hypothesis on
\lp{cbn (R N) V} gives \lp{typeof V A}.
\end{proof}
The book adds that the same holds for \lp{cbv} (left to the reader), and comments: \enquote{The hard technical
part in such a proof is establishing a substitution lemma \ldots{} Here, that substitution lemma is obtained
essentially for free by referring to well-known properties of intuitionistic logic} \book{192}.

\begin{kanren}[The usual miniKanren rendering: an explicit context]
The standard miniKanren way to write a type inferencer avoids both \lp{pi} and \lp{=>}. It uses de Bruijn
indices, so that bound names do not exist, and it passes the context as an explicit list. This runs on today's
\micro{} (indices are Peano numerals and start from \scm{z} for the innermost binder; the book's de Bruijn
indices in \cref{sec:lambdatree-debruijn} start from 1):
\begin{scheme}
(define-relation (nth° Γ k a)
  (fresh° (b Γ* k*)
    (=° Γ `(,b . ,Γ*))
    (cond° ((=° k 'z) (=° a b))
           ((=° k `(s ,k*)) (nth° Γ* k* a)))))

(define-relation (typeof° Γ m a)
  (cond°
    ((fresh° (k) (=° m `(var ,k)) (nth° Γ k a)))
    ((fresh° (m1 n b) (=° m `(app ,m1 ,n)) (typeof° Γ m1 `(arr ,b ,a)) (typeof° Γ n b)))
    ((fresh° (body a1 b) (=° m `(lam ,body)) (=° a `(arr ,a1 ,b)) (typeof° `(,a1 . ,Γ) body b)))))

(define S '(lam (lam (lam (app (app (var (s (s z))) (var z)) (app (var (s z)) (var z)))))))
(define I '(lam (var z)))
(define ω '(lam (app (var z) (var z))))

(print (μkanren-run (q -1 #t) (typeof° '() S q)))
(print (μkanren-run (q -1 #t) (typeof° '() I q)))
(print (μkanren-run (q -1 #t) (typeof° '() ω q)))
\end{scheme}
\begin{answers}
((arr (arr α (arr β γ)) (arr (arr α β) (arr α γ))))
((arr α α))
()
\end{answers}
These are the book's three answers. The occurs check of \micro{} rejects $\lambda x\,(x\,x)$ exactly as in the
tutorial above. The difference between the two programs is where the context lives. In \scm{typeof°} it is
data, threaded through every call and looked up with a hand-written \scm{nth°}. In \lp{typeof} it is part of
the program, and the logic guarantees the substitution lemma behind \cref{thm:lambdatree-subject}. With an
explicit context, every relation that may meet a bound variable must accept and pass on the context. With
\lp{=>}, a predicate written without any knowledge of binders, such as \lp{atom} in
\cref{sec:lambdatree-fol-comp}, keeps working under them. Supporting the second style is the point of the
proposal (\cref{ch:design}).
\end{kanren}

\subsection{Translating to and from de Bruijn syntax}\label{sec:lambdatree-debruijn}

In de Bruijn's notation \cite{deBruijn1972} a variable is a positive integer $n$ that refers to the $n$-th
enclosing abstraction, counting from 1 \book{192}:
\begin{center}
\begin{tabular}{@{}ll@{\qquad}ll@{}}
\toprule
term & de Bruijn & term & de Bruijn\\
\midrule
$\lambda x\,x$ & $\lambda 1$ & $\lambda x\lambda y\,(y\,x)$ & $\lambda\lambda(1\,2)$\\
$\lambda x\,(x\,x)$ & $\lambda(1\,1)$ & $\lambda x\lambda y\lambda z\,((x\,z)(y\,z))$ & $\lambda\lambda\lambda((3\,1)(2\,1))$\\
$\lambda x\lambda y\,x$ & $\lambda\lambda2$ & $\lambda x\,(x\,(\lambda y\,x))$ & $\lambda(1\,\lambda2)$\\
$\lambda x\lambda y\,y$ & $\lambda\lambda1$ & $\lambda x\,(x\,(\lambda y\,(x\,(\lambda w\,(w\,x)))))$ & $\lambda(1\,(\lambda(2\,(\lambda(1\,3)))))$\\
\bottomrule
\end{tabular}
\end{center}
The same variable receives different numbers at different depths (the $x$ of the last row is 1, 2 and 3).
The encoding \book{193}, in which \lp{deb} is both a kind and a constructor, and the translation are:
\begin{lprolog}
kind deb  type.
type ab   deb -> deb.
type ap   deb -> deb -> deb.
type deb  int -> deb.
\end{lprolog}
\begin{lprolog}
type trans  int -> tm -> deb -> o.
type depth  int -> tm -> o.
trans D (abs M) (ab P) :- pi c\ depth D c =>
                            (E is D + 1, trans E (M c) P).
trans D (app M N) (ap P Q) :- trans D M P, trans D N Q.
trans D X (deb E) :- depth N X, E is (D - N).
\end{lprolog}
In \lp{trans D t d}, \lp{D} is one more than the number of enclosing binders. Each eigenvariable is tagged, by
the assumption \lp{depth D c}, with the level at which it was bound. At an occurrence, the index is the current
level minus the binding level. The relation works in both directions \book{193}:
\begin{lprolog}
?- trans 1 (abs x\ app x (abs y\ app x (abs w\ app w x))) D.
D = ab (ap (deb 1) (ab (ap (deb 2) (ab (ap (deb 1) (deb 3))))))
?- trans 1 P
     (ab (ap (deb 1) (ab (ap (deb 2) (ab (ap (deb 1) (deb 3))))))).
P = abs W1\ app W1 (abs W2\ app W1 (abs W3\ app W3 W1))
\end{lprolog}

\begin{tutorial}[Tracing \texttt{?- trans 1 (abs x\textbackslash\ abs y\textbackslash\ abs z\textbackslash\ y) P.}]
The book traces this query \book{193}--\book{194}.
\begin{enumerate}[nosep]
  \item \lp{P := (ab P1)}. A new eigenvariable \lp{c} is introduced and \lp{depth 1 c} is assumed. Goal:
        \lp{trans 2 (abs y\ abs z\ y) P1}. (The meta-level $\beta$ has replaced the bound \lp{x} by \lp{c}; since
        \lp{x} does not occur in the body, nothing visible changes.)
  \item \lp{P1 := (ab P2)}, eigenvariable \lp{d}, assume \lp{depth 2 d}. Goal: \lp{trans 3 (abs z\ d) P2}.
  \item \lp{P2 := (ab P3)}, eigenvariable \lp{e}, assume \lp{depth 3 e}. Goal: \lp{trans 4 d P3}.
  \item Only the third clause applies: \lp{P3 := (deb E)} and \lp{depth N d, E is (4 - N)}. The assumption gives
        \lp{N = 2}, so \lp{E = 2}.
  \item Answer: \lp{P = (ab (ab (ab (deb 2))))}.
\end{enumerate}
The book's text writes the assumptions as \enquote{depth c 1}, \enquote{depth d 2}, \enquote{depth e 3}, with the
arguments swapped with respect to \lp{type depth int -> tm -> o}. The program assumes \lp{depth 1 c}, and so on.
In the reverse direction the same clauses run with \lp{P} unknown. The third clause then finds the eigenvariable
\lp{X} whose \lp{depth} gives the right index, and the \lp{(M c)} of the first clause is solved as a pattern
problem, as in the path synthesis of \cref{sec:lambdatree-paths}.
\end{tutorial}

A variant (\book{194}) with no \lp{depth} predicate assumes new \lp{trans} clauses for each
eigenvariable. It has \emph{no} static clause for variables:
\begin{lprolog}
trans D (app M N) (ap P Q) :- trans D M P, trans D N Q.
trans D (abs M) (ab P) :-
   pi u\
     (pi N\ pi H\ trans N u (deb H) :- H is (N - D))
       => (D' is D + 1, trans D' (M u) P).
\end{lprolog}
For the query traced above, the assumed clauses are
\begin{lprolog}
pi N\ pi H\ trans N c (deb H) :- H is (N - 1).
pi N\ pi H\ trans N d (deb H) :- H is (N - 2).
pi N\ pi H\ trans N e (deb H) :- H is (N - 3).
\end{lprolog}
and \lp{?- trans 4 d P3.} is solved by the second. Each assumed clause has captured the value of the outer
\lp{D} at the time it was assumed. It behaves as a closure over the context.

\begin{kanren}[The same translation with an explicit environment]
In \micro{} the translation is written with an explicit environment of names. It is the miniKanren analogue of
the \lp{depth} assumptions:
\begin{scheme}
; named terms: a symbol, (app m n), (λ x b); de Bruijn terms as in typeof.scm (z = innermost binder)
(define-relation (index° env x k)
  (fresh° (y env*)
    (=° env `(,y . ,env*))
    (cond° ((=° x y) (=° k 'z))
           ((fresh° (k*) (≠° x y) (=° k `(s ,k*)) (index° env* x k*))))))

(define-relation (trans° env t d)
  (cond°
    ((symbol° t) (fresh° (k) (=° d `(var ,k)) (index° env t k)))
    ((fresh° (m n p r) (=° t `(app ,m ,n)) (=° d `(app ,p ,r)) (trans° env m p) (trans° env n r)))
    ((fresh° (x b p) (=° t `(λ ,x ,b)) (symbol° x) (=° d `(lam ,p)) (trans° `(,x . ,env) b p)))))

; λx (x (λy (x (λw (w x))))), the p.193 example, forwards
(print (μkanren-run (q -1 #t) (trans° '() '(λ x (app x (λ y (app x (λ w (app w x)))))) q)))
; λx λy x backwards: the binder names come back as constrained logic variables
(print (μkanren-run (q -1 #f) (trans° '() q '(lam (lam (var (s z)))))))
\end{scheme}
\begin{answers}
((lam (app (var z) (lam (app (var (s z)) (lam (app (var z) (var (s (s z))))))))))
((λ (α β) (begin (deny (equal? β α))) (assert (every (μ v (symbol? v)) (list α β))) (cons (quote λ) (cons α (cons (cons (quote λ) (cons β (cons α (quote ())))) (quote ()))))))
\end{answers}
The forward answer is the book's \lp{D}, shifted to 0-based indices. The backward answer shows what names cost.
λProlog answers \lp{abs W1\ abs W2\ W1}: there are no names, and the printer invents \lp{W1}, \lp{W2}. The named
miniKanren answer is $(\lambda\,\alpha\,(\lambda\,\beta\,\alpha))$ for \emph{any} two symbols $\alpha\ne\beta$.
The disequality is needed to rule out the shadowed reading $\lambda x\lambda x\,x$, which would be
$\lambda\lambda1$. The names are logic variables with side conditions, and the programmer had to write the
\scm{≠°} that excludes capture. In λ-tree syntax there is no such condition to write.
\end{kanren}

\section{Computations over first-order formulas}\label{sec:lambdatree-fol-comp}

\subsection{Vacuous quantifiers and quantifier-free formulas}

Figure~7.6 (\book{195}):
\begin{lprolog}
type vacuous    form -> form -> o.
type quantfree  form -> o.
vacuous (all  x\ P) P.
vacuous (some x\ P) P.
quantfree tt & quantfree ff.
quantfree F :- atom F.
quantfree (neg F) :- quantfree F.
quantfree (F && G) & quantfree (F !! G) & quantfree (F ==> G) :-
   quantfree F, quantfree G.
\end{lprolog}
\lp{vacuous} \enquote{succeeds exactly when the top-level abstracted variable does not appear in the body of the
formula and, in this case, binds the second argument to the body} \book{195}. There is no occurrence test in
the code. Matching \lp{(all x\ P)} against \lp{(all x\ p a)} succeeds with \lp{P = p a}. Matching it against
\lp{(all x\ p x)} would need \lp{P = p x}, but \lp{P} is quantified outside \lp{x\}, so \lp{x} cannot appear in its
value and unification fails. The scoping constraint performs the occurrence test. \lp{quantfree} relies on
\lp{atom} (as in Figure~7.2), which is defined for each object signature.

\subsection{Negation normal form}\label{sec:lambdatree-nnf}

A formula is in \emph{negation normal form} (NNF) if it contains no implication and negation is applied only to
atoms. It is in \emph{prenex normal form} if no quantifier is in the scope of a connective. In classical
first-order logic every formula is equivalent to one in NNF, to one in prenex form, and to one in both
\book{195}. The NNF transformation orients the equivalences
\[
\begin{aligned}
\neg\neg B&\equiv B, & B_1\imp B_2&\equiv\neg B_1\vee B_2, & \neg\forall x\,B\,x&\equiv\exists x\,\neg B\,x,\\
\neg(B_1\wedge B_2)&\equiv\neg B_1\vee\neg B_2, & \neg(B_1\vee B_2)&\equiv\neg B_1\wedge\neg B_2, &
\neg\exists x\,B\,x&\equiv\forall x\,\neg B\,x
\end{aligned}
\]
from left to right (Figure~7.7, \book{196}):
\begin{lprolog}
type nnf  form -> form -> o.
nnf A A & nnf (neg A) (neg A) :- atom A.
nnf (neg (neg B)) D :- nnf B D.
nnf (neg (B && C)) (D !! E) &
nnf (neg (B !! C)) (D && E) :- nnf (neg B) D, nnf (neg C) E.
nnf (B ==> C) (D !! E) :- nnf (neg B) D, nnf C E.
nnf (B !! C) (D !! E) &
nnf (B && C) (D && E) :- nnf B D, nnf C E.
nnf (neg (all B)) (some D) &
nnf (neg (some B)) (all D) :- pi x\ nnf (neg (B x)) (D x).
nnf (all B) (all D) &
nnf (some B) (some D) :- pi x\ nnf (B x) (D x).
\end{lprolog}
The quantifier clauses are the idiom without \lp{=>}. Under \lp{pi x\} the eigenvariable \lp{x} is an object
\emph{term}, and the only predicate that meets the atoms containing it is \lp{atom}, which does not look inside
terms. The output \lp{(D x)} is again a pattern problem. Its solution abstracts \lp{x} out of the transformed
body, rebuilding the quantifier around it. For example (our trace, assuming \lp{atom} holds of \lp{p} atoms),
\lp{?- nnf (neg (all x\ p x)) D.} skips the two \lp{atom}-guarded clauses, since neither the whole formula nor
\lp{all x\ p x} is atomic. It reaches the \lp{neg (all B)} clause with \lp{B = x\ p x} and \lp{D = some D1}. Under
an eigenvariable \lp{c} the goal is \lp{nnf (neg (p c)) (D1 c)}. The clause \lp{nnf (neg A) (neg A) :- atom A}
gives \lp{(D1 c) = neg (p c)}, so \lp{D1 := x\ neg (p x)} and \lp{D = some x\ neg (p x)}, the formula
$\exists x\,\neg p(x)$.

The book claims that no two heads overlap, so \lp{nnf} is a partial function, and that \enquote{a simple argument
by induction shows also that the function represented is total} \book{196}. Both claims need qualification.

\begin{pitfall}[Non-overlap modulo a guard, and a missing case]
Syntactically, the head \lp{nnf A A} overlaps every other head, since \lp{A} is a variable. The clauses are
disjoint only \emph{modulo the guard} \lp{atom A}, which fails on compound formulas. A miniKanren port must keep
that guard decisive. The \lp{atom} relation must \emph{fail}, not enumerate or suspend, on compound or unbound
input, or \lp{nnf} produces spurious answers. Moreover, \lp{tt}, \lp{ff} and \lp{neg (B ==> C)} have no clause.
The last matches none of \lp{neg (neg _)}, \lp{neg (_ && _)}, \lp{neg (_ !! _)}, \lp{neg (all _)},
\lp{neg (some _)}. Totality holds only for the formulas without these cases. The fix when porting is one more
clause, \lp{nnf (neg (B ==> C)) (D && E) :- nnf B D, nnf (neg C) E.}
\end{pitfall}

\subsection{Prenex normal form}\label{sec:lambdatree-prenex}

For formulas in NNF, the prenex transformation uses the equivalences \book{196}
\[
(\forall x\,B_1x)\wedge(\forall x\,B_2x)\equiv\forall x\,(B_1x\wedge B_2x),\qquad
(\exists x\,B_1x)\vee(\exists x\,B_2x)\equiv\exists x\,(B_1x\vee B_2x),
\]
and, for $Q\in\{\forall,\exists\}$ and $\circ\in\{\wedge,\vee\}$,
$B_1\circ(Q x\,B_2x)\equiv Qx\,(B_1\circ B_2x)$ and $(Qx\,B_2x)\circ B_1\equiv Qx\,(B_2x\circ B_1)$. The usual
side condition, $x$ not free in $B_1$, is not written anywhere in the program, because it is automatic: in the
clauses below \lp{B} and \lp{C} are clause variables quantified outside the \lp{pi x\}. (The restriction to NNF
input \enquote{is not necessary and is used only to shorten the specification}, \book{196}.) Figure~7.8
(\book{197}):
\begin{lprolog}
type prenex, merge  form -> form -> o.
prenex A A & prenex (neg A) (neg A) :- atom A.
prenex (B && C) D :- prenex B U, prenex C V,
                     merge (U && V) D.
prenex (B !! C) D :- prenex B U, prenex C V,
                     merge (U !! V) D.
prenex (all B) (all D) &
prenex (some B) (some D) :- pi x\ prenex (B x) (D x).
merge ((all B) && (all C)) (all D) :-
   pi x\ merge ((B x) && (C x)) (D x).
merge ((all B) && C) (all D) &
merge ((some B) && C) (some D) :-
   pi x\ merge ((B x) && C) (D x).
merge (B && (all C)) (all D) &
merge (B && (some C)) (some D) :-
   pi x\ merge (B && (C x)) (D x).
merge ((some B) !! (some C)) (some D) :-
   pi x\ merge ((B x) !! (C x)) (D x).
merge ((some B) !! C) (some D) &
merge ((all B) !! C) (all D) :-
   pi x\ merge ((B x) !! C) (D x).
merge (B !! (some C)) (some D) &
merge (B !! (all C)) (all D) :-
   pi x\ merge (B !! (C x)) (D x).
merge B B :- quantfree B.
\end{lprolog}
Look at the clause \lp{merge ((all B) && C) (all D) :- pi x\ merge ((B x) && C) (D x).} The formula \lp{C}
moves under the new quantifier without any renaming, because \lp{C} is a closed value with respect to \lp{x}.
Capture is impossible by construction.

Some formulas have a unique prenex form \book{196}. For example
\begin{lprolog}
(all x\ (p x) && (all y\ q x y) && (p (f x))) !! (p a).
\end{lprolog}
has only
\begin{lprolog}
all x\ all y\ ((p x) && (q x y) && (p (f x))) !! (p a).
\end{lprolog}
In general, though, \lp{prenex} is a relation, not a function \book{196}--\book{197}:
\begin{lprolog}
?- prenex ((all x\ q x x) && (all z\ all y\ q z y)) P.
\end{lprolog}
has five answers \book{197}:
\begin{lprolog}
all z\ all y\ (q z z) && (q z y)
all x\ all z\ all y\ (q x x) && (q z y)
all z\ all x\ (q x x) && (q z x)
all z\ all x\ all y\ (q x x) && (q z y)
all z\ all y\ all x\ (q x x) && (q z y)
\end{lprolog}

\begin{tutorial}[Where the five prenex forms come from]
Both conjuncts are already prenex, so the goal reduces to
\lp{merge ((all x\ q x x) && (all z\ all y\ q z y)) P}. The following explanation is our own hand trace
against the clause order of Figure~7.8. It reproduces the book's order under depth-first search.
\begin{enumerate}[nosep]
  \item The $\forall\wedge\forall$ clause shares one eigenvariable between \lp{x} and \lp{z}, leaving
        \lp{merge ((q c c) && (all y\ q c y))}, which pulls \lp{y} out: answer 1.
  \item The \lp{((all B) && C)} clause pulls \lp{x} alone, then \lp{z} and \lp{y} come out of the right side one at
        a time: answer 2.
  \item The \lp{(B && (all C))} clause pulls \lp{z} first, leaving
        \lp{merge ((all x\ q x x) && (all y\ q c y))}. That subproblem again has three ways out: the
        $\forall\wedge\forall$ clause sharing \lp{x} with \lp{y} (answer 3), pulling \lp{x} then \lp{y} (answer
        4), or pulling \lp{y} then \lp{x} (answer 5).
\end{enumerate}
Every answer is a correct prenex form. No other clause applies at any step: the $\vee$ clauses need \lp{!!},
and \lp{merge B B :- quantfree B} fails while a quantifier remains. A miniKanren port with the same clauses
would have to produce exactly these five answers, but its interleaving search may list them in a different
order.
\end{tutorial}

\subsection{Syntactic classes and a richer interpreter}

Figure~7.9 (\book{198}) recognises the goal and clause formulas of \fohc{} and \fohh{} (\cref{ch:horn,ch:hh}):
\begin{lprolog}
type fohcG, fohcD, fohhG, fohhD  form -> o.
fohcG tt.
fohcG A :- atom A.
fohcG (some B) :- pi x\ fohcG (B x).
fohcG (B && C) & fohcG (B !! C) :- fohcG B, fohcG C.
fohcD A :- atom A.
fohcD (G ==> D) :- fohcG G, fohcD D.
fohcD (D1 && D2) :- fohcD D1, fohcD D2.
fohcD (all D) :- pi x\ fohcD (D x).
fohhG tt.
fohhG A :- atom A.
fohhG (D ==> G) :- fohhD D, fohhG G.
fohhG (B && C) & fohhG (B !! C) :- fohhG B, fohhG C.
fohhG (some B) & fohhG (all B) :- pi x\ fohhG (B x).
fohhD A :- atom A.
fohhD (D1 && D2) :- fohhD D1, fohhD D2.
fohhD (G ==> D) :- fohhG G, fohhD D.
fohhD (all D) :- pi x\ fohhD (D x).
\end{lprolog}
Here the book remarks on the \lp{pi x\} without \lp{=>}, a deviation from the idiom \book{198}. The
eigenvariables are object terms and \lp{atom} does not inspect terms, so no assumptions are needed. (The same
silent deviation occurs in \lp{red1}, \lp{nnf} and \lp{prenex}.) The book also notes a free substitution theorem:
\enquote{If $\forall x\,D$ is a first-order Horn clause and t is a term, then $D[t/x]$ is a first-order Horn
clause} \book{198}. It is the analogue of \cref{thm:lambdatree-bnorm-subst} and is proved the same way.

Two clauses extend the interpreter of Figure~7.1 from \fohc{} to \fohh{} \book{198}:
\begin{lprolog}
interp D (D1 ==> G) :- interp (D1 && D) G.
interp D (all G) :- pi x\ interp D (G x).
\end{lprolog}
Object $\wedge$, $\vee$, $\exists$ in goals are handled by meta $\wedge$, $\vee$, $\exists$, and the object
$\forall$ in goals by the meta \lp{pi}. Object implication in goals extends the object program. Since the object
program is an explicit argument, the extension is just \lp{(D1 && D)}, and no meta-level \lp{=>} is used.

\section{Specifying substitution without \texorpdfstring{$\beta$}{β}}\label{sec:lambdatree-copy}

Substitution into a term with an abstraction, by $\beta$, is one clause \book{199}:
\begin{lprolog}
type subst (tm -> tm) -> tm -> tm -> o.
subst R N (R N).
\end{lprolog}
As \cref{ch:lambda} showed, however, \enquote{full $\beta$-conversion is a complex operation on terms} \book{199}. The book therefore shows how to
specify substitution with only the weak $\beta_0$ of \cref{sec:lambdatree-llambda}. The tools are a
\emph{copy} predicate and a hypothetical assumption.

\subsection{Copy clauses}

Figure~7.10 (\book{200}):
\begin{lprolog}
type copy   tm -> tm -> o.
type subst  (tm -> tm) -> tm -> tm -> o.
copy (app M N) (app P Q) :- copy M P, copy N Q.
copy (abs M) (abs N) :- pi x\ copy x x => copy (M x) (N x).
subst M T S :- pi x\ copy x T => copy (M x) S.
\end{lprolog}
\lp{copy} specifies equality: for closed terms, \lp{copy t s} is derivable if and only if $t=s$. A trace
\book{199}:
\begin{lprolog}
?- copy (abs x\ abs y\ app y x) M.
   % M := (abs M1); eigenvariable c; assume (copy c c)
?- copy (abs y\ app y c) (M1 c).
   % (M1 c) := (abs M2); eigenvariable d; assume (copy d d)
?- copy (app d c) (M2 d).
   % three backchainings: app clause, then (copy d d), then (copy c c)
   % finally (M2 d) unified with (app d c)
M2 = u\app u c
M1 = v\abs u\app u v
M  = abs v\abs u\app u v.
\end{lprolog}
The equations \lp{(M1 c) = (abs M2)} and \lp{(M2 d) = (app d c)} are pattern problems, solved by abstracting
the eigenvariable out. So \lp{copy} \emph{rebuilds} its input binder by binder, and the abstractions of the output
are produced by unification.

\begin{tutorial}[Reading \texttt{subst} as a hypothetical copy]
\lp{subst M T S} copies the body of \lp{M} except that the outermost bound variable copies to \lp{T} \book{200}.
The clause does this in three steps:
\begin{enumerate}[nosep]
  \item the binder of \lp{M} becomes a formula-level \lp{pi x\};
  \item the assumption \lp{copy x T} says \enquote{wherever you meet \lp{x}, produce \lp{T}};
  \item \lp{copy (M x) S} copies the body.
\end{enumerate}
Inside the body, the assumptions \lp{copy c c} for inner binders preserve inner variables. The assumption
\lp{copy x T} replaces the substituted one. Since \lp{T} is closed with respect to every eigenvariable created
inside, no capture can happen.
\end{tutorial}

\begin{proposition}[Correctness of \texttt{subst}, \book{200}]\label{prop:lambdatree-subst}
If \lp{(subst M T S)} is derivable, then \lp{S} is equal to \lp{(M T)}.
\end{proposition}
\begin{proof}
\lp{pi x\ copy x T => copy (M x) S} is derivable. Since \lp{copy} is equality in the program without
assumptions, \lp{copy T T} is derivable. Instantiate \lp{x} with \lp{T}: the assumption becomes \lp{copy T T} and can
be discharged by modus ponens. So \lp{copy (M T) S} is derivable from the program without assumptions, hence
\lp{(M T)} $=$ \lp{S}.
\end{proof}

\subsection{Copy clauses for first-order formulas, and where they come from}

For the first-order formulas over the signature
\begin{lprolog}
type a  term.
type f  term -> term.
type g  term -> term -> term.
type p  term -> form.
type q  term -> term -> form.
\end{lprolog}
the copy clauses and substitution are Figure~7.11 (\book{201}):
\begin{lprolog}
kind term, form  type.
type copyterm  term -> term -> o.
type copyform  form -> form -> o.
type subst     (term -> form) -> term -> form -> o.

copyterm a a.
copyterm (f X) (f U) :- copyterm X U.
copyterm (g X Y) (g U V) :- copyterm X U, copyterm Y V.

copyform tt tt.
copyform ff ff.
copyform (neg B) (neg D) :- copyform B D.
copyform (B && C) (D && E) &
copyform (B !! C) (D !! E) &
copyform (B ==> C) (D ==> E) :- copyform B D, copyform C E.
copyform (all B) (all D) &
copyform (some B) (some D) :- pi y\ copyterm y y =>
                                     copyform (B y) (D y).
copyform (p X) (p U) :- copyterm X U.
copyform (q X Y) (q U V) :- copyterm X U, copyterm Y V.

subst M T N :- pi x\ copyterm x T => copyform (M x) N.
\end{lprolog}
These clauses are not ad hoc. They are generated from the signature by a translation that \enquote{lowers} the
equality at a higher type to relations at the base types \book{201}--\book{202}.

\begin{definition}[Type lowering, \book{201}--\book{202}]\label{def:lambdatree-lowering}
For terms $t,s$ of type $\tau$ built from the base types \lp{term} and \lp{form}, define the formulas
$[\![t,s:\tau]\!]^{+}$ (for goals) and $[\![t,s:\tau]\!]^{-}$ (for clauses) by
\[
\begin{aligned}
{}[\![t,s:\cst{term}]\!]^{\pm} &= \cst{copyterm}\ t\ s, \qquad
[\![t,s:\cst{form}]\!]^{\pm} = \cst{copyform}\ t\ s,\\
[\![t,s:\tau\to\sigma]\!]^{+} &= \forall x\,\big([\![x,x:\tau]\!]^{-}\imp[\![t\,x,\ s\,x:\sigma]\!]^{+}\big),\\
[\![t,s:\tau\to\sigma]\!]^{-} &= \forall x\,\forall y\,\big([\![x,y:\tau]\!]^{+}\imp[\![t\,x,\ s\,y:\sigma]\!]^{-}\big).
\end{aligned}
\]
\end{definition}
The positive case is the idiom itself: to compare two functions, apply both to a fresh $x$ and assume that $x$
is related to itself. The negative case, used for clauses, compares arguments pairwise. Polarity flips at each
arrow, just as an arrow-typed argument of a constructor turns a clause into a hypothetical goal.

\begin{example}[Lowering the quantifier \lp{all}, \book{201}--\book{202}]\label{ex:lambdatree-lower-all}
For $\cst{all}:(\cst{term}\to\cst{form})\to\cst{form}$:
\[
\begin{aligned}
&[\![\cst{all},\cst{all}:(\cst{term}\to\cst{form})\to\cst{form}]\!]^{-}\\
&\quad= \forall x\forall y\,\big([\![x,y:\cst{term}\to\cst{form}]\!]^{+}\imp[\![\cst{all}\,x,\cst{all}\,y:\cst{form}]\!]^{-}\big)\\
&\quad= \forall x\forall y\,\big((\forall z\,[\![z,z:\cst{term}]\!]^{-}\imp[\![x\,z,\ y\,z:\cst{form}]\!]^{+})\imp[\![\cst{all}\,x,\cst{all}\,y:\cst{form}]\!]^{-}\big)\\
&\quad= \forall x\forall y\,\big(\forall z\,(\cst{copyterm}\ z\ z\imp\cst{copyform}\ (x\,z)\ (y\,z))\\
&\qquad\qquad\imp\cst{copyform}\ (\cst{all}\,x)\ (\cst{all}\,y)\big).
\end{aligned}
\]
This is the \lp{all} clause of Figure~7.11. \enquote{In fact, the clauses in that figure are exactly the formulas
$[\![c, c : \tau]\!]^-$ for each of the constants c : $\tau$ of the object logic} \book{202}.
\end{example}

\begin{proposition}[\book{202}]
\lp{(copyterm T S)} is derivable if and only if \lp{T} and \lp{S} are equal terms, and \lp{(copyform B C)} is
derivable if and only if \lp{B} and \lp{C} are equal formulas. Hence, if \lp{(subst M T S)} is provable, then
\lp{S} is equal to \lp{(M T)}.
\end{proposition}

Two remarks from \book{202}. First, these programs are \emph{signature-dependent}: the object signature is
reflected into clauses, whereas the one-line $\beta$ version is not. Second, lowering is general. Each primitive
type maps to its own predicate, the relation need not be binary, and the \lp{term} predicate of
\cref{sec:lambdatree-term} is the arity-one lowering for the signature \lp{abs}, \lp{app}.

Substitution can also be defined by direct recursion on the abstraction, with no copy predicate (Figure~7.12,
\book{203}):
\begin{lprolog}
type subst      (term -> form) -> term -> form -> o.
type substterm  (term -> term) -> term -> term -> o.
subst (x\ tt) T tt.
subst (x\ ff) T ff.
subst (x\ neg (B x)) T (neg D) :- subst B T D.
subst (x\ (B x) && (C x))  T (D && E)  &
subst (x\ (B x) !! (C x))  T (D !! E)  &
subst (x\ (B x) ==> (C x)) T (D ==> E) :- subst B T D,
                                         subst C T E.
subst (x\ all  (B x)) T (all  D) &
subst (x\ some (B x)) T (some D) :-
   pi y\ substterm (x\y) T y => subst (x\ B x y) T (D y).
subst (x\ p (X x)) T (p U) :- substterm X T U.
subst (x\ q (X x) (Y x)) T (q U V) :- substterm X T U,
                                      substterm Y T V.
substterm (x\ x) T T.
substterm (x\ a) T a.
substterm (x\ f (F x)) T (f S) :- substterm F T S.
substterm (x\ g (F x) (G x)) T (g S R) :-
   substterm F T S,
   substterm G T R.
\end{lprolog}
The book says only that these clauses are signature dependent, and in its notes that they are loosely based
on Andrews' equational presentation of $\lambda$-conversion \book{210}. Three points are worth seeing. The patterns
\lp{(x\ a)} and \lp{(x\y)} are \emph{vacuous} abstractions: they match only bodies that do not depend on \lp{x}.
Under a nested quantifier the body becomes a two-place abstraction \lp{(x\ B x y)}. The assumption
\lp{substterm (x\y) T y} says that the inner bound \lp{y}, seen as a constant function of \lp{x}, substitutes to
itself.

\subsection{Substitution runs backwards}

Every version of \lp{subst} is a relation (\book{202}). Given \lp{T} and the result, it computes
the abstraction:
\begin{lprolog}
?- subst M (f a) ((p (f a)) && (q a (f a))).
M = x\ (p x) && (q a x)
M = x\ (p x) && (q a (f a))
M = x\ (p (f a)) && (q a x)
M = x\ (p (f a)) && (q a (f a))
\end{lprolog}
Each of the two occurrences of \lp{(f a)} may or may not be abstracted, so there are $2\times2$ answers. They
are the four \emph{existential generalisations} of $p(f(a))\wedge q(a,f(a))$ on the term $f(a)$ \book{203}:
$\exists x\,p(x)\wedge q(a,x)$, $\exists x\,p(x)\wedge q(a,f(a))$, $\exists x\,p(f(a))\wedge q(a,x)$,
$\exists x\,p(f(a))\wedge q(a,f(a))$. The \enquote{abstracted} answer comes first. This is consistent with the
convention of \cref{ch:hh} that clauses added by an implication are tried before the static ones, so
\lp{copyterm x (f a)} is tried before \lp{copyterm (f X) (f U)}.

\begin{pitfall}[Copy is equality logically, not operationally]
\enquote{the operational behavior of the copy predicate can be quite different from that of the equality notion
it encodes} \book{207}. \lp{?- M = N.} binds \lp{M} to \lp{N} at once. \lp{?- copy M N.} builds terms. With
\begin{lprolog}
kind i     type.
type a     i.
type f     i -> i.
type copyi i -> i -> o.
copyi a a.
copyi (f X) (f Y) :- copyi X Y.
\end{lprolog}
\lp{?- copyi M N.} enumerates \lp{a}, \lp{(f a)}, \lp{(f (f a))}, \ldots, and \enquote{if the order of the clauses
for copyi were reversed, then this query would cause the depth-first $\lambda$Prolog interpreter to loop without
providing any solution} \book{207}.
\end{pitfall}

\begin{kanren}[Clause order and interleaving]
The same predicate in \micro, in both clause orders:
\begin{scheme}
(define-relation (copyi° m n)
  (cond° ((=° m 'a) (=° n 'a))
         ((fresh° (x y) (=° m `(f ,x)) (=° n `(f ,y)) (copyi° x y)))))
(define-relation (copyi-rev° m n)
  (cond° ((fresh° (x y) (=° m `(f ,x)) (=° n `(f ,y)) (copyi-rev° x y)))
         ((=° m 'a) (=° n 'a))))

(print (μkanren-run (q 4 #t) (fresh° (m n) (=° q `(,m ,n)) (copyi° m n))))
(print (μkanren-run (q 4 #t) (fresh° (m n) (=° q `(,m ,n)) (copyi-rev° m n))))
\end{scheme}
\begin{answers}
((a a) ((f a) (f a)) ((f (f a)) (f (f a))) ((f (f (f a))) (f (f (f a)))))
((a a) ((f a) (f a)) ((f (f a)) (f (f a))) ((f (f (f a))) (f (f (f a)))))
\end{answers}
The reversed order loops in λProlog but produces the same answers here. \scm{define-relation} delays the
recursive call and \scm{cond°} interleaves its branches, so the base case is reached. Interleaving changes the
\emph{order} of answers (the five prenex forms come out in the order 1~3~2~4~5, \cref{sec:examples-prenex};
the four generalisations happen to keep the book's order, \cref{sec:examples-subst}) and turns some depth-first loops into
productive streams. It does not remove genuine divergence. \lp{cbv} on $(\lambda x\lambda w\,w)\,\Omega$ has no
answer under any strategy, and a conjunction such as \enquote{\scm{copyi°} and then fail} still runs forever,
because it has infinitely many candidates and no finite failure.
\end{kanren}

\section{The \texorpdfstring{$\lambda$}{λ}-tree approach to abstract syntax}\label{sec:lambdatree-hoas}

Concrete syntax (strings, whitespace, parentheses, infix operators, comments) is for humans. Parsing produces
\emph{abstract syntax}, a tree \book{203}. For binders there are several choices \book{204}: keep explicit names;
or map object binding to the meta-language's binding, which is called \emph{higher-order abstract syntax}
(HOAS) \cite{PfenningElliott1988}. The term covers \enquote{a range of encoding techniques}, and the book
distinguishes two of them sharply.
\begin{itemize}[nosep]
  \item In functional programming, and in constructive type theories that compute by evaluation, meta-level
        abstraction builds \emph{rich function spaces}. \enquote{Because of the richness of these spaces, it is not
        possible to compare function-valued expressions for equality in any simple way}. If quantified formulas
        are encoded with such functions, \enquote{we cannot also compare or structurally analyze formulas under
        such an encoding} \book{204}.
  \item Most logic programming languages have no meta-level binders in terms at all.
  \item The $\lambda$-terms of hohh use \enquote{a $\lambda$-calculus with a weak notion of functionality that makes
        it possible to analyze the structures of objects under the encoding} \book{204}.
\end{itemize}
The book contrasts \enquote{the two extremes in the treatment of binding---one using functions that are best
thought of extensionally and the other using $\lambda$-expressions that are easily treated intensionally}.

\begin{definition}[$\lambda$-tree syntax, \book{204}]\label{def:lambdatree-lambdatree}
\enquote{We shall use the term $\lambda$-tree syntax to denote the approach to higher-order abstract syntax that uses
typed $\lambda$-terms modulo equality based on $\alpha$-conversion, some sufficiently weak form of
$\beta$-conversion, and possibly $\eta$-conversion.}
\end{definition}
In this chapter $\beta$ is \enquote{weakened} by simple typing, which makes it terminating. The next section
weakens it much further, to $\beta_0$.

The book's bibliographic notes (\book{208}--\book{210}) trace the idea back a long way. Encoding a binder as a
constant applied to a $\lambda$-abstraction goes back at least to Church's simple theory of types (1940), and
λProlog's \lp{pi} and \lp{sigma} are named after Church's $\Pi$ and $\Sigma$. Computing with such terms
starts with Huet and Lang (1978) and Miller and Nadathur (1987). The phrase \enquote{higher-order abstract
syntax} comes from Pfenning and Elliott \cite{PfenningElliott1988}, and the name \enquote{$\lambda$-tree
syntax} from Miller and Palamidessi (1999). The copy clauses of \cref{sec:lambdatree-copy} go back to Miller
(1991) at second order and to Felty for arbitrary orders. The $\Llam$ fragment of the next section is
Miller's \cite{Miller1991}.

For a miniKanren programmer, the options can be compared directly. The first column is the
\enquote{functional} reading of HOAS that we ran in \cref{sec:lambdatree-beta}. The second is nominal logic
programming as in $\alpha$Kanren \cite{ByrdFriedman2007,Byrd2009}, based on nominal unification
\cite{UrbanPittsGabbay2004}. The third is $\lambda$-tree syntax.
\begin{center}
\small
\begin{tabular}{@{}>{\raggedright\arraybackslash}p{0.27\linewidth}>{\raggedright\arraybackslash}p{0.2\linewidth}>{\raggedright\arraybackslash}p{0.2\linewidth}>{\raggedright\arraybackslash}p{0.23\linewidth}@{}}
\toprule
 & host closures & names + freshness ($\alpha$Kanren) & $\lambda$-tree syntax \\
\midrule
$\alpha$-equivalence & no (\scm{eq?} only) & yes, via swapping & yes, built in \\
substitution & host application & a user relation with freshness side conditions & meta-level $\beta$ (or $\beta_0$ + copy) \\
analysing a body & only by applying to a name & by unification under \texttt{tie} & by mobility (\texttt{pi}, \texttt{=>}) \\
building an abstraction by unification & impossible & yes & yes (pattern unification) \\
\enquote{$x$ not free in $M$} & not expressible & freshness constraint $a\,\#\,M$ & scoping of variables \\
\bottomrule
\end{tabular}
\end{center}
The same notes (\book{209}--\book{210}) name nominal logic as the alternative approach: it is
first-order, with name-swapping and freshness, and binding is derived from these primitives. Porting the book's
programs to $\alpha$Kanren would mean rewriting them, not transliterating them: there is no $\beta$, mobility is replaced by fresh names, and hypothetical
contexts must be threaded explicitly.

\section{The \texorpdfstring{$L_\lambda$}{Lλ} subset}\label{sec:lambdatree-llambda}

Predicate quantification is useful for manipulating $\lambda$-tree syntax, but \enquote{it is certainly not
central to it} \book{205}. Can $\beta$ be weakened further? The book recalls the chapter's two principles: bound
names are a fiction ($\alpha$), and binders are mobile, from term level to formula level to proof level. It then
identifies the variables that carry mobility.

\begin{definition}[Essentially universal and essentially existential occurrences, \book{205}]\label{def:lambdatree-essential}
In a goal formula $G$, a bound variable occurrence is \emph{essentially universal} if it is bound by a positive
$\forall$, a negative $\exists$, or a term-level $\lambda$. It is \emph{essentially existential} otherwise, that is,
if it is bound by a negative $\forall$ or a positive $\exists$. In a program clause the classification is dual: an
occurrence is essentially universal if it is bound by a negative $\forall$, a positive $\exists$, or a term-level
$\lambda$, and essentially existential if it is bound by a positive $\forall$ or a negative $\exists$.
\end{definition}
During proof search, essentially existential variables become logic variables, instantiated by unification
with general terms. Essentially universal ones are instantiated only with eigenvariables. \enquote{It is the
essentially universal bound variable occurrences that provide mobility during proof search} \book{205}. In the
clause \lp{typeof (abs M) (arr A B) :- pi x\ typeof x A => typeof (M x) B.}, the clause variables \lp{M},
\lp{A}, \lp{B} are essentially existential, while the \lp{pi}-bound \lp{x} sits in the negative body of the clause
and is essentially universal.

\begin{definition}[$\Llam$, \book{205}--\book{206}]\label{def:lambdatree-llambda}
$\Llam$ restricts hohh by requiring (1) that all quantification be over non-predicate types, and (2) that
\enquote{Every subterm in B with an essentially existential variable occurrence as its head, that is, a subterm of
the form ($x\,t_1\ldots t_n$) ($n\ge0$), where x is essentially existentially quantified in B, must be such that
$t_1,\ldots,t_n$ are distinct essentially universal variable occurrences bound in the scope of x.}
\end{definition}
Such an application $x\,y_1\ldots y_n$ is a \emph{higher-order pattern}. If $x$ is ever instantiated with a
term $t$, the only redexes created have the form $(t\,y_1\ldots y_n)$ with the $y_i$ not free in $t$. Up to
$\alpha$ and $\eta$, $t=\lambda y_1\ldots\lambda y_n\,t'$, and the reduction just strips the binders
\book{206}. Only this sub-case of $\beta$ is needed.

\begin{definition}[$\beta_0$, \book{206}]\label{def:lambdatree-beta0}
\[
(\abs{x}{s})\,x = s \qquad(\beta_0)
\]
In the presence of $\alpha$-conversion this is equivalent to $(\abs{y}{s})\,x=\subst{s}{x}{y}$, provided $x$ is not
free in $\abs{y}{s}$.
\end{definition}

\begin{example}[Goal versus clause, \book{206}]
With $p:\typ{i}\to\typ{o}$ and $f:\typ{i}\to\typ{i}$, the formula
$\forall_{\typ{i}\to\typ{i}}x\,\forall_{\typ{i}}y\,(p\,(x\,y)\imp p\,(f\,y))$ is an $\Llam$ \emph{goal}: both
quantifiers are positive, so $x$ and $y$ are essentially universal, and nothing existential is applied. As a
\emph{program clause} it is not in $\Llam$: there both $x$ and $y$ are essentially existential, and $(x\,y)$ applies
an existential variable to an existential one. Membership in $\Llam$ depends on polarity.
\end{example}

The book states which programs of the chapter belong to $\Llam$ \book{206}: \enquote{all the program clauses in
Figures 7.7 through 7.12 belong to $L_\lambda$}, and exactly one clause of Figure~7.4 does not, namely
\lp{redex (app (abs R) N) (R N).}, where \lp{R} is applied to the existential \lp{N}. By the same definition, more
programs of the chapter are outside $\Llam$, although the book singles out only \lp{redex}. The table summarises the
situation; the classification of the rows other than \lp{redex} and Figures~7.7--7.12 is ours.
\begin{center}
\small
\begin{tabular}{@{}p{0.47\linewidth}p{0.45\linewidth}@{}}
\toprule
in $\Llam$ (applications are \texttt{(M x)} with \texttt{x} \texttt{pi}-bound) & outside $\Llam$ \\
\midrule
\texttt{term}, \texttt{bnorm}, \texttt{red1}, \texttt{reduce}, \texttt{path}, \texttt{bpath}, \texttt{addbeta},
\texttt{typeof}, \texttt{trans}, \texttt{vacuous}, \texttt{nnf}, \texttt{prenex}, \texttt{fohcG} \ldots,
\texttt{copy}, every \texttt{subst} of Figures 7.10--7.12
& \texttt{redex} (\texttt{(R N)}); \texttt{list\_instan} (\texttt{(B T)}); \texttt{interp}/\texttt{backchain}
(\texttt{(G X)}, \texttt{(D1 X)}); \texttt{cbn}/\texttt{cbv} (\texttt{(R N)}, \texttt{(R U)}); the one-line
\texttt{subst R N (R N)}; \texttt{repeat} and \texttt{foreach} (they quantify over predicates) \\
\bottomrule
\end{tabular}
\end{center}
In every non-$\Llam$ clause of the chapter, the abstraction is applied to an argument of \emph{base} type
(\lp{term} or \lp{tm}). Instantiating \lp{R : tm -> tm} with \lp{N : tm} therefore never creates a new redex. It is
plain \enquote{instantiate the bound variable}, provided \lp{R} is known when the application is formed.

A non-$\Llam$ program can often be rewritten into $\Llam$ \book{206}. With the copy clauses of Figure~7.10, the
offending clause becomes
\begin{lprolog}
redex (app (abs R) N) S :- subst R N S.
\end{lprolog}
\enquote{If we do this, we obtain a specification that is entirely within $L_\lambda$.} The translation has two
costs \book{207}. The signature must be reflected into copy clauses with $[\![\cdot]\!]^{\pm}$. And copy behaves
operationally unlike equality (the \lp{copyi} pitfall above).

Why care? Compare the two rules \book{207}:
\[
(\abs{x}{M})\,N=\subst{M}{N}{x}\quad(\beta)\qquad\qquad(\abs{x}{M})\,x=M\quad(\beta_0).
\]
A $\beta_0$ step removes one abstraction--application pair, and the result is strictly smaller and has no new
redexes. A $\beta$ step may shrink the term (if $x$ does not occur in $M$), may blow it up (if $x$ occurs many times
and $N$ is large), and may create new redexes (if $x$ occurs applied in $M$). The difference carries over to
unification, the subject of \cref{ch:hounif} \book{208}:
\begin{itemize}[nosep]
  \item unification modulo $\beta\eta$ is undecidable and may have several most general unifiers;
  \item $\Llam$ unification is \enquote{much more like first-order unification}: it is decidable, and when a unifier
        exists there is a single most general one \cite{Miller1991,Nipkow1993};
  \item \enquote{an overwhelming majority of clauses in \textquote{typical} $\lambda$Prolog programs are within the
        $L_\lambda$ fragment}.
\end{itemize}

\begin{pitfall}[$\beta_0$ is not the same as $\Llam$]
All $\Llam$ unification problems are problems modulo $\alpha$, $\beta_0$, $\eta$, but not conversely \book{208}.
Finding $F$ with $\forall x\,(F\,x\,x=x)$ needs only $\beta_0$ to justify the solutions $F:=\lambda u\lambda v\,u$
and $F:=\lambda u\lambda v\,v$. It is not an $\Llam$ problem, because $F$ is applied to the same universal variable
twice, and it has two incomparable most general unifiers. A unifier for a miniKanren must decide what to do with
such non-pattern equations: reject them, or postpone them until they become patterns.
\end{pitfall}

Finally, raising (\cref{ch:lambda}), the transformation of $\forall y\,\exists x\,G$ into
$\exists h\,\forall y\,\subst{G}{(h\,y)}{x}$ that implementations use to simplify quantifier alternation,
\enquote{preserves the class of $L_\lambda$ formulas} \book{208}: $(h\,y)$ is a pattern.

\begin{kanren}[A requirements checklist for \micro]
Gathering the boxes of this chapter, a miniKanren that runs the book's programs needs the following. Each item is
justified by a program above. The design that provides them is in \cref{ch:design,ch:module}.
\begin{center}
\small
\begin{tabular}{@{}p{0.34\linewidth}p{0.58\linewidth}@{}}
\toprule
requirement & needed by \\
\midrule
abstractions as data, equal up to $\alpha$ (and $\eta$) & every program; \texttt{(all p)}; answers like \texttt{abs W1\textbackslash\ W1} \\
application of a known abstraction to a term & \texttt{list\_instan}, \texttt{interp}, \texttt{cbn}, \texttt{cbv}, \texttt{redex} \\
eigenvariables with the scoping constraint & every \texttt{pi}; correctness of \texttt{vacuous}, $\eta$-\texttt{redex}, \texttt{prenex} \\
scoped assumed facts, possibly with logic variables & \texttt{term}, \texttt{bnorm}, \texttt{typeof}, \texttt{trans}, \texttt{copy} \\
scoped assumed rules with bodies and own $\forall$ variables & \texttt{bpath}, \texttt{addbeta}, the \texttt{trans} variant \\
pattern unification for output abstractions \texttt{(D x)} & \texttt{red1}, \texttt{path} synthesis, \texttt{nnf}, \texttt{prenex}, \texttt{copy}, \texttt{subst} backwards \\
occurs check, also through abstractions & \texttt{typeof (abs x\textbackslash\ app x x)} \\
printing binders with invented names & every answer containing an abstraction \\
\bottomrule
\end{tabular}
\end{center}
Some of this exists in the miniKanren literature. The miniKanren with symbolic constraints has an
\enquote{eigen} operator with a scope check \cite{HemannByrdFriedman2013}. $\lambda$Kanren adds \texttt{tie},
application, \texttt{all} and assumed relations to $\mu$Kanren \cite{MaLuFriedman2020}, and its successor
discusses the $\alpha$-equivalence surprises of combining higher-order unification with miniKanren
\cite{MaFriedman2021} (see \cref{ch:lkanren}). Today's \micro{} has none of these features. It does have
first-order unification with occurs check and records, which can represent de Bruijn terms, and an interleaving
search, which already changes some of the book's loops into productive streams.
\end{kanren}

\section{Summary}\label{sec:lambdatree-summary}
\begin{itemize}
  \item An object-level binder is a constant applied to a meta-level abstraction: \lp{all (x\ B)},
        \lp{abs (x\ t)}. Equality of the meta-language includes $\alpha$ (bound names are a fiction), $\beta$ and
        $\eta$ (so \lp{(all p)} is a quantifier too).
  \item Object-level substitution is meta-level $\beta$: \lp{(B T)}. Interpreters (\lp{interp}) and evaluators
        (\lp{cbn}, \lp{cbv}) need no renaming code at all.
  \item An abstraction cannot be decomposed into a name and a body. Recursion under a binder uses
        \emph{mobility}: term-level $\lambda$, then formula-level \lp{pi}, then proof-level eigenvariable, with
        \lp{=>} assuming the clauses that describe the new eigenvariable (\cref{def:lambdatree-idiom}).
  \item The eigenvariable's scoping constraint is part of the meaning of \lp{pi}. It silently implements every
        \enquote{$x$ not free in $M$} side condition (\lp{vacuous}, $\eta$, prenex). A fresh constant without it is
        unsound.
  \item Assumptions are scoped, may contain logic variables (\lp{typeof x A}), and may be rules with bodies and
        their own quantifiers (\lp{bpath}, \lp{addbeta}, \lp{trans}). Some predicates have no static clause for
        variables at all.
  \item Metatheory comes cheap: the substitution lemma for \lp{bnorm}, subject reduction for \lp{typeof}, and the
        correctness of \lp{subst} all follow from instantiating a \lp{pi} and discharging a \lp{=>}.
  \item Programs that \emph{build} abstractions (\lp{red1}, \lp{nnf}, \lp{prenex}, \lp{copy}, path synthesis,
        \lp{trans} backwards) rely on unification solving \lp{(D x) = t} by abstracting the eigenvariable out:
        pattern unification.
  \item Copy clauses, generated from the signature by type lowering $[\![\cdot]\!]^{\pm}$, specify substitution
        using only $\beta_0$. They are logically equality but operationally term generators.
  \item The relations are not functions: \lp{prenex} has five answers on the book's example, \lp{subst} run
        backwards enumerates the four existential generalisations.
  \item $\lambda$-tree syntax is HOAS with $\lambda$-terms modulo $\alpha$, a weak $\beta$ and possibly $\eta$, which can
        be analysed intensionally, unlike host-language closures. $\Llam$ restricts existential variables to
        patterns, so that only $\beta_0$ is needed and unification is decidable and unitary (\cref{ch:hounif}).
  \item On today's \micro, abstractions as closures work only as inputs, and explicit contexts with de Bruijn
        indices replace \lp{pi}/\lp{=>} at the cost of threading the context and writing capture conditions by
        hand. The requirements checklist above is the starting point of the proposal (\cref{ch:design}).
\end{itemize}

\chapter{Unification of \texorpdfstring{$\lambda$}{λ}-terms and higher-order patterns}\label{ch:hounif}

% ============================================================================================
\section{What this chapter is about}\label{sec:hounif-intro}

The previous chapters built programs whose execution keeps asking one question: \emph{find values for the
logic variables that make two $\lambda$-terms equal modulo $\alpha$, $\beta$ and $\eta$}. \Cref{ch:lambda} showed
how such questions are written as quantified formulas; \cref{ch:hoprog,ch:lambdatree} produced them by the
thousand, from higher-order predicates such as \lp{mappred} and from the \lp{pi x\ ...} goals that walk under
binders in $\lambda$-tree syntax. \Cref{ch:lambdatree} also isolated the fragment $\Llam$, in which a logic
variable is applied only to distinct bound variables. This chapter, which follows Chapter~8 of the
book~\cite{MillerNadathur2012}, explains how an implementation \emph{solves} these questions, and why $\Llam$ is
the fragment to aim for.

The story has two halves.
\begin{enumerate}[nosep]
  \item \textbf{General higher-order unification} is hard: there may be no most general unifier, complete sets
    of unifiers may be infinite, and unifiability is undecidable (\cref{sec:hounif-hard}). Still, Huet's
    procedure~\cite{Huet1975} organises the search (\cref{sec:hounif-huet}): it simplifies rigid-rigid
    equations, \emph{guesses} imitation or projection substitutions for flexible-rigid ones, and stops when only
    flexible-flexible equations are left. Its search space is a \emph{matching tree}, which can be infinite.
  \item \textbf{Higher-order pattern unification} ($\Llam$-unification, \cref{sec:hounif-llambda,sec:hounif-algo})
    is the same procedure restricted to patterns. There it becomes \emph{deterministic}, has \emph{most general
    unifiers} and, with an occurs-check, \emph{terminates}~\cite{Miller1991}. Its final form, \emph{variable
    elimination}, is \enquote{computationally similar to first-order unification} \book{224}: one walk over the
    rigid side that inverts it, prunes, and checks occurrences and scope.
\end{enumerate}
We then do the algorithm by hand on the book's examples (\cref{sec:hounif-traces}), rebuild Nipkow's de~Bruijn
implementation~\cite{Nipkow1993} and run a Scheme transliteration of it (\cref{sec:hounif-nipkow}), and finish
with what real systems do: universe levels instead of raising, and delayed non-pattern equations
(\cref{sec:hounif-pragmatics}). If you only need the implementer's digest, read
\cref{def:hounif-elim,thm:hounif-mgu} and \cref{sec:hounif-nipkow,sec:hounif-pragmatics}.

\paragraph{Conventions.} As in the book's Chapter~8, uppercase $F,G,H,U,V,W,X,Y$ are \emph{existentially}
quantified variables (logic variables), and lowercase $a,b,c,d,f,g,u,w,x,y,z$ are \emph{universally}
quantified ones (signature constants or eigenvariables). $\bar x$ abbreviates $x_1\dots x_n$, and $(F\,\bar x)$
is $(F\,x_1\dots x_n)$.

% ============================================================================================
\section{Unification problems and their normal form}\label{sec:hounif-problems}

\begin{definition}[Unification problem {\book{211}}]\label{def:hounif-problem}
A \emph{unification problem} is a formula
\[
  Q_1x_1\dots Q_nx_n\,[\,t_1=s_1\wedge\dots\wedge t_m=s_m\,]
\]
where $n,m\ge 0$, each $Q_i$ is $\forall$ or $\exists$, and $t_j$, $s_j$ are $\lambda$-terms of the same type. The
quantifier sequence is the \emph{prefix}, the conjunction is the \emph{matrix}. A \emph{unifier} is a substitution
for the existentially quantified variables that makes every $t_j$ $\lambda$-convertible to $s_j$.
\end{definition}

Reading the problem as a formula is what makes the scoping right. In $\exists X\,\forall y\,[X=y]$ no unifier
exists, because $X$ is chosen before $y$ and cannot mention it; in $\forall y\,\exists X\,[X=y]$ the unifier is
$X:=y$. A miniKanren with only first-order terms never sees the difference because it has no universal
quantifier in goals; a $\lambda$Prolog with \lp{pi} sees it all the time.

\Cref{ch:lambda} distinguished a unifier from a \emph{solution}, which in addition uses only closed terms and
therefore \enquote{embodies an actual proof of the quantified formula}; following common practice the book's
Chapter~8 limits itself to unifiers \book{211}. (The difference matters with empty types; see the last pitfall of
\cref{sec:hounif-dynamic}.)

\subsection{The normal form assumed by the algorithms}

Before solving, the book puts every problem in a fixed shape \book{212}. The four steps are worth knowing one by
one, because an implementation performs each of them lazily.
\begin{enumerate}
  \item \textbf{A $\forall\exists\forall$ prefix.} Using \emph{raising} \book{107--108}, the formulas
    $\forall y\,\exists x\,G$ and $\exists h\,\forall y\,G[(h\,y)/x]$ are equiprovable when $h:\sigma\to\tau$ for
    $y:\sigma$, $x:\tau$. Moving every $\exists$ to the left of every $\forall$ in this way yields a prefix
    $\forall\bar a\;\exists\bar X\;\forall\bar y$. The outer $\forall\bar a$ represent \enquote{the ambient signature
    comprising the declared constants in a $\lambda$Prolog program, and we will interpret them as such}
    \book{212}. The price: \enquote{the type of the variable that is existentially quantified needs to be
    \enquote{raised} by the type of the universally quantified variable} \book{108}.
  \item \textbf{Substitution as $\beta$-redex.} For $\theta=\{\langle x_1,t_1\rangle,\dots,\langle x_n,t_n\rangle\}$,
    the book defines \book{212}
    \[\theta(t):=((\lambda x_1\dots\lambda x_n\,t)\,t_1\dots t_n).\]
    Thus $\theta(t)$ and
    $\theta(t')$ are equal whenever $t$ and $t'$ are $\lambda$-convertible, and we may assume all terms are in
    $\beta$-normal form.
  \item \textbf{Equal binders.} By $\eta$-conversion the two sides of every equation get binders of the same
    length.
  \item \textbf{Binders become universals.} Using $(\lambda x\,t=\lambda x\,s)\equiv\forall x\,(t=s)$ and moving
    such $\forall x$ out of the conjunction, the binders at the heads of the equations turn into universals that
    \emph{extend the inner block}. After this step \enquote{the equations in the matrix are between terms that
    have empty binders. We shall assume that all unification problems are of this kind} \book{212}.
\end{enumerate}

\begin{example}[Normalising a flexible-flexible problem]\label{ex:hounif-ch4}
\Cref{ch:lambda} met the problem \book{113}
\[\exists G\,\exists H\,[\lambda x\lambda y\,(G\,x)=\lambda x\lambda y\,(H\,y)].\]
Step~4 turns it into $\exists G\,\exists H\,\forall x\,\forall y\,[(G\,x)=(H\,y)]$. Now $x$ and $y$ are
universals quantified \emph{inside} the scope of $G$ and $H$. Any common instance must be a term in which
neither $x$ nor $y$ occurs, and the book announces that the \enquote{most general} unifier substitutes
$\lambda w\,T$ for both $G$ and $H$, with $T$ new \book{113}. We derive it in \cref{ex:hounif-ffdiff}.
\end{example}

\begin{remark}
Step~4 is exactly what an implementation does when it meets $\lambda x\,s=\lambda x\,t$: open both binders
with one fresh eigenvariable and continue. Because the new $\forall$ goes at the \emph{end} of the prefix, no
existing logic variable may depend on it, which is what a \enquote{newest level} eigenvariable gives you
(\cref{sec:hounif-pragmatics}).
\end{remark}

\begin{definition}[Rigid and flexible {\book{214}, \book{113}}]\label{def:hounif-rigid}
In a normalised problem, a term on either side of an equation is \emph{rigid} if its head is bound by a
universal quantifier of the prefix, and \emph{flexible} if its head is bound by an existential one. Equations
are accordingly rigid-rigid, rigid-flexible, flexible-rigid or flexible-flexible; a problem is
flexible-flexible if all its equations are.
\end{definition}

The key fact about a rigid term is that \enquote{its head is unaffected by any substitutions that might be made
for existentially quantified variables} \book{214}. A flexible head, on the other hand, can be replaced by any
term, including one that throws away or reorders the head's arguments. This asymmetry drives every rule
below.

\begin{kanren}[first-order unification is the degenerate case]
In \micro{} every Scheme atom, pair, vector or record constructor plays the role of an outer
universal $\forall\bar a$ (a rigid head), and every logic variable is an existential. There is no inner
$\forall$ block and no logic variable is ever applied to arguments, so every flexible term is a bare $X$, which
is a pattern with $n=0$. The unifier \scm{μkanren-state-unify} is precisely the rigid-rigid decomposition
plus the $n=0$ case of variable elimination (bind $X$ after an occurs-check). Everything in this chapter is an
extension of that code path, not a replacement. What a higher-order extension needs is: inner universals
(eigenvariables) as a new kind of rigid head, binders, applications of logic variables, and a scope
discipline; see \cref{ch:design}.
\end{kanren}

% ============================================================================================
\section{Why general higher-order unification is hard}\label{sec:hounif-hard}

Recall from first-order unification (\cref{ch:horn}) that a \emph{most general unifier} (mgu) is a unifier from
which every other unifier is obtained \enquote{by making further substitutions} \book{212}. It satisfies all the
constraints while leaving every other way of satisfying them open, which is what makes Prolog's and
miniKanren's \enquote{one substitution per answer} work. At higher order this breaks down in three ways.

\begin{example}[No mgu {\book{213}}]\label{ex:hounif-nomgu}
Let $a:i$, $g:i\to i\to i$, $F:i\to i$ and consider $\forall a\,\forall g\,\exists F\,[(F\,a)=(g\,a\,a)]$. It has
exactly four unifiers,
\[
  F:=\lambda x\,(g\,x\,x),\qquad F:=\lambda x\,(g\,a\,x),\qquad F:=\lambda x\,(g\,x\,a),\qquad F:=\lambda x\,(g\,a\,a),
\]
because each occurrence of $a$ on the right can come either from the argument of $F$ or from the constant
itself. All four are closed, so none is an instance of another: \emph{there is no mgu}.
\end{example}

\begin{definition}[Complete set of unifiers {\book{213}}]
A set $\mathcal U$ of substitutions is a \emph{complete set of unifiers} (CSU) for a problem if every member of
$\mathcal U$ is a unifier, and every unifier is obtained from some member of $\mathcal U$ by a further
substitution. One would like $\mathcal U$ to be finite.
\end{definition}

\begin{example}[No finite CSU {\book{213}}]\label{ex:hounif-infcsu}
With $w:i$ and $F,u:i\to i$, the problem
\[\forall u\,\exists F\,\forall w\,[(u\,(F\,w))=(F\,(u\,w))]\]
is solved by $F:=\lambda y\,y$, $\lambda y\,(u\,y)$, $\lambda y\,(u\,(u\,y))$, \dots, that is, by every power $u^k$. Each of these is
closed, so every CSU contains all of them.
\end{example}

One might settle for a CSU without redundancy, in which no two members have a common instance; but \enquote{a
general guarantee of this kind of nonredundancy cannot be given in the higher-order setting} \book{213}. Worse,
\emph{unifiability} itself is undecidable: \cref{ch:lambda} encoded the Post correspondence problem as a
unification problem \book{114}, so \enquote{higher-order unification also must be undecidable} \book{214}. The
bibliographic notes add that undecidability already holds at third order (Huet; Lucchesi) and even at second
order (Goldfarb, 1981), while \emph{matching}, where only one side contains existential variables, is decidable
(second order: Huet and Lang 1978; third order: Dowek 1992; general case: Stirling 2009) \book{226--227}.

\begin{pitfall}[do not promise \enquote{the} answer]
A higher-order equation in general has no single best answer, may have infinitely many incomparable ones, and
may run forever without being able to report \enquote{no}. Any design that represents an answer as \emph{one}
substitution, as \micro{} does, can be exact only on a fragment that has mgus. That fragment is $\Llam$.
\end{pitfall}

% ============================================================================================
\section{Huet's procedure}\label{sec:hounif-huet}

Despite all of the above there is a systematic, complete search \book{214}. We present it as a rewrite system on
pairs $\langle\mathcal Q;E\rangle$ of a prefix and a multiset of equations, in the style of Snyder and Gallier
that the book mentions as an alternative presentation \book{227}. The side conditions and substitution terms
are the book's; the rule names are ours.

\subsection{Simplification of rigid-rigid equations}

\begin{mathpar}
\inferrule{\langle\mathcal Q;\ \{(c\,s_1\dots s_n)=(c\,t_1\dots t_n)\}\uplus E\rangle}
          {\langle\mathcal Q;\ \{s_1=t_1,\dots,s_n=t_n\}\cup E\rangle}\;\rulename{Dec}
\and
\inferrule{\langle\mathcal Q;\ \{(c\,\bar s)=(d\,\bar t)\}\uplus E\rangle}{\bot}\;\rulename{Clash}\ \ (c\neq d)
\end{mathpar}
Here $c$, $d$ are universally quantified in $\mathcal Q$; for $n=0$, \rulename{Dec} just deletes the equation.
Since the head of a rigid term can never change, an equation between rigid terms with different heads is
hopeless, and one with equal heads is equivalent to the equations between the arguments. \rulename{Dec} removes
two symbols from the problem \book{214}. \enquote{The process of repeatedly treating all rigid-rigid pairs in this
fashion will be called \emph{simplification}} \book{214}. Rigid-flexible equations are turned around and treated
as flexible-rigid ones.

\subsection{Imitation and projection}

Take a flexible-rigid equation
\[
  (F\,t_1\dots t_n)=(c\,s_1\dots s_m),\qquad F:\tau_1\to\dots\to\tau_n\to\beta,\quad c:\sigma_1\to\dots\to\sigma_m\to\beta,
\]
with $\beta$ primitive. \enquote{This equation can be proved only by making the head of the left term identical
to $c$} \book{215}, and there are only two ways to do that: the term substituted for $F$ has head $c$ itself, or it
uses one of $F$'s arguments as head.

\begin{definition}[Imitation and projection {\book{215--216}}]\label{def:hounif-imitproj}
\leavevmode
\begin{itemize}
  \item The \emph{imitation} substitution, allowed only if \enquote{the quantification over $F$ occurs within the
    scope of the quantifier that binds $c$}, is
    \[\theta_{\mathrm{imit}}=\{\langle F,\ \lambda x_1\dots\lambda x_n\,(c\,(H_1\,x_1\dots x_n)\cdots(H_m\,x_1\dots x_n))\rangle\}\]
    with fresh $H_i:\tau_1\to\dots\to\tau_n\to\sigma_i$.
  \item The $i$-th \emph{projection} substitution, allowed only if the target type of $\tau_i$ is $\beta$, say
    $\tau_i=\varphi_1\to\dots\to\varphi_k\to\beta$, is
    \[\theta_{\mathrm{proj}_i}=\{\langle F,\ \lambda x_1\dots\lambda x_n\,(x_i\,(H_1\,x_1\dots x_n)\cdots(H_k\,x_1\dots x_n))\rangle\}\]
    with fresh $H_j:\tau_1\to\dots\to\tau_n\to\varphi_j$.
\end{itemize}
In both cases $\exists F$ is replaced in the prefix by the new $\exists H_j$ at the same place, $\theta$ is
applied to the matrix, and the result is $\beta$-normalised.
\end{definition}

The design idea is minimal commitment. Imitation fixes only the head $c$ and leaves \enquote{all choices of what
the arguments should be to the determination of substitutions for the $H_j$ variables in subsequent steps}
\book{215}; each $H_j$ receives all of $F$'s arguments, so nothing is lost. Projection commits $F$ \enquote{only to
projecting onto its $i$th argument} \book{216}. For one flexible-rigid equation there may or may not be an
imitation and there are between $0$ and $n$ projections, so \enquote{up to $n+1$ substitutions may have to be
considered} \book{216}.

\begin{pitfall}[imitation is scope-sensitive, projection is type-sensitive]
Imitating $c$ is illegal when $c$ is quantified \emph{inside} the scope of $F$: $F$ cannot mention it. This
side condition is the logical root of every \enquote{eigenvariable escape} check. Projection needs the target
type and the arity of $\tau_i$, and imitation needs the arity $m$ of $c$ and its argument types $\sigma_i$:
\emph{Huet's procedure cannot run without types}. Pattern unification can (\cref{sec:hounif-algo}).
\end{pitfall}

\subsection{A complete run, and the matching tree}

\begin{tutorial}[Huet's procedure on \texttt{(F a) = (g a a)}]
We solve \cref{ex:hounif-nomgu}, following the book \book{216--217}.
\begin{enumerate}
  \item The only equation is flexible-rigid with head $g$, and $\forall g$ is to the left of $\exists F$, so
    imitation is allowed: $F:=\lambda x\,(g\,(H_1\,x)\,(H_2\,x))$ with $H_1,H_2:i\to i$. The single projection is
    $F:=\lambda x\,x$ (the argument type $i$ is the target type).
  \item Projection gives $\forall a\,\forall g\,[a=(g\,a\,a)]$: a rigid-rigid clash, a \emph{dead end}.
  \item Imitation gives $\forall a\,\forall g\,\exists H_1\,\exists H_2\,[(g\,(H_1\,a)\,(H_2\,a))=(g\,a\,a)]$, which
    \rulename{Dec} simplifies to $[(H_1\,a)=a\wedge(H_2\,a)=a]$.
  \item For $(H_1\,a)=a$ both the imitation $\lambda x\,a$ and the projection $\lambda x\,x$ work; either leaves
    $[\top\wedge(H_2\,a)=a]$.
  \item The same two choices for $H_2$ empty the matrix.
\end{enumerate}
Composing the substitutions along the four successful paths gives the four unifiers
$\lambda x\,(g\,a\,a)$, $\lambda x\,(g\,a\,x)$, $\lambda x\,(g\,x\,a)$, $\lambda x\,(g\,x\,x)$ \book{217}, in the order of a
left-to-right depth-first walk of \cref{fig:hounif-tree}.
\end{tutorial}

\begin{figure}[t]
\centering
\begin{tikzpicture}[
  level distance=17mm,
  level 1/.style={sibling distance=78mm},
  level 2/.style={sibling distance=44mm},
  level 3/.style={sibling distance=22mm},
  every node/.style={font=\small},
  lab/.style={font=\scriptsize,inner sep=1.5pt,fill=white}]
\node {$\{(F\,a)=(g\,a\,a)\}$}
  child { node {$\{(H_1\,a)=a,\ (H_2\,a)=a\}$}
    child { node {$\{(H_2\,a)=a\}$}
      child { node {$\{\}$} edge from parent node[lab,left] {$H_2{:=}\lambda x\,a$} }
      child { node {$\{\}$} edge from parent node[lab,right] {$H_2{:=}\lambda x\,x$} }
      edge from parent node[lab,left] {$H_1{:=}\lambda x\,a$} }
    child { node {$\{(H_2\,a)=a\}$}
      child { node {$\{\}$} edge from parent node[lab,left] {$H_2{:=}\lambda x\,a$} }
      child { node {$\{\}$} edge from parent node[lab,right] {$H_2{:=}\lambda x\,x$} }
      edge from parent node[lab,right] {$H_1{:=}\lambda x\,x$} }
    edge from parent node[lab,left,xshift=-2mm] {$F{:=}\lambda x\,(g\,(H_1\,x)\,(H_2\,x))$} }
  child { node {$\mathbf{F}$}
    edge from parent node[lab,right,xshift=2mm] {$F{:=}\lambda x\,x$} };
\end{tikzpicture}
\caption{The matching tree for \mbox{$\forall a\,\forall g\,\exists F\,((F\,a)=(g\,a\,a))$}, after Figure~8.1 of the book
\book{219}. Arcs are labelled by substitutions; $\mathbf{F}$ marks a failure node; the four $\{\}$ leaves are
success nodes.}
\label{fig:hounif-tree}
\end{figure}

\begin{definition}[Matching tree {\book{219}}]\label{def:hounif-tree}
The \emph{matching tree} of a problem has as nodes the simplified matrices, viewed as multisets of equations
(the prefix is implicit), and as arcs the imitation and projection substitutions for one selected
flexible-rigid equation. Leaves are \emph{failure nodes}, \enquote{for which it is immediately apparent that there
are no solutions}, or \emph{success nodes}, whose matrix contains at most flexible-flexible equations. The
\emph{pre-unifier} of a success node is the composition of the arc labels on the path to it.
\end{definition}

\begin{theorem}[Completeness of matching trees {\book{219--220}}]\label{thm:hounif-tree}
If the problem at the root has a unifier, then the tree has a success node at a finite depth, and that unifier
is obtained by composing the pre-unifier of the success node with a unifier of the problem at that node. This
holds however the flexible-rigid equation to expand is chosen at each node.
\end{theorem}

So the choice of \emph{which equation} to work on is don't-care nondeterminism, while the choice of \emph{which
substitution} is don't-know nondeterminism. The bibliographic notes add that Huet's procedure, \enquote{which is
limited to determining unifiability, is nonredundant} \book{227}.

\subsection{Flexible-flexible equations and pre-unification}

Why stop at flexible-flexible equations? Because they always have a unifier, and their unifiers are
uncontrollable.

\begin{example}[A flexible-flexible residue {\book{217}}]\label{ex:hounif-residue}
With $a:i$ and $X,Y,g:i\to i$, take $\forall a\,\forall g\,\exists X\,\exists Y\,[(X\,a)=(g\,(Y\,a))]$. Imitation
$X:=\lambda x\,(g\,(H\,x))$ and simplification give $\forall a\,\forall g\,\exists H\,\exists Y\,[(H\,a)=(Y\,a)]$, which no
rule can reduce further. (The other alternative for $X$, the projection $X:=\lambda x\,x$, gives $a=(g\,(Y\,a))$,
a rigid-rigid clash.) The residue has the unifier $H,Y:=\lambda x\,a$ (making $X=\lambda x\,(g\,a)$), and many more: any
term built from $a$, $g$ and other existentials, abstracted over some of the occurrences of $a$, will do,
\enquote{too large and uncontrolled a set to try to generate explicitly} \book{218}.
\end{example}

\begin{proposition}[Canonical substitution {\book{217--218}}]\label{prop:hounif-canonical}
Associate with each primitive type $\tau$ a fixed new existential variable $H^\tau$, and with each type
$\tau_1\to\dots\to\tau_n\to\sigma$ ($\sigma$ primitive) the \emph{canonical constant term}
$\lambda x_1\dots\lambda x_n\,H^\sigma$. The \emph{canonical substitution} maps every existential variable of the
prefix to the canonical constant term of its type. For a problem with only flexible-flexible equations, the
canonical substitution is a unifier.
\end{proposition}

The reason is simple: under it both sides of every equation become $H^\beta$. For \cref{ex:hounif-residue}, with
$W$ the variable chosen for $i$, the canonical substitution is $\{\langle H,\lambda x\,W\rangle,\langle Y,\lambda x\,W\rangle\}$
\book{217}. Hence: \enquote{The reduction of a unification problem to a flexible-flexible form thus constitutes
success from the perspective of determining unifiability. The process of carrying out such a reduction is known
as \emph{pre-unification}, and the substitution that leads to the corresponding success is called a
\emph{pre-unifier}} \book{218}.

\subsection{Non-termination}

\begin{example}[An infinite branch with infinitely many answers {\book{218}}]\label{ex:hounif-loop}
Let $x:i$ and $g,F:i\to i$: $\forall g\,\exists F\,\forall x\,[(F\,(g\,x))=(g\,(F\,x))]$. The single projection
$F:=\lambda x\,x$ solves it. The imitation $F:=\lambda x\,(g\,(H\,x))$ followed by simplification gives
$\forall g\,\exists H\,\forall x\,[(H\,(g\,x))=(g\,(H\,x))]$: the same problem with $H$ in place of $F$. Repeating
forever, the paths of the tree yield $F:=\lambda x\,x,\ \lambda x\,(g\,x),\ \lambda x\,(g\,(g\,x)),\dots$
\end{example}

\begin{example}[An infinite branch and no answer {\book{220}}]\label{ex:hounif-noanswer}
Add a second equation, with constants $a,b:i$: $\forall g\,\exists F\,\forall x\,[((F\,(g\,x))=(g\,(F\,x)))\wedge(F\,a)=(F\,b)]$.
\enquote{The second equation in the matrix here causes each of the unifiers that we have seen for the first
equation to be rejected} \book{220}: $F=g^k$ gives $g^k\,a=g^k\,b$, a clash. The tree has no success node but has an
infinite path. As unifiability is undecidable, such problems must exist \book{220}.
\end{example}

\begin{kanren}[the matching tree is a miniKanren search tree]
Imitation versus projection is don't-know choice among at most $n+1$ alternatives, i.e.\ a disjunction; the
matching tree is the search tree of those disjunctions. A depth-first interpreter, like $\lambda$Prolog's, can be
trapped in the infinite imitation branch of \cref{ex:hounif-loop} before it ever tries the projection.
miniKanren's interleaving \scm{or°} instead reaches every success node at finite depth, provided each expansion
step is guarded by a suspension (as \scm{define-relation} and \scm{fresh°} provide). The fairness is
skewed: \scm{or°} nests the binary interleaving to the right, so with $k$ alternatives the first one gets
about half of the steps; and conjunction gives no fairness guarantee across an infinite first conjunct.
Two hard limits remain. First, \cref{ex:hounif-noanswer} makes even \scm{(run 1 ...)} diverge, and no search
strategy can fix that. Second, imitation and projection need types at run time, which \micro{} does not have.
A Huet layer is therefore at best an opt-in, typed search over equations that the deterministic core has set
aside; see \cref{sec:roadmap-later-huet} (and \cref{sec:design-space} for why Huet's search is not the meaning of \scm{=°}).
\end{kanren}

% ============================================================================================
\section{Higher-order patterns}\label{sec:hounif-llambda}

\Cref{ch:lambdatree} introduced $\Llam$ as a programming discipline. The book recalls it as follows.

\begin{definition}[The $\Llam$ condition {\book{220}}, {\cite{Miller1991}}]\label{def:hounif-llambda}
A program clause or goal formula $B$ is in $\Llam$ if all quantification in $B$ is over non-predicate types and
\enquote{every subterm of $B$ of the form $(x\,t_1\dots t_n)$ in which $x$ is an essentially existentially
quantified variable is such that $t_1,\dots,t_n$ are distinct variables that are universally quantified within
the scope of the quantifier binding $x$} \book{220}. Transposed to a unification problem: every flexible subterm
$(F\,t_1\dots t_n)$ has as arguments \emph{distinct} universal variables quantified \emph{to the right of $F$} in
the prefix. Such a term is a \emph{(higher-order) pattern}.
\end{definition}

The case $n=0$ is allowed, so a bare logic variable is a pattern, and first-order unification is a special
case. Miller's formal definition~\cite[Fig.~3, p.~10--11]{Miller1991} is a judgement $Q\vdash_0 t:\tau$ relative
to a prefix $Q$, whose key rule says that $x\,y_1\cdots y_n$ is well formed when \enquote{$x$ is existentially
quantified in $Q$ to the left of where the distinct variables $y_1\dots y_n$ are universally quantified}.
\Cref{tab:hounif-patterns} gives a feel for the definition.

\begin{table}[t]
\centering\small
\caption{Patterns and non-patterns among the chapter's examples.}
\label{tab:hounif-patterns}
\begin{tabular}{@{}lll@{}}
\toprule
problem & pattern? & why \\
\midrule
$\exists F\,\forall x\,\forall y\,[(F\,x\,y)=\dots]$ & yes & distinct universals, right of $F$ \\
$\forall a\,\exists F\,[(F\,a)=\dots]$ (\cref{ex:hounif-nomgu}) & no & $a$ is quantified \emph{outside} $F$ \\
$\exists F\,\forall x\,[(F\,x\,x)=x]$ & no & arguments not distinct \\
$\forall g\,\exists F\,\forall x\,[(F\,(g\,x))=\dots]$ (\cref{ex:hounif-loop}) & no & argument is not a variable \\
$\exists F\,\exists N\,[(F\,N)=\dots]$ (the \lp{subst} clause) & no & argument is existential \\
$\forall f\,\exists X\,[X=(f\,X)]$ & yes & $n=0$ \\
\bottomrule
\end{tabular}
\end{table}

\begin{example}[Discharging a constant {\cite[p.~4--5]{Miller1991}}]\label{ex:hounif-discharge}
Miller motivates the fragment with \lp{append}. With the program
\[
  \forall x\,\forall l\,\forall k\,\forall m\,(\mathit{append}\ l\ k\ m\imp\mathit{append}\ (\mathit{cons}\ x\ l)\ k\ (\mathit{cons}\ x\ m)),
  \qquad\forall k\,(\mathit{append}\ \mathit{nil}\ k\ k),
\]
try to find $X$ such that $\forall y\,(\mathit{append}\ (\mathit{cons}\ a\ (\mathit{cons}\ b\ \mathit{nil}))\ y\ X)$ is
provable. The $\forall y$ introduces a new constant $k$, and two backchaining steps succeed only if
$X$ is instantiated with $(\mathit{cons}\ a\ (\mathit{cons}\ b\ k))$, which is impossible because $X$ was chosen
before $k$ existed. \enquote{Such a failure here is quite sensible since the value of $X$ should be independent of the
choice of the constant used to instantiate $\forall y$}. Now ask instead for
$\forall y\,(\mathit{append}\ (\mathit{cons}\ a\ (\mathit{cons}\ b\ \mathit{nil}))\ y\ (H\,y))$. The same steps
produce the pattern equation
\[(H\,k)=(\mathit{cons}\ a\ (\mathit{cons}\ b\ k)).\]
Two $\lambda$-terms for $H$ make it true, $\lambda w\,(\mathit{cons}\ a\ (\mathit{cons}\ b\ k))$ and
$\lambda w\,(\mathit{cons}\ a\ (\mathit{cons}\ b\ w))$, but \enquote{since $H$ cannot contain $k$ free, only the second
of these possible substitutions will succeed}. Miller calls this \emph{discharging} the constant $k$ from the
term: $H$ is the computed answer with the choice of eigenvariable abstracted away. In terms of Huet's procedure: imitation of
$\mathit{cons}$, $a$, $\mathit{cons}$, $b$ is forced, and at the leaf $k$ only the projection is legal, because
imitating $k$ would violate the scope condition. The pattern algorithm below does this in one walk.
\end{example}

\begin{definition}[$\beta_0$ {\book{206}}]
$\beta_0$-conversion is the special case $(\lambda x\,s)\,x=s$ of $\beta$; with $\alpha$ it is equivalent to
$(\lambda y\,s)\,x=s[x/y]$, provided $x$ is not free in $\lambda y\,s$.
\end{definition}

Under the $\Llam$ condition, instantiating a logic variable only ever creates $\beta_0$-redexes: \enquote{simply
replacing the bound variables within the body of an abstraction with a new set of names} \book{220}. A
$\beta_0$-step makes the term smaller, so we may expect its inverse, $\Llam$-unification, to be much better
behaved than inverting full $\beta$ \book{220}.

\begin{pitfall}[$\beta_0$ alone is not enough]
$\Llam$ problems are unification problems modulo $\alpha,\beta_0,\eta$, but not conversely. The problem
$\forall x\,(F\,x\,x=x)$ needs only $\beta_0$, yet has the two incomparable unifiers $\lambda u\lambda v\,u$ and
$\lambda u\lambda v\,v$ \book{208}. The \emph{distinctness} of the arguments is what buys most general unifiers,
not merely the shape of the redexes.
\end{pitfall}

\subsection{Why the procedure becomes deterministic}

Raising preserves the $\Llam$ condition \book{208, 220--221}, so we may still assume a $\forall\bar a\,\exists\bar X\,\forall\bar y$
prefix, and rigid-rigid simplification \enquote{applies unchanged} \book{221}. The interesting case is
flexible-rigid.

\begin{lemma}[At most one useful substitution {\book{221}}]\label{lem:hounif-one}
Let $(F\,c_1\dots c_n)=(c\,t_1\dots t_m)$ with $F$ existential and $c_1,\dots,c_n$ distinct universals quantified
in the scope of $F$. Then at most one imitation or projection substitution avoids an immediate clash.
\end{lemma}
\begin{proof}[Argument]
If $F$ is in the scope of $c$ ($c$ is to the left of $F$, e.g.\ a signature constant), then $c$ differs from every
$c_i$, which are to the right of $F$; every projection produces a head $c_i\ne c$ and clashes, so only
imitation can succeed. If $c$ is in the scope of $F$, imitation is not allowed, and the $i$-th projection
succeeds only if $c_i=c$; since the $c_i$ are distinct, at most one $i$ qualifies, and if $c\notin\bar c$ the
equation fails. \enquote{The preceding observations can be built into the substitution selection process so as to
make it entirely deterministic} \book{221}.
\end{proof}

Moreover the class is preserved: the substitutions considered have the form
\[\lambda x_1\dots\lambda x_n\,(@\,(H_1\,x_1\dots x_n)\cdots(H_k\,x_1\dots x_n))\]
with $@$ either $c$ or some $x_i$ and the
$H_j$ quantified where $F$ was; substituted into $(F\,c_1\dots c_n)$ and $\beta_0$-reduced, it yields again
patterns $(H_j\,c_1\dots c_n)$ \book{221}.

\subsection{Flexible-flexible pattern equations have mgus}\label{sec:hounif-ff}

Here is the second change. In $\Llam$ a flexible-flexible equation has the form
\[(F\,c_1\dots c_n)=(G\,d_1\dots d_m),\]
with the $c_i$ and $d_j$ universals in the scope of $F$ and $G$, and in a $\forall\exists\forall$ prefix $F$ and $G$ are
\enquote{effectively bound \enquote{at the same place}} \book{222}.

\begin{proposition}[Flexible-flexible mgus {\book{222}}]\label{prop:hounif-ff}
Let $H$ be a new existential quantified at the place of $F$.
\begin{itemize}
  \item \textbf{Different heads} ($F\ne G$). Let $e_1,\dots,e_\ell$ list the variables common to $c_1,\dots,c_n$ and
    $d_1,\dots,d_m$. Then
    \[\{\langle F,\ \lambda c_1\dots\lambda c_n\,(H\,e_1\dots e_\ell)\rangle,\ \langle G,\ \lambda d_1\dots\lambda d_m\,(H\,e_1\dots e_\ell)\rangle\}\]
    is an mgu.
  \item \textbf{Same head} ($F=G$, so $n=m$). Let $e_1,\dots,e_\ell$ list the $c_i$ such that $c_i=d_i$.
    Then
    \[\{\langle F,\ \lambda c_1\dots\lambda c_n\,(H\,e_1\dots e_\ell)\rangle\}\]
    is an mgu.
\end{itemize}
\end{proposition}
In the first case a common instance $t$ cannot contain a $c_i$ that is not also among the $d_j$, and conversely;
every such $t$ is an instance of $(H\,\bar e)$. In the second case the \emph{same} substitution acts on both sides,
so a variable survives only if it sits in the \emph{same argument position} on both sides \book{222}. (The
binders reuse the names $c_i$, so in $\lambda c_1\dots\lambda c_n\,(H\,e_1\dots e_\ell)$ the $e$'s are bound.)

\begin{pitfall}[same head: positions, not sets]
For $\forall x\,\forall y\,[(F\,x\,y)=(F\,y\,x)]$ the arguments are the same \emph{set}, but no position agrees, so
$\ell=0$ and the mgu is $F:=\lambda x\lambda y\,H$. Intersecting sets, as in the different-heads case, would
wrongly keep both arguments.
\end{pitfall}

\begin{theorem}[Mgus for higher-order patterns {\book{222}}, {\cite{Miller1991}}]\label{thm:hounif-mgu-exist}
\enquote{By combining the completeness property of matching trees with the observations that there is at most one
productive path in the tree and that most general unifiers can be provided for flexible-flexible pairs in the
$\Llam$ setting, it follows that the higher-order pattern unification problem possesses the property of having
most general unifiers} \book{222}. For $\forall\exists\forall$ problems the mgu is unique up to renaming of
variables~\cite[Prop.~9.2]{Miller1991}.
\end{theorem}

\subsection{Termination needs the occurs-check}

The procedure so far may still loop, even on first-order problems. For $\forall f\,\exists X\,[X=(f\,X)]$ the only
useful substitution is the imitation $X:=(f\,H)$, which produces $\forall f\,\exists H\,[H=(f\,H)]$, \enquote{identical
to the one with which we started} \book{222}. In $\Llam$ such behaviour arises only when the flexible head
occurs in the rigid side; so \enquote{we can ensure termination for the unification procedure by adding an
\enquote{occurs-check} to the substitution-generation step} \book{223}. Note that in $\Llam$ every occurrence of $F$ in
the rigid side is fatal: since pattern arguments are variables, there is no flexible occurrence that a later
substitution could erase.

% ============================================================================================
\section{The pattern unification algorithm}\label{sec:hounif-algo}

The book's final step is to fuse the imitation and projection traversal with the occurs-check traversal into
\enquote{a larger, more efficient substitution generation operation} \book{223}, called \emph{variable elimination}
\enquote{because of its similarity to an operation of the same name discussed for first-order unification}
\book{223}. We now state the whole algorithm precisely. The rules restate the book's \S8.3 and Miller's
transitions~\cite[\S6, p.~14--16]{Miller1991}. Two things the book leaves implicit are made explicit here: binders
\emph{inside} terms, and the order of pruning.

\paragraph{Setting.} The prefix is $\forall\bar a\;\exists\bar X\;\forall\Gamma$, where $\bar a$ are the signature
constants and $\Gamma$ is the growing list of inner universals, or \emph{eigenvariables}. A state is a pair
$\langle E;\theta\rangle$ of a multiset of equations and the substitution computed so far. Equations are symmetric.
$\theta$ is applied eagerly, and terms are kept $\beta_0$-normal.

\begin{mathpar}
\inferrule{\langle\{\lambda x\,s=\lambda x\,t\}\uplus E;\theta\rangle}{\langle\{s[z/x]=t[z/x]\}\cup E;\theta\rangle}\;\rulename{Abs}
\and
\inferrule{\langle\{\lambda x\,s=t\}\uplus E;\theta\rangle}{\langle\{s[z/x]=(t\,z)\}\cup E;\theta\rangle}\;\rulename{Eta}
\\
\inferrule{\langle\{(h\,s_1\dots s_n)=(h\,t_1\dots t_n)\}\uplus E;\theta\rangle}{\langle\{s_1=t_1,\dots,s_n=t_n\}\cup E;\theta\rangle}\;\rulename{Dec}
\and
\inferrule{\langle\{(h\,\bar s)=(k\,\bar t)\}\uplus E;\theta\rangle}{\bot}\;\rulename{Clash}
\\
\inferrule{\langle\{(F\,\bar c)=t\}\uplus E;\theta\rangle}{\langle\rho'(\rho(E));\ \rho'\circ\rho\circ\theta\rangle}\;\rulename{Elim}
\and
\inferrule{\langle\{(F\,\bar c)=(G\,\bar d)\}\uplus E;\theta\rangle}{\langle\sigma(E);\ \sigma\circ\theta\rangle}\;\rulename{FlexFlex}
\end{mathpar}
Side conditions:
\begin{itemize}
  \item \rulename{Abs}, \rulename{Eta}: $z$ is a fresh eigenvariable appended to $\Gamma$; in \rulename{Eta}, $t$ is
    not an abstraction. This is step~4 of the normal form, applied on demand.
  \item \rulename{Dec}, \rulename{Clash}: $h,k$ are rigid heads (signature constants or eigenvariables) and $h\ne k$
    in \rulename{Clash}. Untyped, an arity mismatch for the same head is also a clash~\cite[\S9.3]{Miller1991}.
  \item \rulename{Elim}: $(F\,\bar c)$ is a pattern, $t$ is rigid, and the elimination
    $\mathcal E_{F,\bar c}(t)$ of \cref{def:hounif-elim} succeeds with result $t'$ and pruning substitution $\rho$;
    then $\rho'=\{\langle F,\lambda\bar c.\,t'\rangle\}$. If $\mathcal E$ fails, the state goes to $\bot$.
  \item \rulename{FlexFlex}: both sides are patterns and $\sigma$ is the mgu of \cref{prop:hounif-ff}; for the
    same head with all positions agreeing, the equation is simply deleted.
  \item A flexible term that is \emph{not} a pattern blocks the rules; we come back to it in
    \cref{sec:hounif-dynamic}.
\end{itemize}

\begin{definition}[Variable elimination]\label{def:hounif-elim}
Let $(F\,c_1\dots c_n)$ be a pattern. The elimination $\mathcal E_L(t)$ of a $\beta_0$-normal term $t$, relative to a
set $L$ of variables bound by $\lambda$s \emph{inside} $t$ (initially $L=\emptyset$), either fails or returns a term,
and may record pruning bindings in $\rho$. It is defined by cases on $t$, after applying the pruning bindings
recorded so far:
\[
\begin{array}{@{}l@{\quad}l@{\quad}l@{}}
  \mathcal E_L(a) &= a & a \text{ a signature constant (outer } \forall)\\
  \mathcal E_L(y) &= y & y\in\bar c\cup L\\
  \mathcal E_L(y) &= \text{fail (scope)} & y\in\Gamma\setminus(\bar c\cup L)\\
  \mathcal E_L(\lambda z\,s) &= \lambda z\,\mathcal E_{L\cup\{z\}}(s)\\
  \mathcal E_L(h\,s_1\dots s_k) &= \mathcal E_L(h)\,\mathcal E_L(s_1)\dots\mathcal E_L(s_k) & h \text{ rigid}\\
  \mathcal E_L(F\,\bar s) &= \text{fail (occurs)}\\
  \mathcal E_L(G\,d_1\dots d_m) &= G'\,e_1\dots e_\ell & G\ne F \text{ a pattern; see below}
\end{array}
\]
In the last case $e_1\dots e_\ell$ is the subsequence of $d_1\dots d_m$ of the variables in $\bar c\cup L$. If
$\ell=m$, nothing is pruned and $G'=G$. Otherwise the \emph{pruning substitution}
$\langle G,\ \lambda d_1\dots\lambda d_m\,(G'\,e_1\dots e_\ell)\rangle$, with $G'$ new and quantified at the place of $G$,
is added to $\rho$.
\end{definition}

This is the book's description \book{223}, case by case. \enquote{This pruning substitution replaces $G$ with the term
$\lambda d_1\dots\lambda d_m\,(G'\,e_1\dots e_\ell)$, where $G'$ is a new existentially quantified variable, and
$e_1,\dots,e_\ell$ is some listing of the variables that are common to $c_1,\dots,c_n$ and $d_1,\dots,d_m$. Now let the
result of applying these auxiliary substitutions to the term $t$ be $t'$. Then the main part of variable
elimination checks whether there are occurrences in $t'$ either of $F$ or of variables different from the ones in
$c_1,\dots,c_n$ that are quantified universally within the scope of the quantifier binding $F$. If there are such
occurrences, then it marks the problem as one for which no unifiers exist. Otherwise, it generates the
substitution $\lambda c_1\dots\lambda c_n\,t'$ for $F$ and replaces the equation under consideration by $\top$.}

Read the definition as an \emph{inversion}. $F$ applied to $\bar c$ must produce $t$, so $\lambda\bar c.\,t'$ must
\enquote{undo} the application: every $c_i$ in $t$ becomes the $i$-th bound variable of $F$'s binder. A universal
that $F$ cannot see and that $F$ does not receive as an argument has no preimage. In a rigid position it makes the
problem unsolvable; inside the arguments of another flexible term $G$ it can be removed by \emph{pruning}, that
is, by committing $G$ to ignore that argument. Imitation (rigid heads copied) and projection (a $c_i$ copied as a
bound variable) are both inside this single walk.

\begin{pitfall}[four things the prose leaves implicit]
\leavevmode
\begin{enumerate}[nosep]
  \item An out-of-scope universal is fatal only in a \emph{rigid} position. Inside the argument list of a
    flexible subterm it is pruned. The book's \cref{ex:hounif-8j} below says exactly this: \enquote{in a position that
    is not the argument of a flexible term} \book{224}.
  \item Variables bound by $\lambda$s \emph{inside} $t$ are always allowed, both in rigid positions and as kept
    arguments of pruned terms; this is the set $L$. The book never needs $L$ because its normal form has no inner
    binders at the top, but terms like $h\,(\lambda z\,(G\,z\,w))$ have them below the top.
  \item Pruning \emph{binds other variables} ($G$) as a side effect, although the equation is about $F$. If $G$
    occurs twice, the second occurrence must be examined \emph{after} the first pruning, which may prune again
    (\cref{ex:hounif-two}).
  \item Pruning never meets $F$ itself: that is the occurs failure.
\end{enumerate}
\end{pitfall}

\begin{theorem}[Pattern unification is unitary and terminating]\label{thm:hounif-mgu}
On a $\forall\exists\forall$ problem whose flexible subterms are all patterns, the rules terminate. If they reach $\bot$
the problem has no unifier; otherwise they end in $\langle\emptyset;\theta\rangle$ and $\theta$ is an mgu
(\book{222--224}; \cite[Thm.~7.1, Prop.~7.3]{Miller1991}).
\end{theorem}

Miller proves termination with a lexicographic measure: the number of existentials, the total size of the
equations, the number of equations, a weight of the universals, and the number of possibly prunable
occurrences~\cite[Thm.~7.1, p.~18]{Miller1991}. Intuitively, \rulename{Elim} and the different-heads case of
\rulename{FlexFlex} each remove one existential for good (a pruned $G$ is traded for one $G'$; $F$ and $G$ are
traded for one $H$); the same-head case trades $F$ for one $H$, so the number of existentials stays put, but the
equation disappears (both sides become $H\,\bar e$); and the other rules shrink the problem. The
occurs-check is what prevents the $X=(f\,X)$ loop. Qian later showed the problem can be solved in linear
time~\cite{Qian1996}.

Two remarks complete the picture. First, \enquote{types have no bearing on the shapes of unifiers}: the
algorithm \enquote{makes no use at all of type information in determining unifying substitutions}, though types are
still needed to guarantee correctness, for example that same-head flexible terms have equal arity \book{224}.
Miller's untyped variant replaces the type guarantee by two failures, an arity mismatch in \rulename{Dec} and
same-head flexible terms of different arity, and uses $\eta$-expansion \enquote{only on terms which, if typed, must
have functional type}~\cite[\S9.3]{Miller1991}. That is the variant an untyped Scheme host needs. Second, the
result is a genuine algorithm, as for first-order unification: \enquote{In both cases, term simplification and
variable elimination constitute the critical parts of the computation, although variable elimination has a more
involved structure in the higher-order pattern unification case} \book{224}.

% ============================================================================================
\section{Worked traces}\label{sec:hounif-traces}

Every result below was also checked mechanically with the Scheme transliteration of \cref{sec:hounif-nipkow}.
The book does not number its unification problems. We label them 8.A--8.M in their order of appearance in
the book's Chapter~8, as the proposal's test suite does, except 8.M, the problem $\forall x\,(F\,x\,x=x)$ of
\book{208} in the book's Chapter~7. The main example below is 8.H, the occurs and scope failures are 8.I
and 8.J, and the eager-versus-delayed raising of \cref{ex:hounif-8k} is 8.K.

\begin{example}[The book's main example {\book{223--224}}]\label{ex:hounif-8h}
With $x,y,w:i$, $f,g:i\to i$ and $U,V:i\to i\to i$, solve
\[
  \forall f\,\forall g\,\exists U\,\exists V\,\forall w\,\forall x\,\forall y\,[(f\,(U\,x\,y))=(f\,(g\,(V\,y\,w)))].
\]
\end{example}
\begin{tutorial}[the book's main example, step by step]
\begin{enumerate}
  \item Both sides are rigid with head $f$: \rulename{Dec} leaves $(U\,x\,y)=(g\,(V\,y\,w))$.
  \item The left side is the pattern $U$ applied to $\bar c=(x,y)$, and the right side is rigid: \rulename{Elim}.
    Walk $t=g\,(V\,y\,w)$ with $L=\emptyset$:
    \begin{itemize}[nosep]
      \item $g$ is a signature constant, so it is kept (this is the imitation part);
      \item $(V\,y\,w)$ is a flexible pattern with $\bar d=(y,w)$. Of these, $y\in\bar c$ is kept and $w\notin\bar c$ must
        go: prune $V:=\lambda y\,\lambda w\,(V'\,y)$, so the subterm becomes $(V'\,y)$.
    \end{itemize}
    Now $t'=g\,(V'\,y)$. $U$ does not occur in it, and its only inner universal $y$ is in $\bar c$, so bind
    $U:=\lambda x\,\lambda y\,(g\,(V'\,y))$.
  \item The matrix is empty. The mgu is $\{\langle V,\lambda y\lambda w\,(V'\,y)\rangle,\ \langle U,\lambda x\lambda y\,(g\,(V'\,y))\rangle\}$
    \book{224}.
\end{enumerate}
Check: $U\,x\,y\to_{\beta_0}g\,(V'\,y)$ and $V\,y\,w\to_{\beta_0}V'\,y$, so both sides become $f\,(g\,(V'\,y))$.
\end{tutorial}

\begin{example}[Occurs failure {\book{224}}]\label{ex:hounif-8i}
$\forall f\,\forall g\,\exists U\,\exists V\,\forall w\,\forall x\,\forall y\,[(U\,x\,y)=(g\,(U\,y\,w))]$: the walk meets
$(U\,y\,w)$, whose head is the variable being eliminated, so the problem fails, \enquote{because of the occurrence of
$U$, the head of the flexible term, in the rigid term}. Note that pruning cannot help: whatever $U$ does with its
arguments, the right side is one $g$ larger than the left.
\end{example}

\begin{example}[Scope failure {\book{224}}]\label{ex:hounif-8j}
$\forall f\,\forall g\,\exists U\,\exists V\,\forall w\,\forall x\,\forall y\,[(U\,x\,y)=(g\,w)]$ \enquote{has no unifiers
because the universally quantified variable $w$, which is not one of the arguments in $(U\,x\,y)$, appears on the
right-hand side of the equation and in a position that is not the argument of a flexible term}. Logically: $U$ is
chosen before $w$ and does not receive it, so it cannot produce it.
\end{example}

\begin{example}[Discharging, by elimination]
For \cref{ex:hounif-discharge}, $(H\,k)=(\mathit{cons}\ a\ (\mathit{cons}\ b\ k))$ with $\bar c=(k)$: the walk keeps the
signature constants $\mathit{cons}$, $a$, $b$ (imitation), meets $k\in\bar c$ and maps it to $H$'s bound variable
(projection), finds no flexible subterm to prune and no occurrence of $H$, and returns
$H:=\lambda w\,(\mathit{cons}\ a\ (\mathit{cons}\ b\ w))$, Miller's unique answer.
\end{example}

\begin{kanren}[the scope failure that surprised \texorpdfstring{$\lambda$}{λ}Kanren]
Ma and Friedman~\cite{MaFriedman2021} open with the miniKanren query
\texttt{(run 1 X (== (λ (a) X) (λ (b) b)))} and observe that it \enquote{fails. Under any existing higher-order
unification algorithm, X and a cannot be unified due to the scope restrictions}, while a novice expects
\texttt{X = a}. Read through \rulename{Abs}, the equation is $\exists X\,\forall z\,[X=z]$: \cref{ex:hounif-8j}
with $n=0$, so the failure is the \emph{correct} answer under the quantifier reading (no single $X$ equals every
$z$). Their proposal changes the semantics with de~Bruijn levels so that the query answers \texttt{(d0)};
\cref{ch:lkanren} examines that choice. Whichever semantics a new module picks, it must say so explicitly;
see \cref{ch:design}.
\end{kanren}

\begin{example}[Flexible-flexible, different heads]\label{ex:hounif-ffdiff}
For \cref{ex:hounif-ch4}, $\exists G\,\exists H\,\forall x\,\forall y\,[(G\,x)=(H\,y)]$: here $\bar c=(x)$ and
$\bar d=(y)$ have nothing in common, $\ell=0$, and \cref{prop:hounif-ff} gives
\[\{\langle G,\lambda x\,T\rangle,\ \langle H,\lambda y\,T\rangle\},\]
the book's \enquote{$\lambda w\,T$ for both} \book{113}. Had the
equation been $(G\,x\,y)=(H\,y)$, then $\bar e=(y)$ and the mgu would be $G:=\lambda x\lambda y\,(T\,y)$,
$H:=\lambda y\,(T\,y)$.
\end{example}

\begin{example}[Pruning under an inner binder]\label{ex:hounif-loc}
$\forall h\,\exists F\,\exists G\,\forall w\,\forall x\,[(F\,x)=(h\,(\lambda z\,(G\,z\,w)))]$. \rulename{Elim} with $\bar c=(x)$
walks $h$ (kept), then the abstraction, adding $z$ to $L$, then $(G\,z\,w)$: $z\in L$ is kept and $w$ is pruned,
$G:=\lambda z\lambda w\,(G'\,z)$. Result: $F:=\lambda x\,(h\,(\lambda z\,(G'\,z)))$. The local $z$ is legal everywhere below its
binder, which is item~2 of the pitfall after \cref{def:hounif-elim}.
\end{example}

\begin{example}[Two occurrences, two prunings]\label{ex:hounif-two}
$\forall g\,\exists F\,\exists G\,\forall w\,\forall x\,[(F\,x)=(g\,(G\,x\,w)\,(G\,w\,x))]$, with $\bar c=(x)$. The first
occurrence prunes the second argument: $G:=\lambda d_1\lambda d_2\,(G'\,d_1)$. The second occurrence now reads $(G'\,w)$,
and $w$ must be pruned again: $G':=\lambda d\,G''$. So $G:=\lambda d_1\lambda d_2\,G''$ and $F:=\lambda x\,(g\,G''\,G'')$. This
is correct: $G\,x\,w$ must not mention $w$, so $G$ ignores its second argument, and $G\,w\,x$ must not mention $w$,
so it ignores its first.
\end{example}

\begin{example}[$\eta$ on demand]
$\forall g\,\exists F\,[\lambda x\,(F\,x)=g]$: \rulename{Eta} gives $(F\,z)=(g\,z)$ with $z$ fresh, and \rulename{Elim}
binds $F:=\lambda z\,(g\,z)$, which is $\eta$-equal to $g$. Without \rulename{Eta}, an abstraction facing a constant
would be a spurious clash.
\end{example}

\begin{table}[t]
\centering\footnotesize
\caption{A test battery. The Huet column describes the pre-unification procedure exactly as the book presents it
in its \S8.2: it stops at flexible-flexible equations, and it has no occurs-check, since that is added only for
$\Llam$ \book{223}. The pattern column assumes the delaying discipline of \cref{sec:hounif-dynamic}.}
\label{tab:hounif-battery}
\begin{tabular}{@{}>{\raggedright\arraybackslash}p{0.29\textwidth}l>{\raggedright\arraybackslash}p{0.2\textwidth}>{\raggedright\arraybackslash}p{0.33\textwidth}@{}}
\toprule
problem (matrix) & $\Llam$? & pattern unification & Huet pre-unification \\
\midrule
$(F\,a)=(g\,a\,a)$ & no & not applicable (delay) & 4 answers \\
$(G\,x)=(H\,y)$ & yes & $G,H:=\lambda w\,T$ & already flex-flex; residue kept \\
$X=(f\,X)$ & yes & fail (occurs) & diverges (no occurs-check) \\
$(f\,(U\,x\,y))=(f\,(g\,(V\,y\,w)))$ & yes & mgu of \cref{ex:hounif-8h} & $U:=\lambda x\lambda y\,g\,(H\,x\,y)$, residue $(H\,x\,y)=(V\,y\,w)$ \\
$(U\,x\,y)=(g\,(U\,y\,w))$ & yes & fail (occurs) & diverges \\
$(U\,x\,y)=(g\,w)$ & yes & fail (scope) & fail \\
$(F\,(g\,x))=(g\,(F\,x))$ & no & delay & infinitely many answers \\
$(F\,x\,x)=x$ & no & delay & 2 answers \\
$(F\,x\,y)=(F\,y\,x)$ & yes & $F:=\lambda x\lambda y\,H$ & already flex-flex; residue kept \\
\bottomrule
\end{tabular}
\end{table}

% ============================================================================================
\section{Nipkow's de Bruijn algorithm}\label{sec:hounif-nipkow}

Nipkow~\cite{Nipkow1993} turned the transformation rules into \enquote{a directly executable functional program}
and then adapted it \enquote{to lambda-terms in de Bruijn's notation}, for typed and untyped terms alike. It is short and
complete, so we reconstruct it from the author's ML code.
It assumes a $\forall\exists\forall$ prefix, raised in a preprocessing phase \book{227}.

\paragraph{Encoding the prefix in the terms.} Existentials are global names \texttt{V F}, signature constants are
\texttt{C c}, and an abstraction \texttt{A t} has a body in which bound variables are de~Bruijn indices \texttt{B i}
(\cref{ch:lambdatree}). The inner $\forall$ block is not represented at all. To unify $\lambda w\lambda x\lambda y.\,s$ with
$\lambda w\lambda x\lambda y.\,t$, the algorithm simply strips the binders, which is rule \rulename{Abs}, and the
\emph{loose} indices of $s$ and $t$ then play the inner universals: under three binders, $y$ is \texttt{B 0}, $x$
is \texttt{B 1}, $w$ is \texttt{B 2}. Because every existential is quantified outside every inner universal, the
values bound to the \texttt{V}s are closed terms, and \enquote{fresh eigenvariable} costs nothing.

\paragraph{The algorithm} in pseudo-code; $S$ is a triangular substitution, i.e.\ an association list that is
walked rather than applied.
\begin{Verbatim}[fontsize=\footnotesize]
devar(S, s):                          ; β0 by renaming
  if s = (V F) b1..bk and S binds F to λ^k.body
  then devar(S, body with each bound index renamed to the matching b)  else s
unif(S, s, t):  s := devar(S, s);  t := devar(S, t);  case
  A s', A t'  → unif(S, s', t')                          ; Abs
  A s', t     → unif(S, s', (↑t) (B 0))                  ; Eta (and symmetric)
  F ȳ,  G z̄   → flexflex(F, ȳ, G, z̄, S)
  F ȳ,  t     → flexrigid(F, ȳ, t, S)                    ; (and symmetric)
  a s̄,  b t̄   → if a = b and |s̄| = |t̄| then fold unif over the pairs else FAIL
flexrigid(F, ȳ, t, S):
  if F occurs in t through S then FAIL                   ; occurs-check
  u := t with each loose index i replaced by the position
       of (B i) in ȳ counted from the right, or by •     ; inversion
  return proj(S + [F ↦ λ^|ȳ|. u], u)
proj(S, s):  case devar(S, s) of                         ; scope check and pruning
  A s'               → proj(S, s')
  (C c) s̄, (B i) s̄   → fold proj over s̄
  • s̄                → FAIL                              ; out of scope, rigid position
  (V G) b̄            → S + [G ↦ λ^|b̄|. H (positions of b̄ that are not •)],  H new
flexflex(F, ȳ, F, z̄, S):                                ; same head
  if ȳ = z̄ then S else S + [F ↦ λ^|ȳ|. H (positions i with y_i = z_i)]
flexflex(F, ī, G, j̄, S):                                ; different heads
  if ī ⊆ j̄ then S + [G ↦ λj̄. F ī]  elif j̄ ⊆ ī then S + [F ↦ λī. G j̄]
  else S + [F ↦ λī. H (ī ∩ j̄),  G ↦ λj̄. H (ī ∩ j̄)]
\end{Verbatim}

Compare it with the rules of \cref{sec:hounif-algo}.
\begin{itemize}
  \item \textbf{Inversion first, pruning second.} \texttt{flexrigid} binds $F$ to $\lambda^{|\bar y|}.u$ \emph{before}
    pruning, so $u$ may still contain markers \texttt{•} inside the arguments of other flexible terms. That is
    harmless: \texttt{proj} then binds each such $G$ to a term that drops those positions, and since $S$ is
    triangular, every later \texttt{devar} through $F$ also goes through $G$ and never sees a marker.
    (Nipkow's ML encodes the marker as the index \texttt{B(\textasciitilde 10000)}, which shifting keeps
    negative.)
  \item \textbf{One walk.} Imitation is the copying of \texttt{C} heads, projection the renaming of loose
    indices, the scope check the \texttt{•} in rigid position, and pruning the \texttt{V} case: this is the
    book's \enquote{variable elimination}. Inner binders need no set $L$: below \texttt{A}, indices that refer to
    local binders are never shifted into markers.
  \item \textbf{\texttt{devar} is $\beta_0$.} Arguments of a bound existential are indices, so
    \enquote{substituting} only renames indices; no shifting of arbitrary terms is needed.
  \item \textbf{An optimisation in \texttt{flexflex}.} If one argument list is contained in the other, one variable
    is bound directly to the other instead of introducing $H$; up to renaming, this is the mgu of
    \cref{prop:hounif-ff}.
  \item \textbf{No pattern test.} The code assumes its input is a pattern problem; a non-pattern argument would
    simply not be handled correctly. A dynamic-$\Llam$ implementation adds a test there and delays
    (\cref{sec:hounif-dynamic}).
\end{itemize}

\begin{tutorial}[\cref{ex:hounif-8h} in de Bruijn form]
Under the three binders $w,x,y$ the equation is $f\,(U\,1\,0)=f\,(g\,(V\,0\,2))$. \rulename{Dec} on $f$ leaves
$U\,1\,0=g\,(V\,0\,2)$. Here $\bar y=(1,0)$: index~$1$ (that is $x$) sits at position $1$ from the right, index~$0$
($y$) at position $0$, and index~$2$ ($w$) is absent. So $u=g\,(V\,0\,\text{•})$ and $U:=\lambda\lambda.\,g\,(V\,0\,\text{•})$.
\texttt{proj} meets $V\,0\,\text{•}$ and prunes: $V:=\lambda\lambda.\,H\,1$, i.e.\ $\lambda y\lambda w\,(H\,y)$. Walking $U$'s
binding through $S$ gives $\lambda x\lambda y\,(g\,(H\,y))$, the book's answer.
\end{tutorial}

\subsection{An executable transliteration}

The following plain Scheme is a line-by-line transliteration of the de~Bruijn ML code (file
\texttt{lics93-db.ML}). It is \emph{not} a proposal for the new module's interface (that is \cref{ch:design}); it is
here so that you can run, trace and modify the algorithm. Terms are tagged lists, failure escapes through a
continuation, and \scm{(OUT)} is the marker \texttt{•}.

\begin{scheme}[fontsize=\footnotesize]
(define (tag t) (car t))
(define (V n) (list 'V n))  (define (C n) (list 'C n))  (define (B i) (list 'B i))
(define (A t) (list 'A t))  (define ($ s t) (list '$ s t))  (define OUT '(OUT))
(define escape #f)
(define (fail) (escape 'fail))

(define (strip t)                       ; head and arguments of an application spine
  (let loop ((t t) (ts '()))
    (if (eq? (tag t) '$) (loop (cadr t) (cons (caddr t) ts)) (cons t ts))))
(define (apps h ts) (fold (λ (t h) ($ h t)) h ts))  ; foldl op $
(define (abs n t) (if (zero? n) t (abs (- n 1) (A t))))
(define (hnf n h ts) (abs n (apps h ts)))           ; λ^n. h ts

(define (mapbnd f t)                    ; f on the loose indices of t; #f gives OUT
  (let mpb ((d 0) (t t))
    (case (tag t)
      ((B) (let ((i (cadr t)))
             (if (< i d) t (let ((j (f (- i d)))) (if j (B (+ j d)) OUT)))))
      ((A) (A (mpb (+ d 1) (cadr t))))
      (($) ($ (mpb d (cadr t)) (mpb d (caddr t))))
      (else t))))
(define (incr t) (mapbnd add1 t))

(define (red s bs js)                   ; β0: (λ^k s) applied to indices renames them
  (if (and (eq? (tag s) 'A) (pair? bs) (memq (tag (car bs)) '(B OUT)))
      (red (cadr s) (cdr bs) (cons (and (eq? (tag (car bs)) 'B) (cadr (car bs))) js))
      (apps (mapbnd (λ (i) (list-ref js i)) s) bs)))

(define (devar S s)                     ; walk the head through the triangular S
  (let ((ht (strip s)))
    (cond ((and (eq? (tag (car ht)) 'V) (assoc (cadr (car ht)) S))
           => (λ (p) (devar S (red (cdr p) (cdr ht) '()))))
          (else s))))

(define (occ F S t)
  (let occ ((t t))
    (case (tag t)
      ((V) (or (equal? F (cadr t))
               (cond ((assoc (cadr t) S) => (λ (p) (occ (cdr p)))) (else #f))))
      (($) (or (occ (cadr t)) (occ (caddr t))))
      ((A) (occ (cadr t)))
      (else #f))))

(define counter 0)
(define (newV)
  (set! counter (add1 counter)) (V (string->symbol (format "H~a" counter))))

(define (pos p xs)                      ; indices (from the right) of the xs with p
  (let loop ((xs xs))
    (cond ((null? xs) '())
          ((p (car xs)) (cons (B (length (cdr xs))) (loop (cdr xs))))
          (else (loop (cdr xs))))))
(define (idx bs b)                      ; the index of b in the binder list bs, or #f
  (let loop ((bs bs))
    (cond ((null? bs) #f)
          ((equal? (car bs) b) (length (cdr bs)))
          (else (loop (cdr bs))))))

(define (proj S s)                      ; scope check + pruning, one walk over s
  (let* ((ht (strip (devar S s))) (h (car ht)) (ss (cdr ht)))
    (case (tag h)
      ((A) (proj S (cadr h)))
      ((C B) (fold (λ (t S) (proj S t)) S ss))
      ((OUT) (fail))                    ; out-of-scope universal, rigid position
      ((V) (let ((keep (pos (λ (b) (eq? (tag b) 'B)) ss)))
             (cons (cons (cadr h) (hnf (length ss) (newV) keep)) S))))))

(define (flexflex1 F ym zn S)           ; same head: keep the positions that agree
  (if (equal? ym zn) S
      (let ((keep (pos (λ (p) (equal? (car p) (cdr p))) (map cons ym zn))))
        (cons (cons F (hnf (length ym) (newV) keep)) S))))
(define (lam im G jn) (hnf (length im) G (map (λ (j) (B (idx im j))) jn)))
(define (subset? xs ys) (every (λ (x) (member x ys)) xs))
(define (flexflex2 F im G jn S)         ; different heads
  (cond ((subset? im jn) (cons (cons G (lam jn (V F) im)) S))
        ((subset? jn im) (cons (cons F (lam im (V G) jn)) S))
        (else (let ((kl (filter (λ (i) (member i jn)) im)) (H (newV)))
                (cons (cons F (lam im H kl)) (cons (cons G (lam jn H kl)) S))))))
(define (flexflex F ym G zn S)
  (if (equal? F G) (flexflex1 F ym zn S) (flexflex2 F ym G zn S)))

(define (flexrigid F im t S)            ; occurs check, inversion, then pruning
  (if (occ F S t) (fail)
      (let ((u (mapbnd (λ (i) (idx im (B i))) t)))
        (proj (cons (cons F (abs (length im) u)) S) u))))

(define (unif S s t)
  (let ((s (devar S s)) (t (devar S t)))
    (cond ((and (eq? (tag s) 'A) (eq? (tag t) 'A)) (unif S (cadr s) (cadr t)))
          ((eq? (tag s) 'A) (unif S (cadr s) ($ (incr t) (B 0))))     ; η on demand
          ((eq? (tag t) 'A) (unif S ($ (incr s) (B 0)) (cadr t)))
          (else (cases S s t)))))
(define (cases S s t)
  (let ((hs (strip s)) (ht (strip t)))
    (cond ((and (eq? (tag (car hs)) 'V) (eq? (tag (car ht)) 'V))
           (flexflex (cadr (car hs)) (cdr hs) (cadr (car ht)) (cdr ht) S))
          ((eq? (tag (car hs)) 'V) (flexrigid (cadr (car hs)) (cdr hs) t S))
          ((eq? (tag (car ht)) 'V) (flexrigid (cadr (car ht)) (cdr ht) s S))
          ((and (equal? (car hs) (car ht)) (= (length (cdr hs)) (length (cdr ht))))
           (fold (λ (a b S) (unif S a b)) S (cdr hs) (cdr ht)))
          (else (fail)))))
\end{scheme}

The test driver, a printer that applies $S$ and turns indices back into names (omitted here), and the
problems are in \texttt{writers/ch07-hounif/nipkow.scm}. For instance, \cref{ex:hounif-8h} is posed as
\begin{scheme}
(define f (C 'f)) (define g (C 'g)) (define h (C 'h))
(define (lam3 t) (A (A (A t))))     ; ∀w ∀x ∀y as three binders: y = 0, x = 1, w = 2
;; Example 8.H: ∀f∀g∃U∃V∀w∀x∀y [(f (U x y)) = (f (g (V y w)))]
(run "8.H  f (U x y) = f (g (V y w))" '(U V)
     (lam3 ($ f ($ ($ (V 'U) (B 1)) (B 0))))
     (lam3 ($ f ($ g ($ ($ (V 'V) (B 0)) (B 2))))))
\end{scheme}
and the whole battery prints, with \texttt{H1}, \texttt{H2} the fresh variables and \texttt{V1} standing for the
book's $V'$:
\begin{answers}
8.H  f (U x y) = f (g (V y w))
  U := λx1. λx2. g (H1 x2)
  V := λx1. λx2. H1 x1
8.I  U x y = g (U y w)
  no unifier
8.J  U x y = g w
  no unifier
4.5  λx λy (G x) = λx λy (H y)
  G := λx1. H1
  H := λx1. H1
ff   λx λy (F x y) = λx λy (F y x)
  F := λx1. λx2. H1
8.K  U x = f (V1 x y z)
  U := λx1. f (H1 x1)
  V1 := λx1. λx2. λx3. H1 x1
eta  λx (F x) = g
  F := λx1. g x1
dup  λx (F x) = λx (g x x)
  F := λx1. g x1 x1
loc  F x = h (λz (G z w))
  F := λx1. h (λx2. H1 x2)
  G := λx1. λx2. H1 x1
two  F x = g (G x w) (G w x)
  F := λx1. g H2 H2
  G := λx1. λx2. H2
\end{answers}
Every line agrees with the hand derivations: \cref{ex:hounif-8h,ex:hounif-8i,ex:hounif-8j,ex:hounif-ffdiff,ex:hounif-loc,ex:hounif-two},
the same-head pitfall, and the eagerly raised form of \cref{ex:hounif-8k} below.

\begin{kanren}[what first-order unification does with the same terms]
\micro{} can already store de~Bruijn terms as records, but its unifier treats an application of a logic
variable as a pair of slots. This sometimes gives the right answer by luck, sometimes an ill-scoped one, and
sometimes a wrong failure. With \scm{(define-record lam body)}, \scm{(define-record app rator rand)} and
\scm{(define-record dbi n)}:
\begin{scheme}
(show "X = (f X)"
      (°->list #f (fresh° (q) (=° q (list 'f q)))))
(show "λx (F x) = λx (g x)"
      (°->list #f (fresh° (f) (=° (make-lam (make-app f (make-dbi 0)))
                                  (make-lam (make-app 'g (make-dbi 0)))))))
(show "λx (F x) = λx (g x x)"
      (°->list #f (fresh° (f) (=° (make-lam (make-app f (make-dbi 0)))
                                  (make-lam (make-app (make-app 'g (make-dbi 0)) (make-dbi 0)))))))
(show "λx (F x) = λx (g x 7)"
      (°->list #f (fresh° (f) (=° (make-lam (make-app f (make-dbi 0)))
                                  (make-lam (make-app (make-app 'g (make-dbi 0)) 7))))))
\end{scheme}
\begin{answers}
;; X = (f X)
()
;; λx (F x) = λx (g x)
((λ () 'g))
;; λx (F x) = λx (g x x)
((λ () (make-record-instance 'app 'g (make-record-instance 'dbi 0))))
;; λx (F x) = λx (g x 7)
()
\end{answers}
The occurs-check case is already right. The second answer is right up to $\eta$. The third binds $F$ to
\enquote{$g$ applied to index $0$} \emph{outside} any binder: a dangling index, where the mgu is
$F:=\lambda y\,(g\,y\,y)$ (\texttt{dup} above). The fourth fails, where the mgu is $F:=\lambda y\,(g\,y\,7)$. A
higher-order extension must therefore recognise a flexible application and run \rulename{Elim} on it rather than
decomposing it structurally; see \cref{ch:design}.
\end{kanren}

% ============================================================================================
\section{Mixed prefixes, levels and delayed equations}\label{sec:hounif-pragmatics}

Everything so far assumed a $\forall\exists\forall$ prefix. Real proof search does not deliver one. Each
\lp{pi x\ G} adds a universal, and each backchaining step adds existentials \emph{after} it, giving prefixes such as
$\forall\exists\forall\exists\forall\dots$ \book{224}. Raising them eagerly \enquote{has to be performed dynamically},
may raise existentials \enquote{over long lists of universally quantified variables} (typical of the \lp{pi}-based
recursion of \cref{ch:lambdatree}), and \enquote{often can be redundant and may have to be undone by a pruning
substitution at a later point} \book{224}.

\begin{example}[Eager versus delayed raising {\book{224--225}}]\label{ex:hounif-8k}
Consider
\[\forall f\,\exists U\,\forall x\,\forall y\,\forall z\,\exists V\,[(U\,x)=(f\,V)].\]
\begin{itemize}
  \item \emph{Eagerly}, raise $V$ over $x,y,z$:
    \[\forall f\,\exists U\,\exists V'\,\forall x\,\forall y\,\forall z\,[(U\,x)=(f\,(V'\,x\,y\,z))].\]
    Variable elimination
    immediately prunes $V':=\lambda x\lambda y\lambda z\,(W\,x)$, giving $U:=\lambda x\,(f\,(W\,x))$, so two of the three raised
    arguments were wasted work.
  \item \emph{Delayed}: record that $U$ may be raised over $f$ and $V$ over $x,y,z,f$. When solving, \enquote{noting
    the presence of $x$ as an argument of the flexible term and the possibility of raising $V$ over $x$, we would
    replace $V$ with the term $(V'\,x)$, where $V'$ is a new variable quantified at the same place in the prefix as
    $U$ and would produce the unifying substitution $\lambda x\,(f\,(V'\,x))$ for $U$} \book{225}.
\end{itemize}
\end{example}

The lists of universals that each existential may depend on \enquote{are all initial segments of the sequence of
variables that are universally quantified in the prefix, \dots{} and this information can be maintained and used
efficiently by attaching numerical indices to the quantified variables} \book{225}. This idea has two
literature sources.

\paragraph{Universe indices for \texttt{pi} goals~\cite{NadathurJayaramanKwon1995}.} In a WAM-style interpreter for
first-order hereditary Harrop formulas, \enquote{each time a universal quantifier is encountered, a new constant must
be introduced, giving rise to the next universe in the hierarchy}: the universe index is incremented and the new
constant is tagged with it. A logic variable is tagged with the current index, meaning \enquote{that a term may be
substituted for the variable only if all the constants appearing in the term have a smaller or identical tag}
(p.~16). Binding $X^i:=t$ then requires (1)~the usual occurs-check; (2)~that $t$ contains no constant with tag
greater than $i$; and (3)~that, before binding, the tags of the variables of $t$ greater than $i$ be lowered to
$i$, \enquote{to prevent a later instantiation of these variables from violating the restrictions on instantiations
for $X$} (p.~16).

\begin{example}[Why tags must be lowered {\cite[p.~16]{NadathurJayaramanKwon1995}}]
Program $\forall z\,(q(z)\imp p(d(z)))$, goal $\exists x\,\forall y\,(q(y)\imp p(x))$. After the quantifiers the goal
is $q(c^2)\imp p(X^1)$. Backchaining on the clause with a new $Z^2$ unifies $p(d(Z^2))$ with $p(X^1)$; binding
$X:=d(Z)$ lowers $Z$ to $Z^1$. The remaining goal $q(Z^1)$ must be matched against the hypothesis $q(c^2)$ and
\emph{fails}, because $c$ has tag $2>1$. Without step~(3), $Z$ would be bound to $c$ and $X$ would end up as
$d(c)$, mentioning a constant introduced after $X$.
\end{example}

\paragraph{Prefix erasure for patterns~\cite{NadathurLinnell2005}.} Nadathur and Linnell tag every variable $y$ with
$l(y)$, \enquote{the number of times an existential quantifier is immediately followed by a universal one in the prefix
up to and including the quantifier binding that occurrence}. Then \enquote{a universal variable $u$ is quantified
within the scope of the quantifier for an existential variable $x$ just in case $l(x)<l(u)$} (p.~8), so the pattern
test becomes local. Raising is restricted to what is needed: for $f(\bar y)=g(\bar z)$ with $l(f)\le l(g)$, $g$ needs
to be raised only over $\bar y{\Uparrow}g=\{u\in\bar y\mid l(u)\le l(g)\}$, the universals among $f$'s arguments that
$g$ may depend on. A combined rule (9$'$) binds $f$ and $g$ to a common new $h$ with $l(h)=l(f)$, applied first to
those universals and then to the common arguments (p.~9). Their Theorem~2 states that prefix erasure preserves
termination and mgus (p.~10), and their Figure~2 gives the single-traversal \texttt{mksubst}/\texttt{bnd} that
computes the binding, the pruning and the raising together (p.~11). The algorithm is used in Bedwyr, Abella and
Teyjus~2 \book{228}.

\paragraph{Variable elimination with levels.} Combining the two, \cref{def:hounif-elim} generalises to tagged
variables. This is our reconstruction, checked by hand on \cref{ex:hounif-8h,ex:hounif-8k}; Nadathur and Linnell's
Figure~2 is the authoritative version. Write $l(\cdot)$ for levels, with signature constants at the lowest level.
To eliminate $F$ in $(F\,\bar c)=t$, where pattern means that the $c_i$ are distinct universals with
$l(c_i)>l(F)$:
\begin{itemize}
  \item a universal $u$ in rigid position is kept if $u\in\bar c\cup L$ or $l(u)\le l(F)$ (since $F$ may mention it
    directly), and fails otherwise;
  \item a flexible pattern $(G\,\bar d)$ keeps the $d_j\in\bar c\cup L$ and the $d_j$ with $l(d_j)\le l(F)$, prunes the rest,
    and, if $l(G)>l(F)$, is \emph{also raised} over the $c_i$ with $l(c_i)\le l(G)$ that it does not already receive;
    the new $G'$ gets level $l(F)$. In \cref{ex:hounif-8k}, $l(U)=0$ and $l(x)=l(V)=1$, so $V$ (no arguments) is
    raised over $x$: $V:=(V'\,x)$ with $l(V')=0$, and $U:=\lambda x\,(f\,(V'\,x))$, as in the book;
  \item when no $c_i$ qualifies, raising degenerates to \emph{lowering}: $G:=\lambda\bar d.\,(G'\,\bar e)$ with
    $l(G')=l(F)$ and nothing added, which is step~(3) of the universe-index discipline above.
\end{itemize}

\paragraph{De Bruijn levels in ELPI~\cite{DunchevEtAl2015}.} ELPI makes the two notions of level coincide.
\enquote{When a clause's head is unified with the goal in a context of length $n$, the universally quantified variables
of the clause are instantiated to unification variables $X^n$ where the level $n$ records that $X$ has only
visibility of the initial prefix of length $n$ of the context. If later a fresh constant is added by the
\texttt{pi} rule, the constant occupies position $n+1$ \dots{} and it will not be allowed to occur in instances of the
variable $X^n$} (p.~4). Lowering a level is done by pruning: \enquote{unifying $B_0$ with $(\mathit{arr}\ C_1\ D_1)$
prunes $C$ and $D$ to level $0$} (p.~7), i.e.\ the higher-level variable is bound to a fresh lower-level one.

\begin{kanren}[levels and \micro]
\micro{} already allocates variables from a monotone counter, \scm{vars-count}, so creation order is a
linearisation of the dynamically built prefix. If eigenvariables were allocated from the same counter, the
creation index would give every variable its initial level for free. Three facts of the current engine matter for
any level discipline.
\begin{enumerate}[nosep]
  \item \emph{Level is not index.} Pruning and raising create a \emph{new} variable \enquote{quantified at the same
    place} as an \emph{old} one \book{225}, so a variable's level must be stored separately from its index.
  \item \emph{Orientation of var-var bindings.} Today the older variable is bound to the newer one. With scopes,
    binding a low-level $X$ to a high-level $Y$ lets a later binding of $Y$ smuggle a young eigenvariable into $X$.
    Either the orientation must follow levels, or the binding must lower $Y$, as in step~(3) above.
  \item \emph{Unification allocates.} Pruning, raising and the flexible-flexible $H$ create variables, so the
    unifier must return an updated \scm{vars-count}, and the bindings it reports must include the prunings of
    \emph{other} variables (item~3 of the pitfall after \cref{def:hounif-elim}), because constraint re-verification keys on them.
\end{enumerate}
How these requirements are met is the subject of \cref{ch:design}.
\end{kanren}

\subsection{Dynamic \texorpdfstring{$\Llam$}{Lλ}: delay what is not a pattern}\label{sec:hounif-dynamic}

A static $\Llam$ restriction \enquote{can rule out some useful programming idioms} \book{225}, the standard one
being substitution by meta-level $\beta$-reduction:
\begin{lprolog}
type   subst   (tm -> tm) -> tm -> tm -> o.
subst R N (R N).
\end{lprolog}
Here \lp{R} and \lp{N} are essentially existential, so \lp{(R N)} is not a pattern. \enquote{A common practical
resolution to this dilemma is to not impose the $\Llam$ restriction on program clauses statically but rather to
expect the condition to be satisfied by the unification problems that are encountered dynamically} \book{225--226}:
when \lp{subst} is called, \lp{R} and \lp{N} are expected to be instantiated, and \lp{(R N)} is then just a $\beta$-redex
to reduce. More generally: \enquote{we would treat any unification problem as if it is a higher-order
pattern-unification one. In the course of such a treatment, we might discover that a particular equation does not
respect the $\Llam$ condition. In this case, we might defer further consideration of this equation until such time
that substitutions resulting from solving other parts of the problem have altered its status} \book{226}.

The literature agrees on this design.
\begin{itemize}
  \item Dowek, Hardin, Kirchner and Pfenning~\cite{DowekHardinKirchnerPfenning1996}, \S5.4 (p.~13): equations outside
    the pattern fragment \enquote{should be postponed as constraints in the hope that they will be instantiated
    further}, which \enquote{avoids all branching during constraint simplification}. Their motivation is
    empirical~\cite{MichaylovPfenning1992}: the static restriction is too severe, but \enquote{almost all dynamically
    arising equations fall within the pattern fragment} (p.~2). Their algorithm phrases \rulename{Elim} as the
    inversion of a pattern substitution: $X[s]=b$ is solved by $X=b[s^{-1}]$ when every index of $b$ occurs in $s$
    (Prop.~8), with pruning computed together with the inverse (p.~11--12).
  \item Nadathur and Linnell~\cite{NadathurLinnell2005}, \S6: a non-pattern subterm met inside \texttt{bnd} can be
    replaced by a fresh flexible term, with the equation between them \emph{deferred}.
  \item Systems: Teyjus~1~\cite{NadathurMitchell1999} implemented (pre-)unification for unrestricted terms, while
    Teyjus~2 \enquote{only implements $\Llam$ unification}; Twelf and Minlog also implement only $\Llam$ unification,
    over an unrestricted term language \book{228}. ELPI's default for a problem outside the fragment is an error,
    with a deprecated option to delay it \enquote{for Teyjus compatibility}; its supported mechanism for suspension
    is instead a general syntactic constraint on goals, resumed when a variable is assigned.
\end{itemize}
Going the other way, \lp{subst} can be written inside $\Llam$ by making the object signature explicit with
\lp{copy} clauses (\cref{ch:lambdatree}), and by extension \enquote{it also may be possible to describe general
higher-order unification as a logic program in the $\Llam$ language} \book{226}~\cite{Miller1991}.

\begin{pitfall}[unifier versus solution; delayed is not solved]
A state with suspended equations has succeeded only \emph{conditionally}: its answer is a pre-unifier together
with residual constraints, and those must be printed with it. The canonical substitution of
\cref{prop:hounif-canonical} would make flexible-flexible residues \enquote{true}, but it introduces new existentials
$H^\tau$: it is a unifier, not a solution, and in a typed setting with an empty type there may be no closed solution
at all \book{211}. Untyped Scheme terms avoid that last problem~\cite[\S9.3]{Miller1991}, but not the first.
\end{pitfall}

\begin{kanren}[what delaying asks of a miniKanren]
Following the book's \enquote{defer further consideration} \book{226}, an \micro{}-style engine needs: a
place in the state for suspended equations; a wake-up step after each successful unification that retries the
suspended equations mentioning a variable just bound (the same \enquote{touched variables} mechanism that already
re-verifies disequalities); and reification that prints what is still suspended. Soft-cut, projection and
disequality must be audited, because each of them can observe a state whose success is only conditional.
Disequality needs care even without delays: the mgu of a pattern problem may contain \emph{new} variables
(pruning, the flexible-flexible $H$), and in $\neg\exists H.(X=\lambda x\,H)$ these are existentially quantified
inside the negation (equivalently, universally outside it), which a store of plain bindings would get wrong
(\cref{def:module-trial}). The concrete design is in \cref{ch:design}.
\end{kanren}

% ============================================================================================
\section{Summary}\label{sec:hounif-summary}

\begin{itemize}
  \item A unification problem is a quantified formula; its prefix decides who may depend on whom. The book
    normalises to a $\forall\exists\forall$ prefix by raising, to $\beta$-normal terms, to equal binders by $\eta$, and turns
    top binders into inner universals \book{212}.
  \item General higher-order unification has no mgus, may need infinite complete sets of unifiers, and is
    undecidable \book{213--214}.
  \item Huet's procedure: \rulename{Dec}/\rulename{Clash} on rigid-rigid pairs, imitation (scope-sensitive) or
    projection (type-sensitive) on flexible-rigid ones, success when only flexible-flexible equations remain
    (pre-unification; the canonical substitution unifies them). Its matching tree is complete but may have infinite
    branches, with or without answers \book{214--220}.
  \item In $\Llam$, flexible terms are patterns $F\,\bar c$ with distinct universals bound inside $F$'s scope. Only
    $\beta_0$ arises, at most one imitation or projection is useful, flexible-flexible pairs have explicit mgus
    (common arguments for different heads, agreeing positions for the same head), and with an occurs-check the
    procedure terminates: pattern unification is decidable and unitary \book{220--224}, \cite{Miller1991}.
  \item Variable elimination in one walk: keep rigid heads, map each $c_i$ to $F$'s $i$-th bound variable, keep local
    binders, prune the arguments of other flexible terms that $F$ cannot see, and fail on $F$ itself (occurs) or on an
    unseeable universal in rigid position (scope). Then $F:=\lambda\bar c.\,t'$ \book{223}.
  \item Nipkow's de Bruijn version is short (our Scheme transliteration of its core is under a hundred lines); loose indices are the inner universals, \texttt{devar}
    is $\beta_0$ by renaming, and inversion uses a marker that pruning removes \cite{Nipkow1993}. No types are needed
    \book{224}.
  \item In practice prefixes are mixed and dynamic: replace raising by levels on variables and eigenvariables,
    lower levels when binding, raise lazily only over the arguments needed \book{225},
    \cite{NadathurJayaramanKwon1995,NadathurLinnell2005,DunchevEtAl2015}.
  \item Dynamic $\Llam$: treat every equation as a pattern problem and \emph{delay} the non-patterns until bindings
    make them patterns; report what remains as residual constraints \book{226},
    \cite{DowekHardinKirchnerPfenning1996}.
  \item For \micro{}: first-order unification is the $n=0$ case; the extension needs eigenvariables, binders,
    flexible applications, levels, allocation inside unification, and a store of delayed equations; see
    \cref{ch:design}.
\end{itemize}

\chapter{Applications: provers, functional programs, processes}\label{ch:apps}

The previous chapters assembled three tools: $\lambda$-tree syntax, which represents object-level binders
by meta-level abstractions (\cref{ch:lambdatree}); the universal and implicational goals of hereditary
Harrop formulas, which give fresh constants and scoped hypotheses (\cref{ch:hh}); and the unification of
$\lambda$-terms, decidable and most-general in the pattern fragment $\Llam$ (\cref{ch:hounif}). Chapters
9--11 of the book~\cite{MillerNadathur2012} put the three to work on three families of meta-programs:
\emph{proof systems} (intuitionistic and classical provers, proof checkers, tactics),
\emph{functional programs} (type inference, evaluators, partial evaluation, CPS conversion) and
\emph{processes} (the $\pi$-calculus, its transitions, traces, and an encoding of the $\lambda$-calculus).

All three follow the same recipe, and it is the main idea of this chapter:

\begin{enumerate}[nosep]
  \item a judgment of the object system becomes a predicate, and each inference rule becomes a clause
        (the book calls this a ``reflection of provability relations directly into predicates'', \book{240});
  \item object-level binders are meta-level $\lambda$s, so object-level substitution is $\beta$-reduction;
  \item to go under a binder, the program introduces a fresh constant with \lp{pi x\ ...} and, usually, a
        hypothetical fact about it with \lp{=>}.
\end{enumerate}

\noindent The side conditions that clutter the usual presentations (``$y$ fresh'', ``$y$ not free in the
conclusion'', ``$w\notin fn(P)$'') then disappear: the nesting of quantifiers already enforces them. For
each application we ask three questions, since they decide what a miniKanren needs: which binding
facilities does it use ($\forall$, $\imp$, pattern unification, more than patterns)? Which part runs on
today's \micro{} with a first-order encoding? What changes when depth-first search is replaced by
miniKanren's interleaving search? Four small ports below were actually run on \micro{}. Their outputs
are reproduced verbatim.

% =====================================================================================================
\section{Proof systems as logic programs}\label{sec:apps-provers}

\subsection{Intuitionistic propositional logic: from a loop to a decision procedure}\label{sec:apps-ipl}

Formulas are built from two constants and three connectives \book{229--230}:

\begin{lprolog}
kind  term, form    type.             % types for terms and formulas
type  ff,                             % encoding the false proposition
      tt            form.             % encoding the true proposition
type  &&,                             % encoding conjunction
      !!,                             % encoding disjunction
      ==>           form -> form -> form.   % encoding implication
infixl &&  5.
infixl !!  4.
infixr ==> 3.
\end{lprolog}

\noindent A sequent $\Gamma\longrightarrow G$ has a multiset $\Gamma$ of hypotheses and one conclusion.
The book's starting point is the following calculus \book{230}. (In the printed figure the axiom for
$\top$ is mislabelled $\bot L$.)

\begin{mathpar}
\inferrule{ }{\Gamma, A\longrightarrow A}\;\rulename{init}\ (A\text{ atomic})
\and \inferrule{ }{\Gamma,\bot\longrightarrow G}\;\bot L
\and \inferrule{ }{\Gamma\longrightarrow\top}\;\top R
\and \inferrule{\Gamma,A,B\longrightarrow G}{\Gamma,A\land B\longrightarrow G}\;\land L
\and \inferrule{\Gamma\longrightarrow A \\ \Gamma\longrightarrow B}{\Gamma\longrightarrow A\land B}\;\land R
\and \inferrule{\Gamma,A\longrightarrow G \\ \Gamma,B\longrightarrow G}{\Gamma,A\lor B\longrightarrow G}\;\lor L
\and \inferrule{\Gamma\longrightarrow A_i}{\Gamma\longrightarrow A_1\lor A_2}\;\lor R_i
\and \inferrule{\Gamma,A\imp B\longrightarrow A \\ \Gamma,B\longrightarrow G}{\Gamma,A\imp B\longrightarrow G}\;\imp L
\and \inferrule{\Gamma,A\longrightarrow B}{\Gamma\longrightarrow A\imp B}\;\imp R
\end{mathpar}

\paragraph{Why the direct translation loops.}
Each rule, read from conclusion to premises, removes a connective, except $\imp L$, whose left premise
still contains the principal formula $A\imp B$. The instance
\[
  \inferrule{p\imp q\longrightarrow q \\ q\longrightarrow q}{p\imp q\longrightarrow q}\;\imp L
\]
has a left premise identical to its conclusion. A depth-first interpreter that tries this premise first
recurses forever. The book: ``Depth-first search usually steers the interpreter into endless loops when
trying to work backwards from a given sequent'' \book{230}. Soundness of such direct translations is
easy. Completeness of their \emph{execution} is not.

\paragraph{The cure: a better calculus, not a better interpreter.}
Replace $\imp L$ by six rules that look at the shape of the antecedent~\cite{Dyckhoff1992,Hudelmaier1992}
(\book{231}). (The text says ``four rules''. The figure shows six.)

\begin{mathpar}
\inferrule{\Gamma,A,B\longrightarrow G}{\Gamma,A,A\imp B\longrightarrow G}\;{\imp}L_1\ (A\text{ atomic})
\and \inferrule{\Gamma,C\imp D\imp B\longrightarrow G}{\Gamma,(C\land D)\imp B\longrightarrow G}\;{\imp}L_2
\and \inferrule{\Gamma,C\imp B,D\imp B\longrightarrow G}{\Gamma,(C\lor D)\imp B\longrightarrow G}\;{\imp}L_3
\and \inferrule{\Gamma,D\imp B\longrightarrow C\imp D \\ \Gamma,B\longrightarrow G}{\Gamma,(C\imp D)\imp B\longrightarrow G}\;{\imp}L_4
\and \inferrule{\Gamma\longrightarrow G}{\Gamma,\bot\imp B\longrightarrow G}\;{\imp}L_5
\and \inferrule{\Gamma,B\longrightarrow G}{\Gamma,\top\imp B\longrightarrow G}\;{\imp}L_6
\end{mathpar}

\begin{theorem}[Hudelmaier, Dyckhoff; \book{230--231}]\label{thm:apps-dyckhoff}
The calculus obtained by replacing $\imp L$ with ${\imp}L_1$--${\imp}L_6$ is sound and complete for this
fragment of propositional intuitionistic logic, and every sequence of sequents produced by backward proof
search is finite. A weight can be assigned to connectives and sequents so that every conclusion weighs
more than each of its premises. Consequently the direct translation of the calculus into $\lambda$Prolog
``actually constitutes a decision procedure'' \book{231}.
\end{theorem}

\noindent The book does not give the weight function. Provability in this logic is PSPACE-complete
(Statman, cited on \book{245}). Here is the program \book{232}:

\begin{lprolog}
type  atom     form -> o.
type  seq      list form -> form -> o.

seq Gamma A :- atom A, memb_and_rest A Gamma _.
seq Gamma tt.
seq Gamma (A && B)  :- seq Gamma A, seq Gamma B.
seq Gamma (A !! B)  &
seq Gamma (B !! A)  :- seq Gamma A.
seq Gamma (A ==> B) :- seq (A::Gamma) B.
seq Gamma _ :- memb_and_rest ff Gamma _.
seq Gamma G :- memb_and_rest (A && B) Gamma Gamma',
               seq (A::B::Gamma') G.
seq Gamma G :- memb_and_rest (A !! B) Gamma Gamma',
               (seq (A::Gamma') G; seq (B::Gamma') G).
seq Gamma G :- memb_and_rest (A ==> B) Gamma Gamma',
  ( atom A, memb_and_rest A Gamma' _, seq (B::Gamma') G;
    A = (C && D),  seq ((C ==> D ==> B)::Gamma') G;
    A = (C !! D),  seq ((C ==> B)::(D ==> B)::Gamma') G;
    A = (C ==> D), seq ((D ==> B)::Gamma') A,
                   seq (B::Gamma') G;
    A = ff,        seq Gamma' G;
    A = tt,        seq (B::Gamma') G  ).
\end{lprolog}

\noindent The relation \lp{memb_and_rest X L R} (from \book{58}) holds when \lp{L} contains \lp{X} and
\lp{R} is \lp{L} with that occurrence removed. It plays the role of multiset selection. The head
\lp{H1 & H2 :- B} is two clauses sharing a body, so one line gives both $\lor R$ rules. The last clause is
${\imp}L_1$--${\imp}L_6$ written as a single six-way disjunction on the shape of \lp{A}. The predicate
\lp{atom} is deliberately left open: ``This predicate can be defined simply by adding clauses to the
program that enumerate all the propositional letters'' \book{231}. The book proves the formula
$a\imp(a\imp b)\imp(a\imp b\imp c)\imp c$ by introducing the letters locally \book{231} (the printed formula
has an unbalanced parenthesis):

\begin{lprolog}
?- pi a\ pi b\ pi c\ atom a => atom b => atom c =>
              seq nil (a ==> (a ==> b) ==> (a ==> b ==> c) ==> c).
\end{lprolog}

\noindent The query ``achieves the same effect by introducing three new constants and assuming that each
of them satisfies the atom predicate'' \book{231}. It succeeds. It is a compact showcase of \lp{pi} (fresh
constants) together with \lp{=>} (local clauses).

\begin{tutorial}[the proof the program finds]
Three $\imp R$ steps give the context $[a\imp b\imp c,\ a\imp b,\ a]$ (most recent first, since $\imp R$
conses) and goal $c$. The atom $c$ is not in the context. \lp{memb_and_rest (A ==> B)} first selects
$a\imp(b\imp c)$. Since $a$ is present, ${\imp}L_1$ replaces it by $b\imp c$. On this implication
${\imp}L_1$ fails because $b$ is not yet present, so the search backtracks to $a\imp b$, and ${\imp}L_1$
adds $b$. Now ${\imp}L_1$ on $b\imp c$ adds $c$, and \rulename{init} closes the branch. No rule ever puts
back the formula it consumed. This is the decreasing weight of \cref{thm:apps-dyckhoff} at work.
\end{tutorial}

\begin{pitfall}[a typo that makes the prover unsound]
The $\lor L$ clause printed in the figure reads \lp{(seq (A::Gamma') G; seq (B::Gamma') G)}. It succeeds if
\emph{either} branch succeeds, but $\lor L$ needs \emph{both} premises. With \lp{;} the sequent
$a\lor b\longrightarrow a$ is ``proved''. The body must be the conjunction
\lp{seq (A::Gamma') G, seq (B::Gamma') G}. The page image confirms the semicolon is really there. The run
below shows both versions.
\end{pitfall}

\begin{kanren}[the decision procedure on today's (aux kanren micro)]
Nothing in this program needs binders, so it ports to \micro{} directly. Formulas become tagged lists,
atoms become Scheme symbols, and \lp{atom} becomes the type constraint \scm{symbol°}. The constraint does
not enumerate letters, and it stays pending if its argument is unbound. Because \scm{tt} and \scm{ff} are
symbols too, they are excluded with \scm{≠°}. The parameter \scm{both} chooses how the two $\lor L$
premises combine: \scm{andª} gives the correct rule, \scm{orª} the printed typo.
\begin{scheme}[fontsize=\footnotesize]
(define-relation (atom° a) (symbol° a) (≠° a 'tt) (≠° a 'ff))

(define-relation (memb-and-rest° x l rest)
  (cond° ((fresh° (d) (=° l (cons x d)) (=° rest d)))
         ((fresh° (y d r) (=° l (cons y d)) (=° rest (cons y r)) (memb-and-rest° x d r)))))

;; Figure 9.3, with the ∨L clause written as a conjunction; `both` combines the two ∨L premises
(define ((seq°/with both) Γ G)
  (define seq° (seq°/with both))
  (λ (s)
    (δ ((cond°
          ((atom° G) (fresh° (r) (memb-and-rest° G Γ r)))
          ((=° G 'tt))
          ((fresh° (A B) (=° G `(&& ,A ,B)) (seq° Γ A) (seq° Γ B)))
          ((fresh° (A B) (=° G `(!! ,A ,B)) (or° (seq° Γ A) (seq° Γ B))))
          ((fresh° (A B) (=° G `(==> ,A ,B)) (seq° `(,A . ,Γ) B)))
          ((fresh° (r) (memb-and-rest° 'ff Γ r)))
          ((fresh° (A B r) (memb-and-rest° `(&& ,A ,B) Γ r) (seq° `(,A ,B . ,r) G)))
          ((fresh° (A B r) (memb-and-rest° `(!! ,A ,B) Γ r)
             (both (seq° `(,A . ,r) G) (seq° `(,B . ,r) G))))
          ((fresh° (A B r) (memb-and-rest° `(==> ,A ,B) Γ r)
             (cond°
               ((atom° A) (fresh° (r*) (memb-and-rest° A r r*)) (seq° `(,B . ,r) G))
               ((fresh° (C D) (=° A `(&& ,C ,D)) (seq° `((==> ,C (==> ,D ,B)) . ,r) G)))
               ((fresh° (C D) (=° A `(!! ,C ,D)) (seq° `((==> ,C ,B) (==> ,D ,B) . ,r) G)))
               ((fresh° (C D) (=° A `(==> ,C ,D)) (seq° `((==> ,D ,B) . ,r) A) (seq° `(,B . ,r) G)))
               ((=° A 'ff) (seq° r G))
               ((=° A 'tt) (seq° `(,B . ,r) G))))))
        s))))

(define seq° (seq°/with andª))          ; the correct ∨L: both premises
(define seq-typo° (seq°/with orª))      ; the printed ∨L: either premise

(define (proofs g) (length (μkanren-run (q -1 #t) g (=° q 'yes))))

(define f1 '(==> a (==> (==> a b) (==> (==> a (==> b c)) c))))   ; the formula of p.247
(define em '(!! a (==> a ff)))                                    ; excluded middle
(define peirce '(==> (==> (==> a b) a) a))                        ; Peirce's formula

(pretty-print
  (list (list 'f1      (μkanren-run (q 1 #t) (seq° '() f1) (=° q 'yes)) (proofs (seq° '() f1)))
        (list 'or-left (proofs (seq° '((!! a b)) 'a)) (proofs (seq-typo° '((!! a b)) 'a)))
        (list 'or-comm (proofs (seq° '() '(==> (!! a b) (!! b a)))))
        (list 'em      (proofs (seq° '() em)))
        (list 'peirce  (proofs (seq° '() peirce)))))
\end{scheme}
\begin{answers}
((f1 (yes) 3) (or-left 0 1) (or-comm 1) (em 0) (peirce 0))
\end{answers}
The body of \scm{seq°/with} is wrapped in \scm{(λ (s) (δ ...))}, exactly as \scm{define-relation} expands,
so that building the recursive goal does not loop. Every \scm{run*} above terminates, as
\cref{thm:apps-dyckhoff} promises for a ground query. For the formula of \book{231} the search finds three derivations.
$a\lor b\longrightarrow a$ has none with the correct rule and one with the typo. Excluded middle and
Peirce's formula, which are classically valid but not intuitionistically (\cref{ch:hh}), get the answer
\emph{no}.
\end{kanren}

\begin{kanren}[what interleaving changes for the looping calculus]
Interleaving does not make the naive calculus terminate, but it does change which queries succeed.
Take only the atomic, $\imp R$ and the original $\imp L$ rule:
\begin{scheme}[fontsize=\footnotesize]
(define-relation (naive° Γ G)
  (cond°
    ((symbol° G) (fresh° (r) (memb-and-rest° G Γ r)))
    ((fresh° (A B) (=° G `(==> ,A ,B)) (naive° `(,A . ,Γ) B)))
    ((fresh° (A B r) (memb-and-rest° `(==> ,A ,B) Γ r) (naive° Γ A) (naive° `(,B . ,r) G)))))

(pretty-print (μkanren-run (q 1 #t) (naive° '((==> p q) (==> q r) q) 'r) (=° q 'yes)))
(flush-output)
(pretty-print (μkanren-run (q 1 #t) (naive° '((==> p q)) 'q) (=° q 'yes)))
\end{scheme}
\begin{answers}
(yes)
\end{answers}
The first query is a theorem, and \scm{run 1} answers immediately. A depth-first reading of the same clauses,
in clause order, would select $p\imp q$ first and then loop on its left premise
$\ldots\longrightarrow p$, which reproduces itself. Interleaving gives the branch through $q\imp r$ its
turn. The second query is the book's looping instance $p\imp q\longrightarrow q$, which is not a theorem.
Here the stream is infinite and empty: the run was killed after 30 seconds with no output. In short,
interleaving makes the search \emph{complete for answers}, while termination still has to come from the
calculus.
\end{kanren}

% -----------------------------------------------------------------------------------------------------
\subsection{Natural deduction with proof terms}\label{sec:apps-nd}

For first-order intuitionistic logic, add the quantifiers \lp{type all, some (term -> form) -> form.} of
\cref{ch:lambdatree}. Natural deduction has \emph{hypothetical} premises (in ${\imp}I$, ${\lor}E$, ${\exists}E$
a derivation may start from a discharged assumption) and \emph{eigenvariable} conditions (for ${\forall}I$
and ${\exists}E$, ``This variable must not be free in the conclusion of the rule or in any assumption that
remains undischarged'', \book{232}). Four representative rules out of the fourteen \book{233}:

\begin{mathpar}
\inferrule{[A] \\\\ \vdots \\\\ B}{A\imp B}\;{\imp}I
\and \inferrule{A \\ A\imp B}{B}\;{\imp}E
\and \inferrule{A[y/x]}{\forall x\,A}\;{\forall}I\ (y\text{ fresh})
\and \inferrule{\forall x\,A}{A[t/x]}\;{\forall}E
\end{mathpar}

\noindent The book records whole derivations as $\lambda$-terms of a type \lp{proof}, with one
constructor per rule \book{233}. An excerpt of the fourteen declarations:

\begin{lprolog}
type imp_i           (proof -> proof) -> proof.
type imp_e           form -> proof -> proof -> proof.
type all_e           term  -> (term -> form) -> proof -> proof.
type all_i           (term -> proof) -> proof.
\end{lprolog}

\noindent Three design choices are visible in the types. A rule that \emph{discharges} a hypothesis takes
an abstraction over proofs (\lp{imp_i (proof -> proof)}), and the bound proof variable \emph{is} the
hypothesis label. A rule with an \emph{eigenvariable} abstracts over terms (\lp{all_i (term -> proof)}).
\emph{Elimination} rules store the formulas that appear in a premise but not in the conclusion, such as
the antecedent in \lp{imp_e A}, which makes checking syntax-directed. The checking relation
\lp{P # A}, ``P is a proof of A'', follows Felty~\cite{Felty1993} (\book{234}):

\begin{lprolog}
type  #     proof -> form -> o.
infix # 2.

true_i # tt.
(and_i P1 P2) # (A && B) :- (P1 # A), (P2 # B).
(or_i1 P) # (A !! B)  :- P # A.
(or_i2 P) # (A !! B)  :- P # B.
(imp_i Q) # (A ==> B) :- pi p\ (p # A) => ((Q p) # B).
(some_i T P) # (some A)  :- P # (A T).
(all_i Q) # (all A)   :- pi y\ (Q y) # (A y).
(false_e A P) # A     :- P # ff.
(and_e1 B P) # A      :- P # (A && B).
(and_e2 A P) # B      :- P # (A && B).
(or_e A B P Q1 Q2) # C
   :- (P # (A !! B)),
      (pi p1\ (p1 # A) => ((Q1 p1) # C)),
      (pi p2\ (p2 # B) => ((Q2 p2) # C)).
(imp_e A P1 P2) # B   :- (P1 # A), (P2 # (A ==> B)).
(some_e A P1 Q) # B
   :- (P1 # (some A)),
      pi y\ pi p\ (p # (A y)) => ((Q y p) # B).
(all_e T A P) # (A T) :- (P # (all A)).
\end{lprolog}

\begin{tutorial}[reading the ${\imp}I$ clause]
To check \lp{imp_i Q} against $A\imp B$, the clause picks a fresh proof constant \lp{p} (\lp{pi p\}),
assumes the atomic fact \lp{p # A} (\lp{=>}), and checks \lp{Q p} against $B$. Since \lp{Q} has type
\lp{proof -> proof}, \lp{Q p} $\beta$-reduces and puts the hypothesis label in place. Discharge is
automatic, because the assumption \lp{p # A} disappears when the goal \lp{(Q p) # B} has been proved.
In the \lp{all_i} clause, \lp{pi y\} produces a constant that occurs nowhere else, which is exactly the
eigenvariable condition. In the book's words, ``implicational goals are used to capture the hypothetical
judgments \ldots and universal goals are used to introduce and subsequently abstract over the proofs of
the assumed formulas'' and ``to enforce the newness constraint'' \book{234}.
\end{tutorial}

\noindent The same clauses can be run in four modes. Reading them is the best way to see what
``relational'' means for a meta-program. The vocabulary is \lp{a, b : form}, \lp{q : term -> form},
\lp{c : term}, \lp{f : term -> term} \book{234}.

\paragraph{1. Checking.} All three queries are solvable \book{234}:
\begin{lprolog}
?- (imp_i w\w) # (a ==> a).
?- (imp_i x\ imp_i y\ imp_e a x y) # (a ==> ((a ==> b) ==> b)).
?- (imp_i P\ all_i y\ imp_i Q\
      (imp_e (q (f y))
        (imp_e (q y) Q (all_e y (x\ (q x) ==> (q (f x))) P))
        (all_e (f y) (x\ (q x) ==> (q (f x))) P)))
   # ((all x\ (q x) ==> (q (f x))) ==>
      (all x\ (q x) ==> (q (f (f x))))).
\end{lprolog}
The third proof term reads, from the outside in: assume \lp{P}$\,:\forall x(q\,x\imp q(f\,x))$; for a
fresh \lp{y} assume \lp{Q}$\,:q\,y$; \lp{all_e y _ P} proves $q\,y\imp q(f\,y)$, so the inner
\lp{imp_e (q y) Q (...)} proves $q(f\,y)$; \lp{all_e (f y) _ P} proves $q(f\,y)\imp q(f(f\,y))$, and the
outer \lp{imp_e (q (f y)) ...} concludes $q(f(f\,y))$. Note that the instantiation bodies
\lp{x\ (q x) ==> (q (f x))} are spelled out inside both \lp{all_e} nodes: this is the information an
elimination rule stores so that checking stays syntax-directed.

\paragraph{2. Inference.} With the formula unknown, the checker computes it \book{235}:
\begin{lprolog}
?- (imp_i w\ (and_i (and_e2 a w) (and_e1 b w))) # R.
R = a && b ==> b && a.
\end{lprolog}

\paragraph{3. Completion.} ``Sometimes proofs can be incompletely described by including free variables
in them'' \book{235}. Here the two instantiation bodies of the third checking query are replaced by
logic variables \lp{A} and \lp{A'} (of type \lp{term -> form}), and the formula by \lp{B}. (Here \lp{P} and \lp{Q} are $\lambda$-bound variables, not logic variables, despite the
capitals.)
\begin{lprolog}
?- (imp_i P\ all_i y\ imp_i Q\
      (imp_e (q (f y)) (imp_e (q y) Q (all_e y A P))
             (all_e (f y) A' P))) # B.

A = w\ p w ==> p (f w)
A' = w\ p w ==> p (f w)
B = all (w\ p w ==> p (f w)) ==> all (w\ p w ==> p (f (f w)))
\end{lprolog}
(The printed answer says \lp{p} where the query says \lp{q}. Read \lp{q}.) The two unknowns are found in
different ways. \lp{A} is found by a pattern equation: \lp{(all_e y A P) # (q y ==> q (f y))} yields
\lp{A y = q y ==> q (f y)} with \lp{y} a \lp{pi}-constant, which is in $\Llam$ and has the unique solution
$\abs{w}{q\,w\imp q\,(f\,w)}$. \lp{A'} appears in \lp{A' (f y) = ...}, which is \emph{not} a pattern. It is
fixed only because \lp{P # (all A')} must match the stored hypothesis about \lp{P}, forcing \lp{A' = A}.
A Teyjus-style implementation delays the non-pattern equation until then (\cref{ch:hounif}).

\paragraph{4. Proving fails.} \lp{?- P # (a ==> ((a ==> b) ==> b)).} loops, although the second checking
query above contains its proof \book{235}. After two \lp{imp_i} steps the goal \lp{X # b} matches no
hypothesis. The elimination clauses then accept any formula as conclusion: \lp{(false_e A P) # A} asks for
\lp{P # ff}, which again only \lp{false_e} matches, and so on forever. Elimination rules used backwards have
unconstrained premises.

\begin{pitfall}[mode matters]
One specification serves as checker, inferencer and completer, and fails as a prover. The logical reading
is the same in all four modes. The operational behaviour is not. When a specification is ``relational'',
it is worth asking for each argument whether it is given, partially given, or unknown.
\end{pitfall}

\begin{kanren}[what natural deduction needs]
\textbf{$\forall$.} Fresh proof labels (\lp{pi p\}) and eigenvariables (\lp{pi y\}) are needed, the second
\emph{with a scope check}: a logic variable created before \lp{y} must never be bound to a term
containing \lp{y}. \textbf{$\imp$.} Every hypothesis is an \emph{atomic} fact (\lp{p # A}). An explicit
context argument with a lookup clause, the standard miniKanren technique for type environments, is
therefore enough. \textbf{Patterns.} Checking and inference work with first-order named or locally
nameless binders plus a substitution relation. Completion with an unknown \lp{A} in \lp{all_e} needs
pattern unification (\lp{A y = t} means ``abstract \lp{y} out of \lp{t}''), and \lp{A' (f y)} needs
postponement. \textbf{Search.} Under interleaving the prover mode becomes complete in principle: the
\lp{false_e} branch is only one of several fair branches. The search space of the unconstrained
elimination rules remains huge, however. The sources do not measure it. See \cref{ch:design} for how the
proposal provides these features.
\end{kanren}

% -----------------------------------------------------------------------------------------------------
\subsection{A classical prover with bounded contraction}\label{sec:apps-cl}

For classical first-order logic, the book uses formulas in \emph{negation normal form} (no $\imp$,
negation only on atoms) and a one-sided calculus \textsf{CL} with four \emph{zones} \book{236}:
\[
  \Sigma : \mathcal{L};\ \Gamma;\ \Phi
\]
Here $\Sigma$ holds the eigenvariables, $\mathcal{L}$ is a \emph{set} of literals, $\Gamma$ is a \emph{list}
of formulas still to be introduced, and $\Phi$ is a \emph{list} of existential formulas. The rules are (with
$A$ atomic, $t$ a $\Sigma$-term, $y\notin\Sigma$):

\begin{mathpar}
\inferrule{\Sigma : A,\mathcal{L};\ \Gamma;\ \Phi}{\Sigma : \mathcal{L};\ A,\Gamma;\ \Phi}\;\rulename{lit}
\and \inferrule{\Sigma : \neg A,\mathcal{L};\ \Gamma;\ \Phi}{\Sigma : \mathcal{L};\ \neg A,\Gamma;\ \Phi}\;\rulename{lit}
\and \inferrule{\Sigma : \mathcal{L};\ \Gamma;\ \Phi,\exists x.B}{\Sigma : \mathcal{L};\ \exists x.B,\Gamma;\ \Phi}\;\exists
\and \inferrule{ }{\Sigma : \mathcal{L};\ \top,\Gamma;\ \Phi}\;\top R
\and \inferrule{\Sigma : \mathcal{L};\ B,\Gamma;\ \Phi \\ \Sigma : \mathcal{L};\ C,\Gamma;\ \Phi}{\Sigma : \mathcal{L};\ B\land C,\Gamma;\ \Phi}\;\land R
\and \inferrule{\Sigma : \mathcal{L};\ B,C,\Gamma;\ \Phi}{\Sigma : \mathcal{L};\ B\lor C,\Gamma;\ \Phi}\;\lor R
\and \inferrule{\Sigma\cup\{y\} : \mathcal{L};\ B[y/x],\Gamma;\ \Phi}{\Sigma : \mathcal{L};\ \forall x\,B,\Gamma;\ \Phi}\;\forall R
\and \inferrule{ }{\Sigma : A,\neg A,\mathcal{L};\ \cdot;\ \Phi}\;\rulename{init}
\and \inferrule{\Sigma : \mathcal{L};\ B[t/x];\ \Phi,\exists x\,B}{\Sigma : \mathcal{L};\ \cdot;\ \exists x\,B,\Phi}\;\exists R
\end{mathpar}

\noindent The calculus is a systematic tableau in the sense of Smullyan~\cite{Smullyan1968}. A non-empty
$\Gamma$ determines the rule completely. When $\Gamma$ is empty, the only choices are \rulename{init} versus
$\exists R$, and the witness $t$. The rule $\exists R$ puts $\exists x\,B$ back at the \emph{end} of $\Phi$,
so all existentials get their turn. This alone does not give termination: for
$\cdot:\cdot;\ \exists x.q(x);\ \cdot$, ``proof search \ldots leads to the attempt to prove sequents
of the form $\cdot : q(t_1),\ldots,q(t_n);\cdot;\exists x.q(x)$, for increasingly larger sets of terms''
\book{237}. The remedy is to bound the number of instances of each existential and to retry with a larger
bound (iterative deepening).

\begin{pitfall}[no rule for $\bot$]
``If the first formula in the introduction zone is $\bot$ \ldots that sequent is not provable''
\book{236}. Hence $\bot\lor\top$ is not provable as written. In a one-sided classical calculus $\bot$ should
simply be dropped from $\Gamma$. A port should add \lp{prv (ff::Gamma) Phi :- prv Gamma Phi.} (This is our
notes' analysis. The book does not comment.)
\end{pitfall}

\noindent Literals are \lp{p A} and \lp{n A} for an atom \lp{A : atm}. The existential zone is a queue of
pairs (formula, remaining uses), implemented as a \emph{functional difference list} \book{238}: a list is
represented by a function from its tail to the whole list, so the empty queue is \lp{que x\x} and
\lp{qpush X (que k\ Phi k) (que k\ Phi ((pr X N)::k)) :- bound N.} appends at the end by
$\beta$-reduction. \lp{qpop} takes the front pair and, if uses remain, re-appends it with one use fewer.
All these matches are $\Llam$ patterns. Popping the empty queue fails because \lp{x\x} and
\lp{k\ (pr X 1)::(Phi k)} clash rigidly, so no emptiness test is needed. The prover itself is
\book{239}:

\begin{lprolog}
type lit         form -> o.
type prv         list form -> que form -> o.

prv (tt       :: Gamma) Phi.
prv ((B && C) :: Gamma) Phi :- prv (B::Gamma) Phi,
                               prv (C::Gamma) Phi.
prv ((B !! C) :: Gamma) Phi :- prv (B::C::Gamma) Phi.
prv ((all B)  :: Gamma) Phi :- pi x\ prv ((B x)::Gamma) Phi.
prv ((some B) :: Gamma) Phi :- qpush (some B) Phi Q,
                               prv Gamma Q.
prv (p A      :: Gamma) Phi :- lit (p A) => prv Gamma Phi.
prv (n A      :: Gamma) Phi :- lit (n A) => prv Gamma Phi.
prv nil Phi :- lit (n A), lit (p A).
prv nil Phi :- qpop (some B) Phi Q, prv ((B T)::nil) Q.

posints 1.
posints N :- posints M, N is M + 1.

thm B :- posints N, bound N => prv (B::nil) (que x\x).
\end{lprolog}

\noindent Two of the four zones are not arguments at all. $\Sigma$ is realized by the \lp{pi x\} of the
$\forall R$ clause: ``it uses $\lambda$Prolog-level universal quantifiers and the generic constants they
introduce during computation to implicitly realize the eigenvariables'' \book{238}. $\mathcal{L}$ is
realized by \lp{lit A =>}, a hypothetical fact. In the $\exists R$ clause the witness \lp{T} is a
\emph{fresh logic variable}. It is not guessed: the initial clause \lp{lit (n A), lit (p A)} chooses it
later, by unification. The side condition ``$t$ is a $\Sigma$-term'' is enforced by the unifier, which
refuses to bind \lp{T} to an eigenvariable introduced after \lp{T} was created. Iterative deepening comes
from \lp{posints}: a failure at bound $N$ backtracks into \lp{posints}, which produces $N+1$. ``B
represents a formula in negation normal form that has a CL proof if and only if the goal (thm B) is
provable'' \book{239}. On non-theorems \lp{thm} runs forever, since first-order logic is only
semi-decidable.

\begin{pitfall}[cut plus logic variables loses completeness]
An improvement closes a branch as soon as a complementary literal arrives \book{239}:
\lp{prv (p A :: Gamma) Phi :- lit (n A) ; lit (p A) => prv Gamma Phi.} (and the dual clause). Adding a
cut, \lp{lit (n A), ! ; ...}, ``is acceptable when there are no existential quantifiers'', but with them
``the interaction between such logic variables and the Prolog-cut (!) can lead to incompleteness''
\book{240}. The counterexample is
\begin{lprolog}
((p (r c)) !! (p (r t)) !! (p (g c)) !! (some x\ (n (r x)) && (n (g x)))).
\end{lprolog}
It is a theorem, but ``the logic variable that instantiates x gets bound to t first and the cut then
eliminates the path in which it would get bound to c'' \book{240}. Cut never breaks soundness, since it
only removes alternatives. It keeps completeness only after a test that cannot bind logic variables.
\end{pitfall}

\begin{kanren}[what the classical prover needs, and what interleaving gives]
\textbf{$\forall$ with a scope check is soundness-critical here.} Our notes derive two tests (not in the
book). The drinker formula $\exists x(D x\imp\forall y\,D y)$, in NNF
\lp{some x\ (n (d x)) !! (all y\ p (d y))}, must need bound~2: the first witness is created before $y$ and
cannot become $y$. The invalid $(\forall y\exists x\,R(x,y))\imp\exists x\forall y\,R(x,y)$ must never be
proved. Without the scope check a port proves it. \textbf{$\imp$.} Only atomic facts (\lp{lit A}, \lp{bound N})
are assumed, so an explicit literal-list argument and an explicit bound suffice. \textbf{Patterns.} The
functional difference list needs $\Llam$. A naive first-order open-tailed list is \emph{wrong}: popping an
empty \scm{(dl V V)} unifies \scm{V} with a cons and conjures an item. With a ground spine, an ordinary list
and \scm{append°} is the safe port. \textbf{Search.} Under interleaving the round-robin queue alone already
makes the prover complete, because every node has finitely many children and the infinite $\exists R$ chain
is interleaved with the \rulename{init} alternatives. The bound then serves only to make failing searches
stop. \textbf{Cut.} \micro{}'s \scm{if°} is a \emph{soft} cut: it keeps every answer of the test. Translating
\lp{!} as \scm{(if° (take° 1 test) ...)} reproduces the book's incompleteness exactly. Plain \scm{if°}
happens to prove the formula above but is still incomplete in general. Only the cut-free disjunction is
safe. See \cref{ch:design}.
\end{kanren}

% -----------------------------------------------------------------------------------------------------
\subsection{Tactics and tacticals}\label{sec:apps-tactics}

Direct reflection ``has the benefit of yielding a transparently correct specification. However, it also
means that we must use the control regimen of the interpreter'' \book{240}. The remedy is the LCF
architecture~\cite{GordonMilnerWadsworth1979}, recast as relations~\cite{FeltyMiller1988,Felty1993}. An
inference rule becomes \emph{data}: a relation between the goal of proving the conclusion and the goal of
proving the premises. Search strategies then become ordinary higher-order predicates over such relations.

\begin{definition}[goals, tactics, tacticals; \book{241}]\label{def:apps-tactic}
A \emph{primitive goal} is an object to be proved (a sequent, a typing judgment \lp{P # A}, \ldots).
\emph{Compound goals} are \lp{trueg} (closed), \lp{G1 cc G2} (both) and \lp{allg (x\ G)} (for a generic
\lp{x}). A \emph{tactic} is a predicate of type \lp{goal -> goal -> o} relating a primitive goal to the
goal of its premises. It fails when it does not apply, and it may leave logic variables, such as
quantifier witnesses, in its output. A \emph{tactical} is a higher-order predicate combining tactics.
\end{definition}

\noindent The goal language and its simplifier are \book{241}:

\begin{lprolog}
kind  goal     type.
type  trueg    goal.                    % vacuously true goal
type  cc       goal -> goal -> goal.    % conjunctive goal
type  allg     (A -> goal) -> goal.     % universally quantified goal
infixl cc 3.

type  goalreduce, redex, red1   goal -> goal -> o.
type  primgoal                  goal -> o.

redex (trueg cc G) G & redex (G cc trueg) G.
redex (allg x\trueg) trueg.
red1 G H :- redex G H, !.
red1 (G cc H) (Gx cc H) &
red1 (H cc G) (H cc Gx) :- red1 G Gx.
red1 (allg G) (allg Gx) :- pi x\ red1 (G x) (Gx x).
goalreduce G H :- red1 G Gx, !, goalreduce Gx H.
goalreduce G G.
\end{lprolog}

\noindent \lp{goalreduce} rewrites with $\top_g\mathbin{cc}G\to G$, $G\mathbin{cc}\top_g\to G$ and
$\mathit{allg}\,(\abs{x}{\top_g})\to\top_g$ anywhere in a goal, until no rewrite applies. A goal ``is
proved'' when it reduces to \lp{trueg}. The pattern \lp{allg x\trueg} matches only a \emph{vacuous}
abstraction, one whose body does not mention \lp{x}. The binder type \lp{A} of \lp{allg} is polymorphic,
so the same goal language serves goals over terms and over proofs. The cuts are safe here because the
input is ground and the rewriting is confluent.

\paragraph{A first example: reachability.} The book's smallest instance has nodes, edges as primitive goals
\lp{adj X Y}, and paths as primitive goals \lp{path X Y} \book{242}:

\begin{lprolog}
kind node   type.
type a, b, c, d, e   node.
type adj, path   node -> node -> goal.

primgoal (adj _ _)  &  primgoal (path _ _).

type adj_tac, path_base_tac, path_rec_tac   goal -> goal -> o.

adj_tac (adj a b) trueg  &  adj_tac (adj a c) trueg.
adj_tac (adj b d) trueg  &  adj_tac (adj c d) trueg.
adj_tac (adj d a) trueg  &  adj_tac (adj d e) trueg.

path_base_tac  (path X Y) (adj X Y).
path_rec_tac   (path X Y) ((adj X Z) cc (path Z Y)).
\end{lprolog}

\noindent \lp{adj_tac} closes an edge goal. The two path tactics are the two inference rules for
\lp{path}. In \lp{path_rec_tac} the intermediate node \lp{Z} appears only in the output: a tactic may leave
an unknown in the premises, to be fixed later when \lp{adj_tac} is applied to \lp{adj X Z}. The graph
($a{\to}b$, $a{\to}c$, $b{\to}d$, $c{\to}d$, $d{\to}a$, $d{\to}e$) has the cycle $a{\to}b{\to}d{\to}a$.
The book gives no query for this figure. A run on \micro{} is in the kanren box below.

\paragraph{Sequent calculus as tactics.} The rules of an intuitionistic sequent calculus for $\land$,
$\imp$ and $\forall$ become tactics \book{242}:

\begin{lprolog}
type sq   list form -> form -> goal.
primgoal (sq _ _).

type initial, and_r, imp_r, all_r, and_l, imp_l, all_l, all_l'   goal -> goal -> o.

initial (sq Gamma A) trueg :- memb_and_rest A Gamma _.
and_r   (sq Gamma (A && B)) ((sq Gamma A) cc (sq Gamma B)).
imp_r   (sq Gamma (A ==> B)) (sq (A::Gamma) B).
all_r   (sq Gamma (all A)) (allg x\ sq Gamma (A x)).
and_l   (sq Gamma A) (sq (B::C::Gamma') A) :-
          memb_and_rest (B && C) Gamma Gamma'.
imp_l   (sq Gamma A) ((sq Gamma B) cc (sq (C::Gamma') A)) :-
          memb_and_rest (B ==> C) Gamma Gamma'.
all_l   (sq Gamma A) (sq ((B T)::Gamma) A) :-
          memb_and_rest (all B) Gamma Gamma'.
all_l'  (sq Gamma A) (sq ((B T)::Gamma') A) :-
          memb_and_rest (all B) Gamma Gamma'.
\end{lprolog}

\noindent Each clause relates a conclusion to the goal of its premises. \lp{imp_l} is the plain ${\imp}L$ of
the first calculus: its left premise \lp{sq Gamma B} keeps $B\imp C$. The tactic \lp{all_r} does not
introduce the eigenvariable itself. It returns a quantified \emph{goal} \lp{allg x\ ...}, and the
eigenvariable becomes a real \lp{pi}-constant only when a tactical enters that goal. The tactic \lp{all_l}
keeps the universal hypothesis (contraction: it uses the original \lp{Gamma}), while \lp{all_l'} drops it
(it uses \lp{Gamma'}). In both, \lp{T} is an open witness. The tacticals are \book{244}:

\begin{lprolog}
type maptac   (goal -> goal -> o) -> goal -> goal -> o.

maptac Tac trueg trueg.
maptac Tac (I1 cc I2) (O1 cc O2) :- maptac Tac I1 O1,
                                    maptac Tac I2 O2.
maptac Tac (allg In) (allg Out) :- pi t\ maptac Tac (In t) (Out t).
maptac Tac In Out :- primgoal In, Tac In Out.

type idtac          goal -> goal -> o.
type repeat, try    (goal -> goal -> o) -> goal -> goal -> o.
type then,
     orelse, orelse! (goal -> goal -> o) ->
                     (goal -> goal -> o) -> goal -> goal -> o.

idtac             In In.
then    Tac1 Tac2 In Out :- Tac1 In Mid, maptac Tac2 Mid Out.
orelse  Tac1 Tac2 In Out :- Tac1 In Out  ; Tac2 In Out.
orelse! Tac1 Tac2 In Out :- Tac1 In Out, ! ; Tac2 In Out.
repeat  Tac       In Out :- orelse (then Tac (repeat Tac))
                                   idtac In Out.
try     Tac       In Out :- orelse Tac idtac In Out.
\end{lprolog}

\noindent Read tactics as binary relations. Then \lp{idtac} is the identity, \lp{orelse} is union,
\lp{orelse!} is ``the first answer of the first tactic, else the second'', \lp{try T} is ``\lp{T} if it
applies, else no change'', and \lp{repeat T} applies \lp{T} as often as possible first, with fewer
applications on backtracking. \lp{maptac T} lifts \lp{T} from primitive goals to compound goals. It enters
\lp{allg} with \lp{pi t\} and re-abstracts the result through the pattern \lp{Out t}, and it fails if
\lp{T} fails on any primitive leaf. \lp{then T1 T2} is the relational composition of \lp{T1} with
\lp{maptac T2}. It cannot simply be \lp{T1 In Mid, T2 Mid Out}, because \lp{Mid} may be compound
\book{243}.

\paragraph{Invertible rules first.} A rule is \emph{invertible} when its premises are provable iff its
conclusion is. Such rules ``can be forced to appear at the end of a proof without any loss of
completeness'' \book{243}:

\begin{lprolog}
invertible In Out :-
   repeat (orelse and_r (orelse and_l (orelse imp_r all_r))) In Out.
\end{lprolog}

\noindent With \lp{p : term -> form} and \lp{f : term -> term}, the book reports \book{243--244}:

\begin{lprolog}
?- invertible (sq [] ((a && (a ==> b)) ==> (a && b))) Out.
Out = ((sq (a :: (a ==> b) :: nil) a) cc (sq (a :: (a ==> b) :: nil) b))

?- invertible (sq [] ((all x\ (p x) ==> (p (f x))) ==>
                      (all x\ (p x) ==> (p (f (f x))))))
              Out.
Out = (allg (w\ sq (p w :: (all w\ p w ==> p (f w)) :: nil)
                   (p (f (f w)))))
\end{lprolog}

\noindent What remains in the second answer becomes propositional after two instantiations of the
universal hypothesis. To finish it, the book links the decision procedure of
\cref{sec:apps-ipl} in as a tactic and composes a \emph{script} \book{244--245}:

\begin{lprolog}
ip_decide (sq Gamma A) trueg :- seq Gamma A.

?- then invertible (then all_l (then all_l' ip_decide))
       (sq [] ((all x\ (p x) ==> (p (f x))) ==>
               (all x\ (p x) ==> (p (f (f x))))))
       Out.
Out = allg W1\ trueg
\end{lprolog}

\begin{tutorial}[how the script closes the goal]
\lp{invertible} leaves \lp{allg w\ sq [p w, all (...)] (p (f (f w)))}. \lp{maptac} enters the \lp{allg}
with a fresh constant \lp{w}. \lp{all_l} adds \lp{p T1 ==> p (f T1)} and keeps the $\forall$. \lp{all_l'}
adds \lp{p T2 ==> p (f T2)} and drops it. Because \lp{T1} and \lp{T2} were created \emph{inside} the scope
of \lp{w}, they may be bound to terms containing \lp{w}. \lp{ip_decide} runs \lp{seq}. There, ${\imp}L_1$ on
the first implication needs the atom \lp{p T2} in the context, and \lp{memb_and_rest} \emph{unifies}
\lp{T2 := w}, adding \lp{p (f w)}. The second implication then gets \lp{T1 := f w}, adding
\lp{p (f (f w))}, and \rulename{init} closes the sequent. The witnesses are found by unification at the
leaves, exactly as in the classical prover. (\lp{seq} needs \lp{atom} for first-order atoms, for example
\lp{atom (p _).}, which the book does not show.) The output \lp{allg W1\ trueg} reduces to \lp{trueg}
with the goal-simplifier \lp{goalreduce} of \book{241}.
\end{tutorial}

\begin{kanren}[tacticals in Scheme, and what interleaving does to them]
Tacticals are higher-order relations, and in Scheme they are just procedures from relations to relations.
Everything except \lp{allg} is first order, so the reachability example \book{242} runs on \micro{}
today. Goals are tagged lists, and \scm{maptac°} has no \lp{allg} case:
\begin{scheme}[fontsize=\footnotesize]
(define-relation (primgoal° g)
  (fresh° (x y) (or° (=° g `(adj ,x ,y)) (=° g `(path ,x ,y)))))

(define-relation (adj-tac° in out)
  (=° out 'trueg)
  (cond° ((=° in '(adj a b))) ((=° in '(adj a c))) ((=° in '(adj b d)))
         ((=° in '(adj c d))) ((=° in '(adj d a))) ((=° in '(adj d e)))))
(define-relation (path-base-tac° in out)
  (fresh° (x y) (=° in `(path ,x ,y)) (=° out `(adj ,x ,y))))
(define-relation (path-rec-tac° in out)
  (fresh° (x y z) (=° in `(path ,x ,y)) (=° out `(cc (adj ,x ,z) (path ,z ,y)))))

(define (maptac° tac)
  (define-relation (m in out)
    (cond° ((=° in 'trueg) (=° out 'trueg))
           ((fresh° (i1 i2 o1 o2) (=° in `(cc ,i1 ,i2)) (=° out `(cc ,o1 ,o2)) (m i1 o1) (m i2 o2)))
           ((primgoal° in) (tac in out))))
  m)
(define (idtac° in out) (=° in out))
(define ((then° t1 t2) in out) (fresh° (mid) (t1 in mid) ((maptac° t2) mid out)))
(define ((orelse° t1 t2) in out) (or° (t1 in out) (t2 in out)))
(define ((try° t) in out) ((orelse° t idtac°) in out))
(define ((repeat° t) in out)
  (λ (s) (δ (((orelse° (then° t (repeat° t)) idtac°) in out) s))))

;; stands in for `goalreduce Out trueg`: Out is made of trueg and cc only
(define-relation (closed° g)
  (cond° ((=° g 'trueg))
         ((fresh° (g1 g2) (=° g `(cc ,g1 ,g2)) (closed° g1) (closed° g2)))))

(define step° (orelse° adj-tac° (orelse° path-base-tac° path-rec-tac°)))

(pretty-print (μkanren-run (q -1 #t) ((then° path-rec-tac° (try° adj-tac°)) '(path a e) q)))
(pretty-print (μkanren-run (q 4 #t) ((repeat° step°) '(path a e) q)))
(pretty-print (μkanren-run (q 3 #t) ((repeat° step°) '(path a e) q) (closed° q)))
\end{scheme}
\begin{answers}
((cc (adj a α) (path α e)) (cc trueg (path b e)) (cc trueg (path c e)))
((path a e) (adj a e) (cc (adj a α) (path α e)) (cc trueg (path b e)))
((cc trueg (cc trueg trueg))
 (cc trueg (cc trueg trueg))
 (cc trueg (cc trueg (cc trueg (cc trueg (cc trueg trueg))))))
\end{answers}
The open witness \scm{α} of \scm{path-rec-tac°} is chosen later by \scm{adj-tac°}, just as \lp{T} is in
\lp{all_l}. The first answer of \scm{repeat°} is \emph{zero} applications, \scm{(path a e)}. The book's
depth-first \lp{repeat} would instead answer first with the maximal sequence, which here is the unprovable
leftover \lp{adj a e}. Interleaving enumerates all application lengths fairly. On the cyclic graph it
produces the two three-edge proofs ($a{\to}b{\to}d{\to}e$, $a{\to}c{\to}d{\to}e$) and then a six-edge one
around the cycle. A depth-first enumeration of all answers would never terminate here, and interleaving
does not terminate either, but it keeps producing proofs. A script built from \lp{orelse} and
\lp{repeat} is therefore \emph{complete} under interleaving, while committed choice (\lp{orelse!}, the cuts
of \lp{goalreduce}) needs \scm{(if° (take° 1 ...) ...)}. \textbf{What is missing} is \lp{maptac} over \lp{allg}. It needs a
universal goal (a fresh eigenvariable) and the abstraction of that eigenvariable out of the result
(\lp{Out t}, a pattern). See \cref{ch:design}.
\end{kanren}

% =====================================================================================================
\section{Computations over functional programs}\label{sec:apps-fp}

Chapter 10 of the book turns to programming-language semantics. The object language \textsf{miniFP} is an
untyped call-by-value $\lambda$-calculus with a conditional, a fixed point, booleans, integers and lists
\book{248}:

\[
  M ::= x\mid \lambda x.M\mid M\,N\mid \mathsf{cond}\,M\,N\,P\mid \mathsf{fix}\,x.M\mid \mathsf{cns}\,M\,N
        \mid \mathsf{i}\,n\mid \mathsf{tt}\mid\mathsf{ff}\mid\mathsf{null}\mid \mathsf{cons}\mid\mathsf{car}\mid\ldots
\]
\begin{lprolog}
type abs    (tm -> tm) -> tm.          % function abstraction
type @      tm -> tm -> tm.            % application
infixl @ 4.                            % application is infix
type cond   tm -> tm -> tm -> tm.      % conditional
type fixpt  (tm -> tm) -> tm.          % recursive functions
\end{lprolog}

\noindent The remaining constructors are first order: \lp{cns} is the list \emph{value} constructor, \lp{i}
injects integers, and constants such as \lp{cons}, \lp{car}, \lp{sum}, \lp{equal} have type \lp{tm}. \lp{cons} is a built-in \emph{function}, applied
with \lp{@}. The encoding is untyped, so \lp{cons @ (i 4) @ (i 5)} is a legal \lp{tm} but not a valid
program. Validity is a separate typing judgment. Four sample programs are named by \lp{prog} \book{249}.
For instance:

\begin{lprolog}
prog "map" (fixpt map\ abs f\ abs l\
  cond (nullp @ l) null (cons @ (f @ (car @ l)) @
                                (map @ f @ (cdr @ l)))).
\end{lprolog}

\subsection{Typing: hypotheses about bound variables}\label{sec:apps-typeof}

The four structural rules of simple typing are \book{250}:

\begin{lprolog}
typeof (M @ N) A         :- typeof M (arr B A), typeof N B.
typeof (cond P Q R) A    :- typeof P bool, typeof Q A, typeof R A.
typeof (abs M) (arr A B) :- pi x\ typeof x A => typeof (M x) B.
typeof (fixpt M) A       :- pi x\ typeof x A => typeof (M x) A.
\end{lprolog}

\noindent One axiom per constant follows, such as \lp{typeof car (arr (lst A) A).} There is no clause for
variables. Instead, when the \lp{abs} clause enters the binder, it \emph{assumes} a clause
\lp{typeof x A} for the fresh \lp{x}. The book reports \book{249--250} that \lp{typeof (abs w\ w)} has type
\lp{arr int int}, \lp{arr bool bool}, and even \lp{pi t\ typeof (abs w\ w) (arr t t)}. For the four programs
it reports:

\begin{lprolog}
?- sigma Exp\ prog Name Exp, typeof Exp Ty.
Ty = arr int int.
Ty = arr A (arr (lst A) bool).
Ty = arr (lst A) (arr (lst A) (lst A)).
Ty = arr (arr A B) (arr (lst A) (lst B)).
\end{lprolog}

\begin{pitfall}[polymorphic constants, monomorphic variables]
The type variables of a constant's axiom are renamed at each use of the clause, so each occurrence of
\lp{car} may get its own instance. The assumed clause \lp{typeof x A}, by contrast, contains the
\emph{outer} logic variable \lp{A}, which is not renamed, so all occurrences of \lp{x} share one type.
That is exactly simple typing. It is also why let-polymorphism is out of reach: ``no simple extension to
the program provided in Figure 10.3 can realize this algorithm'' \book{260}.
\end{pitfall}

\begin{kanren}[the classic miniKanren type inferencer is this program]
Since every hypothesis is an atomic fact about a fresh variable, an explicit environment $\Gamma$ can
replace \lp{=>}. Variables are then named by Scheme symbols, and constants are wrapped as \scm{(c name)}
to keep the two apart. Sharing a logic variable inside $\Gamma$ is precisely the non-renaming of \lp{A}.
Wrapping each constant's type variables in \scm{fresh°} is the per-use renaming. Excerpt (the remaining
constant clauses, and the four programs written with a helper \scm{@*} for iterated application, are in
the run file):
\begin{scheme}[fontsize=\footnotesize]
(define-relation (lookup° Γ x t)
  (fresh° (y u r)
    (=° Γ `((,y . ,u) . ,r))
    (cond° ((=° x y) (=° t u))
           ((≠° x y) (lookup° r x t)))))

(define-relation (typeof° Γ m t)
  (cond°
    ((symbol° m) (lookup° Γ m t))
    ((fresh° (M N B) (=° m `(@ ,M ,N)) (typeof° Γ M `(arr ,B ,t)) (typeof° Γ N B)))
    ((fresh° (P Q R) (=° m `(cond ,P ,Q ,R)) (typeof° Γ P 'bool) (typeof° Γ Q t) (typeof° Γ R t)))
    ((fresh° (x M A B) (=° m `(abs ,x ,M)) (=° t `(arr ,A ,B)) (typeof° `((,x . ,A) . ,Γ) M B)))
    ((fresh° (f M) (=° m `(fixpt ,f ,M)) (typeof° `((,f . ,t) . ,Γ) M t)))
    ((fresh° (n) (=° m `(i ,n)) (=° t 'int)))
    ((=° m '(c tt)) (=° t 'bool))
    ;; ... (c ff), (c null), (c equal), (c cons) ...
    ((fresh° (A) (=° m '(c car)) (=° t `(arr (lst ,A) ,A))))
    ;; ... (c cdr), (c nullp), (c consp), (c and), (c zerop), (c minus), (c sum)
    ))

(for-each (λ (p) (pretty-print (cons (car p) (μkanren-run (q -1 #t) (typeof° '() (cdr p) q)))))
  `((id . (abs w w)) (fib . ,fib) (mem . ,mem) (appnd . ,appnd) (map . ,map) (self . (abs w (@ w w)))))
\end{scheme}
\begin{answers}
(id (arr α α))
(fib (arr int int))
(mem (arr α (arr (lst α) bool)))
(appnd (arr (lst α) (arr (lst α) (lst α))))
(map (arr (arr α β) (arr (lst α) (lst β))))
(self)
\end{answers}
These are the book's types. Self-application has none, because \micro{}'s occurs check rejects
$B = B\to A$. Running the relation \emph{backwards} enumerates inhabitants of a type:
\begin{scheme}[fontsize=\footnotesize]
(pretty-print (μkanren-run (q 5 #t) (typeof° '() q '(arr int int))))
\end{scheme}
\begin{answers}
((abs α α)
 (fixpt α α)
 (abs α (i β))
 (@ (abs α α) (abs β β))
 (abs α (fixpt β β)))
\end{answers}
(Grounded answers omit the pending \scm{symbol°} constraints on \scm{α}, \scm{β}.) Note the
well-typed divergent term \scm{(fixpt α α)}. In the book's clause order, depth-first search would try the
application clause first and recurse on an ever larger function type. Interleaving gives every clause its
turn. The named-binder encoding works here because typing never substitutes. Evaluation does, and needs
more (see below).
\end{kanren}

\subsection{Big-step evaluation}\label{sec:apps-eval}

The natural semantics~\cite{Kahn1987} of miniFP is a relation $M\Downarrow V$ \book{251--252}. Its
structural core reads:

\begin{mathpar}
\inferrule{V\text{ a value}}{V\Downarrow V}
\and \inferrule{M\Downarrow F \\ N\Downarrow U \\ \mathrm{apply}(F,U)\Downarrow V}{M\,N\Downarrow V}
\and \inferrule{M[\mathsf{fix}\,x.M/x]\Downarrow V}{\mathsf{fix}\,x.M\Downarrow V}
\and \inferrule{M[U/x]\Downarrow V}{\mathrm{apply}(\lambda x.M,\,U)\Downarrow V}
\end{mathpar}

\begin{lprolog}
val (abs _) & val (i _) & val tt & val ff & val null.
val (cns _ _) & val (spec _ _ _).

eval V V            :- val V.
eval (M @ N) V      :- eval M F, eval N U, apply F U V.
eval (fixpt R) V    :- eval (R (fixpt R)) V.
eval (cond C L R) V :- eval C B, if (B = tt) (eval L V) (eval R V).
eval F (spec I F nil) :- special I F.

apply (abs R) U V         :- eval (R U) V.
apply (spec 1 F Args) U V :- eval_spec F (U::Args) V.
apply (spec C F Args) U (spec D F (U::Args)) :- C > 1, D is C - 1.
\end{lprolog}

\noindent Application of an abstraction ``corresponds to a substitution that is realized elegantly through
$\beta$-conversion in the specification language'' \book{251}: \lp{eval (R U) V}. A built-in \lp{F} of
arity $n$ evaluates to the partial application \lp{spec n F nil}, which collects already evaluated
arguments \emph{in reverse order} until the $\delta$-rule \lp{eval_spec} fires, for example
\lp{eval_spec minus ((i N)::(i M)::nil) (i V) :- V is M - N.} With these clauses
\lp{prog "fib" F, eval (F @ (i 12)) V} gives \lp{V = (i 144)}, and mapping fib over $[9,4]$ gives
\lp{(cns (i 34) (cns (i 3) null))} \book{251--253}. Since \lp{eval_spec equal} compares with
$\lambda$Prolog's \lp{=}, which is equality of $\lambda$-terms up to $\alpha\beta\eta$,
\lp{eval (equal @ (abs x\x) @ (abs y\y)) V} gives \lp{V = tt} \book{253}. This is ``stronger'' than a
functional language, which cannot compare functions. (The book suggests replacing \lp{B = C} by a
first-order \lp{eq B C} that only compares integers, booleans and lists. Since \lp{eq} has no clause for
\lp{abs} and the test sits inside an \lp{if}, the same query then gives \lp{V = ff}, our notes observe.)

\begin{tutorial}[evaluating a small application]
Take \lp{eval ((abs x\ cond x (i 1) (i 2)) @ tt) V}. The clause for \lp{@} evaluates the function, which
is a value and evaluates to itself, so \lp{F = abs x\ cond x (i 1) (i 2)}; it evaluates the argument
\lp{tt} to \lp{U = tt}; then \lp{apply (abs R) U V} calls \lp{eval (R U) V}. Here \lp{R} is the
meta-level function \lp{x\ cond x (i 1) (i 2)}, so \lp{R U} $\beta$-reduces to
\lp{cond tt (i 1) (i 2)}: object-level substitution has been done by the metalanguage. The \lp{cond}
clause evaluates \lp{tt} to \lp{B = tt} and takes the then branch, \lp{V = i 1}.
\end{tutorial}

\noindent New binding constructs fit the same pattern. The book's bibliographic notes add \lp{let}, with
concrete \lp{let x = t in b} represented as \lp{(let t (x\b))}, and one clause \book{259--260}:
\lp{eval (let T R) V :- eval T U, eval (R U) V.}

\begin{pitfall}[\texttt{if} is a commit, not logic]
The conditional is defined with a cut \book{137}:
\begin{lprolog}
type if   o -> o -> o -> o.
if P Q R :- P, !, Q.
if P Q R :- R.
\end{lprolog}
It commits to the first proof of \lp{P}. The book says it is used ``only for convenience'' and can be
rewritten as \lp{(U = tt, G ; U = ff, H)} when \lp{U} is a boolean \book{251}. The two readings differ on a non-boolean test. They also differ, more seriously, on an
\emph{eigenvariable} test: \lp{k = tt} fails for a generic \lp{k}, so \lp{if} silently takes the else
branch. This matters for mixed evaluation below.
\end{pitfall}

\subsection{Evaluation contexts as meta-level functions}\label{sec:apps-evalc}

The reduction semantics of Felleisen and Hieb~\cite{FelleisenHieb1992} splits a term into an
\emph{evaluation context} with a hole and the next redex: $E[R]\longmapsto E[N]$ when $R$ reduces to $N$,
with $E ::= [\,]\mid \mathsf{cond}\,E\,N\,P\mid E\,N\mid V\,E$. In $\lambda$Prolog a context is simply a
function \lp{tm -> tm}, and plugging is application \book{254}:

\begin{lprolog}
context R (x\ x) R :- redex R.
context (cond M N P) (x\ cond (E x) N P) R :-
    non_val M, context M E R.
context (M @ N) (x\ (E x) @ N) R :-
    non_val M, context M E R.
context (V @ M) (x\ V @ (E x)) R :-
    val V, non_val M, context M E R.

evalc V V :- val V.
evalc M V :- context M E R, reduce R N, evalc (E N) V.
\end{lprolog}

\noindent \lp{redex}, \lp{non_val} and \lp{reduce} are first-order tables (for example
\lp{reduce ((abs R) @ N) (R N).}). Each decomposition succeeds exactly once \book{255}:
\begin{lprolog}
?- context (cond ((abs x\ ff) @ tt) (i 2) (i 3)) E R.
%  succeeds exactly once:  R = ((abs x\ ff) @ tt),  E = (W\ cond W (i 2) (i 3))
\end{lprolog}
\noindent The output context \lp{x\ cond (E x) N P} is built from the recursively computed \lp{E} by
abstraction. Every context answer is therefore a higher-order output.

\subsection{Partial evaluation: unfolding and mixed evaluation}\label{sec:apps-mixeval}

Both evaluators unfold \lp{fixpt R} into \lp{R (fixpt R)}. The query
\lp{prog "map" (fixpt Body), Unfold = (Body (fixpt Body))} exhibits this as a source-to-source
transformation \book{255}. (The printed result is mis-parenthesized: \lp{@ f @ (cdr @ l)} must apply the
inner \lp{fixpt}, not lie inside its body.) More interesting is \emph{mixed evaluation}, which evaluates
under a binder by treating the bound variable as an unknown value \book{256}:

\begin{lprolog}
type mixeval  tm -> tm -> o.
mixeval (abs R) (abs S) :- pi k\ val k => eval (R k) (S k).

?- prog "appnd" App,
   eval (App @ (cons @ (i 1) @ (cons @ (i 5) @ null))) R,
   mixeval R S.
\end{lprolog}

\noindent This ``yields the single answer substitution, binding the variable S to''
\lp{abs w\ cons @ (i 1) @ (cons @ (i 5) @ w)} \book{256}. (By the \lp{eval_spec cons} clause the value form
should be \lp{cns (i 1) (cns (i 5) w)}, the same function.) Append has been specialized away. Note
\lp{(S k)}: a logic variable applied to an eigenvariable, the pattern case. Its solution abstracts \lp{k}
out of the computed value.

\begin{proposition}[\book{256}]\label{prop:apps-mixeval}
If \lp{mixeval (abs R) (abs S)} holds and \lp{val T} is provable, then evaluating \lp{(R T)} yields the value
\lp{(S T)}.
\end{proposition}

\noindent The proof idea is the substitution property of $\forall$ and $\imp$: instantiate the generic
derivation of \lp{eval (R k) (S k)} at \lp{T} and discharge \lp{val k} with \lp{val T}.

\begin{pitfall}[the proposition needs the logical reading]
With the cut-based \lp{if}, \cref{prop:apps-mixeval} fails. Our notes' counterexample is
\lp{mixeval (abs b\ cond b (i 1) (i 2)) S}. The test \lp{if (k = tt) ...} cannot unify the
eigenvariable, so the else branch gives \lp{S = b\ i 2}, but \lp{eval (cond tt (i 1) (i 2))} is \lp{i 1}.
With \lp{(B = tt, ... ; B = ff, ...)} the conditional on \lp{k} simply fails, and soundness holds. Failing
on an eigenvariable means ``unknown'', not ``different''.
\end{pitfall}

\subsection{CPS conversion and administrative redexes}\label{sec:apps-cps}

Fischer's call-by-value CPS transformation~\cite{Fischer1972} \book{257} on the pure fragment is
defined by $\mathcal{F}$ on terms and $\Psi$ on values (abstractions and variables). The marked
$\bar\lambda$ are the \emph{administrative} abstractions:
\begin{align*}
  \mathcal{F}[V] &= \bar\lambda k.\,(k\;\Psi[V])\\
  \mathcal{F}[M\,N] &= \bar\lambda k.\,(\mathcal{F}[M]\;\bar\lambda m.\,(\mathcal{F}[N]\;\bar\lambda n.\,(m\,k\,n)))\\
  \Psi[x] &= x\\
  \Psi[\lambda x.M] &= \lambda k.\,\lambda x.\,(\mathcal{F}[M]\,k)
\end{align*}

\begin{lprolog}
type ftrans, phi  tm -> tm -> o.

ftrans (abs V) (abs k\ k @ U) :- phi (abs V) U.
ftrans (M @ N) (abs k\ P @ (abs m\ Q @ (abs n\ m @ k @ n))) :-
    ftrans M P, ftrans N Q.

phi (abs M) (abs k\ abs x\ (P x) @ k) :-
    pi x\ pi y\ ftrans x (abs k\ k @ y) => ftrans (M x) (P y).
\end{lprolog}

\begin{tutorial}[where did the variable case go?]
Neither \lp{ftrans} nor \lp{phi} has a clause for variables. When \lp{phi} enters the binder of
\lp{abs M}, it creates a source variable \lp{x}, a target variable \lp{y}, and the local clause
\lp{ftrans x (abs k\ k @ y)}. That clause \emph{is} the equation $\mathcal{F}[x]=\bar\lambda k.(k\,\Psi[x])$
for this one variable. The translated body is \lp{P y}, a pattern, so \lp{P} is the body abstracted over
\lp{y}, and \lp{abs x\ (P x) @ k} re-binds it. The same device reappears in the $\lambda$-to-$\pi$
translation of \cref{sec:apps-lampi}: ``one base case is assumed for every bound variable encountered''
\book{274}.
\end{tutorial}

\noindent On $(\lambda u.u)(\lambda u.u)$ the book reports \book{257}:
\begin{lprolog}
?- ftrans ((abs u\u) @ (abs u\u)) F.

F = (abs W1\ (abs (W2\ W2 @ abs W3\ abs W4\ abs (W5\ W5 @ W4) @ W3)) @
     (abs W2\ (abs (W3\ W3 @ abs W4\ abs W5\ abs (W6\ W6 @ W5) @ W4)) @
     (abs W3\ W2 @ W1 @ W3))).
\end{lprolog}
\noindent To read it, let $U$ be $\Psi[\lambda u.u]$ before any simplification,
\lp{abs a\ abs b\ ((abs c\ c @ b) @ a)}. Then \lp{F} is
\lp{abs k\ (abs k1\ k1 @ U) @ (abs m\ (abs k2\ k2 @ U) @ (abs n\ m @ k @ n))}, which is the
application equation of $\mathcal{F}$ with both sides instantiated by the value equation. The result
contains many $\beta$-redexes that exist ``solely for the purpose of providing a compositional account of
the transformation''. These are the administrative redexes. The two-phase version \book{258} builds administrative abstractions with a separate constructor
\lp{adm : (tm -> tm) -> tm} and then reduces them away:
\begin{lprolog}
admred ((adm R) @ N) (R N).

red1 M N :- admred M N.
red1 (M @ N) (M' @ N) & red1 (N @ M) (N @ M') :- red1 M M'.
red1 (adm R) (abs S) & red1 (abs R) (abs S) :-
    pi x\ red1 (R x) (S x).

red M N :- red1 M P, !, red P N.
red M M.

?- ftrans ((abs x\x) @ (abs x\x)) T, red T S.
S = (abs W1\ (abs W2\ abs W3\ W2 @ W3) @ W1 @ (abs W2\ abs W3\ W2 @ W3)).
\end{lprolog}
\noindent The \lp{ftrans}/\lp{phi} clauses are those above with \lp{adm} for the marked $\bar\lambda$. The
answer is $\lambda k.\,\Psi'\,k\,\Psi'$ with $\Psi'=\lambda k.\lambda x.(k\,x)$. The cut in \lp{red} is
safe because administrative reduction is confluent and terminating.

\begin{kanren}[what functional meta-programs need]
\textbf{$\forall$ and $\imp$.} Every hypothesis in this chapter is an atomic fact about a fresh variable:
\lp{typeof x A}, \lp{val k}, \lp{ftrans x (abs k\ k @ y)}. A universal goal plus a scoped atomic
hypothesis covers them all. The scope check is mostly invisible here, since no query lets an eigenvariable
escape. \textbf{Patterns.} The \emph{outputs} \lp{(S k)} of \lp{mixeval}, \lp{(P y)} of \lp{phi}, \lp{(S x)}
of \lp{red1} and the contexts \lp{x\ cond (E x) N P} are pattern equations. Each means ``abstract the
eigenvariable out of the result''. \textbf{Beyond patterns.} \lp{(R U)}, \lp{(R (fixpt R))}, \lp{(E N)}
are not patterns, but in the intended modes their heads are known, so they are plain $\beta$-reductions.
Running them backwards would need general higher-order unification. \textbf{First-order alternatives.} The
book cites Hannan and Miller's transformation of $\beta$-based specifications into closure-based
ones~\cite{HannanMiller1992}. The standard relational miniKanren interpreter with environments and
closures is that transformation applied to \lp{eval}. It runs backwards, but values become closures, which
changes \lp{equal} on functions and rules out \lp{mixeval}. Contexts can be defunctionalized into
first-order frames with a structural plug relation. \textbf{Conditionals.} \lp{if (B = C) ...} becomes a
\scm{cond°} of \scm{=°} and \scm{≠°}, which is \emph{more} relational than the book and works with unbound
arguments. Under the module's default generic reading (\cref{def:roadmap-readings}), however,
\scm{≠°} does decide $k\neq\cst{tt}$ for an eigenvariable $k$, so a \scm{cond°} of \scm{=°} and
\scm{≠°} reproduces the \lp{mixeval} unsoundness. Under \scm{(fresh° (q) ...)},
\begin{scheme}
(∀° (k) (cond° ((=° k 'tt) (=° q 1)) ((≠° k 'tt) (=° q 2))))
\end{scheme}
answers \scm{(2)}. With \scm{(μkanren-λ-reading 'extensional)} the disequality is kept and the
conditional on $k$ fails, answering \scm{()}, as soundness requires. Where a literal translation of a
commit is wanted, \micro{}'s soft-cut \scm{if°} needs \scm{take° 1}. See \cref{sec:module-stores} and
\cref{sec:roadmap-lim-readings}.
\end{kanren}

% =====================================================================================================
\section{A process calculus: the \texorpdfstring{$\pi$}{pi}-calculus}\label{sec:apps-pi}

The $\pi$-calculus~\cite{MilnerParrowWalker1992} has the most delicate binders of the three
applications. Names are bound by input and by restriction, and a restricted name can be sent out of its
scope (\emph{scope extrusion}). Chapter 11 of the book shows that $\lambda$-tree syntax treats this
``naturally'' \book{261}, and it also uses the calculus to expose a limit of the logic.

\subsection{Syntax and one-step transitions}\label{sec:apps-pi-one}

Processes are $P ::= 0\mid P\mid P\mid P+P\mid x(y).P\mid \bar x y.P\mid [x=y].P\mid \tau.P\mid (y)P\mid\ !P$,
where $y$ is bound in $x(y).P$ and in $(y)P$. Free names are constants, and bound names are
$\lambda$-variables \book{262--263}:

\begin{lprolog}
kind name  type.
kind proc  type.
type null        proc.
type plus, par   proc -> proc -> proc.
type in          name -> (name -> proc) -> proc.
type out, match  name -> name -> proc -> proc.
type taup        proc -> proc.
type nu          (name -> proc) -> proc.
type bang        proc -> proc.

type a, b, c     name.
type example     int -> proc -> o.

example 1 (par (in b y\ null) (out b a null)).
example 2 (plus (in b y\ out b a null) (out b a (in b y\ null))).
example 3 (nu x\ par (in x y\ null) (out x a null)).
example 4 (nu x\ out a x null).
example 5 (in a y\ par (in y w\ null) (out b b null)).
example 6 (in a y\ plus (in y w\ out b b null)
                         (out b b (in y w\ null))).
example 7 (nu y\ out a y (par (in y w\ null) (out b b null))).
example 8 (nu y\ out a y (plus (in y w\ out b b null)
                               (out b b (in y w\ null)))).
\end{lprolog}

\noindent For example, example 1 is $(b(y).0)\mid(\bar b a.0)$ and example 3 is
$(x)((x(y).0)\mid(\bar x a.0))$. Actions are \lp{tau}, the free output \lp{up x y} ($\bar x y$) and the free
input \lp{dn x y}. A \emph{bound} action is the partial application \lp{up x} or \lp{dn x}, of type
\lp{name -> action}. There are two transition predicates. \lp{one P A P'} has a free action and a process
as result. \lp{onep P A P'} has a bound action and an \emph{abstracted} continuation of type
\lp{name -> proc}. The book's rules with binding side conditions are \book{264--265}:

\begin{mathpar}
\inferrule{P\xrightarrow{\alpha}P'}{(y)P\xrightarrow{\alpha}(y)P'}\;\rulename{res}\ (y\notin n(\alpha))
\and \inferrule{ }{x(z).P\xrightarrow{x(w)}P\{w/z\}}\;\rulename{input}\ (w\notin fn((z)P))
\and \inferrule{P\xrightarrow{\bar x y}P'}{(y)P\xrightarrow{\bar x(w)}P'\{w/y\}}\;\rulename{open}\ (y\neq x,\ w\notin fn((y)P'))
\and \inferrule{P\xrightarrow{\bar x(w)}P' \\ Q\xrightarrow{x(w)}Q'}{P\mid Q\xrightarrow{\tau}(w)(P'\mid Q')}\;\rulename{close}
\end{mathpar}

\noindent Together with the binder-free rules (prefixes, match, sum, par, and the communication rule com),
they become \book{264}:

\begin{lprolog}
kind action  type.
type tau     action.
type up, dn  name -> name -> action.

type one   proc -> action -> proc -> o.
type onep  proc -> (name -> action) -> (name -> proc) -> o.

one  (taup P)      tau      P.
one  (out X Y P)   (up X Y) P.
onep (in X M)      (dn X)   M.
one  (match X X P) A P' :- one  P A P'.
onep (match X X P) A P' :- onep P A P'.
one  (plus P Q) A P'    :- one  P A P'; one  Q A P'.
onep (plus P Q) A P'    :- onep P A P'; onep Q A P'.
one  (par P Q) A (par P' Q) &
one  (par Q P) A (par Q P') :- one P A P'.
onep (par P Q) A (y\ par (P' y) Q) &
onep (par Q P) A (y\ par Q (P' y)) :- onep P A P'.
one  (nu P) A (nu P')           :- pi y\ one  (P y) A (P' y).
onep (nu P) A (x\ nu y\ P' y x) :- pi y\ onep (P y) A (P' y).
onep (nu P) (up X) P'           :-
        pi y\ one (P y) (up X y) (P' y).
one  (par P Q) tau (nu y\ par (P' y) (Q' y)) &
one  (par Q P) tau (nu y\ par (Q' y) (P' y)) :-
        onep P (up X) P', onep Q (dn X) Q'.
one  (par P Q) tau (par S (T Y)) :- one  P (up X Y) S,
                                    onep Q (dn X)   T.
one  (par P Q) tau (par (S Y) T) :- onep P (dn X) S,
                                    one  Q (up X Y) T.
\end{lprolog}

\begin{tutorial}[where the side conditions went]
\emph{res}: in \lp{one (nu P) A (nu P') :- pi y\ one (P y) A (P' y)}, the variable \lp{A} is quantified
over the whole clause, outside the \lp{pi y}. No legal instance can put \lp{y} into \lp{A}: ``the side
condition associated with the res rule is realized declaratively through a proper nesting of
quantification'' \book{265}. \emph{input} and \emph{open}: the bound name is ``represented by an explicit
abstraction in the abstracted processes that result'' \book{266}. The input clause simply returns the
abstraction \lp{M}, so there is nothing to rename. \emph{open}'s $y\neq x$ holds because \lp{y} is
quantified inside the scope of \lp{X}. \emph{close}: the abstraction over the restricted name lies inside
the scope of \lp{P'} and \lp{Q'}, so the name cannot occur free in them \book{266}. The input
continuation stays abstract until a communication plugs the received name in, at \lp{(T Y)} and \lp{(S Y)}.
This is the late semantics.
\end{tutorial}

\noindent The book animates the specification with queries \book{266--267}:

\begin{lprolog}
?- example 1 P, one P A P'.
P = par (in b W\ null) (out b a null)
A = up b a
P' = par (in b W\ null) null;
P = par (in b W\ null) (out b a null)
A = tau
P' = par null null;
no

?- example 1 P, onep P A P'.
P = par (in b W\ null) (out b a null)
A = W\ dn b W
P' = W\ par null (out b a null);
no
\end{lprolog}

\noindent For example 3 the query for \lp{one} has a single answer, \lp{tau} to
\lp{nu W\ par null null} \book{267}: every visible action would mention the restricted name. Our notes also derive, by the same reading of the clauses, that example 4,
$(x)(\bar a x.0)$, has \emph{no} free transition (sending $x$ would leak it) but has the bound output
\lp{A = W\ up a W}, \lp{P' = W\ null} by \emph{open}.

\begin{tutorial}[three transitions, clause by clause]
\emph{The $\tau$ of example 1.} The goal is \lp{one (par (in b y\ null) (out b a null)) tau P'}. The
first communication clause needs the left component to output, which an input cannot do. The second,
\lp{one (par P Q) tau (par (S Y) T) :- onep P (dn X) S, one Q (up X Y) T}, works:
\lp{onep (in b y\ null) (dn X) S} matches the input clause with \lp{X = b} and \lp{S = y\ null}; then
\lp{one (out b a null) (up b Y) T} matches the output clause with \lp{Y = a}, \lp{T = null}. The result
\lp{par (S Y) T} is \lp{par ((y\ null) a) null}, which $\beta$-reduces to \lp{par null null}. The
received name was plugged in only now, by the application \lp{(S Y)}: this is the late semantics.


\smallskip\noindent\emph{Why example 3 has no visible action.} The \emph{res} clause asks
\lp{pi y\ one (P y) A (P' y)}, with \lp{P y = par (in y z\ null) (out y a null)}. The output component
offers \lp{up y a}, so \lp{A = up y a} would be needed. But \lp{A} was created \emph{before} the
\lp{pi y}, and the unifier refuses to bind it to a term that contains \lp{y}. Only the internal
communication, whose label \lp{tau} does not mention \lp{y}, survives.


\smallskip\noindent\emph{Scope extrusion in example 4.} \lp{one (nu x\ out a x null) A P'} fails for the same reason
(\lp{A = up a y} is out of scope). The bound-output clause for \emph{open},
\lp{onep (nu P) (up X) P' :- pi y\ one (P y) (up X y) (P' y)}, succeeds instead: the action pattern
\lp{up X y} already contains the fresh \lp{y}, so only \lp{X = a} is needed, and
\lp{(P' y) = null} is a pattern equation with the vacuous solution \lp{P' = W\ null}. The extruded
name is carried by the abstraction in the result, not by a free constant.
\end{tutorial}

\subsection{Traces, may and must}\label{sec:apps-pi-traces}

Traces record sequences of actions. Bound outputs are continued with a \emph{generic} name, while bound
inputs are continued with a \emph{logic variable}, ``for an input action, we would want to consider
different possible instantiations'' \book{268}:

\begin{lprolog}
kind trace  type.
type empty  trace.
type tr     action -> trace -> trace.
type trp    (name -> action) -> (name -> trace) -> trace.

type trace  proc -> trace -> o.
trace P empty.
trace P (tr Act Tr)     :- one P Act Q, trace Q Tr.
trace P (trp (up X) Tr) :- onep P (up X) Q,
                           pi x\ trace (Q x) (Tr x).
trace P (tr (dn X Y) Tr) :- onep P (dn X) Q, trace (Q Y) Tr.
\end{lprolog}

\noindent For example 1 the query \lp{?- example 1 P, trace P Tr.} produces six traces \book{268}:
\begin{lprolog}
empty
tr (up b a) empty
tr (up b a) (tr (dn b T) empty)
tr tau empty
tr (dn b T) empty
tr (dn b T) (tr (up b a) empty)
\end{lprolog}
\noindent (The book prints \lp{trp (dn b T)} in the last two lines. Its own clauses produce
\lp{tr (dn b T)}, since only bound \emph{outputs} build \lp{trp}, and its later \lp{comptrace} answers
use \lp{tr (dn b T1)}.) Each input leaves a logic variable \lp{T}: the trace says ``some name was
received''. The logic variable can be \emph{narrowed} later by \lp{match}. The query below enumerates exactly
four traces \book{268--269}:
\begin{lprolog}
?- trace (in a Y\ plus (match Y b (out Y Y null))
                       (match Y c (out Y Y null))) Tr.

empty
tr (dn a T) empty
tr (dn a b) (tr (up b b) empty)
tr (dn a c) (tr (up c c) empty)
\end{lprolog}
\noindent After the input, the continuation is \lp{plus (match T b (out T T null)) (match T c ...)}.
The clause \lp{one (match X X P) A P'} can fire only if its two arguments unify, so it binds \lp{T := b}
in one branch and \lp{T := c} in the other. Had the input been continued with a \emph{generic} name
(\lp{pi}), both matches would have failed and the last two traces would be lost: ``this process can do
an unrestricted input and \ldots\ if that input is chosen to be b or c, then further, distinct output
actions can follow'' \book{269}.

\begin{definition}[may- and must-judgments; \book{269}]
A \emph{may-judgment} asserts the existence of a computation path, so it corresponds to the existence of a
proof. A \emph{must-judgment} quantifies over all behaviours: a process can never do some action, all
traces of $P$ are traces of $Q$, $P$ is simulated by $Q$.
\end{definition}

\noindent ``The style of logic specifications that we have described in this book does not, in general,
provide declarative treatments of must-judgments. Negation-as-failure can be used, however, to capture or
approximate some of these judgments'' \book{269}. Negation-as-failure is the cut-based predicate of
\book{137}:
\begin{lprolog}
type not  o -> o.
not P :- P, !, fail.
not P.
\end{lprolog}
\noindent With it the book defines terminal processes, complete traces, separating traces and trace
equivalence \book{270}:
\begin{lprolog}
type possible, terminal  proc -> o.
type comptrace           proc -> trace -> o.
type separating_trace    proc -> proc -> trace -> o.
type trace_equiv         proc -> proc -> o.

possible P :- one P _ _ ; onep P _ _.
terminal P :- not (possible P).

comptrace P empty :- terminal P.
comptrace P (tr Act Tr)      :- one P Act Q, comptrace Q Tr.
comptrace P (tr (dn X Y) Tr) :- onep P (dn X) P',
                                comptrace (P' Y) Tr.
comptrace P (trp (up X) Tr)  :- onep P (up X) P',
                                pi x\ comptrace (P' x) (Tr x).

separating_trace P Q T :- trace P T,
                          not (trace Q T).
trace_equiv P Q :- not (separating_trace P Q _),
                   not (separating_trace Q P _).
\end{lprolog}
\noindent A complete trace is one that cannot be extended: it may end with \lp{empty} only at a process
that can do nothing. Each of the four predicates hides a ``for all'' inside a \lp{not}: \lp{terminal P}
says that \emph{no} transition exists, \lp{trace_equiv} that \emph{no} trace separates. Examples 5 and 6
are separated, in one direction only, by an input of \lp{b} followed by an internal step
\book{270--271}:
\begin{lprolog}
?- example 5 P, example 6 Q, separating_trace P Q T.
T = tr (dn a b) (tr tau empty)
\end{lprolog}
\noindent Example 5 is $a(y).((y(w).0)\mid(\bar b b.0))$. After the input, its two parallel components
can communicate, but only on a common channel, so the communication clause unifies the received logic
variable with \lp{b}. Example 6, $a(y).((y(w).\bar b b.0)+(\bar b b.y(w).0))$, offers the same two
actions only in sequence, inside a choice, and can never make the internal step. The query in the other
direction answers \lp{no}. Examples 7 and 8 replace the input prefix by a bound output. They have no
separating trace in either direction \book{271}, because the extruded name is generic and can never equal
\lp{b}.

\paragraph{An incorrect simulation.} Trace equivalence is too coarse: $\bar x a.(\bar y b.0+\bar y c.0)$
and $(\bar x a.\bar y b.0+\bar x a.\bar y c.0)$ are trace-equivalent, but after $\bar x a$ the first can
still choose between two outputs and the second cannot \book{271}. The finer notion is simulation: ``$P$
is simulated by $Q$ if whenever $P$ can perform an action and evolve into a process $P'$, $Q$ can perform
an identical action that transforms it into a process $Q'$ that simulates $P'$'' \book{271}, read
coinductively. The book attempts it with negation-as-failure and labels its own attempt
\emph{incorrect} \book{272}:
\begin{lprolog}
type foreach2  (A -> B -> o) -> (A -> B -> o) -> o.
type sim       proc -> proc -> o.

foreach2 P Q :- not (sigma X\ sigma Y\ P X Y, not (Q X Y)).

sim P Q :- foreach2 (A\P'\ one P A P')
                    (A\P'\ sigma Q'\ one Q A Q', sim P' Q'),
           foreach2 (A\P'\ onep P A P')
                    (A\P'\ sigma Q'\ onep Q A Q',
                           pi x\ sim (P' x) (Q' x)).
\end{lprolog}
\noindent The query
\begin{lprolog}
?- sim (in a x\ par (in x y\ null) (out c b null))
       (in a x\ plus (in x y\ out c b null) (out c b (in x y\ null))).
solved
\end{lprolog}
\noindent answers \lp{solved}, which is wrong. After receiving \lp{c} on \lp{a}, the first
process can do a $\tau$ and the second cannot. Diagnosis: ``the universal quantifier in goals \ldots has an
\emph{intensional} character: The generic goal must be derived by showing that an \emph{identical}
derivation exists for any of its instances. However, $\pi$-calculus processes may have different evolutions
for different instantiations of an input action'' \book{272}. What is needed is ``a universal quantifier that
has an \emph{extensional} interpretation''. The bibliographic notes point to logics with
definitions~\cite{McDowellMiller2000}, where $\forall$ is extensional, together with a separate
\emph{generic} quantifier $\nabla$ for names~\cite{MillerTiu2005} (\book{275}).

\begin{pitfall}[one \texttt{pi}, two jobs]
In $\lambda$Prolog, \lp{pi} serves both as ``a fresh name'' (binders, where it is right) and as ``for all
instances'' (must-judgments, where it is wrong). Negation-as-failure is ``at best an ad hoc device'' with
uninstantiated variables \book{273}, and it needs finite failure, which replication can destroy.
\end{pitfall}

\subsection{Replication, and the \texorpdfstring{$\lambda$}{lambda}-calculus inside \texorpdfstring{$\pi$}{pi}}\label{sec:apps-lampi}

Replication $!P$, ``the parallel composition of any number of copies of $P$'' \book{262}, adds four
clauses \book{273}:
\begin{lprolog}
one  (bang P) A (par P1 (bang P)) :- one P A P1.
onep (bang P) X (y\ par (M y) (bang P)) :- onep P X M.
one  (bang P) tau (par (par R (M Y)) (bang P)) :-
     onep P (dn X) M, one P (up X Y) R.
one  (bang P) tau (par (nu y\ par (M y) (N y)) (bang P)) :-
     onep P (up X) M, onep P (dn X) N.
\end{lprolog}
\noindent A copy acts with a free action, a copy acts with a bound action, two copies communicate a free
name, and two copies communicate an extruded one. In each case a fresh \lp{bang P} remains. Milner's call-by-name encoding~\cite{Milner1990} sends arguments over a channel $u$
\book{273}:
\begin{align*}
  [\![x]\!](u) &= \bar x u.0 \qquad\qquad
  [\![\lambda x M]\!](u) = u(x).u(v).[\![M]\!](v)\\
  [\![M\,N]\!](u) &= (v).([\![M]\!](v)\mid(x).(\bar v x.\bar v u.!x(w).[\![N]\!](w)))
\end{align*}
\begin{lprolog}
type trans  tm -> (name -> proc) -> o.

trans (abs M) (u\ in u x\ in u v\ P x v) :-
    pi x\ pi y\ trans x (u\ out y u null) => trans (M x) (P y).

trans (app M N)
      (u\ nu v\ par (P v)
                    (nu x\ out v x (out v u (bang (in x Q))))) :-
    trans M P, trans N Q.
\end{lprolog}
\noindent As in \lp{phi}, the base case for variables is assumed as a local clause when the recursion
enters a binder. The book checks the encoding by computing complete traces \book{274}:
\begin{lprolog}
?- trans (abs w\ w) P, comptrace (P b) T.
P = u\ in u x\ in u y\ out x y null
T = tr (dn b T1) (tr (dn b T2) (tr (up T1 T2) empty));
no
\end{lprolog}
\noindent For \lp{(app (abs w\ w) (abs w\ w))} the complete trace is
\lp{tr tau (tr tau (tr tau (tr (dn b T1) (tr (dn b T2) (tr (up T1 T2) empty)))))}: three silent steps
``corresponding to reduction of the term $(\lambda x\,x)(\lambda x\,x)$ to $\lambda x\,x$'', then the trace of the
identity \book{275}. (The printed process term for this query shadows a name and swaps the two \lp{par}
components, but the trace is unaffected.)

\begin{remark}[$\Llam$ and $\pi$I; \book{275}]
``Almost all the one-step transition rules \ldots are in the $\Llam$ fragment''. The exceptions are the last
two clauses of the \lp{one}/\lp{onep} program and the third replication clause, where an abstraction is applied
to a received logic variable (\lp{T Y}, \lp{S Y}, \lp{M Y}). Deleting them gives Sangiorgi's
$\pi$I~\cite{Sangiorgi1996}, in which ``only new names are bound to inputs'', so the metalogic needs only
$\beta_0$-conversion.
\end{remark}

\begin{kanren}[what the $\pi$-calculus needs]
\textbf{$\forall$ with a real scope check.} Here the scope check is the whole point, because \emph{res},
\emph{open} and \emph{close} hold only because outer variables cannot capture the fresh name. Without it,
example 4 would perform the free output $\bar a x$ and example 3 would perform \lp{up W a}. A cheaper
discipline, posting an \scm{absent°} constraint on the clause's own variables after each generic subgoal,
is \emph{not} enough. Our notes give the test \lp{trace (in a z\ nu x\ match z x (out b b null)) Tr}. In
$\lambda$Prolog it yields only \lp{empty} and \lp{tr (dn a T) empty}, because the received \lp{T}, a
variable hidden inside the process, may not become the restricted \lp{x}. An absent-on-clause-variables
port also yields \lp{tr (dn a T) (tr (up b b) empty)}. The check must live in unification.
\textbf{$\imp$.} \lp{trans} assumes atomic clauses about fresh variables, as in \cref{sec:apps-fp}.
\textbf{Patterns.} Abstracted continuations, \lp{(P' y)} under \lp{pi y}, and the trace abstraction
\lp{(Tr x)} are all pattern equations. \textbf{Beyond patterns.} \lp{(T Y)} with \lp{Y} a logic variable is
outside $\Llam$. A faithful engine postpones it until \lp{T} is known. A first-order locally nameless
encoding turns it into a harmless structural ``open'' relation. \textbf{Negation.} \lp{not G} is
\scm{(if° (take° 1 G) ✗° ✓°)}. Without \scm{take° 1}, \micro{}'s soft cut would map \scm{✗°} over every
answer of \lp{G} and diverge whenever \lp{G} has infinitely many. \textbf{Search.} With \lp{bang},
\lp{trace} becomes an infinite but productive stream under interleaving, but negation still needs finite
failure. \textbf{Must-judgments.} Our notes propose, beyond the book, to make the bound-input case of
\lp{sim} extensional by case analysis over the free names plus one fresh name. This is a finite substitute
for the definitions-and-$\nabla$ reading. See \cref{ch:design} for which of these the proposal adopts.
\end{kanren}

% =====================================================================================================
\section{Summary}\label{sec:apps-summary}

\begin{itemize}[nosep]
  \item \textbf{One recipe for three domains.} A judgment becomes a predicate. Object binders become meta
        $\lambda$s. Going under a binder means \lp{pi} plus a hypothetical fact. Freshness and non-capture
        side conditions are then discharged by quantifier nesting.
  \item \textbf{Specification versus search.} Direct translations are sound, but depth-first execution can
        loop (${\imp}L$, \lp{#} as a prover). A better calculus (Dyckhoff's) gives a decision procedure.
        Programmable control (tactics and tacticals) moves the strategy into the program.
  \item \textbf{Logic variables as delayed choices.} Witnesses of $\exists R$, \lp{all_l} and
        \lp{path_rec_tac}, and received names in traces, are chosen by unification at the leaves. The unifier
        must respect eigenvariable scope.
  \item \textbf{Cut and \texttt{if} are dangerous with logic variables and eigenvariables.} They can lose
        completeness (the CL counterexample) or break soundness theorems (\lp{mixeval}).
  \item \textbf{Must-judgments are beyond hohh.} \lp{pi} is intensional. Simulation needs an extensional
        $\forall$ (definitions) and a generic $\nabla$.
  \item \textbf{Interleaving search} makes naive calculi, prover modes and tactic scripts complete for
        answers, and it makes the classical prover complete without iterative deepening. It never provides
        termination.
\end{itemize}

\noindent The requirements, application by application (``atomic'' means every assumed clause is an atomic
fact, so an explicit context could replace \lp{=>}):

\begin{center}
\small
\begin{tabularx}{\linewidth}{@{}>{\raggedright\arraybackslash}p{3.1cm}>{\raggedright\arraybackslash}X>{\raggedright\arraybackslash}X>{\raggedright\arraybackslash}X@{}}
\toprule
application & $\forall$ (eigenvariables) & $\imp$ (hypotheses) & $\lambda$-term unification\\
\midrule
Dyckhoff \texttt{seq} & only in the query & atomic, query only & none: \textbf{runs on \micro{} today}\\
natural deduction \texttt{\#} & proof labels, $\forall I$, $\exists E$ & atomic (\texttt{p \# A}) & patterns for completion, \mbox{\texttt{A' (f y)}} delayed\\
classical \texttt{prv} & $\forall R$, scope check vital & atomic (\texttt{lit}, \texttt{bound}) & patterns (difference-list queue)\\
tactics & \texttt{maptac} under \texttt{allg} & none & pattern \texttt{Out t}; rest \textbf{runs today}\\
\texttt{typeof} & under \texttt{abs}, \texttt{fixpt} & atomic & none with named binders: \textbf{runs today}\\
\texttt{eval}, \texttt{evalc} & none & none & $\beta$ with known heads; contexts as outputs\\
\texttt{mixeval}, CPS & under the binder & atomic & patterns \texttt{(S k)}, \texttt{(P y)}\\
$\pi$ transitions, traces & \emph{res}/\emph{open}/\emph{close}, scope check vital & atomic (\texttt{trans}) & patterns, plus non-$\Llam$ \mbox{\texttt{(T Y)}}\\
\bottomrule
\end{tabularx}
\end{center}


\part{Higher-order miniKanren}
\chapter{Higher-order miniKanren: the literature}\label{ch:lkanren}

\providecommand{\lkclo}[2]{\langle #1;\,#2\rangle}
\providecommand{\lkq}{\stackrel{?}{=}}
\providecommand{\lkJust}{\mathsf{Just}}
\providecommand{\lkNothing}{\mathsf{Nothing}}
\providecommand{\lkyes}{{\ttfamily ✓}}
\providecommand{\lkno}{{\ttfamily ✗}}

The previous chapters described what λProlog is. It has a logic, higher-order hereditary Harrop
formulas (\cref{ch:hh}). It has a term language, simply typed λ-terms used as λ-tree syntax
(\cref{ch:lambda,ch:lambdatree}). It has a unifier, Huet's procedure or its decidable core, pattern
unification (\cref{ch:hounif}). And it has a style of programming (\cref{ch:apps}). This chapter
changes the point of view. Several people have tried to put these ideas into a miniKanren, and
several systems implement λProlog efficiently. What did they actually do, and what went wrong?

The user named one paper as the key reference: Ma and Friedman, \emph{A New Higher-order
Unification Algorithm for λKanren} \cite{MaFriedman2021}. We present it in full in
\cref{sec:lkanren-mf21}, followed by a critical assessment against pattern unification and nominal
unification (\cref{sec:lkanren-assess}). Then we survey the rest of the literature. The earlier
λKanren paper \cite{MaLuFriedman2020} comes first (\cref{sec:lkanren-mk20}). Then nominal unification
and αKanren \cite{UrbanPittsGabbay2004,ByrdFriedman2007} (\cref{sec:lkanren-nominal}). Then
eigenvariables, negation and ∀ inside miniKanren
\cite{HemannByrdFriedman2013,Moiseenko2019,JinEtAl2021} (\cref{sec:lkanren-eigen}). Last come the
techniques used by real λProlog systems, Teyjus, ELPI and Makam (\cref{sec:lkanren-impl}). The chapter
ends with a comparison table of the design options (\cref{sec:lkanren-compare}). It does
\emph{not} choose between them: the choice is made in \cref{ch:design}.

\section{The map of the territory}\label{sec:lkanren-map}

The book gives the logic and an operational semantics. It has the search rules \rulename{and},
\rulename{augment}, \rulename{generic} and \rulename{backchain}, plus unification of λ-terms. A
miniKanren that wants to host this must answer six concrete questions. The rest of the chapter is
organized around them.

\begin{enumerate}[label=\textbf{Q\arabic*.}]
  \item \textbf{Representation of λ-terms.} The options are named binders, de Bruijn indices, de
    Bruijn levels, locally nameless terms, nominal atoms, or explicit substitutions. The goal is to
    make α-equivalence cheap and to avoid β as far as possible.
  \item \textbf{Which unification runs inside \scm{=°}.} There are five candidates: first-order (today),
    nominal (αKanren), Miller's pattern unification $\Llam$ (λKanren 2020, Teyjus~2, ELPI), Huet's full
    pre-unification, or the \enquote{static closures} of Ma and Friedman 2021.
  \item \textbf{How to implement \rulename{generic} ($\forall x.G$).} You need fresh eigen-constants
    plus a \emph{scope check}. The check can use an explicit quantifier prefix
    \cite{Miller1991,MaLuFriedman2020}, numeric universe indices or levels
    \cite{NadathurJayaramanKwon1995,DunchevEtAl2015}, or an absento-like constraint over the
    variables in scope \cite{HemannByrdFriedman2013}.
  \item \textbf{How to implement \rulename{augment} ($D\imp G$).} A hypothetical clause must be visible
    to every relation call made during the proof of $G$, and only there. In a shallow embedding this
    clashes with the lexical scope of host-language procedures.
  \item \textbf{What to do outside the decidable fragment.} You can fail, delay the equation as a
    constraint, or enumerate unifiers.
  \item \textbf{How the new features meet the other constraints} \scm{≠°}, \scm{absent°} and the
    type tags. Eigen-constants and disequality do not mix naively.
\end{enumerate}

The historical arc, in the notes' order: Miller 1991 ($\Llam$, mixed prefixes, raising and pruning)
\cite{Miller1991}; Nipkow 1993 (functional, de Bruijn) \cite{Nipkow1993}; Nadathur, Jayaraman and Kwon
1995 (universe indices) \cite{NadathurJayaramanKwon1995}; Dowek, Hardin, Kirchner and Pfenning 1996
(explicit substitutions) \cite{DowekHardinKirchnerPfenning1996}; Urban, Pitts and Gabbay 2004 (nominal
unification) \cite{UrbanPittsGabbay2004}; Nadathur and Linnell 2005 (on-the-fly raising)
\cite{NadathurLinnell2005}; αKanren 2007 \cite{ByrdFriedman2007}; ELPI 2015 \cite{DunchevEtAl2015};
then the miniKanren papers of 2019--2021
\cite{Moiseenko2019,MaLuFriedman2020,MaFriedman2021,JinEtAl2021}.

% =====================================================================================================
\section{Ma and Friedman 2021: static closures}\label{sec:lkanren-mf21}

\subsection{What kind of paper this is}

The paper is a ten-page workshop paper, presented at the miniKanren Workshop 2021. Its abstract states
the thesis (p.\,1):
\begin{quote}
\enquote{λKanren, adapted from λProlog, extends miniKanren with higher-order logic programming.
The existing higher-order unification algorithms, however, are neither sound nor complete w.r.t.\ an
“intuitive” understanding of α-equivalence. This paper shows these gaps and proposes a new method to
handle α-equivalence in higher-order unification.}
\end{quote}
It opens provocatively: \enquote{Is there a world where true is equal to false? Indeed, there is, if
higher-order unification and miniKanren meet each other} (p.\,1).

The background section recalls the one notion from \cref{ch:hounif} that the whole paper argues
about, in its own words (p.\,1): \enquote{For an application form $Fa$, if $F$ is a unification
variable, then $a$ must be a list of zero or more distinct bound variables. Such restricted application
forms are called \emph{patterns} and we hereafter refer to the problem identified by Miller as
\emph{pattern unification}.} The related work (p.\,7--8) adds Huet's vocabulary: a term is \emph{rigid}
\enquote{if its head symbol is a bound variable or a constant}, and \emph{flex} \enquote{if its head symbol
is a unification variable}. Keep in mind one fact for later: in every query of §2--4, the unknown
$X$ sits under a λ but is \emph{not} applied to anything. It is a pattern with zero arguments, so
pattern unification lets it depend on none of the bound variables.

\begin{pitfall}[a preview, not an algorithm]
The conclusion says: \enquote{This paper shows a preview of static closures} (p.\,9). The paper has
\textbf{no algorithm listing, no inference rules for unification, no theorems and no proofs}. Its
formal content is two grammars, three operations on scopes and three worked examples. Its code is a
set of REPL transcripts. Soundness and completeness are explicitly future work: \enquote{We
anticipate extending the formulation of higher-order unification with static closures and derive a
sound and complete algorithm for α-equivalence} (p.\,9). The rules in \cref{sec:lkanren-recon} below
are our own \emph{reconstruction}, not the paper's, and we flag them as such. (Where this chapter
says \enquote{the notes}, it means the unpublished working notes from which this survey was written.)
\end{pitfall}

The paper has seven sections. §1 gives background. §2 asks \enquote{grafting or substitution?}
(p.\,2--4). §3 shows how one can \enquote{turn false to true} (p.\,4). §4 introduces static closures
(p.\,5--6). §5 describes a shallow-embedding implementation of the language (p.\,6--7). §6 is related
work (p.\,7--8) and §7 the conclusion. The paper uses two box styles for examples: \enquote{solid
boxes for examples that are complete and dash boxes for examples that are incomplete or somewhat
questionable} (footnote 2, p.\,2). In every \scm{(run 1 X ...)} of the paper, \scm{X} is simply
existentially quantified at the top, outside every λ (footnote 1, p.\,2).

\subsection{Grafting versus substitution}\label{sec:lkanren-graft}

First-order miniKanren already has a word for what unification produces. For

\begin{scheme}
> (run 1 X
    (== X 'cat))
'(cat)
\end{scheme}

\noindent the paper says \enquote{internally, the answer is a grafting \dots\ written as $\{X\mapsto
\texttt{'cat}\}$} (p.\,2).

\begin{definition}[Grafting and substitution {\cite[p.\,2--3]{MaFriedman2021}}]\label{def:lkanren-graft}
A \emph{grafting} $\{X\mapsto t\}$ is a purely textual replacement of the unification variable $X$
by $t$. It is \emph{not} capture-avoiding: grafting $X\mapsto a$ into $\abs{a}{X}$ yields
$\abs{a}{a}$, so the name $a$ gets captured. A \emph{substitution} is the capture-avoiding
instantiation of standard higher-order unification: \enquote{Standard higher-order unification relies
on substitution, but not grafting, to instantiate unification variables. Substitution is
capture-avoiding, whereas grafting is not} (p.\,3).
\end{definition}

In first-order unification the two notions coincide, because there are no binders. They part ways
as soon as terms have binders. The paper follows a \enquote{novice} through three queries. The first
one, α-equivalence of closed terms, is unproblematic:

\begin{scheme}
> (run 1 X
    (== (λ (a) a) (λ (b) b)))
'(_0)
\end{scheme}

\noindent (\texttt{\_0} means that $X$ is unrestricted). Then, \enquote{out of curiosity}, the novice
puts the query variable under a binder and asks \enquote{Does this query succeed?}:

\begin{scheme}
> (run 1 X
    (== (λ (a) X) (λ (b) b)))
'()
\end{scheme}

\noindent \enquote{Surprisingly enough, with higher-order unification, there is no answer.} The novice
expects $X = a$. The reason is the \emph{scope restriction}.

\begin{definition}[Scope restriction {\cite[p.\,3]{MaFriedman2021}}]\label{def:lkanren-scope}
\enquote{Unification variables are ∃-quantified and lambda variables are ∀-quantified. A scope is a
list of those quantifiers. Further, a ∀-quantified variable may instantiate (or occurs in the
instantiation of) an ∃-quantified variable, only if the ∀ variable precedes the ∃ variable in the
scope.} As a side condition on a binding in a mixed prefix $Q$:
\[
  \inferrule{X \text{ is } \exists\text{-quantified in } Q \\ \forall y\in FV_\forall(t).\; y \prec_Q X}
            {X\mapsto t \text{ is allowed}}
\]
\end{definition}

For \lp{λa.X =? λb.b} the scope is \scm{'((∃ X) (∀ a))}. The variable $a$ comes \emph{after} $X$, so
it cannot instantiate $X$. This is exactly λProlog's behaviour: it is the eigenvariable condition of
\rulename{generic} (\cref{ch:hh}), and in unification it is what raising and pruning enforce
(\cref{ch:hounif}). The paper draws the opposite conclusion (p.\,3): \enquote{It is fair to say that
the existing pattern unification algorithms are not complete w.r.t.\ the natural sense of
α-equivalence.} It also says (p.\,3): \enquote{For beginners who are not familiar with higher-order
unification, α-equivalence is only decided by the binder positions. But pattern unification imposes
an appendage limitation by taking unification variables into account.}

A conjunction shows why the restriction has a reason to exist (dashed box, p.\,3):

\begin{scheme}
> (run 1 X
    (== (λ (a) X) (λ (b) b))
    (== (λ (c) X) (λ (d) d)))
\end{scheme}

\noindent \enquote{Associating X with a cannot work, neither does c. If we text-replace X with either,
then the other one fails.} The paper's first proposal is to stop treating λ-variables as names:
\enquote{We first propose to treat unification variables and lambda variables differently. In
particular, lambda variables are effectively constants if they are represented by de Bruijn levels}
(p.\,3).

\begin{definition}[de Bruijn levels and indices {\cite[footnote 3, p.\,4]{MaFriedman2021}}]\label{def:lkanren-levels}
\enquote{There are two ways to implement de Bruijn's nameless syntax, indices and levels. Both use
numbers to represent variables} \cite{deBruijn1972}. A \emph{level} counts binders from the
\emph{outside}: the outermost binder is $0$. An \emph{index} counts binders from the \emph{inside}:
the nearest binder is $0$.
\end{definition}

\begin{example}
In $\abs{c}{\abs{d}{c}}$, the variable $c$ has level $0$ and index $1$, and $d$ has level $1$ and
index $0$. The term is $\lambda.\lambda.\,0$ with levels and $\lambda.\lambda.\,1$ with indices. The
paper writes the level $\ell$ as \texttt{d}$\ell$. So \texttt{d0} names the outermost enclosing
binder, whatever it is called, and the answer \texttt{'(d0)} reads: \enquote{these two terms can be made
α-equal, if X corresponds to de Bruijn level 0} (p.\,3).
\end{example}

With levels, the two problematic queries succeed (solid boxes, p.\,3):

\begin{scheme}
> (run 1 X
    (== (λ (a) X) (λ (b) b)))
'(d0)
> (run 1 X
    (== (λ (a) X) (λ (b) b))
    (== (λ (c) X) (λ (d) d)))
'(d0)
\end{scheme}

\noindent \enquote{If we instantiate X with d0 and reify these levels w.r.t.\ their binders, we have
the problem solved with no surprise}: the two equations become \scm{(== (λ (a) a) (λ (b) b))} and
\scm{(== (λ (c) c) (λ (d) d))}. The paper goes on (p.\,3--4): \enquote{using grafting with de Bruijn
levels captures the essence of α-equivalence: names don't matter. A unification algorithm based on de
Bruijn levels would also be more straightforward because it's practically first-order. But, we cannot
yet lift the scope restriction on lambda variables and drop in de Bruijn levels, because sometimes
names do matter and \dots}. The sentence continues in the title of §3.

\begin{kanren}[(aux kanren micro) already grafts]
Grafting is what a first-order unifier does on nameless terms. Today's \micro{} unifies records slot
by slot, so a \scm{define-record} is a new term constructor for free. If we write λ-terms with de
Bruijn indices as records, \scm{=°} grafts. It reproduces the paper's answers, and also the problems
discussed below. The demonstration uses indices, not levels; for the examples here only the
outermost binder matters, and it is index~$0$ under one λ.
\begin{scheme}
(import scheme (chicken base) (chicken pretty-print)
        (aux base) (aux stream) (aux kanren micro))

(define-syntax-rule (show label expr)
  (begin (display ";; ") (display label) (newline)
         (pretty-print expr)))

;; 1. Named binders as plain lists: syntactic unification ignores alpha.
(show "named: (λ a a) = (λ b b)"
  (°->list #f (fresh° (q) (=° '(λ a a) '(λ b b)))))

;; 2. Nameless binders as records with de Bruijn indices.
(define-record lam body)
(define-record dbi n)
(define (λ. t) (make-lam t))
(define (ix n) (make-dbi n))

(show "de Bruijn: λ.0 = λ.0"
  (°->list #f (fresh° (q) (=° (λ. (ix 0)) (λ. (ix 0))))))
(show "grafting: λa.X = λb.b"
  (°->list #f (fresh° (x) (=° (λ. x) (λ. (ix 0))))))
(show "grafting: λa.X = λb.b  and  λc.X = λd.d"
  (°->list #f (fresh° (x) (=° (λ. x) (λ. (ix 0)))
                          (=° (λ. x) (λ. (ix 0))))))
(show "true = false: λa.λa.X = λc.λd.c"
  (°->list #f (fresh° (x) (=° (λ. (λ. x)) (λ. (λ. (ix 1)))))))
(show "dangling: λa.X = λb.b, then Q = X at top level"
  (°->list #f (fresh° (q) (fresh° (x) (=° (λ. x) (λ. (ix 0)))
                                      (=° q x)))))
\end{scheme}
\begin{answers}
;; named: (λ a a) = (λ b b)
()
;; de Bruijn: λ.0 = λ.0
((λ (α) α))
;; grafting: λa.X = λb.b
((λ () (make-record-instance 'dbi 0)))
;; grafting: λa.X = λb.b  and  λc.X = λd.d
((λ () (make-record-instance 'dbi 0)))
;; true = false: λa.λa.X = λc.λd.c
((λ () (make-record-instance 'dbi 1)))
;; dangling: λa.X = λb.b, then Q = X at top level
((λ () (make-record-instance 'dbi 0)))
\end{answers}
The first query shows that named binders as plain data are useless: α-equivalent terms do not
unify. The next three reproduce the paper's \texttt{'(\_0)}, \texttt{'(d0)}, \texttt{'(d0)}. The fifth
is the \enquote{true = false} problem of the next subsection. The answer \texttt{(dbi 1)} is
perfectly correct for the nameless terms, but it denotes the \emph{outer} binder. In the user's
named term $\abs{a}{\abs{a}{X}}$ that binder is shadowed. The last query shows a \emph{dangling}
index: $Q$ lives at top level, outside every binder, yet it is bound to \enquote{the nearest
binder}. \Cref{sec:lkanren-assess} returns to it.
\end{kanren}

\subsection{Binder-merging and \enquote{true = false}}\label{sec:lkanren-merge}

\begin{definition}[Binder-merging {\cite[p.\,4]{MaFriedman2021}}]\label{def:lkanren-merge}
\enquote{When two terms are unified, one term's binders are replaced by the other's. Then
α-equivalence is as simple as structural equivalence.} The paper says that it is \enquote{used in
all pattern unification algorithms, except for those with de Bruijn numbers}.
\end{definition}

\begin{example}[binder-merging at work, p.\,4]
\[
  \abs{x}{\abs{y}{X\,x\,y}} \lkq \abs{a}{\abs{b}{a\,b}}
  \;\longrightarrow\;
  \abs{a}{\abs{b}{X\,a\,b}} \lkq \abs{a}{\abs{b}{a\,b}}
\]
This is solved by $\{X\mapsto \abs{a}{\abs{b}{a\,b}}\}$. Instantiating gives
$\abs{a}{\abs{b}{((\abs{a}{\abs{b}{a\,b}})\,a\,b)}}$, which β-reduces to $\abs{a}{\abs{b}{a\,b}}$.
\end{example}

Now take a binder name that is reused, i.e.\ \emph{shadowing} (dashed box, p.\,4):

\begin{scheme}
> (run 1 X
    (== (λ (a) (λ (a) X)) (λ (c) (λ (d) c))))
\end{scheme}

\noindent \enquote{With either binder-merging or de Bruijn syntax, the unification problem is
effectively reduced to} $\abs{c}{\abs{d}{X}}\lkq\abs{c}{\abs{d}{c}}$. \enquote{Then it appears to be
solvable by $\{X\mapsto c\}$, or equivalently $\{X\mapsto \texttt{d0}\}$, that associates $X$ with the
outer-binder.} Naive reification, which writes each level back as the name of its binder, then gives

\begin{scheme}
(== (λ (a) (λ (a) a)) (λ (c) (λ (d) c)))
\end{scheme}

\noindent \enquote{which claims that the λ-encoded false is true!} With Church booleans
$\cst{true}=\abs{t}{\abs{f}{t}}$ and $\cst{false}=\abs{t}{\abs{f}{f}}$, the left-hand side
$\abs{a}{\abs{a}{a}}$ is α-equivalent to $\cst{false}$, because the occurrence $a$ refers to the inner
binder. The right-hand side is $\cst{true}$.

The paper argues that binder-merging goes wrong \enquote{even without lifting the scope
restrictions}:
\[
  \abs{x}{\abs{y}{x\,y}} \lkq \abs{a}{\abs{a}{a\,a}}
  \quad\text{is converted to}\quad
  \abs{a}{\abs{a}{a\,a}} \lkq \abs{a}{\abs{a}{a\,a}},
\]
\enquote{which then ends up with another false positive}. The correct answer is failure:
$\abs{a}{\abs{a}{a\,a}}\equiv_\alpha\abs{a}{\abs{b}{b\,b}}$ applies the \emph{inner} variable to itself.
The paper's claim: \enquote{binder-merging does not handle the name shadowing problem. This is
overlooked by all pattern unification algorithms. Algorithms that use what we call binder-merging state
upfront that they don't consider name shadowing}. Nipkow--Prehofer and Qian are cited as examples. The
existing remedies are \enquote{a pre-processor that adds gensyms to lambda variables to make each
individual unique or a smart reifier that detects name shadowing} (p.\,4). Their downsides are
unreadable gensyms and extra complexity. The related work gives a real case (p.\,8): Agda uses de Bruijn
numbers in its core language, so it \enquote{must make all variable names unique in a “gensym” manner},
which \enquote{bloats type checking with randomly generated symbols and makes the error messages painful to
read}.

\subsection{Static closures}\label{sec:lkanren-static}

\enquote{We propose a simple and efficient alternative, static closures. Static closures are
integrated into the unification process and they decide α-equivalence while correctly handling name
shadowing. We obtain a new algorithm by extending traditional pattern unification with static
closures} (p.\,5). No efficiency evidence or complexity bound is given.

\begin{definition}[Terms, scopes, static closures {\cite[Fig.~1, p.\,5]{MaFriedman2021}}]\label{def:lkanren-closure}
\[
\begin{array}{rcll}
t & ::= & X & \text{unification variable}\\
  & \mid & t\;t & \text{function application}\\
  & \mid & \abs{a}{t} & \text{named binder}\\
  & \mid & a & \text{name}\\
  & \mid & \lambda\,t & \text{nameless binder (derived)}\\
  & \mid & \ell & \text{level (derived)}\\[2pt]
\phi & ::= & \epsilon \mid a,\phi & \text{scope}\\
\kappa & ::= & \lkclo{\phi}{t} & \text{static closure}
\end{array}
\]
\enquote{The nameless syntax, i.e., nameless-binders and levels are derived during unification}
(p.\,5): the user writes only $X$, $t\,t$, $\abs{a}{t}$ and $a$. \enquote{A static closure is a pair
of a scope and a unification term. A scope is a list of binders. In essence, a static closure
captures the surrounding binders of a term. E.g., the term λa.λb.X reduces to the static closure
$\lkclo{a,b}{X}$.}
\end{definition}

Two remarks from our notes. First, the grammar has no constants, although the examples use them
(\texttt{'cat}); read \enquote{name} as covering free constants too. Second, the grammar conses on
the left ($a,\phi$), but the example writes the scope of $\abs{a}{\abs{b}{X}}$ as $a,b$, outermost
first. We adopt \emph{outermost first}: then the level of a binder is its position in the scope.

\begin{definition}[Operations on scopes {\cite[p.\,5]{MaFriedman2021}}]\label{def:lkanren-extdb}
In the paper's words:
\begin{itemize}
  \item \enquote{$\mathrm{ext}(\phi,a)$ adds the name $a$ to the scope $\phi$. If $a$ has been
    previously added to $\phi$, then this older one is shadowed and its corresponding level becomes
    unusable.}
  \item \enquote{$\mathrm{db}(\phi,a)$ is $\lkJust\;\ell$ where $\ell$ is the de Bruijn level of $a$
    w.r.t.\ the scope $\phi$, if $a$ is not shadowed in $\phi$; otherwise, the result is
    $\lkNothing$.}
  \item \enquote{$\mathrm{db}(\phi,\ell)$ is $\lkJust\;\ell$, if $\ell$ is available in $\phi$;
    otherwise, the result is $\lkNothing$.}
\end{itemize}
\end{definition}

So $\mathrm{db}$ is overloaded. On a name it converts the name to a level; on a level it checks that
the level is available. The notes make this precise. Write $\phi=a_0,\dots,a_{n-1}$ outermost first:
\[
  \mathrm{ext}(\phi,a)=\phi,a\qquad
  \mathrm{avail}(\phi)=\{\ell<|\phi| \;:\; a_\ell\notin\{a_{\ell+1},\dots,a_{n-1}\}\}
\]
\[
  \mathrm{db}(\phi,a)=\begin{cases}\lkJust\;\max\{\ell : a_\ell=a\} & a\in\phi\\ \lkNothing & a\notin\phi\end{cases}
  \qquad
  \mathrm{db}(\phi,\ell)=\begin{cases}\lkJust\;\ell & \ell\in\mathrm{avail}(\phi)\\ \lkNothing&\text{otherwise}\end{cases}
\]
In words: a level is \emph{available} if its binder's name is not rebound further in. The
name-to-level case needs an interpretation. An \emph{occurrence} of a name always refers to its
innermost binder, which is never shadowed, so the paper's \enquote{if $a$ is not shadowed} can only
mean \enquote{the innermost $a$}. The other reading would make $\abs{a}{\abs{a}{a}}$ impossible to
convert, which is absurd. This remains an interpretation.

\begin{tutorial}[the three examples of §4, step by step]
\textbf{Example 1} (p.\,5). Both sides are turned into closures:
\[
  \abs{a}{\abs{b}{X}} \lkq \abs{c}{\abs{d}{c}}
  \;\longrightarrow\; \lkclo{a,b}{X} \lkq \lkclo{c,d}{c}.
\]
\enquote{We first apply db on the right-hand side}: $\mathrm{db}((c,d),c)=\lkJust\;0$, so the problem
is $\lkclo{a,b}{X}\lkq\lkJust\;0$. \enquote{Then we try to solve X by applying db on the scope of the
left-hand side and 0}: $\mathrm{db}((a,b),0)=\lkJust\;0$, because $a$ is not rebound. Then
$\lkJust\;0\lkq\lkJust\;0$. \enquote{It works. So, the solution is $\{X\mapsto\texttt{d0}\}$.}

\medskip
\textbf{Example 2} (p.\,5--6), the shadowing case:
\[
  \abs{a}{\abs{a}{X}} \lkq \abs{c}{\abs{d}{c}}
  \;\longrightarrow\; \lkclo{a,a}{X} \lkq \lkclo{c,d}{c}.
\]
As before, the right side gives $\lkJust\;0$. \enquote{This time, db($\langle a, a; 0\rangle$) is
Nothing because the outer binder a is shadowed}. Here the paper puts the $0$ inside the closure
brackets, an abuse of notation for $\mathrm{db}((a,a),0)$. We get $\lkNothing\lkq\lkJust\;0$,
\enquote{and we know the unification fails}. True is not false any more.

\medskip
\textbf{Example 3} (p.\,6), two unknowns:
\[
  \abs{a}{\abs{b}{X}} \lkq \abs{c}{\abs{d}{Y}}
  \;\longrightarrow\; \lkclo{a,b}{X} \lkq \lkclo{c,d}{Y}.
\]
\enquote{Here, the grafting itself is not able to describe the solution. Now our solution needs a
context: the set of equations between static closures of unification variables. This problem has the
solution: the grafting being empty and the context being $\{\lkclo{a,b}{X}\lkq\lkclo{c,d}{Y}\}$.}
The paper adds that, in practice, \enquote{the accumulated equations between closures of one
unification problem are usually solved and reduced to grafting, during the unification on other
problems}.
\end{tutorial}

So a solution is a pair $(\theta,C)$: a grafting $\theta$ and a \emph{context} $C$ of postponed
flex-flex equations between closures. This is analogous to the flex-flex pairs that Huet's procedure
leaves unsolved (\cref{ch:hounif}).

\subsection{A reconstruction of the rules (not in the paper)}\label{sec:lkanren-recon}

To check that the three examples fit together, we wrote a minimal rule system that
reproduces all of them. It was implemented as a standalone CHICKEN program of about 150 lines, which
does not use \micro{}. We repeat that this is \textbf{our reconstruction, not the authors'
algorithm}.

\paragraph{Level form.} A surface term $t$ under scope $\phi$ is converted to internal syntax
$⟦t⟧_\phi$:
\[
  ⟦X⟧_\phi=\lkclo{\phi}{X}\qquad
  ⟦a⟧_\phi=\begin{cases}\ell & \mathrm{db}(\phi,a)=\lkJust\;\ell\\ a&\text{otherwise (a constant)}\end{cases}
\]
\[
  ⟦\abs{a}{t}⟧_\phi=\lambda\,⟦t⟧_{\mathrm{ext}(\phi,a)}\qquad
  ⟦t\,u⟧_\phi=⟦t⟧_\phi\,⟦u⟧_\phi
\]
Both sides of an equation are converted from $\epsilon$, so they are always compared at equal depth.
A binding is $\theta(X)=(\phi_0,t)$, where $t$ is in level form at depth $|\phi_0|$. Levels below
$|\phi_0|$ are \emph{free} (scope) levels; the others are bound inside $t$. To \emph{instantiate}
$\theta(X)$ at an occurrence $\lkclo{\phi}{X}$:
\begin{itemize}
  \item a free level $\ell$ must satisfy $\mathrm{db}(\phi,\ell)=\lkJust\;\ell$, or else the
    occurrence fails;
  \item a local level $\ell\ge|\phi_0|$ is shifted to $\ell+|\phi|-|\phi_0|$;
  \item a nested closure is re-rooted at $\phi$.
\end{itemize}

\paragraph{Rules.} The judgement is $\theta;C\vdash t\lkq u\Rightarrow\theta';C'$. Terms are walked
first, i.e.\ a bound closure is replaced by its instance. Any pair of head constructors not listed
fails.
\begin{mathpar}
\inferrule{ }{\theta;C\vdash \ell\lkq\ell \Rightarrow \theta;C}\;\rulename{Lvl}
\and
\inferrule{c\text{ a constant}}{\theta;C\vdash c\lkq c \Rightarrow \theta;C}\;\rulename{Const}
\and
\inferrule{\theta;C\vdash t\lkq u\Rightarrow\theta';C'}{\theta;C\vdash \lambda t\lkq\lambda u\Rightarrow\theta';C'}\;\rulename{Lam}
\and
\inferrule{\theta;C\vdash t_1\lkq u_1\Rightarrow\theta_1;C_1 \\ \theta_1;C_1\vdash t_2\lkq u_2\Rightarrow\theta_2;C_2}
          {\theta;C\vdash t_1\,t_2\lkq u_1\,u_2\Rightarrow\theta_2;C_2}\;\rulename{App}
\and
\inferrule{X\notin\mathrm{dom}\,\theta \\ FL_{<|\phi|}(t)\subseteq\mathrm{avail}(\phi) \\ X\notin FV(t) \\
           \theta[X\mapsto(\phi,t)];\emptyset\vdash C\Rightarrow\theta';C'}
          {\theta;C\vdash\lkclo{\phi}{X}\lkq t\Rightarrow\theta';C'}\;\rulename{Graft}
\and
\inferrule{ }{\theta;C\vdash\lkclo{\phi}{X}\lkq\lkclo{\phi}{X}\Rightarrow\theta;C}\;\rulename{Refl}
\and
\inferrule{X,Y\notin\mathrm{dom}\,\theta \\ (X,\phi)\neq(Y,\psi)}
          {\theta;C\vdash\lkclo{\phi}{X}\lkq\lkclo{\psi}{Y}\Rightarrow\theta;C\cup\{\lkclo{\phi}{X}\lkq\lkclo{\psi}{Y}\}}\;\rulename{Postpone}
\end{mathpar}
In \rulename{Graft}, $t$ is not an unbound closure. $FV(t)$ is the set of unification variables of
$t$ (the occurs check). $FL_{<|\phi|}(t)$ is the set of levels below $|\phi|$ that occur in $t$: those
that point into the scope $\phi$ rather than to a nameless binder inside $t$. The availability premise
says that a solution may mention only scope binders that are visible, i.e.\ not shadowed, at the
occurrence of $X$. The premise \enquote{$\vdash C$} means that every postponed equation is solved
again, the \emph{wake-up}. Example~2 fails in the availability premise of \rulename{Graft}:
$FL(0)=\{0\}\not\subseteq\mathrm{avail}(a,a)=\{1\}$.

\begin{tutorial}[a wake-up, step by step]
Take Example~3 and then bind $Y$, in two variants. Both start from
$\abs{a}{\abs{b}{X}}\lkq\abs{c}{\abs{d}{Y}}$ (first variant) or
$\abs{a}{\abs{a}{X}}\lkq\abs{c}{\abs{d}{Y}}$ (second variant).
\begin{enumerate}
  \item Conversion gives $\lambda\lambda\lkclo{a,b}{X}\lkq\lambda\lambda\lkclo{c,d}{Y}$. \rulename{Lam}
    twice, then \rulename{Postpone}: $\theta=\emptyset$, $C=\{\lkclo{a,b}{X}\lkq\lkclo{c,d}{Y}\}$. This is the
    line \texttt{(X ctx: 1)} below.
  \item The second equation $\abs{c}{\abs{d}{Y}}\lkq\abs{c}{\abs{d}{c}}$ converts to
    $\lambda\lambda\lkclo{c,d}{Y}\lkq\lambda\lambda\,0$. \rulename{Graft} checks
    $FL(0)=\{0\}\subseteq\mathrm{avail}(c,d)=\{0,1\}$ and binds $Y\mapsto((c,d),0)$.
  \item Wake-up: $C$ is solved again. Walking $\lkclo{c,d}{Y}$ instantiates $Y$'s solution at
    $(c,d)$, giving $0$. The equation becomes $\lkclo{a,b}{X}\lkq 0$, and \rulename{Graft} checks
    $0\in\mathrm{avail}(a,b)$. It holds, so $X\mapsto((a,b),0)$, printed \texttt{d0}.
  \item In the second variant the same wake-up yields $\lkclo{a,a}{X}\lkq 0$, but
    $\mathrm{avail}(a,a)=\{1\}$, so the whole conjunction fails.
\end{enumerate}
The second variant is \enquote{true = false} again, smuggled through the context. A wake-up that only
re-unified, without the availability check, would accept it.
\end{tutorial}

Running the prototype gives the following results. In \texttt{(λ d1)} the level is printed absolute,
at the depth where $X$ was solved.

\begin{answers}
X = cat                               => (cat ctx: 0)
λa.a = λb.b                           => (X ctx: 0)
λa.X = λb.b                           => (d0 ctx: 0)
λa.X = λb.b ∧ λc.X = λd.d             => (d0 ctx: 0)
λa.λa.X = λc.λd.c  (true=false)       => FAIL
λx.λy.(x y) = λa.λa.(a a)             => FAIL
Ex1 λa.λb.X = λc.λd.c                 => (d0 ctx: 0)
Ex2 λa.λa.X = λc.λd.c                 => FAIL
Ex3 λa.λb.X = λc.λd.Y                 => (X ctx: 1)
Ex3 then λc.λd.Y = λc.λd.c            => (d0 ctx: 0)
λa.λa.X = λc.λd.Y then Y = d0         => FAIL
λa.X = λb.λc.c ; λa.λb.X = λ_.λ_.λc.c => ((λ d1) ctx: 0)
λa.X = λb.λc.b                        => ((λ d0) ctx: 0)
λa.X = λb.b ∧ X = cat                 => FAIL
λa.X = λb.b ∧ λa.λa.X = λc.λd.d       => FAIL
λa.λb.X = λa.λb.a ∧ λb.λa.X = λc.λd.c => (d0 ctx: 0)
\end{answers}

\noindent Three lines are worth reading closely. In \enquote{λa.λa.X = λc.λd.Y then Y = d0}, the
wake-up must re-run the availability check, or the shadowing bug comes back through the context. In
\enquote{λa.X = λb.λc.c ; …}, a binder \emph{inside} a grafted solution must be shifted when the
solution is used at a different depth. The last line shows that the semantics is \emph{positional}:
the binder names above $X$ are irrelevant as long as no level is shadowed.

\subsection{The language: Horn clauses, \texttt{assume}, \texttt{all}}\label{sec:lkanren-mf21-lang}

§5 of the paper (p.\,6--7) recalls the logic from the 2020 paper \cite{MaLuFriedman2020}. Its Fig.~2
contrasts Horn clauses with hereditary Harrop formulas. In the figure the shared rows are
$D::=G\imp D\mid D\wedge D\mid\forall x\,D$ and $G::=G\wedge G\mid G\vee G\mid\exists x\,G$, and the
HH-only rows are $D\imp G$ and $\forall x\,G$ for goals. The atoms $A$ and $\top$ are missing from the
figure. With them it is exactly the book's grammar (3.1) \book{91}:
\[
  G ::= \top \mid A \mid G\wedge G \mid G\vee G \mid \exists x\,G \mid D\imp G \mid \forall x\,G\qquad
  D ::= A \mid G\imp D \mid D\wedge D \mid \forall x\,D .
\]
The four code boxes are these (all solid, p.\,6--7):

\begin{scheme}
> (run* (X Y)
    (conde
      [(== X Y) (== X 'cat)]
      [(== X 'cat) (== Y 'dog)]))
'((cat cat) (cat dog))

(define (empty? X)
  (== X '()))

>(run 2 X
   (assume (== 'apple 'orange)
     (== 'apple X)))
'(apple orange)

>(run 1 F
   (all (A)
     (== A (F A))))
'((lambda (_0) _0))
\end{scheme}

\noindent The paper prints \texttt{conde} with a superscript \textit{e}. The first box is
$\exists X\exists Y((X=Y\wedge X=\texttt{'cat})\vee(X=\texttt{'cat}\wedge Y=\texttt{'dog}))$. The
\texttt{define} is the clause $\forall X\,(X=\texttt{'()})\imp(\texttt{empty?}\;X)$. \texttt{assume}
\enquote{introduces hypotheses in its clause. As an example, we may use the assume operator to
temporarily extend the == relation}. Our reading, inferred because the paper gives no operational
account: the first answer comes from ordinary unification, and the second from the hypothesis used as
a clause whose head matches \scm{(== 'apple X)} with $X=\texttt{orange}$. For \texttt{all},
\enquote{we may synthesize a function that yields whatever it takes in}. Also, \enquote{the ==
operator recognizes β-equivalence. Indeed, it will be powered by the unification algorithm proposed
in the last section and recognize αβ-equivalence} (p.\,7). Note that \scm{(F A)} is a Miller pattern
(\cref{ch:hounif}): a variable applied to a distinct eigenvariable.

% =====================================================================================================
\section{A critical assessment of static closures}\label{sec:lkanren-assess}

The paper attacks pattern unification. The fair comparison also includes nominal unification
(\cref{sec:lkanren-nominal}), which the paper mentions only in related work. This section gives our own
analysis. In the λProlog column, $X$ (and $Y$) are quantified outside the λ, as in the paper's
queries.

\begin{table}[ht]
\centering\small
\caption{The paper's examples under three semantics (our analysis, after \cite{MaFriedman2021}).}
\label{tab:lkanren-three}
\begin{tabularx}{\textwidth}{>{\raggedright\arraybackslash}p{4.1cm}>{\raggedright\arraybackslash}X>{\raggedright\arraybackslash}X>{\raggedright\arraybackslash}X}
\toprule
problem & λProlog (patterns) & αKanren (nominal) & static closures\\
\midrule
$\abs{a}{a}\lkq\abs{b}{b}$ & \lkyes{} & \lkyes{} & \lkyes{}\\
$\abs{a}{X}\lkq\abs{b}{b}$ & \lkno{} (scope) & \lkyes{} $X=a$ & \lkyes{} $X=\texttt{d0}$\\
\dots\ $\wedge\ \abs{c}{X}\lkq\abs{d}{d}$ & \lkno{} & \lkno{} ($a\neq c$) & \lkyes{} \texttt{d0}\\
$\abs{a}{\abs{a}{X}}\lkq\abs{c}{\abs{d}{c}}$ & \lkno{} & \lkno{} (needs $a\#a$) & \lkno{}\\
$\abs{x}{\abs{y}{x\,y}}\lkq\abs{a}{\abs{a}{a\,a}}$ & \lkno{} & \lkno{} & \lkno{}\\
$\abs{a}{\abs{b}{X}}\lkq\abs{c}{\abs{d}{Y}}$ & \lkyes{} $X=Y$ (neither may depend on $a,b,c,d$) &
  $X=(a\,c)(b\,d)\cdot Y$ with $a\#Y$, $b\#Y$ & postponed\\
\bottomrule
\end{tabularx}
\end{table}

\begin{tutorial}[what λProlog offers the novice instead: make the dependency explicit]
The paper calls the failure of $\abs{a}{X}\lkq\abs{b}{b}$ counter-intuitive. Miller's 1991 paper
\cite[p.\,4--5]{Miller1991} explains why it is the intended behaviour, and how to get the novice's
answer anyway. Its program is the usual \texttt{append}, and its goal is
$\forall y\,(\mathtt{append}\;(\mathtt{cons}\;a\;(\mathtt{cons}\;b\;\mathtt{nil}))\;y\;X)$.
\rulename{generic} replaces $y$ by a new constant $k$, and backchaining twice would need
$X=\mathtt{cons}\;a\;(\mathtt{cons}\;b\;k)$. This is rejected, and Miller says why: \enquote{Such a failure
here is quite sensible since the value of $X$ should be independent of the choice of the constant used
to instantiate $\forall y$.} If we do want a value that depends on $y$, we must say so, by writing
$(H\;y)$ instead of $X$. Then $(H\;k)$ must equal $\mathtt{cons}\;a\;(\mathtt{cons}\;b\;k)$. Two λ-terms
would do, $\abs{w}{\mathtt{cons}\;a\;(\mathtt{cons}\;b\;k)}$ and
$\abs{w}{\mathtt{cons}\;a\;(\mathtt{cons}\;b\;w)}$, but \enquote{since $H$ cannot contain $k$ free, only
the second} is legal. Miller calls this \enquote{discharging the constant $k$}.

The same move answers the novice. Write $\abs{a}{(X\;a)}\lkq\abs{b}{b}$: this is a pattern ($X$ applied
to one bound variable), and its most general unifier is $X\mapsto\abs{x}{x}$. The conjunction with
$\abs{c}{(X\;c)}\lkq\abs{d}{d}$ is solved by the same $X$, and a top-level $Q=X$ gets the closed term
$\abs{x}{x}$, with nothing dangling. Both λKanren papers show this style of query:
\scm{(== (tie (a b) (app a b)) (tie (x y) (app f x y)))} of 2020 and
\scm{(all (A) (== A (F A)))} of 2021 (\cref{sec:lkanren-mf21-lang,sec:lkanren-mk20}). Static closures
make this dependency \emph{implicit}: $\lkclo{\phi}{X}$ behaves like $X$ applied to the available
levels of $\phi$ (observation~1 below).
\end{tutorial}

\subsection{Three counter-examples}\label{sec:lkanren-counter}

The notes isolate three examples that pin down what static closures really are. They are worth
keeping as regression tests for any implementation, whatever it chooses.

\begin{tutorial}[Counter-example 1: the conjunction separates static closures from nominal unification]
Problem: $\abs{a}{X}\lkq\abs{b}{b}\ \wedge\ \abs{c}{X}\lkq\abs{d}{d}$.

\emph{Nominal}, using the rule $a.t\approx^? b.s\Rightarrow t\approx^?(a\,b)\cdot s,\ a\#^? s$ of
\cref{sec:lkanren-nominal}. The first equation gives $X\approx^?(a\,b)\cdot b=a$ with $a\#b$, which
holds, so $X:=a$. The second becomes $c.a\approx^?d.d$, i.e.\ $a\approx^?(c\,d)\cdot d=c$, which
fails. Nominal unification thus agrees with static closures on the single equation ($X$ is the bound
name) and disagrees on the conjunction.

\emph{Pattern}: both equations fail at once, by the scope restriction.

\emph{Static closures}: $X\mapsto\texttt{d0}$ satisfies both.

So static closures are \emph{neither} pattern unification \emph{nor} nominal unification. $X$ denotes
an open term \emph{relative to the binder depth}. This is grafting on de Bruijn terms \emph{without}
the \enquote{pre-cooking} step of Dowek et al.\ (\cref{sec:lkanren-dhkp}).
\end{tutorial}

\begin{tutorial}[Counter-example 2: the same variable at two depths dangles]
Problem: \scm{(== Q X)} together with \scm{(== (λ (a) X) (λ (b) b))}.

Under static closures, $X\mapsto\texttt{d0}$, so the top-level answer is $Q=\texttt{d0}$. That is a
level which is \emph{not available} in the empty scope where $Q$ lives. The answer is meaningful only
relative to the occurrence of $X$ under one binder. The reading becomes incoherent as soon as the same
variable occurs at different depths. Our \micro{} run above shows the same thing with indices. The
prototype chooses to check availability inside unification, so \enquote{λa.X = λb.b then top-level
Y = X} fails there. But it skips the check when printing, which is how the paper can print
\texttt{'(d0)}. Dowek et al.\ fix this with $X[\uparrow^n]$; Teyjus and ELPI fix it with levels on
variables. The 2021 paper does not address it. It also gives no rules for application and β, nor for
flex-flex closures beyond postponing them.
\end{tutorial}

\begin{tutorial}[Counter-example 3: \enquote{true = false} is a bug of binder-merging, not of pattern unification]
Problem: $\abs{a}{\abs{a}{X}}\lkq\abs{c}{\abs{d}{c}}$.

\emph{Pattern}: $X$ cannot depend on the bound variables at all, so this fails.

\emph{Nominal}: $a.(a.X)\approx^? c.(d.c)$ gives $a.X\approx^?(a\,c)\cdot(d.c)=d.a$, with
$a\#(d.c)$, which holds. Then $a.X\approx^?d.a$ gives $X\approx^?(a\,d)\cdot a=d$, with
\textbf{$a\#a$}, which is false. So it fails.

\emph{Static closures}: fails (Example~2).

All three agree. The \enquote{false = true} danger is specific to \emph{binder-merging}
implementations that do not rename apart. The earlier λKanren paper renames apart: its ξ step uses
gensyms (\cref{sec:lkanren-mk20}), so it does not have the bug.
\end{tutorial}

\subsection{Further observations}

The notes add several observations that a porter must know. The paper does not say them.

\begin{enumerate}
  \item \textbf{The meaning of a unification variable changes.} In λProlog, $X$ in $\abs{a}{X}$
    cannot mention $a$. With static closures, $X$ implicitly depends on the \emph{positions} of its
    enclosing binders. Logically, $\lkclo{\phi}{X}$ behaves like Miller's raising
    $X\,\ell_{i_1}\cdots\ell_{i_k}$ over the available levels of $\phi$, with one twist: the
    \enquote{function} is indexed by absolute position, so occurrences under scopes of different
    lengths share it. The paper's related work quotes Huet's \enquote{elementary substitution}:
    \enquote{Each unification variable should capture all its visible bound variables} (p.\,8).
    Static closures make that idea implicit, minus the shadowed binders.
  \item \textbf{α-equivalence is no longer a congruence on open terms.} $\abs{a}{\abs{a}{X}}$ and
    $\abs{a}{\abs{b}{X}}$ are α-equivalent in the usual sense, because $X$ is a metavariable with no
    free names. Yet against the same $\abs{c}{\abs{d}{c}}$, Example~2 fails and Example~1 succeeds.
    With grafting, the names of the binders above a metavariable matter. Nominal logic manages exactly
    this with suspended permutations $\pi\cdot X$ and freshness $a\#X$. The paper's \enquote{names
    don't matter} holds for the levels, not for the closures.
  \item \textbf{Levels are stable for scope variables, not for binders inside solutions.} A grafted
    $\lambda.\ell_1$ solved at depth~1 must become $\lambda.\ell_2$ at depth~2. A locally nameless
    representation avoids the shift \cite{McBrideMcKinna2004,Chargueraud2012}: levels for scope
    variables, indices for binders inside solutions.
  \item \textbf{Flex-flex need not be postponed.} With levels, $\lkclo{\phi}{X}\lkq\lkclo{\psi}{Y}$
    at equal depth has a most general solution: $X\mapsto Y$, plus the restriction that $Y$ may use
    only levels in $\mathrm{avail}(\phi)\cap\mathrm{avail}(\psi)$. This is Miller-style pruning,
    phrased as a freshness-like constraint.
  \item \textbf{The trick must not reach \texttt{all}.} For ∀-goals the eigenvariable must respect the
    scope restriction, or \scm{(run 1 X (all (A) (== X A)))} would succeed. That is unsound for ∀.
    The level semantics is meant only for λ-binders inside terms.
  \item \textbf{β brings undecidability back.} Once \scm{==} \enquote{recognizes αβ-equivalence},
    unification outside the pattern fragment is undecidable (\cref{ch:hounif}). One must restrict to
    patterns, delay, or enumerate.
\end{enumerate}

The paper's related work (p.\,8) recalls that \enquote{nominal unification and pattern unification are
equally expressive}. Cheney reduces pattern unification to nominal unification \cite{Cheney2005}, and
Levy and Villaret give the converse reduction \cite{LevyVillaret2008}. Pattern unification is linear
(Qian, although the paper notes that Qian's algorithm \enquote{has not been implemented}); for nominal
unification \enquote{$O(n^2)$ is the state-of-the-art} (p.\,8). The notes summarize the correspondence like this:
the static closure $\lkclo{\phi}{X}$ is close to the nominal suspension $\pi\cdot X$ (both record the
context around an unknown); $\mathrm{db}$ is close to the action of $\pi$ on names; and \enquote{level
unavailable in $\phi$} is close to the freshness $\ell\#X$. But static closures are
\emph{positional} where nominal terms are \emph{name-based}, as counter-example~1 shows. The design is
closer to de Bruijn-level \enquote{contextual} metavariables than to nominal logic.

\begin{kanren}[what static closures would demand of a miniKanren]
The notes list what a static-closure layer on \micro{} would need. This is a requirements list, not a
design (see \cref{ch:design}):
\begin{itemize}
  \item \textbf{New term kinds, dispatched early.} λ, level and closure records must be handled
    \emph{before} the generic record branch. Today \scm{μkanren-structure?} matches any record
    instance and unifies it slot by slot, names included. The same applies to every walker: the occurs
    check, \scm{μkanren-state-find/value}, reification, \scm{μkanren-hasvar?}, and so on.
  \item \textbf{Conversion at the boundary.} Converting both sides of \scm{=°} to level form under
    $\epsilon$ leaves a unifier that is \enquote{practically first-order} except for closures.
  \item \textbf{A walk that can fail.} Walking $\lkclo{\phi}{X}$ instantiates $X$ at $\phi$, which can
    fail on an unavailable level. So the walk needs a third outcome besides \enquote{value} and
    \enquote{unbound}. The first-order hot path should stay untouched.
  \item \textbf{A context store with wake-up.} The postponed equations must be re-checked when their
    variables get bound. The \scm{μkanren-verify-D/touched} idiom, which re-runs only the entries that
    mention a just-bound variable, already does this for disequalities.
  \item \textbf{Disequality.} With postponement, a tentative unification inside \scm{≠°} may succeed
    by only adding a context equation, producing no bindings. An empty prefix in D means
    \enquote{already equal}, so this would be wrong. A disequality entry would then have to carry
    its postponed part as well.
\end{itemize}
\end{kanren}

% =====================================================================================================
\section{λKanren 2020: Miller's algorithm in a shallow embedding}\label{sec:lkanren-mk20}

The earlier paper, \emph{Higher-order Logic Programming with λKanren} \cite{MaLuFriedman2020}, is the
\enquote{[17]} of the 2021 paper. Its workshop title is \enquote{λKanren: Higher-order Logic
Programming with Shallow Embedding}. It adds four operators to μKanren \cite{HemannFriedman2013}:
\scm{tie} (a binder), \scm{app} (application), \scm{assume-rel} (an implication goal $D\imp G$) and
\scm{all} (a universal goal $\forall x\,G$). In it, \scm{==} \enquote{recognizes αβ-conversions}. The
paper claims \enquote{about 40 lines} added to μKanren (p.\,7); the repository's \texttt{unify.rkt}
has 648 lines.

\paragraph{Terms and state.} The terms are
$\mathit{Term}::=\mathit{Constant}\mid\mathit{Var}\mid\mathsf{App}(\mathit{rator},\mathit{rand})\mid
\mathsf{Tie}(\mathit{var},\mathit{body})\mid(\mathit{Term}\,.\,\mathit{Term})$. Pairs stay a separate
first-order constructor. The state is the list \texttt{(qs es σ name->idx)}. Here \texttt{qs} is an
\textbf{explicit quantifier prefix} of \texttt{(∃ x)} and \texttt{(∀ x)} entries, appended at the end
by \texttt{fresh} and \texttt{all}. \texttt{es} holds postponed equations and σ is the substitution.

\paragraph{The unifier} is a direct transcription of Miller 1991 §6 \cite{Miller1991}: ξ, rigid-rigid,
raising, pruning, flex-flex and flex-rigid (\cref{sec:lkanren-miller}). Its ξ step strips the
\texttt{Tie}s on both sides and renames both binder lists to the same \emph{fresh} gensyms by
capture-avoiding substitution. It then pushes a \texttt{(∀ w)} onto the prefix for each name. So
shadowing is handled by renaming apart, and the 2021 problem does not arise. Non-pattern applications
are \textbf{not delayed}; they make the unifier fail.

\begin{scheme}
(define-syntax tie
  (syntax-rules ()
    [(_ () t) t]
    [(_ (x0 x ...) body)
     (let ([x0 (Var 'x0)])
       (Tie x0 (tie (x ...) body)))]))

(define-syntax app
  (syntax-rules ()
    [(_ rator rand) (App rator rand)]
    [(_ rator rand0 rand ...)
     (app (App rator rand0) rand ...)]))
\end{scheme}

\noindent (The paper prints \texttt{[(\_ () t) body]}; the repository has \texttt{t}.) The paper reports
these unification queries (p.\,2--3):

\begin{scheme}
(run 1 q (== (tie (a b) (app a b)) (tie (x y) (app x y))))      ; ⇒ '(_0)
(run 2 f (== (tie (a b) (app a b)) (tie (x y) (app f x y))))    ; ⇒ '((tie (_0) (tie (_1) (app _0 _1))))
(run 1 f (== (tie (a) a) (tie (b) (app f b b))))                ; ⇒ '()   not a pattern
(run 1 q (== (tie (a b) (app a b)) (tie (x y) (app x q))))      ; ⇒ '()   scope restriction
\end{scheme}

\noindent The second query is a pattern, so \scm{run 2} still finds one answer, the most general
unifier: \enquote{There is only one instantiation} (p.\,3). Binders and unification variables are both
reified as \texttt{\_}$n$ symbols. The third is not a pattern ($f$ is applied to $b$ twice) and fails,
even though it has solutions ($\abs{x}{\abs{y}{x}}$ and $\abs{x}{\abs{y}{y}}$); in the paper's words,
\enquote{we cannot decide which b takes control}. The fourth fails by the scope restriction, which the
paper phrases as visibility: \enquote{A binding variable x is visible to a unification variable q if the
introduction of x lexically precedes that of q} (p.\,3). Here $q$ would have to be the bound $y$, which
is introduced after $q$.

\paragraph{Implication and the shallow-embedding problem.} The motivating program is a curriculum:

\begin{scheme}
(defrel (taken name class)
  (conde
    [(== 'Josh name) (== 'B521 class)]
    [(== 'Josh name) (== 'B522 class)]))
(defrel (pl-major name)
  (taken name 'B521)
  (taken name 'B523)
  (taken name 'B522))
> (run 1 q
    (assume-rel [(taken name class) (== 'Josh name) (== q class)]
      (pl-major 'Josh)))
'(B523)
\end{scheme}

\noindent The hypothetical clause extends \texttt{taken}, and \texttt{pl-major}, which was defined
\emph{before}, must see it. In a shallow embedding, relation definitions live in the host language's
lexical scope, so \texttt{pl-major} cannot see the new clause. The paper describes a
\textbf{dynamic-scope layer}: a map \texttt{name->idx} threaded through the query, a global
\texttt{idx->def}, and calls that go through both. (Our notes write \scm{(== q class)}; the extracted
text has a quote, \texttt{'q}, which could not produce \texttt{B523}.) Concretely, before the query
\texttt{name->idx} is \texttt{((pl-major . 2) (taken . 1))}; \texttt{assume-rel} pushes a new index, so
inside it the map is \texttt{((taken . 3) (pl-major . 2) (taken . 1))}, and the newer pair shadows the
older one when \texttt{pl-major} calls \texttt{taken} (p.\,5).

The paper also offers \texttt{assume}, \enquote{a shorter version of the assume-rel operator} that
\enquote{only takes an atomic hypothesis} (p.\,5), and uses it on a relation \texttt{eq} that is
reflexive, transitive and symmetric. Plain \scm{(run 1 q (eq 'apple 'orange))} does not terminate in
μKanren, because the recursive clauses never run out (p.\,5). With hypotheses (p.\,5--6):

\begin{scheme}
(defrel (eq x y)
  (conde
    [(== x y)]
    [(fresh (z) (eq x z) (eq z y))]
    [(eq y x)]))
> (run 5 q (assume (eq 'orange 'apple) (assume (eq 'orange 'dog) (eq 'orange q))))
'(dog apple orange orange dog)
> (run 1 q (assume (eq 'orange q) (eq 'apple 'orange)))
'(apple)
\end{scheme}

\noindent In the first query the answers come from the two hypotheses (\texttt{dog}, \texttt{apple}) and
from reflexivity (\texttt{orange}); \texttt{run 5} reports repetitions, because the recursive clauses
derive the same facts again along different paths. The second query \enquote{runs backward}: \enquote{the
hypothesis can be inferred as well} (p.\,6). The hypothesis \scm{(eq 'orange q)} mentions the query
variable $q$, and it is \emph{not} renamed when it is used. The goal \scm{(eq 'apple 'orange)} can use the
symmetry clause to become \scm{(eq 'orange 'apple)}, which matches the hypothesis with $q=\texttt{apple}$.
This is the book's semantics for a clause with a free logic variable (\cref{ch:hh}): only a clause's own
∀-bound variables are freshened per use.

\paragraph{Universal goals.} \enquote{the all operator effectively creates a constant \dots\ we create
an all variable as a free binding variable. Thus, the all variable cannot be unified with anything
but itself} (p.\,6). The scope restriction is enforced by the prefix. The examples (p.\,6--7):

\begin{scheme}
> (run 1 q (all (x)
             (assume-rel (taken x 'B521)
               (assume-rel (taken x 'B522)
                 (assume-rel (taken x 'B523)
                   (pl-major x))))))
'(_0)
> (run 1 f (all (x) (== x (app f x))))
'((tie (_0) _0))
(defrel (proved x) (== x 'trivial))
> (run 1 g (all (p) (proved p)))            '()
> (run 1 g (proved g))                      '(trivial)
> (run 1 g (all (p q) (assume ((proved p) (proved q)) (proved q) (proved p))))   '(_0)
> (run 1 goal (all (p q) (assume ((proved p)) (conde [(proved p)] [(proved q)]))))  '(_0)
\end{scheme}

\noindent The query \scm{(all (p) (proved p))} shows \rulename{generic} at work. \texttt{proved} holds only of
\texttt{trivial}, and a generic $p$ is not \texttt{trivial}, so $\forall p.\,\texttt{proved}\;p$ is not
provable, whereas the same goal without \scm{all} succeeds with $g=\texttt{trivial}$. The last two queries
use λKanren \enquote{as a theorem prover} (§5, p.\,6--7). The first proves the commutativity of
conjunction, $\forall p,q\,(p\wedge q)\imp(q\wedge p)$: \rulename{generic} makes $p$ and $q$ new
constants, \rulename{augment} adds \scm{(proved p)} and \scm{(proved q)} as hypotheses, and the two
conjuncts of the body are then proved by backchaining on them. The second proves disjunction introduction,
$\forall p,q\;p\imp(p\vee q)$, by the first \scm{conde} branch. The answers \texttt{'(\_0)} only say that
the query variable was left unconstrained: the information is in the success.

\begin{pitfall}[hypotheses in a host dynamic parameter get lost]
The repository's active code does not use the described state-threaded design. It uses Racket
parameters: \texttt{defrel} makes each body a \texttt{make-parameter}, and \texttt{assume} is a
\texttt{parameterize} that extends it. We reproduced this structure in CHICKEN. The body
of \texttt{pl-major} is only \emph{constructed} when the goal is applied to a state, which happens
after \texttt{parameterize} has returned. So the inner call to \texttt{taken} reads the
\emph{original} parameter, and the output was \texttt{(orig Josh B523)}. Relation calls nested inside
other relations' delayed bodies therefore lose the hypotheses. The described \texttt{name->idx} design
does not have this problem.
\end{pitfall}

\begin{kanren}[hypotheses belong to the state]
\micro{}'s \scm{define-relation} is inverse-η-delayed: it expands to
\scm{(λ (s) (δ ((and° g ...) s)))}. This is exactly the situation of the experiment above. A
hypothetical clause database for \rulename{augment} must therefore be a \emph{field of the state},
pushed on entry to $G$ and invisible outside it. It can never be a host parameter. A further remark
from the notes: a hypothetical clause $\forall\bar x(G\imp A)$ can be a host closure whose body runs
\scm{fresh°}. Each use then renames the clause's own variables, while the variables captured from
the enclosing scope stay shared. This is exactly λProlog's semantics for clauses with free logic
variables. How the module does this is shown in \cref{sec:module-hh}.
\end{kanren}

% =====================================================================================================
\section{Nominal unification and αKanren}\label{sec:lkanren-nominal}

Nominal unification \cite{UrbanPittsGabbay2004} takes a different route to α-equivalence. Bound names
stay names, called \emph{atoms}. Metavariables are holes, and substitution of a term for a
metavariable is \emph{capturing}, i.e.\ grafting. Capture is controlled by two devices:
\emph{permutations} suspended on variables, and \emph{freshness} constraints. Our presentation of the
rules follows Urban's survey \cite{Urban2010}.

\begin{definition}[Nominal terms and permutations]\label{def:lkanren-nominal}
\[
  t ::= \langle\rangle \mid \langle t_1,t_2\rangle \mid f\,t \mid a \mid a.t \mid \pi\cdot X
\]
where $a.t$ abstracts the atom $a$ in $t$, and $\pi\cdot X$ is a variable with a suspended
permutation. A permutation $\pi$ is a list of swappings $(a\,b)$:
\[
  []\cdot a=a,\qquad ((a_1\,a_2)::\pi)\cdot a=\begin{cases}a_2 & \pi\cdot a=a_1\\ a_1 & \pi\cdot a=a_2\\ \pi\cdot a&\text{otherwise.}\end{cases}
\]
Composition is list append and the inverse is reversal. The \emph{disagreement set} is
$ds(\pi,\pi')=\{a\mid \pi\cdot a\neq\pi'\cdot a\}$. A permutation acts on an abstraction on
\emph{both} the binder and the body, and it suspends on variables:
$\pi\cdot(\pi'\cdot X)=(\pi @ \pi')\cdot X$. Substitution is capturing: $\sigma(a.t)=a.\sigma(t)$ with
no side condition, and $\sigma(\pi\cdot X)=\pi\cdot\sigma(X)$.
\end{definition}

\begin{definition}[Freshness and α-equality {\cite[Fig.~1]{Urban2010}}]\label{def:lkanren-fresh}
With a freshness context $\nabla$, a set of pairs $a\#X$:
\begin{mathpar}
\inferrule{ }{\nabla\vdash a\#\langle\rangle}
\and
\inferrule{\nabla\vdash a\#t_1 \\ \nabla\vdash a\#t_2}{\nabla\vdash a\#\langle t_1,t_2\rangle}
\and
\inferrule{\nabla\vdash a\#t}{\nabla\vdash a\#f\,t}
\and
\inferrule{ }{\nabla\vdash a\#a.t}
\and
\inferrule{\nabla\vdash a\#t \\ a\neq b}{\nabla\vdash a\#b.t}
\and
\inferrule{a\neq b}{\nabla\vdash a\#b}
\and
\inferrule{(\pi^{-1}\cdot a,X)\in\nabla}{\nabla\vdash a\#\pi\cdot X}
\\
\inferrule{\nabla\vdash t_1\approx t_2}{\nabla\vdash a.t_1\approx a.t_2}
\and
\inferrule{a\neq b \\ \nabla\vdash a\# t_2 \\ \nabla\vdash t_1\approx(a\,b)\cdot t_2}{\nabla\vdash a.t_1\approx b.t_2}
\and
\inferrule{\forall c\in ds(\pi,\pi').\ (c,X)\in\nabla}{\nabla\vdash\pi\cdot X\approx\pi'\cdot X}
\end{mathpar}
plus the congruence rules for unit, pairs, function symbols and atoms.
\end{definition}

The intuition: $a\#t$ says that $a$ does not occur \emph{free} in $t$, also after any later
instantiation of the variables of $t$. To compare $a.t_1$ with $b.t_2$, rename $b$ to $a$ in $t_2$ by
a \emph{swapping}, which is always safe because it is a bijection, and require that $a$ was not
already free in $t_2$.

\begin{definition}[Nominal unification {\cite[p.\,9]{Urban2010}}]
\[
  \{\langle t_1,t_2\rangle\approx^?\langle s_1,s_2\rangle,\dots\}\Rightarrow\{t_1\approx^?s_1,\ t_2\approx^?s_2,\dots\}\qquad
  \{a.t\approx^?a.s,\dots\}\Rightarrow\{t\approx^?s,\dots\}
\]
\[
  \{a.t\approx^?b.s,\dots\}\Rightarrow\{t\approx^?(a\,b)\cdot s,\ a\#^?s,\dots\}\qquad
  \{\pi\cdot X\approx^?\pi'\cdot X,\dots\}\Rightarrow\{a\#^?X\mid a\in ds(\pi,\pi')\}\cup\{\dots\}
\]
and $\pi\cdot X\approx^?t$, with $X$ not in $t$, gives $X:=\pi^{-1}\cdot t$. A solution is a pair
$(\nabla,\sigma)$ with $\nabla\vdash\sigma(t_1)\approx\sigma(t_2)$. For example, $a.X\approx^?b.X$ has
the solution $(\{a\#X,\,b\#X\},[\,])$.
\end{definition}

Most general solutions exist, and \textbf{no fresh names are ever generated}: every atom of a solution
occurs in the problem. Naive rules are exponential; quadratic algorithms exist.

\subsection{αKanren}

αKanren \cite{ByrdFriedman2007,Byrd2009} embeds nominal unification in miniKanren. \scm{fresh}
introduces \emph{noms}: \enquote{a nom behaves more like a constant than a variable, since it only
unifies with itself or with an unassociated variable}. The constructor \scm{(tie a t)}, typeset
$\bowtie$, binds the nom $a$ in $t$; it delimits scope and does not create noms. \scm{(# a t)},
entered \scm{hash}, asserts that $a$ does not occur free in $t$. Ordinary logic variables are
introduced by \scm{exist}. The thesis examples (ch.~9; \texttt{◃▹} is the thesis glyph for tie):

\begin{scheme}
(run* (q) (fresh (a b) (≡ (◃▹ a a) (◃▹ b b))))      ⇒ (_.0)
(run* (q) (fresh (a b) (≡ (◃▹ a q) (◃▹ b b))))      ⇒ (a.0)
(run* (q) (fresh (a b) (# a (◃▹ b a))))             ⇒ ()
(run* (q) (fresh (a) (# a (◃▹ a a))))               ⇒ (_.0)
(run* (q) (fresh (a b) (exist (x y) (≡ (◃▹ a (◃▹ a x)) (◃▹ a (◃▹ b y))) (≡ `(,x ,y) q))))
  ⇒ ((((susp ((a.0 a.1)) _.0) _.0) : ((a.0 . _.0))))
\end{scheme}

\noindent The second line is counter-example~1's single equation: $q$ becomes the bound nom. The last
line shows a suspension in an answer, with a freshness constraint printed after the colon. The
capture-avoiding substitution relation shows the programming style:

\begin{scheme}
(define substo
  (λ (e new a out)
    (matche (e out)
      (((var ,a) ,new))
      (((var ,y) (var ,y)) (# a y))
      (((app ,rator ,rand) (app ,rator-res ,rand-res))
       (substo rator new a rator-res)
       (substo rand new a rand-res))
      (((lam ,(◃▹ @c body)) (lam ,(◃▹ c body-res)))
       (# c a)
       (# c new)
       (substo body new a body-res)))))
(run* (q) (fresh (a b) (substo `(lam ,(◃▹ a `(app (var ,a) (var ,b)))) `(var ,b) a q)))
  ⇒ ((lam (tie a.0 (app (var a.0) (var a.1)))))     ; [b/a]λa.ab = λc.cb
\end{scheme}

\noindent (\texttt{@c} in a \texttt{matche} pattern introduces $c$ as a nom.) The unifier works in two
phases. First the σ-rules turn equations into a substitution plus pending freshness constraints. Then
the ∇-rules reduce freshness constraints until each is $a\#x$ with $x$ an unbound variable. Variables
are suspensions with an empty permutation. A triangular-substitution version, due to Joe Near, is
reported \enquote{as much as ten times faster}. The thesis states a limitation (p.\,178): \enquote{neither
currently supports disequality constraints}. αKanren also powers αleanTAP \cite{NearByrdFriedman2008},
where ∀-bound variables of the object logic are noms. There, \enquote{Universally quantified formulas
are used as templates, rather than instantiated directly}.

\begin{kanren}[what nominal unification asks of a miniKanren]
It needs three new term kinds (nom, \scm{tie}, suspension $\pi\cdot X$) and a freshness store ∇ of
$a\#X$ pairs, re-reduced after each unification. That re-reduction works like \micro{}'s
post-\scm{=°} verification of D, A and T. Note that \scm{absent°} is \emph{not} freshness: freshness is
about \emph{free} occurrences, so \scm{(# a (◃▹ a a))} succeeds, whereas absento forbids any
occurrence. Nominal unification has no function variables and no β, so it is a different logic from
λProlog's. Noms are distinct constants, not eigenvariables, and ∀-goals still need their own
mechanism (see \cref{ch:design}).
\end{kanren}

% =====================================================================================================
\section{Eigenvariables, negation and ∀ inside miniKanren}\label{sec:lkanren-eigen}

\subsection{The scope-extrusion problem}

\rulename{generic} proves $\forall x.G$ by proving $G[c/x]$ for a \emph{new} constant $c$. If the
constant is only a gensym, nothing stops an older logic variable from being bound to it later. That
\enquote{proves} $\exists X\forall y.\,X=y$, which must fail. Contrast $\forall y\exists X.\,X=y$,
which must succeed. The fixes in the literature are exactly four:
(i)~an explicit prefix \cite{Miller1991,MaLuFriedman2020}; (ii)~universe indices or levels \cite{NadathurJayaramanKwon1995,DunchevEtAl2015}; (iii)~an absento-like eigen constraint over the variables in scope \cite{HemannByrdFriedman2013}; (iv)~an end-of-scope check.
The notes call (iv) the weakest. It is sound only if nothing created inside the scope survives it in a
form that can be bound later, and its failures are detected late.

\subsection{\texttt{eigen} in miniKanren with symbolic constraints}

Hemann, Byrd and Friedman added \scm{eigen} to miniKanren-with-symbolic-constraints in 2013
\cite{HemannByrdFriedman2013}. The header comment of \texttt{mk.scm} reads: \enquote{E = (e* . x*)*,
where e* is a list of eigens and x* is a list of variables. Each e in e* is checked for any of its
eigens be in any of its x*. Then it fails.} The state carries B, the list of all variables created so
far. \scm{eigen} creates fresh eigen objects and installs an eigen-absento constraint $(e^*\,.\,B)$:
\emph{none of the variables that existed before the eigen may ever walk to a term containing it}.

\begin{scheme}
(define-syntax eigen
  (syntax-rules ()
    ((_ (x ...) g0 g ...)
     (lambdag@ (c : B E S)
       (let ((x (eigen-var)) ...)
         ((fresh () (eigen-absento `(,x ...) B) g0 g ...) c))))))

(define eigen-absento
  (lambda (e* x*)
    (lambdag@ (c : B E S D Y N T)
      (cond
        ((eigen-occurs-check e* x* S) (mzero))
        (else (unit `(,B ((,e* . ,x*) . ,E) ,S ,D ,Y ,N ,T)))))))
\end{scheme}

\noindent This is the universe-index check implemented with explicit sets instead of numbers. Its
cost is $O(|B|\cdot|\text{term}|)$ per check. The README warns: \enquote{Beware: this implementation
does \emph{not} support use of \texttt{eigen} with constraints other than \texttt{==}.}

\begin{kanren}[the eigen-absento idea, by hand, on today's engine]
\micro{} already has \scm{absent°}, which takes any datum as its tag, including a record. So the
Hemann--Byrd--Friedman check can be written by hand. Here there is one outer variable $Q$, and the
eigen is a fresh record that $Q$ may never contain. This illustrates the idea; it is not an API
proposal.
\begin{scheme}
(import scheme (chicken base) (chicken pretty-print)
        (aux base) (aux stream) (aux kanren micro))

(define-syntax-rule (show label expr)
  (begin (display ";; ") (display label) (newline)
         (pretty-print expr)))

;; A generic constant as a fresh, unique record: no scope check at all.
(define-record eigen id)
(show "naive: exists Q forall x. Q = x"
  (°->list #f (fresh° (q) (=° q (make-eigen 1)))))

;; The eigen-absento idea, by hand with absent°: the constant
;; may never occur in a variable that existed before it.
(show "with absent°: exists Q forall x. Q = x"
  (°->list #f
    (fresh° (q)
      (let ((c (make-eigen 1)))
        (and° (absent° c q) (=° q c))))))
(show "with absent°: forall x exists Y. Y = x, then Q = ok"
  (°->list #f
    (fresh° (q)
      (let ((c (make-eigen 1)))
        (and° (absent° c q) (fresh° (y) (=° y c)) (=° q 'ok))))))
(show "with absent°: escape through an alias is caught"
  (°->list #f
    (fresh° (q)
      (let ((c (make-eigen 1)))
        (and° (absent° c q)
              (fresh° (y) (=° q (list y)) (=° y c)))))))
\end{scheme}
\begin{answers}
;; naive: exists Q forall x. Q = x
((λ () (make-record-instance 'eigen 1)))
;; with absent°: exists Q forall x. Q = x
()
;; with absent°: forall x exists Y. Y = x, then Q = ok
((λ () 'ok))
;; with absent°: escape through an alias is caught
()
\end{answers}
Three cautions come from the notes on \micro{}. First, \scm{μkanren-state-find/value} rebuilds records,
so eigens must be compared by an identifier (with \scm{equal?}), never by \scm{eq?}; and two eigens
must never share an identifier. Second, \micro{} binds the \emph{older} variable to the \emph{newer}
one in a variable-variable equation. So \enquote{rank = creation index} is not a scope discipline by
itself, and creation time is not scope depth. Third, the check is first-order only: there is no
higher-order unification here.
\end{kanren}

\subsection{Moiseenko 2019: constructive negation}

Moiseenko \cite{Moiseenko2019} builds the answers to $\neg g$ by collecting \emph{all} answers of $g$
and negating them. To make this closed under negation, states become formulas (eq.~(1), p.\,6):
\[
  \exists\bar x.\Big(\bigwedge_i v_i=t_i\ \wedge\ \bigwedge_j\forall\bar y_j.\bigvee_k w_{jk}\neq u_{jk}\Big).
\]
This needs a solver for \emph{universally quantified disequalities}:
\begin{lemma}[{\cite[Lemma 3, p.\,9]{Moiseenko2019}}]
$t\neq u$ is satisfiable iff $\mathrm{mgu}(t,u)\neq\top$.
\end{lemma}
\noindent This is exactly what \micro{}'s \scm{≠°} implements.
\begin{lemma}[{\cite[Def.~11, Lemma 4, p.\,10]{Moiseenko2019}}]
Let $\theta=\mathrm{mgu}(t,u)$, oriented so that $\bar x\not\subseteq\mathrm{codom}\,\theta$. Then
$\forall\bar x.\,t\neq u$ is satisfiable iff $\theta=\bot$ or $\mathrm{dom}(\theta)\not\subseteq\bar x$,
and unsatisfiable iff $\mathrm{dom}(\theta)\subseteq\bar x$.
\end{lemma}
\noindent The orientation is a normalization: if the mgu maps an outer variable $v$ to some $x\in\bar x$,
flip that binding to $x\mapsto v$. In words: $\forall\bar x.\,t\neq u$ is unsatisfiable exactly when $t$
and $u$ can be made equal by choosing values for the universal variables alone. Example: for
$\forall x.\,x\neq\alpha$, the mgu is $\{x\mapsto\alpha\}$, its domain is $\{x\}\subseteq\bar x$, so the
constraint is unsatisfiable. For $\forall x.\,\alpha\neq(1,x)$ the mgu $\{\alpha\mapsto(1,x)\}$ binds the
outer $\alpha$, so it is satisfiable (by any $\alpha$ that is not a pair with first component~$1$).
\begin{lemma}[{\cite[Def.~12, Lemma 5, p.\,10--11]{Moiseenko2019}}]
For $\forall\bar x\exists\bar y.\,t\equiv u$, let $\theta$ be idempotent and $\hat\theta$ its restriction
to the non-$\bar y$ variables. The formula is satisfiable iff $\mathrm{dom}(\hat\theta)\cap\bar x=\emptyset$
and no $p\in\mathrm{codom}(\hat\theta)$ has $\mathrm{Vars}(p)\cap\bar x\neq\emptyset$.
\end{lemma}
\noindent The notes observe that Lemma 5 is exactly the eigenvariable condition. Universal variables
may not be bound, and they may not appear in the bindings of outer variables. With negation one gets
$\mathit{ifte}\;c\;t\;e=(c\wedge t)\vee(\neg c\wedge e)$, $g_1\Rightarrow g_2=\neg g_1\vee(g_1\wedge g_2)$
and $\mathit{forall}\;g=\neg\,\mathit{fresh}\,(x)\,(\neg\,g\,x)$. The price: \emph{stratification} is
required, negating a goal with infinitely many answers diverges, and infinitely many constructors of
each arity are assumed. A remark on p.\,4: \enquote{there is no sound implementation of eigen variables
with the support of disequality constraints. We observed that our implementation of universal
quantifications through double negation works nicely with disequalities}.

\begin{pitfall}[eigen-constants and disequality]
λProlog's generic ∀ is sound because λProlog has no negation. Add \scm{≠°} and consider
$\forall x.\,x\neq\alpha$, with $\alpha$ an outer variable. It must \emph{fail}: take $x:=\alpha$.
Moiseenko's Lemma~4 agrees, since the mgu $\{x\mapsto\alpha\}$ has domain $\subseteq\{x\}$. A naive
eigen-constant $c$ makes $c\neq\alpha$ hold forever, because $\alpha$ can never be bound to $c$, and
the query wrongly succeeds. The same issue arises for type constraints: $c$ stands for an arbitrary
term, so \scm{symbol°} of $c$ should not be provable unless assumed. It also arises for
\scm{absent°}.
\end{pitfall}

\subsection{Jin, Rosenblatt, Might and Zhang 2021: ∀ and ⇒ by relative complement}

Jin et al.\ \cite{JinEtAl2021} build on the \emph{first-order miniKanren representation}
\cite{RosenblattEtAl2019}, where goals are data that can be inspected. \scm{(forall (v) domain goal)}
works in four steps: it checks the domain's satisfiability; it extracts candidates by running
\scm{(exists (v) (conj domain body))} and splitting the states into disjunction-free ones; it
computes a \emph{relative complement}, recursing on \scm{(forall (v) (conj relcomp domain) body)};
and it applies a domain relative negation. This is an \emph{extensional} ∀ over the universe of
literal values: the classical \enquote{for every value} reading, not λProlog's generic one.
\scm{(implies A B)} combines \emph{semantic} solving, which rewrites to $\neg A\vee B$ when $A$ is
negatable, with \emph{syntactic} solving. In the latter, an assumptions list is threaded through, a
user-relation call can be solved by an assumed call, and relations are unfolded under a
\emph{closed-world} reading. The paper's opening examples (p.\,1):

\begin{scheme}
> (run 1 (a) (forall (v) () (== v a)))                                   '()
> (run 1 () (forall (a) () (exists (v) () (== v a))))                    '(())
> (run 1 (a b) (forall (v) () (disj* (== v a) (=/= v b))))              '((_.0 _.0))
> (run 1 (b) (forall (a) () (disj* (not-symbolo a) (=/= a b))))         '((#f))
> (run 1 (f) (implies (evalo f 1) (exists (a) (evalo f a))))            '((_.0))
\end{scheme}

\noindent The third query is the instructive one. $\forall v.\,(v=a\vee v\neq b)$ holds when $a=b$.
That is a classical case split on $v$, which λProlog's generic ∀ cannot do; λProlog has no $\neq$ at
all. The authors list the limitations themselves (§7): the system is slow, produces duplicates, is
incomplete, and soundness is unproven. Their related work notes that λKanren 2020 supports only
relation calls and equalities in antecedents, with \enquote{no unfolding \dots\ consistent with an
open-world interpretation}.

\begin{pitfall}[three different ∀]
λProlog's ∀ is \rulename{generic}: intuitionistic, proved by a new constant, with no case analysis
(\cref{ch:hh}). Jin et al.'s ∀ is extensional: \enquote{for every value in the universe}, by relative
complements. Moiseenko's ∀ is $\neg\exists\neg$: classical, stratified, and it needs finitely many
answers. A query that succeeds in one may fail or diverge in another. For instance, Jin et al.'s §6.6
shows that \scm{(forall (expr) () (disj (neg (halto expr)) (halto expr)))} does not terminate. Only the
first ∀ is the book's.
\end{pitfall}

% =====================================================================================================
\section{How λProlog systems do it}\label{sec:lkanren-impl}

\subsection{Miller 1991: \texorpdfstring{$\Llam$}{L-lambda} by transitions over a mixed prefix}\label{sec:lkanren-miller}

\Cref{ch:hounif} presents pattern unification as the book does. Here we recall only the shape of
Miller's original formulation \cite{Miller1991}, because both λKanren 2020 and Nipkow transcribe it.
The book's statement of the restriction is: \enquote{a program clause or goal formula $B$ is in
$\Llam$ if all quantification in $B$ is over nonpredicate types and every subterm of $B$ of the form
$(x\;t_1\dots t_n)$ in which $x$ is an essentially existentially quantified variable is such that
$t_1,\dots,t_n$ are distinct variables that are universally quantified within the scope of the
quantifier binding $x$} \book{220}. Then only β$_0$ is needed: $(\abs{y}{s})\,x=s[x/y]$ when $x$ is a
variable not free in $\abs{y}{s}$.

Miller's interpreter turns \rulename{generic} into a \emph{meta-level} universal quantifier, and
\rulename{backchain} turns the clause's ∀s into ∃s placed at the point of use. This is why mixed
prefixes $\forall\exists\forall\dots$ arise dynamically. The unification transitions are ξ (descend
under binders, renaming apart), rigid-rigid, \emph{raising}, \emph{pruning}, flex-flex and flex-rigid.
The algorithm terminates (Theorem 7.1), and when it succeeds it yields a most general unifier
(Proposition 7.3). One of Miller's worked examples (p.\,14--15):
\[
  Q\exists u\exists v\forall y\,[f(\abs{x}{f(u\,x\,y)})=f(\abs{w}{v\,y})]
  \xrightarrow{\text{rigid}}\abs{x}{f(u\,x\,y)}=\abs{w}{v\,y}
  \xrightarrow{\xi}\forall y\forall x\,[f(u\,x\,y)=v\,y]
\]
\[
  \xrightarrow{\text{pruning}}[u\mapsto\abs{x}{\abs{y}{u'\,y}}]
  \xrightarrow{\text{flex-rigid}}[v\mapsto\abs{y}{f(u'\,y)}].
\]
Read the last two steps with the prefix in mind, which is now $Q\exists u\exists v\forall y\forall x$. The
equation is flex-rigid: $v\,y$ against $f(u\,x\,y)$. The universal $x$ is quantified after $v$ and is not
among $v$'s arguments, so $v$'s solution may not mention it. Every occurrence of $x$ sits under the
flexible head $u$, so the problem is not hopeless: \emph{pruning} removes $u$'s dependency on its first
argument, $u\mapsto\abs{x}{\abs{y}{u'\,y}}$. After that the rigid side is $f(u'\,y)$, and flex-rigid
abstracts $v$'s argument: $v\mapsto\abs{y}{f(u'\,y)}$. Had $x$ occurred \emph{outside} any flexible
head, pruning would fail and so would the equation.

In the untyped version (§9.3), rigid-rigid with different arities fails, and η-expansion is used only
on terms that must have functional type. This matters for an untyped Scheme term language. Miller also
shows how to stay inside $\Llam$ by replacing meta-level β with \emph{copy clauses} and hypothetical
reasoning, as in the book's Figure~7.10 \book{200}:

\begin{lprolog}
subst M T S :- pi x\ copy x T => copy (M x) S.
\end{lprolog}

\noindent The same trick encodes a \emph{non-pattern} problem inside $\Llam$ \cite[§10.3, p.\,30--31]{Miller1991}.
With copy clauses \lp{copyterm a a}, \lp{copyterm (f X) (f U) :- copyterm X U.} and so on for the
signature, define
\begin{lprolog}
substterm M T N :- pi c\(copyterm c T => copyterm (M c) N).
\end{lprolog}
\noindent The unifiers of $F\,a\lkq g\,a\,a$, which is not a pattern because $a$ is a constant, are exactly
the $F$ that make \lp{substterm F a (g a a)} provable. Miller lists the four answers an $\Llam$
interpreter should return: \lp{w\(g w w)}, \lp{w\(g w a)}, \lp{w\(g a w)}, \lp{w\(g a a)}. The choice
\enquote{is this $a$ the argument or the constant?} is made by backchaining, either on the hypothesis
\lp{copyterm c a} or on the clause \lp{copyterm a a}. Huet's branching has become ordinary proof
search, which a miniKanren would interleave.

\noindent The book's position on programs outside $\Llam$ is pragmatic. A common resolution is
\enquote{to not impose the $\Llam$ restriction on program clauses statically but rather to expect the
condition to be satisfied by the unification problems that are encountered dynamically}. One
\enquote{might defer further consideration of this equation until such time that substitutions
resulting from solving other parts of the problem have altered its status} \book{225--226}.

\subsection{Nipkow 1993: functional pattern unification on de Bruijn terms}

Nipkow \cite{Nipkow1993} restricts attention to $\forall\exists\forall$ problems, where raising is never
needed. Existential variables are global (\texttt{V F}), outer universals are constants (\texttt{C c}),
and inner universals are loose de Bruijn indices (\texttt{B i}). The substitution is triangular. From
the author's ML code, the algorithm is:
\begin{enumerate}
  \item head-normalize both sides through the substitution, with β$_0$ done by renaming indices;
  \item $\lambda s=\lambda t\Rightarrow s=t$; η on demand, $\lambda s=t\Rightarrow s=(\uparrow t)\,0$;
  \item rigid-rigid: equal heads give pairwise unification, different heads fail;
  \item flex-rigid $F\,\bar y=t$: occurs check, then \emph{invert} $t$ with respect to $\bar y$ and
    prune the flexible subterms that mention indices outside $\bar y$;
  \item flex-flex, same head: $F\,\bar y=F\,\bar z\Rightarrow F:=\lambda^{|\bar y|}.H\,(\text{positions where }y_i=z_i)$;
  \item flex-flex, different heads: if $\bar\imath\subseteq\bar\jmath$ then $G:=\lambda\bar\jmath.F\,\bar\imath$,
    symmetrically in the other case, and otherwise both are bound to a fresh $H$ applied to
    $\bar\imath\cap\bar\jmath$.
\end{enumerate}

\subsection{Dowek, Hardin, Kirchner, Pfenning 1996: explicit substitutions}\label{sec:lkanren-dhkp}

DHKP \cite{DowekHardinKirchnerPfenning1996} work in the λσ-calculus, where substitution $a[s]$ is
syntax. Meta-variables are \emph{grafted}, i.e.\ replaced first-order. To recover the logical,
capture-avoiding reading, terms are \textbf{pre-cooked}:
$F(\lambda a,n)=\lambda F(a,n+1)$, $F(a\,b,n)=F(a,n)\,F(b,n)$, $F(k,n)=1[\uparrow^{k-1}]$,
$F(X,n)=X[\uparrow^n]$. A metavariable under $n$ binders is explicitly \emph{shifted past them}. This
is precisely the step that makes grafting agree with substitution, and that static closures omit
(counter-example~1). In a λσ-pattern, every $X[s]\,b_1\dots b_q$ has $s=a_1\dots a_p\cdot\uparrow^n$
with the $a$'s and $b$'s distinct indices. Pattern substitutions are injective and have a right
inverse, so $X[s]=b$ is solved by $X:=b[s^{-1}]$ when every index of $b$ occurs in $s$ (Prop.~8,
inversion). Unification of λσ-patterns is decidable and unitary (Theorem~17), and \enquote{our
presentation does not need mixed prefixes}: raising is just a rule. Outside the fragment, §5.4 says
that equations \enquote{should be postponed as constraints in the hope that they will be instantiated
further}, which \enquote{avoids all branching during constraint simplification}. Empirically,
\enquote{almost all dynamically arising equations fall within the pattern fragment} (p.\,2).

\subsection{Universe indices: Nadathur, Jayaraman and Kwon 1995}

The idea behind every efficient implementation of \rulename{generic} is numeric
\cite{NadathurJayaramanKwon1995}: \enquote{Each time a universal quantifier is encountered, a new
constant must be introduced, giving rise to the next universe in the hierarchy. This requirement can be
accounted for by increasing the “universe index” by 1 and introducing a new constant tagged with this
index. \dots\ When an existential quantifier is encountered, it is instantiated by a logic variable.
This variable may be tagged with the current value of the universe index, the intended interpretation
of the tag being that a term may be substituted for the variable only if all the constants appearing in
the term have a smaller or identical tag} (p.\,15--16).

\begin{definition}[Binding with universe indices {\cite[p.\,16]{NadathurJayaramanKwon1995}}]\label{def:lkanren-ui}
To bind a variable $X$ with tag $i$ to $t$:
(1)~perform the usual occurs check; (2)~check that $t$ contains no constant with tag $>i$; (3)~\emph{before binding}, change to $i$ the tags of the variables of $t$ whose tag is $>i$.
\end{definition}

Step (3) is \emph{lowering}. It is \enquote{required in order to prevent a later instantiation of
these variables from violating the restrictions on instantiations for X}.

\begin{tutorial}[why lowering is needed]
Program: $\forall z(q(z)\imp p(d(z)))$. Goal: $\exists x\forall y(q(y)\imp p(x))$.
\begin{enumerate}
  \item $\exists x$ creates $X$ with tag~1. $\forall y$ raises the index to 2 and creates $c$ with
    tag~2. \rulename{augment} adds $q(c)$. The goal is now $p(X^1)$.
  \item \rulename{backchain} on the clause gives a new variable $Z$ tagged with the current index, 2.
    Unifying $p(X^1)=p(d(Z^2))$ binds $X:=d(Z)$, and \textbf{lowers} $Z$ to tag~1.
  \item The subgoal $q(Z^1)$ meets the hypothesis $q(c^2)$. Binding $Z^1:=c^2$ violates step~(2), so
    it fails.
\end{enumerate}
Without lowering, $Z^2:=c$ would be accepted and $X$ would become $d(c)$. That smuggles the
eigen-constant into a variable quantified outside its scope. The failure is correct: the goal would
need $x=d(y)$, but $x$ is quantified before $y$.
\end{tutorial}

The abstract interpreter works on tuples $\langle G,P,I\rangle$. A ∃-goal tags its new variable with
$I$; a ∀-goal goes to $\langle[c/x]G,P,I{+}1\rangle$; and backchaining instantiates the clause's
variables with tag $I$. The interpreter is sound and complete with respect to intuitionistic
provability, and goal selection order is irrelevant (Prop.~4.1). In the WAM, a \emph{UI register} is
incremented on entering a universal goal and decremented on its completion, and choice points save
it.

\subsection{Teyjus: suspensions and on-the-fly raising}

Teyjus represents term-bound variables by de Bruijn indices and logic variables by mutable cells
\cite{Nadathur2005}. β-reduction is made lazy with the \emph{suspension notation}
$[[t,ol,nl,e]]$: \enquote{a term $t$ that used to occur inside $ol$ abstractions but now appears within
$nl$ of them}, with an environment $e$. Some of the rules:
\begin{gather*}
  (\beta_s)\ (\lambda t_1)\,t_2\to[[t_1,1,0,(t_2,0)::nil]]\qquad
  (r_9)\ [[t,0,0,nil]]\to t\\
  (r_7)\ [[\lambda t,ol,nl,e]]\to\lambda[[t,ol{+}1,nl{+}1,@nl::e]]
\end{gather*}
Terms also carry a closed/open annotation bit; empirically, substitution walks were reduced to
\enquote{a third to an eighth}.

Nadathur and Linnell \cite{NadathurLinnell2005} eliminate the explicit prefix. Each variable gets a
tag $l(y)$, the number of ∃-followed-by-∀ alternations up to its quantifier, and \enquote{a universal
variable $u$ is quantified within the scope of the quantifier for an existential variable $x$ just in
case $l(x)<l(u)$}. Raising is restricted to the universals that actually appear as arguments. A single
traversal of the rigid side (\texttt{mksubst}/\texttt{bnd}) computes the binding, i.e.\ the inversion,
and does pruning and raising on the fly. A non-pattern subterm met during the traversal can be replaced
by a fresh $h(\bar y)$ with the equation $h(\bar y)=t$ \textbf{deferred}. The book describes the same
idea of \enquote{attaching numerical indices to the quantified variables} \book{225}. It also notes that
Teyjus~2 implements only $\Llam$ unification \book{228}.

\begin{example}[no mgu outside the fragment {\cite[p.\,5]{NadathurLinnell2005}}]\leavevmode\\
The problem $\forall u_1\forall u_2\exists x\,(x(u_2)=u_1(u_2))$ is not a pattern: $u_2$ is quantified
\emph{outside} the scope of $x$, so it is a constant for $x$. It has \enquote{two incomparable solutions},
$x\mapsto\abs{w}{u_1(w)}$ and $x\mapsto\abs{w}{u_1(u_2)}$ (written $\lambda(1,u_1(1))$ and
$\lambda(1,u_1(u_2))$ in the paper's de Bruijn notation), and no most general one. Move the quantifier,
$\forall u_1\exists x\forall u_2\,(x(u_2)=u_1(u_2))$, and it becomes a pattern with the single solution
$x\mapsto\abs{w}{u_1(w)}$. A unifier that must return one mgu can only defer the first problem. An
enumerating one would return both answers, as Miller's \texttt{substterm} does for
$F\,a\lkq g\,a\,a$ (\cref{sec:lkanren-miller}).
\end{example}

\subsection{ELPI: de Bruijn levels and a reduction-free fragment}

ELPI \cite{DunchevEtAl2015} starts from an observation about the two roles of β:

\begin{lprolog}
of (lam F) (arr A B) :- pi x\ of x A => of (F x) B.
cbn (app (lam F) N) M :- cbn (F N) M.
\end{lprolog}

\noindent In \lp{of}, β only replaces a bound name with a fresh constant. In \lp{cbn}, β is genuine
object-level substitution. The design has three parts.

\paragraph{Levels for the prefix.} \enquote{When a clause's head is unified with the goal in a context
of length $n$, the universally quantified variables of the clause are instantiated to unification
variables $X^n$ where the level $n$ records that $X$ has only visibility of the initial prefix of length
$n$ of the context. If later a fresh constant is added by the pi rule, the constant occupies position
$n+1$ \dots\ and it will not be allowed to occur in instances of the variable $X^n$} (p.\,4). Level
lowering is done by pruning: the higher-level variable is bound to a fresh lower-level one.

\paragraph{De Bruijn levels, not indices.} The example is $\abs{x}{(\abs{y}{\abs{z}{f\,x\,y\,z}})\,x}$:
\begin{align*}
  \text{indices:}\quad& \lambda x.(\lambda y.\lambda z.f\,x_2\,y_1\,z_0)\,x_0\to_\beta\lambda x.\lambda z.f\,x_1\,x_1\,z_0
    && (x \text{ is renumbered}),\\
  \text{levels:}\quad& \lambda x.(\lambda y.\lambda z.f\,x_0\,y_1\,z_2)\,x_0\to_\beta\lambda x.\lambda z.f\,x_0\,x_0\,z_1
    && (x \text{ keeps its number}).
\end{align*} Three properties follow. A variable bound in Γ keeps its name in any extension
of Γ (DBL1) and in a β-reduct (DBL2), and pushing a binder into the context needs no lifting (DBL3).
\enquote{variables already pushed in the context are treated exactly as constants, and \dots\ the two
notions of level --- De Bruijn's and the position in the context --- coincide}. Terms are
$t::=x_i\mid X^j\mid\lambda x_i.t\mid t\,t$, with global constants as negative $i$; $X^j$ sees the names
below $j$.

\begin{definition}[Reduction-free fragment $L^\beta_\lambda$ {\cite[Def.~1]{DunchevEtAl2015}}]
Every occurrence of a unification variable $X^j$ is applied to $x_j\dots x_{j+k-1}$, for $k\ge0$: exactly
the next $k$ context names, in order.
\end{definition}

\noindent This fragment is contained in $\Llam$ (Property 1). Head reductions of
$(X^j\,x_j\dots x_{j+k-1})\sigma$ take constant time (Property 2), and so does solving
$X^j\,x_j\dots x_{j+k-1}\equiv t$ when no occurs check is needed and the maximum level of $t$ is cached
(Property 3). But $L^\beta_\lambda$ is \textbf{not stable under resolution} (p.\,7), so the fast path
must fall back to general pattern unification. ELPI reports large speed-ups: typeof in 0.27\,s versus
5.64\,s for Teyjus, reduce-cbv in 0.15\,s versus 11.11\,s.

\paragraph{Outside the fragment,} ELPI's default is an \textbf{error}: \enquote{Unification problem
outside the pattern fragment. \dots\ Pass -delay-problems-outside-pattern-fragment \dots\ in order to
delay (deprecated, for Teyjus compatibility)}. Delay is offered instead through \emph{syntactic
constraints}. \lp{declare_constraint (even X) [X]} suspends a goal on variables, and the goal is resumed
when any of them is assigned, so \lp{declare_constraint (even X) [X], X = 1.} fails. Suspended goals keep
their hypothetical context. The documentation's example (\texttt{ELPI.md}, \enquote{Syntactic
constraints}):
\begin{answers}
goal> pi x\ sigma Y\ even x => declare_constraint (even Y) [Y].
Success:
Constraints:
 {x} : even x ?- even (W x)  /* suspended on W */
\end{answers}
\noindent The residual constraint remembers the eigenvariable \texttt{x} and the hypothesis
\lp{even x}. Note also the automatic raising: \texttt{Y} was created inside the scope of \texttt{x}, so
it is printed as \texttt{W x}. Constraint Handling Rules can operate on suspended goals.

\subsection{Makam}

Makam \cite{StampoulisChlipala2018} is a λProlog dialect with polymorphism, used to prototype type
systems. Two idioms are relevant for us. \texttt{bindmany} is a binder over a \emph{list} of
variables, with \texttt{openmany}/\texttt{applymany} to open many binders at once. And reflective
predicates such as \texttt{refl.isunif} test whether a term is an uninstantiated unification
variable. That test is non-logical, and is the analogue of \micro{}'s \scm{project°}.

\begin{kanren}[what the implementations ask of a miniKanren]
The real systems agree on three mechanisms:
(i)~a \emph{level} on every logic variable and eigen-constant, with the binding check and lowering
    of \cref{def:lkanren-ui}; (ii)~pattern unification with inversion and pruning (Nipkow, Nadathur--Linnell, ELPI); (iii)~a store of \emph{delayed} equations or goals, woken on assignment (DHKP §5.4, Teyjus, ELPI's
    syntactic constraints).
\micro{} has no notion of level. Its creation index is a time stamp, not a scope depth: a variable
created after a ∀-scope has closed has a larger index, but it must not see that scope's constant. On
the other hand, the \scm{assocs}-driven re-verification of D is exactly the wake-up hook that a delay
store needs. Search is unchanged: \scm{or°}'s fair interleaving can host Huet's branching. That fits
Miller's remark that a breadth-first search could compute the converse of his \texttt{prenex}
relation, which the depth-first interpretation cannot \cite{Miller1991}. The design itself is in \cref{ch:design}.
\end{kanren}

% =====================================================================================================
\section{Comparison of the design options}\label{sec:lkanren-compare}

\Cref{tab:lkanren-options} sets the four candidate meanings of \scm{=°} on λ-terms side by side. The
rows summarize the analysis of \cref{sec:lkanren-mf21,sec:lkanren-assess,sec:lkanren-nominal,sec:lkanren-impl}. The table records trade-offs and does not rank the options: the choice,
and the reasons for it, are the subject of \cref{ch:design}.

\begin{table}[ht]
\centering\footnotesize
\caption{Design options for higher-order unification in a miniKanren. No option is
chosen here; see \cref{ch:design}.}
\label{tab:lkanren-options}
\renewcommand{\arraystretch}{1.25}
\begin{tabularx}{\textwidth}{>{\raggedright\arraybackslash}p{2.0cm}*{4}{>{\raggedright\arraybackslash}X}}
\toprule
& \textbf{A.\ Pattern unification, de Bruijn + levels} (Miller, Nipkow, Nadathur--Linnell, ELPI)
& \textbf{B.\ Nominal} (Urban--Pitts--Gabbay, αKanren)
& \textbf{C.\ Static closures} (Ma--Friedman 2021)
& \textbf{D.\ Huet pre-unification, enumerated by the stream}\\
\midrule
Fidelity to the book & exact: $\Llam$ is what Teyjus~2 and ELPI run & a different logic (nominal); no
function variables, no β & non-standard grafting semantics: $\abs{a}{X}\lkq\abs{b}{b}$ succeeds, unlike
λProlog & exact for full higher-order unification (pre-unifiers)\\
Decidable / unitary & yes / yes (mgu) & yes / yes, with a freshness context & claimed; no full
algorithm published (a \enquote{preview}) & semi-decidable; possibly infinitely many pre-unifiers\\
New term forms & λ node (indices), eigen-constant with level, flexible application $X\,\bar y$ & nom,
\texttt{tie}, suspension $\pi\cdot X$ & named and nameless binders, levels, closures & as A, plus
imitation and projection\\
New state & a level per variable, current depth, delayed-equation store & freshness store ∇ of $a\#X$ &
a \enquote{context} of closure equations & as A; branching via \texttt{or°}\\
Cost of α & free (syntactic on indices) & swappings on the fly, suspensions & levels & free\\
β & only β$_0$; general β via copy/subst relations, or delayed & none (write \texttt{substo}) &
\enquote{will be powered} (not given) & full\\
∀-goals (\rulename{generic}) & eigen-constant + level check + lowering & a separate mechanism is still
needed: noms are distinct constants, not eigenvariables & \texttt{all} from λKanren 2020 (prefix) & as
A\\
With \texttt{≠°}, \texttt{absent°}, types & eigens need a ∀-aware check (Moiseenko) & αKanren has no
disequality at all & not discussed & as A\\
Effort in \aux{} & medium--high: a Nipkow-size core, plus levels and a delay store & medium: Byrd's
ch.~11 is about 300 lines; a triangular version exists & unclear & high, plus control of infinite
branching\\
\bottomrule
\end{tabularx}
\end{table}

Some combinations are possible and appear in the literature. Option D can sit on top of A as a
separate, explicitly nondeterministic relation; the book itself presents Huet's procedure as a
search over a matching tree (\cref{ch:hounif}), and miniKanren's interleaving could explore it
fairly. Option B can live beside A, for users who want names as data. Option C behaves differently
from all the others on the three counter-examples of \cref{sec:lkanren-counter}; any implementation
should state its answers to them explicitly. The smaller alternatives for ∀ alone are the explicit
prefix of λKanren 2020 (simplest, but each check walks the prefix and the prefix only grows), the
eigen-absento of Hemann--Byrd--Friedman (smallest change, $O(|B|)$ per check, first-order only), and
the end-of-scope check (weakest). The constructive ∀/⇒/¬ of Moiseenko and Jin et al.\ is orthogonal:
it has a classical, extensional reading and does not implement the book's \rulename{generic} and
\rulename{augment}.

% =====================================================================================================
\clearpage
\section{Summary}\label{sec:lkanren-summary}

\begin{itemize}
  \item Ma and Friedman 2021 \cite{MaFriedman2021} is a workshop \emph{preview}. It has two grammars,
    the operations $\mathrm{ext}$ and $\mathrm{db}$, three examples, and no algorithm, rules or
    theorems.
  \item Its first complaint is that pattern unification rejects $\abs{a}{X}\lkq\abs{b}{b}$. That is
    the scope restriction, which λProlog needs for the soundness of ∀. Its fix is to graft de Bruijn
    levels, giving $X\mapsto\texttt{d0}$. λProlog's own answer is to make the dependency explicit:
    $\abs{a}{(X\,a)}\lkq\abs{b}{b}$ is a pattern with mgu $X\mapsto\abs{x}{x}$ (Miller's
    \enquote{discharging a constant}).
  \item Its second complaint is that binder-merging plus naive reification makes
    $\abs{a}{\abs{a}{X}}\lkq\abs{c}{\abs{d}{c}}$ succeed, so \enquote{true = false}. Static closures
    $\lkclo{\phi}{t}$ make shadowed levels unavailable, and the equation fails. Flex-flex closure
    equations are postponed in a \emph{context}.
  \item The assessment uses three counter-examples. The conjunction separates static closures from
    nominal unification. The dangling \texttt{d0} at top level shows that $X$ has no
    depth-independent meaning. \enquote{True = false} also fails under pattern and nominal
    unification, so it is a bug of binder-merging only. Static closures are grafting on de Bruijn
    terms without DHKP's pre-cooking, and α is no longer a congruence on open terms.
  \item λKanren 2020 transcribes Miller 1991 with an explicit prefix and gensym renaming. Non-patterns
    fail. Hypotheses must be threaded in the state, because a host parameter loses them under delayed
    relation bodies.
  \item Nominal unification (αKanren) handles α with swappings and freshness, using capturing
    substitution. It has no function variables, no β and no disequality, and noms are not
    eigenvariables.
  \item Eigenvariables in miniKanren: a fresh constant without a scope check is unsound. Hemann, Byrd
    and Friedman implement the check as an absento over the earlier variables. Disequality and
    eigen-constants conflict (Moiseenko). Jin et al.'s ∀ is extensional, not generic.
  \item Real λProlog systems use universe indices or levels with lowering (Nadathur et al., Teyjus,
    ELPI), pattern unification with inversion and pruning (Nipkow, Nadathur--Linnell), and delayed
    non-pattern equations (DHKP, Teyjus, ELPI's syntactic constraints). Teyjus adds the suspension
    calculus; ELPI uses de Bruijn levels and the $L^\beta_\lambda$ fast path.
  \item The options are pattern+levels, nominal, static closures, and Huet enumerated by the stream
    (\cref{tab:lkanren-options}). The choice is made in \cref{ch:design}.
\end{itemize}


\part{A proposal for \texorpdfstring{\texttt{(aux kanren micro)}}{(aux kanren micro)}}
\chapter{Anatomy of \texorpdfstring{\texttt{(aux kanren micro)}}{(aux kanren micro)}}\label{ch:micro}

% ------------------------------------------------------------------------------------------------
\section*{What and why}

The chapters so far studied what a higher-order logic programming language \emph{is}: the uniform
proofs of hereditary Harrop formulas (\cref{ch:hh}), λ-terms as data (\cref{ch:lambda,ch:lambdatree}),
higher-order pattern unification (\cref{ch:hounif}), and, in \cref{ch:lkanren}, how λKanren grafts
these features onto a miniKanren~\cite{MaFriedman2021,MaLuFriedman2020}. This chapter opens the
\emph{proposal} part of the document. Before deciding how \micro{} should grow, we need an exact
picture of the engine as it is today: what a variable is, what a state contains, how unification
and the constraint stores interact, how the search interleaves, and how answers are printed. The
point is not to explain the code to its author but to put next to each other the facts that matter
when the engine is extended: which invariants a new feature must preserve, which parts can be
reused from outside, and which parts are closed.

The engine is a μKanren in the sense of Hemann and Friedman~\cite{HemannFriedman2013} (a goal maps a
state to a stream of states) with a cKanren-like constraint store holding disequalities, disjoint
type tags and \emph{absento} constraints, in the style of \emph{The Reasoned
Schemer}~\cite{FriedmanEtAl2018}. All references \texttt{:N} are line numbers of
\texttt{src/aux.kanren.micro.scm} (713 lines) at commit \texttt{f53bdcb}. Every Scheme snippet in this
chapter was run on that engine; its output is pasted verbatim below it. The snippets import
\scm{(aux base)}, \scm{(aux stream)}, \scm{(aux fds sbral)} and \micro.

\begin{table}[h]
\centering\small
\caption{A map of \texttt{aux.kanren.micro.scm}.}\label{tab:micro-map}
\begin{tabular}{@{}lll@{}}
\toprule
lines & content & section \\
\midrule
\texttt{3–26}    & module header, explicit export list, imports & \\
\texttt{28–59}   & sentinel, tags, variables & \cref{sec:micro-terms} \\
\texttt{60–107}  & the state record, sbral helpers & \cref{sec:micro-state,sec:micro-subst} \\
\texttt{110–204} & occurs check and binding, walks, unification & \cref{sec:micro-subst,sec:micro-unify} \\
\texttt{206–326} & reification and projection & \cref{sec:micro-reify} \\
\texttt{330–523} & D/A/T verification and subsumption & \cref{sec:micro-stores} \\
\texttt{525–629} & goals, combinators, constraint goals & \cref{sec:micro-search} \\
\texttt{631–684} & aggregation & \cref{sec:micro-aggr} \\
\texttt{686–713} & \texttt{define-relation}, runners & \cref{sec:micro-search,sec:micro-reify} \\
\bottomrule
\end{tabular}
\end{table}

% ------------------------------------------------------------------------------------------------
\section{Terms and variables}\label{sec:micro-terms}

There is no term datatype: a term is any Scheme value, and the engine gives structure to some of
them.

\begin{definition}[Terms of \micro]\label{def:micro-term}
A \emph{term} is one of: a logic variable; an atom (symbol, boolean, \scm{'()}, character, number,
string, or any other Scheme object); a pair \scm{(t . u)} of terms; a vector of terms; an instance of a
CHICKEN record type (\scm{define-record}) whose slots are terms. Vectors and record instances
together are the \emph{structures} (\scm{μkanren-structure?}, \texttt{:110}).
\end{definition}

A logic variable is a record with a single integer slot (\texttt{:41}):
\begin{scheme}
(define-record μkanren-var index)
(define (μkanren-var-working? α) (and (μkanren-var? α) (<= 0 (μkanren-var-index α))))
\end{scheme}
The sign of the index separates two kinds of variable:
\begin{itemize}
\item a \emph{working} variable has index $i \ge 0$; it lives during the search, its index is its
  slot in the substitution, and it prints as \scm{_i};
\item a \emph{reified} variable has index $-1,-2,\dots$; it exists only in answers, is created only by
  reification (\cref{sec:micro-reify}), and prints as a Greek letter \scm{α}, \scm{β}, \dots,
  \scm{ω}, then \scm{α1}, \dots (\texttt{:52–59}).
\end{itemize}
The helper \scm{μkanren-var-index>=0} (\texttt{:50}) maps both kinds onto $\mathbb{N}$:
$i\mapsto i$ for working variables, $-1\mapsto 0$, $-2\mapsto 1$, \dots{} for reified ones.

\begin{scheme}
(write (map make-μkanren-var '(0 1 7 -1 -2 -24 -25))) (newline)
\end{scheme}
\begin{answers}
(_0 _1 _7 α β ω α1)
\end{answers}

Indices are allocated by \scm{freshª} from a counter in the state (\cref{sec:micro-search}), so along
one branch of the search they increase monotonically: \emph{the index of a variable is its age},
lower meaning older. This fact, invisible to a user, is the most useful one for the λProlog
extension.

\begin{pitfall}[A variable is a record]
A \scm{μkanren-var} is itself a \scm{record-instance?}. Every structural walker of the engine
therefore tests \scm{μkanren-var?} \emph{before} the generic record case (e.g.\ \texttt{:117} before
\texttt{:124}, and the explicit \scm{(not (μkanren-var? w*))} at \texttt{:190}). Any new walker,
such as a scope check, must do the same, or it will descend into a variable as if it were a
two-slot structure.
\end{pitfall}

\begin{kanren}[Indices as levels]
λProlog's GENERIC rule reduces $\seq{\Sig}{\Prog}{\forall x\,G}$ to
$\seq{c{:}\tau,\Sig}{\Prog}{G[c/x]}$ for a token $c$ not in $\Sig$ \book{39}, and the query
\lp{sigma x\ pi y\ x = y} must fail \book{85--86}: with logic variables, a variable created
\emph{before} $c$ must never be bound to a term containing it. Implementations enforce this
with universe indices or levels attached to variables and constants~\cite{NadathurJayaramanKwon1995}.
In \micro{} the monotone index already is a level boundary: if the eigenvariable is created when the
counter is $n$, then ``created before $c$'' is exactly ``index $< n$''. The prototype of
\cref{sec:micro-forall} uses precisely this. Index order alone is not enough, though: a variable created
after the $\forall$ has exited must not see $c$ either, and lowering changes a variable's level after its
creation. The proposal therefore keeps a level per variable, in an \texttt{sbral} indexed like S and padded
only when the depth changes or the unifier creates a variable of another level (\cref{def:design-levels},
\cref{def:module-state}, \cref{sec:module-state}).
\end{kanren}

% ------------------------------------------------------------------------------------------------
\section{The state}\label{sec:micro-state}

\begin{definition}[State]\label{def:micro-state}
A state is a record of six fields (\texttt{:69}):
\begin{scheme}
(define-record μkanren-state vars-count S D A T tags)
(define μkanren-state-empty (make-μkanren-state 0 empty/sbral empty/sbral empty/sbral empty/sbral '()))
\end{scheme}
\end{definition}

\begin{longtable}{@{}>{\raggedright\arraybackslash}p{0.14\fullwidthlen}>{\raggedright\arraybackslash}p{0.36\fullwidthlen}>{\raggedright\arraybackslash}p{0.42\fullwidthlen}@{}}
\toprule
field & content & written by / read by \\
\midrule\endhead
\texttt{vars-count} & the next fresh index, i.e.\ the number of variables introduced on this branch
  & written and read by \scm{freshª} (\texttt{:532–537}) \\
\texttt{S} & the substitution: a skew-binary random-access list (sbral) of terms or of the sentinel
  \scm{*μkanren-unbound*} & written only by \scm{μkanren-state-update} (\texttt{:113–134}); read by every walk \\
\texttt{D} & disequality store: a list of \emph{prefixes}, each a list of \scm{(var . term)} pairs
  & \scm{≠°} (\texttt{:554–567}), re-verification after \scm{=°} \\
\texttt{A} & type store: a list of \scm{(var . tag)}; the tags are pairwise disjoint
  & \scm{symbol°}, \scm{number°}, user tags; verified after \scm{=°} \\
\texttt{T} & absento store: a list of \scm{(term . tag)}, the term a variable or a compound
  & \scm{absent°} (\texttt{:617–624}); verified after \scm{=°} \\
\texttt{tags} & tags asserted through A & written at \texttt{:481}, read by nobody in the engine (a dead field) \\
\bottomrule
\end{longtable}

Only S is a real sbral; D, A and T are plain lists (the \scm{empty/sbral} in the empty state is just
\scm{'()}). States are immutable and copied by functional update with keyword arguments
(\texttt{:76–83}),
\begin{scheme}
(μkanren-state-with s D: (append D* D))          ; a copy of s with a new D
\end{scheme}
except at the two hottest sites, \scm{μkanren-state-update} and \scm{freshª}, which call the
positional constructor directly. A positional destructuring macro \scm{μkanren-state-match} of fixed
arity six (\texttt{:96–103}) is used by six goals (\texttt{:130}, \texttt{:515}, \texttt{:533},
\texttt{:540}, \texttt{:555}, \texttt{:618}).

To see states, we use a small printer (not part of the engine) and run goals on the empty state:
\begin{scheme}
(define (dump s)   ; the fields of a state on one line; a tag is shown by its name
  (for-each (λ (k v) (display k) (display " ") (write v) (display "  "))
            '(vars-count S D A T)
            (list (μkanren-state-vars-count s)
                  (sbral->list (μkanren-state-S s))
                  (μkanren-state-D s)
                  (map (λ (p) (cons (car p) (μkanren-tag-name (cdr p)))) (μkanren-state-A s))
                  (map (λ (p) (cons (car p) (μkanren-tag-name (cdr p)))) (μkanren-state-T s))))
  (newline))
(define (states g) (§->list (g μkanren-state-empty)))   ; every answer state of g

(for-each dump (states (fresh° (q x y) (=° q (list x y)) (≠° x 1) (symbol° y) (absent° 'k x))))
\end{scheme}
\begin{answers}
vars-count 3  S ((_1 _2))  D (((_1 . 1)))  A ((_2 . μkanren-tag/sym87))  T ((_1 . k))
\end{answers}
The three variables \scm{q}, \scm{x}, \scm{y} are \scm{_0}, \scm{_1}, \scm{_2}; S holds one binding
(\scm{_0} is bound to the list \scm{(_1 _2)}); each store holds one constraint. The numeric suffix
of a tag name such as \scm{μkanren-tag/sym87} is a \scm{gensym} counter (\texttt{:33–34}): it identifies
the tag but its value depends on the program that loaded the module, so the same tag prints as
\scm{μkanren-tag/sym98} in another script of this chapter.

\begin{proposition}[The eq?-invariant]\label{prop:micro-eq}
States are compared by identity: \scm{(define (μkanren-state-equal? s s*) (eq? s s*))} (\texttt{:73}).
Unification returns \emph{the very same state object} when it adds no binding, and a freshly
allocated one otherwise (by inspection of \texttt{:160–181}: the non-binding cases call
\scm{return}, which closes over the input state, and only \scm{μkanren-state-update} allocates).
\end{proposition}

This is not an optimisation but semantics: every entailment test of the engine reads ``unification
gave back \scm{s}'' as ``the two terms are already equal''. The users are \scm{≠°} (\texttt{:559}: the
terms are already equal, so the disequality fails), \scm{=°} (\texttt{:544}: nothing new, so skip
re-verification), the re-verification of one disequality (\texttt{:460}), the absento test
\scm{μkanren-tag-unify} (\texttt{:394}) and subsumption (\texttt{:304}).

\begin{pitfall}[Preserving the eq?-invariant]
Any extension that runs inside unification, or after it, must return the input state itself when it
changes nothing. A hook that rebuilds the state with \scm{μkanren-state-with}, which always allocates,
on a no-op unification would make \scm{(≠° 1 1)} look satisfiable.
\end{pitfall}

\paragraph{How the record has grown.} The state record was widened in place several times: from
\scm{(substitution counter)} to \scm{(S D A T)} (commit \texttt{61bb868}), then \texttt{vars-count}
(\texttt{1e7994e}), then \texttt{tags} (\texttt{d118ddd}). Each step edited micro and the positional
state constructions of \texttt{test/microkanren.scm}. No sibling module (arith, date, show, unionfind)
ever builds a state: every store so far was added \emph{inside} micro. This is the precedent the
proposal of \cref{ch:design} has to reckon with.

% ------------------------------------------------------------------------------------------------
\section{The substitution and the walk}\label{sec:micro-subst}

\paragraph{S as an sbral.} \texttt{(aux fds sbral)} implements Okasaki's skew-binary random-access
lists: \scm{cons/sbral}, \scm{car/sbral}, \scm{cdr/sbral} in $O(1)$, \scm{sbral-ref} and
\scm{update/sbral} in $O(\log n)$, \scm{length/sbral} in $O(\log n)$. S stores the \emph{most recent
variable first}: the slot of the working variable of index $i$ sits at position $\ell-1-i$, where
$\ell$ is the length of S (\scm{μkanren-var-index/sbral}, \texttt{:86}). Extending S at the front
with \scm{cons/sbral} therefore keeps the positions of older variables valid when counted from the
end.

\paragraph{Lazy extension.} \scm{freshª} does not touch S. Binding the variable of index $i$
first pads S with the sentinel \scm{*μkanren-unbound*} until it has a slot $i$
(\scm{μkanren-var-extend/sbral}, \texttt{:89–94}). Hence
$\ell = 1+\max\{\text{bound index}\}$, and a variable whose index is $\ge \ell$ is unbound
\emph{without} having a slot.

\begin{scheme}
(for-each dump (states (fresh° (a b c d) (=° c 5))))
\end{scheme}
\begin{answers}
vars-count 4  S (5 #<aux.kanren.micro#μkanren-unbound> #<aux.kanren.micro#μkanren-unbound>)  D ()  A ()  T ()
\end{answers}
Four variables were introduced but S has three slots: \scm{c} (index 2) is bound to 5 at the front,
\scm{a} and \scm{b} hold the sentinel, \scm{d} (index 3) lies beyond S.

\begin{definition}[Walk]\label{def:micro-walk}
Let $\ell$ be the length of S. The \emph{walk} of a term $t$ (\scm{μkanren-state-find/SL},
\texttt{:139–145}) is computed by the loop $F(\beta_0,\beta)$, started with $\beta_0=\beta=t$:
\begin{enumerate}
\item if $\beta$ is a variable with $\mathit{idx}(\beta) \ge \ell$ (with $\mathit{idx}$ $=$ \scm{μkanren-var-index>=0}), answer $\beta$ (unbound, no slot);
\item if $\beta$ is a working variable, continue with $F(\beta, S[\ell-1-i])$;
\item if $\beta$ is the sentinel, answer $\beta_0$, the last variable visited;
\item otherwise answer $\beta$.
\end{enumerate}
\end{definition}

Case~4 catches non-variables but also \emph{reified} variables, which are not working: to the walk
a reified variable is a constant. Bindings are triangular (S maps a variable to a term that may
contain bound variables) and the walk never compresses paths. The variant with explicit \scm{S} and
\scm{l} hoists the $O(\log n)$ \scm{length/sbral} out of deep traversals; \scm{μkanren-state-find}
(\texttt{:147–148}) computes it once.

\paragraph{Deep walks.} Three exported procedures resolve a term completely:
\begin{itemize}
\item \scm{μkanren-state-find/value} (\texttt{:183–191}) returns a Scheme value, rebuilding pairs with
  \scm{cons}, vectors with \scm{list->vector} and records with \scm{(apply make-record-instance slots)};
  unbound variables stay \scm{μkanren-var} records;
\item \scm{μkanren-state-find/values} (\texttt{:632}) maps it over a list (the key of an aggregation group);
\item \scm{μkanren-state-find/repr} (\texttt{:193–204}) returns an \emph{expression} that rebuilds the
  value: a variable becomes its symbol, a symbol \scm{'sym}, a pair \scm{(cons a d)}, a vector
  \scm{(vector ...)}, a record \scm{(make-record-instance 'type ...)}.
\end{itemize}

\begin{pitfall}[Deep walks rebuild records]
Because \scm{find/value} rebuilds every record, the value it returns is \scm{equal?} but not
\scm{eq?} to the record that was unified. Here a record \scm{c} of the type \scm{eigen} of
\cref{sec:micro-forall} is bound to \scm{q} and read back:
\begin{scheme}
(let* ((c (make-eigen 99)) (r (car (§->list ((fresh° (q) (=° q c)) μkanren-state-empty))))
       (v (μkanren-state-find/value (make-μkanren-var 0) r)))
  (pretty-print (list (eq? v c) (equal? v c) (eigen? v))))
\end{scheme}
\begin{answers}
(#f #t #t)
\end{answers}
Constructed terms must be compared with predicates and \scm{equal?} (or by an id slot), never with
\scm{eq?}; \cref{sec:micro-forall} shows a scope check that goes wrong exactly this way.
\end{pitfall}

% ------------------------------------------------------------------------------------------------
\section{Unification}\label{sec:micro-unify}

\subsection{Binding and the occurs check}

Exactly one procedure creates a binding (\texttt{:113–134}, abridged):
\begin{scheme}
(define (μkanren-state-update α v s assocs #!key (unsafe #f))
  (define (occur? v*) ...)                  ; does α occur in the walk of v*?
  (and
    (or unsafe (not (occur? v)))
    (μkanren-state-match ((vc S D A T tags) s)
      (let* ((S* (μkanren-update/sbral α v S))
             (s* (make-μkanren-state vc S* D A T tags))
             (assocs* (cons `(,α . ,v) assocs)))
        `(,s* ,assocs*)))))
\end{scheme}
It answers \scm{#f} if $\alpha$ occurs in $v$ (the occurs check walks $v$, its pairs and its structures
through S), otherwise the pair \scm{(s* assocs*)}. The list \scm{assocs} accumulates \emph{the bindings
made by the current unification} as \scm{(var . term)} pairs: this is cKanren's ``prefix''. The
\scm{unsafe:} key is never passed, so the occurs check is always on. The procedure is not exported;
it is the only place where a binding can be refused before it enters the state, e.g.\ for violating a
scope; a check run right after each \scm{=°}, before the next goal, works as well (the route the proposal
takes, \cref{sec:design-hooks}).

\subsection{The unification algorithm}

\begin{definition}[Unification, \texttt{:160–181}]\label{def:micro-unify}
\scm{(μkanren-state-unify u v s assocs)} walks both sides to $u^*$, $v^*$ and tries these cases
\emph{in order}:
\begin{enumerate}
\item $u^* \mathrel{\texttt{eq?}} v^*$: success, nothing bound;
\item two strings that are \scm{string=?}: success;
\item two numbers that are \scm{μkanren-number-eqv?} (\texttt{:153–158}: \scm{eqv?}, but two NaN
  flonums are equal and non-real numbers are compared part by part): success;
\item two working variables of indices $i, j$: if $i=j$ success; if $i<j$ bind $u^* := v^*$, else
  $v^* := u^*$, i.e.\ \emph{the older variable is bound to the newer one};
\item one working variable: bind it to the other side;
\item two structures: unify the lists of their slots (\scm{μkanren-structure->list}, \texttt{:111}),
  where \scm{record->vector} puts the record's type name first;
\item two pairs: unify the cars, then the cdrs in the resulting state;
\item otherwise fail with \scm{#f}.
\end{enumerate}
The result is \scm{(s* assocs*)} or \scm{#f}. With \scm{assocs} $=$ \scm{'()}, the second component is
exactly the list of new bindings.
\end{definition}

Read as rules, the algorithm is the standard first-order syntactic one with occurs check; what is
specific to \micro{} is (a) the orientation of variable–variable bindings, (b) structures, (c) the
treatment of atoms. The exported procedure can be called on a raw state:
\begin{scheme}
(define q (make-μkanren-var 0))
; a state where _0 exists
(define s0 (μkanren-state-with μkanren-state-empty vars-count: 1))

(define r (μkanren-state-unify q '(1 2) s0 '()))
(pretty-print (cadr r))                                         ; the new bindings
(define r2 (μkanren-state-unify q '(1 2) (car r) '()))
; nothing new: the very same state
(pretty-print (list (eq? (car r2) (car r)) (cadr r2)))
(pretty-print (μkanren-state-unify q (list 1 q) s0 '()))        ; occurs check
\end{scheme}
\begin{answers}
((_0 1 2))
(#t ())
#f
\end{answers}
The first line is the one-element prefix \scm{((_0 . (1 2)))}, printed as Scheme prints it; the second
shows \cref{prop:micro-eq} at work; the third is the occurs check.

\paragraph{Orientation.} Unifying two fresh variables binds the older to the newer:
\begin{scheme}
(for-each dump (states (fresh° (a b) (=° a b))))
\end{scheme}
\begin{answers}
vars-count 2  S (_1)  D ()  A ()  T ()
\end{answers}
S has one slot, the one of \scm{_0}, and it holds \scm{_1}. A level discipline in the style of
λProlog would prefer the opposite orientation for a variable–variable pair (bind the younger variable
to the older one, so that such a binding never moves a variable out of its scope). Flipping the
orientation would not be enough, though: in \lp{sigma X\ pi y\ sigma Z\ (X = f Z, Z = y)} the binding
$X := f\,Z$ makes the young $Z$ reachable from the old $X$, so a real level discipline must also lower
the level of every variable of a term bound to an older variable (this is the shape of the fourth
query of \cref{sec:micro-forall}, with \scm{list} for $f$). With the orientation of \micro{} an
old variable can moreover be aliased to a young one, and a scope check must look through such chains
(\cref{sec:micro-forall}).

\paragraph{Structures and atoms.} Records unify as the list of their slots, so any
\scm{define-record} is a new term constructor for free; but a record and a vector with the same slot
list unify too. Numbers are compared with \scm{eqv?}, so an exact and an inexact number never unify.
\begin{scheme}
(define-record dbi n)
(pretty-print (record->vector (make-dbi 3)))
(pretty-print (°->list #f (fresh° (q) (=° (make-dbi q) (vector 'dbi 7)))))
(pretty-print (°->list #f (fresh° (q) (=° (vector 1 q) '(1 2)))))
(pretty-print (°->list #f (fresh° (q) (=° (/ 0. 0.) (/ 0. 0.)) (=° "ab" (string #\a #\b)) (=° q 'ok))))
(pretty-print (°->list #f (fresh° (q) (=° 1 1.0) (=° q 'ok))))
\end{scheme}
\begin{answers}
#(dbi 3)
((λ () 7))
()
((λ () 'ok))
()
\end{answers}
In order: the slot list of a record starts with its type name; a \scm{dbi} record unifies with the
vector \scm{#(dbi 7)}; a vector never unifies with a list; NaN equals NaN and equal strings unify
although they are distinct objects; \scm{1} and \scm{1.0} do not unify. Anything else (procedures,
bytevectors, \dots) unifies only with itself under \scm{eq?}.

\begin{pitfall}[Procedures are opaque]
A Scheme closure is an atom compared with \scm{eq?}. Higher-order abstract syntax with real Scheme
procedures (the λProlog \lp{abs x\ ...} rendered as \scm{(λ (x) ...)}) cannot be unified, walked,
constrained or reified. λ-terms must be first-order data, e.g.\ de Bruijn records such as \scm{dbi}
above (\cref{ch:lambdatree}).
\end{pitfall}

\begin{kanren}[Where higher-order unification would go]
The \scm{cond} of \cref{def:micro-unify} is closed: there is no dispatch table and no hook. A new term
kind with \emph{syntactic} equality needs nothing (it is a record); a term kind with a different
equality, such as $\Llam$ flexible terms $F\,x_1\dots x_n$ solved by pattern unification
(\cref{ch:hounif}) or nominal terms with suspended swappings~\cite{ByrdFriedman2007,UrbanPittsGabbay2004},
needs either a new clause here (before case~6) or a goal that solves the problem outside and emits
ordinary bindings (\cref{sec:micro-noedit}). The choice between them is made in \cref{ch:design}.
\end{kanren}

% ------------------------------------------------------------------------------------------------
\section{The constraint stores and their re-verification}\label{sec:micro-stores}

Unification only binds; the three stores keep \emph{negative} and \emph{type} information, and must
be re-checked whenever a binding is made.

\subsection{D: disequalities}

\begin{definition}[Prefix]\label{def:micro-prefix}
A \emph{prefix} is a list $d = ((x_1 . t_1) \dots (x_k . t_k))$ of variable–term pairs; it denotes
the constraint $\neg(x_1 = t_1 \wedge \dots \wedge x_k = t_k)$. The store D is a list of prefixes,
read as their conjunction.
\end{definition}

The goal \scm{(≠° u v)} (\texttt{:554–567}) runs a \emph{trial} unification of $u$ and $v$ with an
empty \scm{assocs} and reads the result:
\begin{itemize}
\item \scm{#f}: $u$ and $v$ can never be equal, so the constraint is dropped and the goal succeeds
  with the state unchanged;
\item the very same state (\cref{prop:micro-eq}): $u$ and $v$ are already equal, so the goal fails;
\item a new state with bindings $d$: the goal succeeds with the \emph{old} state extended by the
  prefix $d$ in D, unless a type tag of A or an absento of T makes $d$ vacuous (subsumption,
  below).
\end{itemize}
A tuple disequality gives one prefix with several pairs. After a later \scm{=°}, a prefix is
re-verified by trial unification of its pairs (\scm{μkanren-verify-D+}, \texttt{:458–461}): if the
trial fails the prefix is satisfied forever and dropped; if it adds nothing the prefix is violated and
the \scm{=°} fails; otherwise the prefix is replaced by its normalised version, the new trial bindings.

\begin{tutorial}[The life of a tuple disequality]
We assert $(x,y) \ne (1,2)$ and then bind $x$, watching D.
\begin{scheme}
(for-each dump (states (fresh° (q x y) (=° q (list x y)) (≠° (list x y) '(1 2)))))
(for-each dump (states (fresh° (q x y) (=° q (list x y)) (≠° (list x y) '(1 2)) (=° x 1))))
(for-each dump (states (fresh° (q x y) (=° q (list x y)) (≠° (list x y) '(1 2)) (=° x 3))))
(write (states (fresh° (q x y) (=° q (list x y)) (≠° (list x y) '(1 2)) (=° x 1) (=° y 2)))) (newline)
\end{scheme}
\begin{answers}
vars-count 3  S ((_1 _2))  D (((_2 . 2) (_1 . 1)))  A ()  T ()
vars-count 3  S (1 (_1 _2))  D (((_2 . 2)))  A ()  T ()
vars-count 3  S (3 (_1 _2))  D ()  A ()  T ()
()
\end{answers}
\begin{enumerate}
\item The trial unification of \scm{(_1 _2)} with \scm{(1 2)} would bind both variables, so D holds the
  single prefix \scm{((_2 . 2) (_1 . 1))}: ``not both''.
\item After \scm{(=° x 1)}, \scm{_1} is touched; the trial now binds only \scm{_2}, and the prefix is
  normalised to \scm{((_2 . 2))}: the constraint has become $y \ne 2$.
\item After \scm{(=° x 3)} the trial fails: the prefix can never be violated and disappears.
\item After \scm{(=° x 1)} and \scm{(=° y 2)} the trial binds nothing: the prefix is violated and the
  goal has no answer state.
\end{enumerate}
\end{tutorial}

\paragraph{Incremental re-verification.} After an \scm{=°}, only the prefixes that mention a
variable just bound (on either side of a pair) are re-verified (\scm{μkanren-verify-D/touched},
\texttt{:472–478}); the \scm{assocs} returned by unification is exactly the list needed to decide
which prefixes are touched. An untouched prefix cannot have become violated: its keys are unbound
variables, and since none of them was just bound, its trial unification would still bind its first key
(the reasoning of the comment at \texttt{:466–471}). It may, however, have become satisfied through the
occurs check deep in its terms; projection runs one last full \scm{μkanren-verify-D} to drop
those (\texttt{:466–471}, \texttt{:317}).

\subsection{A: disjoint type tags}

A \emph{tag} is a record \scm{(name def pred)} (\texttt{:31}): an identity, the quoted source of a
predicate (printed in answers), and the predicate itself. Type tags are made with
\scm{define-μkanren-tag} and turned into goals by \scm{μkanren-make-tag-A} (\texttt{:498–504}); the
built-in ones are (\texttt{:611–615})
\begin{scheme}
(define-μkanren-tag (μkanren-tag/sym v) (symbol? v))
(define-μkanren-tag (μkanren-tag/num v) (number? v))
(define symbol° (μkanren-make-tag-A μkanren-tag/sym))
(define number° (μkanren-make-tag-A μkanren-tag/num))
\end{scheme}
On a bound term the goal just tests the predicate; on an unbound variable it extends A. Tags are
\emph{mutually exclusive}: \scm{μkanren-ext-A} (\texttt{:431–436}) answers \scm{#f} when the variable
already carries a different tag. Adding a tag also drops the prefixes of D that it makes vacuous
(\scm{μkanren-subsume}, \texttt{:370–371}: a pair \scm{(α . u)} whose $\alpha$ is tagged by a predicate
rejecting $u$).
\begin{scheme}
(for-each dump (states (fresh° (q) (≠° q 'a) (number° q))))
(write (states (fresh° (q) (number° q) (symbol° q)))) (newline)
\end{scheme}
\begin{answers}
vars-count 1  S ()  D ()  A ((_0 . μkanren-tag/num99))  T ()
()
\end{answers}
A number can never be the symbol \scm{a}, so the disequality is gone; and no term is both a number and a
symbol.

\subsection{T: absento}

\leavevmode\scm{(absent° datum term)} (\texttt{:617–624}; note the argument order) states that \emph{datum}, which
may be any Scheme value, atomic or compound, does not occur in \emph{term}. The store holds
\scm{(term . tag)} pairs where the tag is an ad-hoc tag whose name is the datum itself. The term is a
variable, or a compound that could still become the datum (it is then ``watched whole'' and
decomposed only when it no longer can, \scm{μkanren-verify-T+}, \texttt{:410–419}). When a watched
variable is bound, the constraint moves to the parts:
\begin{scheme}
(for-each dump (states (fresh° (q x y) (absent° 'k q) (=° q (list x y)))))
(write (states (fresh° (q x y) (absent° 'k q) (=° q (list x y)) (=° y 'k)))) (newline)
\end{scheme}
\begin{answers}
vars-count 3  S ((_1 _2))  D ()  A ()  T ((_2 . k) (_1 . k))
()
\end{answers}
When a variable carrying an absento gets a type tag, the absento is converted into a disequality
(e.g.\ \scm{symbol° α} plus \scm{absent° 'closure α} becomes the prefix \scm{((α . closure))}), by
\scm{μkanren-subsume-A} and \scm{μkanren-update-D/T} (\texttt{:340–356}, \texttt{:480–485}). The
conversion happens for \emph{any} tag, and because the subsumption of D by the new tag runs before it
(\texttt{:493}), a converted prefix that the tag makes vacuous stays in the state and is dropped only at
projection:
\begin{scheme}
(for-each dump (states (fresh° (q) (absent° 'closure q) (symbol° q))))
(for-each dump (states (fresh° (q) (absent° 'closure q) (number° q))))
(pretty-print (°->list #f (fresh° (q) (absent° 'closure q) (number° q))))
\end{scheme}
\begin{answers}
vars-count 1  S ()  D (((_0 . closure)))  A ((_0 . μkanren-tag/sym98))  T ()
vars-count 1  S ()  D (((_0 . closure)))  A ((_0 . μkanren-tag/num99))  T ()
((λ (α) (assert (every (μ v (number? v)) (list α))) α))
\end{answers}
For a symbol the prefix is meaningful ($q \ne \cst{closure}$); for a number it is vacuous, and the
answer does not show it. A new store that reads D directly must expect such leftovers.

\subsection{The post-unification chain}

The goal \scm{=°} (\texttt{:539–552}, re-indented to fit the page) is where everything meets:
\begin{scheme}
(define ((=° u v) s)
  (μkanren-state-match ((vc S D A T tags) s)
    (define post°
      (λ1-match/first
        ((,s* ,assocs)
         (cond
           ((μkanren-state-equal? s s*) (✓° s*))          ; nothing bound: done
           ; D, touched prefixes only
           ((μkanren-verify-D/touched D s* assocs)
            => (μ D* (cond                                 ; then A, T, subsumptions
                       ((μkanren-verify-D/post D* A T s*) => ✓°)
                       (else (✗° s*)))))
           (else (✗° s*))))))
    (cond
      ((μkanren-state-unify u v s '()) => post°)
      (else (✗° s)))))
\end{scheme}
After a unification that binds something, the stores are re-verified in a hard-wired order:
\begin{center}
D (touched prefixes only) $\to$ A (in full) $\to$ D subsumed by A $\to$ T (in full) $\to$ D subsumed
by T $\to$ absentos of typed variables turned into D,
\end{center}
\begingroup\emergencystretch=1.5em
through \scm{μkanren-verify-D/post} (\texttt{:450–453}), \scm{μkanren-verify-A/post}
(\texttt{:424–\allowbreak 429}), \scm{μkanren-verify-T/post} (\texttt{:373}) and \scm{μkanren-subsume-T}
(\texttt{:352–356}), which rebuilds the state with the new D, A and T. Only D is incremental; A and T
are re-verified in full after every binding \scm{=°}. The whole chain is skipped only when nothing was
bound (\cref{prop:micro-eq}).\par\endgroup

\begin{kanren}[A new store has nowhere to plug in]
There is no list of verifiers to which a new store could append itself: the chain is a fixed sequence
of unexported procedures. λProlog features that need state of their own all hit this wall: a dynamic
program for AUGMENT \book{39}, scope information for GENERIC, a store of postponed non-pattern
equations that must be woken after each binding (\cref{ch:hounif}). What such a store requires of the
engine is (i) a place in the state that every goal copies forward, (ii) a call after each binding with
the new \scm{assocs}, (iii) a say in reification. \Cref{sec:micro-noedit} shows how far one gets
without (ii) and (iii); \cref{ch:design} proposes how to provide them.
\end{kanren}

% ------------------------------------------------------------------------------------------------
\section{Goals and the interleaving search}\label{sec:micro-search}

\subsection{Streams and goals}

\begin{definition}[Streams and goals]\label{def:micro-goal}
A stream (\texttt{(aux stream)}) is \scm{'()}, a \emph{mature} pair \scm{(state . stream)}, or an
\emph{immature} CHICKEN promise made with \scm{δ} (\scm{delay}, memoising). A \emph{goal} is a procedure
from a state to a stream of states. A \emph{relation} is a procedure from terms to a goal.
\end{definition}

The primitives (\texttt{:529–570}; an excerpt, comments ours) are few:
\begin{scheme}
; succeed once with the given state
(define ✓° list)
(define ✗° (K '()))                                    ; fail
(define ((freshª f) s)
  (μkanren-state-match ((vc S D A T tags) s)
    (let* ((α   (make-μkanren-var vc))
           (g   (f α))                                 ; the body is built only now
           (s*  (make-μkanren-state (add1 vc) S D A T tags)))
      (δ (g s*)))))                                    ; and suspended
(define ((orª f g) s) (append§/interleaved/2 (δ (f s)) (δ (g s))))
(define ((andª f g) s) (append-map§ g (δ (f s))))
\end{scheme}
The suffix \scm{ª} marks the \emph{applicative} procedure behind a macro; \scm{°} marks a goal, a
goal-making macro or a relation. Interleaving lives in \texttt{(aux stream)} (\texttt{aux.stream.scm:80–85}):
\begin{scheme}
(define append§/interleaved/2
  (λ (s1 s2)
    (cond
      ((null? s1) s2)
      ; SWAP on suspension
      ((promise? s1) (δ (append§/interleaved/2 s2 (force s1))))
      (else (cons§ (car s1) (append§/interleaved/2 (cdr s1) s2))))))
\end{scheme}
When the left stream is suspended, the two streams swap places; this is what makes disjunction
complete, and every \scm{δ} is a point where control can pass to another branch.

\paragraph{The macros} (\texttt{:575–593}, \texttt{:626–629}). \scm{(fresh° (x y) g ...)} nests one
\scm{freshª} per variable, each suspending once; \scm{(fresh° q ← (a b) g ...)} also binds \scm{q} to
\scm{(list a b)}. \scm{and°} is \emph{left}-nested,
\scm{(and° g1 g2 g3)} $=$ \scm{(andª (andª g1 g2) g3)}; \scm{or°} is \emph{right}-nested,
\scm{(or° g1 g2 g3)} $=$ \scm{(orª g1 (orª g2 g3))}. \scm{cond°} is \texttt{conde}:
\scm{(or° (and° g ...) ...)}, with no \texttt{else}. \scm{null°}, \scm{boolean°}, \scm{cons°} are
one-liners over \scm{=°} and \scm{or°}.

\begin{tutorial}[Fairness, and its shape]
Two infinite generators in a disjunction share the answers:
\begin{scheme}
(define-relation (nat° n) (or° (=° n 'z) (fresh° (m) (=° n `(s ,m)) (nat° m))))
(write (°->list/ground (take° 6 (fresh° (q) (or° (fresh° (n) (nat° n) (=° q `(left ,n)))
                                                (fresh° (n) (nat° n) (=° q `(right ,n)))))))) (newline)
\end{scheme}
\begin{answers}
((left z) (right z) (left (s z)) (right (s z)) (left (s (s z))) (right (s (s z))))
\end{answers}
The interleaving is fair between two branches, but \scm{or°} nests to the right, so with three
branches the first one competes against the \emph{pair} of the other two:
\begin{scheme}
(define-relation (numb° n k) (or° (=° n k) (numb° n (+ k 1))))
(write (°->list/ground (take° 12 (fresh° (q) (fresh° (n)
          (or° (and° (numb° n 0) (=° q (list 'a n)))
               (and° (numb° n 0) (=° q (list 'b n)))
               (and° (numb° n 0) (=° q (list 'c n))))))))) (newline)
\end{scheme}
\begin{answers}
((a 0) (a 1) (b 0) (a 2) (c 0) (a 3) (b 1) (a 4) (c 1) (a 5) (b 2) (a 6))
\end{answers}
In the long run branch \scm{a} gets half of the answers, \scm{b} and \scm{c} a quarter each (here 7, 3
and 2 of the first 12). Conjunction, on the other hand, is not fair:
\scm{(and° generator failing-test)} diverges, so as in every miniKanren generators go before tests and recursive calls last.
For instance, the \scm{append°} of \cref{sec:micro-reify} with its recursive call placed \emph{before}
\scm{(cons° a c rs)} still finds the four splits of \scm{(1 2 3)}, but then searches forever for a fifth
one under an unbounded run, while the order shown there terminates.
\end{tutorial}

\subsection{\texttt{define-relation} and the inverse-η delay}

Because \scm{and°} and \scm{or°} expand into \emph{procedure calls} of \scm{andª} and \scm{orª}, their
goal arguments are evaluated when the goal tree is built, before any state exists. A recursive relation
written with plain \scm{define} whose recursive call is not under a \scm{fresh°} therefore builds an
infinite tree:
\begin{scheme}
(define (ones° x) (or° (=° x 1) (ones° x)))          ; plain define: no delay
(write (°->list/ground (take° 3 (fresh° (q) (ones° q))))) (newline)
\end{scheme}
Run under \texttt{timeout 10}, this prints nothing and is killed (\texttt{exit=124}). The macro of
\texttt{:690} applies an inverse-η expansion and a suspension:
\begin{scheme}
(define-syntax-rule (define-relation (name arg ...) g ...) (define (name arg ...) (λ (s) (δ ((and° g ...) s)))))
\end{scheme}
so the body is built only when the goal meets a state, and then suspended:
\begin{scheme}
(define-relation (ones° x) (or° (=° x 1) (ones° x)))
(write (°->list/ground (take° 3 (fresh° (q) (ones° q))))) (newline)
\end{scheme}
\begin{answers}
(1 1 1)
\end{answers}
Plain \scm{define} remains fine when every recursive call sits under a \scm{fresh°}, whose body is under
\scm{λ (α)}; the relational interpreter of \texttt{test/microkanren-untagged.scm} is written that way.

\begin{pitfall}[Answer order is observable]
Each \scm{fresh°} variable and each relation call adds a suspension, \scm{or°} adds two. The test
suites pin the raw answers, \emph{including their order}. A new goal that adds or removes a single
\scm{δ} (for example a new \scm{=°}-like goal) can permute the answers of existing suites.
\end{pitfall}

\subsection{Soft cut, bounded runs, escape to Scheme}

\leavevmode\scm{(if° g? gt gf)} (\texttt{:596–603}) is LogicT's \texttt{ifte}, i.e.\ miniKanren's \texttt{conda}
for one clause: if \scm{g?} has at least one answer, \emph{all} its answers are continued with
\scm{gt}; otherwise \scm{gf} runs on the original state. \scm{(take° n g)} (\texttt{:605}) keeps the first
$n$ answers of \scm{g} (all of them if $n<0$); \scm{(take° 1 g)} is \texttt{once}. With \scm{gt} $=$ \scm{✗°}
and \scm{gf} $=$ \scm{✓°}, \scm{if°} is negation as failure:
\begin{scheme}
(write (°->list/ground (fresh° (q) (if° (or° (=° q 1) (=° q 2)) ✓° (=° q 3)))))
(newline)
(write (°->list/ground (fresh° (q) (if° ✗° ✓° (=° q 3))))) (newline)
(write (list (°->list/ground (fresh° (q) (=° q 1) (if° (=° q 2) ✗° ✓°)))
             (°->list/ground (fresh° (q) (=° q 2) (if° (=° q 2) ✗° ✓°)))
             (°->list/ground (fresh° (q) (if° (=° q 2) ✗° ✓°))))) (newline)
(write (°->list/ground (fresh° (q) (take° 2 (nat° q))))) (newline)
\end{scheme}
\begin{answers}
(1 2)
(3)
((1) () ())
(z (s z))
\end{answers}
The third line exhibits the classic unsoundness: with \scm{q} unbound, ``\scm{q} is not 2'' fails,
because the test succeeds by binding \scm{q}.

λProlog's \lp{if} and \lp{not} are defined with the cut \book{137}:
\begin{lprolog}
type if   o -> o -> o -> o.
if P Q R :- P, !, Q.
if P Q R :- R.

type not  o -> o.
not P :- P, !, fail.
not P.
\end{lprolog}
and the book's example of their non-logical character, ``\lp{?- X = 2, not (1 = X).} succeeds, whereas
\lp{?- not (1 = X), X = 2.} fails'' \book{137--138}, behaves the same way with \scm{if°}. There is one
difference: the cut keeps only the \emph{first} proof of \lp{P}, while \scm{if°} keeps all the answers
of its test, so λProlog's \lp{if P Q R} is \scm{(if° (take° 1 P) Q R)}:
\begin{scheme}
; X = 2, not (1 = X)
(write (list (°->list/ground (fresh° (q) (=° q 2) (if° (=° 1 q) ✗° ✓°)))
             ; the converse
             (°->list/ground (fresh° (q) (if° (=° 1 q) ✗° ✓°) (=° q 2)))))
(newline)
(write (list (°->list/ground (fresh° (q) (if° (or° (=° q 1) (=° q 2)) ✓° ✗°)))
             (°->list/ground
               (fresh° (q) (if° (take° 1 (or° (=° q 1) (=° q 2))) ✓° ✗°)))))
(newline)
\end{scheme}
\begin{answers}
((2) ())
((1 2) (1))
\end{answers}

\scm{(project° ((x* x) ...) g ...)} (\texttt{:626–627}) binds Scheme variables to the deep-walked values
of logic terms and runs \scm{g ...}: the non-relational escape hatch. An unbound variable is seen as a
\scm{μkanren-var} record.
\begin{scheme}
(write (°->list/ground (fresh° (q x) (=° x 20) (project° ((x* x)) (=° q (* 2 x*)))))) (newline)
(pretty-print (°->list #f (fresh° (q x) (project° ((x* x)) (=° q (list (μkanren-var? x*) (μkanren-var-index x*)))))))
\end{scheme}
\begin{answers}
(40)
((λ () (cons #t (cons 1 '()))))
\end{answers}

\begin{kanren}[AUGMENT and GENERIC as goals]
A goal sees the whole state and returns a stream of states, so it can do more than \scm{=°} does: it
can change the state on entry and post-process every answer state on exit. That is the shape that
both $D\imp G$ (add $D$, solve $G$, then restore the program on every answer of $G$) and $\forall x\,G$
(create $c$, solve $G$, then check that $c$ did not escape) require \book{39}. The two
prototypes of \cref{sec:micro-noedit} are written exactly this way, with \scm{append-map§} over the
answers of $G$.
\end{kanren}

% ------------------------------------------------------------------------------------------------
\section{Reification and the runners}\label{sec:micro-reify}

\subsection{From a state to a λ-expression}

An answer is not a substitution but a \emph{closed Scheme expression}
\begin{center}
\scm{(λ (α β ...) constraint-form ... term)}
\end{center}
which, applied to values for the free variables, re-checks the constraints with \scm{assert} and
\scm{deny} and returns the term. It is built by \scm{μkanren-project} (\texttt{:310–326}) in six steps:
\begin{enumerate}
\item deep-walk the query term;
\item name its free variables in order of first occurrence (left to right, car before cdr) by
  \emph{binding} each of them to a fresh reified variable $-1,-2,\dots$ in a scratch state
  (\scm{μkanren-state-reify}, \texttt{:207–222});
\item re-verify D in full;
\item \emph{purify}: drop every constraint that mentions a variable unreachable from the answer (still
  working in the scratch state); this is sound for disequalities, since such a variable can always be
  chosen to satisfy them;
\item drop the prefixes subsumed by others or made vacuous by a tag;
\item render with \scm{μkanren-form} (\texttt{:243–249}), sorting everything lexicographically.
\end{enumerate}

\begin{longtable}{@{}>{\raggedright\arraybackslash}p{0.3\fullwidthlen}>{\raggedright\arraybackslash}p{0.62\fullwidthlen}@{}}
\toprule
store & form in the answer \\
\midrule\endhead
D, prefix of one pair & \scm{(begin (deny (equal? α <repr>)))} \\*
D, prefix of $n$ pairs & \scm{(begin (deny (and (equal? α ...) (equal? β ...))))} \\*
A, grouped by tag & \scm{(assert (every (μ v (symbol? v)) (list α γ)))} \\*
T & \scm{(assert (absent? '<datum> <repr>))} \\
\bottomrule
\end{longtable}

\begin{scheme}
(pretty-print (°->list #f (fresh° (q x y) (=° q (list x y)) (≠° x 'a) (symbol° x) (absent° 'z y))))
(pretty-print (°->list #f (fresh° (q x y) (=° q (list x y)) (≠° (list x y) '(1 2)))))
(pretty-print (°->list #f (fresh° (q) (fresh° (h) (≠° h 3) (=° q 7)))))
(pretty-print (°->list/ground (fresh° (q x) (=° q (list x 1)) (≠° x 2))))
\end{scheme}
\begin{answers}
((λ (α β)
    (begin (deny (equal? α 'a)))
    (assert (every (μ v (symbol? v)) (list α)))
    (assert (absent? 'z β))
    (cons α (cons β '()))))
((λ (α β)
    (begin (deny (and (equal? α 1) (equal? β 2))))
    (cons α (cons β '()))))
((λ () 7))
((α 1))
\end{answers}
The third answer shows purification: the disequality on the unreachable \scm{h} is gone. The fourth
shows grounding: \scm{°->list/ground} evaluates only \scm{(lambda vars term)}, \emph{dropping the
constraints}, and applies it to the variable symbols themselves, so the free variable reads back as the
symbol \scm{α}. The show module \texttt{(aux kanren micro show)} parses the answer forms back without
\scm{eval} (\texttt{aux.kanren.micro.show.scm:58–72}), so a new kind of constraint in answers needs a
case there too.

\subsection{Running}

The runners are at \texttt{:686–713}:
\begin{scheme}
(define (°->§ g)
  (let* ((§ (δ (g μkanren-state-empty)))
         (P (μkanren-project (make-μkanren-var 0))))
    (map§ P §)))
(define-syntax-rule (μkanren-run (v n grounded) g ...) (°->list grounded (take° n (fresh° (v) g ...))))
\end{scheme}
\scm{(°->list grounded g)} forces \scm{°->§} and, if \scm{grounded}, \scm{eval}s each answer in micro's
own module environment; \scm{°->list/ground} is \scm{(°->list #t g)}. \scm{μkanren-run} is the
TRS-like \texttt{run}, with $n<0$ meaning ``all''. The query variable is always \emph{the variable of
index 0}, i.e.\ the first one the goal introduces. With the usual \scm{append°}:
\begin{scheme}
(define-relation (append° r s rs)
  (cond°
    ((null° r) (=° s rs))
    ((fresh° (a d c) (cons° a d r) (cons° a c rs) (append° d s c)))))
(pretty-print (μkanren-run (q -1 #t) (fresh° (x y) (append° x y '(1 2 3)) (=° q (list x y)))))
(pretty-print (μkanren-run (q 3 #f) (fresh° (x y z) (append° x y z) (=° q (list x y z)))))
\end{scheme}
\begin{answers}
((() (1 2 3)) ((1) (2 3)) ((1 2) (3)) ((1 2 3) ()))
((λ (α) (cons '() (cons α (cons α '()))))
 (λ (α β) (cons (cons α '()) (cons β (cons (cons α β) '()))))
 (λ (α β γ)
    (cons (cons α (cons β '())) (cons γ (cons (cons α (cons β γ)) '())))))
\end{answers}

\begin{pitfall}[The query variable is index 0]
\scm{°->§} hard-wires the empty initial state and projects variable 0. An extension that reserves
indices (for eigenvariables, registers, \dots) \emph{before} the first \scm{fresh°} shifts the query
variable; a module that needs a richer initial state needs its own runner, which is easy since
\scm{make-μkanren-var} and \scm{μkanren-project} are exported.
\end{pitfall}

\begin{kanren}[What answers will have to show]
Eigenvariables and a dynamic program never appear in today's answers. If a \scm{≠°} mentioning an
eigenvariable survives its scope, or a postponed higher-order equation is left at the end, reification
must either drop it or print it; purification (step 4) does not know about scopes. See
\cref{sec:micro-forall} for an example and \cref{ch:design} for the proposal.
\end{kanren}

% ------------------------------------------------------------------------------------------------
\section{Aggregation}\label{sec:micro-aggr}

Four SQL-like goals (\texttt{:631–684}) consume the \emph{whole} answer stream of a sub-goal with
\scm{foldr§} and turn it into Scheme data, so the sub-goal must be finite and the goals are
non-relational (\texttt{findall}/\texttt{bagof}-like). \scm{groupby°} groups the answers of \scm{g} by
the values of the keys and runs the continuation once per group \emph{on the original state}, with the
aggregate and the key names bound as Scheme variables; \scm{window°} is \texttt{OVER (PARTITION BY ...)}
and runs the continuation on each answer state of \scm{g}; \scm{set°} folds per key; \scm{enumerate°}
numbers the answers in stream order.
\begin{scheme}
(pretty-print (°->list/ground (fresh° (q) (fresh° (k v)
   (groupby° (((vs (λ (l) l)) v)) over (k) from
             (or° (and° (=° k 'a) (=° v 1)) (and° (=° k 'b) (=° v 2)) (and° (=° k 'a) (=° v 3)))
             => (=° q (list k vs)))))))
\end{scheme}
\begin{answers}
((b (2)) (a (1 3)))
\end{answers}
They are the precedent for goals that run a sub-search to completion, the mechanism a
\texttt{findall}-style predicate would use.

% ------------------------------------------------------------------------------------------------
\section{Extension points}\label{sec:micro-ext}

\subsection{Open versus closed}

\Cref{tab:micro-open} sorts the engine's parts by whether an importing module can use or extend them
without editing \texttt{aux.kanren.micro.scm}.

\begin{longtable}{@{}>{\raggedright\arraybackslash}p{0.47\fullwidthlen}>{\raggedright\arraybackslash}p{0.47\fullwidthlen}@{}}
\caption{The inventory of \micro{} for an extension.}\label{tab:micro-open}\\
\toprule
\textbf{open} (usable from outside) & \textbf{closed} (needs an edit of micro) \\
\midrule\endfirsthead
\toprule
\textbf{open} & \textbf{closed} \\
\midrule\endhead
Goals are plain procedures \emph{state} $\to$ \emph{stream}; all field accessors and
\scm{μkanren-state-with} are exported.
& The field list of the state (\texttt{:69}), its positional constructions (\texttt{:71}, \texttt{:83},
\texttt{:132}, \texttt{:536}), the six-ary \scm{μkanren-state-match} and its six users. \\
Raw unification \scm{μkanren-state-unify} returns the new state and the exact new bindings (but does
not re-verify D/A/T).
& The unification \scm{cond} (\texttt{:160–181}): no hook for a term kind with non-syntactic equality. \\
Walks \scm{μkanren-state-find}, \scm{/value}, \scm{/values}, \scm{/repr}.
& The binding site with occurs check, \scm{μkanren-state-update}, unexported: the only place to refuse a binding before it enters the state (a check right after each
\scm{=°} works as well, \cref{sec:design-hooks}). \\
Variables: \scm{make-μkanren-var}, \scm{μkanren-var-index}, \dots; a fresh unique integer by bumping
\texttt{vars-count}.
& The post-\scm{=°} chain (\texttt{:539–552}, \texttt{:373–478}), wired for D/A/T, its pieces unexported. \\
New first-order constructors: any \scm{define-record} unifies, reifies, works with \scm{≠°} and
\scm{absent°}.
& Reification (\texttt{:310–326}, \texttt{:243–249}): knows only D/A/T; purification and subsumption
are store-specific. \\
New disjoint types: \scm{define-μkanren-tag} with \scm{μkanren-make-tag-A}.
& \scm{°->§}: fixed empty initial state, fixed query index 0. \\
Projection \scm{μkanren-project}; the combinators \scm{freshª}, \scm{andª}, \scm{orª}, \scm{if°},
\scm{take°}, \scm{project°}; all of \texttt{(aux stream)}; \scm{define-relation}. & \\
\bottomrule
\end{longtable}

Two of the open entries in action: a user-defined disjoint type, and a goal written as a bare lambda
that reads the state.
\begin{scheme}
(define-μkanren-tag (tag/str v) (string? v))
(define string° (μkanren-make-tag-A tag/str))
(pretty-print (°->list #f (fresh° (q) (string° q))))
(pretty-print (map (λ (g) (°->list #f g))
                   (list (fresh° (q) (string° q) (=° q "hi"))
                         (fresh° (q) (string° q) (=° q 3))
                         (fresh° (q) (string° q) (symbol° q))
                         (fresh° (q) (string° q) (≠° q 3)))))
(pretty-print (°->list/ground (fresh° (q x y) (λ (s) ((=° q (μkanren-state-vars-count s)) s)))))
\end{scheme}
\begin{answers}
((λ (α) (assert (every (μ v (string? v)) (list α))) α))
(((λ () "hi")) () () ((λ (α) (assert (every (μ v (string? v)) (list α))) α)))
(3)
\end{answers}
The new tag is printed in answers, is disjoint from \scm{symbol°}, and subsumes \scm{(≠° q 3)}. The
last goal sees \texttt{vars-count} $=3$.

A CHICKEN rule constrains every new module: a macro exported from a module whose expansion refers to
an unexported binding fails at the use site. This is why \scm{freshª}, \scm{andª}, \scm{orª},
\scm{✓°}, \scm{✗°}, \scm{=°} and \scm{μkanren-group§} are exported (\texttt{:4–17}); macros such as
a future \texttt{∀°} or \texttt{⊃°} must follow the same discipline.

\subsection{Three tricks that need no edit, and their limits}\label{sec:micro-noedit}

\paragraph{Trick 1: wrapper goals that emit \texttt{=°}.} A solver written outside the engine (say,
higher-order pattern unification computing a most general unifier $\{F\mapsto t_1,\dots\}$) should
express its result as ordinary \scm{=°} goals, \scm{(and° (=° F t1) ...)}, so that each binding runs
the D/A/T chain. The tempting shortcut, calling the exported raw unification and returning its state,
silently bypasses the stores:
\begin{scheme}
; WRONG: bypasses the D/A/T re-verification
(define ((bind/raw x t) s)
  (let1 (r (μkanren-state-unify x t s '())) (if r (✓° (car r)) (✗° s))))
; right: the result is expressed as an ordinary =°
(define ((bind/emit x t) s) ((=° x t) s))

(define raw (§->list ((fresh° (q) (≠° q 'a) (bind/raw q 'a)) μkanren-state-empty)))
; q = a and q ≠ a!
(pretty-print (map (λ (s) (list (sbral->list (μkanren-state-S s)) (μkanren-state-D s))) raw))
(pretty-print (handle-exceptions e (list 'error ((condition-property-accessor 'exn 'message) e))
                (°->list #f (fresh° (q) (≠° q 'a) (bind/raw q 'a)))))
(pretty-print (°->list #f (fresh° (q) (≠° q 'a) (bind/emit q 'a))))
\end{scheme}
\begin{answers}
(((a) (((_0 . a)))))
(error "bad argument type - not a list")
()
\end{answers}
The raw version produces a state where S says $q=\cst{a}$ and D says $q\ne\cst{a}$; reification then
re-verifies D, gets \scm{#f} for ``violated'', and crashes. The emitting version fails, as it should.
\emph{Limits}: extra suspensions and repeated verification; and plain \scm{=°} applied by a user to
terms that need the special solver would still unify them syntactically, so a module must forbid that
or shadow \scm{=°}.

\paragraph{Trick 2: S-registers.} There is no way to add a field from outside: every goal of micro
rebuilds the state positionally, so a wrapper record would be lost at the next \scm{=°}. Dynamic
scoping (\scm{parameterize}) fails too, because promises are forced outside its extent, and so does a
side table keyed by the state, because a new state is allocated at every binding. But S itself is
copied forward by every goal and is persistent and branch-local. So: reserve a fresh index $n$ (bump
\texttt{vars-count}), bind variable $n$ by raw unification to an opaque record carrying the data, and
find the current value as the register of highest index. Raw unification is harmless here because the
register variable is fresh and unconstrained, and registers are never reachable from the query variable,
so they are never reified. Here the data is a λProlog-style program context, giving a prototype of
AUGMENT:
\begin{scheme}
; a register: an opaque record bound in S
(define-record ctx clauses)
; reserve index n, bind var n (raw unify) to a ctx record
(define (ctx-push s Γ)
  (let* ((n (μkanren-state-vars-count s))
         (s1 (μkanren-state-with s vars-count: (add1 n))))
    (car (μkanren-state-unify (make-μkanren-var n) (make-ctx Γ) s1 '()))))
; the most recent register, by a shallow walk downwards
(define (ctx-current s)
  (let loop ((i (sub1 (μkanren-state-vars-count s))))
    (cond ((< i 0) '())
          (else (let1 (v (μkanren-state-find (make-μkanren-var i) s))
                  (if (ctx? v) (ctx-clauses v) (loop (sub1 i))))))))
; D ⊃ G, D a list of clause procedures (λ (x) goal)
(define ((⊃ª D g) s)
  (let1 (Γ (ctx-current s))
    (append-map§ (λ (s2) (✓° (ctx-push s2 Γ))) (δ (g (ctx-push s (append D Γ)))))))
; a dynamic predicate: try every clause of the context
(define ((hyp° x) s)
  (let1 (Γ (ctx-current s)) ((fold-right (λ (c g) (orª (c x) g)) ✗° Γ) s)))

(define (is° v) (λ (x) (=° x v)))            ; the clause  hyp X :- X = v
(for-each (λ (g) (write (°->list/ground g)) (newline))
  (list (fresh° (q) (⊃ª (list (is° 'a)) (hyp° q)))
        (fresh° (q) (⊃ª (list (is° 'a)) ✓°) (hyp° q))
        (fresh° (q) (⊃ª (list (is° 'a)) (⊃ª (list (is° 'b)) (hyp° q))))
        (fresh° (q) (or° (⊃ª (list (is° 'a)) (hyp° q)) (hyp° q)))
        (fresh° (q) (fresh° (y) (⊃ª (list (λ (x) (=° x y))) (fresh° (u) (hyp° u) (=° u 7))) (=° q y)))))
\end{scheme}
\begin{answers}
(a)
()
(b a)
(a)
(7)
\end{answers}
In order: the assumed clause is used; outside its scope it is gone; nested assumptions stack, newest
first; two branches of a disjunction do not see each other's assumptions; and a free logic variable
\scm{y} of an assumed clause is \emph{shared} with the context, not renamed, as in λProlog
\book{86--88}. \emph{Limits}: the lookup scans the indices, $O(\texttt{vars-count}\cdot\log)$ per call;
registers must be allocated after the query variable; it is a library convention, not an engine
guarantee.

\subsection{Trick 3: a \texttt{∀} prototype with a post-filter scope check}\label{sec:micro-forall}

The eigenvariable is a record whose id is a reserved index, hence unique; distinct eigenvariables never
unify (their id slots differ). The scope check is a post-filter on each answer state of the body: every
variable of index $< n$, i.e.\ older than the eigenvariable, is deep-walked and must not contain it.
\begin{scheme}
(define-record eigen id)
(define (contains? c t)
  (cond
    ; find/value REBUILDS records: compare ids, not identity
    ((eigen? t) (= (eigen-id t) (eigen-id c)))
    ((pair? t) (or (contains? c (car t)) (contains? c (cdr t))))
    ((vector? t) (any (λ (u) (contains? c u)) (vector->list t)))
    ((and (record-instance? t) (not (μkanren-var? t)))
     (any (λ (u) (contains? c u)) (vector->list (record->vector t))))
    (else #f)))
(define ((∀ª f) s)
  (let* ((n (μkanren-state-vars-count s))
         (c (make-eigen n))                        ; unique: index n is reserved
         (s* (μkanren-state-with s vars-count: (add1 n)))
         (escapes? (λ (s1) (let loop ((i 0))
                             (and (< i n)
                                  (or (contains? c (μkanren-state-find/value (make-μkanren-var i) s1))
                                      (loop (add1 i))))))))
    (append-map§ (λ (s1) (if (escapes? s1) '() (✓° s1))) (δ ((f c) s*)))))
(define-syntax-rule (∀° (x) g ...) (∀ª (λ (x) (and° g ...))))

(for-each (λ (g) (pretty-print (°->list #f g)))
  (list (fresh° (q) (∀° (x) (=° q x)))                              ; 1. ∃Q ∀x. Q = x
        (fresh° (q) (∀° (x) (fresh° (y) (=° y x))) (=° q 'ok))      ; 2. ∀x ∃Y. Y = x
        (fresh° (q) (∀° (x) (∀° (y) (=° x y))))                     ; 3. ∀x ∀y. x = y
        (fresh° (q) (∀° (x) (fresh° (y) (=° q (list y)) (=° y x)))) ; 4. alias
        (fresh° (q) (∀° (x) (fresh° (z) (=° q (list 'f z)))))       ; 5.
        (fresh° (q) (∀° (x) (≠° q x)))))                            ; 6. ≠ survives
\end{scheme}
\begin{answers}
()
((λ () 'ok))
()
()
((λ (α) (cons 'f (cons α '()))))
((λ (α) (begin (deny (equal? α (make-record-instance 'eigen 1)))) α))
\end{answers}

\begin{tutorial}[Reading the six answers]
\begin{enumerate}
\item $\exists Q\,\forall x.\,Q=x$ is the book's \lp{sigma x\ pi y\ x = y}, ``not provable'' \book{85}.
  Here \scm{q} is \scm{_0}; \scm{∀°} reserves index 1 and makes \scm{c} $=$ \scm{(eigen 1)}; the body
  binds \scm{_0 := c}; the post-filter walks \scm{_0}, finds \scm{c}, and discards the state.
\item $\forall x\,\exists Y.\,Y=x$ succeeds: \scm{y} has index 2 $\ge n=1$, so it is inside the scope.
\item Two eigenvariables \scm{(eigen 1)} and \scm{(eigen 2)} differ in their id slot: unification fails.
\item The escape is indirect: \scm{_0 := (list _2)} and \scm{_2 := c}. Because the check deep-walks the
  old variable, it sees \scm{(list c)} and rejects the state. This is the chain created by the
  older-to-newer orientation of \cref{sec:micro-unify}.
\item A younger variable left unbound is harmless and is reified as \scm{α}.
\item A disequality mentioning \scm{c} survives in D and is reified with the eigenvariable in it: the
  check filters bindings, not constraints.
\end{enumerate}
\end{tutorial}

Why is checking at the end of the body enough for soundness? After the scope closes, \scm{c} is
lexically unreachable and, by the check, no term reachable from an older variable contains it, so no
later goal can bring it back.\sidenote{This argument is informal; the prototype was tested, not proved.} What
is lost is \emph{early pruning}: a bad binding is detected only when the body has produced an answer,
at a cost of $O(n\cdot|\text{term}|)$ per answer. Eager checking needs the scope test at the binding
site (closed), or a post-unification hook with a level map, or a module-level \scm{=°} that shadows
micro's and checks its bindings against the open scopes kept in an S-register; the last option covers
only goals written with the new module's \scm{=°}.

\begin{pitfall}[The scope check must not use \texttt{eq?}]
With the first line of \scm{contains?} written as \scm{((eigen? t) (eq? t c))}, the same query as
answer~1 wrongly succeeds, because \scm{find/value} rebuilt the record (\cref{sec:micro-subst}):
\begin{scheme}
(pretty-print (°->list #f (fresh° (q) (∀° (x) (=° q x)))))
\end{scheme}
\begin{answers}
((λ () (make-record-instance 'eigen 1)))
\end{answers}
\end{pitfall}

\begin{kanren}[What the tricks leave open]
The three tricks show that the shape of AUGMENT and GENERIC fits \micro{} today: a goal that changes
the state on entry and filters or restores it on every answer. What they cannot provide is what only
the engine can do: a \emph{named} place in the state (instead of a register scan), a call after every
binding (for eager scope checks and for waking postponed equations), and a say in reification (to drop
out-of-scope constraints and print new ones). These are the requirements \cref{ch:design} turns into a
proposal, and \cref{ch:module} into a module.
\end{kanren}

% ------------------------------------------------------------------------------------------------
\section{Summary}\label{sec:micro-summary}

\begin{itemize}
\item Terms are Scheme data; logic variables are records with an integer index, working ($\ge 0$,
  printed \scm{_i}) or reified ($<0$, printed \scm{α}, \scm{β}, \dots). Indices grow monotonically along a
  branch, so an index is an age and doubles as a level boundary.
\item A state has six fields \texttt{vars-count}, S, D, A, T, \texttt{tags} (the last one dead). States
  are compared by \scm{eq?}: unification returns the same state iff it binds nothing, and every
  entailment test relies on this.
\item S is an sbral, most recent variable first, extended lazily; the walk follows triangular bindings
  and treats reified variables as constants; deep walks rebuild records, losing their identity.
\item Unification is first-order with occurs check; the only binding site is
  \scm{μkanren-state-update}; the older variable is bound to the newer; vectors and records unify as
  slot lists (so records and same-shaped vectors unify); numbers by a NaN-reflexive \scm{eqv?}; closures
  only by \scm{eq?}.
\item D holds prefixes (negated conjunctions of bindings), re-verified incrementally on the touched
  prefixes; A holds disjoint type tags, T absento constraints, both re-verified in full; the chain after
  \scm{=°} is fixed and has no slot for a new store.
\item Goals are \emph{state} $\to$ \emph{stream}; \scm{orª} interleaves by swapping on suspensions
  (\scm{or°} nests right, so three branches get $\frac12,\frac14,\frac14$); conjunction is not fair;
  \scm{define-relation} delays recursive relations; \scm{if°} is \texttt{conda} (it keeps every answer
  of its test, so λProlog's cut-based \lp{if} is \scm{if°} on a \scm{(take° 1 ...)} test), \scm{take°} bounds a run,
  \scm{project°} escapes to Scheme; answer order is pinned by the tests.
\item Answers are closed expressions \scm{(λ (vars) constraints ... term)}, purified of unreachable
  variables; grounding drops the constraints; the query variable is index 0.
\item Open to extension: goals as procedures, raw unification, walks, records as constructors, user
  type tags, projection, the combinators. Closed: the state's fields, the unification \scm{cond}, the
  binding site, the post-\scm{=°} chain, reification, the runner.
\item Without editing micro one can: emit \scm{=°} from external solvers (never return raw-unified
  states), carry branch-local data in S-registers (a working AUGMENT prototype), and check
  eigenvariable scope by post-filtering the answers of a \scm{∀} body (a working GENERIC prototype). Their
  limits — scans, late pruning, constraints outliving scopes — are the requirements of \cref{ch:design}.
\end{itemize}

\chapter{Bringing λProlog to \texorpdfstring{\texttt{(aux kanren micro)}}{(aux kanren micro)}: the design}\label{ch:design}

% ------------------------------------------------------------------------------------------------
\section*{What and why}

\Cref{ch:micro} ended with a precise list of what an importing module cannot obtain from \micro{}
today: a named place in the state for its own data, a call after every binding (to check the scope of
eigenvariables eagerly and to wake postponed equations), and a say in reification (to print new
constraints and to hide what must not escape). The three tricks of \cref{sec:micro-noedit} show that
\rulename{augment} and \rulename{generic} \emph{fit} the shape of a μKanren goal; they do not make them sound under every
goal of the engine, nor cheap.

This chapter turns that observation into a design. It proceeds in four steps.
\begin{enumerate}
\item \textbf{Requirements} (\cref{sec:design-req}). The nine acceptance items A1--A9 against which the
  proposal was built and tested, stated as properties of a run of \scm{°->list}.
\item \textbf{The design space} (\cref{sec:design-space}). The four meanings of unification on λ-terms
  compared in \cref{sec:lkanren-compare}, and why higher-order pattern unification with universe levels is
  the one retained.
\item \textbf{Three prototypes} (\cref{sec:design-prototypes}). Three complete implementations of the
  same requirements, which differ in where the code sits with respect to \micro: a layer that does not
  touch the engine, a small generic protocol in the engine, and a native engine. They were compared by
  three reviews (engineering, fidelity to the book, pedagogy); the protocol won.
\item \textbf{The extension protocol} (\cref{sec:design-protocol}). The edit to \micro{} in full: the
  seventh state field, the registry, every hook with its contract and call site, the two new primitives,
  the one generic change to \scm{absent°}, and the two invariants a client must keep.
\end{enumerate}
The chapter closes with what the edit costs a first-order program (\cref{sec:design-cost}) and with
the history of the review rounds that hardened the proposal (\cref{sec:design-reviews}).
The λ layer that sits on the protocol, \scm{(aux kanren lambda)}, with its unifier, its \scm{∀°} and
\scm{⊃°} and its treatment of micro's stores, is the subject of \cref{ch:module}; the book's examples
run on it are in \cref{ch:examples}; the plan for merging is in \cref{ch:roadmap}. In the repository the
proposal lives under \texttt{doc/hol-kanren/prototype/} (the module, \texttt{micro.diff},
\texttt{show.diff}, the design notes \texttt{DESIGN.md}, the tests, the examples, the benchmarks and the
scripts that build and run them, \cref{sec:roadmap-run}). After merging, the module and its suites
become \texttt{src/}\allowbreak\texttt{aux.kanren.lambda.scm} and
\texttt{src/}\allowbreak\texttt{test/}\allowbreak\texttt{microkanren-lambda*.scm}.

Every Scheme snippet of this chapter with an answer was run against the proposal (the edited micro and
\scm{(aux kanren lambda)}) unless it says otherwise, and its output is pasted verbatim. Snippets of the
form \scm{(show "label" query)} are quoted from the proposal's \texttt{examples.scm}, where \scm{show}
prints the label as a comment and then the answers; their outputs are quoted from
\texttt{examples.out}.

% ================================================================================================
\section{Requirements}\label{sec:design-req}

The proposal was written against nine acceptance items, A1 to A9. Each one is a property of the
answers that \scm{°->list} returns for a class of goals, illustrated by the book's own examples with the
book's own answers. This section states each item in a form precise enough to be tested, then gives
its intuition and one worked example. The general statements are \emph{tested}, by 83 hand-written cases
with 320 assertions, not proved.

\subsection{Setting and notation}\label{sec:design-notation}

\paragraph{Goals and terms.} The goal language is that of \micro{} (\scm{=°}, \scm{≠°}, \scm{absent°},
\scm{symbol°}, \scm{number°}, \scm{fresh°}, \scm{and°}, \scm{or°}, \scm{cond°}, relations) extended with
three forms: \scm{(∀° (x) G)}, the universal goal \lp{pi x\ G}; \scm{(⊃° (D ...) G)}, the implication
\lp{D => G}; and relations declared with \scm{define-dynamic-relation}, the only ones that take
hypotheses. The term language is micro's (Scheme data and logic variables) extended with abstractions
\scm{(ƛ (x) t)}, applications \scm{(@ t u ...)} and eigenvariables. We write $\mathrm{Ans}(G)$ for the list
returned by \scm{(°->list #f G)}, and say that $G$ \emph{fails} when $\mathrm{Ans}(G)$ is empty.

\paragraph{Reference semantics.} The specification is the goal-directed search of the book for
hereditary Harrop formulas, that is, the six rules \rulename{and}, \rulename{or}, \rulename{instan}, \rulename{generic}, \rulename{augment} and \rulename{backchain}
(\cref{fig:rules-right}; \book{77}, \book{82}, \book{124}; see \cref{sec:hh-rules} and \cref{app:rules}). Uniform proofs are complete for intuitionistic provability of
these formulas (\cref{thm:hh-intuitionistic}), so implementing these rules faithfully \emph{is} the
semantics. Where λProlog says nothing (disequality, \emph{absento}, types), the reading adopted extends
micro's own semantics over λProlog's term model (item A8).

\begin{definition}[Universe levels]\label{def:design-levels}
A state carries a \emph{depth} $d\in\mathbb N$, the number of \scm{∀°} goals whose body is running.
Every logic variable $X$ has a \emph{level} $\ell(X)$, the depth at its creation unless later lowered.
An eigenvariable $c$ created by a \scm{∀°} at depth $d$ has $\ell(c)=d+1$. The names that the unifier
itself opens when it goes under a binder have level $\infty$.
\end{definition}
These levels are Nadathur, Jayaraman and Kwon's tags minus one: their query starts at tag 1
(\cref{fig:rules-ui}), ours at depth 0.

The level is the linear trace of Miller's mixed prefix: a variable of level $d$ stands for an
existential quantifier that follows the first $d$ universal quantifiers of its branch. The intended
reading is Nadathur, Jayaraman and Kwon's~\cite{NadathurJayaramanKwon1995} (\cref{def:lkanren-ui}): a term may be substituted for $X$ only
if every constant in it has a level $\le \ell(X)$.

\begin{example}
In \scm{(fresh° (q) (∀° (x) (fresh° (y) G)))}, \scm{q} has level 0, the eigenvariable \scm{x} level 1,
and \scm{y} level 1. So \scm{y} may be bound to a term containing \scm{x}, while \scm{q} may not. This
is $\exists q\,\forall x\,\exists y.\,G$ read off the creation depths.
\end{example}

\begin{definition}[Order independence]\label{def:design-order}
A conjunction is \emph{order independent} when permuting its conjuncts changes at most the order of its
answers, not the set of answers (taken as data, with their residual constraints).
\end{definition}

Order independence is not an item of its own. It runs through all of them, because it is what
Nadathur, Jayaraman and Kwon prove for their tagged interpreter: the choice of the goal to reduce next
is irrelevant (their Prop.~4.1, \cref{sec:lkanren-impl}). A miniKanren conjunction is fair to no
conjunct, so any order-dependence is visible to the user. Most defects found by the reviews of
\cref{sec:design-reviews} were violations of it, \enquote{in one goal order of three}.

\subsection{The nine items}\label{sec:design-items}

\begin{definition}[A1: the scope of \texorpdfstring{$\forall$}{∀}]\label{def:design-a1}\leavevmode
\begin{enumerate}[label=(\alph*)]
\item \textbf{\rulename{generic}.} In a state of depth $d$, \scm{(∀° (x) G)} runs $G[c/x]$ at depth $d+1$, with
  $c$ a new eigenvariable of level $d+1$, and gives every answer state of $G$ back the depth $d$.
\item \textbf{Scope invariant.} In every state reached by the search, every binding $X:=t$ of the
  substitution satisfies the three binding conditions of \cref{def:lkanren-ui} (stated also in \cref{sec:hh-scope} and
  \cref{def:rules-binding}):
  (1)~$X$ does not occur in $t$; (2)~every eigenvariable occurring in $t$ has a level $\le \ell(X)$;
  (3)~every unbound variable $Y$ of $t$ has $\ell(Y)\le\ell(X)$ (\emph{lowering}; for a flexible
  subterm $Y\,\bar a$, the arguments not visible at $\ell(X)$ are pruned).
\item \textbf{Universality of the invariant.} (b) holds for every binding, whatever goal made it: the
  goals of the new module, micro's own \scm{=°}, every relation built on it (\scm{cons°}, the user's
  relations), and the unifier's internal bindings.
\end{enumerate}
\end{definition}

Intuitively, (a) is the book's \rulename{generic} rule \book{82} (\cref{def:hh-eigen}) and (b) is the reason why
it is sound: no eigenvariable can reach a variable created before its \scm{∀°}, directly or through an
alias, so a proof never assumes that a fresh name equals an older term. Restoring the depth on each
answer matters: a variable created \emph{after} the \scm{∀°} has exited must not see $c$, which a creation
timestamp alone would not ensure.

The invariant (b) is stated on single bindings, but it controls whole values. Suppose $X:=t$ and $Y$
occurs unbound in $t$; by (3), $\ell(Y)\le\ell(X)$, so if $Y$ is bound later to some $t'$, condition
(2) for $Y:=t'$ admits only eigenvariables of level $\le\ell(Y)\le\ell(X)$. By induction on the
chain of bindings, the fully walked value of $X$ never holds an eigenvariable of level $>\ell(X)$:
checking each binding once, when it is made, is enough. Condition (3) is what makes the induction go
through; without lowering, an alias created before the eigenvariable could receive it later.

Item (c) is what separates a sound design from a layered one, as
\cref{sec:design-d1} shows. The book's two classical queries \book{85} and two escapes through aliases:
\begin{scheme}
(show "∃Y∀x. Y = x  and  ∀x∃Y. Y = x"
      (list (°->list #f (fresh° (q) (∀° (x) (=° q x))))
            (°->list #f (fresh° (q) (∀° (x) (fresh° (y) (=° y x))) (=° q 'ok)))))

(show "escape through an alias (q = (y), y = x), and through micro's cons°: both fail"
      (list (°->list #f (fresh° (q) (∀° (x) (fresh° (y) (=° q (list y)) (=° y x)))))
            (°->list #f (fresh° (q) (∀° (x) (cons° x '() q))))))
\end{scheme}
\begin{answers}
;; ∃Y∀x. Y = x  and  ∀x∃Y. Y = x
(() ((λ () 'ok)))

;; escape through an alias (q = (y), y = x), and through micro's cons°: both fail
(() ())
\end{answers}
In the third query the binding \scm{y := x} is legal (\scm{y} has level 1), but \scm{q := (y)} was made
earlier with \scm{y} unbound; condition (3) lowered \scm{y} to level 0 at that moment (by binding it
to a fresh variable of level 0), so the later
\scm{y := x} violates (2). The fourth query never mentions the new module's unifier: \scm{cons°} is
micro's, and its \scm{=°} is checked all the same. The suite adds Nadathur, Jayaraman and Kwon's own
example (\cref{sec:hh-scope}, and the tutorial after \cref{def:lkanren-ui}), where lowering reaches a variable created by
backchaining on an assumed clause: $\exists x\forall y(q\,y\imp p\,x)$ with the clause
$\forall z(q\,z\imp p\,(d\,z))$ fails, as it must.

\begin{definition}[A2: \rulename{augment} and \rulename{backchain}]\label{def:design-a2}\leavevmode
\begin{enumerate}[label=(\alph*)]
\item \textbf{\rulename{augment}.} $\Sig;\Prog\longrightarrow D\imp G$ reduces to $\Sig;\Prog\cup\{D\}\longrightarrow
  G$: \scm{(⊃° (D ...) G)} runs $G$ with the clauses $D$ added to a context $\Gamma$ kept in the state, and
  gives every answer state of $G$ back the old $\Gamma$. Hypotheses are visible exactly during $G$, and
  the branches of an \scm{or°} do not see each other's hypotheses.
\item \textbf{Clause order.} $\Gamma$ is a list, extended at the front. A dynamic relation $p$ called on
  $\bar t$ is the disjunction of its clauses in $\Gamma$, most recent first, followed by its static body
  \book{78}.
\item \textbf{Variables of a clause.} The logic variables free in $D$ belong to the context: they are
  shared by every use of the clause and with the goal around it \book{86--88}. The clause's own
  universally quantified variables, \scm{(∀ (z) ...)}, are renamed at each use, at the current level
  (Nadathur, Jayaraman and Kwon's \enquote{new tagged instance}).
\item \textbf{Static relations are closed.} Assuming a clause for a relation not declared with
  \scm{define-dynamic-relation} is an error, not a silent no-op.
\end{enumerate}
\end{definition}

Point (b) is a statement about \emph{order}, which a miniKanren usually leaves to interleaving. The book
fixes it with two queries against the program \lp{p 1} \book{78}:
\begin{lprolog}
?- p 2 => p 3 => p X.
X = 3;
X = 2;
X = 1
?- (p 2 & p 3) => p X.
X = 2;
X = 3;
X = 1
\end{lprolog}
\begin{scheme}
(define-dynamic-relation (p° x) (=° x 1))
(show "clause order under augmentation (p.94)"
      (list (°->list/datum (fresh° (q) (⊃° ((p° 2)) (⊃° ((p° 3)) (p° q)))))
            (°->list/datum (fresh° (q) (⊃° ((p° 2) (p° 3)) (p° q))))))
\end{scheme}
\begin{answers}
;; clause order under augmentation (p.94)
((3 2 1) (2 3 1))
\end{answers}
(The label's \enquote{p.94} is the PDF page, \book{78}; see \cref{sec:intro-conventions}.) Each clause answers without
suspending, so the flat disjunction returns its disjuncts in order, and the order is the book's. Point
(c) is the subtle one: in \lp{pi z\ (q z => p (d z))} the variable \lp{z} is the clause's own, while a
free \lp{Y} in an assumed \lp{p Y} is the goal's. \Cref{sec:hh-kanren-context} explains why $\Gamma$ must
live in the state rather than in a Scheme parameter.

\begin{definition}[A3: \texorpdfstring{$\forall$}{∀} and \texorpdfstring{$\imp$}{⊃} together]\label{def:design-a3}
The book's type inference for untyped λ-terms (\cref{sec:lambdatree-typeof}), whose abstraction case is
\lp{typeof (abs M) (arr A B) :- pi x\ typeof x A => typeof (M x) B}, computes principal types: for a
closed term it answers the book's type up to the names of type variables, and it fails on terms that
have no simple type.
\end{definition}

This item exercises all of A1, A2 and A4 in one clause: \scm{∀°} makes the bound variable a scoped
constant, \scm{⊃°} assumes its type, \scm{(@ R x)} β-reduces the object-level body, and the type
variables \scm{A}, \scm{B}, created outside the \scm{∀°}, must never be bound to a type containing
\scm{x}. In the module's syntax:
\begin{scheme}
(define-dynamic-relation (typeof° m t)
  (cond°
    ((fresh° (M N B) (=° m `(app ,M ,N)) (typeof° M `(arr ,B ,t)) (typeof° N B)))
    ((fresh° (R A B) (=° m `(abs ,R)) (=° t `(arr ,A ,B))
       (∀° (x) (⊃° ((typeof° x A)) (typeof° (@ R x) B)))))))
\end{scheme}
Run on the combinators S and I and on $\lambda x.\,x\,x$, written in the object syntax with
\scm{ƛ} for the meta-level binder:
\begin{scheme}
(show "typeof S, I and λx.x x"
      (list (°->list/datum (fresh° (q) (typeof° `(abs ,(ƛ (x) `(abs ,(ƛ (y) `(abs ,(ƛ (z)
                                          `(app (app ,x ,z) (app ,y ,z)))))))) q)))
            (°->list/datum (fresh° (q) (typeof° `(abs ,(ƛ (x) x)) q)))
            (°->list/datum (fresh° (q) (typeof° `(abs ,(ƛ (x) `(app ,x ,x))) q)))))
\end{scheme}
\begin{answers}
;; typeof S, I and λx.x x
(((arr (arr α (arr β γ)) (arr (arr α β) (arr α γ)))) ((arr α α)) ())
\end{answers}
The third query fails by the occurs check: the type of \scm{x} would have to be an arrow type
containing itself.

\begin{definition}[A4: λ-tree syntax and β]\label{def:design-a4}\leavevmode
\begin{enumerate}[label=(\alph*)]
\item \textbf{α.} \scm{(ƛ (x) t)} builds an abstraction whose bound variable is a Scheme variable;
  terms built alike are \scm{equal?}, and unification is modulo α.
\item \textbf{β.} \scm{(@ (ƛ (x) t) s)} is $t[s/x]$ at construction; an application whose head becomes an
  abstraction through a later binding is reduced on the fly, in unification and in every deep walk.
\item \textbf{η.} Terms are kept η-short; η-expansion happens on demand when one side of an equation is
  an abstraction.
\item \textbf{Object-level substitution is meta-level β} (\cref{sec:lambdatree-beta}): the book's
  evaluators, \texttt{copy} and \texttt{subst} give the book's results.
\end{enumerate}
\end{definition}

\begin{scheme}
(show "α-equivalence: λa.a = λb.b, and the β-redex (λx.g x x) a"
      (list (°->list #f (fresh° (q) (=° (ƛ (a) a) (ƛ (b) b))))
            (°->list/datum (fresh° (q f) (=° f (ƛ (x) `(g ,x ,x))) (=° q (@ f 'a))))))
\end{scheme}
\begin{answers}
;; α-equivalence: λa.a = λb.b, and the β-redex (λx.g x x) a
(((λ (α) α)) ((g a a)))
\end{answers}
The first unification succeeds and binds nothing, so \scm{q} stays fresh. In the second, \scm{f} is
bound first and the application \scm{(@ f 'a)} is β-reduced when the deep walk of the answer meets it.
With call-by-name evaluation, the book's $(\lambda x\lambda w.w)\,\Omega$ \book{183} evaluates to
\lp{abs w\w} without ever evaluating Ω, because the substitution of the argument is a meta-level β that
drops it (\cref{ch:examples}).

\begin{definition}[A5: \texorpdfstring{$\Llam$}{Lλ} unification and delay]\label{def:design-a5}
Call \emph{pattern} for a variable $F$ an argument list $\bar a$ of $F\,\bar a$ whose members
head-normalise to distinct eigenvariables of level $>\ell(F)$ (\cref{def:hounif-llambda}, with levels in
place of the prefix).
\begin{enumerate}[label=(\alph*)]
\item \textbf{Patterns.} For an equation all of whose flexible subterms are patterns, unification
  terminates; it fails if and only if the equation has no unifier, and otherwise it binds a most general
  unifier (\cref{thm:hounif-mgu}). On the problems 8.G--8.K it gives the book's answers and mgus
  \book{222--225}. (The book does not number these problems; the labels 8.A--8.M are those of
  \cref{sec:hounif-traces}, in their order of appearance, as the test suite uses them; 8.M is the problem
  of \book{208}.)
\item \textbf{Dynamic $\Llam$} \book{225--226}. An equation outside $\Llam$ is neither failed nor
  enumerated: it is suspended, retried after every later binding, solved as soon as it becomes a
  pattern, and otherwise reported in the answer as a residual \scm{(assert (≐ u v))}, closed over the
  eigenvariables whose scope it outlived.
\end{enumerate}
\end{definition}

The level condition is Miller's \enquote{universally quantified inside the scope of $F$}. To see it
at work, compare $\exists F\,\forall x.\,F\,x=x$ with $\forall x\,\exists F.\,F\,x=x$. In the first,
$\ell(F)=0<1=\ell(x)$, so \scm{(@ F x)} is a pattern, and the only unifier is $F:=\abs{x}{x}$. In the second, $F$ is created
under the \scm{∀°}, so $\ell(F)=1=\ell(x)$ and \scm{(@ F x)} is not a pattern: $F$ may mention
$x$ as a constant, and both $\abs{z}{z}$ and $\abs{z}{x}$ solve the equation, so no unifier is most
general. The proposal solves the first; the second it suspends, accepting either solution when a
later goal supplies it:
\begin{scheme}
(pretty-print
 (list (°->list/datum (fresh° (q) (∀° (x) (=° (@ q x) x))))
       (°->list/datum (fresh° (q) (∀° (x) (fresh° (F) (=° (@ F x) x)
                                                    (=° F (ƛ (z) z)) (=° q 'id)))))
       (°->list/datum (fresh° (q) (∀° (x) (fresh° (F) (=° (@ F x) x)
                                                    (=° F (ƛ (z) x)) (=° q 'const)))))))
(pretty-print
 (°->pretty (fresh° (q) (∀° (x) (fresh° (F) (=° (@ F x) x) (=° q 'ok))))))
\end{scheme}
\begin{answers}
(((ƛ (x) x)) (id) (const))
((ok : (= (ƛ (x) (@ α x x)) (ƛ (x) x))))
\end{answers}
The last residual is instructive. When the \scm{∀°} exits, $F$ (of level 1) is \emph{raised},
$F:=\alpha\,x$ with $\ell(\alpha)=0$, and the equation is closed over $x$: $\forall x\,\exists F.\,F\,x=x$
has become $\exists\alpha\,\forall x.\,\alpha\,x\,x=x$, which is problem 8.M below.

Part (a) is Miller's theorem~\cite{Miller1991}, whose algorithm (\cref{def:hounif-elim}) prunes and raises; part (b) is
what Teyjus~2 does with the equations it cannot solve, after Dowek, Hardin, Kirchner and
Pfenning~\cite{DowekHardinKirchnerPfenning1996} (\cref{sec:hounif-dynamic}). The book's main
example, problem 8.H \book{223--224} (\cref{ex:hounif-8h}), $\forall w\,\forall x\,\forall y.\ f\,(U\,x\,y)=f\,(g\,(V\,y\,w))$
with $U$, $V$ quantified outside, is the λProlog query
\lp{pi w\ pi x\ pi y\ f (U x y) = f (g (V y w))}; it has the mgu
$U=\abs{x}{\abs{y}{g\,(V'\,y)}}$, $V=\abs{y}{\abs{w}{V'\,y}}$ (\cref{ex:hounif-8h}):
\begin{scheme}
(show "Example 8.H, pruning: (U V)" (°->list/datum (fresh° (q u v) (=° q (list u v)) (∀° (w x y) (=° `(f ,(@ u x y)) `(f (g ,(@ v y w))))))))
\end{scheme}
\begin{answers}
;; Example 8.H, pruning: (U V)
(((ƛ (x y) (g (@ α y))) (ƛ (y w) (@ α y))))
\end{answers}
The fresh \scm{α} is the book's $V'$. A non-pattern such as problem 8.M, \lp{pi x\ F x x = x}, has two
incomparable unifiers, so it waits; it is reported, and a later binding wakes it. Problem 8.A,
$F\,a=g\,a\,a$ with $a$ a constant, is not a pattern either (its argument is no eigenvariable):
\begin{scheme}
(show "non-patterns: 8.M and 8.A, suspended; 8.M woken by F := λa.λb.b"
      (list (°->pretty (fresh° (q) (∀° (x) (=° (@ q x x) x))))
            (°->pretty (fresh° (q) (=° (@ q 'a) '(g a a))))
            (°->list/datum (fresh° (q) (∀° (x) (=° (@ q x x) x))
                                       (=° q (ƛ (a b) b))))))
\end{scheme}
\begin{answers}
;; non-patterns: 8.M and 8.A, suspended; 8.M woken by F := λa.λb.b
(((α : (= (ƛ (x) (@ α x x)) (ƛ (x) x))))
 ((α : (= (@ α a) (g a a))))
 ((ƛ (a b) b)))
\end{answers}
(The first two answers, 8.M and 8.A, are printed by \scm{°->pretty} of the show module, the third
by \scm{°->list/datum}; the residual of 8.M is closed over \scm{x}, which has left its scope.)

\begin{definition}[A6: λProlog's semantics where Ma and Friedman differ]\label{def:design-a6}
On the equations by which Ma and Friedman~\cite{MaFriedman2021} motivate static closures (\cref{sec:lkanren-mf21}), the
answer is λProlog's:
(1)~$\abs{a}{X}=\abs{b}{b}$ and their conjunction containing it fail, since $\exists X\forall a.\,X=a$ is
false (\cref{ex:hounif-8j}); the paper answers $X=\mathtt{d0}$;
(2)~$\abs{a}{\abs{a}{X}}=\abs{c}{\abs{d}{c}}$ fails, with no false \enquote{true = false}, because shadowing
is lexical;
(3)~their Example~3, $\abs{a}{\abs{b}{X}}=\abs{c}{\abs{d}{Y}}$, gives $X=Y$;
(4)~an assumption is tried before the program: with a dynamic \scm{eq°}, their \texttt{assume} query gives
\scm{(orange apple)} where the paper prints \scm{(apple orange)};
(5)~their E12, $\forall A.\,A=F\,A$, gives $F=\abs{A}{A}$, as in the paper.
\end{definition}

\begin{scheme}
(show "λa.X = λb.b fails (the paper: X = d0); λa.λa.X = λc.λd.c fails; ∀A. A = F A gives F = λA.A"
      (list (°->list #f (fresh° (q) (=° (ƛ (a) q) (ƛ (b) b))))
            (°->list #f (fresh° (q) (=° (ƛ (a) (ƛ (a) q))
                                        (ƛ (c) (ƛ (d) c)))))
            (°->list/datum (fresh° (q) (∀° (a) (=° a (@ q a)))))))
\end{scheme}
\begin{answers}
;; λa.X = λb.b fails (the paper: X = d0); λa.λa.X = λc.λd.c fails; ∀A. A = F A gives F = λA.A
(() () ((ƛ (a) a)))
\end{answers}
\Cref{sec:lkanren-counter} discusses these equations; the point here is only that the proposal answers
them the way λProlog's quantifier reading dictates.

\begin{definition}[A7: the book's programs]\label{def:design-a7}
The book's programs over λ-tree syntax give the book's answers, in the book's order where the
interleaving search allows it: de Bruijn translation in both directions \book{193}; prenex normal forms,
the unique one and the five of a non-unique case \book{196--197}; \texttt{mappred} (Figs.~5.3--5.4);
palindromes over functional difference lists, the book's first nine answers \book{142}; Fischer's CPS
transformation, exactly the book's term \book{257}.
\end{definition}

These programs are the integration tests: each needs several of A1--A5 at once, and each has an answer
printed in the book to compare with. \Cref{ch:examples} walks through them.

\begin{definition}[A8: micro's stores on λ-terms]\label{def:design-a8}\leavevmode
\begin{enumerate}[label=(\alph*)]
\item \scm{≠°} decides modulo αβη and respects scope; its outcome does not depend on whether a binding
  comes before or after it. Under \scm{∀°} it reads an eigenvariable as \rulename{generic}'s new name by default
  (\lp{pi x\ not (q = x)} succeeds), and as any term of its universe when the parameter
  \scm{μkanren-λ-reading} is \scm{'extensional} (the reading of Moiseenko and of Jin et
  al.~\cite{Moiseenko2019,JinEtAl2021}).
\item \scm{absent°} looks for its datum in the parts of a λ-term, never in binder names, de Bruijn
  indices or the internals of an eigenvariable; on a flexible term it waits.
\item \scm{symbol°}, \scm{number°} and the user's types wait on a flexible term, and decide a rigid
  λ-term or an eigenvariable by the type's predicate.
\item No eigenvariable is ever printed free in an answer; residual constraints print as \scm{≐} and
  \scm{(∀ (x) ...)}; binders print renamed apart; ground answers evaluate back to live λ-terms.
\end{enumerate}
\end{definition}

λProlog has no disequality, \emph{absento} or types, so this item is not fidelity but coherence: micro's
stores must keep meaning what they mean once terms have binders. For instance:
\begin{scheme}
(show "≠ modulo αβη, flex-flex with an existential, a flexible absento, a type on a flexible term"
      (list (°->pretty (fresh° (q) (≠° q (ƛ (x) x))))
            (°->pretty (fresh° (q f g) (=° q (list f g))
                         (≠° (ƛ (x y) (@ f x)) (ƛ (x y) (@ g y)))))
            (°->pretty (fresh° (q F) (=° q F) (absent° 'c (@ F 'c))))
            (°->list/datum (fresh° (q F) (=° q F) (absent° 'c (@ F 'c))
                             (=° F (ƛ (z) 'b))))
            (°->list/datum (fresh° (q F) (=° q (@ F 'a)) (symbol° q)
                             (=° F (ƛ (z) z))))))
\end{scheme}
\begin{answers}
;; ≠ modulo αβη, flex-flex with an existential, a flexible absento, a type on a flexible term
(((α : (≠ α (ƛ (x) x))))
 (((α β) : (∀ (_3) (≠ (α β) ((ƛ (x) _3) (ƛ (y) _3))))))
 ((α : (absento c (@ α c))))
 ((ƛ (z) b))
 (a))
\end{answers}
Line by line. (1) $q\ne\abs{x}{x}$ is kept as a constraint, as micro keeps $q\ne 5$. (2) The trial
unification of $\abs{x}{\abs{y}{f\,x}}$ with $\abs{x}{\abs{y}{g\,y}}$ prunes both $f$ and $g$ to a
common constant function: the line says $\neg\exists H.\,f=\abs{x}{H}\wedge g=\abs{y}{H}$, printed in
the equivalent form $\forall H.\,(f\;g)\ne(\abs{x}{H}\;\abs{y}{H})$. The variable $H$ (\scm{_3}), which
the trial created, is quantified \emph{inside} the negation, which micro's prefixes cannot express.
(3) Whether \scm{c} is absent from $F\,\cst c$ depends on $F$, so the constraint waits on the whole
flexible term. (4) Once $F:=\abs{z}{\cst b}$, the argument is dropped and the constraint holds.
(5) \scm{(symbol° q)} waits on $F\,\cst a$ and holds once $F:=\abs{z}{z}$ makes it \scm{a}.
\Cref{ch:module} explains each of these answers in detail.

\begin{definition}[A9: conservativity]\label{def:design-a9}\leavevmode
\begin{enumerate}[label=(\alph*)]
\item micro's six suites (\texttt{microkanren}, \texttt{-aggregation}, \texttt{-untagged},
  \texttt{-show}, \texttt{-arith}, \texttt{-date}; 125 cases) pass with unchanged expectations, answer
  order included, against the edited micro, both with and without \scm{(aux kanren lambda)} loaded.
\item A first-order program pays at most a few per cent for the edit (measured in
  \cref{sec:design-cost}).
\end{enumerate}
\end{definition}

\Cref{tab:design-items} summarises the nine items and the evidence for each.

\begin{table}[ht]
\centering\footnotesize
\caption{The acceptance items and the hand-written cases that pin them (83 cases, 320 assertions in
\texttt{microkanren-lambda.scm}; 6 more cases in \texttt{microkanren-lambda-limits.scm} pin the known
limits).}
\label{tab:design-items}
\renewcommand{\arraystretch}{1.2}
\begin{tabularx}{\textwidth}{@{}l>{\raggedright\arraybackslash}X>{\raggedright\arraybackslash}p{2.3cm}r@{}}
\toprule
item & requirement & book & cases / assertions\\
\midrule
A1 & scope of ∀: \rulename{generic}, NJK's binding conditions on every binding & \book{82, 85} & 9 / 22\\
A2 & ⊃: \rulename{augment}, clause order, shared and renamed variables, closed static relations & \book{77--78, 86--88} & 8 / 21\\
A3 & ∀ and ⊃ together: type inference & \book{191} & 3 / 8\\
A4 & λ-tree syntax: α, β, η; object substitution by meta β & \book{177--199} & 4 / 9\\
A5 & $\Llam$ unification with mgus; dynamic $\Llam$ (delay) & \book{220--226} & 10 / 40\\
A6 & λProlog's answers where Ma--Friedman differ & --- & 5 / 13\\
A7 & the book's programs, the book's answers & \book{142, 193--197, 257} & 5 / 12\\
A8 & ≠, absento, types and reification on λ-terms & --- & 11 / 65\\
A9 & micro's suites unchanged; first-order cost & --- & 1 / 2, + 125 cases ×2\\
\midrule
R, R2, R3 & regressions of the three review rounds (\cref{sec:design-reviews}) & & 27 / 128\\
\bottomrule
\end{tabularx}
\end{table}

% ================================================================================================
\section{The design space}\label{sec:design-space}

\Cref{sec:lkanren-compare} put four candidate meanings of \scm{=°} on λ-terms side by side, without
choosing (\cref{tab:lkanren-options}). In one paragraph each:

\begin{description}
\item[A.] \textbf{Pattern unification with de Bruijn terms and levels}, after Miller, Nipkow, and
  Nadathur and Linnell~\cite{Miller1991,Nipkow1993,NadathurLinnell2005}; what Teyjus~2 and
  ELPI~\cite{DunchevEtAl2015} run. Decidable and unitary on $\Llam$; non-patterns are delayed.
\item[B.] \textbf{Nominal unification} (Urban--Pitts--Gabbay~\cite{UrbanPittsGabbay2004}; αKanren~\cite{ByrdFriedman2007}). Decidable, unitary with a freshness
  context. A different logic: names as data, no function variables, no β; noms are distinct constants,
  not eigenvariables, so \rulename{generic} still needs a separate mechanism; αKanren has no disequality.
\item[C.] \textbf{Static closures} (Ma--Friedman~\cite{MaFriedman2021}). Grafting semantics: $\abs{a}{X}=\abs{b}{b}$ succeeds,
  unlike λProlog; a workshop preview with no published algorithm.
\item[D.] \textbf{Huet's pre-unification, enumerated by the stream} (Huet~\cite{Huet1975}; Teyjus~1). Exact for full higher-order
  unification, but semi-decidable, with possibly infinitely many pre-unifiers.
\end{description}

\noindent The proposal retains A, for three reasons.

\begin{enumerate}
\item \textbf{λProlog's semantics.} The requirements are the book's. Option A is what the book's own
  implementations run, and with delayed non-patterns its behaviour on the book's examples is the book's
  (A5). Option C gives different answers on A6's equations by design; option B answers a different
  question.
\item \textbf{Decidability and most general unifiers.} Miller's theorem (\cref{thm:hounif-mgu}) makes
  unification on patterns terminating and unitary. This matches the shape of micro's \scm{=°}, which
  maps a state to \emph{at most one} state (\scm{✓°} or \scm{✗°}). Option D needs branching inside
  unification, so it cannot be an \scm{=°}; it can be added later as an explicit search over the
  suspended equations, which are its natural agenda (\cref{ch:roadmap}).
\item \textbf{Universe levels make \rulename{generic} cheap and exact.} The tagged interpreter of Nadathur,
  Jayaraman and Kwon is sound and complete, and needs one integer per variable and one comparison per
  binding (\cref{def:lkanren-ui}). The alternatives for ∀ alone are weaker or costlier: the explicit
  prefix of λKanren 2020~\cite{MaLuFriedman2020} is walked at each check and only grows (\cref{sec:lkanren-mk20}); the
  eigen-absento of Hemann, Byrd and Friedman~\cite{HemannByrdFriedman2013} costs $O(|B|)$ per check and is first-order only
  (\cref{sec:lkanren-eigen}); the end-of-scope check of \cref{sec:micro-forall} prunes late and lets
  constraints outlive their scope.
\end{enumerate}

Three secondary choices follow from A and are shared by every prototype below. Terms are
\emph{locally nameless}~\cite{McBrideMcKinna2004,Chargueraud2012} (de Bruijn indices inside bodies,
eigenvariables for free names), so α costs
nothing and terms built alike are \scm{equal?}. Raising is done lazily through levels rather than
eagerly as in the book's §4.4.1 \book{107}; the two differ in the shape of the mgu, not in its meaning
(\cref{ex:hounif-8k}). And non-patterns are delayed rather than rejected: rejecting them, as ELPI does
by default, would lose \texttt{subst R N (R N)} with \texttt{R} unknown, and the palindrome program.

What A leaves open is \emph{where} the code lives with respect to \micro, and that question was
answered by building it three times.

% ================================================================================================
\section{Three prototypes}\label{sec:design-prototypes}

Three complete prototypes were built against A1--A9. They share the representation (locally nameless
abstractions, applications, eigenvariables with levels) and the algorithmic core (variable elimination
with pruning and lowering, a store of suspended equations). They differ in architecture.
\Cref{tab:design-prototypes} compares them at a glance; the subsections describe each one, its strengths,
and the concrete defects the three reviews found.

\begin{table}[ht]
\centering\footnotesize
\caption{The three prototypes.}
\label{tab:design-prototypes}
\renewcommand{\arraystretch}{1.2}
\begin{tabularx}{\textwidth}{@{}p{2.3cm}>{\raggedright\arraybackslash}X>{\raggedright\arraybackslash}X>{\raggedright\arraybackslash}X@{}}
\toprule
& \textbf{d1, layered} & \textbf{d2, hooks} & \textbf{d3, native}\\
\midrule
edit to micro & none (byte-identical binary) & +105/−25 lines: a field and seven hooks & micro grows from 713 to 1217 lines; a 931-line diff\\
λ state & a record bound to the reserved variable 0 & a seventh field \texttt{X} & \texttt{X} holds U, E, P; levels on the variable record\\
unifier & a second one, \texttt{=λ°}, beside micro's & micro's, plus an \texttt{unify} hook & micro's, rewritten\\
runner & its own, \texttt{λrun} & micro's \texttt{°->list} & micro's \texttt{°->list}\\
suite & 29 cases, 137 assertions, hand-written & 42 cases, 102 assertions, hand-written & 14 cases, 116 assertions, a generated snapshot\\
rank & last for all three reviews & first for two reviews & first for the pedagogy review\\
\bottomrule
\end{tabularx}
\end{table}

\subsection{d1: a layer with no edit to micro}\label{sec:design-d1}

\paragraph{Architecture.} d1 takes \cref{sec:micro-noedit} at its word: it imports micro unchanged and
adds a pure library. Its whole context (current level, level boundaries and overrides, the clause
database, delayed equations, pending constraints) is a record bound to the reserved variable 0, the
\enquote{S-register} of Trick~2, overwritten by a persistent \scm{update/sbral} in $O(\log|S|)$, so
each branch of an \scm{or°} keeps its own context; the query variable is variable 1. Since micro's binding site is closed, d1 brings its own unifier \scm{=λ°}, its own
constraints \scm{≠λ°} and \scm{absentλ°}, and its own runner \scm{λrun}. Levels need no slot in the
variables: micro allocates indices monotonically, so a variable's level at creation is a step function
of its index, recorded as boundaries pushed by \scm{∀°}; lowering is recorded as overrides.

\paragraph{Strengths.} Micro's behaviour and speed are untouched by construction. d1's \scm{absent°}
was the most precise on flexible terms, it offered the extensional reading of ≠ under ∀, it had
an αKanren-style printer, a data runner comparable with \scm{equal?}, and a table of every place where
purity leaks, each pinned by a test. All these were grafted into the final proposal.

\paragraph{Defects.} The leaks are the architecture's, not the code's. Micro's \scm{=°}, and every
relation built on it, binds without d1 seeing the binding, so A1(c) fails:
\begin{scheme}
; d1-layered, its own runner
(pretty-print (λrun (q -1) (∀° (x) (cons° x '() q))))
(pretty-print (λrun (q -1) (∀° (x) (=λ° q x))))
\end{scheme}
\begin{answers}
((x₂))
()
\end{answers}
With d1's own unifier the escape is caught; with micro's \scm{cons°} the eigenvariable \scm{x₂} leaks
into the answer. The proposal answers \scm{()} to the first query (A1, above). The reviews found four
more defects: the module's goals do not run under \scm{°->list} (they need the register);
conjunction is not commutative, because d1 read a list headed by a symbol as a rigid application in
spine notation, conflating lists with λ-tree applications, and together with η this made the answers of
some conjunctions depend on the order of the conjuncts; d1 duplicates a large part of micro; and \scm{=λ°} costs
4.4~times micro's \scm{=°} on search. Its clause database was a linear list, on which the reviews
measured slowdowns of 18~times. d1 was ranked last by all three reviews.

\subsection{d2: an extension protocol in micro}\label{sec:design-d2}

\paragraph{Architecture.} d2 makes the smallest principled edit that opens the closed points of
\cref{tab:micro-open}, and builds the whole λ layer as a client module on top of it. The state gets a
seventh field \texttt{X}, which micro copies and never reads. A record of procedures, registered
globally, has hooks for \texttt{term?}, \texttt{unify}, \texttt{post}, \texttt{diseq}, \texttt{value},
\texttt{repr} and \texttt{forms}; micro calls them at the matching points. The edit was +105/−25 lines
(counted on its \texttt{micro.diff}), the module 708 lines.

\paragraph{Strengths.} The reviews found no soundness defect in d2. It kept terms η-short as an
invariant, so both $\abs{x}{q\,x}=q$ and $F\,x=G\,x\Rightarrow F=G$ hold; it had a \enquote{soft}
occurs check that suspends $F\,x=G\,(F\,x)$ instead of failing it (it holds with $G:=\abs{z}{z}$); its
suite was hand-written with book page references; its printer was the only one that renamed binders
correctly; and the edit cost at most 2.8\% on micro's benchmarks.

\paragraph{Defects.} The reviews found seven, listed in \cref{tab:design-d2}. None is a soundness
defect; most are imprecisions of the stores on λ-terms and one is a cost.

\begin{table}[ht]
\centering\footnotesize
\caption{The defects of d2 found by the reviews, and their fixes (design notes, §10.2).}
\label{tab:design-d2}
\renewcommand{\arraystretch}{1.2}
\begin{tabularx}{\textwidth}{@{}l>{\raggedright\arraybackslash}X>{\raggedright\arraybackslash}X@{}}
\toprule
& defect & fix in the proposal\\
\midrule
D1 & a suspended disequality printed as \texttt{(deny (and))} & the denied equation, \texttt{(deny (≐ u v))}\\
D2 & \texttt{absent°} over-strict on flexible terms & the \texttt{parts} hook: a flexible term is watched whole\\
D3 & trivially true residuals kept ($F\,x\,x=F\,x\,x$) & delay drops α-equal sides\\
D4 & Miller's n7 left suspended and printed with \texttt{\_2} & discharge at ∀ exit; the \texttt{terms} hook names residual variables\\
D5 & registering the extension slowed ≠°-heavy first-order code by 45--50\% & migration only when a new binding holds a λ-term\\
D6 & goal suffix on term constructors, levels hidden in property lists, λ-terms not \texttt{equal?} & \texttt{ƛ}, \texttt{@}; plain records thanks to \texttt{parts}, \texttt{value}, \texttt{atom?}; interned hints\\
D7 & \texttt{symbol°}/\texttt{number°} failed eagerly on flexible terms; \texttt{⊃°} on a static relation silently ignored & the \texttt{tag} hook; an error for non-dynamic relations\\
\bottomrule
\end{tabularx}
\end{table}

One more defect turned up while grafting: $F=\abs{x}{G\,(F\,x)}$, reached through η-short construction,
failed micro's occurs check before any hook could see it. That is why the final protocol has an
\texttt{occurs} hook.

\subsection{d3: a native engine}\label{sec:design-d3}

\paragraph{Architecture.} d3 integrates λ-terms into micro itself, in the style of ELPI and Teyjus. The
variable record gets a level slot; \scm{μkanren-unify} knows abstractions, eigenvariables and flexible
applications; the binding site does the occurs check, the scope check and lowering in one walk; \scm{∀°}
and \scm{⊃°} are core goals; the current universe, the delayed equations and the program are packed in
one extra state field. Micro grew from 713 to 1217 lines.

\paragraph{Strengths.} The pedagogy review ranked d3 first for three reasons: one vocabulary, readable
answers in \scm{ƛ}, \scm{@} and \scm{≐}, and answers usable as code. d3 also discharged residual
constraints at ∀ exit by raising and closing them, and it was the fastest: type inference of the Church
numeral 400 took 105~ms, against d2's 441~ms.

\paragraph{Defects.} The reviews found a crash in \scm{μkanren-verify-D+} on a disequality prefix whose
trial unification delays, and an η-incompleteness that made \scm{≠°} unsound: after
$\abs{x}{F\,x}\ne\abs{x}{G\,x}$, the binding $G:=F$ was accepted.
\begin{scheme}
(pretty-print (°->list #f (fresh° (q F G) (=° q (list F G)) (≠° (ƛ (x) (@ F x)) (ƛ (x) (@ G x))) (=° G F))))
\end{scheme}
Run against d3, then against the proposal:
\begin{answers}
((λ (α) (cons α (cons α '()))))
\end{answers}
\begin{answers}
()
\end{answers}
The answer of d3 is a pair $(F,F)$, which violates the disequality it was asked to keep. Besides, its
suite was a snapshot generated from the actual answers rather than written from the book, and its diff
carried trailing-whitespace churn.

\subsection{The verdict}\label{sec:design-verdict}

The three reviews gave aggregate scores of \textbf{23} to d2, \textbf{21} to d3 and \textbf{15.5} to
d1. Two reviews ranked d2 first, the pedagogy review ranked d3 first, and all three ranked d1 last. The
choice of d2 as the base rests on four facts.

\begin{itemize}
\item \textbf{Size and scope of the edit.} The final protocol is +235/−41 lines of micro (713 to 907
  lines; 214 of the added lines are not blank), comments included, all generic: nothing in micro mentions λ, levels or clauses. d3 special-cases every walker.
  The same protocol could carry another store, such as nominal terms or finite domains.
\item \textbf{Soundness} is equal. d3's eager binding rule and the λ layer's eager \texttt{post} check
  enforce the same binding conditions; every scope attack of the reviews fails correctly on both, and on
  the final version.
\item \textbf{Correctness} favours d2: d3 had a crash and an unsound \scm{≠°}, d2 neither.
\item \textbf{Pedagogy.} d3's three advantages were grafted onto d2 (one vocabulary of \scm{ƛ}, \scm{@},
  \scm{≐}; readable residuals; ground answers as live terms through the \texttt{environment} hook), so
  the pedagogical argument for the native design no longer holds.
\end{itemize}

The reviews also recommended about twenty grafts from d1 and d3. Each was either done or explicitly
declined; \cref{tab:design-grafts} lists the main ones.

\begin{table}[ht]
\centering\footnotesize
\caption{Grafts recommended by the reviews (a selection from the design notes, §10.3).}
\label{tab:design-grafts}
\renewcommand{\arraystretch}{1.2}
\begin{tabularx}{\textwidth}{@{}l>{\raggedright\arraybackslash}X>{\raggedright\arraybackslash}p{4.6cm}@{}}
\toprule
from & graft & status\\
\midrule
d3 & discharge residual constraints at ∀ exit by raising and closing & done\\
d3 & \texttt{≐} exported, residuals as \texttt{(assert (≐ …))}, merged binders \texttt{(ƛ (x y) …)} & done\\
d3 & a first-order fast path for the scope check & done (\texttt{μkanren-plain?}, \texttt{top}, the \texttt{atom?} hook)\\
d3 & ground answers as live terms & done (\texttt{environment} hook)\\
d3 & the 2--4× faster λ machinery & partly: Church 400 from 441 to 254 ms; d3: 105 ms\\
d3 & var--var bindings oriented by level & declined: a hook on micro's second-hottest path\\
d3 & levels on the variable record & declined: \texttt{freshª} would need the universe, so micro would become λ-aware\\
d1 & precise \texttt{absent°} on flexible terms & done\\
d1 & the extensional reading of ≠ under ∀, as an option & done (\texttt{μkanren-λ-reading})\\
d1 & a data runner, an αKanren-style view & done (\texttt{°->list/datum}, \texttt{°->pretty})\\
d1 & a limits table, each row pinned by a case & done\\
d1 & truly nameless λ-records, \texttt{equal?} modulo α & partly: interned hints; α-variants with different hints are not \texttt{equal?}\\
\bottomrule
\end{tabularx}
\end{table}

The two declined grafts share a reason. Both would put knowledge of levels into micro's core: the
var--var clause of the unifier, its hottest path after pair--pair, and the variable constructor
\scm{freshª}, which would need the current universe kept in the extension's data. The protocol's
discipline is that micro stays λ-agnostic; for the orientation, the lowering done by the scope check
yields the same mgu up to renaming.

% ================================================================================================
\section{The extension protocol}\label{sec:design-protocol}

The protocol is the part of the proposal that edits \micro. It is generic: it lets a client module own
a class of records (its \emph{terms}), keep per-branch data in the state, and be consulted at each point
where micro would otherwise treat those records as plain slot vectors. This section gives it in full.
Excerpts are verbatim from \texttt{micro.diff}; lines starting with \texttt{+} are added, with
\texttt{-} removed, and with a space unchanged context. Comment lines of the diff are left out of some
excerpts, as is the removed body of \scm{=°}.

\subsection{The seventh field}\label{sec:design-X}

The state gets a field \texttt{X}, the extensions' slot. Micro never reads it: it copies it at its two
positional constructions, the binding site \scm{μkanren-state-update} and \scm{freshª}, and
\scm{μkanren-state-with} accepts \texttt{X:}. The positional constructor keeps its six-argument
signature, with \texttt{X} optional, so the positional states of \texttt{test/microkanren.scm} still
work.
\begin{scheme}
-  (define-record μkanren-state vars-count S D A T tags)
+  (define-record-type μkanren-state
+    (make-μkanren-state/X vars-count S D A T tags X)
+    μkanren-state?
+    (vars-count μkanren-state-vars-count) (S μkanren-state-S) (D μkanren-state-D) (A μkanren-state-A)
+    (T μkanren-state-T) (tags μkanren-state-tags) (X μkanren-state-X))
+
+  (define (make-μkanren-state vars-count S D A T tags #!optional (X '())) (make-μkanren-state/X vars-count S D A T tags X))
\end{scheme}
and in \scm{freshª}:
\begin{scheme}
-             (s*  (make-μkanren-state (add1 vc) S D A T tags)))
+             (s*  (make-μkanren-state/X (add1 vc) S D A T tags (μkanren-state-X s))))
\end{scheme}
\texttt{X} is \scm{'()} until an extension writes it, which the λ layer reads as \enquote{everything
empty}; so \scm{°->§} keeps its fixed empty initial state. The λ layer stores in \texttt{X} one record
with the current depth, the highest depth reached, the variables' levels, the context $\Gamma$, the
suspended constraints E and the λ-disequalities N (\cref{ch:module}). Since S is persistent and
\texttt{X} is copied forward by every goal, the data is branch-local for free: each branch of an
\scm{or°} sees its own $\Gamma$, as A2(a) requires.

\subsection{Extensions and the registry}\label{sec:design-registry}

An extension is a record built by
\begin{scheme}
  (define (make-μkanren-extension name #!key (term? #f) (atom? #f) (unify #f) (occurs #f) (post #f) (diseq #f) (tag #f) (parts #f)
                                             (value #f) (repr #f) (terms #f) (forms #f) (environment #f))
\end{scheme}
and \scm{μkanren-extension-register!} adds it to a global list. Registering again under the same name
replaces the old extension, so reloading a module is harmless. Dispatch is by term: the extension that
owns a pair of terms is the first whose \texttt{term?} holds of either.
\begin{scheme}
+  (define (μkanren-extension/term u v)
+    (let loop ((E μkanren-extensions))
+      (cond
+        ((null? E) #f)
+        ((let1 (term? (μkanren-extension-term? (car E))) (and term? (or (term? u) (term? v)))) (car E))
+        (else (loop (cdr E))))))
\end{scheme}
The call sites are guarded by \scm{(pair? μkanren-extensions)} (those of \texttt{forms} and
\texttt{environment}, in reification and in the runner, just map over the empty list), so with no
extension registered a hook costs one \scm{pair?} test, and none sits on micro's hottest path (the variable clauses and the
pair--pair clause of the unifier). The registry is global rather than per state because \scm{°->list}
starts from \scm{μkanren-state-empty}: a per-state hook list would have to be carried by every initial
state.

The protocol's whole exported surface is small: \scm{μkanren-state-X} (and the \texttt{X:} keyword of
\scm{μkanren-state-with}), \scm{make-μkanren-extension}, \scm{μkanren-extension?},
\scm{μkanren-extension-name}, \scm{μkanren-extension-register!}, and the two primitives
\scm{μkanren-post-unify} and \scm{μkanren-state-forget} of \cref{sec:design-forget}. The hook
accessors and the dispatchers (\scm{μkanren-extension/term} and its siblings) stay internal to micro.

The λ layer registers itself once, at load time:
\begin{scheme}
  (define μkanren-λ-extension
    (make-μkanren-extension 'λProlog
      term?: μkanren-λ-term? atom?: (λ (u) (or (μkanren-eigen? u) (μkanren-bvar? u))) unify: μkanren-λ-unify occurs: μkanren-λ-occurs post: μkanren-λ-post diseq: μkanren-λ-diseq
      tag: μkanren-λ-tag parts: μkanren-λ-parts value: μkanren-λ-value repr: μkanren-λ-repr
      terms: μkanren-λ-terms forms: μkanren-λ-forms environment: μkanren-λ-environment))

  (μkanren-extension-register! μkanren-λ-extension)
\end{scheme}

\subsection{The hooks}\label{sec:design-hooks}

\Cref{tab:design-hooks} lists the thirteen hooks. Compared with d2's seven, the final protocol adds
\texttt{atom?}, \texttt{occurs}, \texttt{tag}, \texttt{parts}, \texttt{terms} and \texttt{environment},
each answering a defect or a graft of the reviews (\texttt{atom?} the first-order fast path,
\texttt{environment} live ground answers, the others the defects of \cref{tab:design-d2} and the
occurs failure found while grafting). The last column says which closed point of
\cref{tab:micro-open} each hook opens. Then a paragraph per hook gives its contract and its call site.
\newpage
\begin{longtable}{@{}>{\raggedright\arraybackslash}p{2.0cm}>{\raggedright\arraybackslash}p{5.5cm}>{\raggedright\arraybackslash}p{3.3cm}>{\raggedright\arraybackslash}p{2.5cm}@{}}
\caption{The hooks of an extension. $u,v$: walked terms; $s$: a state; \emph{as}: the bindings made so
far, as \texttt{(α . t)} pairs; $\to$ gives the result, \textbar\ separates alternatives.}\label{tab:design-hooks}\\
\toprule
hook & arguments $\to$ result; contract & called by & opens\\
\midrule\endfirsthead
\toprule
hook & arguments $\to$ result; contract & called by & opens\\
\midrule\endhead
\texttt{term?} & \texttt{(u)}: $u$ is one of the extension's records & dispatch & ---\\
\texttt{atom?} & \texttt{(u)}: $u$ holds no variable & micro's occurs check, which then does not look inside $u$ & the binding site\\
\texttt{unify} & \texttt{(u v s as)} $\to$ \texttt{(s* as*)} \textbar\ \texttt{\#f}: a pair with one extension term, neither a variable & \texttt{μkanren-state-unify}: the structure--structure clause and the final \texttt{else} & the unification \texttt{cond}\\
\texttt{occurs} & \texttt{(α v s as)} $\to$ \texttt{(s* as*)} \textbar\ \texttt{\#f}: micro's occurs check failed on $α:=v$ & \texttt{μkanren-state-update}, failure path only & the binding site\\
\texttt{post} & \texttt{(s* as)} $\to$ \texttt{(s** as**)} \textbar\ \texttt{\#f}: runs after the bindings of \texttt{=°}, before D/A/T are verified; may bind more & \texttt{μkanren-post-unify}, called by \texttt{=°} & the post-\texttt{=°} chain\\
\texttt{diseq} & \texttt{(u v s)} $\to$ \texttt{\#f} \textbar\ state \textbar\ \texttt{'fail}: may claim a disequality & \texttt{≠°} & \texttt{≠°} (store D)\\
\texttt{tag} & \texttt{(u tag s)} $\to$ state \textbar\ \texttt{\#f}: a type on its term & \texttt{symbol°}, \texttt{number°}, user types & the A store\\
\texttt{parts} & \texttt{(u s [projecting])} $\to$ \texttt{'stuck} \textbar\ list: where \texttt{absent°} looks into its term; when projecting, the parts whose constraints an answer keeps & \texttt{absent°}, the re-verification of T, \texttt{μkanren-project-T} & the T store; reification\\
\texttt{value}, \texttt{repr} & \texttt{(u walk s)}: the deep value; the expression that rebuilds it & \texttt{μkanren-state-find/value}, \texttt{/repr} & reification\\
\texttt{terms} & \texttt{(s)} $\to$ terms: the terms of its residual constraints & \texttt{μkanren-project} & reification\\
\texttt{forms} & \texttt{(s s*)} $\to$ forms: the forms its constraints add to an answer & \texttt{μkanren-project} & reification\\
\texttt{environment} & \texttt{()}: where \texttt{°->list/ground} evaluates answers & \texttt{°->list} & the runner\\
\bottomrule
\end{longtable}

\paragraph{\texttt{term?}} identifies the extension's records. Every dispatch by term goes through it
(\texttt{unify}, \texttt{tag}, \texttt{parts}, \texttt{value}, \texttt{repr}); the other hooks are
consulted on every registered extension, so an extension without \texttt{term?} could still use them.
For the λ layer, the four records are abstractions,
applications, eigenvariables and de Bruijn indices.

\paragraph{\texttt{unify}} is consulted when micro's unifier meets a pair of walked terms, one of which
is the extension's and neither of which is a variable. Two clauses of the unification \scm{cond} hand
such pairs over: the structure--structure clause (two records, or a record and a vector) and the final
\scm{else} (an extension term against a pair or an atom).
\begin{scheme}
         ((and (μkanren-structure? u*) (μkanren-structure? v*))
-          (μkanren-state-unify (μkanren-structure->list u*) (μkanren-structure->list v*) s assocs))
+          (cond
+            ((and (pair? μkanren-extensions) (μkanren-extension/term u* v*)) => (λ (e) ((μkanren-extension-unify e) u* v* s assocs)))
+            (else (μkanren-state-unify (μkanren-structure->list u*) (μkanren-structure->list v*) s assocs))))
         ; the hottest path: a unification answers `(s* assocs*)` or #f, no need for a checked match.
         ((and (pair? u*) (pair? v*))  (let1 (r (μkanren-state-unify (car u*) (car v*) s assocs))
                                         (and r (μkanren-state-unify (cdr u*) (cdr v*) (car r) (cadr r)))))
+        ; an extension term against a pair or an atom.
+        ((and (pair? μkanren-extensions) (μkanren-extension/term u* v*)) => (λ (e) ((μkanren-extension-unify e) u* v* s assocs)))
         (else #f))))
\end{scheme}
The variable clauses, above this excerpt, and the pair--pair clause are untouched. The hook may call
\scm{μkanren-state-unify} back on parts: the λ layer's ξ rule opens both bodies with one new name and
unifies them with micro's own unifier. So a list nested in a λ-term is still unified by micro, and a
λ-term nested in a list is still handed to the hook. Before the edit, two records unified slot by
slot, so an abstraction would have been compared by its hint and its raw body.

\paragraph{\texttt{atom?}} tells micro's occurs check that a record holds no variable, so the check need
not look into its slots. For the λ layer, eigenvariables and indices are atoms. Without it micro rebuilt
every eigenvariable as a slot list at each binding; adding it cut type inference of the Church numeral
400 from 311 to 254~ms (an 18\% gain on the λ workloads).
\begin{scheme}
-                ((μkanren-structure? w) (O (μkanren-structure->list w)))
+                ((μkanren-structure? w) (and (not (μkanren-extension/atom? w)) (O (μkanren-structure->list w))))
\end{scheme}

\paragraph{\texttt{occurs}} gets the last word when micro's occurs check fails on $α:=v$ and $v$ may hold
extension terms. Inside a λ-term an occurrence may be only apparent: $q = F\,q$ holds with
$F:=\abs{z}{z}$, and $F=\abs{x}{G\,(F\,x)}$ with $G:=\abs{z}{z}$. The binding site, closed in
\cref{tab:micro-open}, is opened on its failure path only:
\begin{scheme}
+    (cond
+      ((or unsafe (not (occur? v)))
+        (μkanren-state-match ((vc S D A T tags) s)
+          (let* ((S* (μkanren-update/sbral α v S))
+                 (s* (make-μkanren-state/X vc S* D A T tags (μkanren-state-X s)))
+                 (assocs* (cons `(,α . ,v) assocs)))
+            `(,s* ,assocs*))))
+      ; an occurrence inside an extension's term may be only apparent: the extension decides.
+      ((pair? μkanren-extensions) (μkanren-extension/occurs α v s assocs))
+      (else #f)))
\end{scheme}
The λ layer answers by β-normalising $v$ (if it is $α$ itself the equation is trivial) and otherwise by
variable elimination, which fails on a rigid occurrence and suspends one in the argument of a flexible
term (\cref{ch:module}).

\paragraph{\texttt{post}} is the hook that carries the scope check. It runs after the bindings that one
\scm{=°} made and \emph{before} D, A and T are verified again; it receives those bindings, may bind more
(lowering, pruning, waking suspended equations), and returns the extended bindings, so that micro
re-verifies the stores they touch. The post-\scm{=°} chain, closed and wired for D/A/T before the edit,
becomes an exported procedure:
\begin{scheme}
+  (define (μkanren-post-unify s* assocs)
+    (let loop ((E μkanren-extensions) (s* s*) (assocs assocs))
+      (cond
+        ((null? E) (cond
+                     ((μkanren-verify-D/touched (μkanren-state-D s*) s* assocs)
+                       => (μ D* (μkanren-verify-D/post D* (μkanren-state-A s*) (μkanren-state-T s*) s*)))
+                     (else #f)))
+        ((μkanren-extension-post (car E)) => (λ (post)
+                                                (match/first (post s* assocs)
+                                                  (#f #f)
+                                                  ((,s** ,assocs*) (loop (cdr E) s** assocs*)))))
+        (else (loop (cdr E) s* assocs)))))
\end{scheme}
and \scm{=°} calls it:
\begin{scheme}
   (define ((=° u v) s)
+    (define post°
+      (λ1-match/first
+        ((,s* ,assocs)  (cond
+                          ((μkanren-state-equal? s s*) (✓° s*))
+                          ((μkanren-post-unify s* assocs) => ✓°)
+                          (else (✗° s*))))))
+    (cond
+      ((μkanren-state-unify u v s '()) => post°)
+      (else (✗° s))))
\end{scheme}
Read \scm{post°} clause by clause. If raw unification returned the very state it was given
(\scm{μkanren-state-equal?} is the \scm{eq?} test of \cref{prop:micro-eq}), nothing was bound, so no
hook runs and no store is re-verified; this is why invariant (I1) of \cref{sec:design-invariants}
matters. Otherwise the bindings \scm{assocs} go through every extension's \texttt{post}, in
registration order, each one receiving the state and bindings the previous one answered; then D is
re-verified on the touched prefixes and A and T after it. A \scm{#f} anywhere fails the \scm{=°}.

This is how A1(c) is met without opening the binding site: \emph{every} binding reaches the store
through some \scm{=°}, micro's included, and every \scm{=°} ends in \scm{μkanren-post-unify}. The binding
is checked a moment after it is made, but before any other goal can see it. Running \texttt{post}
before the D/A/T verification, not after, matters for two reasons: its own bindings must reach
\scm{μkanren-verify-D/touched}, and the λ layer moves a disequality prefix that has come to hold a
λ-term out of D \emph{before} micro re-unifies it.

\begin{pitfall}[A goal that unifies on its own must end in \texttt{μkanren-post-unify}]
The guarantee \enquote{every binding is checked} holds because every goal of micro that commits a
binding to the search state is an \scm{=°}; micro's other calls of raw unification are trials whose
states are thrown away (in \scm{≠°}, in \scm{absent°}, in subsumption) or reification. A goal of an extension that calls the raw \scm{μkanren-state-unify} and
returns its state bypasses \texttt{post} exactly as Trick~1's \scm{bind/raw} bypasses D/A/T
(\cref{sec:micro-noedit}): an eigenvariable could reach an older variable unchecked. Such a goal must
pass the state and the exact new bindings to \scm{μkanren-post-unify}, which is why it is exported; the
λ layer does so at \scm{∀°} exit (\cref{sec:design-forget}).
\end{pitfall}

\paragraph{\texttt{diseq}} lets an extension claim a disequality; micro's original body becomes
\scm{μkanren-≠}.
\begin{scheme}
+  ; an extension may take the disequality over, when its terms are involved.
   (define ((≠° u v) s)
+    (cond
+      ((and (pair? μkanren-extensions) (μkanren-extension/diseq u v s)) => (λ (r) (if (eq? r 'fail) (✗° s) (✓° r))))
+      (else (μkanren-≠ u v s))))
+
+  (define (μkanren-≠ u v s)
\end{scheme}
\texttt{\#f} leaves $u\ne v$ to the store D, a state keeps it in the extension's store, \scm{'fail}
says that $u$ and $v$ are already equal. A separate hook is needed because a D prefix is a list of
bindings: the mgu of a pattern trial can contain variables the trial itself created, which are
existential inside the negation (the \scm{∃} of item A8), and a prefix cannot express that.

\paragraph{\texttt{tag}} decides a type goal (\scm{symbol°}, \scm{number°}, a type made by
\scm{define-μkanren-tag}) on an extension term. Before the edit such a goal failed on every record that
is not of the type, so \scm{(symbol° (@ F 'a))} failed even though $F:=\abs{z}{z}$ would make it a
symbol; now the λ layer makes it wait.
\begin{scheme}
+      ((,u* ⊣ (μkanren-extension/hook u* μkanren-extension-tag)) (cond
+                                                                  (((μkanren-extension/hook u* μkanren-extension-tag) u* tag s) => ✓°)
+                                                                  (else (✗° s))))
\end{scheme}

\paragraph{\texttt{parts}} says where \scm{absent°} looks into an extension term: a list of parts, or
\scm{'stuck} when the term cannot be looked into yet, in which case \scm{absent°} watches it whole, like a
variable. Micro walks record slots, so without the hook \scm{(absent° 0 q)} would reject λx.x because
of the index 0 in its body, as it did in d1. With \texttt{projecting} true, the same hook tells the projection of T on an answer which
parts' constraints the answer keeps, instead of splitting a λ-term into its raw slots.
\begin{scheme}
+  (define (μkanren-absent-parts u s)
+    (cond
+      ((μkanren-extension/hook u μkanren-extension-parts) => (λ (parts) (parts u s)))
+      (else (μkanren-compound->list u))))
\end{scheme}
Its call sites are \scm{absent°}, \scm{μkanren-verify-T+} (both read \scm{'stuck} as \enquote{watch}) and
\scm{μkanren-project-T}, through the twin \scm{μkanren-project-parts}.

\paragraph{\texttt{value} and \texttt{repr}} give the deep value of an extension term and the Scheme
expression that rebuilds it; \texttt{walk} is passed so that the hook can recurse on the parts. They are
called by \scm{μkanren-state-find/value} and \scm{/repr}, so β-normalisation applies wherever micro
deep-walks: answers, the stores in answers, \scm{project°} and aggregation.
\begin{scheme}
+            ((and (record-instance? w*) (not (μkanren-var? w*)))
+              (cond
+                ((and (pair? μkanren-extensions) (μkanren-extension/term w* w*)) => (λ (e) ((μkanren-extension-value e) w* A s)))
+                (else (apply make-record-instance (map A (μkanren-structure->list w*))))))
\end{scheme}

\paragraph{\texttt{terms} and \texttt{forms}} give the extension a say in the answer. \texttt{terms}
returns the terms of its residual constraints whose variables the answer should name: micro reifies them
after the answer, so a residual prints \scm{α}, \scm{β} instead of \scm{_2}. \texttt{forms} returns the
forms its constraints add to the answer's body, after micro's own.
\begin{scheme}
+    (let* ((w* (μkanren-state-find/value w s))
+           (extra (μkanren-extension/terms s)))
+      ; the variables of the extensions' residual constraints are named too, after those of the answer.
+      (match1/first ((,s* ,c ,vars-reversed ,assocs) (μkanren-state-reify (if (null? extra) w* (cons w* extra)) s))
\end{scheme}
\begin{scheme}
+               (X**   (μkanren-extension/forms s s*))
+               (body  (μkanren-form w** D** A** T** X** s*)))
\end{scheme}

\paragraph{\texttt{environment}} names the environment in which \scm{°->list/ground} evaluates answers.
The λ layer gives its own module environment, where \scm{ƛ}, \scm{@} and \scm{≐} are bound, so a ground
answer evaluates back to a live λ-term that can be applied again (the pedagogy graft from d3).
\begin{scheme}
-                         (let ((E (eval `(lambda ,(cadr expr) ,(last expr)) (μkanren-answer-environment)))
+                         (let ((E (eval `(lambda ,(cadr expr) ,(last expr)) (μkanren-extension/environment)))
\end{scheme}

\subsection{Two primitives: forget and post-unify}\label{sec:design-forget}

The protocol exports two procedures besides the registry.

\scm{(μkanren-state-forget α s)} unbinds $α$:
\begin{scheme}
+  (define (μkanren-state-forget α s)
+    (let1 (S (μkanren-state-S s))
+      (μkanren-state-with s S: (update/sbral (μkanren-var-index/sbral S α) *μkanren-unbound* S))))
\end{scheme}
A \texttt{post} hook uses it to take back a binding that \scm{=°} has just made, turning it into a
constraint of its own, and leaves that binding out of the bindings it answers. Since \texttt{post} runs
before D, A and T are verified, those stores never see the binding. The λ layer needs this when the scope
check of $X:=t$ cannot be decided yet, because $t$ contains a non-pattern that only a later binding can
resolve: the binding becomes the suspended equation $X\doteq t$, as Teyjus would delay it. For instance
$\exists q\,\exists G\,\exists H\,\forall y.\ q = G\,(H\,y)$ leaves \scm{q} unbound with a residual; a later
$G:=\abs{a}{\cst k}$ solves it with $q=\cst k$ (\cref{ch:module}).

\scm{μkanren-post-unify}, shown above, is exported for goals of an extension that unify on their own
and must then run the same chain as \scm{=°}. The λ layer calls it at ∀ exit, where the raising of the
variables of a residual constraint makes new bindings that must wake the suspended equations and be
checked against D, A and T.

\subsection{One generic change to absent°}\label{sec:design-absent}

The third review found a defect of micro itself, independent of λ-terms. Micro decided an atom against
an \emph{absento} datum with \scm{equal?} only, although a datum that holds a variable may still become
that atom. With the unedited micro:
\begin{scheme}
(pretty-print (°->list #f (fresh° (q D) (absent° D q) (=° q 'a) (=° D 'a))))
(pretty-print (°->list #f (fresh° (q D) (=° q 'a) (=° D 'a) (absent° D q))))
(pretty-print (°->list #f (fresh° (q D) (absent° D q) (=° q 'a))))
\end{scheme}
\begin{answers}
((λ () 'a))
((λ () 'a))
((λ () 'a))
\end{answers}
The first two answers claim that \scm{a} is absent from \scm{a}. With the edited micro, which decides
such an atom by \scm{μkanren-tag-unify} (as it already did for a compound) and watches it while it is
pending:
\begin{answers}
()
()
((λ () 'a))
\end{answers}
The third answer stays, correctly: \scm{D} is not in the answer and can be chosen different from
\scm{a}. Two smaller generic changes come with this one. An answer keeps a pending atom constraint when
its datum names a variable of the answer, and prints that datum as a term,
\scm{(absent? (@ α 'b) 'a)}. And when a type turns an \emph{absento} into a disequality prefix of D (micro's
subsumption), a datum that holds an extension term stays in T instead (the predicate
\scm{μkanren-extension/holds-term?}), so D never holds one:
\begin{scheme}
+      ((((,α* . ,tag) . ,T*) ⊣ (and (equal? α α*) (not (μkanren-extension/holds-term? (μkanren-tag-name tag)))))
+        (let1 (D* (μkanren-ext-D α tag D s)) (μkanren-update-D/T+ α T+ D* T* s)))
\end{scheme}

\subsection{The two inherited invariants}\label{sec:design-invariants}

A client of the protocol must respect two invariants that micro already relies on.

\begin{definition}[Protocol invariants]\label{def:design-invariants}\leavevmode
\begin{enumerate}[label=(I\arabic*)]
\item \textbf{eq?-entailment.} A hook that changes nothing returns the very state object it was given.
\item \textbf{Extended bindings.} A hook that binds returns the list of bindings it received extended
  with every binding it made, pruning and lowering included.
\end{enumerate}
\end{definition}

(I1) is \cref{prop:micro-eq} carried over to the hooks. Micro reads \enquote{unification gave back $s$}
as \enquote{the terms are already equal}: \scm{≠°} then fails, \scm{=°} skips re-verification, and the
re-verification of a D prefix declares it violated. A hook that rebuilt an unchanged state would make
\scm{(≠° (ƛ (a) a) (ƛ (b) b))} look satisfiable (the λ analogue of \scm{(≠° 1 1)} in the pitfall
after \cref{prop:micro-eq}). For the λ
layer this has a concrete consequence: the names that the ξ rule opens under binders live outside the
state, so a ξ step that binds nothing returns $s$ itself. So the λ disequality below is recognised
as an equality that already holds, and fails, exactly as micro's \scm{(≠° 1 1)} does:
\begin{scheme}
(pretty-print (°->list #f (fresh° (q) (≠° (ƛ (a) a) (ƛ (b) b)))))
(pretty-print (°->list #f (fresh° (q) (≠° 1 1))))
\end{scheme}
\begin{answers}
()
()
\end{answers}

(I2) is what keeps the stores honest. \scm{μkanren-verify-D/touched} re-verifies only the prefixes
whose variables appear among the new bindings; a binding made by a hook and not reported would leave a
violated disequality in D. This is Trick~1's lesson (\cref{sec:micro-noedit}): a state produced by raw
unification must go through the D/A/T chain with its exact bindings.

\begin{pitfall}[A hook that forgets a binding must also drop it from the bindings]
(I2) has a converse. When \texttt{post} takes a binding back with \scm{μkanren-state-forget}, it must
remove it from the bindings it answers; otherwise the stores would be re-verified against a binding that
S no longer holds. The λ layer's \scm{μkanren-enforce} does both, through \scm{μkanren-drop}.
\end{pitfall}

\subsection{From closed points to hooks}\label{sec:design-map}

\Cref{tab:design-map} reads the protocol against \cref{tab:micro-open}: for each closed point of
\micro, what the edit opens and what it leaves closed on purpose.

\begin{table}[ht]
\centering\footnotesize
\caption{Closed points of \micro{} and the hooks that open them.}
\label{tab:design-map}
\renewcommand{\arraystretch}{1.2}
\begin{tabularx}{\textwidth}{@{}>{\raggedright\arraybackslash}p{3.8cm}>{\raggedright\arraybackslash}X>{\raggedright\arraybackslash}p{3.6cm}@{}}
\toprule
closed point (\cref{tab:micro-open}) & opened by & still closed\\
\midrule
the field list of the state, its positional constructions & the field \texttt{X}, copied at the two constructions & the six-ary \texttt{μkanren-state-match}\\
the unification \texttt{cond} & \texttt{unify} in the structure--structure and \texttt{else} clauses & the variable clauses, pair--pair\\
the binding site \texttt{μkanren-state-update} & \texttt{atom?} inside the occurs check; \texttt{occurs} on the failure path; the scope check moved to \texttt{post} & var--var orientation (older to newer)\\
the post-\texttt{=°} chain & \texttt{post}, threaded; \texttt{μkanren-post-unify} exported; \texttt{μkanren-state-forget} & ---\\
reification (D/A/T only) & \texttt{value}, \texttt{repr}, \texttt{terms}, \texttt{forms}; \texttt{parts} in projection mode & purification and subsumption of micro's stores\\
\texttt{°->§}: fixed empty initial state & not needed: \texttt{X = '()} means empty; \texttt{environment} for ground answers & fixed query index 0\\
\midrule
(listed as open: \texttt{≠°}, \texttt{absent°} and types work on any record, slot by slot) & \texttt{diseq}, \texttt{parts}, \texttt{tag} & ---\\
\bottomrule
\end{tabularx}
\end{table}

The three needs listed at the end of \cref{ch:micro} correspond to three rows: \enquote{a named place in
the state} is \texttt{X}, \enquote{a call after every binding} is \texttt{post}, and \enquote{a say in
reification} is \texttt{value}, \texttt{repr}, \texttt{terms} and \texttt{forms}. The rest of the
protocol exists because micro's stores and occurs check, written for first-order data, look into records
slot by slot.

\begin{tutorial}[One \texttt{=°} on a λ-term, through micro and the hooks]
To watch the protocol at work we re-register the λ extension under its own name with some hooks wrapped
by a printer. Registering under the same name replaces it (\cref{sec:design-registry}). Micro exports no
accessor of the hooks, so the script reads the slots of the registered record directly, with
\scm{record-instance-slot} of \scm{(chicken memory representation)} (imported with \scm{srfi-1},
besides \scm{(aux base)}, \scm{(aux stream)}, \micro{} and \scm{(aux kanren lambda)}); this is for tracing only.
\begin{scheme}
(define fields '(name term? atom? unify occurs post diseq tag parts
                 value repr terms forms environment))
(define (orig h)
  (record-instance-slot μkanren-λ-extension (list-index (λ (f) (eq? f h)) fields)))
(define (D t) (if (μkanren-λ-term? t) (μkanren-datum t) t))
(define (B as) (map (λ (p) (list (car p) ':= (D (cdr p)))) as))
(define (say . xs) (display "  ") (for-each display xs) (newline))
(define traced
  (make-μkanren-extension 'λProlog     ; same name: replaces the λ extension
    term?: (orig 'term?)
    atom?: (λ (u) (let1 (r ((orig 'atom?) u))
                    (say "atom?  " (D u) " => " r) r))
    unify: (λ (u v s as)
             (say "unify  " (D u) " =? " (D v))
             (let1 (r ((orig 'unify) u v s as))
               (say "unify  => " (and r (B (cadr r)))) r))
    occurs: (orig 'occurs)
    post: (λ (s as)
            (say "post   " (B as) " at depth " (μkanren-λ-depth s))
            (let1 (r ((orig 'post) s as))
              (say "post   => " (and r (B (cadr r)))) r))
    diseq: (orig 'diseq) tag: (orig 'tag) parts: (orig 'parts)
    value: (orig 'value) repr: (orig 'repr)
    terms: (λ (s) (let1 (r ((orig 'terms) s)) (say "terms  => " r) r))
    forms: (λ (s s*) (let1 (r ((orig 'forms) s s*)) (say "forms  => " r) r))
    environment: (orig 'environment)))
(μkanren-extension-register! traced)
(pretty-print (°->list/datum (fresh° (q) (∀° (x) (=° (@ q x) `(g ,x ,x))))))
(pretty-print (°->list #f (fresh° (q) (∀° (x) (=° q x)))))
\end{scheme}
\begin{answers}
  unify  (@ _0 (eigen x)) =? (g (eigen x) (eigen x))
  atom?  (ƛ (x) (g x x)) => #f
  atom?  (bvar 0) => #t
  atom?  (bvar 0) => #t
  unify  => ((_0 := (ƛ (x) (g x x))))
  post   ((_0 := (ƛ (x) (g x x)))) at depth 1
  post   => ((_0 := (ƛ (x) (g x x))))
  terms  => ()
  forms  => ()
((ƛ (x) (g x x)))
  atom?  (eigen x) => #t
  post   ((_0 := (eigen x))) at depth 1
  post   => #f
()
\end{answers}
\textbf{The first query}, \lp{pi x\ F x = g x x}, step by step.
\begin{enumerate}
\item \scm{∀°} raises the depth to 1 and makes the eigenvariable \scm{x} of level 1; \scm{q} is
  \scm{_0}, of level 0.
\item Micro's \scm{=°} calls \scm{μkanren-state-unify}. Both sides are walked: the left one is an
  application record, the right one a pair. No variable clause applies, the left side is a structure
  but the right one is not, and the pair--pair clause does not apply: the final \scm{else} finds the
  owning extension and calls \texttt{unify}.
\item The hook head-normalises both sides. The left one is flexible, \scm{_0} applied to the pattern
  \scm{(x)}: one eigenvariable, whose level 1 exceeds the level 0 of \scm{q}. So this is flex--rigid. Variable
  elimination replaces \scm{x} by the bound variable of the new abstraction, and the hook calls micro
  back to bind \scm{_0 := (ƛ (x) (g x x))}.
\item That binding goes through micro's binding site and its occurs check. The check looks into the
  abstraction (\texttt{atom?} is false) and stops at each \scm{(bvar 0)} (\texttt{atom?} is true). No
  occurrence of \scm{_0}: the binding is made, and the hook returns it.
\item Back in \scm{=°}, the state has changed, so \scm{μkanren-post-unify} runs the \texttt{post} hook
  on the new binding. Levels exist, and \scm{_0} is below the highest level reached, so the scope check
  runs: the bound term holds no eigenvariable of level $>0$ (\scm{x} has become an index), and no
  variable to lower. The binding passes; E is empty, so nothing is woken. Then micro verifies D, A and
  T, all empty.
\item The \scm{∀°} body has answered; \scm{∀°} gives the state back depth 0 and discharges the
  constraints mentioning \scm{x} (there are none).
\item Reification asks the extension for the terms of its residual constraints (none) and for its
  forms (none); the \texttt{value} and \texttt{repr} hooks, not traced here, print the abstraction.
\end{enumerate}
\textbf{The second query}, \lp{sigma Y\ pi x\ Y = x}, never calls \texttt{unify}: \scm{q} is a variable,
so micro's variable clause binds it directly. The occurs check meets the eigenvariable, an atom, and
does not look into it. Micro makes the binding \scm{_0 := x}; but the \texttt{post} hook sees that
\scm{x} has level 1 and \scm{_0} level 0, which violates condition (2) of A1, and answers \scm{#f}.
\scm{=°} fails before any other goal could see the binding. This is A1(c) at work: the binding was
made by micro's own clause, and the extension still had the last word.
\end{tutorial}

\subsection{Cost}\label{sec:design-cost}

\Cref{tab:roadmap-perf-protocol} in \cref{ch:roadmap} gives the cost of the protocol on first-order
search, compiled with \texttt{csc -O2 -d1}: the unedited micro against the edited micro, without and
with \scm{(aux kanren lambda)} loaded but unused.
The protocol costs at most 1\% end to end on \scm{append°} (5923 to 5959~ms in round 2), within the
2.5\% round-to-round noise of the unedited binary. With the λ extension loaded, the disequality-heavy
perm6 costs 8--12\% (d2 cost 45--50\%): the remainder is the \texttt{post} hook's look at the new
bindings while D is not empty. One cost is not removed and is documented: once any \scm{∀°} has run,
levels exist, and every later binding below the top level gets the scope check, which costs first-order
search 1.7 to 2.3~times (\cref{ch:roadmap} plans one walk for the occurs and the scope check).

% ================================================================================================
\section{Review history}\label{sec:design-reviews}

After the choice of d2 and the grafts, the proposal went through three adversarial review rounds. Each
round attacked the proposal itself with hand-written probes and with randomly generated goals, run
with their conjunctions in both orders. Every report was first reproduced, then fixed, a code fix
with a regression case (in round 3 each new case was checked to fail on the proposal before the fix).

\begin{table}[ht]
\centering\footnotesize
\caption{The three adversarial review rounds (design notes, §10.4--10.6).}
\label{tab:design-reviews}
\renewcommand{\arraystretch}{1.2}
\begin{tabularx}{\textwidth}{@{}lrr>{\raggedright\arraybackslash}X@{}}
\toprule
round & reports & regression cases & typical defects\\
\midrule
1 & 13 & 12 (44 assertions) & disequality prefixes holding an eigenvariable left in D (a free \texttt{(eigen x)} printed, a crash); pruning inside an argument that a non-pattern may drop; occurs check not modulo β; raw record slots in projected \emph{absento}; a binder hint capturing \texttt{cons} in ground answers\\
2 & 8 & 6 (31 assertions) & raising at ∀ exit not closing every constraint; η-redexes created by later bindings; occurs failure on non-rigid arguments; the extensional reading absorbed by T; first-order cost after a \texttt{∀°}\\
3 & 12 & 9 (53 assertions) & non-patterns with functional arguments; disequalities over unreachable arguments; \emph{absento} of an eigenvariable under a type; user types on λ-terms; η-short occurs check; a datum reducing to an atom; two types on one flexible term; first-order cost with D or A non-empty; λ-terms in the show module's cells view\\
\bottomrule
\end{tabularx}
\end{table}

All 33 reports were real. In round 1 the twelve code defects were fixed, each pinned by a case, and one
documentation error was corrected. In round 2 two reports were one defect; six were fixed with
regression cases, and the seventh, the cost of first-order search after a \scm{∀°}, was measured and
documented rather than changed, since the cheap fix would give up the guarantee that host code cannot
smuggle an eigenvariable into an older variable. In round 3 three pairs of reports were one defect each;
all nine defects were fixed with a regression case each. Round 3 also produced the generic change to
\scm{absent°} of \cref{sec:design-absent}. The defects themselves belong to the λ layer and are discussed
where it is described (\cref{ch:module}).

% ================================================================================================
\section{Summary}\label{sec:design-summary}

\begin{itemize}
\item The requirements are nine acceptance items: the scope of ∀ under Nadathur, Jayaraman and Kwon's
  binding conditions, for \emph{every} binding (A1); \rulename{augment} with the book's clause order and variable
  sharing (A2); type inference as the test of ∀ and ⊃ together (A3); λ-tree syntax with α, β and η (A4);
  $\Llam$ unification with mgus and delayed non-patterns (A5); λProlog's answers where Ma and Friedman
  differ (A6); the book's programs (A7); micro's stores on λ-terms (A8); micro's suites unchanged
  (A9). Order independence runs through all of them. They are pinned by 83 cases and 320 assertions.
\item Of the four meanings of \scm{=°} on λ-terms, pattern unification with universe levels is
  retained: it is λProlog's, it is decidable and unitary (one state per \scm{=°}), and levels give
  \rulename{generic} for one integer per variable and one comparison per binding.
\item Three prototypes were built: a layer with no edit (d1), a hooks protocol (d2), a native engine
  (d3). d1 cannot check bindings made by micro's \scm{=°}, so eigenvariables leak through \scm{cons°};
  d3 had a crash and an η bug that made \scm{≠°} unsound; d2 had seven non-soundness defects. The
  reviews scored them 15.5, 23 and 21; d2 became the base, and d3's pedagogical strengths and d1's
  precision were grafted onto it.
\item The protocol adds a seventh state field \texttt{X} that micro copies and never reads, a global
  registry of extensions dispatched by \texttt{term?}, and thirteen hooks: \texttt{atom?} and
  \texttt{occurs} in the binding site, \texttt{unify} in the unification \scm{cond}, \texttt{post} in the
  chain after \scm{=°}, \texttt{diseq}, \texttt{tag} and \texttt{parts} in the stores, \texttt{value},
  \texttt{repr}, \texttt{terms} and \texttt{forms} in reification, \texttt{environment} in the runner.
\item The scope check lives in \texttt{post}: every binding passes through some \scm{=°}, and every
  \scm{=°} ends in \scm{μkanren-post-unify}, which runs the hooks before D, A and T are verified.
  An extension goal that unifies on its own must end in \scm{μkanren-post-unify} too.
  \scm{μkanren-state-forget} lets \texttt{post} turn a binding into a suspended equation.
\item A client must return the very state when it changes nothing (eq?-entailment) and report every
  binding it makes (extended bindings).
\item The edit is +235/−41 lines, λ-agnostic, and costs about 1\% on first-order search; micro's 125
  cases pass unchanged with and without the λ layer. Three adversarial rounds sent 13, 8 and 12 reports,
  all real; every defect but one documented cost was fixed with a regression case.
\end{itemize}

\chapter{The module \texorpdfstring{\texttt{(aux kanren lambda)}}{(aux kanren lambda)}}\label{ch:module}

% ------------------------------------------------------------------------------------------------
\section*{What and why}

\Cref{ch:micro} took \micro{} apart and ended, in \cref{sec:micro-ext}, with the list of points an extension
cannot reach from outside: the state has no free slot, and the unifier, the walkers and the reifier are closed.
\Cref{ch:design} opened exactly those points with a small, λ-agnostic protocol of hooks and explained why the
proposal is a client module of that protocol rather than a native engine. This chapter is about the client
itself: the module \scm{(aux kanren lambda)}, which brings into \micro{} the three ingredients of Part~I,
\begin{itemize}
  \item λ-tree syntax (\cref{ch:lambdatree}): object-level binders are meta-level abstractions;
  \item higher-order pattern unification, $\Llam$ (\cref{ch:hounif}), with universe levels in place of the
        mixed prefix and with non-patterns suspended (\enquote{dynamic $\Llam$});
  \item universal and hypothetical goals, \rulename{generic} and \rulename{augment}, over hereditary Harrop formulas (\cref{ch:hh}).
\end{itemize}

The chapter is the formal heart of the proposal. It states what the terms are (\cref{sec:module-terms}),
what the module keeps in the state (\cref{sec:module-state}), the unification rules with their side
conditions and the termination argument (\cref{sec:module-unif}), the rules for \scm{∀°} and \scm{⊃°} and
why levels make them sound (\cref{sec:module-hh}), how micro's disequality, absence and type constraints
read λ-terms (\cref{sec:module-stores}), how answers are printed (\cref{sec:module-reify}), and, as a
reference, the exported API (\cref{sec:module-api}). The curated book examples are the subject of
\cref{ch:examples}; the known limits, the costs and the merge plan are in \cref{ch:roadmap}.

{\emergencystretch=2em
The module is 1255 lines, 885 of them neither blank nor comment. In the repository it lives, with its
tests and examples, under \texttt{doc/\allowbreak hol-kanren/\allowbreak prototype/}, and after merging as
\texttt{src/\allowbreak aux.kanren.lambda.scm}, with the suites \texttt{src/\allowbreak test/\allowbreak microkanren-lambda.scm} (83 cases, 320
assertions) and \texttt{src/\allowbreak test/\allowbreak microkanren-lambda-limits.scm} (6 cases). Every Scheme snippet below was run
against the module and micro with the protocol applied; its output is pasted verbatim under it. Before
the imports below work, the edited micro must be built: \cref{sec:roadmap-run} shows how, without
installing anything, and \cref{sec:roadmap-merge} how to merge it. The snippets
import\par}
\begin{scheme}
(import scheme (chicken base) (chicken pretty-print)
        (chicken condition) (chicken memory representation)
        (aux base) (aux stream) (aux kanren micro) (aux kanren micro show)
        (aux kanren lambda))
\end{scheme}
Several of them use a \emph{probe}: a goal that prints part of the state and succeeds. A goal of \micro{} is a
procedure from a state to a stream of states (\cref{def:micro-goal}), so a probe is three lines; each snippet
shows the one it uses.

% ------------------------------------------------------------------------------------------------
\section{Terms}\label{sec:module-terms}

\subsection{Four records}

First-order structure stays Scheme data, exactly as in \micro: pairs, vectors, foreign records and atoms
(\cref{def:micro-term}). The module adds four record types, defined at the top of the file:
\begin{scheme}
(define-record μkanren-lam hint body)        ; λ.body, locally nameless
(define-record μkanren-app head args)        ; head applied to a list of arguments
(define-record μkanren-eigen name level id)  ; a scoped constant of GENERIC
(define-record μkanren-bvar index)           ; a de Bruijn index, only inside bodies
\end{scheme}

\begin{definition}[λ-terms of the module]\label{def:module-terms}
Extend the terms of \micro{} (\cref{def:micro-term}) with
\[
  t \;::=\; \dots \;\mid\; c^{\ell} \;\mid\; \#i \;\mid\; \lambda^{h}.\,t \;\mid\; t\;@\;\langle t_1,\dots,t_n\rangle
\]
where $c^\ell$ is an eigenvariable of level $\ell\in\mathbb N\cup\{\infty\}$ (\scm{μkanren-eigen}), $\#i$ a
bound variable with de Bruijn index $i\ge 0$ (\scm{μkanren-bvar}), $\lambda^h.t$ an abstraction whose binder
carries the printing hint $h$, a symbol (\scm{μkanren-lam}), and $t\,@\,\bar t$ an application with $n\ge 1$
(\scm{μkanren-app}). A term is \emph{locally closed} when every $\#i$ in it lies below at least $i+1$
abstractions. Only locally closed terms are ever bound to variables or stored in constraints.
\end{definition}

The representation is \emph{locally nameless} \cite{Chargueraud2012}: bound variables are de Bruijn indices
\cite{deBruijn1972}, free ones are names (logic variables and eigenvariables). Three design choices follow
from it.
\begin{itemize}
  \item \textbf{The hint is only for printing}, and it is an \emph{interned} symbol. The binder's identity is
    its position. Two terms built alike are therefore \scm{equal?}, so λ-terms can be grouped by
    \scm{groupby°} and compared by the host.
  \item \textbf{An eigenvariable's identity is its \texttt{id}}, an uninterned symbol, compared with \scm{eq?}.
    Its \texttt{name} is a hint and its \texttt{level} is the universe it lives in (\cref{sec:module-state}).
    Level $\infty$ (\scm{+inf.0}) marks the names that the unifier opens when it descends under binders and
    the placeholders that \scm{ƛ} uses while building a body.
  \item \textbf{The first 64 bound variables are preallocated} in a vector, so the common indices allocate
    nothing.
\end{itemize}
The de Bruijn operations (shift, instantiate, close) never walk the substitution: a term bound to a variable
is locally closed, so a binding never contributes a loose index. Instantiation shifts the substituted term
only when it has loose indices itself, which is decided once per instantiation.

\subsection{The constructors \texorpdfstring{\texttt{ƛ}}{ƛ} and \texorpdfstring{\texttt{@}}{@}}

The user never builds these records by hand. Two constructors give the surface syntax of λ-tree terms:
\begin{scheme}
(define-syntax ƛ
  (syntax-rules ()
    ((_ () body) body)
    ((_ (x y ...) body) (μkanren-abstract 'x (λ (x) (ƛ (y ...) body))))))

; λ.(f p), with the placeholder p closed into #0
(define (μkanren-abstract name f)
  (let1 (p (μkanren-eigen* name +inf.0))
    (μkanren-lam/η name (μkanren-close (f p) p))))

(define (@ f . args) (μkanren-beta f args))   ; β at once on an abstraction
\end{scheme}
\scm{(ƛ (x) body)} binds the \emph{Scheme} variable \scm{x} to a fresh placeholder eigenvariable while
\scm{body} is evaluated, then closes that placeholder into the index \#0. The object binder is thus a Scheme
binder: shadowing, lexical scope and quasiquotation all behave as a Scheme programmer expects, which is the
point of λ-tree syntax \book{177}. Two invariants are enforced at construction.
\begin{itemize}
  \item \textbf{η-short.} \scm{μkanren-lam/η} builds $\lambda x.\,(h\,\bar a\,x)$ as $h\,\bar a$ when $x$ is not
    free in $h\,\bar a$. The unifier builds its abstractions the same way. This matters because micro's own
    occurs check runs before any hook can see an equation: $\lambda x.\,(F\,x) = F$ holds by η, but in
    η-long form micro would see $F$ inside the value it binds $F$ to and fail; built η-short, the left side
    \emph{is} $F$ and the equation is trivial.
  \item \textbf{β at construction.} \scm{@} applied to an abstraction β-reduces at once, and nested
    applications are flattened: \scm{(@ (@ f a) b)} is \scm{(@ f a b)}.
\end{itemize}

\begin{example}[Construction and printing]\label{ex:module-construct}
Every record has a printer that shows its term as data. The raw record of $\lambda x.\,(f\ x)$ shows the hint
\texttt{x} and the index in the body:
\begin{scheme}
(ƛ (x y) (cons x y))
(record->vector (ƛ (x) `(f ,x)))     ; from (chicken memory representation)
(equal? (ƛ (x) `(f ,x)) (ƛ (x) `(f ,x)))
(equal? (ƛ (x) x) (ƛ (y) y))
(μkanren-α=? (ƛ (x) x) (ƛ (y) y))
(@ (ƛ (x) `(g ,x ,x)) 'a)
(ƛ (x) (@ 'g x))
(ƛ (x) `(g ,x))
(@ 'g 'a)
(equal? (@ 'g 'a) '(g a))
(ƛ (a) (ƛ (a) a))
\end{scheme}
\begin{answers}
  => (ƛ (x y) (x . y))
  => #(aux.kanren.lambda#μkanren-lam x (f #<aux.kanren.lambda#μkanren-bvar>))
  => #t
  => #f
  => #t
  => (g a a)
  => g
  => (ƛ (x) (g x))
  => (@ g a)
  => #f
  => (ƛ (a a1) a1)
\end{answers}
Line by line: nested binders print merged, \scm{(ƛ (x y) ...)}; the body of a \texttt{μkanren-lam} holds an
index, not the name; terms built alike are \scm{equal?}; α-variants with \emph{different} hints are not
\scm{equal?} but are α-equal; \scm{@} reduced the redex; $\lambda x.\,(g\,x)$ written with \scm{@} is the
η-redex of $g$ and was built as \texttt{g}; written with a quasiquoted list it is no η-redex (the list
\scm{(g x)} is data, not an application), so it stays an abstraction; the λ-tree application of a constant
prints with \texttt{@} and is not the list \scm{(g a)}; shadowing is lexical, and the printer renames the
inner binder apart (\cref{sec:module-reify}).
\end{example}

\subsection{Three kinds of application}

An application is classified by its head \emph{after walking it} in the current substitution:
\begin{definition}[Flexible, β-redex, rigid; head normal form]\label{def:module-hnf}
An application $h\,@\,\bar t$ is \emph{flexible} if $h$ walks to an unbound logic variable, a
\emph{β-redex} if $h$ walks to an abstraction, and \emph{rigid} otherwise (a constant, an eigenvariable, a
bound variable). The \emph{head normal form} $\mathrm{hnf}(t)$ (\scm{μkanren-hnorm}) walks $t$, flattens
nested applications, contracts head β-redexes, and contracts an abstraction whose body has \emph{become} an
η-redex through a binding (e.g.\ $\lambda x.\,F\,(G\,x)$ once $G:=\lambda z.z$).
\end{definition}
Under the $\Llam$ restriction every contraction hnf performs is a $\beta_0$ renaming
(\cref{def:lambdatree-beta0}). On instantiated heads \scm{@} gives full β, which is what the book calls
dynamic $\Llam$ \book{225--226}.
\begin{scheme}
(μkanren-hnorm (@ (ƛ (x) (@ x 'b)) (ƛ (y) `(g ,y ,y))) μkanren-state-empty)
(°->list #f (fresh° (q F) (=° q (@ F 'a))))
(°->list/datum (fresh° (q F) (=° q (list F (@ F 'a))) (=° F (ƛ (z) `(h ,z)))))
\end{scheme}
\begin{answers}
(g b b)
((λ (α) (@ α 'a)))
(((ƛ (z) (h z)) (h a)))
\end{answers}
The second answer is a flexible term, printed as \scm{(@ α 'a)}; once $F$ is bound (third line) the same
term is a β-redex and its value is \scm{(h a)}.

\subsection{Scheme lists versus λ-tree application}

The object language of the book's programs is written with Scheme lists, as \micro{} always did:
$\mathtt{abs}\ M$ is \scm{`(abs ,M)}, $\mathtt{app}\ M\ N$ is \scm{`(app ,M ,N)}, with $M$ a meta-level
abstraction for a binder. The λ-tree application \scm{(@ 'g 'a)} and the list \scm{'(g a)} are distinct terms
and never unify:
\begin{scheme}
(°->list #f (fresh° (q) (=° (@ 'g 'a) '(g a))))
\end{scheme}
\begin{answers}
  => ()
\end{answers}
This is deliberate. One of the three prototypes identified the two (a spine notation for lists) and lost
commutativity of conjunction: with η, the same pair of goals gave different answers in the two orders. Keeping
lists as data and \scm{@} as the only application keeps \micro's first-order unification untouched for lists,
and confines the higher-order rules to terms that contain a λ-record.

\begin{pitfall}[What is and is not an η-redex]
η-contraction applies to λ-tree applications only. \scm{(ƛ (x) (@ 'g x))} \emph{is} \texttt{g}, but
\scm{(ƛ (x) `(g ,x))} is an abstraction over a two-element list, and stays one. The module is untyped, so it
cannot η-\emph{expand} a term against its type: η-expansion happens only on demand, when one side of an
equation is an abstraction (rule \rulename{η} below), as in Miller's untyped treatment
\cite[§9.3]{Miller1991}. A rigid list such as \scm{'(g a)} therefore equals an abstraction only when the
abstraction's body equals the λ-tree application \scm{(@ '(g a) x)} of the list to the bound variable (and
\scm{ƛ} already builds such an η-redex as the list itself); never when the body is a list such as
\scm{`(g a ,x)}.
\end{pitfall}

% ------------------------------------------------------------------------------------------------
\section{The extension's state}\label{sec:module-state}

\subsection{The record in \texorpdfstring{\texttt{X}}{X}}

The protocol gives \micro's state a seventh slot \texttt{X} that micro copies but never reads
(\cref{ch:design}). The module keeps one record in it:
\begin{scheme}
(define-record μkanren-hh depth top L Ln Γ E N)
(define μkanren-hh0 (make-μkanren-hh 0 0 empty/sbral 0 '() '() '()))
(define (μkanren-hh s) (let1 (x (μkanren-state-X s)) (if (null? x) μkanren-hh0 x)))
\end{scheme}
\texttt{X} = \scm{'()}, the value in \scm{μkanren-state-empty}, reads as \enquote{everything empty}, so the
initial state of every runner works unchanged.

\begin{definition}[The λ state]\label{def:module-state}
The record $\langle d,\ \mathit{top},\ L,\ L_n,\ \Gamma,\ E,\ N\rangle$ holds:
\begin{description}[style=nextline,leftmargin=2.2em]
  \item[$d$ (\texttt{depth})] the current universe level, raised by \scm{∀°} for the extent of its body;
  \item[$\mathit{top}$] the highest level reached so far in this branch;
  \item[$L,\,L_n$] the levels of the logic variables: $L$ is a skew-binary random-access list (\aux's
    \texttt{sbral}) indexed like the substitution $S$, holding the levels of the variables with index
    $<L_n$. The level of variable $i$ is
    \[
      \ell(i) \;=\; \begin{cases} 0 & i<0 \text{ (a reified variable)}\\
                                   L[L_n-1-i] & 0\le i<L_n\\
                                   d & i \ge L_n \end{cases}
    \]
    so a variable created since $L$ was last extended has the current depth. $L$ is padded with $d$ exactly
    when the depth changes and whenever the unifier creates a variable of another level; lookup is
    $O(\log n)$;
  \item[$\Gamma$] the assumed clauses: an alist from a relation to its clauses, most recent first, one entry
    per distinct relation;
  \item[$E$] the suspended constraints: equations \scm{(= u v)} outside $\Llam$ (or bindings whose scope check
    a non-pattern decides), types \scm{(tag tag u)} on flexible terms, and absentos \scm{(absent d u)} whose
    datum outlived its eigenvariable's scope;
  \item[$N$] the disequalities that involve λ-terms or eigenvariables, as \scm{(reading u . v)}. micro's D never
    holds one (\cref{sec:module-stores}).
\end{description}
\end{definition}

Levels are the module's whole account of quantifier order. A logic variable created at depth $d$ stands for
an $\exists$ placed after the first $d$ universals of its branch; an eigenvariable created by the $(d{+}1)$-th
\scm{∀°} has level $d+1$. The binding conditions of Nadathur, Jayaraman and Kwon \cite{NadathurJayaramanKwon1995}
(\cref{def:rules-binding,def:lkanren-ui}) are then stated on levels: a variable of level $\ell$ may be bound
only to a term whose eigenvariables have level $\le\ell$, and the variables of level $>\ell$ in that term are
\emph{lowered} to $\ell$.

\subsection{Copied, never mutated}

The record is immutable. Every change builds a new one with \scm{μkanren-hh-with}, a keyword-argument copy, and
installs it with \scm{(μkanren-state-with s X: x)}. Since a state is a value and \scm{or°} gives each branch the
state it was given, two branches can never see each other's depth, Γ or suspended constraints. What needs
care is the \emph{end} of a scope: \scm{∀°} and \scm{⊃°} change the record for their body only, and each
answer of the body must give the change back. Both do so on every answer state (\cref{sec:module-hh}). The
following probe shows depth and levels through two nested \scm{∀°}:
\begin{scheme}
(define ((probe label . vars) s)
  (display label) (display ": depth ") (display (μkanren-λ-depth s))
  (display ", levels ") (write (map (λ (v) (μkanren-λ-level v s)) vars))
  (display ", E ") (write (μkanren-λ-delayed s))
  (display ", N ") (write (μkanren-λ-disequalities s)) (newline)
  (✓° s))
(°->list #f
  (fresh° (q)
    (probe "start" q)
    (∀° (x)
      (fresh° (Y)
        (probe "inside ∀x" q Y)
        (∀° (z) (fresh° (W) (probe "inside ∀x∀z" q Y W)))
        (probe "back in ∀x" q Y)
        (=° (@ q x) (@ Y x x))
        (probe "after q x = Y x x" q Y)))
    (probe "after ∀x" q)))
(°->list #f (fresh° (q) (or° (∀° (x) (probe "branch 1" q) (=° q 1))
                              (and° (probe "branch 2" q) (=° q 2)))))
\end{scheme}
\begin{answers}
start: depth 0, levels (0), E (), N ()
inside ∀x: depth 1, levels (0 1), E (), N ()
inside ∀x∀z: depth 2, levels (0 1 2), E (), N ()
back in ∀x: depth 1, levels (0 1), E (), N ()
after q x = Y x x: depth 1, levels (0 1), E (), N ()
after ∀x: depth 0, levels (0), E (), N ()
((λ (α) (ƛ (x) (@ α x x x))))
branch 1: depth 1, levels (0), E (), N ()
branch 2: depth 0, levels (0), E (), N ()
((λ () 1) (λ () 2))
\end{answers}

\begin{tutorial}[Reading the probe]
$q$ is created at depth~0, $Y$ inside \scm{∀° (x)} at depth~1, $W$ at depth~2; after each \scm{∀°} the depth
is back where it was. Written as a mixed prefix, the goal is $\exists q\,\forall x\,\exists Y\,\forall z\,\exists W$,
and each level is the number of universals to the left of the quantifier: $q{:}0$, $x{:}1$, $Y{:}1$, $z{:}2$,
$W{:}2$. The equation $q\,x \doteq Y\,x\,x$ is solved inside the scope, and its answer shows the
levels at work. $q$ has level~0 and its argument $x$ (level~1) is a pattern for it, so the unifier eliminates
$q$: $q := \lambda x.\,t'$ where $t'$ is $Y\,x\,x$ with $x$ turned into the bound variable. But $Y$ has
level~1 and $q$ level~0: $q$'s value may not mention a variable that could see $x$ unless $x$ is passed
explicitly. $Y\,x\,x$ is no pattern ($x$ is repeated), but both its arguments are $q$'s own argument $x$,
so none of them has to be dropped. $Y$ is therefore \emph{lowered} and \emph{raised} at once, $Y := \lambda a\,b.\,Y'\,a\,b\,x$ with
$\ell(Y')=0$, and $q := \lambda x.\,Y'\,x\,x\,x$. The answer \scm{(ƛ (x) (@ α x x x))} is that binding; the
rules responsible are \rulename{FR} and case~(6c) of variable elimination, \cref{def:module-elim}. The last
two lines show the branches: \scm{∀°} raised the depth in branch~1 only.
\end{tutorial}

The module's hooks keep the two protocol invariants of \cref{def:design-invariants}; in particular, every
function of the module that may leave the state unchanged is written to return it \scm{eq?}.

% ------------------------------------------------------------------------------------------------
\section{Unification}\label{sec:module-unif}

\subsection{Where the rules run}

\micro's unifier \scm{μkanren-state-unify} still does all first-order work: variables, pairs, vectors, foreign
records and atoms. Through the protocol's \texttt{unify} hook it hands any pair with a walked λ-record on one
side, neither side a variable, to \scm{μkanren-λ-unify}, which applies the rules of \cref{fig:module-rules}.
The rules call back into \scm{μkanren-state-unify} for subterms, so the two unifiers interleave freely. After
\scm{=°} has made its bindings, the \texttt{post} hook \scm{μkanren-λ-post} runs before micro verifies D, A
and T:
\begin{enumerate}
  \item \emph{settle}: the scope check \textbf{Check} of every new binding, then the wake-up of $E$, in rounds
    while a round binds (\cref{sec:module-wake});
  \item \emph{migrate}: the prefixes of D that now hold a λ-term or an eigenvariable move to N, and the types of
    A on variables just bound to λ-terms move to the \texttt{tag} hook (\cref{sec:module-stores});
  \item \emph{verify N}: every λ-disequality is tried again.
\end{enumerate}

\subsection{Notation}

Write $\ell(X)$ for the level of a logic variable and $\ell(c)$ for that of an eigenvariable; the names the
unifier opens under binders have $\ell(c)=\infty$. A \emph{name} is an eigenvariable or, inside variable
elimination, a bound variable of the term being walked.
\begin{definition}[Pattern, read off levels]\label{def:module-pattern}
A list of arguments $\bar a = a_1\dots a_n$ is a \emph{pattern for} a variable $F$ if every $\mathrm{hnf}(a_i)$
is an eigenvariable, the $a_i$ are pairwise distinct, and $\ell(a_i) > \ell(F)$ for every $i$
(\scm{μkanren-pattern?}). A flexible term $F\,\bar a$ is a pattern when $\bar a$ is a pattern for $F$.
\end{definition}
For example, inside \scm{(∀° (x y) ...)} with $\ell(F)=0$, $\ell(x)=1$, $\ell(y)=2$, the lists $\langle x,y\rangle$
and $\langle y\rangle$ are patterns for $F$, while $\langle x,x\rangle$ (repeated), $\langle \mathtt{a}\rangle$ (a
constant) and $\langle Y\rangle$ (a variable) are not. The same definition, with a bound variable of the
term being walked counted as an eigenvariable of level $\infty$, is used inside variable elimination (case (6b)
below).
This is Miller's $\Llam$ condition (\cref{def:hounif-llambda}): the arguments are distinct variables
\enquote{quantified inside the scope of $F$}. With a prefix, \enquote{inside the scope} means \enquote{to the
right of $F$}; with levels it means \enquote{of a level above $F$'s}, which is how Nadathur and Linnell read the
condition without a prefix \cite[§4]{NadathurLinnell2005}. An eigenvariable $F$ could already see is not a
pattern argument: in $\forall x\,\exists F.\ F\,x = \dots$ the term $F\,x$ is no pattern, since $F$ may use
$x$ directly and the equation has several incomparable solutions.

Throughout, $\lambda.b$ is an abstraction with body $b$, $b[c/0]$ its body with \#0 instantiated to $c$,
$h\,\bar u$ an application, and $F,G,H,X,Y$ logic variables. Both sides are put in hnf before a rule is
chosen.

\subsection{The rules}

The rules of \cref{fig:module-rules} are transformations, read top-down as in \cref{fig:rules-fo,fig:rules-pattern}: the equation above the
line is replaced by the equations, bindings or suspension below it. They are the rules of \cref{fig:rules-pattern}
with the prefix erased into levels; \cref{sec:hounif-algo} gives the same rules under other names, recalled
in parentheses. \rulename{ξ} is \rulename{abs} (\rulename{Abs}), \rulename{η} is \rulename{eta}
(\rulename{Eta}), \rulename{Rigid} is \rulename{rr-dec} and \rulename{rr-clash} (\rulename{Dec} and
\rulename{Clash}), \rulename{FF-same} is \rulename{ff-same} (the same-head case of \rulename{FlexFlex}),
\rulename{Delay} is \rulename{delay}. \rulename{FR} covers \rulename{elim} (\rulename{Elim}) with its pruning,
\rulename{occurs} and \rulename{scope}, and also \rulename{ff-diff} (\rulename{FlexFlex} with distinct heads): instead of binding both heads to a new $H$,
it eliminates $F$ and prunes $G$ only where needed, which gives the same mgu up to the name of $H$
(\cref{ex:module-ff}). \rulename{Var} and \rulename{Occurs} are \rulename{elim} and \rulename{occurs} for a
variable with no arguments, a case that micro's first-order unifier meets first, and they are where the
eager scope check, \textbf{Check}, enters; \rulename{FF-one} orients a flexible–flexible pair as
\rulename{swap} orients a flexible–rigid one.

\begin{figure*}[!t]
\begin{mathpar}
\inferrule{X \doteq t}{X := t\ \ \text{(micro, occurs check)}\quad\textbf{Check}(X := t)}\;\rulename{Var}
\and
\inferrule{\lambda.b \doteq \lambda.b'}{b[c/0] \doteq b'[c/0]\quad \textbf{Check}\ \text{its bindings}}\;\rulename{ξ}
\and
\inferrule{\lambda.b \doteq t}{b[c/0] \doteq t\;@\;c\quad \textbf{Check}\ \text{its bindings}}\;\rulename{η}
\\
\inferrule{h\,\bar u \doteq h'\,\bar v}{h \doteq h' \quad \bar u \doteq \bar v}\;\rulename{Rigid}
\and
\inferrule{F\,\bar a \doteq F\,\bar b}{F := \lambda\bar x.\,H\,\langle x_i \mid a_i = b_i\rangle}\;\rulename{FF-same}
\\
\inferrule{F\,\bar a \doteq t}{F := \lambda\bar x.\,t'\qquad t' = \mathrm{Elim}(\ell(F),F,\bar a,t)}\;\rulename{FR}
\and
\inferrule{F\,\bar a \doteq G\,\bar b}{G\,\bar b \doteq F\,\bar a\ \ \text{by \rulename{FR}}}\;\rulename{FF-one}
\\
\inferrule{u \doteq v}{E := E \cup \{\lambda L.\,u \doteq \lambda L.\,v\}}\;\rulename{Delay}
\and
\inferrule{X \doteq t \\ X \text{ occurs in } t}{\top\ \text{if } \mathrm{hnf}(t)=X;\ \text{else } X := \mathrm{Elim}(\ell(X),X,\langle\rangle,t)}\;\rulename{Occurs}
\end{mathpar}
\caption{The unification rules of \texttt{μkanren-λ-unify}, applied after hnf of both sides; above the line
the equation, below what is done with it. In \rulename{ξ} and \rulename{η}, $c$ is a new eigenvariable of level
$\infty$. In \rulename{FF-same} and \rulename{FR}, $\bar x$ are new binders, one per argument $a_i$, and Elim
replaces each $a_i$ in $t$ by the bound variable of $x_i$. Variable elimination, Elim, is
\cref{def:module-elim}; \textbf{Check} is \cref{def:module-check}. Each rule may also answer $\bot$ or suspend,
as its side conditions say.}\label{fig:module-rules}
\end{figure*}

The side conditions, in the order the code tests them (\scm{μkanren-λ-unify} and the functions it calls):
\begin{description}[style=nextline,leftmargin=5.6em]
  \item[\rulename{Var}] One side is an unbound variable. micro binds it, with its occurs check; the scope check
    \textbf{Check} runs afterwards, in the \texttt{post} hook. Variable–variable pairs are micro's as well.
  \item[\rulename{ξ}] Both sides are abstractions. $c$ is a new eigenvariable of level $\infty$, so no variable
    can capture it except through its own pattern arguments. Because $c$ is local to this very unification,
    the bindings made under it are checked at once (\scm{μkanren-scoped}), not at the end of \scm{=°}.
  \item[\rulename{η}] One side is an abstraction and the other is not. The other side is applied to the new
    $c$. This is the only η-expansion the module performs.
  \item[\rulename{Rigid}] Both sides are rigid applications with the same number of arguments: heads and
    argument lists are unified by micro. Two eigenvariables unify iff they are the same (\scm{eq?} ids), two
    bound variables iff their indices agree. Any other pair with a λ-term on one side fails: different rigid
    heads, different arities (untyped), a λ-record against a pair or an atom.
  \item[\rulename{FF-same}] Same head on both sides, both argument lists patterns. Equal lists: success with
    no binding. Different lengths: $\bot$. Otherwise $F$ keeps the positions where the arguments agree, with a
    new $H$ of level $\ell(F)$ \book{222}. If $\bar b$ is no pattern, \rulename{Delay}.
  \item[\rulename{FR}] One side is a pattern $F\,\bar a$; the other side $t$ is rigid, or flexible with another
    head (\enquote{FF-diff}: variable elimination prunes $G$ as needed \book{222}). If Elim fails the rule
    fails; if Elim suspends, \rulename{Delay}. The binding $\lambda\bar x.\,t'$ is built η-short, its binders
    named after the eigenvariables $\bar a$.
  \item[\rulename{FF-one}] Two flexible terms, only the right one a pattern: \rulename{FR} on $G$.
  \item[\rulename{Delay}] No side is a usable pattern. The equation is pushed on $E$, closed over the names $L$
    that the unifier had opened so far ($\lambda L.u \doteq \lambda L.v$, so no $\infty$-level name escapes the
    unification that created it), \emph{unless} $u$ and $v$ are already α-equal, in which case it is dropped.
  \item[\rulename{Occurs}] Not a rule of \scm{μkanren-λ-unify} but of the protocol's \texttt{occurs} hook: micro's
    occurs check failed on $X := t$ and $t$ holds a λ-term. If $\mathrm{hnf}(t)$ is $X$ itself the equation
    holds ($q = F\,q$ with $F := \lambda z.z$); otherwise $X$ is eliminated with no arguments, which may fail,
    succeed, or suspend ($F = \lambda x.\,G\,(F\,x)$ holds with $G := \lambda z.z$). A first-order cycle,
    $X = f\,X$, holds no λ-term and fails in micro as before.
\end{description}

\subsection{Variable elimination}

\rulename{FR} is the book's variable elimination \book{223} (\cref{def:hounif-elim}) with levels in place of the
prefix. It walks the right-hand side once and builds the body of $F$'s binding, pruning, lowering and raising the
variables it meets on the way. Compared with \cref{def:hounif-elim}, three things are added: \emph{lowering}
with delayed raising (case~(4)), where \cref{def:hounif-elim} assumes a prefix already raised into
$\forall\exists$ form (the eager raising of \cref{ex:hounif-8k}); the treatment of \emph{non-patterns} inside the right-hand side (cases (6c) and (7)), which the book never meets since
it stays in $\Llam$; and a finer \emph{occurs} test (case (6a)), since an occurrence of $F$ is fatal only when no
instance can make it smaller. The contexts below exist to decide these three cases.

\begin{definition}[Variable elimination]\label{def:module-elim}
$\mathrm{Elim}(\ell, F, \bar a, t)$, with $\bar a$ a pattern for $F$ and $\ell=\ell(F)$ (or $F$ absent and
$\bar a$ empty, for \textbf{Check}), walks $\mathrm{hnf}(t)$ at a local depth $k$ (the binders crossed inside
$t$) and in a \emph{context}: \emph{strict} (\texttt{top}: below binders only; \texttt{rigid}: below a first rigid
constructor that is an application headed by a bound variable or by an $a_i$; \texttt{firm}: below a first rigid
constructor that is a constant, an eigenvariable not in $\bar a$, a pair, a vector or a record, a root no instance
of $F$ can project onto) or \emph{soft} (inside an argument of a flexible term that is no pattern, or
\texttt{weak}: below a flexible head inside such an argument). It answers a term $t'$ with new bindings,
$\bot$, or \emph{suspend}. By cases on the walked subterm:
\begin{enumerate}[label=(\arabic*)]
  \item \emph{Eigenvariable $c$}: if $c = a_i$, the bound variable of $a_i$; else if $\ell(c)\le\ell$, $c$
    itself; else $\bot$ (scope failure, \cref{ex:hounif-8j}). In a weak context the last case is trouble, not
    $\bot$.
  \item \emph{Bound variable}: itself.
  \item \emph{Abstraction} $\lambda.b$: $\lambda.\mathrm{Elim}(b)$ at depth $k+1$.
  \item \emph{Unbound variable $Y$}: if $Y=F$, $\bot$ when $\bar a$ is empty (the first-order occurs check) and
    suspend otherwise; if $\ell(Y)\le\ell$, $Y$; in a soft context, trouble; otherwise \emph{lower and raise}:
    $Y := Y'\,\bar a_Y$ with $\ell(Y')=\ell$ and $\bar a_Y = \langle a_i \mid \ell(a_i)\le\ell(Y)\rangle$, the
    arguments of $F$ that $Y$ could see; the result is $Y'$ applied to their bound variables
    (\enquote{delayed raising}, \cref{ex:hounif-8k}).
  \item \emph{Rigid application, pair, vector, record}: homomorphic. A \texttt{top} context becomes
    \texttt{rigid} below an application whose head is a bound variable or in $\bar a$, and \texttt{firm} below
    any other head, a pair, a vector or a record; any other context is passed on unchanged.
  \item \emph{Flexible $G\,\bar b$ in a strict context}:
    \begin{enumerate}[label=(\alph*)]
      \item if $G=F$: $\bot$ (occurs failure) when the context is not \texttt{top}, no $b_j$ is
        \emph{functional} (an abstraction, a variable or a flexible term, which may become an abstraction), and
        either the context is \texttt{firm} or $|\bar b| = |\bar a| + k_{\mathrm{top}}$, with $k_{\mathrm{top}}$
        the binders on top of $t$; otherwise suspend;
      \item if $\bar b$ is a pattern for $G$ (each $b_j$ a local bound variable or an eigenvariable of level
        $>\ell(G)$, pairwise distinct): keep $b_j$ iff it is local, or in $\bar a$, or $\ell(b_j)\le\ell$; let
        $\mathit{extra}=\langle a_i \mid \ell(a_i)\le\ell(G)\rangle$ if $\ell(G)>\ell$, else empty. If every
        $b_j$ is kept, $\mathit{extra}$ is empty and $\ell(G)\le\ell$, $G\,\bar b$ is unchanged. Otherwise
        \emph{prune}: $G := \lambda\bar y.\,G'\,\bar y|_{\mathit{keep}}\,\mathit{extra}$ with
        $\ell(G')=\min(\ell,\ell(G))$, and the result is $G'$ applied to the images of the kept arguments and of
        $\mathit{extra}$ (\cref{ex:hounif-8h});
      \item if $\bar b$ is no pattern: each $b_j$ is walked on its own in a soft context, binding nothing; it is
        \emph{kept}, as the walk rebuilt it (its $a_i$ turned into bound variables), if the walk succeeds, \emph{dropped} if it fails by scope on a rigid path (it holds a
        forbidden eigenvariable that every instance keeps), and any other trouble suspends the whole equation.
        Then, if nothing was dropped and $G$ needs no lowering, $G\,\bar b$ stays as it is; otherwise, if some
        $b_j$ is functional, suspend; else prune and lower as in (b).
    \end{enumerate}
  \item \emph{Flexible $G\,\bar b$ in a soft context}: $G=F$ or $\ell(G)>\ell$ is trouble; otherwise the
    arguments are walked in a weak context and kept as they are. Inside an argument, the application of a
    variable bound inside that argument is treated as flexible too.
\end{enumerate}
Inside an argument of a non-pattern, case~(6c) reinterprets every failure: a scope failure of case~(1) in
a soft, non-weak context drops the argument, and any other $\bot$ or \emph{trouble} suspends the equation.
\end{definition}

\begin{tutorial}[Reading the contexts]
The context answers one question about the subterm being walked: \emph{will every instance of the right-hand
side still contain it?} At \texttt{top} and below rigid constructors (\texttt{rigid}, \texttt{firm}) the answer is
yes, because no substitution can remove a rigid constructor; so a forbidden eigenvariable there is a definite
failure, and a variable there can be lowered or pruned without losing solutions. Inside an argument of a
non-pattern $G\,\bar b$ the answer is \enquote{not necessarily}, since $G$ may become a function that ignores
that argument; so nothing is bound there and doubts turn into a suspension. \texttt{firm} and \texttt{rigid}
differ only in the occurs test (6a): below a constant, an eigenvariable outside $\bar a$ or a pair, $F$'s instance
cannot be a projection, so an inner $F\,\bar b$ with rigid arguments is fatal whatever their number; below a
bound variable or an $a_i$ at the head, the number matters. The next two tutorials give the reasons in detail.
\end{tutorial}

\begin{tutorial}[Strict and non-strict occurrences]
Pruning removes an argument from a variable; lowering makes a variable forget the eigenvariables it could see.
Both are sound only on \emph{strict} occurrences, those not inside an argument of a flexible term that is no
pattern (Dowek, Hardin, Kirchner and Pfenning \cite{DowekHardinKirchnerPfenning1996}). The reason is that
the flexible head $G$ may later turn out to drop that argument. Take
$\exists q\,\exists G\,\forall y\,\exists Y.\ q = G\,Y \wedge Y = y \wedge G = \lambda a.\,k$: it has the solution
$q=k$. If elimination of $q$ lowered $Y$ inside the argument of $G\,Y$, the later $Y=y$ would fail by scope
and the answer would be lost in that goal order, which breaks the commutativity of conjunction. So nothing is
bound inside such an argument: a variable to lower (case~(4)) or a flexible term to prune (case~(7)) is trouble
there, and trouble suspends the equation. The exception of case
(6c) is the one pruning that \emph{is} strict: an argument that holds, on a rigid path, an eigenvariable that $F$
may not see is held by every instance of that argument, so $G$ must drop it. This makes
$\exists q\,\exists F\,\forall x.\ q = F\,x\,x$ solvable as $F := \lambda a\,b.\,F'$, $q := F'$. Even that
pruning waits when an argument is, or may become, an abstraction: $G$ could apply it to another argument and
drop that one (\cref{sec:module-traces} shows $\exists q\,\forall x\,\exists G.\ q = G\,(\lambda w.\,k)\wedge G =
\lambda y.\,y\,x$, solved in both orders).
\end{tutorial}

\begin{tutorial}[When an occurrence of $F$ is fatal]
Case (6a) is Nipkow's size argument \cite{Nipkow1993}: $F\,\bar b$ inside the body of $F\,\bar a$ makes the
equation unsolvable only if no instance of $F\,\bar b$ can be smaller than the instance of $F\,\bar a$. That
holds when the $b_j$ are rigid (no substitution makes them reduce) and are as many as $\bar a$ plus the
binders on top ($F\,\bar a = \lambda\bar y.\,t'$ is $F\,\bar a\,\bar y = t'$ by η), as in
$F = \lambda x.\,g\,(F\,x)$ and \cref{ex:hounif-8i}. The count does not matter when the root of the body is a
constant, an eigenvariable not in $\bar a$ or a pair (context \texttt{firm}): no instance $\lambda\bar x.M$ of
$F$ can then be a projection, so $M$ has that root and $F\,\bar b$ sits strictly inside it. So the η-short
$F = \lambda x\,y.\,g\,(F\,y)$ fails as Miller and Nipkow fail it. Otherwise the instance may shrink and the
equation is outside $\Llam$: it is suspended. For instance $\forall x.\ F\,x = x\,(F\,(\lambda z.z))$ holds
with $F := \lambda a.\,a\,c$, and $\forall x\,\forall y.\ F\,x\,y = x\,(F\,y)$ holds with $F := \lambda w.w$.
\end{tutorial}

\subsection{The eager scope check}

\begin{definition}[Check]\label{def:module-check}
For a binding $X := t$ that micro has just made, $\textbf{Check}(X:=t)$ is $\mathrm{Elim}(\ell(X), -,
\langle\rangle, t)$ with no head to eliminate (\scm{μkanren-enforce}). If it answers $\bot$, \scm{=°} fails.
If it suspends, the binding is \emph{taken back}: \scm{μkanren-state-forget} unbinds $X$ and the equation
$X \doteq t$ is pushed on $E$. Otherwise its own bindings (lowerings, prunings) are kept.
\end{definition}
\textbf{Check} is NJK's binding rule with its lowering (\cref{def:rules-binding}): \enquote{$t$ contains no
constant with tag $>i$; change to $i$ the tags of the variables in $t$ whose tag is $>i$}
\cite[p.\,16]{NadathurJayaramanKwon1995}, extended from variables to flexible terms by pruning. Taking a binding
back is what Teyjus does with a binding whose check depends on a non-pattern. Its effect is an invariant:
\begin{proposition}[No captured eigenvariable]\label{prop:module-nocapture}
After every \scm{=°}, no binding $X := t$ in the substitution has an eigenvariable of level $>\ell(X)$ in the
value of $t$.
\end{proposition}
{\emergencystretch=2em
\noindent\emph{Argument.} Every binding made by micro goes through \textbf{Check} unless the fast paths of
\cref{sec:module-fast} prove it unnecessary; \textbf{Check} either fails, or rewrites the variables of $t$ so
that it holds, or unbinds $X$. The bindings that \textbf{Check} itself makes are lowerings and prunings,
which by construction respect the levels. The remaining gap is host code (\scm{project°} and \scm{set!}),
discussed in \cref{ch:roadmap}. \hfill\qedsymbol\par}

\subsection{Fast paths}\label{sec:module-fast}

\textbf{Check} on every binding would be costly on first-order search, so three cheap tests come first.
\begin{itemize}
  \item \textbf{No \scm{∀°} yet.} While no \scm{∀°} has run, the depth is 0 and $L$ is empty
    (\scm{μkanren-levels?} is false): every level is 0, no check can fail, none runs.
  \item \textbf{Plain terms.} \scm{μkanren-plain?} is an allocation-free walk, like micro's occurs check. It
    accepts $t$ when $t$ holds no eigenvariable or variable of a level above $\ell(X)$ and no application with a
    variable or abstraction head (which Elim might prune, lower or reduce). Only otherwise does Elim run.
  \item \textbf{Top-level variables.} A variable of level $\mathit{top}$, the highest level reached, can see
    every eigenvariable and every variable there is, so its bindings are skipped. The bindings made under
    \rulename{ξ} or \rulename{η} are checked anyway, by \scm{μkanren-scoped}.
\end{itemize}
The protocol's \texttt{atom?} hook belongs here as well: it tells micro's occurs check that an eigenvariable or
a bound variable holds no variable, so micro does not rebuild it as a slot list. On the type inference of the
Church numeral 400 these steps took the time from 463~ms to 330~ms (\scm{plain?}), 311~ms (the \texttt{top}
skip) and 254~ms (\texttt{atom?}).

\begin{pitfall}[The first \texttt{∀°} switches the check on for good]
The first fast path is about the whole branch, not the current scope: once a \scm{∀°} has run, levels exist
and every binding below \texttt{top} is checked, even after the \scm{∀°} has exited. First-order search that
follows a \scm{∀°} pays for it, measured at $+66\%$ on \scm{append°} backwards and $+127\%$ on open
\scm{append°}. Skipping the check at depth 0 would be unsound against host escapes; the fix planned is one walk
for micro's occurs check and the scope check together (\cref{ch:roadmap}).
\end{pitfall}

\subsection{Wake-up and termination}\label{sec:module-wake}

When \scm{=°} binds something and $E$ is not empty, \scm{μkanren-settle} retries every suspended item, oldest
first (\scm{μkanren-wake}). An item is retried by its kind:
\begin{itemize}
  \item an equation $u\doteq v$: if one side is now an unbound variable $X$ and the other no variable, $X$ is
    bound \emph{only once its check passes} (\scm{μkanren-retry/bind}); otherwise the equation is unified
    again, which may solve it, fail, or suspend it again;
  \item a type on a flexible term: decided if the head is now known, else suspended again;
  \item an absento: decided, kept, or dropped when levels entail it (\cref{sec:module-stores}).
\end{itemize}
The new bindings of a round are checked like those of \scm{=°}, and rounds repeat while a round binds. A
constraint still undecided is suspended again \emph{without} binding anything.

\begin{proposition}[Wake-up terminates]\label{prop:module-wake}
Suppose every hnf computed during the wake-up terminates. Then the loop of \scm{μkanren-settle} terminates.
\end{proposition}
\noindent\emph{Argument} (a sketch, not a mechanised proof). A round that binds nothing ends the loop, so it
suffices that only finitely many rounds bind. A round that binds does so for good: the only binding ever taken back is the one
that \textbf{Check} has just examined, and a retried $X \doteq t$ binds $X$ only after its check passed, so the
retry does not bind and unbind the same variable forever. The variables a round may create come from lowering
(case~(4)), whose new variable has a strictly smaller level than the one it replaces, and from pruning
(cases~(6b) and~(6c)), whose new variable either has a strictly smaller level (when $G$ is lowered, possibly with
\emph{extra} arguments) or the same level and strictly fewer arguments (\rulename{FF-same} likewise creates
$H$ at the level of $F$ with fewer arguments). So every new variable \emph{replaces} exactly one variable that
is bound for good, and is below it in the lexicographic order on (level, number of arguments), which is well
founded. A chain of replacements (a variable, the one that replaces it, the one that replaces that, \dots) is
therefore finite; with finitely many variables at the start, finitely many variables are ever created, each is
bound at most once, and only finitely many rounds can bind. This is the decreasing measure the design relies on.
One point a full proof would have to settle: the number of arguments is that of the occurrence being pruned,
which an untyped term does not fix once and for all. \hfill\qedsymbol

The hypothesis on hnf is not a formality: \scm{@} gives full β, and on untyped meta-terms full β may diverge
(\cref{ch:roadmap}). On pattern problems no such reduction happens (every contraction is a $\beta_0$ renaming),
and the rules are Miller's and Nipkow's algorithm with the prefix erased into levels: it terminates, fails only
when there is no unifier, and otherwise computes a most general one \cite[Thm~7.1, Prop.~7.3]{Miller1991},
\cite[Thm~2]{NadathurLinnell2005}. The suite checks the book's mgus verbatim: $U = \lambda x\,y.\,g\,(V'\,y)$ and
$V = \lambda y\,w.\,V'\,y$ for \cref{ex:hounif-8h}, $U = \lambda x.\,f\,(V'\,x)$ for \cref{ex:hounif-8k}, and
$G, H := \lambda w.\,T$ for $\lambda x\,\lambda y.\,G\,x = \lambda x\,\lambda y.\,H\,y$ (\cref{ex:hounif-ch4}).

\begin{pitfall}[Suspended is not solved]
Non-patterns are suspended, never enumerated: there is no Huet search \book{213}. An answer may therefore be
conditional on residual equations that are jointly unsatisfiable. $F\,a = b$ and $F\,a = c$ ($a$ a constant,
so no pattern) are both suspended, and the answer is printed with the two residuals; only a binding of $F$
reveals the conflict:
\begin{scheme}
(°->list #f (fresh° (q F) (=° q F) (=° (@ F 'a) 'b) (=° (@ F 'a) 'c)))
(°->list/datum (fresh° (q F) (=° q F) (=° (@ F 'a) 'b) (=° (@ F 'a) 'c)
                        (=° F (ƛ (z) 'b))))
\end{scheme}
\begin{answers}
((λ (α) (assert (≐ (@ α 'a) 'b)) (assert (≐ (@ α 'a) 'c)) α))
()
\end{answers}
This is floundering reported, not detected. The residuals are the honest output: they say exactly which
equations the search did not decide.
\end{pitfall}

\subsection{Worked traces}\label{sec:module-traces}

Each trace below poses a problem, follows the rules, and ends with the module's actual output. The probe used
here prints the depth, the levels of the given variables and $E$:
\begin{scheme}
(define ((probe label . vars) s)
  (display label) (display ": depth ") (display (μkanren-λ-depth s))
  (display ", levels ") (write (map (λ (v) (μkanren-λ-level v s)) vars))
  (display ", E ") (write (μkanren-λ-delayed s)) (newline)
  (✓° s))
\end{scheme}

\begin{example}[Warm-up: flexible–rigid, and \rulename{ξ}]\label{ex:module-warmup}
\emergencystretch=2em
Two equations that touch no pruning or lowering: $\exists q\,\forall x\,\forall y.\ q\,x\,y = g\,y\,x$, and
$\exists q.\ \lambda x.\,(g\ x) = \lambda y.\,(g\ (q\,y))$, where both $(g\ \dots)$ are Scheme lists.
\begin{scheme}
(°->list/datum (fresh° (q) (∀° (x y) (=° (@ q x y) `(g ,y ,x)))))
(°->list/datum (fresh° (q) (=° (ƛ (x) `(g ,x)) (ƛ (y) `(g ,(@ q y))))))
\end{scheme}
\begin{answers}
((ƛ (x y) (g y x)))
((ƛ (x) x))
\end{answers}
\emph{Trace of the first.} $\ell(q)=0$, $\ell(x)=1$, $\ell(y)=2$. The left side is flexible and $\langle x,y\rangle$
is a pattern for $q$; the right side is the list \scm{(g y x)}, rigid: \rulename{FR}. Elim walks the list in
context \texttt{firm}: \scm{g} is an atom and stays; $y$ is $a_2$ and becomes the bound variable of the second
binder; $x$ is $a_1$ and becomes that of the first (case~(1)). Nothing was bound on the way, and
$q := \lambda x\,y.\,(g\ y\ x)$, the binders named after the eigenvariables.

\emph{Trace of the second.} Both sides are abstractions: \rulename{ξ} opens them with one new eigenvariable $c$
of level $\infty$ (hint \texttt{x}), leaving \scm{(g c)} $\doteq$ \scm{(g (@ q c))}. These are lists, so micro
decomposes them and hands $c \doteq q\,c$ back to the hook. $\langle c\rangle$ is a pattern for $q$
($\infty > 0$), and Elim turns $c$ into the bound variable: $q := \lambda x.\,x$, a projection. The binding
was made under \rulename{ξ}, so \scm{μkanren-scoped} checks it at once; it holds no eigenvariable and passes.
\end{example}

\begin{example}[Pruning: problem 8.H \book{223--224}, \cref{ex:hounif-8h}]\label{ex:module-8h}
The problem (\cref{ex:hounif-8h}) is
\[
  \exists U\,\exists V\,\forall w\,\forall x\,\forall y.\ f\,(U\,x\,y) = f\,(g\,(V\,y\,w)).
\]
\begin{scheme}
(°->list/datum (fresh° (q u v) (=° q (list u v))
  (∀° (w x y) (probe "before" u v)
              (=° `(f ,(@ u x y)) `(f (g ,(@ v y w))))
              (probe "after" u v))))
\end{scheme}
\begin{answers}
before: depth 3, levels (0 0), E ()
after: depth 3, levels (0 0), E ()
(((ƛ (x y) (g (@ α y))) (ƛ (y w) (@ α y))))
\end{answers}
\emph{Trace.} $U$, $V$ have level 0 and $w,x,y$ levels 1, 2, 3. Both sides are lists, so micro decomposes
them: \scm{f} against \scm{f}, then $U\,x\,y \doteq g\,(V\,y\,w)$ with a λ-record on the left, handed to the
hook. The left side is a pattern ($x$, $y$ distinct, levels $2,3>0$) and the right is rigid: \rulename{FR},
$\mathrm{Elim}(0, U, \langle x,y\rangle, g\,(V\,y\,w))$. The root is the list \scm{(g ...)}: case (5),
context \texttt{firm}. Inside, $V\,y\,w$ is flexible with a pattern for $V$ (levels $3,1>0$): case~(6b). Keep $y$
(it is in $\bar a$, it becomes the bound variable of $y$); drop $w$ (not local, not in $\bar a$, and
$\ell(w)=1>0$); $\mathit{extra}$ is empty since $\ell(V)=0$ is not above $\ell=0$. Something was dropped, so
prune: $V := \lambda y\,w.\,V'\,y$, $\ell(V')=0$, and the body of $U$ is $g\,(V'\,y)$. Finally
$U := \lambda x\,y.\,g\,(V'\,y)$. Both are the book's mgu, with \texttt{α} for $V'$. The binders of $V$ are
named after the eigenvariables $y$, $w$ it was applied to.
\end{example}

\begin{example}[Flexible–flexible]\label{ex:module-ff}
Different heads, $\forall x\,\forall y.\ F\,x\,y = G\,y\,x$, and the same head, $\forall x\,\forall y.\ F\,x\,y =
F\,y\,x$:
\begin{scheme}
(°->list #f (fresh° (q f g) (=° q (list f g)) (∀° (x y) (=° (@ f x y) (@ g y x)))))
(°->list #f (fresh° (q) (∀° (x y) (=° (@ q x y) (@ q y x)))))
\end{scheme}
\begin{answers}
((λ (α) (cons (ƛ (x y) (@ α y x)) (cons α '()))))
((λ (α) (ƛ (x y) α)))
\end{answers}
\emph{Trace.} In the first, both sides are patterns with different heads: \rulename{FR} on $F$. In Elim, $G\,y\,x$
is case (6b) with both arguments in $\bar a$, kept; $\ell(G)=0$ is not above $\ell(F)=0$, so nothing is
pruned or lowered and $G$ stays unbound: $F := \lambda x\,y.\,G\,y\,x$. That is the book's mgu for different
heads \book{222} with the common variables listed as $y,x$: $F := \lambda x\,\lambda y.\,H\,y\,x$ and
$G := \lambda y\,\lambda x.\,H\,y\,x$, which is $H$ by η, up to the name of $H$. In the second, \rulename{FF-same}: the
positions where $x\,y$ and $y\,x$ agree are none, so $F := \lambda x\,y.\,H$ with $H$ new.
\end{example}

\begin{example}[A non-pattern, suspended and woken]\label{ex:module-nonpattern}
$\exists F\,\forall x.\ F\,(g\,x) = h\,(g\,x)$ is not in $\Llam$: the argument $g\,x$ is not an eigenvariable. It
has two solutions, $F := \lambda z.\,h\,z$ and $F := \lambda z.\,h\,(g\,x)$, but the second is illegal here
because $F$ is quantified outside $x$. The module neither guesses nor fails:
\begin{scheme}
(°->list #f (fresh° (q) (∀° (x) (=° (@ q `(g ,x)) `(h (g ,x))) (probe "inside" q))
                        (probe "after ∀" q)))
(°->list/datum (fresh° (q) (∀° (x) (=° (@ q `(g ,x)) `(h (g ,x))))
                           (probe "before F :=" q)
                           (=° q (ƛ (z) `(h ,z)))
                           (probe "after F :=" q)))
(°->list/datum (fresh° (q) (∀° (x) (=° (@ q `(g ,x)) `(h (g ,x))))
                           (=° q (ƛ (z) 'k))))
(°->list/datum (fresh° (q) (=° q (ƛ (z) `(h ,z)))
                           (∀° (x) (=° (@ q `(g ,x)) `(h (g ,x))))))
\end{scheme}
\begin{answers}
inside: depth 1, levels (0), E ((= (@ _0 (g (eigen x))) (h (g (eigen x)))))
after ∀: depth 0, levels (0), E ((= (ƛ (x) (@ _0 (g x))) (ƛ (x) (h (g x)))))
((λ (α)
    (assert
      (≐ (ƛ (x) (@ α (cons 'g (cons x '()))))
         (ƛ (x) (cons 'h (cons (cons 'g (cons x '())) '())))))
    α))
before F :=: depth 0, levels (0), E ((= (ƛ (x) (@ _0 (g x))) (ƛ (x) (h (g x)))))
after F :=: depth 0, levels (0), E ()
((ƛ (z) (h z)))
()
((ƛ (z) (h z)))
\end{answers}
\emph{Trace.} Inside the scope the left side is flexible but no pattern and the right is rigid:
\rulename{Delay}, and $E$ holds the equation, with the eigenvariable $x$ in it. At the exit of \scm{∀° (x)}
the equation mentions $x$, so it is \emph{discharged} (\cref{sec:module-discharge}): no variable in it has a
level $\ge\ell(x)=1$, so nothing is raised, and the equation is closed into
$\lambda x.\,F\,(g\,x) \doteq \lambda x.\,h\,(g\,x)$, equivalent by \rulename{ξ}. It is retried at once, opens
$x$ again as an $\infty$-level name, is still no pattern, and is suspended again. The run ends with it, and
the answer carries it as a residual \scm{(assert (≐ ...))}. In the second run the binding $F := \lambda z.\,h\,z$
wakes it: \rulename{ξ}, then β turns the left side into $h\,(g\,c)$, which equals the right side, and $E$ is
empty. In the third run $F := \lambda z.\,k$ makes the retry compare $k$ with $h\,(g\,c)$: a clash, so \scm{=°}
fails and there is no answer. In the fourth run the binding comes first and the equation is a plain β-redex,
solved with no suspension. The answer is the same in both goal orders.
\end{example}

\begin{example}[A binding taken back]\label{ex:module-forget}
$\exists q\,\exists G\,\exists H\,\forall y.\ q = G\,(H\,y)$. micro binds $q := G\,(H\,y)$ at once, as
\rulename{Var}. \textbf{Check} at level $\ell(q)=0$ then meets $G\,(H\,y)$, a flexible term whose argument
$H\,y$ is no eigenvariable: case (6c), walking $H\,y$ in a soft context. There $y$ (level 1, not allowed for a
level-0 variable) sits below the flexible head $H$, a weak context: trouble, not $\bot$, since $H$ may drop
$y$, and $G$ may drop all of $H\,y$. So the check suspends, the binding is taken back, and $q \doteq G\,(H\,y)$
waits in $E$:
\begin{scheme}
(°->list #f (fresh° (q G H) (∀° (y) (=° q (@ G (@ H y))) (probe "inside" q G H))))
(°->list/datum (fresh° (q G H) (∀° (y) (=° q (@ G (@ H y)))) (=° G (ƛ (a) 'k))))
\end{scheme}
\begin{answers}
inside: depth 1, levels (0 0 0), E ((= _0 (@ _1 (@ _2 (eigen y)))))
((λ (α β γ) (assert (≐ (ƛ (y) α) (ƛ (y) (@ β (@ γ y))))) α))
(k)
\end{answers}
At the exit of the scope the equation is closed over $y$ and reported, $q$ unbound, which is all one can say.
With $G := \lambda a.\,k$ the retry binds $q := k$, a plain term, and the residual is gone. The same happens in
every order of the goals.
\end{example}

\begin{example}[Occurs and scope failures, and a soft occurrence]\label{ex:module-occurs}\leavevmode
\begin{scheme}
(°->list #f (fresh° (q) (=° q `(f ,q))))                              ; X = f X
(°->list #f (fresh° (q) (∀° (w x y) (=° (@ q x y) `(g ,(@ q y w))))))  ; 8.I
(°->list #f (fresh° (q) (∀° (w x y) (=° (@ q x y) `(g ,w)))))          ; 8.J
(°->list #f (fresh° (F) (=° F (ƛ (x) `(g ,(@ F x))))))         ; F = λx.g (F x)
(°->list #f (fresh° (q F G) (=° q (list F G))
              (∀° (x) (=° (@ F x) (@ G (@ F x))))))
\end{scheme}
\begin{answers}
()
()
()
()
((λ (α β) (assert (≐ α (ƛ (x) (@ β (@ α x))))) (cons α (cons β '()))))
\end{answers}
The first is micro's occurs check, unchanged. In 8.I, $U\,y\,w$ occurs below the list \scm{(g ...)}, a
\texttt{firm} context, and its arguments are rigid: case (6a), $\bot$ (\cref{ex:hounif-8i}); it also has as
many arguments as $U\,x\,y$, so the count alone would fail it. In 8.J, $w$ has level~1, is not
an argument of $U$, and $U$ has level~0: case (1), $\bot$ by scope (\cref{ex:hounif-8j}). The fourth fails
micro's occurs check, the \texttt{occurs} hook passes it to Elim with no arguments, and $F\,x$ below \scm{g}
is case (6a) again. The last is the \emph{soft} occurs check: $F\,x$ occurs inside the argument of $G$, which may
be $\lambda z.z$, so the equation is suspended (and is closed over $x$ at the exit of the scope) instead of
failing.
\end{example}

\begin{example}[An argument that may be an abstraction]\label{ex:module-functional}
$\exists q\,\forall x\,\exists G.\ q = G\,(\lambda w.\,k) \wedge G = \lambda y.\,y\,x$ in both orders:
\begin{scheme}
(°->list/datum (fresh° (q) (∀° (x) (fresh° (G) (=° q (@ G (ƛ (w) 'k)))
                                               (=° G (ƛ (y) (@ y x)))))))
(°->list/datum (fresh° (q) (∀° (x) (fresh° (G) (=° G (ƛ (y) (@ y x)))
                                               (=° q (@ G (ƛ (w) 'k)))))))
\end{scheme}
\begin{answers}
(k)
(k)
\end{answers}
In the first order, \textbf{Check} of $q := G\,(\lambda w.k)$ at level 0 finds $G$ of level 1, which must be
lowered; but its argument is an abstraction that $G$ may apply to $x$ and so drop $G$'s own $x$. Lowering $G$
would not be strict, so case (6c) suspends. Then $G := \lambda y.\,y\,x$ wakes the equation, β gives
$(\lambda w.k)\,x = k$, and $q := k$. In the other order β happens at once. A version that lowered $G$ here
lost the answer in the first order.\sidenote{Third adversarial review.}
\end{example}

% ------------------------------------------------------------------------------------------------
\section{\texorpdfstring{\texttt{∀°}}{∀°} and \texorpdfstring{\texttt{⊃°}}{⊃°}}\label{sec:module-hh}

The book's goal-directed search for hereditary Harrop formulas consists of the rules \rulename{and}, \rulename{or}, \rulename{instan},
\rulename{generic}, \rulename{augment} and \rulename{backchain} \book{77}, \book{82}, \book{124} (\cref{fig:hh-rules}; the names are those of \cref{fig:rules-right}). Uniform proofs are
complete for intuitionistic provability of these formulas \cite{MillerNadathurPfenningScedrov1991}
(\cref{thm:hh-intuitionistic}), so implementing the rules faithfully \emph{is} the semantics. \rulename{and}, \rulename{or} and
$\exists$ are micro's \scm{and°}, \scm{or°} (interleaving) and \scm{fresh°}. This section gives the other
three as the module implements them.

\subsection{\texorpdfstring{\rulename{generic}}{GENERIC}}

\begin{mathpar}
\inferrule{\seq{\Sigma\cup\{c\}}{\Prog}{G[c/x]}}{\seq{\Sigma}{\Prog}{\forall x.\,G}}\;\rulename{generic}\quad(c\notin\Sigma)
\end{mathpar}
\scm{(∀° (x) g ...)} expands into \scm{(∀ª 'x (λ (x) (and° g ...)))}, and \scm{(∀° (x y) g ...)} into one
\scm{∀ª} per name, nested, so \scm{(∀° (w x y) ...)} gives $w,x,y$ the levels $d+1$, $d+2$, $d+3$. The Scheme
variable \scm{x} is bound to the eigenvariable, as \scm{ƛ} binds its placeholder. \scm{∀ª} takes a state $s$
whose λ record has depth $d$ and does the following.
\begin{enumerate}
  \item Record the level $d$ of every variable created so far (pad $L$ up to the variable count).
  \item Run the body \scm{(f c)} for a new eigenvariable $c$ of level $d+1$, in the state with depth $d+1$ and
    $\mathit{top} := \max(\mathit{top}, d+1)$.
  \item For each answer state $s_1$ of the body: record the level $d+1$ of the variables created inside (pad
    $L$ again), set the depth back to $d$, and \emph{discharge} the constraints on $c$
    (\cref{sec:module-discharge}); if discharge fails, drop the answer.
\end{enumerate}
\begin{scheme}
(define ((∀ª name f) s)
  (let* ((x  (μkanren-hh s))
         (d  (μkanren-hh-depth x))
         (x* (μkanren-hh-with (μkanren-hh-pad x (μkanren-state-vars-count s))
                              depth: (add1 d) top: (max (add1 d) (μkanren-hh-top x))))
         (c  (μkanren-eigen* name (add1 d))))
    (append-map§
      (λ (s1)
        (let1 (s2 (μkanren-with s1
                    (μkanren-hh-with (μkanren-hh-pad (μkanren-hh s1)
                                                     (μkanren-state-vars-count s1))
                                     depth: d)))
          (match/first (μkanren-discharge s2 c d)
            (#f (✗° s2))
            (,s3 (✓° s3)))))
      (δ ((f c) (μkanren-with s x*))))))
\end{scheme}
The shape is that of micro's \scm{andª} (\cref{ch:micro}): the body's stream is delayed with \scm{δ}, and
\scm{append-map§} post-processes it answer by answer, so the restoring step runs on every answer, however
late the stream produces it. Steps 1 and 3 pad $L$ because a variable's level is \emph{implicit} (the current
depth) until $L$ is extended past it: padding fixes it before the depth changes.

The signature $\Sigma$ is not carried. What carries the quantifier structure is the level of each variable:
a variable created at depth $d$ stands for an $\exists$ after the first $d$ universals of its branch (Miller's
mixed prefix \cite{Miller1991}, linearised by creation), and \textbf{Check} forbids exactly the substitutions
that would violate that prefix. NJK prove this tagging sound and complete (\cref{prop:rules-ui},
\cite[Prop.\,4.1]{NadathurJayaramanKwon1995}).

\begin{tutorial}[Why the depth must be restored on every answer]
A creation timestamp would tell a variable created before $c$ from one created after it, and one might think
that suffices: older variables may not see $c$. But a variable created \emph{after} the \scm{∀°} has exited
may not see $c$ either. With the depth restored on every answer, such a variable gets level $d < \ell(c)$, so
it cannot take $c$. This is λProlog's universe register \cite{NadathurJayaramanKwon1995}. The two classic
cases, \book{85}:
\begin{scheme}
(list (°->list #f (fresh° (q) (∀° (x) (=° q x))))         ; ∃Y∀x. Y = x
      (°->list #f (fresh° (q) (∀° (x) (fresh° (y) (=° y x)))
                              (=° q 'ok))))                ; ∀x∃Y. Y = x
\end{scheme}
\begin{answers}
(() ((λ () 'ok)))
\end{answers}
In the first, $q$ has level~0 and $x$
level~1: \textbf{Check} fails by case (1). In the second, $y$ is created inside, at level~1, and may take $x$.
\end{tutorial}

\begin{example}[Lowering, seen from the state]\label{ex:module-lower}
A variable created inside the scope and placed inside the value of an outer one must forget the scope:
\begin{scheme}
(define ((probe . vars) s)
  (for-each (λ (v) (let1 (w (μkanren-state-find v s))
                     (write v) (display " ↦ ") (write w)
                     (when (μkanren-var? w)
                       (display " level ") (display (μkanren-λ-level w s)))
                     (display "; ")))
            vars)
  (newline) (✓° s))
(°->list #f (fresh° (A) (∀° (x) (fresh° (B) (probe A B)
                                             (=° A `(arr ,B c))
                                             (probe A B)))))
\end{scheme}
\begin{answers}
_0 ↦ _0 level 0; _1 ↦ _1 level 1;
_0 ↦ (arr _1 c); _1 ↦ _2 level 0;
((λ (α) (cons 'arr (cons α (cons 'c '())))))
\end{answers}
$A$ (\texttt{\_0}) has level~0 and $B$ (\texttt{\_1}) level~1. Binding $A := (\mathtt{arr}\ B\ c)$ triggers
\textbf{Check} at level~0, which meets the unbound $B$ of level~$1>0$: case (4), lower. A new variable
\texttt{\_2} of level~0 is created and $B := \mathtt{\_2}$; \textbf{Check} eliminates no head and has no
arguments $\bar a$ to raise over, so the raising part is empty. This is NJK's \enquote{change to $i$ the tags of the variables in $t$ whose tag is $>i$}, done by a
binding instead of a mutation, so that it is undone by backtracking for free.
\end{example}

\subsection{Discharge at the exit of a scope}\label{sec:module-discharge}

When \scm{∀°} exits, some constraints may still mention $c$: suspended equations, generic disequalities of N,
absentos of micro's T. A residual that mentions $c$ outside its scope would be meaningless, and the printer
must never show an eigenvariable free. \scm{μkanren-discharge} closes them, back at depth $d$:
\begin{enumerate}
  \item \textbf{Raise.} Every unbound variable $Y$ of the equations and disequalities to be closed with
    $\ell(Y)\ge\ell(c)$ is raised:
    $Y := Y'\,c$ with $\ell(Y')=d$. Inside the scope $Y$ could depend on $c$; $Y'$ lives outside it, so the
    dependency becomes an explicit argument: $\forall c\,\exists Y$ is turned into $\exists Y'\,\forall c$
    with $Y = Y'\,c$. This is Miller's raising \book{107} (\cref{prop:lambda-raising}), a bijection
    on solutions.
  \item \textbf{Close.} A suspended equation $u \doteq v$ that mentions $c$ stands for
    $\forall c\,(u = v)$ and becomes $\lambda c.u \doteq \lambda c.v$, equivalent by \rulename{ξ}. A generic
    disequality is closed the same way, under \rulename{generic}'s reading of $c$ as a new name. An absento of T whose
    \emph{term} mentions $c$ is closed too, since absence in $\lambda c.u$ is absence in $u$. An absento whose
    \emph{datum} mentions $c$ moves to $E$ as \scm{(absent d u)}, unless the levels entail it: if $u$ does not
    hold $c$ and all its variables have levels below $\ell(c)$, no instance of $u$ can ever hold $c$, and the
    constraint is dropped.
  \item \textbf{Retry.} The raise bindings go through \scm{μkanren-post-unify}, which wakes $E$, so the closed
    constraints are retried. With no raise, the closed equations are retried anyway.
  \item \textbf{Repeat.} A raise may bring $c$ into more constraints (an absento of T on a raised $Y$ now
    watches $Y'\,c$; a prefix of D on $Y$ moves to N holding $c$). Discharge repeats until no raise is left.
\end{enumerate}
The constraints that cannot be closed by an abstraction (a type on $F\,x$; an absento whose datum is $c$ on a
term that may still hold $c$) keep $c$ and print under a \scm{∀} (\cref{sec:module-reify}).

\begin{example}[Discharge: raising, closing, retrying]\label{ex:module-discharge}\leavevmode
\begin{scheme}
(define ((probe label . vars) s)   ; prints values and the levels of unbound variables
  (display label) (display ": depth ") (display (μkanren-λ-depth s))
  (for-each (λ (v) (display ", ") (write v) (display " ↦ ")
                   (write (μkanren-state-find/value v s))
                   (let1 (w (μkanren-state-find v s))
                     (when (μkanren-var? w)
                       (display " level ") (display (μkanren-λ-level w s)))))
            vars)
  (display ", E ") (write (μkanren-λ-delayed s)) (newline)
  (✓° s))
(°->list #f (fresh° (q) (∀° (a) (fresh° (f) (=° (@ f a) `(g ,a ,a))
                                             (probe "inside" f)))
                        (=° q 'ok)))
(°->list #f (fresh° (q) (∀° (x) (=° (@ q x) `(g ,q)) (probe "inside" q))))
\end{scheme}
\begin{answers}
inside: depth 1, _1 ↦ _1 level 1, E ((= (@ _1 (eigen a)) (g (eigen a) (eigen a))))
((λ (α)
    (assert (≐ (ƛ (a) (@ α a a)) (ƛ (a) (cons 'g (cons a (cons a '()))))))
    'ok))
inside: depth 1, _0 ↦ _0 level 0, E ((= (@ _0 (eigen x)) (g _0)))
()
\end{answers}
In the first, $f$ is created inside the scope of $a$ (level 1), so $f\,a$ is no pattern (the book's 8.A, four
unifiers and no mgu, \book{213}) and is suspended. At the exit, $f$ has level $1\ge\ell(a)$: it is raised,
$f := f'\,a$ with $\ell(f')=0$, and the equation closes into $\lambda a.\,f'\,a\,a \doteq \lambda a.\,g\,a\,a$. That is
still no pattern (the argument $a$ is repeated), so it is reported. In the second, $F\,x \doteq g\,F$ is a pattern
equation, but $F$ itself occurs with no argument in its own body, case (4) with $\bar a \neq \langle\rangle$:
suspend. At the exit nothing needs raising ($\ell(F)=0$), the equation closes into
$\lambda x.\,F\,x \doteq \lambda x.\,g\,F$, which by η is $F \doteq \lambda x.\,g\,F$, and the retry fails the occurs
check: no answer, as it must be.
\end{example}

\subsection{\texorpdfstring{\rulename{augment}}{AUGMENT}}

\begin{mathpar}
\inferrule{\seq{\Sigma}{\Prog\cup\{D\}}{G}}{\seq{\Sigma}{\Prog}{D\imp G}}\;\rulename{augment}
\end{mathpar}
\scm{(⊃° (D ...) g ...)} expands into \scm{⊃ª}, which pushes the clauses on Γ for the body and gives the old Γ
back to every answer of the body:
\begin{scheme}
(define ((⊃ª clauses g) s)
  (let* ((x (μkanren-hh s))
         (Γ (μkanren-hh-Γ x))
         (Γ* (fold-right
               (λ (c Γ*)
                 (let1 (old (cond ((assq (car c) Γ*) => cdr) (else '())))
                   (cons (cons (car c) (cons (cdr c) old))
                         (if (null? old) Γ* (alist-delete (car c) Γ* eq?)))))
               Γ clauses)))
    (append-map§
      (λ (s1) (✓° (μkanren-with s1 (μkanren-hh-with (μkanren-hh s1) Γ: Γ))))
      (δ (g (μkanren-with s (μkanren-hh-with x Γ: Γ*)))))))
\end{scheme}
A clause is a pair \scm{(relation . procedure)}: the relation is the dynamic relation's own procedure, used
as an \scm{eq?} key, and the procedure maps the list of a call's arguments to the goal that the clause makes of
them. Since Γ lives in the state, \scm{or°} branches cannot interfere, and restoring it on every answer is the
stack discipline of the book \book{78}, \book{83}. Assuming a clause for a relation already in Γ replaces that
relation's entry by one with the new clause in front of the old ones, so a lookup costs the number of distinct
assumed relations, not the number of clauses. The \scm{fold-right} puts the first clause of a list on top.

\subsection{\texorpdfstring{\rulename{backchain}}{BACKCHAIN} and \texorpdfstring{\texttt{define-dynamic-relation}}{define-dynamic-relation}}

In λProlog every predicate can be extended by an assumption. Here only relations made by
\scm{define-dynamic-relation} can:
\begin{scheme}
(define-syntax define-dynamic-relation
  (syntax-rules ()
    ((_ (name arg ...))
     (define name
       (μkanren-dynamic
         (λ (arg ...) (λ (s) (δ ((μkanren-hyps° name (list arg ...) #f) s)))))))
    ((_ (name arg ...) g ...)
     (define name
       (μkanren-dynamic
         (λ (arg ...)
           (λ (s) (δ ((μkanren-hyps° name (list arg ...) (and° g ...)) s)))))))))

(define ((μkanren-hyps° key args static) s)
  (let* ((gs (map (λ (c) (c args))
                  (cond ((assq key (μkanren-hh-Γ (μkanren-hh s))) => cdr) (else '()))))
         (gs (if static (append gs (list static)) gs)))
    (if (null? gs) '() ((reduce-right orª #f gs) s))))
\end{scheme}
A call disjoins, in one flat \scm{or°}, the assumed clauses for the relation, \emph{most recent first}, and then
the static body. If $C_1,\dots,C_m$ are the clauses assumed for $r$ in the current Γ, in the order they were
assumed, and $B_r$ is the static body of $r$ (absent when it has none), then
\[
  r\,\bar t \;\rightsquigarrow\; C_m(\bar t) \;\vee\; \dots \;\vee\; C_1(\bar t) \;\vee\; B_r(\bar t),
\]
where $C_i(\bar t)$ unifies $\bar t$ with the head of $C_i$, its own \scm{∀} variables renamed and its free
variables shared (see the clause syntax below), and runs its body. This is
\rulename{backchain} (\cref{fig:hh-rules}) with the choice of clause made by \scm{or°} instead of by a selection rule; a
relation with neither hypotheses nor static body fails. When every clause answers at once, the answers come in λProlog's order; otherwise
interleaving decides. The relation is tagged with CHICKEN's \scm{extend-procedure}, so \scm{μkanren-dynamic?}
is an $O(1)$ test and no global registry is needed. Assuming a clause for a relation that is not dynamic is an
error, not a silent no-op:
\begin{scheme}
(define-dynamic-relation (n° x) (=° x 1))
(list (°->list/datum (fresh° (q) (⊃° ((n° 2)) (⊃° ((n° 3)) (n° q)))))
      (°->list/datum (fresh° (q) (⊃° ((n° 2) (n° 3)) (n° q)))))
(define (error-message thunk)     ; the message of the error `thunk` raises
  (handle-exceptions exn ((condition-property-accessor 'exn 'message #f) exn)
    (begin (thunk) #f)))
(define-relation (r° x) (=° x 1))
(error-message (λ () (°->list #f (fresh° (q) (⊃° ((r° 2)) (r° q))))))
\end{scheme}
\begin{answers}
((3 2 1) (2 3 1))
"not a dynamic relation, see define-dynamic-relation"
\end{answers}
The two orders are the book's \book{78}: \lp{p 2 => p 3 => p X} gives 3, 2, 1, and
\lp{(p 2 & p 3) => p X} gives 2, 3, 1. In the first query the inner \scm{⊃°} runs last, so Γ maps \scm{n°} to
the clauses for 3 and 2, in this order, and the call disjoins $q=3$, $q=2$ and the static $q=1$. In the second,
\scm{fold-right} puts the first clause of the list on top, so the clause for 2 comes before the one for 3.

\subsection{Clause syntax}

\begin{longtable}{@{}p{0.36\fullwidthlen}p{0.58\fullwidthlen}@{}}
\toprule
clause $D$ & meaning and expansion \\
\midrule
\endhead
\scm{(rel t ...)} & a fact; \scm{(λ (args) (fresh° (v ...) (=° args (list t ...))))} \\
\scm{((rel t ...) ← g ...)} & a rule, infix \\
\scm{(← (rel t ...) g ...)} & the same rule, prefix; rewritten to the infix form \\
\scm{(∀ (w ...) D)} & $D$ with its own variables \scm{w ...}, accumulated into the \scm{fresh°} of the expansion \\
\bottomrule
\end{longtable}
\noindent A rule expands to \scm{(μkanren-clause* rel (λ (args) (fresh° (v ...) (=° args (list t ...)) g ...)))},
and \scm{μkanren-clause*} raises the error above when \scm{rel} is not dynamic. Two kinds of variables appear
in a clause, as in the book \book{86--88}:
\begin{itemize}
  \item its \emph{free} logic variables are those of the closure: shared with the context, and among all uses of
    the clause;
  \item its \emph{own} \scm{∀} variables are made by the \scm{fresh°} inside the procedure, so they are renamed at
    each use, at the current level. That is NJK's \enquote{new tagged instance} of a clause
    \cite[p.\,14]{NadathurJayaramanKwon1995}.
\end{itemize}
\begin{scheme}
(define-dynamic-relation (h° x))
(list (°->list/datum (fresh° (q y) (⊃° ((h° `(,y ,y))) (h° '(1 1)) (h° '(2 2)))))
      (°->list/datum (fresh° (q) (⊃° ((∀ (z) (h° `(,z ,z)))) (h° '(1 1)) (h° '(2 2)))
                                 (=° q 'ok))))
\end{scheme}
\begin{answers}
(() (ok))
\end{answers}
With the shared $y$, the first use binds $y:=1$ and the second use cannot match \scm{(2 2)}; with the clause's
own $z$, each use gets a new $z$.

\begin{example}[Levels reach into clauses {\cite[p.\,16]{NadathurJayaramanKwon1995}}]\label{ex:module-njk}
\emergencystretch=2em
With the program $\forall z\,(q\,z \imp p\,(d\,z))$, the goal $\exists x\,\forall y\,(q\,y \imp p\,x)$ must fail
and $\forall y\,\exists x\,(q\,y \imp p\,x)$ must succeed (\cref{ex:rules-scope}):
\begin{scheme}
(define-dynamic-relation (q° z))
(define-dynamic-relation (p° x) (fresh° (z) (=° x `(d ,z)) (q° z)))
(list (°->list #f (fresh° (x) (∀° (y) (⊃° ((q° y)) (p° x)))))
      (°->list #f (fresh° (r) (∀° (y) (fresh° (x) (⊃° ((q° y)) (p° x)))) (=° r 'ok))))
\end{scheme}
\begin{answers}
(() ((λ () 'ok)))
\end{answers}
In the first goal, $x$ has level~0. Backchaining on \scm{p°} creates the clause variable $z$ inside the scope of
$y$ (level~1) and binds $x := (d\ z)$; \textbf{Check} lowers $z$ to level~0. The hypothesis \scm{(q° y)} then
asks $z = y$, a level-0 variable against a level-1 eigenvariable: scope failure. Without lowering, $z$ would
keep level~1 and take $y$, and $x = (d\ y)$ would leak $y$ out of its scope.
\end{example}

\subsection{Worked example: type inference}\label{sec:module-typeof}

The book's type inference for untyped λ-terms (\cref{sec:lambdatree-typeof}, \book{190--191}) uses all three
rules at once:
\begin{lprolog}
typeof (app M N) A :- typeof M (arr B A), typeof N B.
typeof (abs M) (arr A B) :- pi x\ typeof x A => typeof (M x) B.
\end{lprolog}
In the module's syntax ({\small\texttt{test/programs.scm}} of {\small\texttt{doc/hol-kanren/prototype/}}):
\begin{scheme}
(define-dynamic-relation (typeof° m t)
  (cond°
    ((fresh° (M N B) (=° m `(app ,M ,N)) (typeof° M `(arr ,B ,t)) (typeof° N B)))
    ((fresh° (R A B) (=° m `(abs ,R)) (=° t `(arr ,A ,B))
       (∀° (x) (⊃° ((typeof° x A)) (typeof° (@ R x) B)))))))
\end{scheme}
The object term \lp{abs x\ M} is \scm{`(abs ,(ƛ (x) M))}; \lp{(M x)} is \scm{(@ R x)}, a β-redex once $R$ is
bound. The answers for S, I and $\lambda x.\,x\,x$ (from \texttt{examples.out} in the same directory):
\begin{answers}
(((arr (arr α (arr β γ)) (arr (arr α β) (arr α γ)))) ((arr α α)) ())
\end{answers}
To see the rules at work, here is the same program with a trace goal that prints the depth and the walked
arguments, with the level of each unbound variable after \texttt{@}. The trace in the static body runs only
when the static clauses are tried, after the hypotheses:
\begin{scheme}
(define ((trace label . ts) s)
  (display label) (display " [depth ") (display (μkanren-λ-depth s)) (display "]")
  (for-each (λ (t) (display " ") (write (μkanren-state-find/value t s))
                   (let1 (w (μkanren-state-find t s))
                     (when (μkanren-var? w)
                       (display "@") (display (μkanren-λ-level w s)))))
            ts)
  (newline) (✓° s))
(define-dynamic-relation (typeof° m t)
  (trace "static typeof°" m t)
  (cond°
    ((fresh° (M N B) (=° m `(app ,M ,N)) (trace " app" M N B)
       (typeof° M `(arr ,B ,t)) (typeof° N B)))
    ((fresh° (R A B) (=° m `(abs ,R)) (=° t `(arr ,A ,B)) (trace " abs" R A B)
       (∀° (x) (⊃° ((typeof° x A)) (typeof° (@ R x) B)))))))
(°->list/datum (fresh° (q) (typeof° `(abs ,(ƛ (x) `(abs ,(ƛ (y) x)))) q)))
(°->list/datum (fresh° (q) (typeof° `(abs ,(ƛ (x) `(app ,x ,x))) q)))
\end{scheme}
\begin{answers}
static typeof° [depth 0] (abs (ƛ (x) (abs (ƛ (y) x)))) _0@0
 abs [depth 0] (ƛ (x) (abs (ƛ (y) x))) _2@0 _3@0
static typeof° [depth 1] (abs (ƛ (y) (eigen x))) _3@0
 abs [depth 1] (ƛ (y) (eigen x)) _7@0 _8@0
static typeof° [depth 2] (eigen x) _8@0
((arr α (arr β α)))
static typeof° [depth 0] (abs (ƛ (x) (app x x))) _0@0
 abs [depth 0] (ƛ (x) (app x x)) _2@0 _3@0
static typeof° [depth 1] (app (eigen x) (eigen x)) _3@0
 app [depth 1] (eigen x) (eigen x) _6@1
static typeof° [depth 1] (eigen x) (arr _6 _3)
static typeof° [depth 1] (eigen x) _7@0
()
\end{answers}

\begin{tutorial}[$\lambda x.\lambda y.\,x$, step by step]
\begin{enumerate}
  \item \emph{Depth 0, Γ empty.} The query $T$ is \texttt{\_0}. No hypothesis for \scm{typeof°}, so the static
    body runs. The \texttt{abs} clause creates $R,A,B$ (\texttt{\_1}, \texttt{\_2}, \texttt{\_3}, level~0),
    binds $R := \lambda x.\,\mathtt{abs}\,(\lambda y.\,x)$ and $T := (\mathtt{arr}\ A\ B)$.
  \item \rulename{generic}. \scm{∀° (x)} records levels, raises the depth to~1 and makes the eigenvariable $x$ of
    level~1.
  \item \rulename{augment}. \scm{⊃°} pushes $\langle\texttt{typeof°}, \lambda\mathit{args}.\ \mathit{args} \doteq (x\ A)\rangle$ on Γ.
  \item \emph{β.} The goal is \scm{(typeof° (@ R x) B)}; \scm{@} reduces the redex to
    $\mathtt{abs}\,(\lambda y.\,x)$, with $x$ now the eigenvariable, printed \scm{(eigen x)}.
  \item \rulename{backchain}, depth 1. The hypothesis is tried first: $(\mathtt{abs}\ \dots) \doteq x$ clashes.
    Then the static body: the \texttt{abs} clause creates $R',A',B'$ at level~1 and binds
    $B := (\mathtt{arr}\ A'\ B')$. Here \textbf{Check} runs: $B$ has level~0, $A'$ and $B'$ level~1, so both are
    \emph{lowered}, to new variables \texttt{\_7}, \texttt{\_8} of level~0 (the trace shows
    \texttt{\_7@0 \_8@0}, the walked values). This is right: the type of the inner abstraction is part of the
    outer type, which lives outside the scope of $x$.
  \item \rulename{generic}, \rulename{augment} again. Depth~2, eigenvariable $y$ of level~2, Γ now holds two clauses for
    \scm{typeof°}, $(y\ A')$ on top of $(x\ A)$.
  \item \rulename{backchain}, depth 2. The goal is \scm{(typeof° x B')}. The most recent hypothesis gives
    $x \doteq y$, a clash of two eigenvariables. The older one gives $B' \doteq A$: a level-0 variable bound to
    another, fine. The static body is tried as well (the last trace line) and fails, since $x$ is neither an
    application nor an abstraction.
  \item \emph{Exits.} Each \scm{∀°} gives its answers back their depth and discharges nothing (no constraint
    mentions $x$ or $y$); each \scm{⊃°} gives Γ back. The answer is
    $T = \mathtt{arr}\ A\ (\mathtt{arr}\ A'\ A)$, printed \scm{(arr α (arr β α))}.
\end{enumerate}
\end{tutorial}

\begin{tutorial}[$\lambda x.\,x\,x$ fails by the occurs check]
At depth~1, with the hypothesis $\mathtt{typeof}\ x\ A$, the goal is
$\mathtt{typeof}\ (\mathtt{app}\ x\ x)\ B$. The \texttt{app} clause binds $M,N := x$ and creates its $B''$
(\texttt{\_6}) at level~1. The first subgoal $\mathtt{typeof}\ x\ (\mathtt{arr}\ B''\ B)$ uses the hypothesis:
$A := (\mathtt{arr}\ B''\ B)$, and since $\ell(A)=0$, \textbf{Check} lowers $B''$ to the level-0
\texttt{\_7}. The second subgoal $\mathtt{typeof}\ x\ B''$ (printed with $B''$ walked to \texttt{\_7}) asks
$\mathtt{\_7} \doteq A = (\mathtt{arr}\ \mathtt{\_7}\ B)$: micro's occurs check fails, and since the term holds
no λ-term, the \texttt{occurs} hook does not intervene. The static attempts fail too, and there is no answer,
the book's \enquote{no} \book{191}.
\end{tutorial}

% ------------------------------------------------------------------------------------------------
\section{micro's stores on λ-terms}\label{sec:module-stores}

λProlog has no disequality, no absento and no type constraints such as \scm{symbol°} (its types are
static). \micro{} has all three (\cref{sec:micro-stores}), and
every relation built on them must keep working when its arguments are λ-terms. The readings adopted here extend
micro's semantics over λProlog's term model: equality is $\alpha\beta\eta$-equality, a flexible term is
undecided until its head is known, and an eigenvariable is, by default, a new name.

\subsection{Disequality: the \texorpdfstring{\texttt{diseq}}{diseq} hook and the store N}

\begin{definition}[Trial]\label{def:module-trial}
The module \emph{claims} a disequality $u \neq v$ when either side, walked, holds a λ-record or an
eigenvariable, and also, under the extensional reading once a \scm{∀°} has run, a first-order one (in D an
absento of T could absorb it, and absence reads eigenvariables generically). Its \emph{trial} unifies $u \doteq v$ in the current state, with \textbf{Check}, and classifies
the result (\scm{μkanren-trial}):
\begin{description}[style=nextline,leftmargin=2.2em]
  \item[\texttt{unsat}] no unifier: the disequality holds for good and is dropped;
  \item[\texttt{entailed}] a unifier that binds no variable older than the trial (or binds older variables only
    to distinct variables the trial itself created) and suspends nothing: $u=v$ already holds, and \scm{≠°}
    fails;
  \item[kept] otherwise: the pair is stored in N as \scm{(reading u . v)} and re-tried after every \scm{=°}.
\end{description}
\end{definition}
The variables a trial creates (by pruning, by \rulename{FF-same}, by lowering) are \emph{existential inside the
negation}: $\lambda x\,y.\,F\,x \neq \lambda x\,y.\,G\,y$ means $\neg\exists H.\ F = \lambda x.H \wedge G =
\lambda y.H$. A prefix of micro's D would wrongly store them as free variables, which is one reason N is a
separate store.
\begin{scheme}
(°->list #f (fresh° (q) (=° q (ƛ (x) x)) (≠° q (ƛ (y) y))))     ; entailed mod α
(°->list #f (fresh° (q) (≠° (@ (ƛ (x) `(g ,x)) 'a) '(g a))))    ; entailed mod β
(°->list #f (fresh° (q) (≠° (ƛ (x) x) (ƛ (x) 'k)) (=° q 'ok)))  ; holds for good
(°->list/datum (fresh° (q F) (=° q F) (≠° (@ F 'a) 'a)
                        (=° F (ƛ (z) z))))                 ; woken, violated
(°->list/datum (fresh° (q F) (=° q F) (≠° (@ F 'a) 'a)
                        (=° F (ƛ (z) 'b))))                ; woken, holds
(°->list #f (fresh° (q f g) (=° q (list f g))
              (≠° (ƛ (x y) (@ f x)) (ƛ (x y) (@ g y)))))
\end{scheme}
\begin{answers}
()
()
((λ () 'ok))
()
((ƛ (z) b))
((λ (α β)
    (begin (deny (∃ (_3) (and (equal? α (ƛ (x) _3)) (equal? β (ƛ (y) _3))))))
    (cons α (cons β '()))))
\end{answers}
Line by line: the first trial succeeds under \rulename{ξ} binding nothing, so it is \texttt{entailed} and
\scm{≠°} fails; the second never reaches the module, since \scm{@} has already β-reduced its left side to the
list \scm{(g a)} and micro's own \scm{≠°} fails on two equal terms; the third trial clashes ($x$ against $k$ under \rulename{ξ}), so the disequality is \texttt{unsat} and
dropped. In the fourth and fifth the trial $F\,a \doteq a$ is no pattern ($a$ is a constant) and suspends, so the
pair is kept in N; the binding of $F$ retries it, and it is violated by $\lambda z.z$ and holds for good
with $\lambda z.b$. In the last, the trial is the flexible–flexible problem of \cref{ex:module-ff} under two
binders: it binds the older $F$ and $G$ to $\lambda x.H$ and $\lambda y.H$ with a new $H$, so the pair is kept,
and the answer denies exactly that mgu, with $H$ (\texttt{\_3}) existential inside the \scm{deny}.

\paragraph{Migration from D to N.} A prefix of D born first-order may come to hold a λ-term or an eigenvariable
later, at any depth below its variables. The \texttt{post} hook moves it to N, as the tuple disequality it
stands for, \emph{before} micro re-verifies D. So D never holds a λ-term, and micro's plain trial never decides a
λ-equation without the scope check and the reading. The probe prints N, D and E:
\begin{scheme}
(define ((probe label) s)
  (display label) (display ": N ") (write (μkanren-λ-disequalities s))
  (display ", D ") (write (μkanren-state-D s))
  (display ", E ") (write (μkanren-λ-delayed s))
  (newline) (✓° s))
(°->list #f (fresh° (q Y) (≠° q (list Y)) (probe "before")
                          (=° Y (ƛ (z) z)) (probe "after")))
\end{scheme}
\begin{answers}
before: N (), D (((_0 _1))), E ()
after: N ((generic (_0) (_1))), D (), E ()
((λ (α) (begin (deny (equal? α (cons (ƛ (z) z) '())))) α))
\end{answers}
To keep first-order search cheap, the hook looks at no more than 64 nodes of the new bindings to decide whether
they may hold a λ-term; past that bound it walks D and A instead.\sidenote{\Cref{ch:roadmap} has the measurements.}

\paragraph{Subsumption by A.} A trial that binds a variable typed in A to a rigid λ-term or an eigenvariable
of which the type's predicate is false has no solution; \scm{symbol°} and \scm{number°} hold of neither. So the disequality is dropped, as micro
drops \scm{(≠° q 5)} under \scm{(symbol° q)}:
\begin{scheme}
(°->list #f (fresh° (q) (symbol° q) (≠° q (ƛ (x) x))))
\end{scheme}
\begin{answers}
((λ (α) (assert (every (μ v (symbol? v)) (list α))) α))
\end{answers}

\paragraph{Two readings of $\forall$.} What does $\forall x.\ q \neq x$ mean? λProlog has no $\neq$, but it has
\lp{pi x\ not (q = x)}, which succeeds: $x$ is a new name, distinct from everything. The parameter
\scm{μkanren-λ-reading} chooses, and each disequality keeps the reading in force when it was posted.
\begin{itemize}
  \item \scm{'generic} (the default) reads an eigenvariable as \rulename{generic}'s new name, ∇-like, as αKanren's noms
    \cite{ByrdFriedman2007}. At the exit of the scope the constraint is closed: $\forall x.\ F\,x \neq x$
    becomes $F \neq \lambda x.x$, consistently whichever comes first, the binding or the $\neq$.
  \item \scm{'extensional} reads an eigenvariable as any term of its universe, as in the miniKanren work on
    universal quantification \cite{Moiseenko2019,JinEtAl2021}. The trial replaces each eigenvariable by a new
    variable of its level, after raising the variables that may depend on it, and counts as entailed when it
    binds older variables only to distinct such variables. Then $\forall x.\ q \neq x$ fails, since $x$ could
    be $q$.
\end{itemize}
\begin{scheme}
(°->list #f (fresh° (q) (∀° (x) (≠° q x))))
(°->list #f (fresh° (q F) (=° q F) (∀° (x) (≠° (@ F x) x))))
(°->list #f (fresh° (q F) (=° q F) (∀° (x) (≠° (@ F x) x)) (=° F (ƛ (z) z))))
(°->list #f (fresh° (q F) (=° q F) (∀° (x) (≠° (@ F x) x)) (=° F (ƛ (z) 'c))))
(parameterize ((μkanren-λ-reading 'extensional))
  (list (°->list #f (fresh° (q) (∀° (x) (≠° q x))))
        (°->list #f (fresh° (q F) (=° q F) (∀° (x) (≠° (@ F x) x))))
        (°->list #f (fresh° (q F) (=° q F) (∀° (x) (≠° (@ F x) x))
                                (=° F (ƛ (z) 'c))))))
\end{scheme}
\begin{answers}
((λ (α) α))
((λ (α) (begin (deny (equal? α (ƛ (x) x)))) α))
()
((λ () (ƛ (z) 'c)))
(() ((λ (α) (∀ (x) (begin (deny (≐ (@ α x) x)))) α)) ())
\end{answers}
Under the generic reading,
$\forall x.\ F\,x\neq x$ is the disequality $F \neq \lambda x.x$: violated by $F := \lambda z.z$, satisfied by
$F:=\lambda z.c$. Under the extensional reading it says that $F\,t \neq t$ for \emph{every} $t$ of the
universe, so $F := \lambda z.c$ violates it ($t = c$); since it cannot be closed by an abstraction, it is printed
under a \scm{∀}.

\paragraph{Projection.} micro drops a disequality over a variable that is out of the answer, since such a
variable can always be chosen to keep the sides apart. That fails for a variable inside an argument of a flexible
term: $\exists Y.\ G \neq F\,Y$ is $F \neq \lambda z.G$, since $F$ may ignore its argument. Such an entry is kept,
and its variable prints as a working variable, read existentially:
\begin{scheme}
(°->pretty (fresh° (q) (fresh° (F G Y) (=° q (list F G)) (≠° G (@ F Y)))))
\end{scheme}
\begin{answers}
(((α β) : (≠ β (@ α _3))))
\end{answers}

\subsection{Absence: the \texorpdfstring{\texttt{parts}}{parts} hook}

\scm{absent°} looks for its datum in the \emph{parts} of a λ-term (\scm{μkanren-λ-parts}), never in binder names
or indices. A flexible term has no parts yet: it is \scm{'stuck} and watched whole, like a variable, and decided
once its head is known. Below a binder, the parts that mention the binder stay closed ($\lambda x.(a\,.\,b)$ is
looked into as $\lambda x.a$ and $\lambda x.b$), so a watched term never has a loose index.
\begin{scheme}
(°->list #f (fresh° (q F) (=° q F) (absent° 'c (@ F 'c))))
(°->list/datum (fresh° (q F) (=° q F) (absent° 'c (@ F 'c)) (=° F (ƛ (z) 'b))))
(°->list #f (fresh° (q) (absent° 'x (ƛ (x) `(g ,x))) (=° q 'ok)))
(°->list #f (fresh° (q F) (=° q F) (absent° 'c (ƛ (x) `(,(@ F x) c)))))
(°->list #f (fresh° (q F) (=° q F) (∀° (x) (absent° x (@ F x)))))
(°->list #f (fresh° (q) (∀° (x) (absent° x q))))
\end{scheme}
\begin{answers}
((λ (α) (assert (absent? 'c (@ α 'c))) α))
((ƛ (z) b))
((λ () 'ok))
()
((λ (α) (∀ (x) (assert (absent? x (@ α x)))) α))
((λ (α) α))
\end{answers}
In order: $c \notin F\,c$ waits; $F := \lambda z.b$ drops the argument, so it holds; the binder name \texttt{x}
is not the symbol \scm{'x}; $c$ is found at once outside the flexible part; an absento of an eigenvariable on a
term that may still hold it keeps $x$ and prints under \scm{∀}; and $\forall x.\ x \notin q$ is entailed by the
levels ($q$ has level $0<1$, so no instance of $q$ can hold $x$) and dropped. The datum may itself be a λ-term,
or reduce to an atom (\scm{(@ G 'b)} is \scm{a} once $G := \lambda z.a$): it is then unified, modulo βη, as micro
unifies a compound datum.

\subsection{Types: the \texorpdfstring{\texttt{tag}}{tag} hook}

A type on a flexible term waits in E as \scm{(tag tag u)}; the \texttt{post} hook does the same for a variable of
A just bound to one. A rigid λ-term or an eigenvariable is decided by the type's predicate: abstractions and
eigenvariables are neither symbols nor numbers, while a type of the user's made with \scm{define-μkanren-tag}
may hold of them. micro's types are disjoint, so two different types on the same flexible term fail at once:
\begin{scheme}
(°->list #f (fresh° (q F) (=° q (@ F 'a)) (symbol° q)))
(°->list/datum (fresh° (q F) (=° q (@ F 'a)) (symbol° q) (=° F (ƛ (z) z))))
(°->list #f (fresh° (q F) (=° q (@ F 'a)) (symbol° q) (number° q)))
(°->list #f (fresh° (q) (∀° (x) (symbol° x))))
(°->list #f (fresh° (q F) (=° q F) (∀° (x) (symbol° (@ F x)))))
\end{scheme}
\begin{answers}
((λ (α) (assert (every (μ v (symbol? v)) (list (@ α 'a)))) (@ α 'a)))
(a)
()
()
((λ (α) (∀ (x) (assert (every (μ v (symbol? v)) (list (@ α x))))) α))
\end{answers}

\begin{pitfall}[Readings are per constraint, and only for $\neq$]
The extensional reading applies to \scm{≠°} only; \scm{absent°} and the types always read an eigenvariable as
\rulename{generic} does. A disequality keeps the reading in force when it was posted; a first-order disequality that
migrates from D to N later takes the reading in force when it migrates, since D records none. The two can differ
only if the program changes \scm{μkanren-λ-reading} between the two goals.
\end{pitfall}

% ------------------------------------------------------------------------------------------------
\section{Reification}\label{sec:module-reify}

micro reifies an answer as a λ-expression whose parameters name the answer's variables and whose body lists the
constraints and ends with the value (\cref{sec:micro-reify}). The module plugs into that through the hooks
\texttt{value}, \texttt{repr}, \texttt{terms}, \texttt{forms} and \texttt{environment} (\cref{ch:design}).

\paragraph{Values.} Deep values are β-normal and η-short (\scm{μkanren-λ-value}). They are used by answers, by
D, A and T, by \scm{project°} and by aggregation.

\paragraph{Binders.} An abstraction prints as \scm{(ƛ (x y) body)}, nested binders merged. A binder is named
after its hint, renamed apart from the binders around it, from every symbol of its body, and from the operators
of the rebuilding code (\scm{cons}, \scm{quote}, \scm{vector}, \scm{make-record-instance}, \scm{ƛ}, \scm{@},
\scm{eigen}, \scm{bvar}). So neither capture nor shadowing can mislead the reader (Ma and Friedman's
\enquote{true = false}, \cref{sec:lkanren-merge}), and a ground answer evaluates back to the same term:
\begin{scheme}
(°->list/datum (fresh° (q) (=° q (ƛ (a) (ƛ (a) a)))))
(μkanren-datum (μkanren-abstract 'cons (λ (c) `(a ,c))))
\end{scheme}
\begin{answers}
((ƛ (a a1) a1))
(ƛ (cons1) (a cons1))
\end{answers}

\paragraph{Applications} print as \scm{(@ h a ...)}, so a flexible application \scm{(@ α x)} cannot be mistaken
for the list \scm{(α x)}.

\paragraph{Residual forms.} The constraints of the λ layer print as Scheme forms that can be evaluated:
\begin{fullwidth}
\medskip\small\centering
\begin{tabular}{@{}p{0.36\linewidth}p{0.58\linewidth}@{}}
\toprule
constraint & form in the answer \\
\midrule
suspended equation (E) & \texttt{(assert (≐ u v))} \\
type on a flexible term (E) & \texttt{(assert (every (μ v (symbol? v)) (list u)))} \\
absento with a pending term (E, T) & \texttt{(assert (absent? d u))} \\
disequality whose trial binds (N) & \texttt{(begin (deny (equal? α t)))}, with \texttt{(∃ (\_3) ...)} inside \texttt{deny} when the trial creates variables \\
disequality whose trial only suspends (N) & \texttt{(begin (deny (≐ u v)))} \\
any of them still holding an eigenvariable & wrapped in \texttt{(∀ (x) form)} \\
\bottomrule
\end{tabular}\par\bigskip
\end{fullwidth}
\scm{≐} is exported: it decides αβη-equality of closed terms by unifying them in the empty state, so the forms
evaluate. Through the \texttt{terms} hook, the variables of residual equations are reified \emph{after} those of
the answer, so they are named \texttt{α}, \texttt{β}, \dots{} rather than \texttt{\_2}
(\cref{ex:module-forget}). micro's show module gives the αKanren-style view \scm{°->pretty}:
\begin{scheme}
(°->pretty (fresh° (q) (∀° (x) (=° (@ q `(g ,x)) `(h (g ,x))))))
(°->pretty (fresh° (q) (∀° (x) (=° (@ q x x) x))))
\end{scheme}
\begin{answers}
((α : (= (ƛ (x) (@ α (g x))) (ƛ (x) (h (g x))))))
((α : (= (ƛ (x) (@ α x x)) (ƛ (x) x))))
\end{answers}

\paragraph{No free eigenvariable.} Every constraint that holds an eigenvariable past its scope is closed, by
\scm{ƛ} or by \scm{∀}, so no eigenvariable ever appears free in an answer. The one value that could hold one is the
answer itself, when the variable that \scm{°->list} reifies is created inside a \scm{∀°}. λProlog reports no binding
for a variable quantified inside \lp{pi x\ }, and this answer cannot be told without $x$, so reifying it is an error:
\begin{scheme}
(error-message (λ () (°->list #f (∀° (x) (fresh° (q) (=° q x))))))
\end{scheme}
\begin{answers}
"the answer holds the eigenvariable of a ∀° it was created in: create the variable it reifies outside every ∀°"
\end{answers}

\paragraph{Live ground answers.} \scm{ƛ}, \scm{@} and \scm{≐} are bound in the module's environment, which the
\texttt{environment} hook hands to \scm{°->list/ground}. A ground answer is therefore a λ-term one can apply:
\begin{scheme}
(@ (car (°->list/ground (fresh° (q) (=° q (ƛ (x y) `(g ,x ,y)))))) 'b 'b)
\end{scheme}
\begin{answers}
(g b b)
\end{answers}
\scm{°->list/datum} answers plain data instead, comparable with \scm{equal?} against quoted book answers; that is what
the suite uses, and \scm{μkanren-repr->datum} turns micro's rebuilding code into such data.

\begin{pitfall}[Hints are not part of equality]
\scm{equal?} sees hints, the unifier does not. \scm{(ƛ (x) x)} and \scm{(ƛ (y) y)} unify and are \scm{μkanren-α=?},
but they are not \scm{equal?}: \scm{equal?} compares records slot by slot, hint included, and the hint is
kept in the record so that answers print the user's binder names. When comparing λ-terms in host code, use \scm{μkanren-α=?} or \scm{≐}.
A related trap is Scheme's own: \scm{(ƛ (quote) 'a)} shadows \scm{quote} while the body is built; terms built through
\scm{μkanren-abstract} with such a hint print and evaluate correctly.
\end{pitfall}

% ------------------------------------------------------------------------------------------------
\section{API reference}\label{sec:module-api}

\begin{longtable}{@{}>{\raggedright\arraybackslash}p{0.40\fullwidthlen}>{\raggedright\arraybackslash}p{0.56\fullwidthlen}@{}}
\toprule
name & meaning \\
\midrule
\endfirsthead
\toprule
name & meaning \\
\midrule
\endhead
\multicolumn{2}{@{}l}{\emph{Terms}}\\
\scm{(ƛ (x ...) body)} & λ-abstraction over Scheme binders, η-short \\
\scm{(@ f a ...)} & application; β at once on an abstraction, nested applications flattened \\
\scm{(≐ u v)} & αβη-equality of closed terms, for residual forms \\
\scm{(μkanren-abstract name f)} & the abstraction of the term \scm{(f p)}, \scm{p} a placeholder; binder hinted \scm{name} \\
\scm{(μkanren-apply f args)} & \scm{@} on a list of arguments \\
\scm{(μkanren-hnorm t s)} & head normal form of \scm{t} in state \scm{s} (\cref{def:module-hnf}) \\
\scm{(μkanren-α=? u v)} & α-equivalence of deep values, hints ignored \\
\scm{(μkanren-datum t)} & a term as data, printed without a state \\
\scm{μkanren-lam?}\newline \scm{μkanren-app?}\newline \scm{μkanren-eigen?}\newline \scm{μkanren-bvar?}\newline \scm{μkanren-λ-term?} & record predicates; the last is their union \\
\scm{μkanren-lam-hint}\newline \scm{μkanren-eigen-name}\newline \scm{μkanren-eigen-level} & accessors \\
\multicolumn{2}{@{}l}{\emph{Goals}}\\
\scm{(∀° (x ...) g ...)} & \rulename{generic}; \scm{(∀ª name f)} its procedural form, \scm{f} taking the eigenvariable \\
\scm{(⊃° (D ...) g ...)} & \rulename{augment}; \scm{(⊃ª clauses g)} its procedural form \\
\scm{(define-dynamic-relation}\newline \scm{  (r a ...) g ...)} & a relation that hypotheses can extend; with no \scm{g}, no static clauses \\
\scm{(μkanren-dynamic f)}\newline \scm{(μkanren-dynamic? f)} & tag a procedure as dynamic, test the tag \\
\scm{μkanren-hyps°}\newline \scm{μkanren-clause}\newline \scm{μkanren-clause/vars}\newline \scm{μkanren-clause*} & what the macros expand to \\
\multicolumn{2}{@{}l}{\emph{Clauses $D$}}\\
\scm{(r t ...)}\newline \scm{((r t ...) ← g ...)}\newline \scm{(← (r t ...) g ...)}\newline \scm{(∀ (v ...) D)} & fact, rule (infix and prefix), own variables \\
\multicolumn{2}{@{}l}{\emph{Answers}}\\
\scm{(°->list/datum g)} & the answers of \scm{g} as plain data, constraints dropped \\
\scm{(μkanren-repr->datum r)} & micro's rebuilding code as data \\
\multicolumn{2}{@{}l}{\emph{Engine}}\\
\scm{μkanren-λ-reading} & parameter, \scm{'generic} (default) or \scm{'extensional} \\
\scm{(μkanren-λ-depth s)} & the current universe level \\
\scm{(μkanren-λ-level α s)} & the level of variable \scm{α} \\
\scm{(μkanren-λ-delayed s)} & the suspended constraints E \\
\scm{(μkanren-λ-disequalities s)} & the λ-disequalities N \\
\scm{μkanren-λ-extension} & the registered extension record \\
\multicolumn{2}{@{}l}{\emph{In \texttt{(aux kanren micro show)}}}\\
\scm{(°->pretty g)}\newline \scm{μkanren-answer->pretty} & αKanren-style answers, \scm{(value : constraint ...)} \\
\bottomrule
\end{longtable}

% ------------------------------------------------------------------------------------------------
\section{Summary}\label{sec:module-summary}

\begin{itemize}\emergencystretch=2em
  \item \textbf{Terms.} Four records, locally nameless: \texttt{μkanren-\allowbreak lam} (interned hint, body),
    \texttt{μkanren-\allowbreak app}, \texttt{μkanren-\allowbreak eigen} (name, level, uninterned id), \texttt{μkanren-\allowbreak bvar}. \scm{ƛ} uses
    Scheme binders and builds η-short; \scm{@} β-reduces at once. Terms built alike are \scm{equal?}; α-equality is
    \scm{μkanren-α=?} and \scm{≐}. Lists remain data: \scm{(@ 'g 'a)} is not \scm{'(g a)}.
  \item \textbf{State.} One immutable record in micro's slot \texttt{X}: depth, top, the levels $L$/$L_n$ (an
    sbral, lazily padded), Γ (one entry per relation), E (suspended constraints), N (λ-disequalities). \scm{∀°}
    and \scm{⊃°} restore depth and Γ on every answer; branches never share it.
  \item \textbf{Unification.} micro does the first-order part; the hook applies \rulename{ξ}, \rulename{η},
    \rulename{Rigid}, \rulename{FF-same}, \rulename{FR}/FF-diff by variable elimination, and suspends
    non-patterns: the rules of \cref{fig:rules-pattern} with the prefix erased into levels. Patterns are read
    off levels. Elim prunes and lowers only strict occurrences, fails an
    occurrence of $F$ only by Nipkow's size argument, and suspends whatever a later binding could decide.
  \item \textbf{Check.} Every binding is scope-checked at its variable's level (NJK's binding rule with pruning),
    or taken back into E when a non-pattern decides it. Fast paths skip it before any \scm{∀°}, on plain terms and
    for top-level variables. Wake-up retries E in rounds while a round binds, and terminates when hnf does.
  \item \textbf{∀° and ⊃°.} \rulename{generic} runs its body one universe deeper with a new eigenvariable and discharges
    the residuals on exit (raise, close by ξ, retry, repeat). \rulename{augment} pushes closures on Γ. Dynamic relations
    backchain on hypotheses first, most recent first, then on their static body; a clause's free variables are
    shared, its \scm{∀} variables renamed per use.
  \item \textbf{Stores.} \scm{≠°} on λ-terms goes to N with a trial whose new variables are existential;
    D never holds a λ-term; two readings of $\forall$, generic (default) and extensional. \scm{absent°} looks into
    parts and watches flexible terms; types wait on flexible terms and are decided by their predicate.
  \item \textbf{Answers.} β-normal, η-short, binders renamed apart; residuals as \scm{(assert (≐ u v))},
    \scm{(deny ...)}, \scm{(absent? ...)}, closed with \scm{∀} when they keep an eigenvariable; never a free
    eigenvariable; ground answers are live terms.
\end{itemize}

\chapter{The book's examples in \texorpdfstring{\texttt{(aux kanren lambda)}}{(aux kanren lambda)}}\label{ch:examples}

\Cref{ch:module} presented the module \texttt{(aux kanren lambda)} rule by rule: the terms \scm{ƛ} and
\scm{@}, the levels that police the scope of eigenvariables,
the unifier for $\Llam$ with its delayed equations, and the goals \scm{∀°} and \scm{⊃°}. This chapter
runs that machinery on the programs of Part~I. Each example was introduced in an earlier chapter, where
we explained what it means in λProlog. Here we ask a narrower question: what does the module
\emph{do} with it? Which rule fires, why does the answer have the shape it has, and does it come out in
the order λProlog gives?

The examples follow the book's progression. We start with the terms themselves (α, β, shadowing), go
on to the two new goals (\scm{∀°}, then \scm{⊃°}, then both together in type inference), then to
programs over λ-tree syntax (substitution, de~Bruijn indices, prenex forms), higher-order programming
(\lp{mappred}, functional difference lists, Fischer's CPS transformation), and finally to the
unification problems of the book's Chapter~8 one by one. The last section takes the queries of Ma and
Friedman's 2021 paper~\cite{MaFriedman2021} and answers them as λProlog would, explaining each point
where the answers differ.

% =====================================================================================================
\section{How to read the examples}\label{sec:examples-howto}

Every example has four parts.
\begingroup\emergencystretch=2em
\begin{enumerate}
  \item The \textbf{λProlog original}, copied from the book, with its printed page
    \book{n}.\sidenote{The \enquote{p.N} in the labels of \texttt{examples.scm} are PDF pages; the prose uses printed pages
(see \cref{sec:intro-conventions}).}
  \item The \textbf{Scheme version}, copied verbatim from the curated script \texttt{examples.scm}, or
    from the programs it includes, \texttt{test/programs.scm}. In the repository both live under
    \texttt{doc/hol-kanren/prototype/}.\sidenote{After the merge, the programs move with the test suite to
    \texttt{src/test/} (\cref{sec:roadmap-merge}).}
  \item The \textbf{answers}, pasted verbatim from \texttt{examples.out}, the actual output of the
    script. A few examples that the script does not contain were run separately with the same imports;
    each of them says so.
  \item A \textbf{walk-through}: the rules of \cref{ch:module} that fire, the state changes they make,
    and how the order of the answers compares with λProlog.
\end{enumerate}
\endgroup

\pagebreak
\subsection{The dictionary}

The header of \texttt{examples.scm} gives the translation that every example uses:
\begin{scheme}
;;; The dictionary, λProlog → here:
;;;   x\ t                 (ƛ (x) t)               object binder by a Scheme binder (λ-tree syntax)
;;;   (F t)                (@ F t)                 application of a variable or an abstraction
;;;   pi x\ G              (∀° (x) G)              GENERIC
;;;   D => G               (⊃° (D) G)              AUGMENT, for a relation made by define-dynamic-relation
;;;   pi X\ H :- B         (∀ (X) (H ← B))         a hypothetical clause with its own variables
;;;   sigma X\ G           (fresh° (X) G)
;;;   G1, G2   G1 ; G2     (and° G1 G2)  (or° G1 G2)
;;; Object-level syntax is Scheme data: `abs M` is `(abs ,M), `app M N` is `(app ,M ,N).
\end{scheme}
Object-level syntax stays first-order Scheme data, as everywhere in \micro{}. Only the binders are
meta-level: \lp{abs x\ app x x} becomes \scm{`(abs ,(ƛ (x) `(app ,x ,x)))}, a list whose second element is
a λ-term. Inside the body of \scm{ƛ}, the Scheme variable \scm{x} is bound to a placeholder, so Scheme
scoping gives the object binders, shadowing included (\cref{sec:lambdatree-encoding}).

\subsection{Three ways to print an answer}

The script prints each query with a small macro, its label first:
\begin{scheme}
(define-syntax show
  (syntax-rules ()
    ((_ label e) (begin (display ";; ") (display label) (newline) (pretty-print e) (newline)))))
\end{scheme}
and uses three runners.
\begin{itemize}
  \item \scm{(°->list/datum g)} gives each answer as plain data, comparable with \scm{equal?} against the
    book's answers. Abstractions print as \scm{(ƛ (x y) body)}, with nested binders merged and renamed
    apart. Applications whose head is not an abstraction print as \scm{(@ h a ...)}. Unbound logic
    variables print as \scm{α}, \scm{β}, \scm{γ}, \dots, in order of first occurrence.
  \item \scm{(°->list #f g)} gives micro's own answer form, a closed expression
    \scm{(λ (α ...) constraint ... term)} (\cref{sec:micro-reify}). The empty list \scm{()} means
    that the query \emph{failed}. \scm{((λ () 'ok))} means one answer, with the ground value \scm{ok}. A
    residual equation prints as \scm{(assert (≐ u v))}.
  \item \scm{(°->pretty g)}, from \texttt{(aux kanren micro show)}, gives the αKanren-style view
    \scm{(α : constraint ...)}, which is the most readable one for suspended equations.
\end{itemize}
Many queries end in \scm{(=° q 'ok)} or \scm{(=° q 'solved)}. A query whose interesting part is a
success or a failure has no answer variable of its own, so these goals give it one to print.

\subsection{The rules, in one table}

The walk-throughs name the rules of \cref{ch:module}: the unification rules of
\cref{fig:module-rules}, variable elimination (\cref{def:module-elim}), the scope check
\textbf{Check} (\cref{def:module-check}), and the three search rules of \cref{sec:module-hh}. We recall
them here in one line each. Two notions are needed first.
\begin{itemize}
  \item The \emph{level} $\ell(X)$ of a logic variable is the number of \scm{∀°} goals around the point
    where it was created. The level $\ell(c)$ of an eigenvariable is the depth of the \scm{∀°} that
    created it, counting that \scm{∀°}. The names that the unifier opens below binders have level
    $\infty$. A variable of level $\ell$ may be bound only to a term whose eigenvariables have level
    $\le\ell$: it stands for an $\exists$ placed after the first $\ell$ universals of its branch.
  \item A \emph{pattern} for $F$ is a list of arguments that reduce to distinct eigenvariables, each of
    level greater than $\ell(F)$. This is Miller's condition \enquote{quantified inside the scope of
    $F$}, read off the levels (\cref{sec:hounif-llambda,sec:rules-levels,def:module-pattern}).
\end{itemize}
For instance, in \scm{(fresh° (q) (∀° (x) (fresh° (y) ...)))} we have $\ell(q)=0$, $\ell(x)=1$ and
$\ell(y)=1$: the list $(x)$ is a pattern for $q$ but not for $y$, and $y$ may be bound to $x$ while $q$
may not.

\begingroup\small
\begin{longtable}{@{}>{\raggedright\arraybackslash}p{0.17\fullwidthlen}>{\raggedright\arraybackslash}p{0.27\fullwidthlen}>{\raggedright\arraybackslash}p{0.5\fullwidthlen}@{}}
\toprule
rule & fires on & effect \\
\midrule\endhead
\rulename{Var}, \textbf{Check} & $X\doteq t$, $X$ unbound & micro binds $X:=t$ after its occurs check. Then the scope check $\mathrm{Elim}(\ell(X),-,\langle\rangle,t)$ fails on an eigenvariable of level $>\ell(X)$, \emph{lowers} the younger variables of $t$ to $\ell(X)$, and prunes flexible subterms. A non-pattern it cannot decide turns the binding back into a suspended equation. \\
\rulename{ξ} & $\lambda.b\doteq\lambda.b'$ & open both bodies with one new name $c$ (level $\infty$), unify them, and \textbf{Check} the bindings made under it \\
\rulename{η} & $\lambda.b\doteq t$, $t$ no abstraction & unify $b[c/0]$ with $(t\ c)$, $c$ new (level $\infty$); only when one side is an abstraction \\
\rulename{Rigid} & $h\,\bar u\doteq h'\,\bar v$ & same head and arity: unify the arguments; otherwise fail \\
\rulename{FF-same} & $F\,\bar a\doteq F\,\bar b$, both patterns & $\bar a=\bar b$: success; arities differ: fail; else $F:=\lambda\bar x.\,H\,\langle x_i\mid a_i=b_i\rangle$, $\ell(H)=\ell(F)$ \\
\rulename{FR} & $F\,\bar a\doteq t$, $\bar a$ a pattern, $t$ rigid or flexible with another head (the case \enquote{FF-diff}) & $F:=\lambda\bar x.\,\mathrm{Elim}(\ell(F),F,\bar a,t)$. In $t$: each $a_i$ becomes the $i$-th bound variable (projection); constants, and eigenvariables of level $\le\ell(F)$, are kept (imitation); an eigenvariable out of reach fails (8.J); a flexible pattern $G\,\bar b$ loses the arguments $F$ may not see (\emph{pruning}, 8.H); a younger variable is lowered to $\ell(F)$ and raised over the $a_i$ it may see (8.K); an occurrence of $F$ fails when no instance can make it smaller (8.I), and otherwise suspends \\
\rulename{FF-one} & $F\,\bar a\doteq G\,\bar b$, only $\bar b$ a pattern & \rulename{FR} on $G$ \\
\rulename{Delay} & no side a pattern & unless the sides are α-equal, suspend $\lambda L.u\doteq\lambda L.v$ in the store E ($L$: the names opened so far); every later binding retries it; at the end it prints as \scm{(assert (≐ u v))} \\
\rulename{Occurs} & $X\doteq t$ with $X$ inside a λ-term of $t$ & succeed if $t$ reduces to $X$, else eliminate $X$ with no arguments \\
\rulename{generic} & \scm{(∀° (x) g)} & a new eigenvariable at depth${}+1$; $g$ runs one level deeper; every answer gets the depth back \\
\rulename{discharge} & leaving \scm{∀°} (back to depth $d$) with residuals on $c$ & raise $Y:=Y'\,c$, $\ell(Y')=d$, for each variable $Y$ of the residual with $\ell(Y)\ge\ell(c)$; close the residual into $\lambda c.u\doteq\lambda c.v$; retry it; repeat while a raise brings $c$ into more constraints (\cref{sec:module-discharge}) \\
\rulename{augment} & \scm{(⊃° (D ...) g ...)} & push the clauses on Γ for $g$ only; every answer gets Γ back \\
\rulename{backchain} & a call of a dynamic relation & one flat \scm{or°}: the hypotheses, most recent first, then the static body \\
\bottomrule
\end{longtable}
\endgroup

% =====================================================================================================
\section{λ-tree syntax: α, β and shadowing}\label{sec:examples-alpha}

The first examples involve no search at all. They check that the terms behave as the book's
λ-terms do: equal up to the names of bound variables (\cref{sec:lambda-conversion}), with β as part of
equality, and with binders that obey lexical scope (\cref{sec:lambdatree-properties}).

\begin{scheme}
(show "λx.λy.(x . y), as data and as micro's rebuilding code"
      (list (°->list/datum (fresh° (q) (=° q (ƛ (x y) (cons x y)))))
            (°->list #f (fresh° (q) (=° q (ƛ (x y) (cons x y)))))))

(show "α-equivalence: λa.a = λb.b, and the β-redex (λx.g x x) a"
      (list (°->list #f (fresh° (q) (=° (ƛ (a) a) (ƛ (b) b))))
            (°->list/datum (fresh° (q f) (=° f (ƛ (x) `(g ,x ,x))) (=° q (@ f 'a))))))

(show "shadowing, λa.λa.a, printed λa.λa1.a1 (Ma–Friedman's \"true = false\" hazard)"
      (°->list/datum (fresh° (q) (=° q (ƛ (a) (ƛ (a) a))))))
\end{scheme}
\begin{answers}
;; λx.λy.(x . y), as data and as micro's rebuilding code
(((ƛ (x y) (x . y))) ((λ () (ƛ (x y) (cons x y)))))

;; α-equivalence: λa.a = λb.b, and the β-redex (λx.g x x) a
(((λ (α) α)) ((g a a)))

;; shadowing, λa.λa.a, printed λa.λa1.a1 (Ma–Friedman's "true = false" hazard)
((ƛ (a a1) a1))
\end{answers}

\begin{tutorial}[what happens]
\begin{enumerate}
  \item \scm{(ƛ (x y) (cons x y))} is built at once, before any goal runs. The body is a pair whose
    two cells hold the de~Bruijn indices 1 and 0, and the names \scm{x}, \scm{y} survive only as
    printing hints. \scm{=°} binds \scm{q} to it by \rulename{Var}: no eigenvariable is involved, so
    \textbf{Check} has nothing to do. The rebuilding code \scm{(λ () (ƛ (x y) (cons x y)))} is an
    expression that rebuilds the same term when evaluated, since \scm{ƛ} is bound in the module.
  \item \scm{(=° (ƛ (a) a) (ƛ (b) b))} is decided by \rulename{ξ}: both bodies are opened with the same
    new name $c$, and $c\doteq c$ holds. Nothing is bound, so \scm{q} stays free and prints as
    \scm{α}. In fact the two terms are already the same record up to the hint: locally nameless
    terms make α-equivalence a structural fact.
  \item \scm{(@ f 'a)} receives the logic variable \scm{f}, not its value: \scm{@} is a Scheme
    procedure and does not consult the state. So it builds the flexible application $f\,a$. Once \scm{f} is bound to $\abs{x}{g\,x\,x}$, the reifier's deep walk contracts the
    β-redex, and the answer is \scm{(g a a)}, as \lp{(x\ g x x) a} is in λProlog.
  \item λa.λa.a is the term λa.λb.b: the inner binder shadows the outer one, because the inner Scheme
    \scm{a} shadows the outer \scm{a}. The printer renames apart, so the answer reads
    \scm{(ƛ (a a1) a1)}: Church's \emph{false}, not \emph{true}. \Cref{sec:examples-mf} returns to
    this.
\end{enumerate}
\end{tutorial}

% =====================================================================================================
\section{Universal goals and the scope of eigenvariables}\label{sec:examples-scope}

\Cref{sec:hh-scope} introduced the simplest example of quantifier scope: \lp{sigma x\ pi y\ x = y} is
not provable, while \lp{pi y\ sigma x\ x = y} is. In the book's words:

\begin{lprolog}
?- sigma x\ pi y\ x = y.
\end{lprolog}
\enquote{This query is not provable because (pi y\textbackslash{} t = y) is not provable for any term t in which y
does not appear} \book{85}. \texttt{examples.scm} writes the two queries with an answer variable:
\begin{scheme}
;; ?- pi x\ (Y = x).        no                ∃Y∀x. Y = x
;; ?- pi x\ sigma Y\ Y = x.  yes               ∀x∃Y. Y = x
(show "∃Y∀x. Y = x  and  ∀x∃Y. Y = x"
      (list (°->list #f (fresh° (q) (∀° (x) (=° q x))))
            (°->list #f (fresh° (q) (∀° (x) (fresh° (y) (=° y x))) (=° q 'ok)))))

(show "escape through an alias (q = (y), y = x), and through micro's cons°: both fail"
      (list (°->list #f (fresh° (q) (∀° (x) (fresh° (y) (=° q (list y)) (=° y x)))))
            (°->list #f (fresh° (q) (∀° (x) (cons° x '() q))))))
\end{scheme}
\begin{answers}
;; ∃Y∀x. Y = x  and  ∀x∃Y. Y = x
(() ((λ () 'ok)))

;; escape through an alias (q = (y), y = x), and through micro's cons°: both fail
(() ())
\end{answers}

\begin{tutorial}[the levels at work]
Write $\ell$ for levels. In each query \scm{q} is created by the outer \scm{fresh°}, at depth 0, so
$\ell(q)=0$.
\begin{enumerate}
  \item \textbf{$\exists Y\forall x.\,Y=x$.} \rulename{generic} creates the eigenvariable $x$ with
    $\ell(x)=1$ and runs \scm{(=° q x)} at depth 1. micro binds $q:=x$, then the \texttt{post} hook runs
    \textbf{Check}: $x$ is an eigenvariable with $\ell(x)=1>\ell(q)=0$, so the binding fails. The
    query has no answer: \scm{()}.
  \item \textbf{$\forall x\exists Y.\,Y=x$.} Now \scm{y} is created inside the \scm{∀°}, so
    $\ell(y)=1$ and $y:=x$ passes \textbf{Check}. When the \scm{∀°} exits, the depth is set back to 0
    for the rest of the query, and \scm{(=° q 'ok)} gives the answer. The binding of \scm{y} is not
    part of the answer, which is what λProlog does for a variable quantified inside \lp{pi x\}.
  \item \textbf{The alias.} \scm{y} has level 1 again. The first goal binds $q:=(y)$. \textbf{Check}
    of this binding finds the younger variable $y$ inside a term bound to the older $q$, and
    \emph{lowers} it: $y:=y'$ with $\ell(y')=0$. This is condition (3) of Nadathur, Jayaraman and Kwon's
    binding rule, quoted in \cref{sec:hh-scope}: the variables of $t$ take the level of $X$. The second
    goal walks $y$ to $y'$ and binds $y':=x$, which \textbf{Check} now rejects, since
    $\ell(x)=1>0$. Without lowering, $x$ would escape into $q$ through $y$. The suite checks both orders
    of the two goals (A1/alias).
  \item \textbf{micro's \scm{cons°}.} \scm{cons°} is an ordinary relation of micro, defined with
    \scm{=°}. Since the check runs in \scm{=°}'s \texttt{post} hook, every relation built on \scm{=°}
    inherits it. Of the three prototypes, the one with its own unifier failed exactly here
    (\cref{sec:design-d1}).
\end{enumerate}
\end{tutorial}

The same check applies to logic variables that reach a clause. \Cref{sec:hh-link} quoted Nadathur,
Jayaraman and Kwon's example: with the clause $\forall z\,(q\,z\imp p\,(d\,z))$, the goal
$\exists x\,\forall y\,(q\,y\imp p\,x)$ must fail, while $\forall y\,\exists x\,(q\,y\imp p\,x)$
succeeds~\cite[p.\,16]{NadathurJayaramanKwon1995}. This is the suite's case A1/levels-reach-clauses,
shown in \cref{ex:module-njk} as well; we ran it separately:
\begin{scheme}
(show "∃x∀y(q y ⊃ p x) with ∀z(q z ⊃ p (d z))"
      (let ()
        (define-dynamic-relation (q° z))
        (define-dynamic-relation (p° x) (fresh° (z) (=° x `(d ,z)) (q° z)))
        (list (°->list #f (fresh° (x) (∀° (y) (⊃° ((q° y)) (p° x)))))
              (°->list #f (fresh° (q) (∀° (y) (fresh° (x) (⊃° ((q° y)) (p° x)))) (=° q 'ok))))))
\end{scheme}
\begin{answers}
;; ∃x∀y(q y ⊃ p x) with ∀z(q z ⊃ p (d z))
(() ((λ () 'ok)))
\end{answers}
In the first query $\ell(x)=0$ and $\ell(y)=1$. \rulename{backchain} on \scm{(p° x)} runs the static
body, which creates the clause's own \scm{z} at depth 1, so $\ell(z)=1$. The binding $x:=(d\,z)$ puts
the younger $z$ inside the older $x$, and \textbf{Check} lowers it: $z:=z'$ with $\ell(z')=0$. The call
\scm{(q° z')} then finds the hypothesis \scm{(q° y)}, and $z'\doteq y$ fails, since $\ell(y)=1>0$.
In the second query $x$ is created inside the \scm{∀°}, so $\ell(x)=1$; $z$ is not lowered, and
$z\doteq y$ succeeds. This is NJK's \enquote{new tagged instance} of a clause: its own variables are
created at the current level, at each use, and then obey the binding rule like any other variable.

% =====================================================================================================
\section{Hypothetical goals and the program context}\label{sec:examples-augment}

\Cref{ch:hh} introduced \rulename{augment}: to prove $D\imp G$, add $D$ to the program and prove
$G$. In the module the program context Γ is a field of the state. \scm{⊃°} extends it for the extent of
its body and gives every answer the old Γ back. Only relations made by \scm{define-dynamic-relation}
read Γ (\cref{sec:roadmap-lim-dynamic}); assuming a clause for any other relation raises an error.

\subsection{Clause order under augmentation}\label{sec:examples-order}

\enquote{λProlog uses the rule that when an implicational goal augments the current program, it does so by
adding the new clauses at the front of the program} \book{78}. With the program \lp{p 1}:
\begin{lprolog}
?- p 2 => p 3 => p X.
X = 3;
X = 2;
X = 1
?- (p 2 & p 3) => p X.
X = 2;
X = 3;
X = 1
\end{lprolog}
\begin{scheme}
;; p 1.
;; ?- p 2 => p 3 => p X.          X = 3; X = 2; X = 1
;; ?- (p 2 & p 3) => p X.         X = 2; X = 3; X = 1
(define-dynamic-relation (p° x) (=° x 1))
(show "clause order under augmentation (p.94)"
      (list (°->list/datum (fresh° (q) (⊃° ((p° 2)) (⊃° ((p° 3)) (p° q)))))
            (°->list/datum (fresh° (q) (⊃° ((p° 2) (p° 3)) (p° q))))))
\end{scheme}
\begin{answers}
;; clause order under augmentation (p.94)
((3 2 1) (2 3 1))
\end{answers}

\begin{tutorial}[the state of Γ, step by step]
\begin{enumerate}
  \item \scm{(define-dynamic-relation (p° x) (=° x 1))} makes \scm{p°} a dynamic relation whose static
    body is \scm{(=° x 1)}, the clause \lp{p 1}.
  \item First query. The outer \scm{⊃°} sets $\Gamma=[$\scm{p°}${}\mapsto(p\,2)]$. The inner one puts
    its clause in front: $\Gamma=[$\scm{p°}${}\mapsto(p\,3,\ p\,2)]$.
  \item The call \scm{(p° q)} is \rulename{backchain}: one flat disjunction
    $(q\doteq3)\vee(q\doteq2)\vee(q\doteq1)$, the hypotheses most recent first, then the static body.
    Each disjunct succeeds at once, so micro's interleaving delivers them in order: \scm{(3 2 1)}.
  \item Second query. One \scm{⊃°} assumes both clauses, which are kept in the order written:
    $\Gamma=[$\scm{p°}${}\mapsto(p\,2,\ p\,3)]$, and the answers are \scm{(2 3 1)}.
\end{enumerate}
Both orders agree with the book. The agreement is not an accident of the search: the dynamic relation
is built as \emph{one} flat \scm{or°}, so when every clause answers immediately, the answers come in
clause order. When the clauses do more work before answering, interleaving decides the order instead
(\cref{sec:examples-prenex}).
\end{tutorial}

\subsection{The course database (book Fig.~3.1)}

\Cref{sec:hh-db} presented the database and its four hypothetical queries:
\begin{lprolog}
fact (finished kim 102) & fact (finished dana 101).
fact (finished kim 210) & fact (finished dana 250).
fact (cs_major X) :-
   (fact (finished X 101); fact (finished X 102)),
   fact (finished X 250), fact (finished X 301).
fact (graduates X) :-
   (fact (finished X 101); fact (finished X 102)),
   (fact (finished X 210); fact (finished X 250)),
   fact (finished X 301).
false :- fact (finished X 210), fact (finished X 250).

?- fact (finished dana X) => false.
X = 210;
no
?- fact (finished X Y) => (fact (graduates X), fact (cs_major X)).
X = dana
Y = 301;
no
?- fact (finished kim X) => fact (finished kim Y) =>
   (fact (graduates kim), fact (cs_major kim)).
X = 250
Y = 301
?- fact (finished kim 250) => fact (finished kim 301) => false.
solved
\end{lprolog}
(Figure 3.1 without its type declarations, and the four sessions, \book{80--81}.) In \texttt{test/programs.scm} the program is one dynamic
relation \scm{fact°}, since every query assumes \lp{fact} clauses:
\begin{scheme}
(define-dynamic-relation (fact° e)
  (fresh° (x)
    (cond°
      ((=° e '(finished kim 102))) ((=° e '(finished dana 101)))
      ((=° e '(finished kim 210))) ((=° e '(finished dana 250)))
      ((=° e `(cs_major ,x))
       (or° (fact° `(finished ,x 101)) (fact° `(finished ,x 102)))
       (fact° `(finished ,x 250))
       (fact° `(finished ,x 301)))
      ((=° e `(graduates ,x))
       (or° (fact° `(finished ,x 101)) (fact° `(finished ,x 102)))
       (or° (fact° `(finished ,x 210)) (fact° `(finished ,x 250)))
       (fact° `(finished ,x 301))))))

(define-dynamic-relation (false°) (fresh° (x) (fact° `(finished ,x 210)) (fact° `(finished ,x 250))))
\end{scheme}
and the queries are
\begin{scheme}
(show "courses: which course would make Dana's record inconsistent"
      (°->list/datum (fresh° (q) (⊃° ((fact° `(finished dana ,q))) (false°)))))
(show "courses: who graduates in CS after one more course, and which"
      (°->list/datum (fresh° (q x y) (=° q (list x y)) (⊃° ((fact° `(finished ,x ,y))) (fact° `(graduates ,x)) (fact° `(cs_major ,x))))))
(show "courses: Kim's two courses (each order is a proof, each found twice; the book shows the first)"
      (sorted (°->list/datum (fresh° (q x y) (=° q (list x y))
                               (⊃° ((fact° `(finished kim ,x))) (⊃° ((fact° `(finished kim ,y))) (fact° '(graduates kim)) (fact° '(cs_major kim))))))))
(show "courses: no consistent way for Kim"
      (°->list/datum (fresh° (q) (⊃° ((fact° '(finished kim 250))) (⊃° ((fact° '(finished kim 301))) (false°))) (=° q 'solved))))
\end{scheme}
\begin{answers}
;; courses: which course would make Dana's record inconsistent
(210)

;; courses: who graduates in CS after one more course, and which
((dana 301))

;; courses: Kim's two courses (each order is a proof, each found twice; the book shows the first)
((250 301) (250 301) (301 250) (301 250))

;; courses: no consistent way for Kim
(solved)
\end{answers}

\begin{tutorial}[a logic variable inside a hypothesis]
Take the first query. \scm{q} is created before \scm{⊃°} runs, so the assumed clause
\scm{(fact° `(finished dana ,q))} is \emph{one} still unknown fact, not a schema: \scm{q} is a free
variable of the clause's closure, shared with the query (\book{86--88}; \rulename{backchain} in \cref{sec:module-hh}).
\begin{enumerate}
  \item \scm{(false°)} runs its body: \scm{x} is new, and \scm{(fact° `(finished ,x 210))} is a
    \rulename{backchain} on Γ first. The hypothesis unifies \scm{(finished dana q)} with
    \scm{(finished x 210)}: $x:=\cst{dana}$, $q:=210$.
  \item \scm{(fact° '(finished dana 250))} succeeds from the static body.
  \item The static clauses also prove \scm{(finished x 210)}, with $x:=\cst{kim}$. But then
    \scm{(finished kim 250)} must hold, and it is neither a fact nor an instance of the hypothesis,
    which is about Dana. So the only answer is \scm{(210)}, as in the book.
\end{enumerate}
In the third query there are four answers instead of the book's one: the book shows only the first
answer. They are proofs of the same goal. The symmetric pair appears because either hypothesis can
supply either course. Each answer appears twice because the disjunct
\scm{(or° (fact° `(finished ,x 210)) (fact° `(finished ,x 250)))} in \scm{graduates} has two proofs once
250 is assumed: from Kim's static 210, and from the hypothesis. The script sorts them, since the order
among them is the interleaving's. \Cref{sec:hh-db} reached the same count with the throwaway interpreter
of \cref{sec:hh-interp}.
\end{tutorial}

\subsection{Jars and bugs}

The example of \cref{sec:hh-generic} combines \rulename{generic} and \rulename{augment}:
\begin{lprolog}
sterile J :- pi x\ bug x => in x J => dead x.
dead B    :- heated J, in B J, bug B.
heated j.
\end{lprolog}
This is the program of \book{83}; its signature has no constant of type \lp{bug}. In Scheme:
\begin{scheme}
;; p.98-99:  sterile J :- pi x\ bug x => in x J => dead x.
;;           dead B :- heated J, in B J, bug B.        heated j.
(define-dynamic-relation (bug° b))
(define-dynamic-relation (in° b j))
(define-relation (heated° j) (=° j 'j))
(define-dynamic-relation (dead° b) (fresh° (j) (heated° j) (in° b j) (bug° b)))
(define-relation (sterile° j) (∀° (x) (⊃° ((bug° x)) (⊃° ((in° x j)) (dead° x)))))
(show "jars and bugs: ?- sterile Q." (°->list/datum (fresh° (q) (sterile° q))))
\end{scheme}
\begin{answers}
;; jars and bugs: ?- sterile Q.
(j)
\end{answers}

\begin{tutorial}[trace of \texttt{sterile Q}]
\begin{enumerate}
  \item \scm{(sterile° q)}: $\ell(q)=0$. \rulename{generic} creates the eigenvariable $x$, $\ell(x)=1$,
    and the depth becomes 1.
  \item Two \rulename{augment}s: $\Gamma=[$\scm{in°}${}\mapsto(\cst{in}\,x\,q)$, \scm{bug°}${}\mapsto(\cst{bug}\,x)]$.
    The clause \scm{(in° x j)} closes over the \scm{j} of \scm{sterile°}, which is \scm{q}.
  \item \scm{(dead° x)} has no hypothesis, so \rulename{backchain} runs its static body. It creates
    \scm{j'} at depth 1, and \scm{(heated° j')} binds $j':=\cst{j}$.
  \item \scm{(in° x j')} backchains on the hypothesis: it unifies $(x\ \cst{j})$ with $(x\ q)$, so
    $q:=\cst{j}$. \textbf{Check} accepts it: the constant \scm{j} contains no eigenvariable.
  \item \scm{(bug° x)} succeeds from the hypothesis. The \scm{∀°} exits: depth 0 and the old Γ come
    back, and there is no residual to discharge. The answer is \scm{(j)}.
\end{enumerate}
The proof is about an arbitrary bug: \scm{bug°} and \scm{in°} have no static clauses at all, and the only
facts about $x$ are the two hypotheses. Had a step tried to bind \scm{q} to something containing $x$,
\textbf{Check} would have rejected it.
\end{tutorial}

\subsection{Reverse with assumed clauses (book Fig.~3.4)}

\Cref{sec:hh-link} showed how an assumed clause can \enquote{tie} a variable of the goal:
\begin{lprolog}
reverse L K :-
   (rev nil K &
    (pi X\ pi L\ pi K\ rev (X::L) K :- rev L (X::K)))
        => rev L nil.
\end{lprolog}
(Figure 3.4, \book{87}). The first clause mentions the \lp{K} of \lp{reverse}, the second has its own
universally quantified \lp{X}, \lp{L}, \lp{K}. The Scheme version says the same with \scm{∀}:
\begin{scheme}
;; Figure 3.4 (p.104):
;; reverse L K :- (rev nil K & (pi X\ pi L\ pi K\ rev (X::L) K :- rev L (X::K))) => rev L nil.
;; the free K of `rev nil K` is the caller's, the clause's own X, L, K are renamed at each use.
(define-dynamic-relation (rev° l k))
(define-relation (reverse° l k)
  (⊃° ((rev° '() k) (∀ (X L K) ((rev° `(,X . ,L) K) ← (rev° L `(,X . ,K))))) (rev° l '())))
(show "reverse with assumed clauses, forwards and backwards"
      (list (°->list/datum (fresh° (q) (reverse° '(1 2 3) q)))
            (°->list/datum (take° 1 (fresh° (q) (reverse° q '(1 2 3)))))))
\end{scheme}
\begin{answers}
;; reverse with assumed clauses, forwards and backwards
(((3 2 1)) ((3 2 1)))
\end{answers}

\begin{tutorial}[two kinds of clause variables]
\scm{⊃°} expands each clause into a pair \scm{(rev° . procedure)}, the procedure taking the list of
arguments of a call. For \scm{(rev° '() k)} the procedure is \scm{(=° args (list '() k))}: \scm{k} is a
free variable of the closure, the caller's \scm{q}. For the \scm{∀} clause the procedure is
\scm{(fresh° (X L K) (=° args ...) (rev° L ...))}: its variables are created anew at each call, which is
how λProlog renames a clause's own variables (\cref{sec:module-hh}).

Forwards, \scm{(rev° '(1 2 3) '())} uses the second clause three times, building the accumulator
\scm{(3 2 1)}, and then \scm{(rev° '() '(3 2 1))} uses the first one, which binds $k:=(3\,2\,1)$.
Backwards, \scm{k} is \scm{(1 2 3)} and \scm{l} is unknown. The second clause guesses longer and longer
lists, and the first clause succeeds only when the accumulator equals \scm{(1 2 3)}, after three steps.
That is the only answer, but nothing tells the search so. In a separate run,
\scm{(take° 2 (fresh° (q) (reverse° q '(1 2 3))))} had not returned after 20 seconds, so the script
takes one answer.
\end{tutorial}

% =====================================================================================================
\section{Type inference: \texorpdfstring{$\forall$}{∀} and \texorpdfstring{$\supset$}{⊃} together}\label{sec:examples-typeof}

Type inference (\cref{sec:lambdatree-typeof,sec:apps-typeof}) is the book's standard example of a
universal goal and a hypothesis working together. To type $\abs{x}{M}$, take a new constant $x$,
assume that it has type $A$, and type the body:
\begin{lprolog}
typeof (app M N) A :- typeof M (arr B A), typeof N B.
typeof (abs M) (arr A B) :- pi x\ typeof x A => typeof (M x) B.

?- typeof (abs x\ abs y\ abs z\ app (app x z) (app y z)) Ty.
Ty = arr (arr T1 (arr T2 T3)) (arr (arr T1 T2) (arr T1 T3))
?- typeof (abs x\x) Ty.
Ty = arr T1 T1
?- typeof (abs x\ app x x) Ty.
no
\end{lprolog}
These are the clauses and the session of \book{190--191}. In \texttt{test/programs.scm}:
\begin{scheme}
; §7.4.3, type inference:  typeof (abs M) (arr A B) :- pi x\ typeof x A => typeof (M x) B.
(define-dynamic-relation (typeof° m t)
  (cond°
    ((fresh° (M N B) (=° m `(app ,M ,N)) (typeof° M `(arr ,B ,t)) (typeof° N B)))
    ((fresh° (R A B) (=° m `(abs ,R)) (=° t `(arr ,A ,B))
       (∀° (x) (⊃° ((typeof° x A)) (typeof° (@ R x) B)))))))
\end{scheme}
and the queries:
\begin{scheme}
(show "typeof S, I and λx.x x"
      (list (°->list/datum (fresh° (q) (typeof° `(abs ,(ƛ (x) `(abs ,(ƛ (y) `(abs ,(ƛ (z) `(app (app ,x ,z) (app ,y ,z)))))))) q)))
            (°->list/datum (fresh° (q) (typeof° `(abs ,(ƛ (x) x)) q)))
            (°->list/datum (fresh° (q) (typeof° `(abs ,(ƛ (x) `(app ,x ,x))) q)))))
(show "typeof backwards: the first terms of type i → i (interleaving finds them where DFS loops)"
      (°->list/datum (take° 3 (fresh° (q) (typeof° q '(arr i i))))))
\end{scheme}
\begin{answers}
;; typeof S, I and λx.x x
(((arr (arr α (arr β γ)) (arr (arr α β) (arr α γ)))) ((arr α α)) ())

;; typeof backwards: the first terms of type i → i (interleaving finds them where DFS loops)
((abs (ƛ (x) x))
 (abs (ƛ (x) (app (abs (ƛ (x1) x1)) x)))
 (abs (ƛ (x) (app (abs (ƛ (x1) x)) x))))
\end{answers}

\begin{tutorial}[typing \texttt{abs x\textbackslash x}, state by state]
\begin{enumerate}
  \item \scm{(typeof° m q)} with $m={}$\scm{(abs ,(ƛ (x) x))}. \rulename{backchain}: Γ has no hypothesis
    yet, so only the static \scm{cond°} runs. The \scm{app} branch fails on \scm{(=° m `(app ,M ,N))}.
  \item The \scm{abs} branch binds $R:=\abs{x}{x}$ and $q:=(\cst{arr}\ A\ B)$; $R$, $A$, $B$ have
    level 0.
  \item \rulename{generic}: an eigenvariable \scm{x}, call it $c$, with $\ell(c)=1$.
    \rulename{augment}: $\Gamma=[$\scm{typeof°}${}\mapsto(\cst{typeof}\ c\ A)]$.
  \item The goal \scm{(typeof° (@ R x) B)}. As in the previous example, \scm{(@ R x)} receives the
    logic variable $R$, not its value, so it is the flexible term $R\,c$. Since $R$ is bound, every
    walk reduces it to $c$.
  \item \rulename{backchain} on Γ first: the hypothesis unifies $(c\ B)$ with $(c\ A)$, so
    $A\doteq B$, a variable--variable binding with no eigenvariable in it (micro binds the older
    variable to the newer one). Then the static clauses fail: $c$ is neither
    an \scm{app} nor an \scm{abs} list.
  \item The \scm{∀°} exits: depth 0, Γ empty. The answer is $q=(\cst{arr}\ A\ A)$, printed
    \scm{(arr α α)}.
\end{enumerate}
For $\abs{x}{x\,x}$ the hypothesis is used twice: once for the head, giving $A:=(\cst{arr}\ B'\ B)$,
once for the argument, giving $B'\doteq(\cst{arr}\ B'\ B)$, which fails micro's ordinary occurs check.
For S, three nested \scm{∀°} goals put three hypotheses in Γ, all for \scm{typeof°}, most recent first. The
answer has the book's shape, with α, β, γ in place of T1, T2, T3. \Cref{sec:module-typeof} prints the
same steps with a probe, for $\abs{x}{\abs{y}{x}}$ and for $\abs{x}{x\,x}$.
\end{tutorial}

\begin{kanren}[running type inference backwards]
The book never asks \lp{typeof Ty (arr i i)} with the term unknown, and for a good reason: depth-first
search tries the \lp{app} clause first, which asks again for a term of unknown type, forever. micro's
interleaving search alternates between the branches, so it finds terms of a given type, and even
(term, type) pairs. We ran the suite's case A3/backwards separately:
\begin{scheme}
(show "typeof backwards: (term type) pairs"
      (°->list/datum (take° 4 (fresh° (q m t) (=° q (list m t)) (typeof° m t)))))
\end{scheme}
\begin{answers}
;; typeof backwards: (term type) pairs
(((abs (ƛ (x) x)) (arr α α))
 ((abs (ƛ (x) (abs (ƛ (x1) x1)))) (arr α (arr β β)))
 ((abs (ƛ (x) (abs (ƛ (x1) x)))) (arr α (arr β α)))
 ((app (abs (ƛ (x) x)) (abs (ƛ (x) x))) (arr α α)))
\end{answers}
The interesting step is $R\,c\doteq c$ with $R$ unknown: a pattern, solved by \rulename{FR} with the
projection $R:=\abs{x}{x}$. The binder of the answer is named \scm{x} after the eigenvariable of the
\scm{∀°}, the name the clause gave it. This is one of the places where the module finds answers that
λProlog misses (\cref{sec:roadmap-lim-order}).
\end{kanren}

% =====================================================================================================
\section{Object-level substitution}\label{sec:examples-subst}

\Cref{sec:lambdatree-beta} explained the central idea of λ-tree syntax: object-level substitution is
meta-level β. \Cref{sec:lambdatree-copy} gave the alternative that avoids β altogether, \lp{copy}
clauses assumed under \lp{pi}. Both run here.

\subsection{Call-by-name evaluation}

\begin{lprolog}
cbn (abs R) (abs R).
cbn (app M N) V :- cbn M (abs R), cbn (R N) V.

?- cbn (app (abs x\ abs w\w) (app (abs x\ app x x) (abs x\ app x x))) V.
\end{lprolog}
gives \lp{V = abs w\w}, while call-by-value loops (Figure 7.3, \book{182--183}).
\begin{scheme}
; Figure 7.3, call-by-name evaluation: β of the object language is β of the meta-language.
(define-relation (cbn° m v)
  (cond°
    ((fresh° (R) (=° m `(abs ,R)) (=° v m)))
    ((fresh° (M N R) (=° m `(app ,M ,N)) (cbn° M `(abs ,R)) (cbn° (@ R N) v)))))
\end{scheme}
\begin{scheme}
(show "call by name: (λx.λw.w) Ω"
      (let* ((ω `(abs ,(ƛ (x) `(app ,x ,x)))) (Ω `(app ,ω ,ω)))
        (°->list/datum (fresh° (q) (cbn° `(app (abs ,(ƛ (x) `(abs ,(ƛ (w) w)))) ,Ω) q)))))
\end{scheme}
\begin{answers}
;; call by name: (λx.λw.w) Ω
((abs (ƛ (w) w)))
\end{answers}

\begin{tutorial}[β on demand]
In the \scm{app} branch, \scm{(@ R N)} is built from the logic variable \scm{R}, as in the type-inference
trace: it is a flexible application.
\scm{(cbn° M `(abs ,R))} evaluates the operator \scm{(abs ,(ƛ (x) ...))}, which is already a value:
$R:=\abs{x}{\cst{abs}\,(\abs{w}{w})}$. When the second call \scm{(cbn° (@ R N) v)} unifies its first
argument with \scm{`(abs ,R')}, the walk of $R\,N$ contracts the redex to \scm{(abs ,(ƛ (w) w))}. The
argument $N=\Omega$ is dropped by that β step and never evaluated, so the search terminates, as
call-by-name should. No rule of the unifier is needed beyond β: every step is a pattern-free,
first-order unification on a reduced term.
\end{tutorial}

\subsection{\texorpdfstring{\texttt{subst R N (R N)}}{subst R N (R N)}: dynamic \texorpdfstring{$\Llam$}{L-lambda}}

The one-line \lp{subst R N (R N)} is the book's example of a clause outside $\Llam$ that is
\enquote{meant only to compute the effect of substitution through reduction} \book{225--226}, since
\lp{R} is expected to be known when the clause is used (\cref{sec:hounif-dynamic}). The module accepts the
non-pattern $R\,N$ in either order of the goals. This query is the suite's case A4/beta, run separately,
with $R$ bound \emph{after} the equation:
\begin{scheme}
(show "subst R N (R N): R bound after"
      (°->list/datum (fresh° (q r) (=° q (@ r `(abs ,(ƛ (y) y)))) (=° r (ƛ (x) `(app ,x ,x))))))
\end{scheme}
\begin{answers}
;; subst R N (R N): R bound after
((app (abs (ƛ (y) y)) (abs (ƛ (y) y))))
\end{answers}
The first goal binds $q$ to the flexible term $r\,(\cst{abs}\ldots)$. The argument is a list, not an
eigenvariable, so this is no pattern, but no \scm{∀°} has run, so \textbf{Check} has nothing to
forbid and the binding stands. The second goal binds $r$, and the reifier's deep walk contracts the
redex. Had the first binding involved an eigenvariable that only the value of $r$ could keep or drop,
\textbf{Check} would have turned it into a suspended equation (\rulename{Delay}), retried when $r$ is
bound (\cref{sec:examples-nonpattern,ex:module-forget}).

\subsection{Copy and substitution without β (book Fig.~7.10)}

\begin{lprolog}
copy (app M N) (app P Q) :- copy M P, copy N Q.
copy (abs M) (abs N) :- pi x\ copy x x => copy (M x) (N x).
subst M T S :- pi x\ copy x T => copy (M x) S.

?- copy (abs x\ abs y\ app y x) M.
\end{lprolog}
\enquote{Solving for all the variables yields the unique answer substitution} \lp{M = abs v\abs u\app u v}
(Figure 7.10 and the trace, \book{199--200}).
\begin{scheme}
; Figure 7.10, copy and subst on untyped λ-terms.
(define-dynamic-relation (copy° m n)
  (cond°
    ((fresh° (M N P Q) (=° m `(app ,M ,N)) (=° n `(app ,P ,Q)) (copy° M P) (copy° N Q)))
    ((fresh° (M N) (=° m `(abs ,M)) (=° n `(abs ,N))
       (∀° (x) (⊃° ((copy° x x)) (copy° (@ M x) (@ N x))))))))
(define-relation (subst° m t s) (∀° (x) (⊃° ((copy° x t)) (copy° (@ m x) s))))
\end{scheme}
\begin{scheme}
(show "copy (p.215), and subst (x\\ app x x) (abs y\\ y)"
      (list (°->list/datum (fresh° (q) (copy° `(abs ,(ƛ (x) `(abs ,(ƛ (y) `(app ,y ,x))))) q)))
            (°->list/datum (fresh° (q) (subst° (ƛ (x) `(app ,x ,x)) `(abs ,(ƛ (y) y)) q)))))
\end{scheme}
\begin{answers}
;; copy (p.215), and subst (x\ app x x) (abs y\ y)
(((abs (ƛ (x) (abs (ƛ (x1) (app x1 x))))))
 ((app (abs (ƛ (y) y)) (abs (ƛ (y) y)))))
\end{answers}

\begin{tutorial}[the book's trace, replayed]
The book traces this query with the eigenvariables $c$ and $d$ \book{199}. In the module:
\begin{enumerate}
  \item The \scm{abs} clause binds $M_0:=\abs{x}{\cst{abs}\,(\abs{y}{\cst{app}\ y\ x})}$ and
    $q:=(\cst{abs}\ N_0)$. \rulename{generic} gives $c$ at level 1, \rulename{augment} assumes
    \scm{(copy° c c)}, and the goal is \scm{copy°}~$(\cst{abs}\,(\abs{y}{\cst{app}\ y\ c}))\ (N_0\,c)$.
  \item The \scm{abs} clause again. The equation $N_0\,c\doteq(\cst{abs}\ N_1)$ is a pattern for $N_0$:
    \rulename{FR} binds $N_0:=\abs{x}{\cst{abs}\,(N_1'\,x)}$. $N_1$ was created at level 1, inside the
    first \scm{∀°}, so it is lowered and raised over $c$: $N_1:=N_1'\,c$ (the 8.K step of
    \cref{sec:examples-unif}). The book's $M_1$ is our $N_0$.
  \item A second \rulename{generic} gives $d$ at level 2, and \scm{(copy° d d)} is assumed. The goal is
    \scm{copy°}~$(\cst{app}\ d\ c)\ (N_1'\,c\,d)$. The \scm{app} clause binds
    $N_1'\,c\,d\doteq(\cst{app}\ P\ Q)$ by \rulename{FR}, and the hypotheses give $P=d$, $Q=c$, by
    projection.
  \item Each \scm{∀°} exits with no residual. The answer is
    $\cst{abs}\,(\abs{x}{\cst{abs}\,(\abs{x_1}{\cst{app}\ x_1\ x})})$, the book's
    \lp{abs v\abs u\app u v} up to the names.
\end{enumerate}
The binders are named \scm{x} and \scm{x1}: the unifier names the abstractions it builds after the
eigenvariables they abstract, and both eigenvariables come from the clause's \scm{(∀° (x) ...)}. The
printer renames the second one apart.
\end{tutorial}

\subsection{Existential generalisation (book Fig.~7.11)}

Figure 7.11 specifies substitution into first-order formulas with \lp{copyterm} and \lp{copyform}
clauses \book{201}. Because \lp{subst} is a relation, it can run backwards and abstract a term out of a
formula: \enquote{given a term T and a formula S, it is possible to compute an abstraction M such that
(subst M T S) is derivable} \book{202}:
\begin{lprolog}
?- subst M (f a) ((p (f a)) && (q a (f a))).
M = x\ (p x) && (q a x)
M = x\ (p x) && (q a (f a))
M = x\ (p (f a)) && (q a x)
M = x\ (p (f a)) && (q a (f a))
\end{lprolog}
These are the four existential generalisations of the formula on $f\,a$ (\cref{sec:lambdatree-copy}).
\texttt{test/programs.scm} has the clauses for the constants used here:
\begin{scheme}
; Figure 7.11, copy and subst on first-order formulas.
(define-dynamic-relation (copyterm° t s)
  (cond°
    ((=° t 'a) (=° s 'a))
    ((fresh° (x u) (=° t `(f ,x)) (=° s `(f ,u)) (copyterm° x u)))
    ((fresh° (x y u v) (=° t `(g ,x ,y)) (=° s `(g ,u ,v)) (copyterm° x u) (copyterm° y v)))))
(define-relation (copyform° b d)
  (cond°
    ((fresh° (B C D E op) (=° b `(,op ,B ,C)) (=° d `(,op ,D ,E)) (or° (=° op 'and) (=° op 'or) (=° op 'imp))
       (copyform° B D) (copyform° C E)))
    ((fresh° (B D op) (=° b `(,op ,B)) (=° d `(,op ,D)) (or° (=° op 'all) (=° op 'some))
       (∀° (y) (⊃° ((copyterm° y y)) (copyform° (@ B y) (@ D y))))))
    ((fresh° (x u) (=° b `(p ,x)) (=° d `(p ,u)) (copyterm° x u)))
    ((fresh° (x y u v) (=° b `(q ,x ,y)) (=° d `(q ,u ,v)) (copyterm° x u) (copyterm° y v)))))
(define-relation (substform° m t n) (∀° (x) (⊃° ((copyterm° x t)) (copyform° (@ m x) n))))
\end{scheme}
\begin{scheme}
(show "existential generalization: the book's four answers, in its order"
      (°->list/datum (fresh° (q) (substform° q '(f a) '(and (p (f a)) (q a (f a)))))))
\end{scheme}
\begin{answers}
;; existential generalization: the book's four answers, in its order
((ƛ (x) (and (p x) (q a x)))
 (ƛ (x) (and (p x) (q a (f a))))
 (ƛ (x) (and (p (f a)) (q a x)))
 (ƛ (x) (and (p (f a)) (q a (f a)))))
\end{answers}

\begin{tutorial}[how an unknown abstraction is built, piece by piece]
\begin{enumerate}
  \item \scm{substform°} creates the eigenvariable $x$ (level 1), assumes \scm{(copyterm° x '(f a))},
    and calls \scm{copyform°} on $q\,x$, a flexible pattern, against the formula.
  \item The first \scm{copyform°} clause matches \scm{(op B C)} on the left: $q\,x\doteq(\cst{op}\ B\ C)$
    is solved by \rulename{FR}. The variables $\cst{op}$, $B$, $C$ were created at level 1, so
    elimination lowers them and raises them over $x$: $q:=\abs{x}{(\cst{op}'\,x\ \ B'\,x\ \ C'\,x)}$.
    The next goal unifies the formula \scm{(and ...)} with \scm{(op D E)}, so $\cst{op}'\,x\doteq\cst{and}$
    gives $\cst{op}':=\abs{x}{\cst{and}}$ by imitation, and the \scm{or°} then accepts \scm{and}. The two
    halves are copied in turn.
  \item For the first half, the \scm{p} clause gives $B'\,x\doteq(\cst{p}\ X)$ and then
    \scm{(copyterm° (X' x) '(f a))}. \rulename{backchain} tries the hypothesis \emph{first}:
    $X'\,x\doteq x$ gives $X':=\abs{x}{x}$ by projection, which is the answer $p\,x$. Then the static
    clauses: $X'\,x\doteq(\cst{f}\ X'')$ and \scm{(copyterm° (X'' x) 'a)}, where only the constant
    \scm{a} clause applies, and $X'':=\abs{x}{a}$ by imitation: the answer $p\,(f\,a)$.
  \item The second half does the same for $q\,a\,(f\,a)$. The occurrence of \scm{a} has only the
    static answer, since the hypothesis copies $x$ to $f\,a$, not to $a$.
\end{enumerate}
The order is the book's: the hypothesis comes first for each occurrence, and the first occurrence
varies slowest.
\end{tutorial}

% =====================================================================================================
\section{Recursion over binders: de Bruijn terms and prenex forms}\label{sec:examples-binders}

\subsection{Translating to and from de Bruijn syntax}

\Cref{sec:lambdatree-debruijn} translated λ-tree syntax into de Bruijn's nameless terms by assuming,
for each binder crossed, the depth at which it was crossed:
\begin{lprolog}
trans D (abs M) (ab P) :- pi c\ depth D c => (E is D + 1, trans E (M c) P).
trans D (app M N) (ap P Q) :- trans D M P, trans D N Q.
trans D X (deb E) :- depth N X, E is (D - N).

?- trans 1 (abs x\ app x (abs y\ app x (abs w\ app w x))) D.
D = ab (ap (deb 1) (ab (ap (deb 2) (ab (ap (deb 1) (deb 3))))))
?- trans 1 P (ab (ap (deb 1) (ab (ap (deb 2) (ab (ap (deb 1) (deb 3))))))).
P = abs W1\ app W1 (abs W2\ app W1 (abs W3\ app W3 W1))
\end{lprolog}
These are the clauses and the session of \book{193}. micro has no \lp{is}, so the Scheme version
computes with \scm{project°}:
\begin{scheme}
; §7.4.4, translation to and from de Bruijn syntax, with the `depth` assumptions.
(define-dynamic-relation (depth° n x))
(define-relation (trans° d t b)
  (cond°
    ((fresh° (M P) (=° t `(abs ,M)) (=° b `(ab ,P))
       (∀° (c) (⊃° ((depth° d c)) (fresh° (e) (project° ((d* d)) (=° e (+ d* 1))) (trans° e (@ M c) P))))))
    ((fresh° (M N P R) (=° t `(app ,M ,N)) (=° b `(ap ,P ,R)) (trans° d M P) (trans° d N R)))
    ((fresh° (n e) (=° b `(deb ,e)) (depth° n t) (project° ((d* d) (n* n)) (=° e (- d* n*)))))))
\end{scheme}
\begin{scheme}
(show "de Bruijn, both directions"
      (list (°->list/datum (fresh° (q) (trans° 1 `(abs ,(ƛ (x) `(app ,x (abs ,(ƛ (y) `(app ,x (abs ,(ƛ (w) `(app ,w ,x))))))))) q)))
            (°->list/datum (fresh° (q) (trans° 1 q '(ab (ap (deb 1) (ab (ap (deb 2) (ab (ap (deb 1) (deb 3))))))))))))
\end{scheme}
\begin{answers}
;; de Bruijn, both directions
(((ab (ap (deb 1) (ab (ap (deb 2) (ab (ap (deb 1) (deb 3))))))))
 ((abs (ƛ (c) (app c (abs (ƛ (c1) (app c (abs (ƛ (c2) (app c2 c)))))))))))
\end{answers}
Both answers are the book's terms. The backward one is \lp{abs W1\ app W1 (abs W2\ app W1 (abs W3\ app W3 W1))}
with \scm{c}, \scm{c1}, \scm{c2} for \lp{W1}, \lp{W2}, \lp{W3}: the unifier named the abstractions after
the eigenvariable \scm{c} of the clause.

The book traces a smaller query step by step, \lp{trans 1 (abs x\ abs y\ abs z\ y) P}, and ends with
\enquote{the instantiation of P with (ab (ab (ab (deb 2))))} \book{193--194}. We ran it separately:
\begin{scheme}
(define t3 `(abs ,(ƛ (x) `(abs ,(ƛ (y) `(abs ,(ƛ (z) y)))))))
(show "trans 1 (abs x\\ abs y\\ abs z\\ y) P"
      (°->list/datum (fresh° (q) (trans° 1 t3 q))))
\end{scheme}
\begin{answers}
;; trans 1 (abs x\ abs y\ abs z\ y) P
((ab (ab (ab (deb 2)))))
\end{answers}

\begin{tutorial}[the book's trace as state changes]
Read \texttt{depth} for the level of the \scm{∀°}, and Γ for the \scm{depth°} hypotheses.
\begin{center}
\begin{tabular}{@{}lll@{}}
\toprule
goal & depth & Γ for \scm{depth°} (most recent first) \\
\midrule
\lp{trans 1 (abs x\ abs y\ abs z\ y) P} & 0 & --- \\
\lp{trans 2 (abs y\ abs z\ y) P1} & 1 & $(1\ c_1)$ \\
\lp{trans 3 (abs z\ c2) P2} & 2 & $(2\ c_2),\ (1\ c_1)$ \\
\lp{trans 4 c2 P3} & 3 & $(3\ c_3),\ (2\ c_2),\ (1\ c_1)$ \\
\bottomrule
\end{tabular}
\end{center}
(We write $c_1,c_2,c_3$ for the three eigenvariables, all hinted \scm{c}.) In the last goal only the
third clause applies, since $c_2$ is neither an \scm{abs} nor an \scm{app} list. \scm{(depth° n c2)}
backchains on Γ, most recent first: $(3\ c_3)$ and $(1\ c_1)$ fail on $c_3\neq c_2$ and $c_1\neq c_2$,
while $(2\ c_2)$ gives $n=2$, and \scm{project°} computes $e=4-2=2$. Each \scm{∀°} then exits,
restoring the depth and Γ. The book's \lp{depth c 1} is the same fact with its arguments in the
order of the trace; the clause and our code write \lp{depth D c}.
\end{tutorial}

In the backward direction the first argument of \scm{trans°} is $M\,c$ with $M$ unknown, a pattern.
The \scm{app} clause solves $M\,c\doteq(\cst{app}\ M_1\ N_1)$ by \rulename{FR}, lowering $M_1,N_1$ and
raising them over $c$. The \scm{deb} clause then meets \scm{(depth° n (M1' c))}: the hypothesis
\scm{(depth° 1 c)} gives $M_1'\,c\doteq c$, a projection. Every equation in this program is a pattern.

\subsection{Prenex normal forms (book Fig.~7.8)}\label{sec:examples-prenex}

\Cref{sec:lambdatree-prenex} showed that \lp{prenex} is a relation, not a function:
\begin{lprolog}
?- prenex ((all x\ q x x) && (all z\ all y\ q z y)) P.
\end{lprolog}
\enquote{will generate the following five answer substitutions for P} \book{197}:
\begin{lprolog}
all z\ all y\ (q z z) && (q z y)
all x\ all z\ all y\ (q x x) && (q z y)
all z\ all x\ (q x x) && (q z x)
all z\ all x\ all y\ (q x x) && (q z y)
all z\ all y\ all x\ (q x x) && (q z y)
\end{lprolog}
The program of \texttt{test/programs.scm} transcribes Figure 7.8 in the book's order, one \scm{cond°}
branch per clause, except that clauses differing only in their connective or quantifier (such as
\lp{prenex (B && C) D} and \lp{prenex (B !! C) D}) share one branch through an \scm{or°}. We show the
main relation and the first two of the seven \scm{merge°} branches:
\begin{scheme}
(define-relation (prenex° b d)
  (cond°
    ((atom° b) (=° d b))
    ((fresh° (a) (=° b `(neg ,a)) (atom° a) (=° d b)))
    ((fresh° (op B C U V) (=° b `(,op ,B ,C)) (or° (=° op 'and) (=° op 'or)) (prenex° B U) (prenex° C V) (merge° `(,op ,U ,V) d)))
    ((fresh° (q B D) (=° b `(,q ,B)) (=° d `(,q ,D)) (or° (=° q 'all) (=° q 'some)) (∀° (x) (prenex° (@ B x) (@ D x)))))))
(define-relation (merge° b d)
  (cond°
    ((fresh° (B C D) (=° b `(and (all ,B) (all ,C))) (=° d `(all ,D))
       (∀° (x) (merge° `(and ,(@ B x) ,(@ C x)) (@ D x)))))
    ((fresh° (q B C D) (=° b `(and (,q ,B) ,C)) (=° d `(,q ,D)) (or° (=° q 'all) (=° q 'some))
       (∀° (x) (merge° `(and ,(@ B x) ,C) (@ D x)))))
\end{scheme}
\begin{scheme}
(show "prenex is a relation: the book's five answers (here in the order 1 3 2 4 5)"
      (°->list/datum (fresh° (q) (prenex° `(and (all ,(ƛ (x) `(q ,x ,x))) (all ,(ƛ (z) `(all ,(ƛ (y) `(q ,z ,y)))))) q))))
\end{scheme}
\begin{answers}
;; prenex is a relation: the book's five answers (here in the order 1 3 2 4 5)
((all (ƛ (x) (all (ƛ (x1) (and (q x x) (q x x1))))))
 (all (ƛ (x) (all (ƛ (x1) (and (q x1 x1) (q x x1))))))
 (all (ƛ (x) (all (ƛ (x1) (all (ƛ (x2) (and (q x x) (q x1 x2))))))))
 (all (ƛ (x) (all (ƛ (x1) (all (ƛ (x2) (and (q x1 x1) (q x x2))))))))
 (all (ƛ (x) (all (ƛ (x1) (all (ƛ (x2) (and (q x2 x2) (q x x1)))))))))
\end{answers}

\begin{tutorial}[matching the answers to the book's]
Rename the binders: our second answer, $\forall x\forall x_1.\,(q\,x_1\,x_1)\wedge(q\,x\,x_1)$, is the
book's third, \lp{all z\ all x\ (q x x) && (q z x)}; our third is the book's second. The other three
are in the book's position. So all five forms are found, in the order 1~3~2~4~5.

Why the difference? \scm{merge°} has three clauses that apply to a conjunction of two universals: the
one that merges both quantifiers into one (\lp{(all B) && (all C)}), and the two that pull out only the
left or only the right one. λProlog tries them depth first, so all the answers of the first clause come
before those of the second. micro interleaves the branches of \scm{cond°}, so an answer that a later
branch reaches in fewer steps can come first (\cref{sec:horn-dfs,sec:micro-search}). The set of
answers is the book's, and each answer is found once. The unique prenex form of the book's first example
\book{196} comes out unique too (we ran it separately):
\begin{scheme}
(define f1 `(or (all ,(ƛ (x) `(and (and (p ,x) (all ,(ƛ (y) `(q ,x ,y)))) (p (f ,x))))) (p a)))
(show "prenex: the book's formula with a unique prenex form"
      (°->list/datum (fresh° (q) (prenex° f1 q))))
\end{scheme}
\begin{answers}
;; prenex: the book's formula with a unique prenex form
((all (ƛ (x) (all (ƛ (x1) (or (and (and (p x) (q x x1)) (p (f x))) (p a)))))))
\end{answers}
(We read the book's \lp{&&} as left-associative.) The answer is the book's
\lp{all x\ all y\ ((p x) && (q x y) && (p (f x))) !! (p a)}. Every equation in this program is a pattern: \scm{(@ B x)} and \scm{(@ D x)} apply a variable to the
eigenvariable of the enclosing \scm{∀°}.
\end{tutorial}

% =====================================================================================================
\section{Higher-order programming}\label{sec:examples-hoprog}

\subsection{\texorpdfstring{\texttt{mappred}}{mappred}}

\Cref{ex:hoprog-mappred} introduced \lp{mappred} (Figures 5.3--5.4):
\begin{lprolog}
mappred P nil nil.
mappred P (X::L) (Y::K) :- P X Y, mappred P L K.
age bob 23 & age sue 24 & age ned 23.

?- mappred (x\y\ age x y) (ned::bob::sue::nil) L.
L = (23::23::24::nil)
?- mappred (x\y\ age y x) (23::24::nil) K.
K = (bob::sue::nil);
K = (ned::sue::nil)
\end{lprolog}
These are the clauses of \book{127} and the session of \book{129}. A predicate argument is a Scheme
procedure here: \scm{age°} itself for
\lp{x\y\ age x y}, which is η-equal to \lp{age}, and a Scheme closure for the swapped version.
\begin{scheme}
; §5.3, Figures 5.3-5.4 (p.143): mappred with a relation argument, a Scheme closure here.
(define-relation (age° n a) (cond° ((=° n 'bob) (=° a 23)) ((=° n 'sue) (=° a 24)) ((=° n 'ned) (=° a 23))))
(define-relation (mappred° p l k)
  (cond°
    ((null° l) (null° k))
    ((fresh° (x xs y ys) (=° l `(,x . ,xs)) (=° k `(,y . ,ys)) (p x y) (mappred° p xs ys)))))
\end{scheme}
\begin{scheme}
(show "mappred (p.143, p.145)"
      (list (°->list/datum (fresh° (q) (mappred° age° '(ned bob sue) q)))
            (°->list/datum (fresh° (q) (mappred° (λ (x y) (age° y x)) '(23 24) q)))))
\end{scheme}
\begin{answers}
;; mappred (p.143, p.145)
(((23 23 24)) ((bob sue) (ned sue)))
\end{answers}
Both answers, and their order, are the book's. Nothing of the λ layer is involved: this is the
higher-order programming that any Scheme-hosted miniKanren has for free (\cref{sec:hoprog-free}). The
price is the one stated in \cref{sec:roadmap-lim-dynamic}: a predicate is a closure, not a term, so a program cannot
unify predicates, quantify over them (\lp{pi rev\}), or assume clauses for a predicate variable.

\subsection{Palindromes with functional difference lists (book Fig.~5.11)}

\Cref{sec:hoprog-fdl} represented a list by an abstraction over its tail, so that both ends can be
reached in one unification step:
\begin{lprolog}
palindrome (fdl x\x).
palindrome (fdl x\ Y::x).
palindrome (fdl x\ Y::(F (Y::x))) :- palindrome (fdl F).

?- palindrome (fdl x\ a::b::c::b::a::x).
solved
?- palindrome (fdl x\ a::b::a::(F x)).
F = x\ x;
F = x\ b::a::x;
F = x\ a::b::a::x;
F = x\ Y::a::b::a::x;
F = x\ Z::Z::a::b::a::x;
F = x\ Y::Z::Y::a::b::a::x;
F = x\ Y::Z::Z::Y::a::b::a::x;
F = x\ Y::Z::U::Z::Y::a::b::a::x;
F = x\ Y::Z::U::U::Z::Y::a::b::a::x
\end{lprolog}
This is the program and the session of \book{142--143}. The third clause is outside $\Llam$: $F$ is applied to \lp{Y::x}, which is not an
eigenvariable.
\begin{scheme}
; §5.8.2, Figure 5.11 (p.158): a functional difference list is an abstraction over its tail, and the
; palindrome clauses destructure one by unification, `x\ Y::(F (Y::x))` being no pattern: it waits.
(define-relation (palindrome° f)
  (cond°
    ((=° f `(fdl ,(ƛ (x) x))))
    ((fresh° (Y) (=° f `(fdl ,(ƛ (x) `(,Y . ,x))))))
    ((fresh° (Y F) (=° f `(fdl ,(ƛ (x) `(,Y . ,(@ F `(,Y . ,x)))))) (palindrome° `(fdl ,F))))))
\end{scheme}
\begin{scheme}
(show "functional difference lists: palindromes (p.158), the book's first nine answers in its order"
      (list (°->list/datum (fresh° (q) (palindrome° `(fdl ,(ƛ (x) `(a b c b a . ,x)))) (=° q 'solved)))
            (°->list/datum (take° 9 (fresh° (q) (palindrome° `(fdl ,(ƛ (x) `(a b a . ,(@ q x))))))))))
\end{scheme}
\begin{answers}
;; functional difference lists: palindromes (p.158), the book's first nine answers in its order
((solved)
 ((ƛ (x) x)
  (ƛ (x) (b a . x))
  (ƛ (x) (a b a . x))
  (ƛ (x) (α a b a . x))
  (ƛ (x) (α α a b a . x))
  (ƛ (x) (α β α a b a . x))
  (ƛ (x) (α β β α a b a . x))
  (ƛ (x) (α β γ β α a b a . x))
  (ƛ (x) (α β γ γ β α a b a . x))))
\end{answers}

\begin{tutorial}[a non-pattern that waits for the recursive call]
Follow the third clause on the query \lp{fdl x\ a::b::a::(F x)}, our \scm{q}.
\begin{enumerate}
  \item micro unifies the two lists \scm{(fdl ...)} cell by cell, and hands the two abstractions to
    the hook. \rulename{ξ} opens them with one name $c$ (level $\infty$):
    $a\mathbin{::}b\mathbin{::}a\mathbin{::}(q\,c)\doteq Y\mathbin{::}(F\,(Y\mathbin{::}c))$.
  \item The pairs are unified cell by cell: $Y:=a$, and $b\mathbin{::}a\mathbin{::}(q\,c)\doteq F\,(a\mathbin{::}c)$. The right side
    applies $F$ to a list, so it is no pattern; the left side is rigid. No side is a pattern, so
    \rulename{FR} cannot fire, and \rulename{Delay} suspends the equation, closed over the opened name
    $c$: $\abs{c}{b\mathbin{::}a\mathbin{::}(q\,c)}\doteq\abs{c}{F\,(a\mathbin{::}c)}$.
  \item The recursive call \scm{(palindrome° `(fdl ,F))} binds $F$ in each of its branches. Each
    binding retries the suspended equation, which has become a pattern equation for $q$, or a
    failure. The second clause, for instance, binds $F:=\abs{x}{Y'\mathbin{::}x}$. The equation becomes
    $b\mathbin{::}a\mathbin{::}(q\,c)\doteq Y'\mathbin{::}a\mathbin{::}c$, so $Y':=b$ and $q\,c\doteq c$, and \rulename{FR} gives
    $q:=\abs{x}{x}$: the first answer, the list \lp{a::b::a}.
\end{enumerate}
The first clause, $F:=\abs{x}{x}$, fails on $b\doteq a$; deeper recursions give the longer answers.
The answers are the book's nine, in the book's order, with α, β, γ for the book's variables in
order of first occurrence (so the fifth answer's α is the book's \lp{Z}). A
procedure that rejected non-patterns, as ELPI does by default, would lose this program (\cref{sec:design-space}).
\end{tutorial}

% =====================================================================================================
\section{Fischer's CPS transformation (book Fig.~10.6)}\label{sec:examples-cps}

\Cref{sec:apps-cps} presented the book's relational CPS transformation. It has no clause for
variables: going under a binder, \lp{phi} assumes how the bound variable translates.
\begin{lprolog}
ftrans (abs V) (abs k\ k @ U) :- phi (abs V) U.
ftrans (M @ N) (abs k\ P @ (abs m\ Q @ (abs n\ m @ k @ n))) :-
   ftrans M P, ftrans N Q.
phi (abs M) (abs k\ abs x\ (P x) @ k) :-
   pi x\ pi y\ ftrans x (abs k\ k @ y) => ftrans (M x) (P y).

?- ftrans ((abs u\u) @ (abs u\u)) F.
F = (abs W1\ (abs (W2\ W2 @ abs W3\ abs W4\ abs (W5\ W5 @ W4) @ W3)) @
     (abs W2\ (abs (W3\ W3 @ abs W4\ abs W5\ abs (W6\ W6 @ W5) @ W4)) @
     (abs W3\ W2 @ W1 @ W3))).
\end{lprolog}
This is Figure 10.6 and the session of \book{257}. The object application \lp{M @ N} is the list
\scm{(app M N)}:
\begin{scheme}
; §10.3.2, Figure 10.6 (p.273): Fischer's CPS transformation; object application `M @ N` is `(app M N)`.
; There is no clause for variables: going under a binder assumes ftrans x (abs k\ k @ y).
(define-dynamic-relation (ftrans° m f)
  (cond°
    ((fresh° (V U) (=° m `(abs ,V)) (=° f `(abs ,(ƛ (k) `(app ,k ,U)))) (phi° m U)))
    ((fresh° (M N P Q) (=° m `(app ,M ,N)) (ftrans° M P) (ftrans° N Q)
       (=° f `(abs ,(ƛ (k) `(app ,P (abs ,(ƛ (m) `(app ,Q (abs ,(ƛ (n) `(app (app ,m ,k) ,n))))))))))))))
(define-relation (phi° m u)
  (fresh° (M P) (=° m `(abs ,M)) (=° u `(abs ,(ƛ (k) `(abs ,(ƛ (x) `(app ,(@ P x) ,k))))))
    (∀° (x y) (⊃° ((ftrans° x `(abs ,(ƛ (k) `(app ,k ,y))))) (ftrans° (@ M x) (@ P y))))))
\end{scheme}
\begin{scheme}
(show "ftrans ((abs u\\u) @ (abs u\\u)) F" (°->list/datum (fresh° (q) (ftrans° `(app (abs ,(ƛ (u) u)) (abs ,(ƛ (u) u))) q))))
\end{scheme}
\begin{answers}
;; ftrans ((abs u\u) @ (abs u\u)) F
((abs (ƛ (k)
         (app (abs (ƛ (k1)
                      (app k1
                           (abs (ƛ (k2)
                                   (abs (ƛ (x) (app (abs (ƛ (k3) (app k3 x))) k2))))))))
              (abs (ƛ (m)
                      (app (abs (ƛ (k1)
                                   (app k1
                                        (abs (ƛ (k2)
                                                (abs (ƛ (x) (app (abs (ƛ (k3) (app k3 x))) k2))))))))
                           (abs (ƛ (n) (app (app m k) n))))))))))
\end{answers}

\begin{tutorial}[reading the answer against the book's term]
Map \scm{k} to \lp{W1} and \scm{m} to the outer \lp{W2}. The operator part
\scm{(abs (ƛ (k1) (app k1 (abs (ƛ (k2) (abs (ƛ (x) (app (abs (ƛ (k3) (app k3 x))) k2))))))))} is the book's
\lp{abs (W2\ W2 @ abs W3\ abs W4\ abs (W5\ W5 @ W4) @ W3)} with \scm{k1}, \scm{k2}, \scm{x}, \scm{k3} for
\lp{W2}, \lp{W3}, \lp{W4}, \lp{W5}. The final continuation \scm{(abs (ƛ (n) (app (app m k) n)))} is
\lp{abs W3\ W2 @ W1 @ W3}. So the answer is exactly the book's term up to α. The suite compares it with
\scm{equal?} against this datum (A7/cps).

The steps that involve the λ layer are in \scm{phi°}, called on \lp{abs u\u}:
\begin{enumerate}
  \item $M:=\abs{u}{u}$, and the second argument of \scm{phi°}, the variable \scm{u} of the clause, is bound to
    \scm{`(abs ,(ƛ (k) `(abs ,(ƛ (x) `(app ,(@ P x) ,k)))))}. Here \scm{(@ P x)} is built inside \scm{ƛ}, with
    \scm{x} the placeholder of an object binder: it is the flexible term $P$ applied to a bound
    variable.
  \item \rulename{generic} twice: eigenvariables $x$, $y$. \rulename{augment}: the clause
    \scm{(ftrans° x `(abs ,(ƛ (k) `(app ,k ,y))))}.
  \item The goal \scm{(ftrans° (@ M x) (@ P y))}, that is, \scm{ftrans°}~$x\ (P\,y)$, since $M\,x$
    reduces to $x$. \rulename{backchain} on Γ first: $P\,y\doteq\cst{abs}\,(\abs{k}{\cst{app}\ k\ y})$
    is a pattern, $\ell(y)=2>\ell(P)$. \rulename{FR} projects $y$ and imitates the rest:
    $P:=\abs{y}{\cst{abs}\,(\abs{k}{\cst{app}\ k\ y})}$. The static clauses fail on $x$, which is
    neither an \scm{abs} nor an \scm{app} list.
  \item Once $P$ is bound, the subterm $P\,x$ of \scm{u}, under the object binder $x$, reduces to
    \scm{(abs (ƛ (k3) (app k3 x)))}, the book's \lp{abs (W5\ W5 @ W4)}.
\end{enumerate}
\end{tutorial}

% =====================================================================================================
\section{Higher-order pattern unification}\label{sec:examples-unif}

The book's Chapter~8 develops pattern unification on small problems with explicit quantifier prefixes
(\cref{ch:hounif}). In the module a prefix is written with goals: a variable created by \scm{fresh°}
outside a \scm{∀°} is an outer $\exists$, an eigenvariable is a $\forall$. So
$\exists F\,\forall x\,\forall y\,[\dots]$ is \scm{(fresh° (F) (∀° (x y) ...))}. We use \scm{q} for the
first variable, since it is the one the runners print. The labels 8.A--8.M are those of
\cref{sec:hounif-traces}.

\subsection{Flex-rigid and flex-flex}

\begin{scheme}
;; ?- pi x\ F x = g x x.                     F = x\ g x x
;; ?- pi x\ pi y\ F x y = G y x.             F = x\ y\ H y x,  G = H
;; ?- pi x\ pi y\ F x y = F y x.             F = x\ y\ H
(show "flex-rigid and flex-flex"
      (list (°->list/datum (fresh° (q) (∀° (x) (=° (@ q x) `(g ,x ,x)))))
            (°->list/datum (fresh° (q f g) (=° q (list f g)) (∀° (x y) (=° (@ f x y) (@ g y x)))))
            (°->list/datum (fresh° (q) (∀° (x y) (=° (@ q x y) (@ q y x)))))))
\end{scheme}
\begin{answers}
;; flex-rigid and flex-flex
(((ƛ (x) (g x x))) (((ƛ (x y) (@ α y x)) α)) ((ƛ (x y) α)))
\end{answers}
\begin{tutorial}[three rules, three answers]
\begin{enumerate}
  \item $q\,x\doteq g\,x\,x$: the left side is a pattern ($\ell(x)=1>\ell(q)=0$), the right side is
    rigid. \rulename{FR} replaces each $x$ by the bound variable: $q:=\abs{x}{g\,x\,x}$.
  \item $f\,x\,y\doteq g\,y\,x$: two patterns with different heads, so \rulename{FR} on $f$ (the case FF-diff). Elimination of $f$
    keeps $g$ with both its arguments, since both are among $f$'s, and nothing needs pruning:
    $f:=\abs{x\,y}{g\,y\,x}$, and $g$ stays free. The book's $H$ is our $g$, printed \scm{α}.
  \item $q\,x\,y\doteq q\,y\,x$: \rulename{FF-same}. The arguments agree in no position, so
    $q:=\abs{x\,y}{H}$: a constant function (\book{222}; \cref{sec:hounif-ff}).
\end{enumerate}
\end{tutorial}

\subsection{Example 8.H: pruning}

\begin{lprolog}
∀f ∀g ∃U ∃V ∀w ∀x ∀y [(f (U x y)) = (f (g (V y w)))]
\end{lprolog}
\enquote{variable elimination becomes applicable. Using it first generates the pruning substitution
λy λw (V′ y) for V and subsequently the substitution λx λy (g (V′ y)) for U} \book{223--224}. We
solved it by hand in \cref{ex:hounif-8h}.
\begin{scheme}
;; Example 8.H (p.239-240): ?- pi w\ pi x\ pi y\ f (U x y) = f (g (V y w)).
;; U = x\ y\ g (V' y),  V = y\ w\ V' y
(show "Example 8.H, pruning: (U V)" (°->list/datum (fresh° (q u v) (=° q (list u v)) (∀° (w x y) (=° `(f ,(@ u x y)) `(f (g ,(@ v y w))))))))
\end{scheme}
\begin{answers}
;; Example 8.H, pruning: (U V)
(((ƛ (x y) (g (@ α y))) (ƛ (y w) (@ α y))))
\end{answers}
\begin{tutorial}[8.H with levels]
The prefix gives $\ell(u)=\ell(v)=0$ and $\ell(w)=1$, $\ell(x)=2$, $\ell(y)=3$. Here \scm{f} is the
symbol \scm{'f} at the head of a list, a constant of the signature. The book quantifies $f$ and $g$
universally at the front of the prefix, before $U$ and $V$; an eigenvariable that every variable may
see behaves like a constant, so nothing is lost.
\begin{enumerate}
  \item micro unifies the two lists \scm{(f ...)}: \scm{f} against \scm{f}, then the hook gets
    $u\,x\,y\doteq g\,(v\,y\,w)$. The left side is a pattern, so \rulename{FR} runs
    $\mathrm{Elim}(0,u,(x,y),\,g\,(v\,y\,w))$.
  \item The walk keeps the constant \scm{g}. It meets the pattern $v\,y\,w$: $y$ is among $u$'s
    arguments and is kept; $w$ is not, and $\ell(w)=1>0$, so it must go. \emph{Pruning} binds
    $v:=\abs{y\,w}{v'\,y}$, with $\ell(v')=0$.
  \item The result $g\,(v'\,y)$ is abstracted over $u$'s arguments: $u:=\abs{x\,y}{g\,(v'\,y)}$.
\end{enumerate}
The answer is the book's mgu, with \scm{α} for $V'$. The binders of $v$'s value are named \scm{y} and
\scm{w}, after the eigenvariables of the query.
\end{tutorial}

\subsection{Example 8.K: delayed raising}

\begin{lprolog}
∀f ∃U ∀x ∀y ∀z ∃V [(U x) = (f V)]
\end{lprolog}
\enquote{we would replace V with the term (V′ x), where V′ is a new variable quantified at the same
place in the prefix as U and would produce the unifying substitution λx (f (V′ x)) for U}
\book{225}; see \cref{ex:hounif-8k}.
\begin{scheme}
;; Example 8.K (p.240-241): ?- pi x\ pi y\ pi z\ sigma V\ U x = f V.      U = x\ f (V' x)
(show "Example 8.K, delayed raising" (°->list/datum (fresh° (q) (∀° (x y z) (fresh° (v) (=° (@ q x) `(f ,v)))))))
\end{scheme}
\begin{answers}
;; Example 8.K, delayed raising
((ƛ (x) (f (@ α x))))
\end{answers}
Here $\ell(q)=0$ and $v$ is created after three \scm{∀°}, so $\ell(v)=3$. Elimination of $q$ over the
pattern $(x)$ meets $v$, a variable younger than $q$. It is \emph{lowered} to $q$'s level and
\emph{raised} over those arguments of $q$ that $v$ may see: $\ell(x)=1\le\ell(v)$, so
$v:=v'\,x$ with $\ell(v')=0$, and $q:=\abs{x}{f\,(v'\,x)}$. No raising over $y$ and $z$ ever happens:
the levels record the information that the book proposes to keep \enquote{with each existentially
quantified variable}, and elimination reads it only when it is needed.

\subsection{Failures: occurs check, scope (8.G, 8.I, 8.J)}

\begin{scheme}
;; Examples 8.G, 8.I, 8.J: X = f X;  U x y = g (U y w);  U x y = g w.      no, no, no
(show "Examples 8.G (occurs), 8.I (occurs under a flexible head), 8.J (scope)"
      (list (°->list #f (fresh° (q) (=° q `(f ,q))))
            (°->list #f (fresh° (q) (∀° (w x y) (=° (@ q x y) `(g ,(@ q y w))))))
            (°->list #f (fresh° (q) (∀° (w x y) (=° (@ q x y) `(g ,w)))))))

;; ∃F∀x. F = x fails; ∃F∀x. F x = x gives F = λz.z.
(show "scope: F = x versus F x = x"
      (list (°->list #f (fresh° (q) (∀° (x) (=° q x))))
            (°->list/datum (fresh° (q) (∀° (x) (=° (@ q x) x))))))
\end{scheme}
\begin{answers}
;; Examples 8.G (occurs), 8.I (occurs under a flexible head), 8.J (scope)
(() () ())

;; scope: F = x versus F x = x
(() ((ƛ (x) x)))
\end{answers}
\begin{itemize}
  \item 8.G, $X\doteq f\,X$, is first-order: micro's own occurs check rejects it.
  \item 8.I, $U\,x\,y\doteq g\,(U\,y\,w)$ (\cref{ex:hounif-8i}): elimination of $U$ meets $U\,y\,w$
    below the rigid constant $g$, with as many rigid arguments as $U\,x\,y$ has. No instance of $U$ can
    make the right side as small as the left, so this is an occurs failure (Nipkow's size argument,
    case (6a) of \cref{def:module-elim}).
  \item 8.J, $U\,x\,y\doteq g\,w$ (\cref{ex:hounif-8j}): $w$ is not among $U$'s arguments and
    $\ell(w)=1>\ell(U)=0$. It appears \enquote{in a position that is not the argument of a flexible
    term} \book{224}, so no pruning can remove it: a scope failure.
  \item $\exists F\forall x.\,F=x$ is the first query of \cref{sec:examples-scope} again, while
    $\exists F\forall x.\,F\,x=x$ is a pattern problem with the projection $F:=\abs{x}{x}$ as its
    answer. The first variable cannot \emph{be} $x$, but it can be a function that \emph{returns}
    it.
\end{itemize}

\subsection{Non-patterns: suspended, reported, woken}\label{sec:examples-nonpattern}

Two of the book's problems have no most general unifier. Example 8.M, $\forall x\,(F\,x\,x=x)$, has the
two solutions $\abs{u\,v}{u}$ and $\abs{u\,v}{v}$ \book{208}. Example 8.A,
$\forall a\,\forall g\,\exists F\,[(F\,a)=(g\,a\,a)]$, has four \book{213}
(\cref{ex:hounif-nomgu}). Without Huet's procedure the module cannot enumerate them, as Teyjus~1 does.
It does what Teyjus~2 does: it suspends the equation and reports it (dynamic $\Llam$,
\cref{sec:hounif-dynamic}).
\begin{scheme}
;; Example 8.M (p.224): pi x\ F x x = x has two incomparable unifiers: it waits, and is reported.
;; Example 8.A (p.229): F a = g a a, four unifiers: likewise (no Huet enumeration).
(show "non-patterns: 8.M and 8.A, suspended; 8.M woken by F := λa.λb.b"
      (list (°->pretty (fresh° (q) (∀° (x) (=° (@ q x x) x))))
            (°->pretty (fresh° (q) (=° (@ q 'a) '(g a a))))
            (°->list/datum (fresh° (q) (∀° (x) (=° (@ q x x) x)) (=° q (ƛ (a b) b))))))
\end{scheme}
\begin{answers}
;; non-patterns: 8.M and 8.A, suspended; 8.M woken by F := λa.λb.b
(((α : (= (ƛ (x) (@ α x x)) (ƛ (x) x))))
 ((α : (= (@ α a) (g a a))))
 ((ƛ (a b) b)))
\end{answers}
\begin{tutorial}[from suspension to answer]
\begin{enumerate}
  \item $q\,x\,x\doteq x$ repeats $x$, so it is no pattern: \rulename{Delay} stores it in E. When the
    \scm{∀°} exits, \rulename{discharge} closes it over $x$, since a residual that mentions $x$ stands
    for $\forall x\,(u=v)$. It is then retried, and stays suspended. The answer is $q$ unconstrained
    except for this equation, printed as \scm{(= (ƛ (x) (@ α x x)) (ƛ (x) x))}: no eigenvariable
    appears free in an answer.
  \item In the second query \scm{'a} is a constant, not an eigenvariable, so $q\,a$ is no pattern
    either. This is 8.A with \lp{a} in the signature, and it is reported as it stands.
  \item In the third query the suspended 8.M equation is woken by the binding
    $q:=\abs{a\,b}{b}$: $(\abs{a\,b}{b})\,x\,x$ reduces to $x$, and the equation holds. The answer is
    the value of $q$, with no residual left.
\end{enumerate}
A program can thus solve a non-pattern problem by \emph{generating} candidates for the unknown
function with ordinary goals. The suspended equation then acts as a test. The book's form of 8.A, with
$a$ universal and $F$ inside its scope, shows \rulename{discharge} at work. This is the suite's case A5/non-pattern,
run separately:
\begin{scheme}
(show "8.A under ∀ a, raised at exit"
      (°->list #f (fresh° (q) (∀° (a) (fresh° (f) (=° (@ f a) `(g ,a ,a)))) (=° q 'ok))))
\end{scheme}
\begin{answers}
;; 8.A under ∀ a, raised at exit
((λ (α)
    (assert (≐ (ƛ (a) (@ α a a)) (ƛ (a) (cons 'g (cons a (cons a '()))))))
    'ok))
\end{answers}
Here $f$ is created inside the scope of $a$, so $\ell(f)=1$, and $f\,a$ is no pattern for it ($a$ is
not \emph{younger} than $f$). At the exit of \scm{∀°}, \rulename{discharge} raises $f:=f'\,a$ with
$\ell(f')=0$, Miller's raising (\cref{sec:lambda-raising}), and closes the equation into
$\abs{a}{f'\,a\,a}\doteq\abs{a}{g\,a\,a}$, over $f'$, which prints as \scm{α}.
\end{tutorial}

\subsection{Occurrences that must not be pruned}

The subtle part of dynamic $\Llam$ is what elimination may do \emph{inside} the arguments of a
non-pattern. The book's advice is to \enquote{defer further consideration of this equation until such
time that substitutions resulting from solving other parts of the problem have altered its status}
\book{226}. Two examples of \texttt{examples.scm} show why deferring is necessary.
\begin{scheme}
;; Dynamic L_λ (p.241-242): G may drop its argument, so nothing inside it is pruned or lowered. The
;; equation ∃q∃G∃H∀y. q = G (H y) waits, closed over y at ∀ exit, and is solved by G := λa.k. Likewise
;; pi y\ F a = F y holds with any constant F: it waits, then F := λz.c solves it, and F := λz.z does not.
(show "non-strict occurrences: q = G (H y), and F a = F y"
      (list (°->pretty (fresh° (q G H) (∀° (y) (=° q (@ G (@ H y))))))
            (°->list/datum (fresh° (q G H) (∀° (y) (=° q (@ G (@ H y)))) (=° G (ƛ (a) 'k))))
            (°->pretty (fresh° (q) (∀° (y) (=° (@ q 'a) (@ q y)))))
            (°->list/datum (fresh° (q) (∀° (y) (=° (@ q 'a) (@ q y))) (=° q (ƛ (z) 'c))))
            (°->list/datum (fresh° (q) (∀° (y) (=° (@ q 'a) (@ q y))) (=° q (ƛ (z) z))))))
\end{scheme}
\begin{answers}
;; non-strict occurrences: q = G (H y), and F a = F y
(((α : (= (ƛ (y) α) (ƛ (y) (@ β (@ γ y))))))
 (k)
 ((α : (= α (ƛ (y) (@ α a)))))
 ((ƛ (z) c))
 ())
\end{answers}
\begin{tutorial}[why \texttt{q = G (H y)} must wait]
All of $q$, $G$, $H$ have level 0, and $\ell(y)=1$.
\begin{enumerate}
  \item micro binds $q:=G\,(H\,y)$, and \textbf{Check} looks at the bound term. $G\,(H\,y)$ is no
    pattern: its argument is a flexible term, not an eigenvariable. Inside that argument is $y$, which
    $q$ may not see. Pruning $H$'s argument ($H:=\abs{z}{H'}$) would be a guess: if $G$ ignores its
    argument, $y$ disappears with no pruning at all, so pruning would lose solutions. The occurrence of
    $y$ is not \emph{strict} (\cref{def:module-elim} and the tutorial after it, after Dowek, Hardin, Kirchner and
    Pfenning~\cite{DowekHardinKirchnerPfenning1996}).
  \item So \textbf{Check} takes the binding back (\scm{μkanren-state-forget}) and suspends the equation
    $q\doteq G\,(H\,y)$. At the exit of \scm{∀°} it is closed over $y$, and prints as
    $\abs{y}{q}\doteq\abs{y}{G\,(H\,y)}$.
  \item In the second query $G:=\abs{a}{k}$ wakes it: $\abs{y}{q}\doteq\abs{y}{k}$, so $q:=k$, and the
    answer is \scm{(k)}.
\end{enumerate}
The equation $q\,a\doteq q\,y$ is similar. It holds for every constant function $q$, so failing it as an
occurs error, or pruning, would be wrong. Only the right side $q\,y$ is a pattern, so \rulename{FF-one}
runs \rulename{FR} on it, with the sides swapped; elimination meets $q$ itself below binders only, and
suspends. At the exit the equation is closed over $y$ into $\abs{y}{q\,y}\doteq\abs{y}{q\,a}$, whose left
side η-contracts to $q$: that is the printed \scm{(= α (ƛ (y) (@ α a)))}. $q:=\abs{z}{c}$ satisfies it; $q:=\abs{z}{z}$ would need $a=y$, and fails.
\end{tutorial}

\begin{pitfall}[a residual is a condition, not a proof]
An answer with a suspended equation means: \emph{this} is a solution \emph{if} the residual has one.
The module does not decide that, since general higher-order unification is undecidable
(\cref{sec:lambda-undecidable}). Several residuals may even be jointly unsatisfiable, and the answer
would still be printed (\cref{sec:roadmap-lim-huet}, the case limits/no-huet). Treat \scm{(assert (≐ ...))} in an answer
as a question returned to the program. Generating candidates, as in the 8.M query above, is one way
to answer it.
\end{pitfall}

% =====================================================================================================
\section{Ma and Friedman 2021, answered by λProlog's semantics}\label{sec:examples-mf}

\Cref{sec:lkanren-mf21} described Ma and Friedman's proposal~\cite{MaFriedman2021}. Their unifier
treats a logic variable under binders as a \emph{static closure} that may refer to the binders by
position, so that \scm{(== (λ (a) X) (λ (b) b))} succeeds with $X$ = \scm{d0}, the level of the
binder. \Cref{sec:lkanren-counter} gave three counter-examples that compare this reading with
pattern unification and with nominal unification. The module takes the other side: a logic variable
outside a binder cannot depend on it, as in λProlog. This section runs the paper's queries and states
each difference. All the outputs below are from \texttt{examples.out}, except the ones marked as run
separately.

The paper does not number its code boxes. We call them E1--E12 in their order of appearance, with
the paper's page; \enquote{Example 1--3} are the paper's own numbered examples of its Section~4.

\subsection{Where the answers agree}

\enlargethispage{2\baselineskip}
The paper's E2 (p.\,2) and E12 (p.\,7):
\begin{answers}
> (run 1 X
    (== (λ (a) a) (λ (b) b)))
'(_0)
>(run 1 F
   (all (A)
     (== A (F A))))
'((lambda (_0) _0))
\end{answers}
(the paper's transcripts). In the module, run separately (the second query is also in
\texttt{examples.out}):
\begin{scheme}
(show "E2: λa.a = λb.b;  E8: λx.λy.x y = λa.λa.a a"
      (list (°->list #f (fresh° (q) (=° (ƛ (a) a) (ƛ (b) b))))
            (°->list #f (fresh° (q) (=° (ƛ (x y) (@ x y)) (ƛ (a a) (@ a a)))))))
(show "E12: (run 1 F (all (A) (== A (F A))))"
      (°->list/datum (fresh° (q) (∀° (a) (=° a (@ q a))))))
\end{scheme}
\begin{answers}
;; E2: λa.a = λb.b;  E8: λx.λy.x y = λa.λa.a a
(((λ (α) α)) ())

;; E12: (run 1 F (all (A) (== A (F A))))
((ƛ (a) a))
\end{answers}
E2 is α-equivalence of closed terms, decided by \rulename{ξ}. E12 is the pattern $F\,a\doteq a$ under
\rulename{generic}, solved by projection, exactly as the paper says: \enquote{the == operator
recognizes β-equivalence} (p.\,7). E8, $\abs{x}{\abs{y}{x\,y}}\doteq\abs{a}{\abs{a}{a\,a}}$, is the
paper's second \enquote{false positive} of binder-merging (p.\,4). Here it fails, as it must: the right
side is $\abs{a}{\abs{b}{b\,b}}$, and \rulename{Rigid} meets $x$ against $b$, two different opened
names.

\subsection{Where the answers differ}

\begin{scheme}
(show "λa.X = λb.b fails (the paper: X = d0); λa.λa.X = λc.λd.c fails; ∀A. A = F A gives F = λA.A"
      (list (°->list #f (fresh° (q) (=° (ƛ (a) q) (ƛ (b) b))))
            (°->list #f (fresh° (q) (=° (ƛ (a) (ƛ (a) q)) (ƛ (c) (ƛ (d) c)))))
            (°->list/datum (fresh° (q) (∀° (a) (=° a (@ q a)))))))
\end{scheme}
\begin{answers}
;; λa.X = λb.b fails (the paper: X = d0); λa.λa.X = λc.λd.c fails; ∀A. A = F A gives F = λA.A
(() () ((ƛ (a) a)))
\end{answers}
and, run separately, the paper's Example~3, the three counter-examples of \cref{sec:lkanren-counter}
(the first is the paper's conjunction, the third is shown again with two ground variants), and the
paper's \texttt{assume} query:
\begin{scheme}
(show "their Example 3: λa.λb.X = λc.λd.Y"
      (°->list/datum (fresh° (q X Y) (=° q (list X Y)) (=° (ƛ (a b) X) (ƛ (c d) Y)))))
(show "counter-example 1: λa.X = λb.b ∧ λc.X = λd.d"
      (°->list #f (fresh° (q) (=° (ƛ (a) q) (ƛ (b) b)) (=° (ƛ (c) q) (ƛ (d) d)))))
(show "counter-example 2: Q = X ∧ λa.X = λb.b"
      (°->list #f (fresh° (q x) (=° q x) (=° (ƛ (a) x) (ƛ (b) b)))))
(show "counter-example 3 and its variants"
      (list (°->list #f (fresh° (q) (=° (ƛ (a) (ƛ (a) q)) (ƛ (c) (ƛ (d) c)))))
            (°->list #f (fresh° (q) (=° (ƛ (a) (ƛ (a) a)) (ƛ (c) (ƛ (d) c)))))
            (°->list #f (fresh° (q) (=° (ƛ (a) (ƛ (a) a)) (ƛ (c) (ƛ (d) d))) (=° q 'ok)))))
(define-dynamic-relation (eq° a b) (=° a b))
(show "(assume (== apple orange) (== apple X)) with a dynamic eq°"
      (°->list/datum (fresh° (q) (⊃° ((eq° 'apple 'orange)) (eq° 'apple q)))))
\end{scheme}
\begin{answers}
;; their Example 3: λa.λb.X = λc.λd.Y
((α α))

;; counter-example 1: λa.X = λb.b ∧ λc.X = λd.d
()

;; counter-example 2: Q = X ∧ λa.X = λb.b
()

;; counter-example 3 and its variants
(() () ((λ () 'ok)))

;; (assume (== apple orange) (== apple X)) with a dynamic eq°
(orange apple)
\end{answers}

For comparison, the paper's transcripts of E5, E6 (p.\,3) and E11 (p.\,7):
\begin{answers}
> (run 1 X
    (== (λ (a) X) (λ (b) b)))
'(d0)
> (run 1 X
    (== (λ (a) X) (λ (b) b))
    (== (λ (c) X) (λ (d) d)))
'(d0)
>(run 2 X
   (assume (== 'apple 'orange)
     (== 'apple X)))
'(apple orange)
\end{answers}

\begin{tutorial}[each difference, and the reason for it]
\begin{enumerate}
  \item \textbf{$\abs{a}{X}\doteq\abs{b}{b}$} (E3--E5, p.\,2--3). The paper answers \scm{(d0)}. Here
    \rulename{ξ} opens both sides with a new name $c$ and leaves $X\doteq c$, which \textbf{Check}
    rejects: $c$ was opened below a binder that $X$ is outside of. Read with quantifiers, the equation
    is $\exists X\,\forall c.\,X=c$, the first query of \cref{sec:examples-scope} and Example 8.J
    with no arguments, which is false: no single $X$ equals every $c$ (\book{85}; \cref{ex:hounif-8j}).
    The paper itself grants that this failure is what every existing algorithm answers:
    \enquote{Under any existing higher-order unification algorithm, X and a cannot be unified due to
    the scope restrictions} (p.\,2). It then calls the failure \enquote{counter-intuitive} (p.\,3).
  \item \textbf{The conjunction} with $\abs{c}{X}\doteq\abs{d}{d}$ (E6, p.\,3). The paper again answers
    \scm{(d0)}, by the same grafting. Both conjuncts fail here.
  \item \textbf{$Q=X\wedge\abs{a}{X}\doteq\abs{b}{b}$} (counter-example~2). With static closures,
    $Q$ would be bound to a level that is not available where $Q$ lives. Here it fails like
    item~1. Under the quantifier reading the variables have no positions to refer to, so this problem
    cannot arise.
  \item \textbf{$\abs{a}{\abs{a}{X}}\doteq\abs{c}{\abs{d}{c}}$} (E7, p.\,4). It fails, in the paper's
    algorithm too (their Example~2). The paper's warning is about binder-merging implementations that
    turn the \emph{ground} term $\abs{a}{\abs{a}{a}}$, which is \emph{false}, into \emph{true}. Here
    the inner \scm{a} shadows the outer one by Scheme's scoping, and the printer renames apart
    (\cref{sec:examples-alpha}). So \emph{false} = \emph{true} fails, and \emph{false} = \emph{false}
    succeeds: the second and third answers.
  \item \textbf{Example~3, $\abs{a}{\abs{b}{X}}\doteq\abs{c}{\abs{d}{Y}}$} (p.\,6). The paper needs a
    \emph{context}, a postponed equation between closures, because $X$ and $Y$ might refer to the
    binders by position. Here neither may depend on the binders, so the equation is simply
    $X\doteq Y$, solved by \rulename{Var}: \scm{((α α))}, with no postponed part.
  \item \textbf{\texttt{assume}} (E11, p.\,7). The paper's query, shown above, gives
    \scm{'(apple orange)}. The paper gives no operational account; the natural reading is ordinary unification first, then the
    hypothesis. λProlog cannot write this query, since \lp{=} is not a
    predicate one can assume clauses for. With a dynamic \scm{eq°} whose static body is \scm{=°}, the
    hypothesis comes \emph{first}, by the rule of \book{78} (\cref{sec:examples-order}):
    \scm{(orange apple)}.
\end{enumerate}
\end{tutorial}

\begin{kanren}[what the comparison shows]
The two systems agree wherever the problem is about closed terms or about patterns under \scm{∀}: α,
the \enquote{true = false} hazard of ground terms, and E12. On unification they differ exactly where a
logic variable occurs under an object binder that it is not applied to (the order of the \texttt{assume}
answers is a separate matter, of clause order). There, Ma and Friedman give the variable
access to the binders by position. The module keeps the quantifier reading: a variable outside a
binder cannot see it, and a variable that should depend on a binder must be applied to it. We ran
\begin{scheme}
(show "λa.F a = λb.b: a variable applied to the binder"
      (°->list/datum (fresh° (q) (=° (ƛ (a) (@ q a)) (ƛ (b) b)))))
\end{scheme}
\begin{answers}
;; λa.F a = λb.b: a variable applied to the binder
((ƛ (b) b))
\end{answers}
Read with quantifiers this is $\exists q\,\forall a.\,q\,a=a$, a pattern problem whose answer is the
projection. The design options are compared in
\cref{sec:lkanren-compare}, and the choice is argued in \cref{ch:design}.
\end{kanren}

% =====================================================================================================
\section{The examples at a glance}\label{sec:examples-glance}

One row per example: the book's page, the place in Part~I where the example was introduced, and how
the module's answers compare with λProlog's. \enquote{Same} means the same answers in the same order.

\begin{longtable}{@{}>{\raggedright\arraybackslash}p{0.18\fullwidthlen}>{\raggedright\arraybackslash}p{0.19\fullwidthlen}>{\raggedright\arraybackslash}p{0.22\fullwidthlen}>{\raggedright\arraybackslash}p{0.3\fullwidthlen}@{}}
\toprule
example & book & Part I & order of the answers against λProlog \\
\midrule\endhead
∃∀ / ∀∃, aliases & \book{85} & \cref{sec:hh-scope} & failure/success as in the book \\
clause order & \book{78} & \cref{sec:hh-rules} & same: \scm{(3 2 1)}, \scm{(2 3 1)} \\
courses & \book{80--81} & \cref{sec:hh-db} & same answers; all proofs listed, the book shows the first \\
jars & \book{83} & \cref{sec:hh-generic} & same \\
reverse & \book{87} & \cref{sec:hh-link} & same; one answer backwards \\
typeof & \book{190--191} & \cref{sec:lambdatree-typeof} & same; backwards finds terms where DFS loops \\
cbn & \book{183} & \cref{sec:lambdatree-interp} & same \\
copy, subst & \book{199--200} & \cref{sec:lambdatree-copy} & same \\
existential generalisation & \book{202} & \cref{sec:lambdatree-copy} & same four, same order \\
de Bruijn & \book{193--194} & \cref{sec:lambdatree-debruijn} & same terms, both directions \\
prenex & \book{196--197} & \cref{sec:lambdatree-prenex} & same five, in the order 1 3 2 4 5 \\
mappred & \book{129} & \cref{ex:hoprog-mappred} & same \\
palindromes & \book{142--143} & \cref{sec:hoprog-fdl} & same nine, same order \\
CPS & \book{257} & \cref{sec:apps-cps} & same term \\
8.H, 8.K & \book{223--225} & \cref{ex:hounif-8h,ex:hounif-8k} & the book's mgus \\
8.G, 8.I, 8.J & \book{222,224} & \cref{ex:hounif-8i,ex:hounif-8j} & failures \\
8.M, 8.A & \book{208,213} & \cref{ex:hounif-nomgu} & suspended and reported (no enumeration) \\
Ma--Friedman & --- & \cref{sec:lkanren-counter} & λProlog's answers; differences above \\
\bottomrule
\end{longtable}

% =====================================================================================================
\section{Summary}\label{sec:examples-summary}

\begin{itemize}
  \item Every core example of the book runs in \texttt{(aux kanren lambda)} with the book's answers:
    the hypothetical queries of Chapter~3, type inference, substitution, the de~Bruijn translation,
    prenex forms, palindromes and the CPS transformation, and the unification problems of Chapter~8
    (the two without a most general unifier are reported, not solved).
  \item The answer \emph{order} is the book's in every example except prenex normal forms (1~3~2~4~5)
    and the courses query, whose duplicate proofs come in the interleaving's order. Answers come in
    clause order when every clause answers at once, because a dynamic relation is one flat
    disjunction: hypotheses most recent first, then the static body.
  \item Scope is enforced by levels at every binding. \scm{∃Y∀x. Y = x} fails, \scm{∀x∃Y. Y = x}
    succeeds, and an eigenvariable cannot escape through an alias, through a relation built on
    \scm{=°}, or through the variables of a clause: lowering gives each such variable the level of the
    variable it is bound into.
  \item Clauses assumed with \scm{⊃°} share their free variables with the context, as \lp{rev nil K} and
    \lp{fact (finished dana X)} need. Their own \scm{∀} variables are renamed at each use.
  \item β is applied on demand: \scm{(@ R N)} built with \scm{R} unknown reduces as soon as \scm{R} is
    bound. Call-by-name evaluation works, and \lp{subst R N (R N)} works with \lp{R} bound before or
    after the equation.
  \item Pattern problems give the book's most general unifiers, read off the levels: projection,
    imitation, pruning (8.H), lowering with delayed raising (8.K), and the occurs and scope failures
    (8.I, 8.J).
  \item Non-patterns are suspended, closed over the eigenvariables that leave their scope, retried at
    every binding, and reported as \scm{(assert (≐ ...))}. That is enough for the palindrome program and
    for generate-and-test over 8.M, but not for enumerating the unifiers of 8.M or 8.A.
  \item Against Ma and Friedman 2021 the module gives λProlog's answers. It agrees on α, on the
    \enquote{true = false} hazard and on E12. It disagrees where a variable occurs under a binder it is
    not applied to (\scm{λa.X = λb.b} fails), and on the order of the answers of \texttt{assume}.
\end{itemize}

\chapter{Limitations, costs and a roadmap}\label{ch:roadmap}

% ------------------------------------------------------------------------------------------------
\section*{What and why}

\Cref{ch:design,ch:module,ch:examples} presented the proposal from the inside. The first gave the
extension protocol added to \micro, the second the module \texttt{(aux kanren lambda)} built on it, and
the third the book's programs running beside their λProlog originals. This chapter looks at the proposal
from the outside, as a reviewer would before a merge. It asks five questions.
\begin{enumerate}
  \item \emph{How do we know it works?} What the test suites check, how a case is written, and what a
    green run does and does not establish (\cref{sec:roadmap-testing}).
  \item \emph{Where does it differ from λProlog?} Nine documented limitations and deviations, each
    pinned by a test case, most with a small run (\cref{sec:roadmap-limits}).
  \item \emph{What does it cost?} The measured overhead on first-order search, after a \texttt{∀°},
    and on λ workloads, and where that time goes (\cref{sec:roadmap-perf}).
  \item \emph{How does it enter the repository?} A staged merge, file by file, with a review order
    (\cref{sec:roadmap-merge}), and six later work items (\cref{sec:roadmap-later}).
  \item \emph{What must the author decide?} The questions that only the maintainer of \aux{} can
    answer (\cref{sec:roadmap-questions}).
\end{enumerate}

Every Scheme snippet below was run against the prototype: micro with the protocol applied, plus the
module. Its output is pasted verbatim. The prototype will live in \texttt{doc/hol-kanren/prototype/}:
the module, \texttt{micro.diff}, \texttt{show.diff}, the design notes \texttt{DESIGN.md}, the tests (with
the two copies of micro's suites, \texttt{micro-tests/} and \texttt{micro-tests-with-lambda/}), the
examples, the benchmarks, and the scripts \texttt{build.sh} and \texttt{run-tests.sh}. After
the merge it moves to \texttt{src/aux.kanren.lambda.scm} and \texttt{src/test/microkanren-lambda*.scm}.
The snippets import \scm{(aux base)}, \scm{(aux stream)}, \micro, \scm{(aux kanren micro show)} and
\scm{(aux kanren lambda)}.

Page numbers \enquote{p.N} in the prototype's comments are PDF pages (see \cref{sec:intro-conventions}).


% ------------------------------------------------------------------------------------------------
\section{How the proposal is tested}\label{sec:roadmap-testing}

\subsection{Three layers of tests}\label{sec:roadmap-layers}

The tests come in three layers, each with its own job.
\begin{description}
  \item[Acceptance (A1--A9).] The acceptance criteria that the three prototypes shared, these cases say what the module must
    do: the scope of \texttt{∀}, hypothetical clauses, type inference, object-level substitution, $\Llam$
    unification, the differences from Ma and Friedman~\cite{MaFriedman2021}, the larger programs of the
    book, and micro's stores on λ-terms. A9 says that micro alone has not changed.
  \item[Regression (R, R2, R3).] The proposal went through three adversarial review rounds, which sent
    thirteen, eight and twelve reports. Every report was reproduced before it was fixed, and every code fix
    came with a regression case; each case of the third round fails on the code before its fix. Several reports turned out to describe the
    same defect, and two were documentation errors or a cost, not code defects. So the rounds left 12, 6
    and 9 cases.
  \item[Limits.] A separate suite, \texttt{microkanren-lambda-limits.scm}, pins what the module does
    \emph{not} do, or does differently from λProlog. Its opening doc string reads \enquote{Known limits of
    (aux kanren lambda), each pinned so that a change of behaviour is noticed.} When one of these cases
    fails, a limitation has changed, perhaps for the better, and the documentation must change with it.
\end{description}

\begin{table}[h]
\centering\small
\caption{The two suites of the module, group by group.}\label{tab:roadmap-suites}
\begin{tabular}{@{}llrr@{}}
\toprule
group & subject & cases & assertions \\
\midrule
A1 & \texttt{∀} scope: escapes, aliases, levels reaching clauses & 9 & 22 \\
A2 & \texttt{⊃} hypothetical clauses: scope, order, sharing; courses, jars, reverse & 8 & 21 \\
A3 & \texttt{∀} and \texttt{⊃} together: type inference, forwards and backwards & 3 & 8 \\
A4 & object-level substitution: cbn, β, copy, subst, existential generalisation & 4 & 9 \\
A5 & $\Llam$: flex--rigid, flex--flex, pruning, 8.G--8.M, non-patterns, η & 10 & 40 \\
A6 & Ma \& Friedman 2021, case by case & 5 & 13 \\
A7 & de Bruijn, prenex, mappred, palindromes, Fischer's CPS & 5 & 12 \\
A8 & micro's stores on λ-terms: \texttt{≠°}, \texttt{absent°}, types, printing, grouping & 11 & 65 \\
A9 & a first-order case inside the suite & 1 & 2 \\
R  & the first review round & 12 & 44 \\
R2 & the second review round & 6 & 31 \\
R3 & the third review round & 9 & 53 \\
\midrule
   & \texttt{microkanren-lambda.scm} & 83 & 320 \\
   & \texttt{microkanren-lambda-limits.scm} & 6 & 13 \\
\bottomrule
\end{tabular}
\end{table}

\subsection{How a case is written}\label{sec:roadmap-case}

A case is a list of assertions \texttt{(}$\assert$\texttt{= expected actual)} from \scm{(aux unittest)}. Every expected
answer was written by hand, from the book or from the logic, and none was captured from a run. This
matters: one prototype had a generated snapshot as its suite, which would have frozen its bugs as
expected answers. Here is the A1 case that tests escapes through an alias:
\begin{scheme}[commandchars=\\\{\}]
((test/A1/alias _)
 ; the escape through a younger variable is caught whatever the order of the bindings and their kind.
 (\ensuremath{\assert}= '() (°->list #f (fresh° (q) (∀° (x) (fresh° (y) (=° q (list y)) (=° y x))))))
 (\ensuremath{\assert}= '() (°->list #f (fresh° (q) (∀° (x) (fresh° (y) (=° y x) (=° q (list y)))))))
 (\ensuremath{\assert}= '() (°->list #f (fresh° (q) (∀° (x) (fresh° (y z) (=° q (list y)) (=° y z) (=° z x))))))
 (\ensuremath{\assert}= '() (°->list #f (fresh° (q) (∀° (x) (fresh° (y) (=° y q) (=° y x))))))
 (\ensuremath{\assert}= '() (°->list #f (fresh° (q) (∀° (x) (=° (vector q) (vector x)))))))
\end{scheme}
Three habits of the suite show here. First, the \emph{same} property is asserted under several goal
orders. Conjunction must be commutative, and many defects found by the reviews broke exactly that: an
answer lost in one order of the goals and found in another. Second, the runner states what is compared.
\scm{(°->list #f g)} gives micro's answer as rebuilding code, constraints included, while
\scm{°->list/datum} gives plain data comparable with \scm{equal?} against an answer quoted from the book.
An error is compared through its message. Third, comments cite the book's page or the review report that
the case pins.

\begin{tutorial}[What the levels do in the second assertion]
Both assertions encode $\exists q\,\forall x\,\exists y.\ q=(y)\wedge y=x$, which has no proof: it would
need $q=(x)$, and $q$ is quantified outside $x$. \Cref{sec:examples-scope} follows the first order, where
\textbf{Check} lowers $y$ when $q:=(y)$ is made and then rejects $y':=x$. Here $\ell(q)=0$ and
$\ell(x)=\ell(y)=1$ (\cref{def:design-levels}), and \scm{(=° y x)} comes first. $\ell(y)=1=\mathit{top}$,
the highest level reached (\cref{def:module-state}), and a variable of level $\mathit{top}$ sees every
eigenvariable that exists, so $y:=x$ is not checked (a fast path, \cref{sec:module-fast}). Then
\scm{(=° q (list y))} binds $q:=(y)$; the check (\cref{def:module-check}) walks $(y)$ to $(x)$ and finds
$\ell(x)=1>\ell(q)=0$. No answer. The two orders fail at different bindings, for the same reason. The
other three assertions vary the route: a chain of two variables, a var--var binding in the other
direction, a vector instead of a list.
\end{tutorial}

\subsection{micro's own suites, run twice}\label{sec:roadmap-a9}

Acceptance item A9 lives mostly outside the module's suite, which holds only one first-order case under
that name. A9 is micro's six existing suites, 125 cases in all, run twice against the edited micro: the
core suite and those of aggregation, untagged terms, show, arithmetic and dates. The first copy is
unchanged. The second differs only in its import line, which also loads the module, so that the
extension is \emph{registered} while first-order code runs. For the core suite:
\begin{answers}
5c5
<   (aux unittest) (aux base) (aux stream) (aux kanren micro) (aux fds sbral))
---
>   (aux unittest) (aux base) (aux stream) (aux kanren micro) (aux kanren lambda) (aux fds sbral))
\end{answers}
The expectations are not touched, answer order included. A seventh suite, the one of union-find, tests
\texttt{(aux kanren unionfind)}, a separate engine that does not import micro, so it was not run.

\subsection{The run}\label{sec:roadmap-run}

The script \texttt{build.sh} compiles the edited micro, its sibling modules (show, arith, date) and the
module with \texttt{csc -O2 -d1 -s -J} into a private repository, \texttt{repo/} beside the script, that
shadows the installed micro. The script \texttt{run-tests.sh} then runs the two suites of the module,
the six suites of micro twice (from \texttt{micro-tests/} and \texttt{micro-tests-with-lambda/}), and the
curated examples, with the private repository first on the search path:
\begin{Verbatim}
cd doc/hol-kanren/prototype
sh build.sh
sh run-tests.sh
\end{Verbatim}
Both scripts set \texttt{CHICKEN\_REPOSITORY\_PATH} to \texttt{repo/} followed by the installed repository,
and \texttt{build.sh} passes \texttt{-I} for the installed headers; both paths are hard-coded to
\texttt{\textasciitilde/.local} and must be adapted on another machine. Nothing is installed. The output of
\texttt{run-tests.sh}, verbatim:
\begin{answers}
test/microkanren-lambda ((ran 83) (failed 0))
test/microkanren-lambda-limits ((ran 6) (failed 0))
micro-tests/microkanren ((ran 24) (failed 0))
micro-tests/microkanren-aggregation ((ran 6) (failed 0))
micro-tests/microkanren-untagged ((ran 7) (failed 0))
micro-tests/microkanren-show ((ran 15) (failed 0))
micro-tests/microkanren-arith ((ran 52) (failed 0))
micro-tests/microkanren-date ((ran 21) (failed 0))
micro-tests-with-lambda/microkanren ((ran 24) (failed 0))
micro-tests-with-lambda/microkanren-aggregation ((ran 6) (failed 0))
micro-tests-with-lambda/microkanren-untagged ((ran 7) (failed 0))
micro-tests-with-lambda/microkanren-show ((ran 15) (failed 0))
micro-tests-with-lambda/microkanren-arith ((ran 52) (failed 0))
micro-tests-with-lambda/microkanren-date ((ran 21) (failed 0))
examples.scm ok
\end{answers}
That is $83+6=89$ cases of the module, plus micro's $24+6+7+15+52+21=125$ cases run twice, $339$ cases
in all, with 0 failures. The aux repository was never written while the prototype was built: its
\texttt{git status} shows only the untracked files it had at the start, and both diffs pass
\texttt{patch --dry-run} against the current \texttt{src/}.

\begin{pitfall}[What a green run does not say]
The suites are example-based. They show that the module gives the book's answers on the book's examples,
and that every defect found by review stays fixed. They do not prove soundness. The guarantees on pattern
problems come from the theorems the algorithm follows: Miller's termination and most-general-unifier
results~\cite{Miller1991}, the soundness and completeness of universe tagging~\cite{NadathurJayaramanKwon1995},
and its levels formulation~\cite{NadathurLinnell2005}. The tests check that the code follows those
algorithms on the cases tried. And a passing limits suite is no good news in itself: it means the
limitations are still there.
\end{pitfall}

% ------------------------------------------------------------------------------------------------
\section{Limitations and deviations from λProlog}\label{sec:roadmap-limits}

λProlog is a typed language with its own implementation of the hohh search. \texttt{(aux kanren lambda)}
is an untyped library over a miniKanren with interleaving search and constraint stores that λProlog
does not have. Some departures follow from that: they are the price of the host. Others are choices that
could be revisited. \Cref{tab:roadmap-limits} lists all nine, with the case that pins each. The
subsections take them in order.

\begin{longtable}{@{}r>{\raggedright\arraybackslash}p{0.52\textwidth}>{\raggedright\arraybackslash}p{0.36\textwidth}@{}}
\caption{The limitations, each pinned by a case.}\label{tab:roadmap-limits}\\
\toprule
& limitation & pinned by \\
\midrule\endfirsthead
\toprule
& limitation & pinned by \\
\midrule\endhead
1 & only dynamic relations take hypotheses; no predicate quantification & \texttt{A2/static-relations} \\
2 & untyped: η only at abstractions, \texttt{@} distinct from lists, β may diverge & \texttt{limits/untyped} \\
3 & non-patterns suspended, never enumerated: floundering reported, not detected & \texttt{limits/no-huet} \\
4 & answer order is micro's interleaving & \texttt{A7/prenex}, \texttt{A3/backwards} \\
5 & \texttt{≠°}, \texttt{absent°} and types are extensions, with a chosen reading of \texttt{∀} & \texttt{A8/diseq-extensional}, \texttt{R2/extensional-absent-order} \\
6 & host escapes: levels police logic variables, not Scheme values & \texttt{limits/host-escape} \\
7 & binder names are hints; some projected constraints keep a working variable & \texttt{A6/alpha}, \texttt{limits/absent-projection} \\
8 & costs: no indexing of E and N; the scope check after any \texttt{∀°} & \texttt{bench/} (\cref{sec:roadmap-perf}) \\
9 & the answer variable must be created outside every \texttt{∀°} & \texttt{R2/answer-inside-forall} \\
\bottomrule
\end{longtable}

\subsection{Only dynamic relations take hypotheses}\label{sec:roadmap-lim-dynamic}

In λProlog every predicate is dynamic: \lp{D => G} may add clauses for any predicate, including one
that has static clauses in the program. Here only a relation defined with \scm{define-dynamic-relation}
consults the context Γ, and assuming a clause for any other relation is an error. A \texttt{⊃°} that
silently ignored such a clause would be worse, and all three prototypes did that.
\begin{scheme}
(define-relation (r° x) (=° x 1))
(define-dynamic-relation (d° x) (=° x 1))
(error-message (λ () (°->list #f (fresh° (q) (⊃° ((r° 2)) (r° q))))))
(°->list/datum (fresh° (q) (⊃° ((d° 2)) (d° q))))
\end{scheme}
\begin{answers}
"not a dynamic relation, see define-dynamic-relation"
(2 1)
\end{answers}
The dynamic relation answers 2 before 1, because hypotheses come first, most recent first, then the
static body (\book{78}, \cref{sec:hh-rules}).

Two consequences follow, both because clauses are Scheme closures and not formulas. A program cannot
build a clause as a term and then assume it. And there is no quantification over predicates, so the
book's \lp{reverse}, which hides its auxiliary \lp{rev} under \lp{pi rev\ ...} \book{125--126}
(\cref{def:hoprog-hohhplus}), cannot be transcribed as it stands. The suite runs instead the \fohh{} \lp{reverse} of
Figure~3.4 \book{87}, which assumes the clauses of its auxiliary predicate but does not hide its name: here
\scm{rev°} is a dynamic relation declared at top level, with no static clauses (\texttt{A2/reverse}). Roadmap item (e) addresses
predicate quantification (\cref{sec:roadmap-later-pred}).

\subsection{Untyped terms}\label{sec:roadmap-lim-untyped}

λProlog's terms are simply typed, and its unifier works on βη-long normal forms. The module has no
types. It keeps terms η-short and η-expands one side only when the other is an abstraction, which is
Miller's treatment of the untyped case~\cite{Miller1991}. Three facts follow.
\begin{scheme}
(°->list #f (fresh° (q) (=° (@ 'g 'a) '(g a))))
(μkanren-datum (ƛ (x) (@ 'g x)))
\end{scheme}
\begin{answers}
()
g
\end{answers}
\begin{itemize}
  \item λ-tree application \scm{(@ 'g 'a)} is not the list \scm{'(g a)}: the first is a term of the
    meta-language, the second Scheme data of the object language. Keeping them apart is what keeps
    conjunction commutative. One prototype conflated them and lost that property.
  \item $\abs{x}{(g\,x)}$ is built η-short as $g$.
  \item Full β through \scm{@} on instantiated heads is not guaranteed to terminate on untyped
    meta-terms. With types, β-reduction would terminate.
\end{itemize}
There is also no check that a solution exists in a possibly empty type (\cref{sec:lambda-solutions}).
Roadmap item (f) addresses types (\cref{sec:roadmap-later-types}).

\subsection{Non-patterns are suspended, never enumerated}\label{sec:roadmap-lim-huet}

Outside $\Llam$ the module does what Teyjus does: it suspends the equation in the store E, retries it after every
binding in case it has become a pattern, and reports it in the answer if it is still suspended (\cref{sec:hounif-pragmatics}).
It does not run Huet's procedure (\cref{sec:hounif-huet}). Two runs show what this means:
\begin{scheme}
(°->list #f (fresh° (q F) (=° q F) (=° (@ F 'a) 'b) (=° (@ F 'a) 'c)))
(°->list #f (fresh° (q F) (=° q F) (=° (@ F 'a) '(g a a))))
\end{scheme}
\begin{answers}
((λ (α) (assert (≐ (@ α 'a) 'b)) (assert (≐ (@ α 'a) 'c)) α))
((λ (α) (assert (≐ (@ α 'a) (cons 'g (cons 'a (cons 'a '()))))) α))
\end{answers}
The second problem is the book's Example 8.A (\cref{ex:hounif-nomgu}), $F\,a=g\,a\,a$, written with list
data for the right-hand side. It has four unifiers and no most general one \book{213}. The module does
not choose: it answers that $F$ is any α satisfying the residual. That answer is correct but conditional.

The first problem shows the cost of not choosing. $F\,a=b$ and $F\,a=c$ have no common solution, since
$b\ne c$. Huet's procedure would find that out, but the module only suspends both equations and returns
an answer whose residuals are jointly unsatisfiable. This is \emph{floundering reported but not detected}:
the answer says what was left undecided, and it is up to the reader, or a later binding, to find out that
it is empty. Roadmap item (d) would add an optional Huet search over E (\cref{sec:roadmap-later-huet}).

\subsection{Answer order is micro's interleaving}\label{sec:roadmap-lim-order}

λProlog searches depth-first, left to right; micro interleaves the branches of a disjunction. Where the
book lists answers, the two orders often agree: augmentation order \book{78}, the existential
generalisations \book{203}, the palindromes \book{142}, reverse, the course database (whose extra,
duplicate proofs come in the interleaving's order), and the CPS term
\book{257}. They differ on \lp{prenex}, whose five answers \book{196--197} come out as the book's first,
third, second, fourth and fifth:
\begin{answers}
((all (ƛ (x) (all (ƛ (x1) (and (q x x) (q x x1))))))
 (all (ƛ (x) (all (ƛ (x1) (and (q x1 x1) (q x x1))))))
 (all (ƛ (x) (all (ƛ (x1) (all (ƛ (x2) (and (q x x) (q x1 x2))))))))
 (all (ƛ (x) (all (ƛ (x1) (all (ƛ (x2) (and (q x1 x1) (q x x2))))))))
 (all (ƛ (x) (all (ƛ (x1) (all (ƛ (x2) (and (q x2 x2) (q x x1)))))))))
\end{answers}
The difference also helps. Where depth-first search loops, interleaving finds answers that λProlog misses.
Asked for terms of type $i\to i$, \lp{typeof} run backwards loops under depth-first search. Here it
enumerates terms:
\begin{answers}
((abs (ƛ (x) x))
 (abs (ƛ (x) (app (abs (ƛ (x1) x1)) x)))
 (abs (ƛ (x) (app (abs (ƛ (x1) x)) x))))
\end{answers}
The suite compares answer \emph{sets}, through a sorted copy, wherever interleaving rather than logic
decides the order. Where it compares a sequence against the book, the book's order is the claim being tested.

\subsection{Disequality, absence and types under \texttt{∀}}\label{sec:roadmap-lim-readings}

λProlog has no disequality, no \emph{absento} and no type constraints. micro has all three, so the
module had to decide what they mean when an eigenvariable is involved. For \scm{≠°} there are two
defensible readings of $\forall x.\ q\ne x$.
\begin{definition}[two readings of an eigenvariable under $\ne$]\label{def:roadmap-readings}
Under the \emph{generic} reading, an eigenvariable is the new constant $c\notin\Sig$ introduced by
\rulename{generic} (\cref{fig:rules-right}), distinct from every term the older variables can denote. Under the \emph{extensional} reading, an eigenvariable stands for an
arbitrary term of its universe~\cite{Moiseenko2019,JinEtAl2021}.
\end{definition}
The generic reading is the default: it matches λProlog, where \lp{pi x\ not (q = x)} succeeds, and it
behaves like αKanren's noms~\cite{ByrdFriedman2007}. The extensional reading is chosen by the parameter
\scm{μkanren-λ-reading}:
\begin{scheme}
(°->list #f (fresh° (q) (∀° (x) (≠° q x))))
(parameterize ((μkanren-λ-reading 'extensional)) (°->list #f (fresh° (q) (∀° (x) (≠° q x)))))
(parameterize ((μkanren-λ-reading 'extensional))
  (list (°->list #f (fresh° (q) (∀° (x) (fresh° (y) (≠° y x) (=° y 'a))) (=° q 'ok)))
        (°->list #f (fresh° (q) (∀° (x) (fresh° (y) (≠° y x) (=° y `(f ,x)))) (=° q 'ok)))))
(°->list #f (fresh° (q) (∀° (x) (fresh° (y) (≠° y x) (=° y 'a))) (=° q 'ok)))
\end{scheme}
\begin{answers}
((λ (α) α))
()
(() ((λ () 'ok)))
((λ () 'ok))
\end{answers}
Read them as statements about all $x$.
\begin{itemize}
  \item Generic, $\forall x.\,q\ne x$: $q$ is older than $x$ and can never denote it, so the
    disequality holds for good and vanishes from the answer.
  \item Extensional, $\forall x.\,q\ne x$: $x$ ranges over every term, $q$'s value included, so the
    disequality cannot hold for all $x$ and the goal fails.
  \item Extensional, $\forall x\,\exists y.\ y\ne x\wedge y=a$: the instance $x=a$ refutes it.
    With $y=(f\;x)$ it holds, since no term $x$ equals the list $(f\;x)$ (occurs check).
  \item Generic, $\forall x\,\exists y.\ y\ne x\wedge y=a$: $a$ is not the new name $x$, so it holds.
\end{itemize}
Both readings were implemented so that the order of a binding and a disequality does not change the
answer. \texttt{A8/diseq-extensional} pins both orders of $\forall x.\,F\,x\ne x$ and $F:=\abs{z}{c}$, and
\texttt{R2/extensional-absent-order} pins four goal orders of $\forall x\,\exists Y.\ Y\ne a\wedge
\mathit{absent}(a,Y)\wedge Y=x$, which fails in every order.

The reading has edges, and the proposal documents them. \scm{absent°} and the types always use the generic
reading. A disequality keeps the reading in force when it is posted, with one exception. A first-order
disequality sits in micro's store D, which records no reading. If it later \emph{migrates} to the
module's store N, because one of its sides came to hold a λ-term or an eigenvariable, it takes the reading
in force at the migration. The two readings then differ only if a program changes
\scm{μkanren-λ-reading} between the two goals. Under the extensional reading, once a \texttt{∀°} has run,
a first-order disequality goes to N at once, with its reading.

\subsection{Host escapes}\label{sec:roadmap-lim-host}

Levels police logic variables. A Scheme value that host code carries out of a scope is not a logic
variable, and the levels know nothing about the route it took. The limits suite pins both halves of the
story:
\begin{scheme}
(define stash #f)
(°->list #f (fresh° (q) (∀° (x) (project° ((v x)) (begin (set! stash v) ✓°)))
                        (λ (s) ((=° q stash) s))))
(°->list #f (fresh° (q) (∀° (x) (project° ((v x)) (begin (set! stash v) ✓°)))
                        (∀° (y) (fresh° (z) (=° z stash) (=° q 'leaked)))))
\end{scheme}
\begin{answers}
()
((λ () 'leaked))
\end{answers}
\begin{tutorial}[Why the first escape is caught and the second is not]
In both goals the first \scm{∀°} raises the depth to 1 and creates $x$ with $\ell(x)=1$. \scm{project°}
stores the eigenvariable record in a Scheme variable, and the scope exits, so the depth is back to 0.
\begin{enumerate}
  \item First goal: $q$ was created at depth 0. Binding $q:=x$ is checked, and $\ell(x)=1>\ell(q)=0$ is
    a scope failure. The smuggled name cannot reach an \emph{older} variable.
  \item Second goal: a second \scm{∀°} raises the depth to 1 again, and $z$ is created inside it with
    $\ell(z)=1$. Binding $z:=x$ passes the check, since $\ell(x)=1\le\ell(z)$. A level is a depth, not the
    identity of a scope, so two sibling scopes at the same depth look alike.
\end{enumerate}
\end{tutorial}
Inside the logic no such program can be written: a variable created in a later scope never sees the
earlier eigenvariable, because it is out of its lexical scope. The escape needs \scm{project°} and
\scm{set!}. The point of the first case is that even this route cannot violate the quantifier prefix of
the older variables. That guarantee has a price, discussed in \cref{sec:roadmap-perf-forall}.

\subsection{Binder names are hints; projected constraints}\label{sec:roadmap-lim-hints}

A λ-record stores an interned hint, the name the user wrote, for printing. Terms built alike are
therefore \scm{equal?}, and they group in \scm{groupby°}. α-variants written with different names are
not \scm{equal?}, although they are α-equivalent, which \scm{≐} and \scm{μkanren-α=?} decide:
\begin{scheme}
(list (equal? (ƛ (x) `(f ,x)) (ƛ (x) `(f ,x)))
      (equal? (ƛ (x) `(f ,x)) (ƛ (y) `(f ,y)))
      (≐ (ƛ (x) `(f ,x)) (ƛ (y) `(f ,y))))
\end{scheme}
\begin{answers}
(#t #f #t)
\end{answers}
A hint that \scm{equal?} ignored could not be stored in the record. The proposal chose to print answers
with the user's binder names, and gave up \scm{equal?} modulo α for that (\cref{sec:roadmap-questions}).

The second half of this item is about answers. An \emph{absento} on a flexible term whose head is in
the answer is kept whole, and an argument out of the answer prints as a working variable:
\begin{scheme}
(°->list #f (fresh° (q) (fresh° (Y) (absent° 'c (@ q Y)))))
(°->list #f (fresh° (q) (fresh° (F) (absent° 'c (@ F q)))))
\end{scheme}
\begin{answers}
((λ (α) (assert (absent? 'c (@ α _1))) α))
((λ (α) α))
\end{answers}
The first answer says: there is some \texttt{\_1} such that $c$ does not occur in $\alpha\,\_1$. It cannot
be weakened to an absento on $\alpha$ alone: $\alpha:=\abs{z}{z\,c}$ contains $c$, yet with
$\_1:=\abs{v}{b}$ the application reduces to $b$. In the second answer the head $F$ is out of the answer
and may be a constant function, so the constraint is dropped. A disequality with such an argument,
\scm{(deny (equal? β (@ α _3)))}, and an absento whose datum may still reduce to an atom of the answer,
\scm{(absent? (@ α 'b) 'a)}, are likewise kept in the answer (\texttt{R3/diseq-unreachable-argument},
\texttt{R3/absento-datum-reduces}).

\subsection{Costs}\label{sec:roadmap-lim-costs}

The engine pays for three things that λProlog implementations avoid. Every binding retries all the
suspended equations of E and the disequalities of N when they are not empty, with no index by variable
(the wake-up loop of \cref{sec:module-wake}).
The scope check walks bound terms below the top level from the first \texttt{∀°} on, for the rest of the
search, even after that \texttt{∀°} has exited. And deeply nested binders scale quadratically.
\Cref{sec:roadmap-perf} gives the figures.

\subsection{The answer variable outside every \texttt{∀°}}\label{sec:roadmap-lim-answer}

\scm{°->list} reifies variable 0, the first one created. If it is created inside a \scm{∀°} and its
value holds that scope's eigenvariable, no answer can be written without the eigenvariable. λProlog
reports no binding for a variable quantified inside \lp{pi x\ }. The module raises an error instead of
printing an eigenvariable free:
\begin{scheme}
(error-message (λ () (°->list #f (∀° (x) (fresh° (q) (=° q x))))))
(°->list #f (∀° (x) (fresh° (q) (=° q 'a))))
(°->list #f (fresh° (q) (∀° (x) (=° q 'a))))
\end{scheme}
\begin{answers}
"the answer holds the eigenvariable of a ∀° it was created in: create the variable it reifies outside every ∀°"
((λ () 'a))
((λ () 'a))
\end{answers}
An answer that does not hold the eigenvariable is reported as usual, wherever its variable was created.
The idiomatic form is the third one: create the query variable first.

% ------------------------------------------------------------------------------------------------
\section{Performance}\label{sec:roadmap-perf}

All figures are in milliseconds, compiled with \texttt{csc -O2 -d1}, the best of three runs per round, with
two rounds. Three binaries are compared. \emph{before} is the installed micro, unedited. \emph{after} is
the edited micro with no extension loaded. \emph{ext} is the edited micro with \texttt{(aux kanren
lambda)} loaded, so its hooks are registered. The first-order benchmark is the one the engineering review
used to compare the three prototypes.

\subsection{The protocol on first-order search}\label{sec:roadmap-perf-protocol}

\begin{table}[h]
\centering\small
\caption{The cost of the protocol on first-order search.}\label{tab:roadmap-perf-protocol}
\begin{tabular}{@{}lrrr@{}}
\toprule
benchmark (ms, round 1 / round 2) & before & after & ext \\
\midrule
\texttt{append°} backwards, 2000 answers, full & 6068 / 5923 & 6121 / 5959 & 6151 / 5929 \\
the same, search only & 672 / 662 & 690 / 669 & 691 / 668 \\
open \texttt{append°}, 1000 answers, search & 1311 / 1314 & 1356 / 1326 & 1327 / 1328 \\
perm6 with 15 pairwise \texttt{≠°}, 720 answers & 176 / 178 & 184 / 178 & 192 / 199 \\
\bottomrule
\end{tabular}
\end{table}

Read \cref{tab:roadmap-perf-protocol} column by column. From \emph{before} to \emph{after}, the protocol
alone costs at most 1\% end to end on \texttt{append°} (5923 to 5959 ms in round 2), within the 2.5\%
round-to-round noise of the unedited binary. That is what one \scm{pair?} test per call site costs when
no extension is registered. All call sites are off micro's hottest path: the variable clauses and the
pair--pair clause of unification are untouched.

From \emph{after} to \emph{ext}, \texttt{append°} does not move, but perm6 does: 192 / 199 ms, 8--12\% more
than the unedited micro's 176 / 178 ms (4--12\% more than \emph{after}). perm6 keeps
15 disequalities in D throughout, so after every \scm{=°} the \texttt{post} hook must decide whether a
new binding brought a λ-term or an eigenvariable into one of D's prefixes (\cref{def:micro-prefix}),
which would then have to migrate to N (\cref{sec:roadmap-lim-readings}). That look costs a few microseconds per binding while D is non-empty.
The d2 prototype, on which the proposal is based, paid 45--50\% here.

\subsection{With D or A non-empty}\label{sec:roadmap-perf-DA}

The third review round found that this look was too expensive in another case: one unrelated constraint
before a long first-order search. Recall micro's stores (\cref{sec:micro-stores}): D holds the
disequalities, A the type constraints such as \scm{symbol°}, and T the absentos. With D or A non-empty,
the proposal walked every new binding through the substitution, and \texttt{append°} backwards binds
long ground tails at every step. That walk is the factor of 1.65 to 1.76 in the middle column below. The last row is a control: its constraint goes to T,
which this walk never concerned.
\begin{table}[h]
\centering\small
\caption{One unrelated constraint before the search (\texttt{t} is not used by \texttt{append°}).}\label{tab:roadmap-perf-DA}
\begin{tabular}{@{}lrrr@{}}
\toprule
\texttt{append°} backwards, 2000 answers, search (ms) & after & ext, before the fix & ext, now \\
\midrule
plain & 673 / 669 & 665 & 657 / 664 \\
+ \texttt{symbol° t} & 674 / 675 & 1097 & 687 / 695 \\
+ \texttt{≠° t 1} & 655 / 657 & 1172 & 670 / 675 \\
+ \texttt{absent° t} & 661 / 664 & --- & 656 / 666 \\
\bottomrule
\end{tabular}
\end{table}
The fix bounds the look: the hook inspects the new bindings for at most 64 nodes, and past that it walks
D and A themselves. The order matters. Walking D and A at every \scm{=°}, the first fix tried, cost perm6
60\% (288 / 302 ms), because perm6's 15 prefixes are walked each time. With the bounded look first,
perm6 stays at 193 / 194 ms, against 181 / 179 for \emph{after}.

\subsection{After a \texttt{∀°}}\label{sec:roadmap-perf-forall}

The second review round measured first-order search that runs after some \texttt{∀°} has run and exited.
In \cref{tab:roadmap-perf-forall}, \texttt{pre} is \scm{(∀° (z) ✓°)}, which binds nothing and exits at
once, and \texttt{typeof°} is one type inference of $\abs{x}{x}$ run before the search.
\begin{table}[h]
\centering\small
\caption{First-order search after a \texttt{∀°} has run. A second run gave 680 / 1125 / 1154 and 1354 / 3055 ms; a third, after the third review's fixes, 663 / 1101 / 1090 and 1310 / 2918 ms.}\label{tab:roadmap-perf-forall}
\begin{tabular}{@{}lrrr@{}}
\toprule
benchmark (ms) & no \texttt{∀°} before & \texttt{pre} before & \texttt{typeof°} before \\
\midrule
\texttt{append°} backwards, 2000 answers, search & 663 & 1104 (+66\%) & 1092 (+65\%) \\
open \texttt{append°}, 1000 answers, search & 1296 & 2940 (+127\%) & --- \\
\bottomrule
\end{tabular}
\end{table}

\begin{tutorial}[Where the time goes after a \texttt{∀°}]
Take one step of \texttt{append°} backwards on the list $L$ of 1999 numbers. It binds a variable to a
pair whose tail is a ground suffix of $L$.
\begin{enumerate}
  \item \emph{No \texttt{∀°} has run.} The level table is empty, no scope check runs, and the only walk of
    the bound term is micro's occurs check.
  \item \emph{\texttt{pre} has run.} The depth went to 1 and back to 0, so the highest level reached is
    $\mathit{top}=1$. The search's variables are at level 0, below $\mathit{top}$. Each binding of such a variable gets the
    scope check, whose allocation-free \scm{μkanren-plain?} walks the bound term through the
    substitution, beside micro's occurs check. On \texttt{append°} the two walk the same ground tail of the
    list, so each binding walks that tail twice.
\end{enumerate}
Why not skip the check when the depth is back to 0 and nothing of a higher level is reachable? For logic
variables that would be safe: an exited scope reaches a depth-0 binding only through constraints, which
are checked when they wake. But an eigenvariable smuggled out by host code would then reach an older
variable unchecked, which is the first case of \cref{sec:roadmap-lim-host}. The cheap fix gives up that
guarantee, so it was not made. The proper fix is one walk that does both checks (\cref{sec:roadmap-later-walk}).
\end{tutorial}

\subsection{λ workloads}\label{sec:roadmap-perf-lambda}

\Cref{tab:roadmap-perf-lambda} compares the final module with the three prototypes. The prototypes'
columns are the engineering review's measurements on the same programs.
\begin{table}[h]
\centering\small
\caption{λ workloads. d1 is the layered prototype, d2 the hooks prototype the proposal grew from, d3
the native one.}\label{tab:roadmap-perf-lambda}
\begin{tabular}{@{}lrrrr@{}}
\toprule
benchmark (ms) & final & d2 & d3 & d1 \\
\midrule
\texttt{typeof} of the Church numeral 400 & 254 & 441 & 105 & 385 \\
\texttt{copy} of the Church numeral 400 & 256 & 442 & 116 & 314 \\
\texttt{typeof} of 150 nested binders & 1306 & --- & 1409 & --- \\
\texttt{copy} of 150 nested binders & 1745 & --- & 1779 & --- \\
\texttt{typeof} backwards, 30 (term, type) pairs & 10 & 11 & --- & --- \\
one pattern equation $F\,x_1\dots x_{50}=g\,x_{50}\dots x_1$ under 50 \texttt{∀} & 2 & 6 & --- & --- \\
\bottomrule
\end{tabular}
\end{table}

The final module is 1.7 times faster than d2, from which it grew, and 2.4 times slower than d3 on the
Church numeral. Its history on \texttt{typeof} of Church 400 shows where the time went:
\begin{center}\small
\begin{tabular}{@{}lr@{}}
\toprule
step & ms \\
\midrule
d2 ported, with instantiation that shifts only open terms & 463 \\
allocation-free \scm{μkanren-plain?} for the scope check & 330 \\
no check for bindings of variables at the top level & 311 \\
the \texttt{atom?} hook: micro's occurs check stops rebuilding eigenvariables & 254 \\
\bottomrule
\end{tabular}
\end{center}
The last step is worth explaining. micro's occurs check rebuilds a foreign record as the list of its
slots to look inside it. An eigenvariable holds no logic variable, so the hook lets micro skip it, and
that alone saved 18\%.

A profile of the remaining time puts it in micro's walks, \scm{μkanren-state-find} and its sbral
lookup. These are the var--var chains that the type variables of \texttt{typeof} form, since micro binds
the older variable to the newer, plus the occurs check of every binding of a subterm. Both are on micro's
hot path, outside the extension's reach. d3's lead comes from its in-core binding, which does the occurs
check and the level check in one walk. That is roadmap item (a). The remaining figures show no gap: on
150 nested binders the final module is 2--8\% faster than d3. The module's cost on nested binders grows
quadratically with the nesting (\cref{sec:roadmap-lim-costs}).

The fixes of the review rounds cost nothing measurable here. After the second round: typeof Church 400
260 ms, copy 259 ms, nested typeof 1292 ms, nested copy 1732 ms. After the third: 253, 256, 1284 and
1729 ms.

% ------------------------------------------------------------------------------------------------
\section{A staged merge into the aux repository}\label{sec:roadmap-merge}

The proposal is three independent pieces of different risk. They should enter the repository in three
stages, each reviewed and tested on its own.

\begin{table}[h]
\centering\small
\caption{Which file goes where.}\label{tab:roadmap-files}
\begin{tabular}{@{}lll@{}}
\toprule
stage & from \texttt{doc/hol-kanren/prototype/} & to \\
\midrule
1 & \texttt{micro.diff} (+235/−41) & \texttt{src/aux.kanren.micro.scm} \\
2 & \texttt{aux.kanren.lambda.scm} (1255 lines) & \texttt{src/aux.kanren.lambda.scm} \\
  & \texttt{test/microkanren-lambda.scm}, \texttt{-limits.scm} & \texttt{src/test/microkanren-lambda*.scm} \\
  & \texttt{test/programs.scm} & \texttt{src/test/programs.scm} \\
  & one \texttt{extension} line & \texttt{src/aux.egg} \\
3 & \texttt{show.diff} (+42/−12) & \texttt{src/aux.kanren.micro.show.scm} \\
\bottomrule
\end{tabular}
\end{table}

\subsection{Stage 1: the protocol}\label{sec:roadmap-stage1}

The diff is self-contained and mentions neither λ-terms, levels nor clauses. It could carry another
store as well, nominal terms or finite domains, so it is worth having even if stage 2 is postponed.
Both diffs were checked with \texttt{patch --dry-run} against the current \texttt{src/}. Applying the first
and, once the egg is rebuilt and installed, running micro's suites (without installing anything, the
private repository of \texttt{build.sh} and \texttt{run-tests.sh}, \cref{sec:roadmap-run}, gives the
same run):
\begin{Verbatim}
cd src
patch aux.kanren.micro.scm < ../doc/hol-kanren/prototype/micro.diff
make test-microkanren test-microkanren-show test-microkanren-arith test-microkanren-date
\end{Verbatim}
These four targets run 112 cases. The aggregation and untagged suites have no target of their own: the
\texttt{test} target runs them with the other four (\texttt{src/Makefile}, lines 46--51), or they run alone
with \texttt{cd test \&\& csi -s microkanren-aggregation.scm}, and likewise for
\texttt{microkanren-untagged.scm}. The expected result is \cref{sec:roadmap-run}'s: 125 cases, unchanged answers. The benchmark's
\emph{after} column (\cref{tab:roadmap-perf-protocol}) is the performance criterion. The diff already
documents each hook in micro's comments; the merge should add to them a short section on the protocol as
a whole.

What to review, in this order:
\begin{enumerate}
  \item \emph{The state.} A seventh field \texttt{X}, the extensions' slot, which micro never reads. It
    is copied at micro's two positional constructions (\scm{μkanren-state-update} and \scm{freshª}),
    \scm{μkanren-state-with} accepts \scm{X:}, and \scm{make-μkanren-state} keeps its six-argument
    signature, so the positional states of \texttt{test/microkanren.scm} still work.
  \item \emph{The registry and the call sites.} \scm{make-μkanren-extension} and
    \scm{μkanren-extension-register!}; each call site costs one \scm{pair?} test when the registry is
    empty. There are thirteen hooks: \texttt{term?}, \texttt{atom?}, \texttt{unify}, \texttt{occurs},
    \texttt{post}, \texttt{diseq}, \texttt{tag}, \texttt{parts}, \texttt{value}, \texttt{repr},
    \texttt{terms}, \texttt{forms} and \texttt{environment}; \cref{ch:design} gives their contracts.
  \item \emph{The two invariants} that every hook must keep (\cref{def:design-invariants}): a hook that
    changes nothing answers the very state it was given, and a hook that binds returns the extended
    \texttt{assocs}.
  \item \emph{The two exported primitives}: \scm{μkanren-state-forget}, which unbinds a variable so that a
    \texttt{post} hook can turn a binding just made back into a constraint before D, A and T see it, and
    \scm{μkanren-post-unify}, the post-\scm{=°} chain that an extension goal unifying on its own must call
    (\cref{sec:design-forget}).
  \item \emph{One change of micro's own behaviour.} micro decided an atom against an absento's datum with
    \scm{equal?} only, so \scm{(absent° D q) (=° q 'a) (=° D 'a)} succeeded in micro alone. The diff
    decides it by unification and watches it while pending. This is a bug fix of micro, independent of λ,
    and deserves its own look.
\end{enumerate}

\subsection{Stage 2: the module and its suites}\label{sec:roadmap-stage2}

Besides \texttt{srfi-1} and CHICKEN's own modules, the module imports \scm{(aux base)}, \scm{(aux stream)},
\scm{(aux fds sbral)} and \micro, and registers
its extension when it is loaded. The egg gets one line after micro's, in the style of its neighbours:
\begin{Verbatim}
  (extension aux.kanren.micro)
  (extension aux.kanren.lambda)
  (extension aux.kanren.arith)
\end{Verbatim}
Other entries of \texttt{src/aux.egg} state their dependencies with \texttt{component-dependencies}. Writing
\texttt{(component-dependencies aux.kanren.micro)} would make this one explicit too. The suite
includes \texttt{programs.scm} from its own directory, so the file goes beside it in
\texttt{src/test/}. The Makefile would follow its existing pattern: two lines in \texttt{test}, after
micro's suites, and a target of its own.
\begin{Verbatim}
	cd test && ${CSI} -s microkanren-lambda.scm
	cd test && ${CSI} -s microkanren-lambda-limits.scm
\end{Verbatim}
Then micro's six suites should be run once more with the module loaded, as in
\cref{sec:roadmap-a9}. The simplest way is to keep the second copy of the six suites under
\texttt{doc/hol-kanren/prototype/}, since they differ from \texttt{src/test/} only in their import line.

What to review, in this order. First the limits suite, 51 lines that state what is not promised. Then
the scope check and variable elimination, which carry soundness (\cref{ch:module}). Then the discharge
at \texttt{∀°} exit, then the store N. The printing code comes last: a mistake there shows up in answers,
not in the search.

\subsection{Stage 3: the show module}\label{sec:roadmap-stage3}

\texttt{show.diff} teaches \scm{(aux kanren micro show)} to read \scm{ƛ}, \scm{@}, \scm{≐}, \texttt{∃} and
\texttt{∀} in answers, draws λ-terms in the mermaid cells view with a node per binder, and adds
\scm{°->pretty}. It reads answer forms syntactically and does not import the module, so it is independent
of stage 2. It is useful only with it.
\begin{Verbatim}
cd src
patch aux.kanren.micro.show.scm < ../doc/hol-kanren/prototype/show.diff
make test-microkanren-show
\end{Verbatim}

% ------------------------------------------------------------------------------------------------
\section{Later work}\label{sec:roadmap-later}

Six items were left out of the proposal on purpose. Each should be measured separately, against the
figures of \cref{sec:roadmap-perf} and the suites of \cref{sec:roadmap-testing}. For each, the paragraphs
below say how it would fit the protocol and which concept of the book it serves.

\subsection{(a) One walk for the occurs and the scope check}\label{sec:roadmap-later-walk}

\emph{What.} micro's binding site, \scm{μkanren-state-update}, runs the occurs check. The module's
\texttt{post} hook then walks the same term again for the scope check. A combined walk does both, as d3's
in-core \texttt{occurs/fast} does.

\emph{How it fits.} A new hook at the binding site, which returns a \emph{verdict} instead of micro's
boolean: no occurrence and nothing out of scope; an occurs failure; or a term that needs variable
elimination. With no extension registered, micro keeps its own check.

\emph{What it buys.} Most of the cost of first-order search after a \texttt{∀°} (+65\% to +127\%,
\cref{tab:roadmap-perf-forall}) and the remaining factor of 2.4 against d3 on Church 400
(\cref{tab:roadmap-perf-lambda}). It keeps the host-escape guarantee, because the check still runs at depth 0.

\emph{Book concept.} The mixed prefix of a unification problem (\cref{sec:lambda-unifprob}), read as
universe levels (\cref{sec:rules-levels}, \cref{sec:hounif-pragmatics}). Nadathur, Jayaraman and Kwon's binding rule~\cite{NadathurJayaramanKwon1995}
is one condition on the bound term, so it is natural to check it in the same walk as the occurs check.

\subsection{(b) Var--var orientation by level}\label{sec:roadmap-later-varvar}

\emph{What.} When two variables are unified, bind the one of higher level to the one of lower level.
micro binds the older to the newer, so a variable of a deep scope may become the value of an outer one,
and the scope check must then lower it (\cref{def:rules-binding}).

\emph{How it fits.} It needs a hook in micro's var--var clause, the hottest path after pair--pair. That
is why the graft from d3 was declined. \textbf{Check} (\cref{def:module-check}) already lowers the
younger variable, which gives the same most general unifier up to renaming, so this item is about cost only: fewer lowering bindings, and
perhaps shorter var--var chains, which the profile of \cref{sec:roadmap-perf-lambda} places among the
remaining costs. It must be measured on first-order search too, where the hook would cost on every
var--var binding.

\emph{Book concept.} Raising and lowering in a mixed prefix (\cref{sec:lambda-raising}). Binding the
deeper variable to the outer one is always within scope. Binding the outer variable to the deeper one is
allowed only after the deeper one has been lowered, which is the extra work this item saves.

\subsection{(c) Indexing E and N by variable}\label{sec:roadmap-later-index}

\emph{What.} Today every \scm{=°} retries all suspended equations of E and all disequalities of N, oldest
first, in rounds, until no round binds (\cref{sec:roadmap-lim-costs}). An index from each variable to the
constraints that mention it would retry only those the new bindings touch.

\emph{How it fits.} Entirely inside the extension. E and N live in the slot \texttt{X}, which micro never
reads, and the \texttt{post} hook already receives the new bindings as \texttt{assocs}. No change to micro.
The care needed is in what \enquote{touches} means: a binding wakes a constraint on a variable it reaches
through the substitution, not only on the variable it binds, and the variables created by lowering and
pruning must be indexed too.

\emph{Book concept.} Dynamic $\Llam$ (\cref{sec:hounif-pragmatics}): an equation that does not respect the
$\Llam$ condition is deferred \enquote{until such time that substitutions resulting from solving other parts of
the problem have altered its status} \book{226}. An index decides cheaply which deferred equations a
substitution may have altered.

\subsection{(d) Huet pre-unification as an optional search}\label{sec:roadmap-later-huet}

\emph{What.} An opt-in goal that takes the suspended equations of E and searches their solutions with
Huet's imitation and projection steps~\cite{Huet1975}, each choice a branch of the interleaving search.
On Example 8.A it would enumerate the four unifiers of \cref{ex:hounif-nomgu}. On the first problem of
\cref{sec:roadmap-lim-huet} it would find no solution and fail, so floundering would be detected, not
only reported.

\emph{How it fits.} E is the natural agenda: it already holds the non-pattern equations, closed over the
eigenvariables they mention, and woken equations are retried by the existing loop. Huet's branching is
don't-know nondeterminism, which a miniKanren stream represents directly. It should be a goal the user
calls, not the default, because the search may not terminate (\cref{sec:hounif-hard}).

\emph{The obstacle.} Projection needs the target type and arities, and imitation the arity and argument
types of the constant: \enquote{Huet's procedure cannot run without types} (\cref{sec:hounif-huet}).
So this item presupposes (f).

\emph{Book concept.} Huet's procedure and its matching trees \book{214--220}, and the limit of
pattern unification that \cref{sec:hoprog-panacea} discusses.

\subsection{(e) Predicate quantification}\label{sec:roadmap-later-pred}

\emph{What.} Goals such as \lp{pi rev\ (... rev ...) => rev L K nil}, which introduce a hidden predicate
local to one call \book{125--126}.

\emph{How it fits.} Through dynamic relations created at run time. Γ is an alist keyed by relation, so a
\texttt{∀} over a predicate would allocate a new key, a new dynamic relation with no static clauses, and
pass it to the body. Clauses assumed for it are then visible only inside that body, and no program text
outside can name it. This is a sketch: the proposal does not implement it.

\emph{Book concept.} \hohh$^+$ (\cref{def:hoprog-hohhplus}): a clause head may be a variable bound by an
essentially universal quantifier. By the time the clause is assumed, $\forall R$ has replaced the variable
by a new constant, so the operational behaviour does not change. Only the visibility of names does.

\subsection{(f) Simple types}\label{sec:roadmap-later-types}

\emph{What.} Optional simple types on λ-terms and on logic variables.

\emph{How it fits.} Through the interplay of the \texttt{atom?} and \texttt{tag} hooks, which the proposal
names for this item without designing it. One reading: a type is a tag on a term, and a typed variable can
be η-expanded by its type. Unification could then work on η-long forms,
as λProlog does, instead of η-expanding only when one side is an abstraction
(\cref{sec:roadmap-lim-untyped}).

\emph{What it buys.} Type-directed η, β that is guaranteed to terminate, the existence check on
possibly empty types, and with them item (d).

\emph{Book concept.} The simply typed terms of \cref{ch:lambda} and the βη-long normal forms on which
the book's unification algorithms are stated (\cref{sec:lambda-normal}).

% ------------------------------------------------------------------------------------------------
\section{Open questions for the author}\label{sec:roadmap-questions}

The proposal makes a choice wherever one was needed to build it and test it. Some of those choices are
the maintainer's to make, and are listed here as questions. Each says what the proposal does now and
what is at stake.

\begin{enumerate}
  \item \emph{The module's name.} The proposal keeps \texttt{(aux kanren lambda)}. The layer is λProlog's
    logic, and λKanren is the name the literature gives this line of work, Ma and Friedman's paper
    included~\cite{MaFriedman2021}. A name such as \texttt{(aux kanren hohh)} would describe only the
    goals, while the module also gives λ-tree syntax and $\Llam$ unification. The names inside follow
    \aux: goals end in \texttt{°}, their applicative forms in \texttt{ª}, term constructors have no
    suffix (\scm{ƛ}, \scm{@}), and internals carry the \texttt{μkanren-} prefix.
  \item \emph{The extensional reading.} \scm{μkanren-λ-reading} exposes a second meaning of
    \scm{≠°} under \texttt{∀°} (\cref{def:roadmap-readings}). It costs code and one edge case,
    the reading of a migrated disequality (\cref{sec:roadmap-lim-readings}). Keep it public, keep it as
    an internal option, or remove it and keep only λProlog's generic reading?
  \item \emph{Should micro own the protocol?} The proposal puts a generic extension protocol in micro,
    a global registry, and one client. The alternatives were a layered module that leaves micro untouched
    (d1, which leaked eigenvariables through micro's own \scm{=°}) and λ-terms native in micro (d3, a
    931-line diff that special-cases every walker). A per-state hook list instead of a global registry
    would require every initial state to carry the hooks, since \scm{°->list} starts from
    \scm{μkanren-state-empty}. Is a protocol in micro, with its contracts and invariants, something the
    maintainer wants to keep stable?
  \item \emph{The cost after a \texttt{∀°}.} First-order search after any \texttt{∀°} takes 1.7 to 2.3
    times as long (\cref{tab:roadmap-perf-forall}). The cheap fix gives up the host-escape guarantee
    (\cref{sec:roadmap-lim-host}); the proper one is roadmap item (a). Should stage 2 be merged now and
    (a) done later, or should (a) be a condition of the merge? And is it worth this price to catch
    escapes made with \scm{project°} and \scm{set!}?
  \item \emph{Which relations are dynamic.} In λProlog every predicate takes hypotheses; here only those
    made by \scm{define-dynamic-relation}, and \texttt{⊃°} on any other is an error
    (\cref{sec:roadmap-lim-dynamic}). Should \scm{define-relation} itself become dynamic when the module
    is loaded? The proposal does not measure what that would cost.
  \item \emph{α-equivalence and \scm{equal?}.} Hints are interned, so terms built alike are
    \scm{equal?}, but α-variants with different names are not (\cref{sec:roadmap-lim-hints}). Truly
    nameless records would make \scm{equal?} decide α, at the price of answers printed with generated
    names instead of the user's. Which matters more for the intended use?
  \item \emph{An answer inside \texttt{∀°}.} Reifying a variable created inside a \texttt{∀°} that holds its
    eigenvariable is an error, where λProlog reports no binding (\cref{sec:roadmap-lim-answer}). Keep the
    error, or answer with no binding as λProlog does?
  \item \emph{A defect of micro left as it is.} Under a type of the user's that holds of pairs, micro's
    subsumption turns \scm{absent° '(a) q} into \scm{q ≠ (a)}, so
    \scm{(absent° '(a) q) (nonnum° q) (=° q '((a)))} succeeds: the subsumption assumes that typed values
    are atoms. It involves no λ-term, so the proposal did not touch it. Fix it in micro, with the protocol
    or separately?
  \item \emph{Two clause syntaxes.} Hypothetical clauses may be written infix, \scm{(H ← B)}, or prefix,
    \scm{(← H B)}; both are accepted. Keep both, or settle on one?
\end{enumerate}

% ------------------------------------------------------------------------------------------------
\section{Summary}\label{sec:roadmap-summary}

\begin{itemize}
  \item The module has two suites: 83 cases with 320 hand-written assertions (acceptance A1--A9 and the
    regressions R, R2, R3 of three adversarial review rounds), and 6 cases that pin its limits. micro's six
    suites, 125 cases, pass with the edited micro both without and with the module loaded, answer order
    unchanged: 339 cases, 0 failures.
  \item The suites test behaviour on examples, under several goal orders; the guarantees on pattern
    problems come from the theorems of Miller, of Nadathur, Jayaraman and Kwon, and of Nadathur and
    Linnell, which the algorithm follows.
  \item Nine limitations are documented and pinned. Only dynamic relations take hypotheses, and there is
    no predicate quantification. Terms are untyped. Non-patterns are suspended, never enumerated, so
    floundering is reported, not detected. Answer order is micro's interleaving. \scm{≠°},
    \scm{absent°} and types use the generic reading of \texttt{∀}, with an extensional option for
    \scm{≠°}. Levels police logic variables, not Scheme values. Binder names are hints. Constraints are
    retried unindexed, and the scope check runs after any \texttt{∀°}. The answer variable must be
    created outside every \texttt{∀°}.
  \item The protocol costs about 1\% on first-order search, and the loaded extension 8--12\% on a
    search that keeps many disequalities. After a \texttt{∀°}, first-order search takes 1.7--2.3 times as long,
    because the scope check walks the bound term beside the occurs check. On λ workloads the module is
    1.7 times faster than the d2 prototype it grew from and 2.4 times slower than the native d3 on
    Church 400, the rest of the time being in micro's walks.
  \item The merge has three stages. First the protocol, into micro's source. Then the module, its
    suites, \texttt{programs.scm} and one egg line. Last the show extension. Review the protocol's invariants
    and its one change to micro's semantics first.
  \item Six later items: one walk for the occurs and scope checks; var--var orientation by level; an
    index of E and N; an optional Huet search, which needs types; predicate quantification through dynamic
    relations made at run time; simple types.
  \item Nine questions are the author's to answer, among them the module's name, the extensional reading,
    whether micro should own a protocol, and whether the cost after a \texttt{∀°} must be paid down before
    the merge.
\end{itemize}


\appendix
\part*{Appendices}
% appA-rules.tex — Appendix: reference card of rules and algorithms (placed after \appendix).
% Local notation (prefixed RL to avoid clashes with other chapters).
\providecommand{\RLfoc}[3]{#1;#2\xrightarrow{\;#3\;}}   % focused sequent  Σ; P --D--> A
\providecommand{\RLst}[2]{\langle #1;\ #2\rangle}         % unification state ⟨Q; E⟩
\providecommand{\RLto}{\;\Longrightarrow\;}

\chapter{Reference card: rules and algorithms}\label{app:rules}

This appendix lists the formal rules used in the document, each with the page where its source states it; \cref{ch:horn,ch:hh,ch:lambda,ch:hoprog,ch:hounif} explain them.
Page references \book{n} use the printed page numbers of \cite{MillerNadathur2012}; papers are cited with their own page numbers.

% =====================================================================================
\section{λProlog concrete syntax and its logical reading}\label{sec:rules-syntax}

\Cref{tab:rules-syntax} maps every λProlog construct used in the document to the logical formula it denotes.

{\small
\begin{longtable}{@{}>{\raggedright\arraybackslash}p{0.30\textwidth}>{\raggedright\arraybackslash}p{0.27\textwidth}>{\raggedright\arraybackslash}p{0.36\textwidth}@{}}
\caption{λProlog syntax and logic.}\label{tab:rules-syntax}\\
\toprule
λProlog & logic & notes \\
\midrule\endfirsthead
\toprule
λProlog & logic & notes \\
\midrule\endhead
\bottomrule\endfoot
\lp{true} & $\top$ & type $\typ{o}$ \book{35} \\
\lp{G1, G2} & $G_1\wedge G_2$ & infix, left assoc.; joins goals \book{35,47--48} \\
\lp{D1 & D2} & $D_1\wedge D_2$ & infix, right assoc.; joins clauses \book{35,47--48} \\
\lp{G1 ; G2} & $G_1\vee G_2$ & infix, left assoc. \book{35} \\
\lp{D => G} & $D\imp G$ & infix, right assoc.; implication in goals \book{35,79} \\
\lp{A :- G} & $G\imp A$ & infix, non-assoc.; \enquote{is implied by}, used in clauses \book{35,47--48} \\
\lp{pi x\ B}, \lp{pi (x:T)\ B} & $\forall x\,B$, $\forall_\tau x\,B$ & $\forall:(A\to\typ{o})\to\typ{o}$ applied to $\abs{x}{B}$ \book{36--37,99} \\
\lp{sigma x\ B} & $\exists x\,B$ & $\exists:(A\to\typ{o})\to\typ{o}$ applied to $\abs{x}{B}$ \book{36--37,99} \\
\lp{x\ t}, \lp{(x:T)\ t} & $\abs{x}{t}$, $\lambda(x{:}\tau)\,t$ & body extends as far right as possible \book{97} \\
\lp{t1 t2} & $(t_1\,t_2)$ & application, left assoc., binds tightest \book{97} \\
\lp{X}, \lp{L1} (free, uppercase) & $\forall$ at clause top; $\exists$ in a query & implicit quantification \book{43,46--47} \\
\lp{_} & a fresh variable per occurrence & anonymous, $\forall$ over the clause \book{46--47} \\
\lp{X = Y} & $X = Y$ & built-in equality, solved by unification \book{106,287} \\
\lp{false} (user constant) & $\bot$ of minimal logic & no ex falso for it; $\neg B$ is $B\imp\bot$ \book{89--90} \\
\lp{A -> B} (in types) & $\tau\to\sigma$ & right assoc. \book{12--13} \\
\lp{o} & $\typ{o}$ & the type of propositions \book{34--35} \\
\lp{kind c type -> type.} & a type constructor of arity 1 & kind declaration \book{11--12} \\
\lp{type c T.} & $c:\tau\in\Sig$ & constant declaration \book{47} \\
\end{longtable}
}

Binding strength, tightest first: application $>$ \lp{=>} $>$ \lp{&} $>$ \lp{,} $>$ \lp{;} $>$ \lp{:-} \book{35}; the scope of \lp{x\ } extends as far right as possible \book{37,97}.

\begin{pitfall}[precedence]
\lp{D => G1, G2} parses as \lp{(D => G1), G2}, so a hypothesis meant for a whole conjunction needs parentheses \book{35}.
\end{pitfall}

% =====================================================================================
\section{Goal and clause grammars}\label{sec:rules-grammars}

The four languages share one top-level shape and differ in which connectives goals may contain and in what an atom may be (\cref{tab:rules-grammars}).

\begin{center}
\small
\captionof{table}{Goal formulas $G$ and program clauses $D$. $A$: atom; $A_r$: rigid atom.}\label{tab:rules-grammars}
\begin{tabular}{@{}lll@{}}
\toprule
language & grammar & source \\
\midrule
\fohc & $G ::= \top \mid A \mid G\wedge G \mid G\vee G \mid \exists_\tau x\,G$ & \book{43} \\
      & $D ::= A \mid G\imp D \mid D\wedge D \mid \forall_\tau x\,D$ & \\[2pt]
\fohh & $G ::= \top \mid A \mid G\wedge G \mid G\vee G \mid \exists_\tau x\,G \mid D\imp G \mid \forall_\tau x\,G$ & \book{75} \\
      & $D ::= A \mid G\imp D \mid D\wedge D \mid \forall_\tau x\,D$ & \\[2pt]
\hohc & as \fohc, with $A\in\mathcal H_1^\Sig$ in goals and heads $A_r\in\mathcal H_1^\Sig$ & \book{122--123} \\[2pt]
\hohh & as \fohh, with $A\in\mathcal H_2^\Sig$ in goals and heads $A_r\in\mathcal H_2^\Sig$ & \book{124} \\[2pt]
$\hohh^+$ & \hohh, plus atomic heads whose head is a variable bound by an & \book{125} \\
          & \quad essentially universal quantifier that has the clause in its scope & \\
\bottomrule
\end{tabular}
\end{center}

\begin{definition}[Herbrand universes \book{122}]\label{def:rules-herbrand}
$\mathcal H_1^\Sig$ is the set of (possibly open) λ-normal terms over $\Sig$ that contain neither $\forall$ nor $\imp$; $\mathcal H_2^\Sig$ is the set of those that do not contain $\imp$.
\end{definition}

\begin{definition}[rigid and flexible atoms \book{119}]\label{def:rules-rigidatom}
An atom is \emph{flexible} if the head $h$ of its λ-normal form $(h\,t_1\dots t_n)$ is a variable, and \emph{rigid} if $h$ is a non-logical constant.
\end{definition}

\begin{definition}[polarity; hereditary Harrop \book{76}]\label{def:rules-polarity}
$B$ is positive in $B$; a positive (negative) occurrence in $B$ stays positive (negative) in $B\wedge B'$, $B'\wedge B$, $B\vee B'$, $B'\vee B$, $B'\imp B$, $\forall x\,B$, $\exists x\,B$, and becomes negative (positive) in $B\imp B'$.
A hereditary Harrop formula is one in which no positive subformula occurrence is a disjunction or an existential quantification.
\end{definition}

\begin{definition}[essentially universal / existential \book{205}]\label{def:rules-essential}
In a goal, a bound-variable occurrence is \emph{essentially universal} if it is bound by a positive $\forall$, a negative $\exists$ or a term-level $\lambda$, and \emph{essentially existential} otherwise; in a clause the classification is dual (negative $\forall$, positive $\exists$ or $\lambda$ are essentially universal).
\end{definition}

\begin{definition}[$\Llam$ \book{205--206,220}]\label{def:rules-llambda}
A goal or clause $B$ is in $\Llam$ if all quantification in $B$ is over non-predicate types and every subterm $(x\,t_1\dots t_n)$ of $B$ with $x$ essentially existential has $t_1,\dots,t_n$ distinct essentially universal variables bound within the scope of $x$.
\end{definition}

% =====================================================================================
\section{Uniform proofs}\label{sec:rules-uniform}

A sequent $\seq{\Sig}{\Prog}{G}$ says that $G$ is provable from program $\Prog$ over signature $\Sig$; $\RLfoc{\Sig}{\Prog}{D}A$ says that $A$ is proved by backchaining on the clause $D$ \book{39,45}.
Read bottom-up, the right rules (goal reduction) and the left rules (backchaining) are the two phases of a \emph{uniform}, i.e.\ goal-directed, proof; a logic in which such proofs are complete is an \emph{abstract logic programming language} \book{73}.

\begin{center}
\small
\begin{mathpar}
\inferrule*[right=\rulename{true}]{ }{\seq{\Sig}{\Prog}{\top}}
\and
\inferrule*[right=\rulename{and}]{\seq{\Sig}{\Prog}{B_1} \\ \seq{\Sig}{\Prog}{B_2}}{\seq{\Sig}{\Prog}{B_1\wedge B_2}}
\and
\inferrule*[right={\rulename{or}\ \normalfont $(i=1,2)$}]{\seq{\Sig}{\Prog}{B_i}}{\seq{\Sig}{\Prog}{B_1\vee B_2}}
\and
\inferrule*[right=\rulename{instan}]{\seq{\Sig}{\Prog}{\subst{B}{t}{x}} \\ \Sig;\emptyset\Vdash t:\tau}{\seq{\Sig}{\Prog}{\exists_\tau x\,B}}
\and
\inferrule*[right=\rulename{augment}]{\seq{\Sig}{\Prog,B_1}{B_2}}{\seq{\Sig}{\Prog}{B_1\imp B_2}}
\and
\inferrule*[right={\rulename{generic}\ \normalfont $(c\notin\Sig)$}]{\seq{c{:}\tau,\Sig}{\Prog}{\subst{B}{c}{x}}}{\seq{\Sig}{\Prog}{\forall_\tau x\,B}}
\end{mathpar}
\captionof{figure}{Right rules (search semantics of the connectives) \book{39--40}. In the book they are $\top R$, $\wedge R$, $\vee R$, $\exists R$, $\imp R$, $\forall R$; $\Prog,B$ is $\Prog\cup\{B\}$.}\label{fig:rules-right}
\end{center}

\begin{center}
\small
\begin{mathpar}
\inferrule*[right={\rulename{decide}\ \normalfont $(D\in\Prog)$}]{\RLfoc{\Sig}{\Prog}{D}A}{\seq{\Sig}{\Prog}{A}}
\and
\inferrule*[right=\rulename{initial}]{ }{\RLfoc{\Sig}{\Prog}{A}A}
\and
\inferrule*[right=${\imp}L$]{\RLfoc{\Sig}{\Prog}{D}A \\ \seq{\Sig}{\Prog}{G}}{\RLfoc{\Sig}{\Prog}{G\imp D}A}
\and
\inferrule*[right={${\wedge}L$\ \normalfont $(i=1,2)$}]{\RLfoc{\Sig}{\Prog}{D_i}A}{\RLfoc{\Sig}{\Prog}{D_1\wedge D_2}A}
\and
\inferrule*[right=${\forall}L$]{\RLfoc{\Sig}{\Prog}{\subst{D}{t}{x}}A \\ \Sig;\emptyset\Vdash t:\tau}{\RLfoc{\Sig}{\Prog}{\forall_\tau x\,D}A}
\\
\inferrule*[right=\rulename{backchain}]{\seq{\Sig}{\Prog}{A_1\theta} \\ \cdots \\ \seq{\Sig}{\Prog}{A_n\theta}}{\RLfoc{\Sig}{\Prog}{\forall\bar x\,(A_1\wedge\dots\wedge A_n\imp A_0)}A}
\\
\text{($\theta$ maps each $x_i$ to a $\Sig$-term of type $\tau_i$, and $A=A_0\theta$)}
\end{mathpar}
\captionof{figure}{Backchaining (left) rules \book{45}; \rulename{backchain} is the compiled form for Prolog-shaped clauses \book{45}.}\label{fig:rules-left}
\end{center}

\begin{definition}[O-proof \book{45,77}]\label{def:rules-oproof}
A derivation built from the typing rules of terms and the rules of \cref{fig:rules-right,fig:rules-left}.
\end{definition}

\paragraph{Logic-variable versions \book{61}.} \rulename{instan} and $\forall L$ replace the guessed $t$ by a new logic variable $X$ with $\Sig;\emptyset\Vdash X:\tau$, and \rulename{initial} becomes $\RLfoc{\Sig}{\Prog}{A'}A$ provided $A$ and $A'$ unify, the unifier being applied to the whole derivation built so far.

\paragraph{Higher-order versions \book{123--124}.} (i) Formulas are taken in λ-normal form when matched against a rule; (ii) \rulename{initial} requires $A'$ and $A$ to be α-convertible; (iii) the term $t$ of \rulename{instan} and $\forall L$ is a λ-term with $\Sig;\emptyset\Vdash t:\tau$ and $t\in\mathcal H_1^\Sig$ (\hohc) or $t\in\mathcal H_2^\Sig$ for the current $\Sig$ (\hohh).

\paragraph{Backchaining by elaboration \cite[p.\,9]{Miller1991}.}
With $\langle\emptyset,D\rangle\in|\Gamma|_\Sig$ for $D\in\Gamma$, conjunctions split, $\langle\Delta,G\imp D\rangle\mapsto\langle\Delta\cup\{G\},D\rangle$ and $\langle\Delta,\forall x\,D\rangle\mapsto\langle\Delta,\lambda\mathrm{norm}(\subst{D}{t}{x})\rangle$:
\begin{mathpar}
\inferrule*[right={\rulename{bc}\ \normalfont ($A$ atomic, $\langle\Delta,A\rangle\in|\Gamma|_\Sig$)}]{\{\seq{\Sig}{\Gamma}{G}\}_{G\in\Delta}}{\seq{\Sig}{\Gamma}{A}}
\end{mathpar}

\paragraph{Search conventions of λProlog \book{60,78}.} Conjunctive goals left to right; the left disjunct of a goal and the left conjunct of a clause first; clauses tried in program order (the program is a list); \rulename{augment} adds its clauses at the \emph{front}; signature and program grow and shrink in a stack discipline.

\paragraph{Flexible goals \book{134}.} A goal $(P\,t_1\dots t_n)$ with $P$ unbound can be (1) suspended until rigid, and if only flexible goals remain each head is bound to $\lambda x_1\dots\lambda x_n\,\top$ (complete for intuitionistic higher-order logic); (2) bound eagerly to $\lambda x_1\dots\lambda x_n\,\top$; or (3) reported as an error; (2) and (3) are incomplete, yet experience favours (3), since a flexible goal is usually a programming mistake such as a capitalised predicate name.

\begin{center}
\small
\captionof{table}{What O-proofs are complete for.}\label{tab:rules-completeness}
\begin{tabular}{@{}lp{0.62\textwidth}l@{}}
\toprule
language & $\seq{\Sig}{\Prog}{G}$ has an O-proof iff it is provable in & source \\
\midrule
\fohc & classical logic, iff intuitionistic logic & \book{45,62} \\
\fohh & minimal logic, equivalently intuitionistic logic (classical logic is not sound for this operational semantics) & \book{89,91} \\
\hohc & classical, equivalently intuitionistic, Simple Theory of Types without extensionality, infinity and choice & \book{124,148} \\
\hohh & higher-order intuitionistic logic & \book{124} \\
\bottomrule
\end{tabular}
\end{center}

\begin{proposition}[flatness of \fohc\ \book{46}]\label{prop:rules-flat}
Every sequent in an O-proof of an \fohc\ sequent $\seq{\Sig}{\Prog}{G}$ has the form $\seq{\Sig}{\Prog}{G'}$ or $\RLfoc{\Sig}{\Prog}{D}A$: signature and program never change.
\end{proposition}

\begin{kanren}[which rules a miniKanren already has]
\micro\ has goals for \rulename{true} (\scm{✓°}), \rulename{and} (\scm{and°}), \rulename{or} (\scm{or°}) and logic-variable \rulename{instan} (\scm{fresh°}); a relation defined as a Scheme procedure plays the role of \rulename{decide}/\rulename{backchain}, and its unifier is the first-order one of \cref{fig:rules-fo}, occurs check included.
Two differences matter: \scm{or°} interleaves its branches instead of following the depth-first order above, and \rulename{augment} and \rulename{generic} have no counterpart, so \micro\ covers \fohc\ but not \fohh; see \cref{ch:design}.
\end{kanren}

% =====================================================================================
\section{Typed λ-terms}\label{sec:rules-typing}

The typing judgement $\Sig;\Gamma\Vdash t:\tau$ assigns types to λ-terms; quantifiers need no rule of their own because they are constants applied to abstractions \book{98--99}.

\begin{center}
\small
\begin{mathpar}
\inferrule*[right={$(\tau\triangleleft\sigma)$}]{c:\sigma\in\Sig}{\Sig;\Gamma\Vdash c:\tau}
\and
\inferrule{x:\tau\in\Gamma}{\Sig;\Gamma\Vdash x:\tau}
\and
\inferrule{\Sig;\Gamma\Vdash g:\tau_1\to\tau_2 \\ \Sig;\Gamma\Vdash t:\tau_1}{\Sig;\Gamma\Vdash (g\,t):\tau_2}
\and
\inferrule*[right=$(\dagger)$]{\Sig;\Gamma,x:\tau\Vdash t:\sigma}{\Sig;\Gamma\Vdash \lambda(x{:}\tau)\,t:\tau\to\sigma}
\and
\inferrule*[right=$(\ddagger)$]{\Sig;\Gamma\Vdash B:\tau}{\Sig;\Gamma\Vdash C:\tau}
\end{mathpar}
\captionof{figure}{Typing rules \book{98}. $\tau\triangleleft\sigma$: $\tau$ is an instance of $\sigma$; ($\dagger$) $x$ not declared in $\Sig$ or $\Gamma$; ($\ddagger$) $B$ and $C$ differ only in bound-variable names.}\label{fig:rules-typing}
\end{center}

\begin{definition}[order of a type \book{13}]\label{def:rules-order}
$\mathrm{ord}(\tau)=0$ if $\tau$ is non-functional; $\mathrm{ord}(\tau_1\to\tau_2)=\max(\mathrm{ord}(\tau_1)+1,\mathrm{ord}(\tau_2))$.
\end{definition}

% =====================================================================================
\section{λ-conversion and normal forms}\label{sec:rules-conversion}

λ-conversion is the equality of the logic, and every typed term has a unique normal form, so equality is decidable \book{101--102}.

\begin{definition}[free for; substitution \book{85,100}]\label{def:rules-freefor}
$t$ is \emph{free for} $x$ in $s$ if no free occurrence of $x$ in $s$ is in the scope of an abstraction (or quantifier) binding a variable free in $t$.
$\subst{s}{t}{x}$ replaces the free occurrences of $x$ in $s$ by $t$, after α-renaming $s$ so that $t$ is free for $x$ (capture-avoiding substitution) \book{101}.
\end{definition}

\begin{center}
\small
\captionof{table}{Rewrite rules on subterms \book{100--101}.}\label{tab:rules-conv}
\begin{tabular}{@{}lll@{}}
\toprule
rule & redex $\Rightarrow$ contractum & proviso \\
\midrule
α-rewriting & $\abs{x}{s} \Rightarrow \abs{y}{\subst{s}{y}{x}}$ & $y$ free for $x$ in $s$, $y$ not free in $s$ \\
β-contraction & $(\abs{x}{s})\,t \Rightarrow \subst{s}{t}{x}$ & none (capture-avoiding $\subst{s}{t}{x}$) \\
η-contraction & $\abs{x}{(s\,x)} \Rightarrow s$ & $x$ not free in $s$ \\
$\beta_0$ \book{206} & $(\abs{x}{s})\,x \Rightarrow s$; up to α: $(\abs{y}{s})\,x\Rightarrow\subst{s}{x}{y}$ & (second form) $x$ not free in $\abs{y}{s}$ \\
\bottomrule
\end{tabular}
\end{center}

\begin{definition}[closures \book{100--101}]\label{def:rules-closures}
α-conversion: reflexive, symmetric, transitive closure of α-rewriting.
β-reduction: reflexive, transitive closure of α-conversion $\cup$ β-contraction; β-conversion: its symmetric, transitive closure.
η-reduction / η-conversion: likewise from η-contraction.
λ-conversion ($\eqb$): transitive closure of α-, β- and η-conversion; it is a congruence and is the equality of higher-order logic.
\end{definition}

\begin{definition}[normal forms \book{101--102,111}]\label{def:rules-nf}
A term is \emph{β-normal} if it has no β-redex, \emph{λ-normal} (βη-normal) if it is β-normal with no η-redex.
Every β-normal term has the shape $\lambda x_1\dots\lambda x_n\,(h\,t_1\dots t_p)$ ($n,p\ge0$), with \emph{binder} $x_1\dots x_n$, \emph{head} $h$ (a constant or variable), \emph{arguments} $t_i$ (β-normal) and \emph{body} $(h\,t_1\dots t_p)$.
It is in \emph{βη-long normal form} if moreover each $t_i$ is βη-long and the body has non-functional (for simple types: primitive) type.
\end{definition}

\begin{theorem}[normalisation \book{102}]\label{thm:rules-norm}
Every typed λ-term has a λ-normal form, unique up to α-conversion; it is computed by β-contracting until β-normal and then η-contracting.
Two terms of the same type are λ-convertible iff their λ-normal forms are α-equivalent.
Every β-normal form can be η-expanded to an equal βη-long form \book{111}.
\end{theorem}

\begin{example}[η-long expansion \book{111}]
With $f:(i\to i)\to i$: $f \;\rightsquigarrow\; \abs{w}{f\,w}$ ($w:i\to i$) $\;\rightsquigarrow\; \abs{w}{f\,(\abs{u}{w\,u})}$ ($u:i$), which is βη-long.
\end{example}

\begin{definition}[rigid and flexible terms \book{111--112,214}]\label{def:rules-rigid}
A β-normal term $\lambda\bar x\,(h\,\bar t)$ in a unification problem is \emph{rigid} if $h$ is universally quantified in the prefix (a constant) or a member of its binder $\bar x$, and \emph{flexible} if $h$ is existentially quantified; equations are rigid-rigid, flexible-rigid, rigid-flexible or flexible-flexible \book{113}.
\end{definition}

% =====================================================================================
\section{Unification problems, raising, Skolemization}\label{sec:rules-raising}

Higher-order unification problems carry a mixed quantifier prefix, which raising normalises to $\forall\exists\forall$ \book{105--108,212}.

\begin{definition}[unification problem \book{105,211}]\label{def:rules-problem}
A formula $Q_1x_1\dots Q_nx_n\,[t_1=s_1\wedge\dots\wedge t_m=s_m]$ with $Q_i\in\{\forall,\exists\}$ and $t_j,s_j$ λ-terms of the same type; the quantifiers form the \emph{prefix}, the conjunction the \emph{matrix}.
\end{definition}

\begin{definition}[unifier, solution \book{109}]\label{def:rules-solution}
For $\forall u_1\dots\forall u_p\,\exists x_1\dots\exists x_n\,\forall w_1\dots\forall w_q\,[t_1=s_1\wedge\dots\wedge t_m=s_m]$, a substitution $\theta$ for $x_1,\dots,x_n$ is \emph{legitimate} if no $w_j$ is free in its range; a \emph{unifier} is a legitimate $\theta$ with $\theta(t_j)\eqb\theta(s_j)$ for all $j$; a \emph{solution} is a unifier whose range has free variables only among $u_1,\dots,u_p$.
Only a solution guarantees provability, since types are not assumed inhabited \book{109--110}.
\end{definition}

\begin{example}[the prefix order matters \book{110,213}]\label{ex:rules-prefix}
With $g:i\to i\to i$, $\forall a\,\exists F\,[(F\,a)=(g\,a\,a)]$ has the four unifiers $F:=\abs{x}{g\,a\,a}$, $\abs{x}{g\,x\,a}$, $\abs{x}{g\,a\,x}$, $\abs{x}{g\,x\,x}$, none an instance of another, so there is no most general unifier; in $\exists F\,\forall a\,[(F\,a)=(g\,a\,a)]$ the solution $F$ cannot mention $a$, and $\abs{a}{g\,a\,a}$ is the only one.
\end{example}

\begin{proposition}[raising \book{107}]\label{prop:rules-raising}
With $x:\tau$, $y:\sigma$ and $h:\sigma\to\tau$: $\forall y\,\exists x\,G$ is provable iff $\exists h\,\forall y\,\subst{G}{(h\,y)}{x}$ is provable.
\end{proposition}
As a rewrite on prefixes, for a block of universals:
\[
\forall y_1{:}\sigma_1\dots\forall y_k{:}\sigma_k\;\exists x{:}\tau\;G \;\rightsquigarrow\; \exists h{:}\sigma_1\to\dots\to\sigma_k\to\tau\;\forall y_1\dots\forall y_k\;\subst{G}{(h\,y_1\dots y_k)}{x}.
\]
Repeated raising turns any prefix into $\forall\exists\forall$ or even $\exists\forall$ form \book{108}, and it preserves $\Llam$ \book{208}.
Miller's transition form \cite[p.\,14--16]{Miller1991}: if $\exists v$ is to the left of $\exists u$ with universals $\bar w$ in between, replace $u$ by $(u'\,\bar w)$ and move $\exists u'$ next to $\exists v$.

\paragraph{Skolemization \book{108}.} An assumption $\forall x\,\exists y\,D$ becomes $\forall x\,\subst{D}{(f\,x)}{y}$ with $f:\sigma\to\tau$ a new constant; dually a goal $\exists x\,\forall y\,G$ becomes $\forall f\,\exists x\,\subst{G}{(f\,x)}{y}$.
Skolemization moves an $\exists$ inwards and introduces a constant of raised type; raising moves a $\forall$ inwards and introduces an existential variable of raised type; relating unifiers of Skolemized problems to the original ones is problematic for higher-order variables, whereas raising is an exact equivalence \book{107--108}.

\paragraph{Normal form of problems \book{212}.} Before solving, (1) the prefix is raised to $\forall\bar a\,\exists\bar X\,\forall\bar y$, with $\forall\bar a$ the signature; (2) $\theta(t)$ is defined as $((\lambda x_1\dots\lambda x_n\,t)\,t_1\dots t_n)$ for $\theta=\{\langle x_i,t_i\rangle\}$, so terms may be taken β-normal; (3) η-expansion makes the two sides of each equation have binders of the same length; (4) by $(\abs{x}{t}=\abs{x}{s})\equiv\forall x\,(t=s)$ the binders are moved into the inner $\forall$ block:
\begin{mathpar}
\inferrule*[right=\rulename{abs}]{\RLst{\mathcal Q}{\{\abs{x}{s}=\abs{x}{t}\}\uplus E}}{\RLst{\mathcal Q\,\forall x}{\{s=t\}\uplus E}}
\and
\inferrule*[right={\rulename{eta}\ \normalfont ($t$ not an abstraction)}]{\RLst{\mathcal Q}{\{\abs{x}{s}=t\}\uplus E}}{\RLst{\mathcal Q}{\{\abs{x}{s}=\abs{x}{(t\,x)}\}\uplus E}}
\end{mathpar}
Here and below a rule is read downwards as one transformation step, $\uplus$ is multiset union, $x$ is fresh, and \rulename{eta} also applies with the two sides exchanged.

% =====================================================================================
\section{First-order unification}\label{sec:rules-fo}

Term reduction, reorientation and variable elimination with occurs check; higher-order pattern unification has the same two critical parts, simplification and variable elimination \book{224}.

\begin{center}
\small
\begin{mathpar}
\inferrule*[right=\rulename{decompose}]{E\uplus\{f\,t_1\dots t_n=f\,s_1\dots s_n\}}{E\uplus\{t_1=s_1,\dots,t_n=s_n\}}
\and
\inferrule*[right={\rulename{clash}\ \normalfont ($f\neq g$ or different types)}]{E\uplus\{f\,\bar t=g\,\bar s\}}{\bot}
\and
\inferrule*[right={\rulename{orient}\ \normalfont ($t$ not a variable)}]{E\uplus\{t=x\}}{E\uplus\{x=t\}}
\and
\inferrule*[right=\rulename{delete}]{E\uplus\{x=x\}}{E}
\and
\inferrule*[right={\rulename{check}\ \normalfont ($x\in\mathrm{fv}(t)$, $t\neq x$)}]{E\uplus\{x=t\}}{\bot}
\and
\inferrule*[right={\rulename{eliminate}\ \normalfont ($x\notin\mathrm{fv}(t)$, $x\in\mathrm{fv}(E)$)}]{E\uplus\{x=t\}}{\subst{E}{t}{x}\uplus\{x=t\}}
\end{mathpar}
\captionof{figure}{Term reduction, reorientation and variable elimination \book{28}; the occurs check is mandatory \book{32}.}\label{fig:rules-fo}
\end{center}

\begin{theorem}[\book{28--29}]\label{thm:rules-fo}
Every sequence of steps (never eliminating twice on the same unchanged equation) terminates in $\bot$ or in a \emph{solved form} (each left side a variable occurring nowhere else), which is a most general unifier, unique up to renaming.
\end{theorem}

% =====================================================================================
\section{Huet's procedure}\label{sec:rules-huet}

On a normalised problem (\cref{sec:rules-raising}) Huet's procedure simplifies rigid-rigid equations and guesses the head of a flexible term facing a rigid one \book{214--216}.

\begin{center}
\small
\begin{mathpar}
\inferrule*[right={\rulename{simp}\ \normalfont ($c$ universal)}]{\RLst{\mathcal Q}{\{c\,s_1\dots s_n=c\,t_1\dots t_n\}\uplus E}}{\RLst{\mathcal Q}{\{s_1=t_1,\dots,s_n=t_n\}\uplus E}}
\and
\inferrule*[right={\rulename{clash}\ \normalfont ($c\neq d$ universal)}]{\RLst{\mathcal Q}{\{c\,\bar s=d\,\bar t\}\uplus E}}{\bot}
\and
\inferrule*[right={\rulename{swap}\ \normalfont ($r$ rigid, $f$ flexible)}]{\RLst{\mathcal Q}{\{r=f\}\uplus E}}{\RLst{\mathcal Q}{\{f=r\}\uplus E}}
\\
\inferrule*[right=\rulename{imitate}]{\RLst{\mathcal Q_1\,\exists F\,\mathcal Q_2}{\{F\,t_1\dots t_n=c\,s_1\dots s_m\}\uplus E}}{\RLst{\mathcal Q_1\,\exists H_1\dots\exists H_m\,\mathcal Q_2}{\theta(\{F\,\bar t=c\,\bar s\}\uplus E)\!\downarrow_\beta}}
\\
\theta=\{\langle F,\ \lambda x_1\dots\lambda x_n\,(c\,(H_1\,x_1\dots x_n)\cdots(H_m\,x_1\dots x_n))\rangle\},\quad \forall c\in\mathcal Q_1,\quad H_i:\tau_1\to\dots\to\tau_n\to\sigma_i
\\
\inferrule*[right=\rulename{project}$_i$]{\RLst{\mathcal Q_1\,\exists F\,\mathcal Q_2}{\{F\,t_1\dots t_n=c\,s_1\dots s_m\}\uplus E}}{\RLst{\mathcal Q_1\,\exists H_1\dots\exists H_k\,\mathcal Q_2}{\theta(\{F\,\bar t=c\,\bar s\}\uplus E)\!\downarrow_\beta}}
\\
\theta=\{\langle F,\ \lambda x_1\dots\lambda x_n\,(x_i\,(H_1\,x_1\dots x_n)\cdots(H_k\,x_1\dots x_n))\rangle\},\quad \tau_i=\varphi_1\to\dots\to\varphi_k\to\beta,\quad H_j:\tau_1\to\dots\to\tau_n\to\varphi_j
\end{mathpar}
\captionof{figure}{Huet's transformation rules \book{214--216}, with $F:\tau_1\to\dots\to\tau_n\to\beta$, $c:\sigma_1\to\dots\to\sigma_m\to\beta$, $\beta$ primitive. \rulename{imitate} needs $F$ in the scope of $c$; \rulename{project}$_i$ needs the target of $\tau_i$ to be $\beta$; so at most $n+1$ choices exist per equation \book{216}. The rewriting format is ours, in the style of the transformation-based presentation of Huet's procedure \cite{Huet1975} mentioned on \book{227}; side conditions and substitution terms are the book's.}\label{fig:rules-huet}
\end{center}

\begin{definition}[pre-unification; canonical substitution \book{217--218}]\label{def:rules-preunif}
Reducing a problem to one with only flexible-flexible equations is \emph{pre-unification}, and the composed substitution is a \emph{pre-unifier}.
The \emph{canonical substitution} maps each existential of type $\tau_1\to\dots\to\tau_n\to\sigma$ ($\sigma$ primitive) to $\lambda x_1\dots\lambda x_n\,H^\sigma$, with one fixed new existential $H^\sigma$ per primitive type; on a flexible-flexible problem it is a unifier.
\end{definition}

\begin{theorem}[matching trees \book{219--220}]\label{thm:rules-matching}
In the matching tree (nodes: simplified matrices; arcs: imitation/projection substitutions for one selected flexible-rigid equation; leaves: failure, or success when only flexible-flexible equations remain), if the root problem has a unifier then a success node exists at finite depth, and that unifier is the composed pre-unifier followed by a unifier of the success node, whichever equation is selected at each node.
Since unifiability is undecidable, some trees have an infinite path and no success node \book{220}.
As a test for unifiability the procedure is nonredundant \book{227}; mgus may not exist, complete sets of unifiers may be infinite, and searching for them must be redundant \book{213,226}.
\end{theorem}

\begin{example}[a matching tree \book{216--219}]\label{ex:rules-huet}
For $\forall a\,\forall g\,\exists F\,[(F\,a)=(g\,a\,a)]$, \rulename{project}$_1$ ($F:=\abs{x}{x}$) gives $a=(g\,a\,a)$, a \rulename{clash}; \rulename{imitate} ($F:=\abs{x}{g\,(H_1\,x)\,(H_2\,x)}$) and \rulename{simp} give $\{(H_1\,a)=a,\ (H_2\,a)=a\}$, and each equation again has an imitation ($\abs{x}{a}$) and a projection ($\abs{x}{x}$), both succeeding.
The four leaves compose to the four unifiers of \cref{ex:rules-prefix}.
In $\forall g\,\exists F\,\forall x\,[(F\,(g\,x))=(g\,(F\,x))]$, imitation reproduces the same problem for a new $H$, so the tree has an infinite path, alongside the unifiers $\abs{x}{x}$, $\abs{x}{g\,x}$, $\abs{x}{g\,(g\,x)}$, \dots\ \book{218}.
\end{example}

% =====================================================================================
\section{Higher-order pattern unification}\label{sec:rules-pattern}

In $\Llam$ every flexible term is a pattern, at most one substitution is useful for each flexible-rigid equation, and the procedure becomes a deterministic algorithm with most general unifiers \book{221--223}.

\begin{definition}[pattern \book{220}]\label{def:rules-pattern}
A flexible term $(F\,c_1\dots c_n)$ ($n\ge0$) is a \emph{pattern} if $c_1,\dots,c_n$ are distinct universal variables quantified to the right of $F$ (within its scope); $\bar c\cap\bar d$ denotes some listing of the variables common to $\bar c$ and $\bar d$.
\end{definition}

\begin{theorem}[\book{208,221--224}]\label{thm:rules-pattern}
On $\Llam$ problems the rules of \cref{fig:rules-pattern} without \rulename{delay} terminate (thanks to the occurs check, \book{222--223}) and return $\bot$ or a most general unifier: $\Llam$ unification is decidable and unitary \cite{Miller1991}.
Types are never consulted to build the unifier \book{224}; the untyped variant fails rigid-rigid equations of different arities \cite[§\,9.3]{Miller1991}.
\end{theorem}

\begin{figure}[tbp]
\centering\small
\begin{mathpar}
\inferrule*[right=\rulename{abs}]{\RLst{\mathcal Q}{\{\abs{x}{s}=\abs{x}{t}\}\uplus E}}{\RLst{\mathcal Q\,\forall x}{\{s=t\}\uplus E}}
\and
\inferrule*[right={\rulename{eta}\ \normalfont ($t$ not an abstraction)}]{\RLst{\mathcal Q}{\{\abs{x}{s}=t\}\uplus E}}{\RLst{\mathcal Q}{\{\abs{x}{s}=\abs{x}{(t\,x)}\}\uplus E}}
\and
\inferrule*[right=\rulename{swap}]{\RLst{\mathcal Q}{\{r=F\,\bar c\}\uplus E}}{\RLst{\mathcal Q}{\{F\,\bar c=r\}\uplus E}}
\\
\inferrule*[right={\rulename{rr-dec}\ \normalfont ($h$ universal)}]{\RLst{\mathcal Q}{\{h\,s_1\dots s_n=h\,t_1\dots t_n\}\uplus E}}{\RLst{\mathcal Q}{\{s_1=t_1,\dots,s_n=t_n\}\uplus E}}
\and
\inferrule*[right={\rulename{rr-clash}\ \normalfont ($h\neq k$)}]{\RLst{\mathcal Q}{\{h\,\bar s=k\,\bar t\}\uplus E}}{\bot}
\\
\inferrule*[right={\rulename{ff-diff}\ \normalfont ($F\neq G$)}]{\RLst{\mathcal Q}{\{F\,\bar c=G\,\bar d\}\uplus E}}{\RLst{\mathcal Q\,\exists H}{\theta(E)}}
\quad \theta=\{\langle F,\lambda\bar c.\,H\,(\bar c\cap\bar d)\rangle,\ \langle G,\lambda\bar d.\,H\,(\bar c\cap\bar d)\rangle\}
\\
\inferrule*[right=\rulename{ff-same}]{\RLst{\mathcal Q}{\{F\,c_1\dots c_n=F\,d_1\dots d_n\}\uplus E}}{\RLst{\mathcal Q\,\exists H}{\theta(E)}}
\quad \theta=\{\langle F,\lambda\bar c.\,H\,\bar e\rangle\},\ \bar e=\text{the } c_i \text{ with } c_i=d_i
\\
\inferrule*[right=\rulename{elim}]{\RLst{\mathcal Q}{\{F\,\bar c=t\}\uplus E}}{\RLst{\mathcal Q\,\exists\bar G'}{\theta(\rho(E))}}
\quad
\begin{array}{l}
t \text{ rigid};\ \rho=\{\langle G,\lambda\bar d.\,G'\,(\bar d\cap\bar c)\rangle \mid (G\,\bar d) \text{ flexible in } t\} \text{ (pruning)};\\
t'=\rho(t);\ F\notin\mathrm{fv}(t');\ \text{every universal of } t' \text{ in the scope of } F \text{ is in } \bar c;\\
\theta=\{\langle F,\lambda\bar c.\,t'\rangle\}
\end{array}
\\
\inferrule*[right=\rulename{occurs}]{\RLst{\mathcal Q}{\{F\,\bar c=t\}\uplus E}}{\bot}
\quad (t \text{ rigid},\ F\in\mathrm{fv}(\rho(t)))
\and
\inferrule*[right=\rulename{scope}]{\RLst{\mathcal Q}{\{F\,\bar c=t\}\uplus E}}{\bot}
\quad (t \text{ rigid},\ \rho(t) \text{ has a universal in the scope of } F \text{ not in } \bar c)
\\
\inferrule*[right=\rulename{delay}]{\RLst{\mathcal Q}{\{s=t\}\uplus E}}{\RLst{\mathcal Q}{E}\ \text{ with } s=t \text{ deferred}}
\\
\text{(a flexible side of $s=t$ is not a pattern; the equation is resumed when substitutions change its status)}
\end{mathpar}
\caption{Pattern unification on $\forall\exists\forall$ problems \book{212,214,221--223}; \rulename{delay} is the dynamic-$\Llam$ extension \book{225--226}. $\mathcal Q\,\exists H$ means a new existential in the $\exists$ block; the bindings $\theta$ (and $\rho$) are composed into the answer substitution. Rule names are ours.}\label{fig:rules-pattern}
\end{figure}


\begin{example}[\rulename{elim}, \rulename{occurs}, \rulename{scope} \book{223--224}]\label{ex:rules-elim}
In $\forall f\,\forall g\,\exists U\,\exists V\,\forall w\,\forall x\,\forall y\,[(f\,(U\,x\,y))=(f\,(g\,(V\,y\,w)))]$, \rulename{rr-dec} leaves $(U\,x\,y)=(g\,(V\,y\,w))$; \rulename{elim} prunes $V:=\abs{y}{\abs{w}{V'\,y}}$ (since $w\notin\{x,y\}$) and binds $U:=\abs{x}{\abs{y}{g\,(V'\,y)}}$, a most general unifier.
With right side $(g\,(U\,y\,w))$ instead, \rulename{occurs} fails; with $(g\,w)$, \rulename{scope} fails, because $w$ is not among the arguments of $U$ and is not inside a flexible subterm.
\end{example}

\begin{example}[\rulename{elim} with pruning \book{224--225}]
In $\forall f\,\exists U\,\forall x\,\forall y\,\forall z\,\exists V\,[(U\,x)=(f\,V)]$, raising $V$ over $x,y,z$ gives $(U\,x)=(f\,(V'\,x\,y\,z))$; \rulename{elim} prunes $V':=\lambda x\lambda y\lambda z\,(W\,x)$ and binds $U:=\lambda x\,(f\,(W\,x))$.
\end{example}

\paragraph{Implementations of the same rules.}
Nipkow's functional version works on de~Bruijn terms with a $\forall\exists\forall$ prefix \cite{Nipkow1993}; Dowek, Hardin, Kirchner and Pfenning recast the rules (Invert, Pruning, Same-variable) in the λσ-calculus of explicit substitutions, where raising is the rule Anti-Exp-λ \cite[p.\,8--9]{DowekHardinKirchnerPfenning1996}; Nadathur and Linnell drop the prefix in favour of level tags \cite{NadathurLinnell2005} (\cref{sec:rules-levels}).

% =====================================================================================
\section{Universe levels}\label{sec:rules-levels}

Tagging every eigenvariable and logic variable with a number replaces the explicit prefix: the tag records which constants a variable may see \cite[p.\,15--18]{NadathurJayaramanKwon1995}.

\begin{center}
\small
\begin{mathpar}
\inferrule*[right=(1)]{\langle G_1\wedge G_2,\Prog,I\rangle}{\langle G_1,\Prog,I\rangle,\ \langle G_2,\Prog,I\rangle}
\and
\inferrule*[right={(2)\ \normalfont $(i=1,2)$}]{\langle G_1\vee G_2,\Prog,I\rangle}{\langle G_i,\Prog,I\rangle}
\and
\inferrule*[right={(3)\ \normalfont ($w$ new, tag $I$)}]{\langle \exists x\,G,\Prog,I\rangle}{\langle \subst{G}{w}{x},\Prog,I\rangle}
\\
\inferrule*[right=(4)]{\langle (D_1\wedge\dots\wedge D_n)\imp G,\Prog,I\rangle}{\langle G,\Prog\cup\{D_1,\dots,D_n\},I\rangle}
\and
\inferrule*[right={(5)\ \normalfont ($c$ new, tag $I+1$)}]{\langle \forall x\,G,\Prog,I\rangle}{\langle \subst{G}{c}{x},\Prog,I+1\rangle}
\\
\inferrule*[right={(6)\ \normalfont (instance $G\imp A'$)}]{\langle A,\Prog,I\rangle}{\sigma(\langle G,\Prog,I\rangle)}
\and
\inferrule*[right={(7)\ \normalfont (instance $A'$)}]{\langle A,\Prog,I\rangle}{\text{solved}}
\end{mathpar}
\captionof{figure}{The abstract interpreter on decorated goals $\langle G,\Prog,I\rangle$ \cite[p.\,17--18]{NadathurJayaramanKwon1995}. In (6) and (7) the clause instance is a new tagged instance relative to $I$ of a clause of $\Prog$, and $\sigma=\mathrm{mgu}(A,A')$ is applied to all remaining goals. A new tagged instance relative to $I$ gives tag 1 to untagged constants and a new variable of tag $I$ to each $\forall x$ of the clause; the original goal starts at $I=1$ with its constants and free variables tagged 1.}\label{fig:rules-ui}
\end{center}

\begin{definition}[{binding condition and lowering \cite[p.\,16]{NadathurJayaramanKwon1995}}]\label{def:rules-binding}
A logic variable $X$ of tag $i$ may be bound to $t$ only if (a) $X$ does not occur in $t$ (occurs check) and (b) no constant in $t$ has a tag greater than $i$.
Before the binding is made, every variable in $t$ with a tag greater than $i$ has its tag \emph{lowered} to $i$, so that its own later instances respect the restriction on $X$.
\end{definition}

\begin{proposition}[{\cite[Prop.\,4.1]{NadathurJayaramanKwon1995}}]\label{prop:rules-ui}
Derivations of the interpreter are sound and complete with respect to intuitionistic provability, and the choice of the next decorated goal to reduce is immaterial.
\end{proposition}

\begin{example}[scope failure \book{85,106}]\label{ex:rules-scope}
\lp{?- sigma x\ pi y\ x = y.} fails: the query gives $X^1$, the $\forall$ gives $c^2$, and $X^1:=c^2$ violates (b).
With the program $\forall z\,(q(z)\imp p(d(z)))$, the goal $\exists x\,\forall y\,(q(y)\imp p(x))$ becomes $q(c^2)\imp p(X^1)$, at $I=2$; backchaining creates $Z^2$, unifying $p(d(Z^2))$ with $p(X^1)$ binds $X^1:=d(Z^2)$ and \emph{lowers} $Z$ to tag 1, so the subgoal $q(Z^1)$ fails against the hypothesis $q(c^2)$; without lowering, $Z^2:=c^2$ would succeed and smuggle $c$ into $X$ \cite[p.\,16]{NadathurJayaramanKwon1995}.
\end{example}

\paragraph{Levels for higher-order patterns \cite[p.\,8--9]{NadathurLinnell2005}.}
Let $l(y)$ count the times an $\exists$ is immediately followed by a $\forall$ in the prefix up to and including the quantifier of $y$; then a universal $u$ is in the scope of an existential $x$ iff $l(x)<l(u)$, which makes the pattern test of \cref{def:rules-pattern} local.
Raising of $g$ is restricted to $\bar y{\Uparrow}g=\{u\in\bar y\mid l(u)\le l(g)\}$.
New existentials created by imitation, projection, pruning or flexible-flexible steps are quantified \enquote{at the same place} as the variable they replace \book{221--222,225}, i.e.\ they take its level, not a new one.
The book describes the same device as lists of universals \enquote{over which it can be raised}, initial segments of the prefix, represented by numerical indices \book{225}.

\begin{remark}[higher-order binding condition; reconstruction in our notes, not in the book]\label{rem:rules-holevel}
In \rulename{elim} for $X\,\bar c=t$ the tag test (b) applies to universals in rigid positions of $t$ not among $\bar c$ (\rulename{scope} in \cref{fig:rules-pattern}); inside the arguments of a flexible $(Y\,\bar d)$ such universals are pruned instead; and lowering a $Y$ of higher level binds $Y$ to $(Y'\,\bar a)$ with $Y'$ at the level of $X$, which is delayed raising \book{225}.
\end{remark}

\begin{kanren}[what levels require]
A miniKanren with \rulename{generic} goals must record a level for every logic variable and eigenvariable and must enforce the binding condition, with lowering, inside its unifier; see \cref{ch:design}.
\end{kanren}

% =====================================================================================
\section{Summary}\label{sec:rules-summary}

\begin{itemize}[nosep]
\item Syntax: \cref{tab:rules-syntax}; grammars of \fohc, \fohh, \hohc, \hohh, $\hohh^+$ and $\Llam$: \cref{tab:rules-grammars,def:rules-llambda}.
\item Proof search: right rules \cref{fig:rules-right}, backchaining \cref{fig:rules-left}, completeness \cref{tab:rules-completeness}.
\item Terms: typing \cref{fig:rules-typing}, conversion \cref{tab:rules-conv}, normal forms \cref{def:rules-nf,thm:rules-norm}.
\item Prefixes: unifier vs solution \cref{def:rules-solution}, prefix order \cref{ex:rules-prefix}, raising \cref{prop:rules-raising}, problem normal form \cref{sec:rules-raising}.
\item Unification: first-order \cref{fig:rules-fo}, Huet \cref{fig:rules-huet,ex:rules-huet}, patterns \cref{fig:rules-pattern,ex:rules-elim}.
\item Scope by numbers: \cref{fig:rules-ui,def:rules-binding}; lowering is what makes the second half of \cref{ex:rules-scope} fail.
\end{itemize}

\chapter{Glossary}\label{app:glossary}

The preceding chapters use vocabulary from three communities. The first is the proof theory of Miller and
Nadathur's book~\cite{MillerNadathur2012}, where you find sequents, uniform proofs and eigenvariables. The second
is the literature on unification of $\lambda$-terms (patterns, raising, pruning, universe levels). The third is
miniKanren, with streams, interleaving, soft cut and reification, and in particular the engine \micro{} that we
want to extend. The same word sometimes means different things in different communities (\enquote{suspension},
\enquote{level}, \enquote{scope extrusion}, \enquote{fresh}), and the same idea sometimes has three names
(GENERIC, \lp{pi}, \texttt{all}). This appendix collects all of them in one alphabetical list. Each entry gives a
short, precise definition and tells you where the concept is treated in this document and in the book.

\section{How to read this glossary}\label{sec:glossary-howto}

\paragraph{Pointers.} Most entries end with a pointer of the form
\enquote{$\rightarrow$ \cref{ch:hh}; \book{82}}. The first part is the chapter of this document where the
concept is developed. The second part is a \emph{printed} page of the book~\cite{MillerNadathur2012}; most of
these come from the book's own index. Entries that come only from the miniKanren literature or from the
\aux{} sources have no book page; they cite the paper instead. Names of \micro{} procedures and macros are set in
\texttt{typewriter} type. We describe only what exists today. The new operators are the subject of
\cref{ch:design,ch:module}.

\paragraph{Notation.} \Cref{tab:glossary-notation} lists the symbols that occur most often in the entries.

\begin{table}[ht]
\centering\small
\caption{Notation used throughout the glossary.}\label{tab:glossary-notation}
\begin{tabularx}{\textwidth}{@{}lX@{}}
\toprule
Symbol & Meaning \\
\midrule
$\seq{\Sig}{\Prog}{G}$ & sequent: signature $\Sig$, program $\Prog$, goal $G$ \book{39} \\
$\Sig;\Prog \xrightarrow{D} A$ & focused sequent: the clause $D$ is selected to prove the atom $A$ \book{44} \\
$G$, $D$ & goal formula, program clause (definite formula) \\
$\top,\wedge,\vee,\imp,\forall,\exists$ & logical constants; concrete syntax \lp{true}, \lp{&} or \lp{,}, \lp{;}, \lp{=>} or \lp{:-}, \lp{pi}, \lp{sigma} \\
$\abs{x}{t}$ & abstraction; λProlog \lp{x\ t} \\
$\subst{s}{t}{x}$ & capture-avoiding substitution of $t$ for $x$ in $s$ \book{85} \\
$=_{\beta\eta}$, $\lambda$-conversion & equality of terms modulo $\alpha$, $\beta$, $\eta$ \book{100} \\
$\tau\to\sigma$, $\typ{o}$ & function type; type of formulas \\
$Q_1x_1\dots Q_nx_n[\,\bigwedge_j t_j=s_j\,]$ & unification problem under a mixed prefix \book{105} \\
$a\mathbin{\#}t$ & nominal freshness: the atom $a$ does not occur free in $t$ \\
$\langle\phi;t\rangle$ & static closure of Ma and Friedman: term $t$ under the scope $\phi$ \\
\texttt{°}, \texttt{ª}, \texttt{§}, \texttt{δ} & \aux{} suffixes and glyphs: goal or relation, applicative form, stream, delay \\
\bottomrule
\end{tabularx}
\end{table}

\begin{pitfall}[one word, several meanings]
Several terms below are overloaded across the sources. Keep the distinctions in mind when you read the papers.
\begin{itemize}[nosep]
\item \textbf{Suspension.} In \micro{}, a suspension is a \texttt{δ}-promise, the immature part of a stream. In
  nominal unification, $\pi\cdot X$ is a permutation suspended on a variable. In Teyjus, the suspension
  $[[t,ol,nl,e]]$ is an explicit-substitution term. In ELPI, a suspended goal is a syntactic constraint.
\item \textbf{Level.} A de Bruijn level numbers a bound variable from the root inwards. A universe level, or
  universe index, is the $\forall$-nesting depth at which a variable or constant was created. ELPI makes the two
  coincide, but they are different ideas.
\item \textbf{Scope extrusion.} The book uses this name for a classical equivalence that moves an implication over
  a disjunction \book{89}. In the $\pi$-calculus it means sending a restricted name out of its scope. In
  implementation discussions it means an eigenvariable escaping into an older variable.
\item \textbf{Fresh.} In miniKanren, \texttt{fresh} is the existential quantifier and introduces logic variables.
  In αKanren, \texttt{fresh} introduces \emph{noms}, which are names rather than variables.
\item \textbf{Cut.} The cut rule of the sequent calculus has no connection to Prolog's pruning operator \lp{!}
  \book{63}.
\item \textbf{Substitution and grafting.} Instantiating a unification variable is capture-avoiding in standard
  higher-order unification. It is capture-permitting (grafting) in the Ma--Friedman and DHKP formulations.
\end{itemize}
\end{pitfall}

\section{Terms from A to Z}\label{sec:glossary-az}

\newlist{glossarylist}{description}{1}
\setlist[glossarylist]{font=\normalfont\bfseries,leftmargin=1.2em,labelindent=0pt,itemsep=0.45em,parsep=0.2em}
\begingroup\sloppy

\subsection*{A}
\begin{glossarylist}
\item[Abstract datatype] A module that exports a type constructor and the predicates over it but hides the data
  constructors, so other code can use only the exported operations (book Ch.~6). $\rightarrow$ \book{165}.
\item[Abstract logic programming language] A logic in which searching for uniform proofs is a complete proof
  procedure. This is the book's criterion for a logic to count as a logic programming language. fohc, fohh, hohc
  and hohh all qualify~\cite{MillerNadathurPfenningScedrov1991}. $\rightarrow$ \cref{ch:horn}; \book{73}.
\item[Abstraction] The term $\abs{x}{t}$ (λProlog \lp{x\ t}), the only binder of the term language. $x$ is the
  bound variable and $t$ the body. The quantifiers are the constants \lp{pi} and \lp{sigma} applied to
  abstractions. $\rightarrow$ \cref{ch:lambda}; \book{97}.
\item[Accumulation (module, signature)] \texttt{accumulate} inlines the code of other modules after renaming
  their private names. \texttt{accum\_sig} inserts the declarations of other signatures, either to re-export them
  or to declare a parameter that a later module supplies. $\rightarrow$ \book{154}, \book{155}.
\item[Administrative redex] A $\beta$-redex that a program transformation (here the CPS transformation) introduces
  itself, as opposed to one present in the source. The two-phase CPS of book Ch.~10 marks these redexes and
  reduces them afterwards. $\rightarrow$ \cref{ch:apps}; \book{257}.
\item[$\alpha$-conversion] Renaming bound variables. Terms that differ only by $\alpha$-conversion are identified,
  so the names of bound variables are \enquote{a fiction} and an implementation may not represent them at all.
  $\rightarrow$ \cref{ch:lambda,ch:lambdatree}; \book{100}, \book{183}, \book{205}.
\item[αKanren] Byrd and Friedman's nominal miniKanren~\cite{ByrdFriedman2007,Byrd2009}. \texttt{fresh}
  introduces noms, \texttt{(tie a t)} binds the nom $a$ in $t$, and \texttt{(hash a t)} asserts the freshness
  $a\mathbin{\#}t$. Unification is nominal unification. The implementation does not support disequality
  constraints. $\rightarrow$ \cref{ch:lkanren}.
\item[AND] The goal-reduction rule for $G_1\wedge G_2$: prove both conjuncts with the same signature and program.
  $\rightarrow$ \cref{ch:horn}; \book{39}.
\item[Anonymous variable] \lp{_}: each occurrence is a distinct variable, and its binding is not reported.
  $\rightarrow$ \cref{ch:horn}; \book{47}.
\item[Answer substitution] The bindings that a successful search reports for the free variables of a query. These
  variables are implicitly existential. The miniKanren counterpart is the reified answer.
  $\rightarrow$ \cref{ch:horn}; \book{43}.
\item[Applicative form] An \aux{} naming convention. The suffix \texttt{ª} marks the procedure behind a goal macro:
  \texttt{freshª} behind \texttt{fresh°}, and \texttt{andª}, \texttt{orª} behind \texttt{and°}, \texttt{or°}.
  $\rightarrow$ \cref{ch:micro}.
\item[\texttt{assume}, \texttt{all}] The λKanren goals for $D\imp G$ and $\forall x\,G$. See \emph{AUGMENT},
  \emph{GENERIC} and \emph{λKanren}. $\rightarrow$ \cref{ch:lkanren}.
\item[Atom (nominal), nom, name] In nominal logic and αKanren, a constant-like name that stands for an
  object-level variable. A nom unifies only with itself or with an unbound logic variable.
  $\rightarrow$ \cref{ch:lkanren}.
\item[Atomic formula] A term of type $\typ{o}$ whose head is a predicate symbol, that is, a constant with target
  type $\typ{o}$. $\rightarrow$ \cref{ch:horn}; \book{34}.
\item[AUGMENT] The goal-reduction rule for $D\imp G$: add $D$ to the program and prove $G$. This rule is the basis
  of hypothetical reasoning; the clause is visible only during the proof of $G$. λProlog writes \lp{D => G};
  λKanren calls it \texttt{assume}.
  \[\inferrule{\seq{\Sig}{\Prog\cup\{D\}}{G}}{\seq{\Sig}{\Prog}{D\imp G}}\;\rulename{augment}\]
  $\rightarrow$ \cref{ch:hh}; \book{39}, \book{77}.
\item[Available level] (Ma--Friedman) A de Bruijn level of a scope $\phi$ whose binder name is not rebound further
  in. A shadowed binder's level is unavailable, so $\lambda a.\lambda a.X \stackrel{?}{=} \lambda c.\lambda d.c$
  fails~\cite{MaFriedman2021}. $\rightarrow$ \cref{ch:lkanren}.
\end{glossarylist}

\subsection*{B}
\begin{glossarylist}
\item[Backchaining] Proving an atomic goal $A$ by selecting a program clause, possibly an assumed one, whose head
  matches $A$, and then proving the clause's body. It is the only phase of search that consults the program.
  The book formalises it with focused sequents and the rules \rulename{decide}, \rulename{initial}, $\imp$L,
  $\wedge$L and $\forall$L. $\rightarrow$ \cref{ch:horn}; \book{3}, \book{44}.
\item[$\beta$-conversion, $\beta$-redex, $\beta$-reduction] A $\beta$-redex is $(\abs{x}{s})\,t$, and it reduces
  to $\subst{s}{t}{x}$. $\beta$-reduction is the computation of functional programming, and it is part of the
  equality of λProlog terms. $\rightarrow$ \cref{ch:lambda}; \book{100}--\book{101}.
\item[$\beta_0$-conversion] The special case $(\abs{x}{s})\,x = s$ of $\beta$. With $\alpha$, it is equivalent to
  $(\abs{y}{s})\,x=\subst{s}{x}{y}$ when $x$ is not free in $\abs{y}{s}$, so it only renames. It is the only $\beta$
  needed in $\Llam$. $\rightarrow$ \cref{ch:lambdatree,ch:hounif}; \book{206}, \book{220}.
\item[$\beta$-normal form] A term without $\beta$-redexes. It has the shape
  $\lambda x_1\dots\lambda x_n.\,(h\,t_1\dots t_m)$, with \emph{binder} $x_1\dots x_n$, \emph{head} $h$ and
  \emph{arguments} $t_1\dots t_m$. $\rightarrow$ \cref{ch:lambda}; \book{101}, \book{111}.
\item[$\beta\eta$-long normal form] A $\beta$-normal form whose arguments are themselves in long form and whose
  body has primitive type. Every $\beta$-normal form can be $\eta$-expanded into an equal long one; for
  $f:(i\to i)\to i$, the long form of $f$ is $\lambda w.\,f\,(\lambda u.\,w\,u)$. Higher-order unification works on
  these forms.
  $\rightarrow$ \cref{ch:lambda}; \book{111}.
\item[Big-step / small-step semantics] A big-step (natural) semantics relates a term to its final value. A
  small-step (structural operational) semantics relates a term to the result of one atomic step.
  $\rightarrow$ \cref{ch:apps}; \book{250}--\book{251}.
\item[Binary clause] A clause whose body is at most one, possibly complex, atom. Every clause can be put in this
  form by threading a continuation goal $K:\typ{o}$ (continuation-passing style).
  $\rightarrow$ \cref{ch:hoprog}; \book{130}.
\item[Binder] A syntactic construct that introduces a locally scoped name: a quantifier, $\lambda$, $\int$,
  $\Sigma$ or a local declaration. In λ-tree syntax, object-level binders are encoded by meta-level $\lambda$.
  $\rightarrow$ \cref{ch:lambdatree}.
\item[Binder mobility] The way a binder moves across levels during computation. A term-level $\lambda$ becomes a
  formula-level $\forall$ (\lp{pi}), which in turn becomes a proof-level eigenvariable. This movement lets
  λProlog programs descend under binders without managing names. $\rightarrow$ \cref{ch:lambdatree};
  \book{184}, \book{205}.
\item[Binder-merging] (Ma--Friedman) The implementation technique of renaming one side's binders to the other
  side's before comparing the bodies structurally. With name shadowing and a naive reading of levels back into
  names, it can make Church \emph{false} $\lambda a.\lambda a.a$ unify with Church \emph{true}
  $\lambda c.\lambda d.c$. $\rightarrow$ \cref{ch:lkanren}.
\item[Bisimulation, simulation] Coinductive behavioural equivalence and preorder on processes. Book Ch.~11 shows
  that a simulation written with \lp{pi} is incorrect, because \lp{pi} is intensional.
  $\rightarrow$ \cref{ch:apps}; \book{271}.
\end{glossarylist}

\subsection*{C}
\begin{glossarylist}
\item[Canonical substitution] For a problem whose equations are all flex-flex: fix one new existential
  variable $H^{\sigma}$ for each primitive type $\sigma$, and map every existential variable of type
  $\tau_1\to\dots\to\tau_n\to\sigma$ to $\lambda x_1\dots\lambda x_n.\,H^{\sigma}$. This substitution always unifies
  such a problem, which is why a flex-flex remainder counts as success. $\rightarrow$ \cref{ch:hounif};
  \book{217}--\book{218}.
\item[Capture-avoiding substitution] $\subst{s}{t}{x}$, carried out after renaming bound variables of $s$ so that
  no free variable of $t$ is captured. It is defined only up to $\alpha$. Contrast grafting.
  $\rightarrow$ \cref{ch:hh,ch:lambda}; \book{85}--\book{86}, \book{100}.
\item[Church numeral] $\ulcorner n\urcorner=\lambda f.\lambda x.\,f^n x$. The book uses Church numerals to show
  $\beta$-reduction as computation and to encode Diophantine equations as unification problems.
  $\rightarrow$ \cref{ch:lambda}; \book{102}, \book{115}.
\item[Clausal order] The nesting depth of implications to the left in a formula, analogous to the order of a type.
  $\rightarrow$ \cref{ch:hh}; \book{76}.
\item[Closed-world / open-world assumption] Whether the world, that is, the signature--program pair, can grow. Under
  the open-world assumption that λProlog adopts, generic and universal judgments coincide.
  $\rightarrow$ \cref{ch:hh}; \book{94}.
\item[Complete set of unifiers (CSU)] A set of unifiers such that every unifier is an instance of some member.
  Higher-order problems have CSUs, possibly infinite, but in general no mgu. One might ask instead for a \emph{non-redundant}
  CSU, in which no two members have a common instance, but the book notes that this cannot be guaranteed in general
  either. $\rightarrow$ \cref{ch:hounif}; \book{213}.
\item[Completeness (of a search strategy)] In the book, completeness is a property of a proof system: uniform proofs
  exist whenever proofs do. A search \emph{strategy} is a different matter. Depth-first search is incomplete.
  miniKanren's interleaving search is complete for answers, meaning every answer with a finite derivation is
  eventually produced~\cite{RozplokhasVyatkinBoulytchev2020}. It does not guarantee termination or finite
  failure. $\rightarrow$ \cref{ch:horn,ch:micro}.
\item[Computation-as-deduction] The view that computation is the construction or transformation of proofs. It
  comes in two forms: \emph{proof normalisation} (functional programming) and \emph{proof search} (logic
  programming). \emph{Computation-as-model} instead uses logic to reason about computations.
  $\rightarrow$ \cref{ch:intro}; \book{2}.
\item[\texttt{cond°}] \micro{}'s \texttt{conde}: an interleaving disjunction of conjunctions of goals. There is no
  \texttt{else} keyword. $\rightarrow$ \cref{ch:micro}.
\item[Consistent / inconsistent program] A program is inconsistent if every formula is provable from it. Every fohc
  and fohh program is consistent. Clause heads must be rigid: with the two clauses \lp{P 5 :- q.} and \lp{q.},
  whose first head is flexible, any formula becomes provable. $\rightarrow$ \cref{ch:hoprog}; \book{119}.
\item[Constant clash] The failure of first-order unification on an equation whose two sides have distinct head
  constants. $\rightarrow$ \cref{ch:intro}; \book{30}.
\item[Context (Ma--Friedman)] The set of postponed equations between static closures of unification variables. A
  solution in their algorithm is a pair (grafting, context), much like Huet's postponed flex-flex pairs.
  $\rightarrow$ \cref{ch:lkanren}.
\item[Constraint store (\micro{})] Besides the substitution \texttt{S}, a state carries three stores:
  \texttt{D} for disequalities (\texttt{≠°}), \texttt{A} for type tags (\texttt{symbol°}, \texttt{number°}) and
  \texttt{T} for absento constraints (\texttt{absent°}). After a unification, the constraints touched by its new
  bindings are checked again. $\rightarrow$ \cref{ch:micro}.
\item[Constructive negation] (Moiseenko) Negation for miniKanren that collects all the answers of $g$ and negates
  their disjunction. The states then contain universally quantified disequalities, which need a dedicated solver.
  It requires stratification, diverges on goals with infinitely many answers, and defines $\forall$ as
  $\neg\exists\neg$~\cite{Moiseenko2019}. $\rightarrow$ \cref{ch:lkanren}.
\item[Contraction] The structural rule that duplicates a hypothesis into a premise. It is the source of
  non-termination in naive provers for propositional intuitionistic logic. Contraction-free calculi such as
  Dyckhoff's terminate. $\rightarrow$ \cref{ch:apps}; \book{245}.
\item[Copy clauses] Clauses \lp{copy} that specify structural equality over an object signature, one clause per
  constructor. Adding the hypothesis \lp{copy c T} for an eigenvariable \lp{c} turns copying into object-level
  substitution, entirely inside $\Llam$. $\rightarrow$ \cref{ch:lambdatree}; \book{199}.
\item[CPS (continuation-passing style)] In book Ch.~5, the threading of a continuation goal that yields binary
  clauses. In book Ch.~10, Fischer's CPS transformation of functional programs.
  $\rightarrow$ \cref{ch:hoprog,ch:apps}; \book{131}, \book{256}.
\item[Cut rule, cut elimination] The cut rule is the lemma rule of the sequent calculus. Gentzen's theorem shows it
  is admissible~\cite{Gentzen1935}. Proof search looks only for cut-free proofs. The cut rule is unrelated to
  Prolog's \lp{!}. $\rightarrow$ \cref{ch:horn}; \book{63}, \book{72}.
\item[Cut, Prolog's \texttt{!}] The pruning operator. λProlog defines \lp{if} and \lp{not} with it, and it affects
  completeness. \micro{} has no \texttt{!}. Its closest tools are the soft cut \texttt{if°} and
  \texttt{(take° 1 g)}. $\rightarrow$ \cref{ch:hoprog,ch:micro}; \book{137}.
\end{glossarylist}

\subsection*{D}
\begin{glossarylist}
\item[D-formula, G-formula] Program clauses and goals. Each language defines them by a pair of grammars. For fohc:
  $G ::= \top\mid A\mid G\wedge G\mid G\vee G\mid \exists x\,G$ and $D ::= A\mid G\imp D\mid D\wedge D\mid\forall x\,D$.
  fohh adds $D\imp G$ and $\forall x\,G$ to the goals. $\rightarrow$ \cref{ch:horn,ch:hh}; \book{43}, \book{76}.
\item[de Bruijn index, de Bruijn level] Nameless representations of bound variables~\cite{deBruijn1972}. An
  \emph{index} counts binders outwards from the occurrence (1-based in the book). A \emph{level} counts binders
  inwards from the root. In $\lambda c.\lambda d.c$, $c$ has level 0 and index 1 when counting from 0. Teyjus uses
  indices in terms and levels in the context. ELPI uses levels throughout.
  $\rightarrow$ \cref{ch:lambdatree,ch:lkanren}; \book{192}.
\item[Decide, initial] The two rules that start and end backchaining. \rulename{decide} picks a clause $D\in\Prog$
  to focus on. \rulename{initial} closes a focused atom whose head matches the goal.
  $\rightarrow$ \cref{ch:horn}; \book{45}.
\item[\texttt{define-relation}] The \micro{} macro that defines a relation. It is inverse-$\eta$-delayed: the body
  is built only when a state arrives, and then it suspends. $\rightarrow$ \cref{ch:micro}.
\item[Definite formula] A program formula that does not contain \enquote{indefinite information of the kind let in
  by disjunctions}. Disjunction (and, by the grammar, top-level $\exists$) is banned from programs.
  $\rightarrow$ \cref{ch:horn}; \book{42}.
\item[Delay, \texttt{δ}] In \micro{}, \texttt{δ} is CHICKEN's memoising \texttt{delay}. A promise is an immature
  stream and is the unit of interleaving. $\rightarrow$ \cref{ch:micro}.
\item[Delayed (postponed) equations] Equations outside the decidable fragment, or flex-flex pairs, kept as
  constraints until further instantiation makes them tractable. Teyjus, DHKP~\cite{DowekHardinKirchnerPfenning1996}
  and the Ma--Friedman \emph{context} postpone them. ELPI raises an error by default. The book's version is
  dynamic $\Llam$. $\rightarrow$ \cref{ch:hounif,ch:lkanren}; \book{225}--\book{226}.
\item[Delayed raising] Instead of raising eagerly, keep with each existential variable the list of universals it
  may be raised over, and raise only when a binding needs it. These lists are initial segments of the prefix, so
  numerical indices represent them~\cite{NadathurLinnell2005}. $\rightarrow$ \cref{ch:hounif}; \book{225}.
\item[Determinate type] A type whose type variables are all transparent, that is, occur in the target type.
  $\rightarrow$ \cref{ch:horn}; \book{67}.
\item[Difference list, functional difference list] \lp{(dl L K)} denotes $L$ minus its suffix $K$.
  \lp{(fdl F)}, with \lp{F : list A -> list A}, denotes \lp{F nil}, and appending is function composition.
  $\rightarrow$ \cref{ch:hoprog}; \book{140}.
\item[Disagreement set] In nominal unification, the set of atoms on which two permutations differ. In Huet-style
  presentations, the multiset of pairs still to be unified. $\rightarrow$ \cref{ch:lkanren}.
\item[Disequality constraint] \texttt{≠°} in \micro{}, stored in the store \texttt{D} as the list of bindings that
  would make the two terms equal. Moiseenko's constructive negation needs universally quantified disequalities
  instead~\cite{Moiseenko2019}. $\rightarrow$ \cref{ch:micro,ch:lkanren}.
\item[Disjunction / existence property] In intuitionistic logic, $\vdash F_1\vee F_2$ implies $\vdash F_i$ for
  some $i$, and $\vdash\exists x\,F$ yields a witness. Harrop formulas guarantee these properties, which is why
  answer substitutions make sense. $\rightarrow$ \cref{ch:hh}; \book{93}.
\item[Dynamic $\Llam$] The practical regime in which $\Llam$ is not imposed statically on clauses. Instead, the
  unification problems met at run time are expected to fall inside it, as in \lp{subst R N (R N)} with \lp{R}
  and \lp{N} already instantiated. $\rightarrow$ \cref{ch:hounif}; \book{225}--\book{226}.
\end{glossarylist}

\subsection*{E}
\begin{glossarylist}
\item[E-formula] $E ::= D\mid E\wedge E\mid\exists x\,E$: a program formula with existentially quantified, local
  names. It gives modules with hidden constants a logical reading. $\rightarrow$ \book{160}--\book{162}.
\item[Eigenvariable, scoped constant] The new constant that GENERIC ($\forall$R) introduces to prove
  $\forall x\,G$. It is not in the current signature and is never instantiated. The \emph{eigenvariable
  condition} says that no logic variable created before it may be bound to a term containing it. The book's
  example is the sterile jar: to prove \lp{pi x\ bug x => in x j => dead x}, pick a new bug constant and reason
  about it alone. $\rightarrow$ \cref{ch:hh,ch:lambda}; \book{82}--\book{83}.
\item[\texttt{eigen}] (Hemann, Byrd, Friedman) The operator of miniKanren-with-symbolic-constraints that creates
  eigen objects and installs an \enquote{eigen-absento} constraint over the list of all variables created so far.
  It is supported only together with \texttt{==}~\cite{HemannByrdFriedman2013}. $\rightarrow$ \cref{ch:lkanren}.
\item[Elementary substitution] Huet's partial bindings for a flex-rigid pair: imitation and projection. In
  Ma--Friedman's reading, each unification variable takes every visible bound variable as an argument.
  $\rightarrow$ \cref{ch:hounif}.
\item[ELPI] An embeddable λProlog interpreter~\cite{DunchevEtAl2015}. It uses de Bruijn levels and annotates
  unification variables with levels $X^j$. It reports an error outside the pattern fragment by default, and it
  offers syntactic constraints. $\rightarrow$ \cref{ch:lkanren}.
\item[Empty type] A type with no closed inhabitant. Higher-order intuitionistic logic does not exclude empty types,
  and they matter for $\exists$ answers and for unification \emph{solutions}. $\rightarrow$ \cref{ch:lambda};
  \book{108}.
\item[\texttt{eq?}-invariant] (\micro{}) A unification that adds no binding returns the very same state object.
  The disequality and entailment checks rely on this to recognise that nothing new was learned.
  $\rightarrow$ \cref{ch:micro}.
\item[Essentially universal / essentially existential] A classification of bound-variable occurrences. In a
  goal, an occurrence is essentially universal if it is bound by a positive $\forall$, a negative $\exists$, or a
  term-level $\lambda$. Otherwise it is essentially existential. Clauses use the dual classification. During search,
  essentially existential variables become logic variables and essentially universal ones become eigenvariables.
  $\rightarrow$ \cref{ch:lambdatree}; \book{205}.
\item[$\eta$-conversion] $\abs{x}{(t\,x)} = t$ when $x$ is not free in $t$. $\rightarrow$ \cref{ch:lambda};
  \book{101}.
\item[Evaluation context] A term with a hole that marks where the next redex is. In book Ch.~10 it is a meta-level
  function \lp{tm -> tm}. $\rightarrow$ \cref{ch:apps}; \book{253}.
\item[Ex falso quodlibet] From $\bot$, infer anything. Intuitionistic logic is minimal logic plus this rule.
  $\rightarrow$ \cref{ch:hh}; \book{90}.
\item[Excluded middle] $B\vee\neg B$: valid classically, rejected intuitionistically. It is one of the
  counterexamples that motivate banning disjunction from programs. $\rightarrow$ \cref{ch:horn}; \book{42}.
\item[Existential generalisation] From $M(t)$ infer $\exists x\,M(x)$. It explains the multiple answers of
  \lp{subst} when run backwards. $\rightarrow$ \cref{ch:lambdatree}; \book{203}.
\item[Explicit substitution] Calculi that make substitution a term constructor, so it can be delayed and composed.
  Examples are $\lambda\sigma$ (DHKP) and Teyjus's suspension notation. $\rightarrow$ \cref{ch:lkanren};
  \book{209}.
\end{glossarylist}

\subsection*{F}
\begin{glossarylist}
\item[Flexible term, rigid term] A $\beta\eta$-normal term is \emph{flexible} if its head is an existential
  (unification) variable. It is \emph{rigid} if its head is a constant or a universally bound or $\lambda$-bound
  variable. The names are suggestive: Ma and Friedman, following Huet, explain that a flexible term is one that
  \enquote{can be removed through a substitution (or a grafting)}, while a substitution never changes the head of
  a rigid term~\cite{MaFriedman2021}. $\rightarrow$ \cref{ch:lambda,ch:hounif}; \book{111}, \book{214}.
\item[Flexible atom, rigid atom, flexible goal] An atom is flexible if its head is a variable and rigid if its head
  is a non-logical constant. Clause heads must be rigid. A flexible goal whose head is still uninstantiated when it
  is called is underconstrained. $\rightarrow$ \cref{ch:hoprog}; \book{119}, \book{132}.
\item[Flex-flex, flex-rigid, rigid-rigid] The classification of equations by the heads of their two sides.
  Rigid-rigid pairs simplify by decomposition or fail by a clash. Flex-rigid pairs call for imitation or
  projection. Flex-flex pairs always have unifiers, though not necessarily solutions, and are postponed.
  $\rightarrow$ \cref{ch:hounif}; \book{113}, \book{217}.
\item[Focused proof] Andreoli's generalisation of the two-phase structure of goal reduction followed by
  backchaining. Focusing groups small rules into large ones~\cite{Andreoli1992}. $\rightarrow$ \cref{ch:horn};
  \book{73}.
\item[fohc, fohh, hohc, hohh, hohh$^{+}$] The book's languages: first-order Horn clauses \book{43}, first-order
  hereditary Harrop formulas \book{76}, higher-order Horn clauses \book{122}, higher-order hereditary Harrop
  formulas \book{124}, and hohh$^+$, which additionally allows atomic clauses headed by an essentially universal
  variable \book{125}. $\rightarrow$ \cref{ch:horn,ch:hh,ch:hoprog}.
\item[Free for] $t$ is free for $x$ in $B$ if substituting $t$ for $x$ in $B$ captures no free variable of $t$.
  $\rightarrow$ \cref{ch:hh}; \book{85}.
\item[Fresh, \texttt{fresh°}] miniKanren's existential quantifier, which introduces logic variables. In \micro{},
  \texttt{fresh°} introduces them one at a time, and each allocation suspends once. For the αKanren meaning, see
  \emph{Atom (nominal)}. $\rightarrow$ \cref{ch:micro}.
\item[Freshness constraint] $a\mathbin{\#}t$: the atom $a$ does not occur free in $t$, also after later
  instantiation of the variables in $t$. It is written \texttt{hash} in αKanren~\cite{UrbanPittsGabbay2004}.
  $\rightarrow$ \cref{ch:lkanren}.
\end{glossarylist}

\subsection*{G}
\begin{glossarylist}
\item[Generic judgment, universal judgment] A generic judgment is parametric in a new object and is proved once,
  with an eigenvariable. A universal judgment holds for all terms. The two coincide under the open-world
  assumption. They correspond to the intensional and the extensional readings of $\forall$.
  $\rightarrow$ \cref{ch:hh}; \book{83}, \book{94}.
\item[GENERIC] The rule for $\forall x\,G$: add a new constant $c$ to the signature and prove $\subst{G}{c}{x}$.
  It is \lp{pi x\ G} in λProlog and \texttt{all} in λKanren.
  \[\inferrule{\seq{\Sig\cup\{c\}}{\Prog}{\subst{G}{c}{x}}}{\seq{\Sig}{\Prog}{\forall x\,G}}\;\rulename{generic}\quad(c\notin\Sig)\]
  $\rightarrow$ \cref{ch:hh}; \book{39}, \book{82}.
\item[Goal (miniKanren)] In \micro{}, a procedure from a state to a stream of states. \texttt{✓°} succeeds once
  and \texttt{✗°} fails. $\rightarrow$ \cref{ch:micro}.
\item[Goal-directed search, goal reduction] A non-atomic goal is decomposed by its top-level connective,
  independently of the program. This fixes the \enquote{search semantics} of the connectives. Only atomic goals
  consult the program. $\rightarrow$ \cref{ch:intro,ch:horn}; \book{3}, \book{39}.
\item[Grafting] First-order, capture-permitting replacement of a metavariable by a term, written $\{X\mapsto t\}$:
  grafting $X\mapsto a$ into $\lambda a.X$ gives $\lambda a.a$. DHKP \enquote{pre-cook} terms so that grafting
  respects scopes. Ma and Friedman combine grafting with de Bruijn levels~\cite{MaFriedman2021}.
  $\rightarrow$ \cref{ch:lkanren}.
\item[Grounding] In \micro{}, \texttt{°->list/ground} evaluates a reified answer, applying its $\lambda$ to its own
  variable names. Constraints are dropped. $\rightarrow$ \cref{ch:micro}.
\end{glossarylist}

\subsection*{H}
\begin{glossarylist}
\item[Harrop formula] $H ::= \top\mid A\mid H\wedge H\mid B\imp H\mid\forall x\,H$. These formulas guarantee the
  disjunction and existence properties at the root of a proof. $\rightarrow$ \cref{ch:hh}; \book{93}.
\item[Herbrand universe $H_1^{\Sig}$, $H_2^{\Sig}$] The terms allowed as quantifier instances. $H_1^{\Sig}$ (for
  hohc) contains the $\lambda$-normal terms over $\Sig$ without $\forall$ and $\imp$. $H_2^{\Sig}$ (for hohh)
  contains those without $\imp$. $\rightarrow$ \cref{ch:hoprog}; \book{122}.
\item[Hereditary Harrop formula] A D-formula of fohh or hohh. Equivalently, a formula with no positive occurrence
  of a disjunction or an existential. Unlike Horn clauses, these formulas let the signature and the program grow
  during search. $\rightarrow$ \cref{ch:hh}; \book{4}, \book{76}.
\item[Higher-order abstract syntax (HOAS)] The principle of encoding object-level binders by meta-level
  binders~\cite{PfenningElliott1988}. It is ambiguous between a functional reading, in which binders are host
  functions that cannot be analysed, and the logical reading, λ-tree syntax. $\rightarrow$ \cref{ch:lambdatree};
  \book{204}.
\item[Higher-order logic] The book distinguishes three senses: second-order model theory; proof theory with
  predicate quantification and comprehension; and Church's simple theory of types, with $\lambda$-terms and
  quantification at higher types. The book works in the third. $\rightarrow$ \cref{ch:intro}; \book{5}.
\item[Higher-order pattern] A term in which every occurrence of an existential variable $F$ is applied to distinct
  variables, either essentially universal or $\lambda$-bound, whose binders lie within the scope of $F$. For
  example, $\abs{x}{\abs{y}{F\,y\,x}}$ is a pattern, while $\abs{x}{F\,x\,x}$ (repeated argument) and $F\,c$ with $c$
  a signature constant (not a bound variable) are not.
  $\rightarrow$ \cref{ch:lambdatree,ch:hounif}; \book{205}, \book{220}.
\item[Higher-order pattern unification ($\Llam$-unification)] Unification restricted to patterns. It is decidable,
  and unifiable problems have an mgu~\cite{Miller1991,Nipkow1993}. Its steps include an occurs check and variable
  elimination with pruning. $\rightarrow$ \cref{ch:hounif}; \book{220}--\book{223}.
\item[Higher-order programming] Passing predicates as arguments, as in \lp{mappred} and \lp{foldl}.
  $\rightarrow$ \cref{ch:hoprog}; \book{126}.
\item[Higher-order unification] Unification of simply typed $\lambda$-terms modulo $\lambda$-conversion. It is
  undecidable (the book reduces Post's correspondence problem and Hilbert's tenth problem to it), has no mgus in
  general, and is semi-decided by Huet's procedure~\cite{Huet1975}. $\rightarrow$ \cref{ch:lambda,ch:hounif};
  \book{111}, \book{114}--\book{115}, \book{214}.
\item[Horn clauses] The class of programs in which the signature and the program stay fixed during search: they are
  \enquote{flat and global}. $\rightarrow$ \cref{ch:horn}; \book{43}, \book{46}.
\item[Huet's procedure, pre-unification] Repeated simplification, imitation and projection, stopping when only
  flex-flex pairs remain~\cite{Huet1975}. The result is a pre-unifier, whose search space is explored as a
  matching tree. $\rightarrow$ \cref{ch:hounif}; \book{218}--\book{219}.
\item[Hypothetical reasoning] Proving a goal under temporarily assumed clauses, by AUGMENT. The book's example is
  the question \enquote{If person X took just one more course Y, will he or she graduate with a CS degree?}, posed as the query \lp{fact (finished X Y) => (fact (graduates X), fact (cs_major X))}.
  $\rightarrow$ \cref{ch:hh}; \book{79}--\book{80}.
\end{glossarylist}

\subsection*{I}
\begin{glossarylist}
\item[\texttt{if}, \texttt{not}] λProlog's conditional and negation-as-failure, both defined with the cut.
  $\rightarrow$ \cref{ch:hoprog}; \book{137}.
\item[\texttt{if°} (soft cut)] \texttt{(if° g? gt gf)}: if \texttt{g?} has at least one answer, every one of
  its answers continues with \texttt{gt}; otherwise \texttt{gf} runs. It corresponds to miniKanren's
  \texttt{conda} and LogicT's \texttt{ifte}. λProlog's \lp{if} becomes \texttt{(if° (take° 1 P) Q R)}.
  $\rightarrow$ \cref{ch:micro}.
\item[Imitation, projection] Huet's partial bindings for a flex-rigid pair. \emph{Imitation} copies the rigid head:
  $F:=\lambda\bar x.\,c\,(H_1\bar x)\cdots(H_m\bar x)$. \emph{Projection} selects an argument:
  $F:=\lambda\bar x.\,x_i\,(H_1\bar x)\cdots(H_k\bar x)$, which requires matching target types.
  $\rightarrow$ \cref{ch:hounif}; \book{215}.
\item[Intensional / extensional equality] Intensional equality compares the structure of $\lambda$-terms.
  Extensional equality compares what they denote. λProlog's \lp{=} is intensional (modulo $\lambda$-conversion).
  $\rightarrow$ \cref{ch:hoprog}; \book{146}.
\item[Intensional / extensional $\forall$] \lp{pi} is intensional: it requires one generic proof, so it cannot
  express \enquote{for all instances}. This is why book Ch.~11's simulation is incorrect. Jin et al.\ implement an
  extensional $\forall$ in miniKanren, which is a different logic~\cite{JinEtAl2021}.
  $\rightarrow$ \cref{ch:hh,ch:apps}; \book{83}, \book{94}.
\item[Interleaving] miniKanren's search discipline. In \micro{}, \texttt{orª} uses
  \texttt{append§/interleaved/2}, which swaps the two streams whenever the first one is a suspension, and
  conjunction (\texttt{append-map§}) interleaves too. The right-nested \texttt{or°} gives the first of three
  infinite branches about half of the answers. $\rightarrow$ \cref{ch:micro}.
\item[Inverse-$\eta$ delay] The wrapping \texttt{(λ (s) (δ ((and° g …) s)))} that \texttt{define-relation} puts
  around a body, so that a recursive relation without \texttt{fresh°} does not loop while the goal tree is built.
  $\rightarrow$ \cref{ch:micro}.
\item[Inversion of a pattern substitution] The core step of pattern unification. To solve $X\,\bar y = t$, where
  $\bar y$ are distinct bound variables, abstract $t$ over $\bar y$, that is, $X:=\lambda\bar y.\,t$. This works only
  if every bound variable free in $t$ is among $\bar y$ and $X$ does not occur in $t$; otherwise pruning or the occurs
  check takes over. DHKP write the inverse as $b[\xi]^{-1}$, and Nipkow's code remaps de Bruijn indices.
  $\rightarrow$ \cref{ch:hounif}.
\item[Invertible rule] A rule whose premises are provable if and only if its conclusion is. Such rules can be
  applied eagerly without losing completeness. $\rightarrow$ \cref{ch:apps}; \book{243}.
\item[Iterative deepening] Retrying a search with bound $N+1$ after failing with bound $N$. In book Ch.~9 it bounds
  the multiplicity of $\exists$-formulas. $\rightarrow$ \cref{ch:apps}; \book{237}, \book{239}.
\end{glossarylist}

\subsection*{K}
\begin{glossarylist}
\item[Kind] The \enquote{type} of a type constructor, of the form \lp{type -> ... -> type}, introduced by a
  \lp{kind} declaration. $\rightarrow$ \cref{ch:intro}; \book{11}--\book{12}.
\item[Kripke model] The possible-worlds semantics of intuitionistic logic. Its worlds correspond to
  signature--program pairs. $\rightarrow$ \cref{ch:hh}; \book{94}.
\end{glossarylist}

\subsection*{L}
\begin{glossarylist}
\item[$\Llam$] The sublanguage of hohh in which quantification is never over predicate types, and every
  essentially existential variable is applied only to distinct essentially universal variables bound within its
  scope. Unification inside $\Llam$ needs only $\beta_0$, is decidable and produces mgus.
  $\rightarrow$ \cref{ch:lambdatree,ch:hounif}; \book{205}--\book{206}, \book{220}.
\item[$\lambda$-conversion, $\lambda$-normal form] $\lambda$-conversion is the closure of $\alpha$, $\beta$ and
  $\eta$, and it is the equality of the logic. The $\lambda$-normal form $\lambda norm(s)$ is the $\beta\eta$-normal
  form. $\rightarrow$ \cref{ch:lambda}; \book{100}--\book{102}.
\item[λ-tree syntax] HOAS with simply typed $\lambda$-terms, taken modulo $\alpha$, a weak $\beta$ ($\beta_0$) and
  possibly $\eta$, used as a data structure that unification can still analyse. $\rightarrow$
  \cref{ch:lambdatree}; \book{204}.
\item[λKanren] miniKanren extended with higher-order logic programming. The 2020 version has terms \texttt{tie}
  and \texttt{app}, goals \texttt{assume} ($D\imp G$) and \texttt{all} ($\forall x\,G$), and Miller's pattern
  unification over an explicit prefix~\cite{MaLuFriedman2020}. The 2021 paper proposes static
  closures~\cite{MaFriedman2021}. $\rightarrow$ \cref{ch:lkanren}.
\item[Left- and right-introduction rules] Right rules introduce a connective on the right of $\longrightarrow$ and
  implement goal reduction. Left rules decompose the clause in focus and implement backchaining.
  $\rightarrow$ \cref{ch:horn}; \book{40}, \book{42}, \book{45}.
\item[Legitimate substitution] For a problem $\forall\bar u\,\exists\bar x\,\forall\bar w\,[\dots]$, a substitution
  for $\bar x$ in whose range no $w_j$ occurs free. Only legitimate substitutions can be unifiers. This is the
  eigenvariable condition, stated for unification. $\rightarrow$ \cref{ch:lambda}; \book{109}.
\item[Leibniz equality] $x=y$ if and only if $\forall P\,(P\,x\imp P\,y)$. $\rightarrow$ \cref{ch:hoprog};
  \book{120}.
\item[Level of a unification variable] The scope depth of a variable. ELPI's $X^j$ may mention only context names
  below $j$. In Teyjus, a variable may be bound only to terms whose constants have a tag no larger than its own.
  Binding a variable must lower the levels of the variables in the bound term.
  $\rightarrow$ \cref{ch:hounif,ch:lkanren}.
\item[Locally nameless] A representation with bound variables as de Bruijn indices and free variables as
  names~\cite{McBrideMcKinna2004,Chargueraud2012}. Eigen constants would play the role of the names.
  $\rightarrow$ \cref{ch:lkanren}.
\item[Logic variable] A placeholder that search instantiates by unification. It is the operational device behind
  the delayed choice of terms in the $\exists$R, $\forall$L and \rulename{initial} rules.
  $\rightarrow$ \cref{ch:horn}; \book{61}.
\end{glossarylist}

\subsection*{M}
\begin{glossarylist}
\item[Matching] A unification problem in which only one side contains existential variables. Higher-order matching
  is decidable. $\rightarrow$ \cref{ch:hounif}; \book{227}.
\item[Matching tree] The search tree of Huet's procedure. Its nodes are multisets of equations and its arcs are
  substitutions. A failure node is a problem that obviously has no solution (for example, a rigid-rigid clash); a success node
  contains at most flex-flex pairs, and reaching one is a pre-unification. Some paths may be infinite.
  $\rightarrow$ \cref{ch:hounif}; \book{219}.
\item[Matrix] The conjunction of equations $t_1=s_1\wedge\dots\wedge t_m=s_m$ under the quantifier prefix of a
  unification problem. $\rightarrow$ \cref{ch:hounif}; \book{211}.
\item[May judgment, must judgment] A may judgment asserts that some computation path exists. A must judgment states
  a property of all paths, and needs negation-as-failure or a meta-level $\forall$.
  $\rightarrow$ \cref{ch:apps}; \book{269}.
\item[Metalanguage, object language] λProlog, the language in which we write, and the language encoded in it.
  $\rightarrow$ \cref{ch:intro}; \book{22}.
\item[Minimal logic] Logic in which $\bot$ is an ordinary constant. Negation is defined as
  $\neg B := B\imp\cst{false}$. $\rightarrow$ \cref{ch:hh}; \book{89}--\book{90}.
\item[miniKanren] A family of small relational languages embedded in a host language~\cite{FriedmanEtAl2018,Byrd2009}.
  $\rightarrow$ \cref{ch:micro}.
\item[Mixed prefix] The quantifier prefix $\forall\exists\forall\dots$ of a unification problem. In a
  $\forall\exists\forall$ problem the outer universals act as the signature and the inner ones as eigenvariables,
  and no raising is needed. $\rightarrow$ \cref{ch:lambda,ch:hounif}; \book{105}, \book{211}.
\item[Module, signature (Teyjus)] A module is a named collection of declarations and clauses. Its signature
  restricts what is visible from outside, and hidden constants are read as existentially quantified (book
  Ch.~6). $\rightarrow$ \cref{ch:intro}; \book{152}--\book{153}, \book{277}.
\item[Most general unifier (mgu)] A unifier $\theta$ such that every other unifier is $\theta$ followed by a
  further substitution. In the first-order case an mgu exists whenever a unifier does, and it is unique up to
  renaming. Pattern problems have mgus too, but general higher-order problems do not.
  $\rightarrow$ \cref{ch:intro,ch:hounif}; \book{28}, \book{213}, \book{222}.
\item[μKanren] Hemann and Friedman's minimal miniKanren, in which goals are functions from a state to a stream of
  states~\cite{HemannFriedman2013}. \micro{} is a μKanren with constraint stores. $\rightarrow$ \cref{ch:micro}.
\item[Multiplicity] In the classical prover of book Ch.~9, the number of times an $\exists$-formula may be
  instantiated. Bounding it, and raising the bound after each failure, gives a complete semi-decision procedure.
  $\rightarrow$ \cref{ch:apps}; \book{237}.
\end{glossarylist}

\subsection*{N}
\begin{glossarylist}
\item[$\nabla$ (nabla)] The Miller--Tiu generic quantifier for fresh names, in a logic with definitions (closed
  world). It is not part of λProlog. $\rightarrow$ \cref{ch:hh,ch:apps}; \book{94}, \book{275}.
\item[Name shadowing] An inner binder that reuses an outer binder's name, so that the outer binder can no longer be
  named. $\rightarrow$ \cref{ch:lkanren}.
\item[Natural deduction] Introduction and elimination rules, with hypotheses discharged by $\imp$I, $\vee$E and
  $\exists$E. Book Ch.~9 encodes it with \lp{=>} and \lp{pi}. $\rightarrow$ \cref{ch:apps}; \book{231}.
\item[Negation-as-failure] \lp{not G} succeeds if and only if \lp{G} finitely fails. It is non-logical and unsound
  on non-ground goals. In \micro{} it is written \texttt{(if° (take° 1 g) ✗° ✓°)}. Without the \texttt{take° 1},
  it diverges whenever the search for further answers of \texttt{g} does not terminate, even if \texttt{g} has a
  single answer (\cref{sec:glossary-bridge} runs an example). $\rightarrow$ \cref{ch:hoprog,ch:apps}; \book{137},
  \book{269}.
\item[Negation normal form, prenex normal form] NNF has no implications and applies negation only to atoms. A
  prenex form has all its quantifiers in front. $\rightarrow$ \cref{ch:lambdatree}; \book{195}.
\item[Nominal logic, nominal unification] A first-order treatment of binding with names, swappings and freshness.
  Nominal unification works modulo $\alpha$ and never generates new names~\cite{UrbanPittsGabbay2004}.
  $\rightarrow$ \cref{ch:lkanren}; \book{210}.
\end{glossarylist}

\subsection*{O}
\begin{glossarylist}
\item[\texttt{o}] The type of propositions, named after Church's omicron. The quantifier names \lp{pi} and
  \lp{sigma} likewise follow Church's $\Pi$ and $\Sigma$. $\rightarrow$ \cref{ch:intro}; \book{74}.
\item[Object-level substitution] Substituting into the body of an encoded binder. In λ-tree syntax it is meta-level
  $\beta$-reduction, as in the clause \lp{subst R N (R N).} with \lp{R : tm -> tm}. That clause lies outside
  $\Llam$, but it is used only with \lp{R} and \lp{N} instantiated (dynamic $\Llam$). Copy clauses specify the same
  relation inside $\Llam$. $\rightarrow$ \cref{ch:lambdatree}; \book{180}, \book{199}, \book{225}.
\item[O-proof] A derivation that uses only goal reduction (right rules) followed by backchaining (focused left
  rules). Every provable fohc sequent has an O-proof. $\rightarrow$ \cref{ch:horn}; \book{45}.
\item[Occurs check] $x=t$ fails when $x$ occurs properly in $t$. In pattern unification, the corresponding failure
  is a flexible head occurring rigidly on the other side. \micro{} performs the check when it extends the
  substitution. $\rightarrow$ \cref{ch:intro,ch:hounif}; \book{30}, \book{223}.
\item[Order of a type] $0$ for non-functional types, and $\max(\mathrm{ord}\,\tau_1+1,\mathrm{ord}\,\tau_2)$ for
  $\tau_1\to\tau_2$. $\rightarrow$ \cref{ch:intro,ch:lambda}; \book{13}.
\end{glossarylist}

\subsection*{P}
\begin{glossarylist}
\item[Pattern] See \emph{Higher-order pattern}. Ma and Friedman's phrasing, for an application form $F\,a$:
  \enquote{if $F$ is a unification variable, then $a$ must be a list of zero or more distinct bound variables}.
\item[Peirce's formula] $((p\imp q)\imp p)\imp p$: valid classically but not intuitionistically. It has no
  O-proof. $\rightarrow$ \cref{ch:hh}; \book{89}.
\item[Pervasive] Built in and globally available: sorts such as \lp{int}, constants such as \lp{nil} and \lp{::},
  and predicates such as \lp{is}. $\rightarrow$ \cref{ch:intro}; \book{14}, \book{49}.
\item[\texttt{pi}, \texttt{sigma}] λProlog's quantifiers, constants of type \lp{(A -> o) -> o}. $\forall x\,G$
  is written \lp{pi x\ G}, the application of \lp{pi} to the abstraction \lp{x\ G}, so the term-level $\lambda$
  is the only binder. $\rightarrow$ \cref{ch:intro,ch:lambda}; \book{36}, \book{99}.
\item[$\pi$-calculus] Milner's calculus of mobile processes, in which names are the only data. Book Ch.~11 encodes
  it in λ-tree syntax, using \lp{in} and \lp{nu} as binders. $\rightarrow$ \cref{ch:apps}; \book{261}.
\item[Polarity, positive and negative occurrence] An occurrence is positive if it lies to the left of an even
  number of implications and negative if odd. Cut-free search preserves polarity, unless implications appear
  inside atoms. $\rightarrow$ \cref{ch:hh,ch:hoprog}; \book{76}, \book{121}--\book{122}.
\item[Pre-cooking] (DHKP) Translating a metavariable that occurs under $n$ binders into $X[\uparrow^n]$, so that
  first-order grafting agrees with capture-avoiding substitution~\cite{DowekHardinKirchnerPfenning1996}.
  $\rightarrow$ \cref{ch:lkanren}.
\item[Pre-unifier] A substitution that reduces a problem to flex-flex pairs only. Such a remainder always has a
  unifier, the canonical substitution. $\rightarrow$ \cref{ch:hounif}; \book{217}--\book{218}.
\item[\texttt{project°}] \micro{}'s escape to Scheme: it binds Scheme variables to deep-walked values, in which
  unbound variables remain variable records, and then runs goals. It is non-relational.
  $\rightarrow$ \cref{ch:micro}.
\item[Proof term] A $\lambda$-term of type \lp{proof} that records a derivation, as in the judgment \lp{P # A}
  (\enquote{$P$ is a proof of $A$}). $\rightarrow$ \cref{ch:apps}; \book{233}.
\item[Pruning] Removing from a variable its dependency on arguments that cannot occur in any solution:
  $G:=\lambda\bar d.\,(G'\,\bar e)$, where $\bar e$ keeps only the allowed arguments. $\rightarrow$
  \cref{ch:hounif}; \book{223}.
\item[Purification] The step of \micro{}'s reification that drops constraints mentioning variables unreachable from
  the answer. $\rightarrow$ \cref{ch:micro}.
\end{glossarylist}

\subsection*{R}
\begin{glossarylist}
\item[Raising] $\forall y\,\exists x\,G \rightsquigarrow \exists h\,\forall y\,\subst{G}{(h\,y)}{x}$: an existential in the
  scope of a universal becomes the application of a new, higher-type existential quantified further out. Raising
  preserves $\Llam$. The book calls it a dual of Skolemization: Skolemization introduces a new \emph{constant} of a
  raised type, raising a new \emph{existential variable} of a raised type. $\rightarrow$ \cref{ch:lambda,ch:hounif};
  \book{108}, \book{208}, \book{212}.
\item[Reduction rules] AND, OR, INSTAN, AUGMENT, GENERIC and TRUE: the search meaning of each logical symbol.
  $\rightarrow$ \cref{ch:horn}; \book{39}.
\item[Reduction-free fragment $L^{\beta}_{\lambda}$] (ELPI) The patterns in which a unification variable $X^j$ is
  applied exactly to the bound variables $x_j\dots x_{j+k-1}$ that follow its level. Binding such a variable needs no
  $\beta$-reduction and takes constant time, but the fragment is not stable under resolution~\cite{DunchevEtAl2015}.
  $\rightarrow$ \cref{ch:lkanren}.
\item[Reification] Turning the final state into a presentable answer. \micro{} produces a closed Scheme expression
  \texttt{(λ (α β …) constraint-form … term)}, naming free variables in order of first occurrence. αKanren prints
  noms as \texttt{a.0}, \texttt{a.1}, \dots{} and adds freshness constraints. $\rightarrow$ \cref{ch:micro}.
\item[Relation (\micro{})] A procedure from terms to a goal, conventionally named \texttt{name°} and defined with
  \texttt{define-relation}. $\rightarrow$ \cref{ch:micro}.
\item[Relative complement] (Jin et al.) A goal that describes the values of a variable not covered by a state. It
  is used to implement an extensional $\forall$~\cite{JinEtAl2021}. $\rightarrow$ \cref{ch:lkanren}.
\item[Representation independence] Replacing a module by an equivalent one does not change the meaning of the
  larger program. $\rightarrow$ \book{151}.
\item[Rigid] See \emph{Flexible term, rigid term} and \emph{Flexible atom}.
\end{glossarylist}

\subsection*{S}
\begin{glossarylist}
\item[Scope check, eigenvariable condition] No variable older than an eigenvariable may reach it. The
  implementations differ: an explicit quantifier prefix (Miller; λKanren 2020), universe indices or levels
  (Teyjus, ELPI), an absento-like constraint over the variables in scope (\texttt{eigen}), or an end-of-scope check,
  which is fragile. $\rightarrow$ \cref{ch:lkanren,ch:micro}.
\item[Scope extrusion] The book's name for the classical equivalence
  $B_1\vee(B_2\imp B_3)\equiv B_2\imp(B_1\vee B_3)$, which moves the scope of an implication over a disjunction. It
  fails intuitionistically: $p\vee(p\imp q)$ has no O-proof, although its classical equivalent $p\imp(p\vee q)$ does. The
  term is borrowed from the $\pi$-calculus, where it means sending a restricted name out and re-closing its scope at
  the receiver. $\rightarrow$ \cref{ch:hh,ch:apps}; \book{89}, \book{95}.
\item[Scope restriction] Ma and Friedman's name for the eigenvariable condition in unification: a $\forall$-bound
  variable may occur in the instantiation of an $\exists$-variable only if it precedes that variable in the
  prefix. $\rightarrow$ \cref{ch:lkanren}.
\item[Shallow embedding] Implementing a logic language as host-language procedures and macros, as miniKanren
  does in Scheme, instead of writing an interpreter or compiler for it. λKanren keeps this
  style~\cite{MaLuFriedman2020}. $\rightarrow$ \cref{ch:lkanren,ch:micro}.
\item[Sequent] $\seq{\Sig}{\Prog}{G}$: signature, program and goal. It is read both as the state of an interpreter
  and as a logical judgment. $\rightarrow$ \cref{ch:horn}; \book{39}.
\item[Signature] A set of typed non-logical constants. In Teyjus, a signature is also a module's interface
  (\lp{sig}). $\rightarrow$ \cref{ch:intro}; \book{17}, \book{153}.
\item[Signature--program pair] $\langle\Sig,\Prog\rangle$, the context in which queries are posed. Under the
  open-world assumption these pairs are the worlds of a Kripke model. $\rightarrow$ \cref{ch:horn,ch:hh};
  \book{39}, \book{94}.
\item[Simplification] Repeated decomposition, or clash, of rigid-rigid pairs in higher-order unification.
  $\rightarrow$ \cref{ch:hounif}; \book{214}.
\item[Skolemization] Replacing an existential in an assumption by a function of the enclosing universals, as in
  $\forall x\,\subst{D}{f\,x}{y}$. It is the dual of raising. $\rightarrow$ \cref{ch:lambda}; \book{108}.
\item[Soft cut] See \texttt{if°}. A soft cut is not complete in general.
\item[Solution] A unifier whose range has free variables only among the outer universals. A solution amounts to a
  proof of the quantified equation. $\rightarrow$ \cref{ch:lambda,ch:hounif}; \book{109}.
\item[Solved form] A set of equations $x=t$ in which each left-hand variable occurs nowhere else. It is read as an
  idempotent mgu. $\rightarrow$ \cref{ch:intro}; \book{29}.
\item[Static closure] $\langle\phi;t\rangle$ (Ma--Friedman): a term paired with the scope, the list of binder names
  around it, manipulated by $\mathrm{ext}$ and $\mathrm{db}$. The authors call the paper \enquote{a preview} and
  leave soundness and completeness to future work~\cite{MaFriedman2021}. $\rightarrow$ \cref{ch:lkanren}.
\item[Stream (\texttt{§})] In \aux{}, a stream is \texttt{'()}, a mature pair \texttt{(state . §)}, or an immature
  promise. $\rightarrow$ \cref{ch:micro}.
\item[Subject reduction] Evaluation preserves types. $\rightarrow$ \cref{ch:lambdatree}; \book{191}.
\item[Substitution (\micro{})] The store \texttt{S}: a triangular substitution kept in a skew-binary random-access
  list indexed by variable number. $\rightarrow$ \cref{ch:micro}.
\item[Subsumption (\micro{})] Dropping a constraint that another constraint implies, or that a type tag or an
  absento constraint makes vacuous. $\rightarrow$ \cref{ch:micro}.
\item[Suspension] Four different things, depending on the source; see the pitfall at the start of this appendix.
\item[Swapping] $(a\;b)$ exchanges the atoms $a$ and $b$ throughout a term. Permutations are lists of swappings.
  $\rightarrow$ \cref{ch:lkanren}.
\item[Syntactic constraint] (ELPI) A goal suspended on some variables with \texttt{declare\_constraint} and resumed
  when one of them is assigned. It keeps its hypothetical context. $\rightarrow$ \cref{ch:lkanren}.
\end{glossarylist}

\subsection*{T}
\begin{glossarylist}
\item[\texttt{take°}] \texttt{(take° n g)} keeps the first \texttt{n} answers of \texttt{g}; a negative
  \texttt{n} keeps them all. It is a goal, so it nests, and \texttt{(take° 1 g)} is miniKanren's \texttt{once}.
  \texttt{μkanren-run} is \texttt{take°} around a \texttt{fresh°} query. $\rightarrow$ \cref{ch:micro}.
\item[Tactic, tactical] A tactic is a relation \lp{goal -> goal -> o} from a primitive goal to the goal of its
  premises. Tacticals such as \lp{then}, \lp{orelse}, \lp{repeat}, \lp{try} and \lp{maptac} combine tactics. Book
  Ch.~9 recasts the LCF architecture relationally. $\rightarrow$ \cref{ch:apps}; \book{241}, \book{243}.
\item[Target type] $\tau_0$ in $\tau_1\to\dots\to\tau_n\to\tau_0$. $\rightarrow$ \cref{ch:intro}; \book{13}.
\item[Term reduction, variable elimination] Two of the three transformations of first-order unification; the third
  is reorientation. Term reduction decomposes equal heads. Variable elimination replaces a variable by its value
  everywhere. $\rightarrow$ \cref{ch:intro}; \book{28}.
\item[Teyjus] The λProlog implementation used by the book, with compiler \texttt{tjcc}, linker \texttt{tjlink} and
  simulator \texttt{tjsim}. Teyjus~2 uses universe-level tags. $\rightarrow$ \cref{ch:intro}; \book{277}.
\item[Trace, complete trace, trace equivalence] A trace is a sequence of actions of a process, and a complete trace
  ends in a terminal process. Two processes are trace equivalent if no trace separates them.
  $\rightarrow$ \cref{ch:apps}; \book{267}, \book{269}, \book{271}.
\item[Transparent type variable] A type variable that occurs in the target type. $\rightarrow$ \cref{ch:horn};
  \book{67}.
\item[Type lowering] A polarity-indexed translation from a relation at a higher type to formulas over base-type
  relations. It generates copy clauses. $\rightarrow$ \cref{ch:lambdatree}; \book{201}--\book{202}.
\end{glossarylist}

\subsection*{U}
\begin{glossarylist}
\item[Unification problem, unifier] In the first-order case, a problem is a finite multiset of equations, and a
  unifier is a type-preserving substitution that makes both sides of each equation equal. In general, a problem
  is a quantified formula $Q_1x_1\dots Q_nx_n[\bigwedge_j t_j=s_j]$, and a unifier is a legitimate substitution
  (no inner universals in its range) that makes each pair $\lambda$-convertible.
  $\rightarrow$ \cref{ch:intro,ch:lambda}; \book{26}, \book{105}, \book{109}.
\item[Unification (\micro{}), \texttt{=°}] First-order syntactic unification with occurs check. Pairs,
  vectors and records are decomposed slot by slot, so a vector and a record with the same slots unify. When two
  variables meet, the one with the lower index is bound to the other. $\rightarrow$ \cref{ch:micro}.
\item[Uniform proof] A cut-free sequent proof in which every sequent with a non-atomic goal is the conclusion of the
  right-introduction rule for that goal's top-level connective. It formalises goal-directed
  search~\cite{MillerNadathurPfenningScedrov1991}. $\rightarrow$ \cref{ch:horn}; \book{73}.
\item[Universe index] (Nadathur, Jayaraman, Kwon) A numeric tag on constants and logic variables that records the
  $\forall$-nesting at their creation. Entering a universal goal increments the index. A variable may be bound only
  to terms whose constants have a tag no larger than its own, and the binding lowers the tags of the variables in
  the term~\cite{NadathurJayaramanKwon1995}. The book mentions the idea as numerical indices on quantified
  variables \book{225} and as universe level tags \book{227}. $\rightarrow$ \cref{ch:hounif,ch:lkanren}.
\item[Universe level] See \emph{Universe index} and \emph{Level of a unification variable}.
\end{glossarylist}

\subsection*{V}
\begin{glossarylist}
\item[Vacuous quantifier] A quantifier whose variable does not occur in its body. $\rightarrow$
  \cref{ch:lambdatree}; \book{195}.
\item[Value constructor] A constant of functional type used to build data, such as \lp{::}.
  $\rightarrow$ \cref{ch:intro}; \book{14}.
\item[Variable elimination (patterns)] The step of pattern unification that combines imitation, projection, the
  occurs check and pruning into one binding $F:=\lambda\bar c.\,t'$. $\rightarrow$ \cref{ch:hounif}; \book{223}.
\end{glossarylist}

\subsection*{W}
\begin{glossarylist}
\item[Walk] Resolving a variable through the substitution. \texttt{μkanren-state-find} is the shallow walk: it
  follows the chain of variable bindings until it reaches a non-variable or an unbound variable.
  \texttt{μkanren-state-find/value} is the deep walk: it also walks inside pairs, vectors and records. $\rightarrow$ \cref{ch:micro}.
\item[WAM] The Warren Abstract Machine, the classical implementation model of Prolog. $\rightarrow$ \book{74}.
\item[Working variable, reified variable] In \micro{}, a working variable has index $\ge 0$, lives during search
  and prints as \texttt{\_i}. A reified variable has a negative index, appears only in answers and prints as α, β,
  \dots{} $\rightarrow$ \cref{ch:micro}.
\end{glossarylist}

\endgroup

\section{A vocabulary bridge}\label{sec:glossary-bridge}

\Cref{tab:glossary-bridge} lines up the names that the three communities use for the same concept. Read it
across a row: a cell left empty means the concept has no counterpart in that column today.

\begin{table}[ht]
\centering\small
\caption{The same concepts under three names, and in the proposed module. The empty cells of the third
column are what \texttt{(aux kanren lambda)} supplies (\cref{ch:module}).}
\label{tab:glossary-bridge}
\begin{tabularx}{\textwidth}{@{}>{\raggedright}X>{\raggedright}X>{\raggedright}X>{\raggedright\arraybackslash}X@{}}
\toprule
Book (λProlog) & miniKanren literature & \micro{} today & \texttt{(aux kanren lambda)} \\
\midrule
logic variable, $\exists$ in goals (\texttt{sigma}) & \texttt{fresh} & \texttt{fresh°} & \texttt{fresh°}, with a level \\
backchaining on a clause & calling a relation & a \texttt{define-relation} call & also a \texttt{define-dynamic-relation} call \\
$\vee$ in goals (\texttt{;}) & \texttt{conde} & \texttt{or°}, \texttt{cond°} & unchanged \\
first-order unification with occurs check & \texttt{==} & \texttt{=°} & \texttt{=°}, through the \texttt{unify} and \texttt{post} hooks \\
\texttt{not}, \texttt{if} (via \texttt{!}) & \texttt{conda}, \texttt{condu} & \texttt{if°}, \texttt{take°} & unchanged \\
\rulename{augment}, \texttt{D => G} & \texttt{assume} (λKanren) & --- & \texttt{⊃°}, \texttt{define-dynamic-relation} \\
\rulename{generic}, \texttt{pi x\textbackslash{} G}, eigenvariable & \texttt{all} (λKanren), \texttt{eigen} & --- & \texttt{∀°}, with universe levels \\
$\lambda$-terms, \texttt{x\textbackslash{} t} & \texttt{tie}/\texttt{app} (λKanren), \texttt{tie} on noms (αKanren) & first-order encodings only (e.g.\ records with de Bruijn indices) & \texttt{ƛ}, \texttt{@} \\
pattern unification, raising, pruning & λKanren 2020 unifier; static closures (2021) & --- & the \texttt{unify} and \texttt{post} hooks with levels; non-patterns delayed in E \\
answer substitution & reified answer & \texttt{μkanren-project}, \texttt{°->list} & \texttt{°->list/datum}, \texttt{°->pretty} \\
\bottomrule
\end{tabularx}
\end{table}

\begin{kanren}[what the glossary says about the missing pieces]
The empty cells of \cref{tab:glossary-bridge} are the λProlog features that \micro{} lacks: a dynamic clause
database for AUGMENT, eigenvariables with a scope check for GENERIC, and terms with binders together with a
unifier that respects them. The entries \emph{Universe index}, \emph{Level of a unification variable}, \emph{Scope
check}, \emph{Higher-order pattern unification} and \emph{Delayed equations} list the known ways to provide them.
\Cref{ch:design} chooses among them and \cref{ch:module} implements the choice.
\end{kanren}

\begin{kanren}[three entries checked on today's engine]
The following file exercises the entries \emph{Interleaving}, \emph{Negation-as-failure} and \emph{Reification}.
The relation \texttt{num°} enumerates $0,1,2,\dots$.
\begin{scheme}
(import (aux kanren micro) (aux base) (aux stream))

; num° enumerates 0, 1, 2, ... (project° escapes to Scheme to add 1)
(define-relation (num° n)
  (or° (=° n 0) (fresh° (m) (num° m) (project° ((m* m)) (=° n (+ 1 m*))))))

; Interleaving: three infinite branches; the right-nested or° favours the first
(print (°->list/ground
         (take° 12 (fresh° (q n)
                     (or° (and° (num° n) (=° q (list 'a n)))
                          (and° (num° n) (=° q (list 'b n)))
                          (and° (num° n) (=° q (list 'c n))))))))

; Negation as failure: (take° 1 g) makes the test finite
(print (°->list/ground (fresh° (q) (=° q 1) (if° (take° 1 (num° q)) ✗° ✓°))))
(print (°->list/ground (fresh° (q) (=° q 1) (if° (=° q 2) ✗° ✓°))))
(print (°->list/ground (fresh° (q) (if° (=° q 2) ✗° ✓°))))

; Reification: free variables become α, β, ...; constraints first
(print (°->list #f (fresh° (q x) (=° q (list x 'a)) (≠° x 'b))))
\end{scheme}
\begin{answers}
((a 0) (b 0) (a 1) (c 0) (a 2) (b 1) (c 1) (a 3) (a 4) (b 2) (c 2) (a 5))
()
(1)
()
((λ (α) (begin (deny (equal? α (quote b)))) (cons α (cons (quote a) (quote ())))))
\end{answers}
Read the lines in order. (1) Of the first twelve answers, six come from branch \texttt{a} and three each from
\texttt{b} and \texttt{c}: the shares are $\frac12,\frac14,\frac14$, as the entry \emph{Interleaving} says.
(2) $q=1$ is a number, so \texttt{not (num° 1)} fails. If you drop the \texttt{take° 1}, the run never returns:
\texttt{num° 1} has only one answer, but the search for a second one never ends. \texttt{if°} continues every
answer of the test with \texttt{✗°}, and over a stream that never ends this continuation never finishes either. (3) and (4) show negation-as-failure on a ground goal and on a
non-ground one. With $q$ unbound, the test \texttt{(=° q 2)} succeeds by binding $q$, so the negation fails,
although $q=1$ would satisfy it. This is the classic unsoundness. (5) The reified answer is a closed
\texttt{λ}-expression whose body lists the constraint before the term.
\end{kanren}

\section{Summary}\label{sec:glossary-summary}
\begin{itemize}
\item Proof-theoretic terms (sequent, uniform proof, AUGMENT, GENERIC, eigenvariable) describe \emph{what} search
  must do. Unification terms (pattern, raising, pruning, level) describe \emph{how} binders and quantifiers are
  respected. miniKanren terms (goal, stream, interleaving, reification) describe the engine the ideas must fit.
\item Watch out for overloaded words: suspension, level, scope extrusion, fresh, cut, substitution.
\item Every entry points to the chapter that develops it and to the book page where it is defined. Use the
  pointers to return to the full treatment.
\end{itemize}


\backmatter
\cleardoublepage
\phantomsection
\addcontentsline{toc}{chapter}{\bibname}
{\RaggedRight % long titles and URLs: ragged right, as tufte sets its own bibliographies
% bibliography.tex — one entry per key cited in chapters/*.tex, alphabetical by first author.
% Sources: the bibliography of Miller & Nadathur 2012 (book PDF pp. 300-311), notes/l01-literature.md §7,
% notes/p01-lambdakanren-2021.md §7, the papers in lit/ and the per-chapter files in bib/.
\begin{thebibliography}{99}

\bibitem{Andreoli1992}
Jean-Marc Andreoli.
\newblock Logic programming with focusing proofs in linear logic.
\newblock \emph{Journal of Logic and Computation} 2(3):297--347, 1992.

\bibitem{BrissetRidoux1991}
Pascal Brisset and Olivier Ridoux.
\newblock Na\"{\i}ve reverse can be linear.
\newblock In \emph{Eighth International Logic Programming Conference}, Paris, France, June 1991. MIT Press,
  Cambridge, MA, 1992.

\bibitem{deBruijn1972}
N.~G. de~Bruijn.
\newblock Lambda calculus notation with nameless dummies, a tool for automatic formula manipulation, with
  application to the Church--Rosser theorem.
\newblock \emph{Indagationes Mathematicae} 34(5):381--392, 1972.

\bibitem{Byrd2009}
William E. Byrd.
\newblock \emph{Relational Programming in miniKanren: Techniques, Applications, and Implementations}.
\newblock PhD thesis, Indiana University, 2009.

\bibitem{ByrdFriedman2007}
William E. Byrd and Daniel P. Friedman.
\newblock αKanren: A fresh name in nominal logic programming.
\newblock In \emph{Proceedings of the 2007 Workshop on Scheme and Functional Programming}, Université Laval
  Technical Report DIUL-RT-0701, pp.~79--90, 2007.

\bibitem{Chargueraud2012}
Arthur Charguéraud.
\newblock The locally nameless representation.
\newblock \emph{Journal of Automated Reasoning} 49(3):363--408, 2012.

\bibitem{Cheney2005}
James Cheney.
\newblock Relating nominal and higher-order pattern unification.
\newblock In \emph{Proceedings of UNIF 2005}, pp.~104--119, 2005.

\bibitem{DowekHardinKirchnerPfenning1996}
Gilles Dowek, Thérèse Hardin, Claude Kirchner, and Frank Pfenning.
\newblock Unification via explicit substitutions: the case of higher-order patterns.
\newblock In M.~Maher, editor, \emph{Joint International Conference and Symposium on Logic Programming
  (JICSLP)}, pp.~259--273. MIT Press, 1996.

\bibitem{DunchevEtAl2015}
Cvetan Dunchev, Ferruccio Guidi, Claudio Sacerdoti Coen, and Enrico Tassi.
\newblock ELPI: fast, embeddable, λProlog interpreter.
\newblock In \emph{Logic for Programming, Artificial Intelligence, and Reasoning (LPAR-20)}, LNCS 9450, pp.~460--468.
  Springer, 2015.

\bibitem{Dyckhoff1992}
Roy Dyckhoff.
\newblock Contraction-free sequent calculi for intuitionistic logic.
\newblock \emph{Journal of Symbolic Logic} 57(3):795--807, September 1992.

\bibitem{FelleisenHieb1992}
Matthias Felleisen and Robert Hieb.
\newblock The revised report on the syntactic theories of sequential control and state.
\newblock \emph{Theoretical Computer Science} 103:235--271, 1992.

\bibitem{Felty1993}
Amy Felty.
\newblock Implementing tactics and tacticals in a higher-order logic programming language.
\newblock \emph{Journal of Automated Reasoning} 11(1):43--81, August 1993.

\bibitem{FeltyMiller1988}
Amy Felty and Dale Miller.
\newblock Specifying theorem provers in a higher-order logic programming language.
\newblock In \emph{Ninth International Conference on Automated Deduction}, LNCS 310, pp.~61--80.
  Springer-Verlag, New York, May 1988.

\bibitem{Fischer1972}
Michael J. Fischer.
\newblock Lambda calculus schemata.
\newblock \emph{ACM SIGPLAN Notices} 7(1):104--109, January 1972.

\bibitem{FriedmanEtAl2018}
Daniel P. Friedman, William E. Byrd, Oleg Kiselyov, and Jason Hemann.
\newblock \emph{The Reasoned Schemer}, second edition.
\newblock MIT Press, 2018.

\bibitem{Gentzen1935}
Gerhard Gentzen.
\newblock Investigations into logical deductions.
\newblock In M.~E. Szabo, editor, \emph{The Collected Papers of Gerhard Gentzen}, pp.~68--131.
  North-Holland, Amsterdam, 1969. Translation of articles that appeared in 1934--35.

\bibitem{Goldfarb1981}
Warren Goldfarb.
\newblock The undecidability of the second-order unification problem.
\newblock \emph{Theoretical Computer Science} 13:225--230, 1981.

\bibitem{GordonMilnerWadsworth1979}
Michael J. Gordon, Arthur J. Milner, and Christopher P. Wadsworth.
\newblock \emph{Edinburgh LCF: A Mechanised Logic of Computation}, LNCS 78.
\newblock Springer-Verlag, New York, 1979.

\bibitem{HannanMiller1992}
John Hannan and Dale Miller.
\newblock From operational semantics to abstract machines.
\newblock \emph{Mathematical Structures in Computer Science} 2(4):415--459, 1992.

\bibitem{Harrop1960}
R.~Harrop.
\newblock Concerning formulas of the types $A\rightarrow B\vee C$, $A\rightarrow(Ex)B(x)$ in
  intuitionistic formal systems.
\newblock \emph{Journal of Symbolic Logic} 25:27--32, 1960.

\bibitem{HemannByrdFriedman2013}
Jason Hemann, William E. Byrd, and Daniel P. Friedman.
\newblock The \texttt{eigen} operator of \emph{miniKanren-with-symbolic-constraints}.
\newblock Software, 2013. \url{https://github.com/webyrd/miniKanren-with-symbolic-constraints}.

\bibitem{HemannFriedman2013}
Jason Hemann and Daniel P. Friedman.
\newblock μKanren: A minimal functional core for relational programming.
\newblock In \emph{Proceedings of the 2013 Workshop on Scheme and Functional Programming (Scheme '13)},
  2013.

\bibitem{Hudelmaier1992}
Jörg Hudelmaier.
\newblock Bounds on cut-elimination in intuitionistic propositional logic.
\newblock \emph{Archive for Mathematical Logic} 31:331--353, 1992.

\bibitem{Huet1973}
Gérard Huet.
\newblock The undecidability of unification in third order logic.
\newblock \emph{Information and Control} 22:257--267, 1973.

\bibitem{Huet1975}
Gérard Huet.
\newblock A unification algorithm for typed λ-calculus.
\newblock \emph{Theoretical Computer Science} 1:27--57, 1975.

\bibitem{JinEtAl2021}
Ende Jin, Gregory Rosenblatt, Matthew Might, and Lisa Zhang.
\newblock Universal quantification and implication in miniKanren.
\newblock In \emph{miniKanren and Relational Programming Workshop}, 2021.

\bibitem{Kahn1987}
Gilles Kahn.
\newblock Natural semantics.
\newblock In \emph{Proceedings of the Symposium on Theoretical Aspects of Computer Science}, LNCS 247,
  pp.~22--39. Springer, New York, March 1987.

\bibitem{LevyVillaret2008}
Jordi Levy and Mateu Villaret.
\newblock Nominal unification from a higher-order perspective.
\newblock In \emph{Rewriting Techniques and Applications (RTA 2008)}, LNCS 5117, pp.~246--260. Springer,
  2008.

\bibitem{MaFriedman2021}
Weixi Ma and Daniel P. Friedman.
\newblock A new higher-order unification algorithm for λKanren.
\newblock In \emph{miniKanren and Relational Programming Workshop}, 2021.

\bibitem{MaLuFriedman2020}
Weixi Ma, Kuang-Chen Lu, and Daniel P. Friedman.
\newblock Higher-order logic programming with λKanren.
\newblock In \emph{miniKanren and Relational Programming Workshop}, 2020.

\bibitem{McBrideMcKinna2004}
Conor McBride and James McKinna.
\newblock Functional pearl: I am not a number --- I am a free variable.
\newblock In \emph{Proceedings of the ACM SIGPLAN Workshop on Haskell}, 2004.

\bibitem{McDowellMiller2000}
Raymond McDowell and Dale Miller.
\newblock Cut-elimination for a logic with definitions and induction.
\newblock \emph{Theoretical Computer Science} 232:91--119, 2000.

\bibitem{MichaylovPfenning1992}
Spiro Michaylov and Frank Pfenning.
\newblock An empirical study of the runtime behavior of higher-order logic programs.
\newblock In \emph{Proceedings of the Workshop on the λProlog Programming Language}, 1992.

\bibitem{Miller1991}
Dale Miller.
\newblock A logic programming language with lambda-abstraction, function variables, and simple unification.
\newblock \emph{Journal of Logic and Computation} 1(4):497--536, 1991.

\bibitem{Miller1992}
Dale Miller.
\newblock Unification under a mixed prefix.
\newblock \emph{Journal of Symbolic Computation} 14(4):321--358, 1992.

\bibitem{MillerNadathur2012}
Dale Miller and Gopalan Nadathur.
\newblock \emph{Programming with Higher-Order Logic}.
\newblock Cambridge University Press, Cambridge, 2012.

\bibitem{MillerNadathurPfenningScedrov1991}
Dale Miller, Gopalan Nadathur, Frank Pfenning, and Andre Scedrov.
\newblock Uniform proofs as a foundation for logic programming.
\newblock \emph{Annals of Pure and Applied Logic} 51:125--157, 1991.

\bibitem{MillerTiu2005}
Dale Miller and Alwen Tiu.
\newblock A proof theory for generic judgments.
\newblock \emph{ACM Transactions on Computational Logic} 6(4):749--783, October 2005.

\bibitem{Milner1990}
Robin Milner.
\newblock Functions as processes.
\newblock In \emph{Automata, Languages and Programming, 17th International Colloquium}, LNCS 443,
  pp.~167--180. Springer-Verlag, New York, July 1990.

\bibitem{MilnerParrowWalker1992}
Robin Milner, Joachim Parrow, and David Walker.
\newblock A calculus of mobile processes, part I.
\newblock \emph{Information and Computation} 100(1):1--40, September 1992.

\bibitem{Moiseenko2019}
Evgenii Moiseenko.
\newblock Constructive negation for miniKanren.
\newblock In \emph{miniKanren and Relational Programming Workshop}, 2019.

\bibitem{Nadathur2005}
Gopalan Nadathur.
\newblock A treatment of higher-order features in logic programming.
\newblock \emph{Theory and Practice of Logic Programming} 5(3):305--354, 2005. arXiv:cs/0404020.

\bibitem{NadathurJayaramanKwon1995}
Gopalan Nadathur, Bharat Jayaraman, and Keehang Kwon.
\newblock Scoping constructs in logic programming: implementation problems and their solution.
\newblock \emph{Journal of Logic Programming} 25(2):119--161, November 1995.

\bibitem{NadathurLinnell2005}
Gopalan Nadathur and Natalie Linnell.
\newblock Practical higher-order pattern unification with on-the-fly raising.
\newblock In \emph{ICLP 2005: 21st International Logic Programming Conference}, LNCS 3668, pp.~371--386,
  Sitges, Spain, October 2005. Springer, 2006.

\bibitem{NadathurMiller1990}
Gopalan Nadathur and Dale Miller.
\newblock Higher-order Horn clauses.
\newblock \emph{Journal of the ACM} 37(4):777--814, October 1990.

\bibitem{NadathurMitchell1999}
Gopalan Nadathur and Dustin J. Mitchell.
\newblock System description: Teyjus --- a compiler and abstract machine based implementation of λProlog.
\newblock In \emph{Automated Deduction (CADE-16)}, LNAI 1632, pp.~287--291. Springer, 1999.

\bibitem{NearByrdFriedman2008}
Joseph P. Near, William E. Byrd, and Daniel P. Friedman.
\newblock αleanTAP: A declarative theorem prover for first-order classical logic.
\newblock In \emph{Logic Programming (ICLP 2008)}, LNCS 5366, pp.~238--252. Springer, 2008.

\bibitem{Nipkow1993}
Tobias Nipkow.
\newblock Functional unification of higher-order patterns.
\newblock In M.~Vardi, editor, \emph{8th Symposium on Logic in Computer Science}, pp.~64--74. IEEE, New
  York, June 1993.

\bibitem{PfenningElliott1988}
Frank Pfenning and Conal Elliott.
\newblock Higher-order abstract syntax.
\newblock In \emph{Proceedings of the ACM-SIGPLAN Conference on Programming Language Design and
  Implementation}, pp.~199--208. ACM Press, New York, June 1988.

\bibitem{Qian1996}
Zhenyu Qian.
\newblock Unification of higher-order patterns in linear time and space.
\newblock \emph{Journal of Logic and Computation} 6(3):315--341, 1996.

\bibitem{RosenblattEtAl2019}
Gregory Rosenblatt, Lisa Zhang, William E. Byrd, and Matthew Might.
\newblock First-order miniKanren representation: great for tooling and search.
\newblock In \emph{miniKanren and Relational Programming Workshop}, 2019.

\bibitem{RozplokhasVyatkinBoulytchev2020}
Dmitry Rozplokhas, Andrey Vyatkin, and Dmitry Boulytchev.
\newblock Certified semantics for relational programming.
\newblock In \emph{Programming Languages and Systems (APLAS 2020)}, LNCS 12470, pp.~167--185. Springer, 2020.

\bibitem{Sangiorgi1996}
Davide Sangiorgi.
\newblock $\pi$-calculus, internal mobility and agent-passing calculi.
\newblock \emph{Theoretical Computer Science} 167(2):235--274, 1996.

\bibitem{Smullyan1968}
Raymond M. Smullyan.
\newblock \emph{First-Order Logic}.
\newblock Springer-Verlag, New York, 1968.

\bibitem{StampoulisChlipala2018}
Antonis Stampoulis and Adam Chlipala.
\newblock Prototyping a functional language using higher-order logic programming: a functional pearl on
  learning the ways of λProlog/Makam.
\newblock \emph{Proceedings of the ACM on Programming Languages} 2(ICFP), Article 93, 2018.
  doi:10.1145/3236788.

\bibitem{Statman1979}
Richard Statman.
\newblock The typed $\lambda$-calculus is not elementary recursive.
\newblock \emph{Theoretical Computer Science} 9:73--81, 1979.

\bibitem{Urban2010}
Christian Urban.
\newblock Nominal unification revisited.
\newblock In \emph{Proceedings of UNIF 2010}, EPTCS 42, pp.~1--11, 2010. doi:10.4204/EPTCS.42.1.

\bibitem{UrbanPittsGabbay2004}
Christian Urban, Andrew M. Pitts, and Murdoch J. Gabbay.
\newblock Nominal unification.
\newblock \emph{Theoretical Computer Science} 323(1--3):473--497, 2004.

\end{thebibliography}
}

\end{document}
